\documentclass[11pt]{article}
\usepackage[utf8]{inputenc}
\usepackage[T1]{fontenc}
\usepackage{amsmath,amssymb}
\usepackage{geometry}
\usepackage{float}
\usepackage{graphicx}
\usepackage[numbers]{natbib}
\usepackage[hidelinks]{hyperref}
\usepackage{parskip}
\usepackage{array}
\usepackage{longtable}
\usepackage[strings]{underscore}
\geometry{a4paper, margin=1in}
\setlength{\parindent}{0pt}
\setlength{\emergencystretch}{2em}
\newcolumntype{L}[1]{>{\raggedright\arraybackslash}p{#1}}
\setcounter{secnumdepth}{0}
\makeatletter
\renewcommand{\@seccntformat}[1]{}
\makeatother

\author{Jacob Chinitz}
\title{Entropic Scalar EFT: From Entanglement Microstructure to Gravity and Cosmic Structure}
\date{\today}

\begin{document}
\maketitle

\begin{abstract}
We propose that empty space is a physical medium with a finite budget of quantum entanglement: the linking structure that allows parts of a quantum system to share state. Matter forms when some of that capacity becomes locked into stable, localized defects of the medium. A particle's mass measures how much entanglement is committed to such a defect. Gravity is the surrounding capacity-strain field: near matter, slightly less entanglement capacity is freely available, and in the weak-field limit the fractional shortfall gives the gravitational potential. We attribute the excess acceleration seen in galaxies to the large-scale continuation of the same capacity response.

The paper specifies a finite cell ensemble, a marked charged-transfer prescription, and a physical matching to the Einstein normalization. Together they determine the cell entropy and ordinary weak-field coefficients. The capacity action is the static scalar sector of the selected Einstein action written in the capacity variable, giving Newton's law and the leading no-slip metric. A separate transverse branch gives the galactic acceleration scale and radial-acceleration relation under its stated microscopic matching conditions.

The electron plays a double role. As the lightest clean charged defect, it fixes the exchange rate between committed entanglement and mass and calibrates the absolute cell scale; Many-Pasts supplies the history-space interpretation of that calibration while preserving Born-rule statistics and no-signaling. Faithful resolution makes local renewal memoryless, a reversible marked vertex fixes the internal response weights, and an additive-generator prescription converts them into the charged transfer energy. For heavier leptons, finite-state recurrence fixes the bare shell factor and a common Gaussian scalar fluctuation realizes the quadratic overlap law; the remaining input identifies this shell with the pole-mass operator.

Simulations of dynamical spacetime test the cell model. The medium's weighting orders the cell states on a stable geometry, a scrambled control tests the ordering, and inserted defects strain the capacity and deform the local geometry. A conserved carrier responds to defects of different strength through one rule without detected screening, and a predicted host-geometry shift follows the cell model while the control follows its own. These tests are small and do not yet measure Newton's constant.

The same cells also carry a quantum spatial geometry. Coherent bonds between neighboring cells leave local probabilities unchanged and carry one bit of entanglement per cut. One angular rule fixes both the cut entropy and the curvature weight, so the horizon area law and the Einstein action share a bare coefficient, which a provisional event-volume calibration places within about three percent of the electron-anchored value. Connected groups of cells reproduce the causal-dynamical-triangulation weight. Directional bond counts reconstruct a spatial metric whose Einstein dynamics give Newton's law, two gravitational-wave polarizations, and Kasner expansion, with an explicit quantum ordering and, for uniform geometries, a positive physical inner product. In a finite model a massive particle crosses changing geometry coherently and keeps its rest mass and clock.

Later sections treat transport, clusters, cosmology, dark energy, strong fields, particles, and complexity at labeled grades. An explicit group-field-theory action selects the finite capacity and record dynamics, and its retained sector reconstructs complex states, reversible evolution, probabilities, measurements, and recorded histories. Other extensions remain conditional or open.
\end{abstract}

\clearpage
\tableofcontents
\newpage

\section*{Part I. Physical Idea and Foundations}
\addcontentsline{toc}{section}{Part I. Physical Idea and Foundations}

\section{1. Introduction: The Physical Claim}

\paragraph{The physical picture.}
Space, in this proposal, is a finite medium of entanglement capacity. The particles we call matter are stable defects that lock away part of that capacity, and the surrounding medium responds to the commitment. Seen at large scales, that response is gravity:
\begin{itemize}
\item \textbf{Matter} is a localized capacity defect of the substrate.
\item \textbf{Mass} is the entanglement that defect commits, read in mass units.
\item \textbf{Gravity} is the extended capacity strain, the fractional capacity deficit, that the medium carries around the defect.
\item \textbf{The gravitational effects usually attributed to dark matter} come from two further regimes of the same medium: the long-range capacity strain on galactic scales and the committed stock of the saturated phase in the early universe.
\item \textbf{General relativity} is the low-energy geometry of this capacity medium.
\end{itemize}

Written as a continuum theory, this becomes a scalar EFT for a vacuum-relative entanglement field $S_{\mathrm{ent}}(x)$ and its deficit $\delta S$ relative to the background capacity. The defect sector is written at continuum scale in ordinary stress-energy variables, but its ontology is unchanged: inertial mass enters through the mass-per-entropy map $\kappa_m$, and the weak-field potential is the fractional deficit $\delta S/S_\infty$.

The gravitational response normally attributed to a dark halo is assigned here to the capacity structure of the vacuum itself. Ordinary weak gravity, the galactic excess, and the homogeneous cosmological mode are different regimes of one medium: general relativity supplies its low-energy geometry, and the capacity variable tracks how localized defects deplete and redistribute the available entanglement. Section~29.9 compares this proposal with recent information-theoretic constructions of gravity, after the paper's own development is complete.

\paragraph{What one finite count fixes.}
The finite-capacity substrate, three postulates, and tetrahedral ensemble define the framework. The selected realization specifies its auxiliary response coordinates, marked routing, additive charged generator, and geometry--capacity matching. Finite enumeration fixes the admissibility entropy and marked weights. The charged transfer and electron anchor then calibrate $L_*$, while source projection and physical matching give the ordinary gravitational normalization.

The most controlled branch is the ordinary static weak-field action and source map
\[
\text{microstructure} \longrightarrow \text{coefficient chain} \longrightarrow \text{static capacity EFT} \longrightarrow \{G,\ \text{baseline metric}\}.
\]
It recovers the Newtonian point-source limit, the leading no-slip metric, and the parametrized post-Newtonian values of general relativity through its Einstein parent. The electron anchor, memoryless dressing, and marked vertex fix the substrate length inside the stated support-to-rate branch. With no continuously adjusted coefficient, the induced Newton scale is $G_*=6.6742890813\times10^{-11}\,\mathrm{m^3\,kg^{-1}\,s^{-2}}$, $-0.073\sigma$ from CODATA, and the same marked vertex gives $m_\mu/m_e=206.768280237$ and $m_\tau/m_e=3477.343310$, at $-0.535\sigma$ and $-0.125\sigma$. In fractional terms these agreements are about $1.6\times10^{-6}$, $1.2\times10^{-8}$, and $6\times10^{-6}$. The shell integers alone, before the marked response, already give $205.71$ and $3454.6$, within $0.51\%$ and $0.65\%$ of the measured ratios. The comparison also runs backward: with the marked weight held at its action value, the measured Newton constant selects the vertex's routing integer, and the unique survivor is the same seven fixed by the tetrahedral alphabet (Section~13.5). Appendix~L separates this action-level closure from the historical fact that the residuals were already known.

\paragraph{One entropy, two cosmic numbers.}
The same sharing entropy reaches galactic and cosmological scales through one coefficient, $\epsilon=g_{\mathrm{share,eff}}/(2\pi)^2\simeq0.188$. The galactic response cell is the tilt mode of the transverse carrier frame, with the baryonic field as its stiffness and the inertia lever $\ell=g_{\mathrm{share,eff}}\,c/H$. Thermalized in the apparent-horizon bath at $k_BT_H=\hbar H/2\pi$, it has the gap $\sqrt{g_{\rm bar}/a_0}$, which sets the galactic acceleration scale $a_0=\epsilon cH_0$: $1.250\times10^{-10}\,\mathrm{m\,s^{-2}}$ at the Hubble constant of the theory's own cosmological fit and $1.231\times10^{-10}\,\mathrm{m\,s^{-2}}$ at the Planck value, against the fitted radial-acceleration scale $1.20\times10^{-10}\,\mathrm{m\,s^{-2}}$. The Bose occupation of this mode gives the radial-acceleration relation, its deep limit $v^4=GMa_0$, and $a_0\propto H(z)$ (Section~15). Read as the allocation cap of the saturated early phase, the same $\epsilon$ fixes the stock-to-baryon ratio at formation, the dark-to-baryon ratio $\Omega_c/\Omega_b$ of standard cosmology, $1/\epsilon=5.32$, or $5.333$ after the free-nucleon and helium conversion (Section~20 and Appendix~O.15). The two numbers obey the parameter-free relation $a_0\,\Omega_c/\Omega_b=cH_0$ for the bare pre-release ratio. Release of that committed stock then gives the present vacuum density $\omega_{\Lambda,0}=\eta_*\,\omega_c^{\rm wb}I_{\eta,\rm rel}$: the closure precision $\eta_*$ fixes the drain per unit release, and the weighted astrophysical history $I_{\eta,\rm rel}$ supplies the remaining normalization (Appendix~O.3). Taking the retention gate from the halo occupation of feedback-resolved simulations gives $\Omega_{\Lambda,0}=0.709\pm0.02$ against the $0.707$ required by the cosmological data, with no input adjusted to the vacuum density (Appendix~O.12). With the reionization redshift set by the fitted optical depth, that frozen history fits the CMB, BAO and supernova data with four cosmological parameters against six for $\Lambda$CDM at $\Delta\chi^2=+1.07$, and $a_0=\epsilon cH_0$ then gives $1.25\times10^{-10}\,\mathrm{m\,s^{-2}}$ with no galactic input. Every term of the present energy budget therefore follows from the measured baryon density: the tie gives the committed density $5.333\,\omega_b$, the release gives $\omega_\Lambda=\eta_*\omega_cI_{\rm rel}$, and flatness gives $h^2$ as their sum. With $\omega_b=0.02238$ the astrophysical chain predicts $H_0=69.3\pm2.1\,\mathrm{km\,s^{-1}\,Mpc^{-1}}$ and $a_0=(1.27\pm0.04)\times10^{-10}\,\mathrm{m\,s^{-2}}$ with no fitted dark-sector parameter (Appendix~O.12). The acceleration scale rests on the lever of Section~15, and the abundance reads the same $\epsilon$ as the allocation cap, a separate premise of Section~20. Among the eleven numbers native to the construction that lie in the radial-acceleration window, only $(2\pi)^2$ also reproduces the abundance (Appendix~L).

\paragraph{Galaxies.}
A specified carrier-resolved transverse branch produces the galactic acceleration scale $a_0$, the radial-acceleration relation, its conservative nonspherical extension, and the leading lensing metric without per-system tuning, subject to the microscopic matching conditions stated in Sections~15--16 and Appendix~N. Evaluated on the joint coarse source of all active carriers, the branch conserves momentum between a galaxy and its satellites, and the response budget set by each carrier's Lagrangian baryons fits the isolation-safe weak-lensing points. With the stock outside each carrier's support left collisionless, the isolated-lens profiles are statistically indistinguishable from those of $\Lambda$CDM halos (Section~15). For 131 SPARC galaxies, its director completion agrees with the algebraic radial-acceleration relation to $0.012$ dex in every radial bin; both have an observational rms residual of $0.147$ dex. It fits the Milky-Way vertical potential as well as NFW controls. Old-disk vertical dispersions separate it from QUMOND: one stellar mass-to-light ratio fits both the rotation curve and the old-disk dispersions of NGC~6946 at $1.0\sigma$ and of NGC~628 at $0.3\sigma$, against $2.6\sigma$ and $2.9\sigma$ for QUMOND (Section~16). The director completion was formulated after the local-contact Milky-Way failure was known, so these comparisons are development-set results; the registered DiskMass old-disk and stream-shape predictions of Section~27.1 are its frozen tests. The long-lived tilt of the Gaia--Sausage--Enceladus debris counts against its aligned Milky-Way solution.

\paragraph{Quantum geometry.}
The capacity/record action splits a primitive shared face into a
seven-state matched channel and a nine-state retained record.
Its coherent spatial encoding preserves the local ensemble.
One angular tensor supplies the bare curvature coefficient;
connected primitive patches realize its covariant history weight.
Directional reconnections give a canonical metric operator with
a homogeneous physical-state construction. The retained cochain
operator joins orbital motion and a positive clock to an enumerated
geometry transfer. These are specified realizations within the
ontology. Their interacting continuum phase and inhomogeneous
physical constraint space remain open. The connected patches reproduce the causal-dynamical-triangulation weight exactly at the integer-channel asymmetry $\alpha=7$, so the lattice host of Part~VII becomes a candidate native measure. Its registered volume-stiffness test agrees with the boundary normalization at 32k, with supporting stationary profiles at 16k and 64k. The phase test of start independence remains undetermined, and the matter response remains open (Section~24.10 and Appendix~J.14).

\paragraph{Cosmology.}
The cosmological release model uses further infrared prescriptions: a post-refresh energy ledger, an ensemble release history evaluated on the capacity clock carried by source-associated flow, and the initial condition $\rho_{\Lambda,i}=0$. Its finite volume projection and renewal identities are exact. A constrained covariant clock action reproduces the background exchange and fixes its lapse, clock, and velocity perturbations for any specified release history. For a closed progenitor cohort under common advection--diffusion, fully resolved position-information loss is bounded and monotone; the physical coarse observable requires its own check. The direct geodesic-clock perturbations admit an exact conserved-fluid representation. Direct geodesic perturbations with matched lensing give $\Delta\chi^2=+0.85$ against $\Lambda$CDM for Planck 2018, DESI DR2 and DES-Dovekie with $\sum m_\nu=0.06\,$eV, with five versus six cosmological parameters ($\Delta{\rm AIC}=-1.15$); the fits are statistically tied. A selected low-density completion preserves this cold linear branch and approaches a positive vacuum-like resting state before local exhaustion. A first self-gravitating progenitor-tagged run supplies a release kernel, whose gate and label sensitivity leaves an independently fixed normalization unresolved.

\paragraph{Scope.}
Later parts treat time-dependent transport, galaxy clusters, cosmology, the saturated early phase, strong fields, and particle and gauge extensions. Part~VIII records their derivational status. The companion Lean project checks mathematical claims across the main text and appendices, including analytic integrals and certified numerical bounds. Its claim ledger distinguishes proved consequences, physical premises, numerical checks, and open constructions (R.1).

Many-Pasts, the third postulate, already does work in this chain: faithful sector resolution selects the memoryless electron-dressing kernel, and Many-Pasts supplies the history space in which that kernel operates. Its consequences for quantum probability, branch realization, and the arrow of time are developed in Section~22 and Appendix~G.

\paragraph{Relation to the companion accounts.}
This manuscript, \emph{A microscopic realization of causal triangulations
with native matter and clock dynamics} (\texttt{MICROSTRUCTURE\_CDT.tex}),
and \emph{Empirical Results: History-dependent gravity and cosmological
release} describe the same selected framework. The CDT account develops
the connected geometry and finite matter--clock construction incorporated
here. The empirical account supplies the current observational and
simulation results. Its effective fits do not measure the CDT ensemble's
interacting Newton coefficient or establish its continuum phase.

\subsection{1.1 Primitive Inputs and Closure Status}

The word ``closure'' is used here in a specific sense. The paper does
not derive the existence of a finite entanglement substrate or the tetrahedral
boundary architecture from a deeper microscopic Hamiltonian; those are
theory-defining inputs, and the closure claim begins only after they are
fixed. There are five such inputs: finite local entanglement capacity;
geometry--capacity equivalence; mass--entropy equivalence, including matter as
localized defects of committed capacity; the Many-Pasts ontology of
Postulate~III, with ordinary probability normalization and finite additivity
for mutually exclusive complete physical records;
and the tetrahedral ultraviolet architecture, including positive oriented
matching of the two primitive descriptions of a shared face. Appendix~B
writes that matching rule as a primitive pair operator. Within the selected
sharp reversible vertex, unitarity proves that the retained carrier is its
orthogonal fusion complement; the hard-core realization, flags, relative
phases, and routing remain selected microscopic details. This refines the
existing ultraviolet input rather than adding a fourth postulate. Maximum
caliber is likewise not a sixth input: it restates, for histories, the same
faithful full-support condition already used to select the admissibility
ensemble: every pass carries the largest path entropy compatible with the
same fixed marginal. The equivalence has to be stated explicitly because
finite capacity by itself does not imply renewal.

Finite counting fixes the admissibility weighting and effective entropy.
Faithful full-support resolution fixes the memoryless local replacement
kernel. The state-weighted determinant supplies the seven-channel
recurrence, while the marked vertex and additive-generator prescription
fix its charged correction. The electron anchor fixes the proper-time
cadence, and \(L_*=c\tau_*\) defines its causal length. The continuum
Einstein reduction gives the Newtonian limit, leading no-slip metric,
and PPN values on the matched branch of C.5.

The coherent completion supplies a spatial state with the same local
capacity marginals. Its angular tensor fixes a common bare curvature
and cut coefficient. Connected primitive patches, normalized by their
full Gram matrices, realize a selected causal triangulation measure.
The canonical area-count construction supplies an Einstein/ADM
continuation with explicit ordering and constraint tests; a positive
rigging form exists for its homogeneous volume--shape constraint.
A finite geometry--matter operator carries the native rest gap,
orbital motion, and mass-derivative clock, with exact preservation
under the specified cochain refinements. The general constraint
space, relativistic phase, and common interacting normalization remain
open. The galactic contact still needs its microscopic carrier
projector and thermal Hamiltonian. In cosmology, the constrained
capacity-clock action fixes exchange for a supplied release history;
the physical coarse mixing observable has a preliminary progenitor-tagged
measurement and a retention gate from feedback-resolved simulations,
with its label resolution and convergence still to be fixed independently.

The five commitments define the physical framework. Its quantitative
realization specifies marked routing, auxiliary response coordinates,
and the charged generator. C.5 uses a normalized geometry--capacity
matching to set \(G=G_*\). The coherent angular model instead gives
\(\hbar G_0/c^3=A_{\rm ch}/(4\ln2)\) before imposing that equality.
The event law, mean event assignment, and regular coarse area--volume
calibration in C.8 give a provisional \(G_0/G_*=1.03002\).
These are distinct normalization prescriptions. Their comparison with a
renormalized bulk measurement must use one common state and action.

In compressed form, the central claim is
\[
\begin{aligned}
&\boxed{\text{primitive UV capacity hypothesis}}
\rightarrow
\boxed{\text{finite counting + admissibility}}
\rightarrow
\boxed{L_*}
\\
&\rightarrow
\boxed{\gamma,\ \kappa/\gamma}
\rightarrow
\boxed{\delta S \leftrightarrow \Phi}
\rightarrow
\boxed{G,\ \text{baseline metric}},\\
&\hspace{3em}
\boxed{a_0,\ \mathrm{RAR},\ \text{galactic lensing}}\ \text{in the carrier-resolved transverse EFT}.
\end{aligned}
\]
The microstructure is an input: a finite ultraviolet counting problem from which the weak-field sector is derived.

The absolute scale calibration uses the electron, the lightest elementary charged defect, as the dimensional anchor. Faithful full-support resolution forces the replacement kernel to be memoryless: a dressing pass can carry the full admissibility entropy only if it retains no memory of the endpoint it replaces (the one-line proof appears in Section~1.2). The lightest-defect functional selects the fermionic ceiling $k=7$ with $\Delta_7=0$, and the state-weighted determinant of the seven renewed clouds is $r=e^{-7g_{\mathrm{share,eff}}}$. The decorated vertex fixes the residual marked-fiber factor and its electron routing $Z_e$. The additive-generator prescription assigns these weights to the charged energy. Positivity gives the raw survival energy $E_{\rm raw}=-(\hbar/\tau_*)\ln(1-r)$, and the dressed electron identification $m_ec^2=(3/2)Z_eE_{\rm raw}$ fixes $\tau_*$. The associated causal length is defined by $L_*=c\tau_*$. This conversion introduces no geometric diameter, but a microscopic matter action must still reproduce the same invariant speed in its spatial dispersion. Appendix~H gives the finite action, the coherence tests, and a joint orbital construction whose general continuum cone remains to be established. Appendix~I separates the arithmetic recurrence statistic that supplies the shell support from the geometric typical recurrence measured by entropy and gives an explicit positive realization of the coherent $N^2$ overlap.

The reduced capacity functional can superficially resemble the scalar sector of a Brans--Dicke theory \citep{BransDicke1961}, but the resemblance is misleading. In the ordinary static branch it is the Einstein constraint action rewritten through $\delta S=-2S_\infty\Phi/c^2$, with no second scalar--tensor action. The open action questions are narrower: the microscopic origin and nonlinear covariant variation of the carrier-resolved transverse contact, the generic covariant capacity observable outside controlled reductions, and the saturation-boundary functional.

The remaining microscopic tasks are to select the volume state and
geometric phase, establish a common matter cone across changing
boundaries and species, and construct the inhomogeneous physical
constraint space. The finite cochain refinement in H.15 preserves a
specified joint operator. The spatial and count limits of the
gravitational constraint algebra require control of metricity boundaries and common domains.
The face-wise native profile/cut mismatch of J.11 applies within
its own ensemble. The connected integer-channel host of
J.12 requires its own interacting bulk and source measurements.
The carrier-resolved transverse contact, finite-loop stiffness
operator, and strong-field boundary action retain their separate
microscopic requirements.

The construction rests on five commitments.

First, the vacuum is a medium with a bounded local capacity for entanglement.

Second, spacetime geometry and that capacity structure are the same substrate seen at different scales, so that in the weak field gravity is the fractional deficit of locally available capacity.

Third, matter is localized committed capacity: a particle is a stable defect of the medium, and its inertial mass is the entanglement content of that defect read in mass units.

Fourth, a recorded present is supported by many compatible microscopic pasts. The operational branch assigns probabilities only to decoherent record histories through the history Gram kernel derived in Appendix~H.11, then conditions them on the realized present. Its finite retained sector gives Born statistics and no-signaling. The reversible renewal dilation supplies a concrete microscopic role for the history degrees of freedom: they receive the previous local state while the present register is renewed.

Fifth, the ultraviolet cell has tetrahedral boundary architecture. Each shared
face begins with two fermionic primitive slots and positive oriented matching.
The matching operator transmits the maximal coupled multiplet. In the selected
sharp reversible vertex, unitarity identifies the complete nine-state marked
carrier with the discarded complement. The three-dimensional closure vector
then fixes $j_0=\tfrac32$ and the seven-state alphabet internally. For a
different microscopic vertex, $j_0=\tfrac32$ remains the surviving branch of
the discrete audit in Appendix~B. Four relational ports, single-copy channel
capacity, and the two orientations then give the 1680-state ensemble. The same
faithful-resolution standard is applied to its histories: the local process
carries the full available path entropy and retains no endpoint memory.
``Maximum caliber'' names this temporal application of the same standard; it
adds no independent commitment to the five listed here.

The finite-capacity substrate is the ultraviolet premise; geometry--capacity equivalence, mass--entropy equivalence, and Many-Pasts are the three postulates; and the tetrahedral ensemble is the ultraviolet architecture. Faithful full-support resolution acts on both states and histories. From these inputs the paper derives the admissibility weighting, $g_{\mathrm{share,eff}}$, the replacement kernel, the one-layer native-vertex update, the marked transfer, the tree edge factor, the weak-field bridge, and the Newtonian metric branch, all inside the displayed decorated action. The separate loop-dressed stiffness, the galactic acceleration scale, and the radial-acceleration relation carry the conditional grades recorded in the closure table.

The architecture therefore contains three postulates, one finite-capacity substrate premise, one minimal ultraviolet ensemble, and a small number of explicitly labeled conditional readings. Applied to the same ensemble, faithful resolution fixes both the state entropy and the maximum-caliber history process.

\subsection{1.2 Physical Motivation for the Primitives}

These commitments are premises, and the derivation begins only after they are fixed. Each addresses a specific open problem in black-hole thermodynamics, quantum information, the equality of inertial and gravitational mass, or the quantum-mechanical role of records. Their motivations do not prove them; their numerical and structural consequences provide the tests.

Three established results motivate a finite-capacity substrate. A black hole's entropy scales with horizon area; the Bekenstein bound limits the information a bounded region can hold; and entanglement-based reconstructions relate spatial geometry to entanglement structure. The present model treats the vacuum as a medium with a finite local entanglement budget and geometry as its large-scale description. In this picture, a black hole exhausts the interior capacity, leaving the active bookkeeping at the boundary between exhausted and available capacity. A finite cell state space also supplies an ultraviolet cutoff because the continuum description ends below the cell scale. The assumption that the empty state maximizes capacity agrees with the thermodynamic direction suggested by gravitational entropy. Three recent continuum results provide compatible limits: the observer-dressed de~Sitter algebra makes empty de~Sitter the maximum-entropy gravitational state \citep{ChandrasekaranEtAl2023}, a horizon modular calculation recovers Einstein curvature from the relative information carried by an excitation \citep{DorauMuch2026}, and a geometric-relative-entropy action yields local bulk information dynamics with an Einstein limit \citep{Bianconi2025,Bianconi2026}. None supplies the finite cell or its coefficients; Section~29.9 states their precise role in the present construction.

The mass--entropy identification addresses a coincidence that general relativity encodes but does not explain. General relativity builds in the equality of inertial and gravitational mass geometrically, through the equivalence principle, but gives no microphysical account of why the mass that resists acceleration and the mass that sources attraction should be one and the same. Here both are readings of a defect's committed entanglement, so the equality follows from the construction instead of being imposed by hand. The same identification bears on why gravity is so weak. The induced gravitational scale contains the exact factor $Z_e^2\ln^2(1-e^{-7g_{\mathrm{share,eff}}})$, whose dominant hierarchy is $e^{-14g_{\mathrm{share,eff}}}$: a large sharing entropy makes gravity exponentially feeble, and the gap between gravity and the other forces becomes a matter of arithmetic rather than fine-tuning. The galactic extension carries a lower closure grade. In its carrier-resolved transverse branch, $a_0=\epsilon cH_0$ with $\epsilon\equiv g_{\mathrm{share,eff}}/(4\pi^2)$ ties the onset of anomalous rotation to the cosmic horizon. The thermal contact fixes the response of each retained carrier, and the auxiliary action relates that response to the resolved baryonic source. The tilt-mode lever, carrier projector, and microscopic metric vertex are the tests of that connection.

The scale-setting proposal separates two questions that sound alike: how hard the medium works to keep a defect bound, and how often the defect fully re-forms. Because the medium never holds still (every instant it is re-drawn from all the ways it could be), a particle is a pattern the medium must continually re-form, and the entanglement it commits is a maintained quantity, re-established at every update. Keeping the knot bound means holding a grip on every strand at once, and separate holds add up, so the maintenance is large. Full re-completion, with every strand falling into alignment at the same instant, is a simultaneous coincidence, and coincidences multiply, so it is exponentially rare. The particle's mass follows the second question, not the first: it is fixed by how rarely the binding fully re-closes. This rare recurrence allows the electron to lie far below the natural substrate scale without an inserted small number: deep binding makes full coherent recurrence exponentially unlikely. Quantitatively, seven channel entropies add, so their effective support multiplies; the state-weighted determinant turns that support into a record-conditioned transfer rate; the positive survival operator gives the rest-energy gap; and the decorated marked vertex fixes the finite closure-response correction while the electron anchor fixes the clock. Section~23 and Appendix~H give this its full form.

A memoryless update rule posits no hidden machinery carrying information forward from tick to tick, and in this construction it is also forced. For any stationary per-channel kernel with marginal $p_{\eta_*}$,
\[
H(B_{t+1}\mid B_t)=g_{\mathrm{share,eff}}-I(B_t;B_{t+1})\le g_{\mathrm{share,eff}},
\]
so a dressing pass carries the full admissibility entropy if and only if the endpoint mutual information vanishes, which fixes the refresh kernel $K(b,b')=p_{\eta_*}(b')$. The allowed local single-label dynamics cannot realize this requirement: they freeze into disconnected sectors that never explore the full space (Appendix~D.4), so the refresh must act nonlocally on the native cell. The decorated vertex prepares the diagonal fresh amplitude and realizes the charged recurrence; selection of an interacting geometric condensate and the physical source of its history capacity remain open. Compatibility with Lorentz tests requires the retained matter sector and the geometric modes to share the same invariant infrared speed and requires substrate loops to suppress dimension-four Lorentz-violating operators. Renewal fixes neither result. Section~4 and Appendix~H.15 state the finite construction and remaining continuum tests.

Many-Pasts holds that the present configuration of the entanglement network is supported not by one definite microscopic past but by a conditional ensemble of compatible decoherent histories. It provides an ontological account of interference, entanglement correlations, and measurement while preserving ordinary quantum dynamics. The interference in a double-slit experiment is the persistence of the unrecorded alternative histories in that ensemble; a durable which-path record conditions the ensemble, and the interference goes away. The correlations of an entangled pair come from weighting the histories of the whole joint system, which reproduces the nonclassical statistics with no signal passing between the two wings. Measurement adds no separate collapse law; it lays down a durable record, after which the relevant histories are the ones compatible with it. The operational construction is therefore ordinary quantum mechanics equipped with a history-space ontology, and it leaves laboratory predictions intact: Born-rule statistics and no-signaling both hold. Its arrow of time follows on finite connected networks from the reduced channel and record root of Appendices~P.8--P.9; the gravitational entropy of the root remains open. The postulate also constrains the construction: faithful sector resolution independently selects the memoryless dressing, and Many-Pasts supplies the history-space setting in which that process operates and in which its exported registers live.

Tetrahedra are standard building blocks in several approaches to quantum geometry \citep{BaezBarrett1999,Oriti2009,AmbjornJurkiewiczLoll2004}. The specific addition here is a primitive matching rule. After the two descriptions of a shared face are transported into one orientation, its positive mismatch operator has the maximal-spin sector as its unique null space. For the physical branch, the independently constructed nine-state marked register matches the complete nonmaximal fusion sector only at $j_0=\tfrac32$. Conditional on identifying those two objects dynamically, three-dimensional closure fixes seven face labels; four relational ports, single-copy channel capacity, and two orientations then give exactly 1680 boundary states. Without that identification the same result is the unique surviving point of a stated discrete branch audit, not a consequence of minimality. The closure structure also permits only three nondegenerate charged-lepton shells, a candidate answer to why the Standard Model contains three charged-lepton generations. Exhaustive enumeration confirms the cutoff without a preferred alphabet: every one- and two-label exclusion keeps a nonzero closure variance, and all 35 three-label exclusions give zero (Appendix~K.3). Appendix~L preserves the historical provenance of the construction. Once the branch is fixed, its entropy contains no adjustable continuous parameter.

Motivation alone carries little evidential weight, because a proposal of this scope can almost always assemble a list of mysteries explained after the fact. The relevant questions are how much freedom remained when the construction was chosen and whether one fixed construction survives measurements it did not anticipate. Appendix~L audits the first question. The following sections derive the weak-field consequences and compare them with observation.


\section{2. Canonical Field Content and Definitions}

Before the symbols, four plain words recur throughout. \emph{Capacity} is the entanglement support locally available in the medium. A \emph{defect} is a stable, localized commitment of that capacity, which we coarse-grain into a particle. A \emph{deficit} is capacity no longer freely available to the surrounding vacuum because a defect has committed it. \emph{Strain} is the extended profile of that deficit reaching out into the medium, whose fractional size the weak-field potential tracks. The field variables below are the precise versions of these words.

We define the fundamental continuum variable as the vacuum-relative coarse-grained entanglement assigned to a UV probe cell of size $L_*$ centered at $x$:
\[
S_{\mathrm{ent}}(x)\in \mathbb{R},
\]
measured in nats and therefore dimensionless. This is not a literal microscopic entropy density at a mathematical point. It is the leading scalar order parameter associated with a vacuum-relative entanglement defect after coarse-graining over a UV cell.

This definition keeps the microscopic and continuum pictures tied together. At continuum level, $S_{\mathrm{ent}}(x)$ is the field that appears in the action and field equations. At the microscopic level it is the coarse variable recording how much local entanglement capacity remains available in the underlying medium after averaging over a UV cell.

The asymptotic vacuum-capacity baseline is denoted $S_\infty$, and the deficit field is
\[
\delta S(x)\equiv S_\infty-S_{\mathrm{ent}}(x).
\]
Positive $\delta S$ denotes reduced available vacuum entanglement capacity in the neighborhood of a localized defect or defect distribution. It is the extended capacity-strain field sourced by the defect sector, not an independent medium acted on by matter from outside. For nonlinear work it is useful to define the bounded occupancy fraction
\[
q(x)\equiv \frac{S_{\mathrm{ent}}(x)}{S_\infty}=1-\frac{\delta S}{S_\infty}\in[0,1].
\]
The variables $S_{\mathrm{ent}}$, $\delta S$, and $q$ therefore describe the same local physics in three closely related ways: available capacity, missing capacity relative to vacuum, and surviving-capacity fraction. Each is used where it is most transparent: $\delta S$ for the weak-field theory, because it maps directly onto the Newtonian potential; $q$ for the nonlinear and strong-field completion, because boundedness is built in from the start; and $S_{\mathrm{ent}}$ itself for the covariant EFT, because it is the field that appears in the action.
The operational meanings are:
\begin{itemize}
\item $q=1$: vacuum capacity fully available in the absence of local defect-induced capacity strain;
\item $0<q<1$: partial local capacity reduction around a defect configuration;
\item $q=0$: complete local exhaustion of available capacity on the physical branch.
\end{itemize}

\paragraph{Fixed-epoch normalization.}
The entropy unit is conventional at a fixed epoch. The physical ratio $\kappa/(\gamma S_\infty)$ is fixed separately by the geometry--capacity matching in Appendix~C.5. Under a constant rescaling
\[
S_{\mathrm{ent}}\mapsto K S_{\mathrm{ent}},
\qquad
S_\infty\mapsto K S_\infty,
\qquad
\delta S\mapsto K\delta S,
\]
the observable bridge
\[
\frac{\Phi}{c^2}=-\frac{\delta S}{2S_\infty}
\]
is unchanged. The source equation is invariant in the same sense: rescaling the entropy field rescales the source coefficient with it, so the observable Newtonian normalization depends on the gauge-invariant combination $\kappa/(\gamma S_\infty)$ rather than on $S_\infty$ alone. A cell-normalized description and a horizon-normalized description can therefore assign different numerical values to $S_\infty$ without changing $\Phi$, $G$, or the PPN limit. This is not a time-dependent gauge symmetry; it is a fixed-epoch entropy-unit convention. Gravity sees fractional capacity depletion.

\paragraph{Substrate length scale.}
The canonical UV cell length is not taken to be the conventional Planck length as an input. Faithful full-support resolution fixes the renewal kernel and its history-space factorization. The state-weighted determinant gives the seven-channel recurrence
\[
r=e^{-7g_{\mathrm{share,eff}}},
\qquad
L_*^{(0)}=-\frac32\lambda_e\ln(1-r),
\qquad \lambda_e=\frac{\hbar}{m_ec}.
\]
The decorated marked-transfer vertex derived in Appendix~H fixes
\[
\zeta_*=9e^{-g_{\mathrm{share,eff}}}
\left(1-\frac{8\eta_*}{49}\right)^{21/2},
\qquad
Z_e=(1+\zeta_*)(1+7\zeta_*^2).
\]
The physical electron-anchored scale is
\[
\boxed{L_*=Z_eL_*^{(0)}}
=1.6162537024\times10^{-35}\ {\rm m}.
\]
The corresponding induced gravitational scale is
\[
G_*:=\frac{c^3L_*^2}{\hbar}
=
\frac94\,\frac{\hbar c}{m_e^2}\,
Z_e^2\ln^2\!\left(1-e^{-7g_{\mathrm{share,eff}}}\right)
=
6.6742890813\times10^{-11}\ {\rm m^3\,kg^{-1}\,s^{-2}}.
\]
The lightest one-bit fermionic defect resolves the seven face sectors once and exports the transverse $2/3$ share of that dressing block. The marked vertex accounts for its finite closure-response fiber and label-return loop. The conventional Planck length $L_P=\sqrt{\hbar G/c^3}$ remains useful for comparison and for standard black-hole thermodynamic notation, but it is not the primitive scale-setting input here.

The principal coefficients and derived quantities used throughout are:
\begin{align}
\gamma &:\ \text{entanglement-field stiffness},\\
\kappa &:\ \text{defect--entropy coupling},\\
\kappa_m(\ell) &:\ \text{mass-per-entropy map at scale }\ell,\\
L_* &:\ \text{causal tick length in the electron-anchored support-to-rate map},\\
G_* &:\ \text{gravitational scale induced by }L_*,\\
g_{\mathrm{share,max}} &= \ln(1680),\\
g_{\mathrm{share,eff}} &:\ \text{admissibility-weighted effective sharing entropy},\\
J_{\mathrm{bare}},\,J_{\mathrm{eff}}^{\mathrm{tree}},\,J_{\mathrm{eff}}^{(\mathrm{ren})} &:\ \text{UV edge-kernel couplings},\\
a_0 &= \frac{cH_0g_{\mathrm{share,eff}}}{4\pi^2}
\quad\text{in the conditional compact two-phase normalization}.
\end{align}
The gravitational potentials are denoted $\Phi$ and $\Psi$, and the canonical weak-field bridge will be written as
\[
\frac{\Phi}{c^2}=-\frac{\delta S}{2S_\infty}.
\]
These same symbols reappear in the ultraviolet closure chain, the continuum action, and the phenomenology sections. From this point onward each one keeps the same meaning, so the later derivations build on a single notation rather than shifting between parallel conventions.


\section{3. The Three Postulates}

The three postulates answer distinct physical questions. Postulate~I identifies spacetime geometry with the long-wavelength expression of the capacity substrate. Postulate~II identifies matter with localized committed capacity and mass with its inertial reading. Postulate~III describes the realized present through a history-space ontology over the compatible decoherent pasts recorded in it. Appendix~H.11 derives the finite retained-sector probability and history machinery; its global continuum/Fock and cosmological realization remains open. The arrow of time follows on finite connected networks from the reduced channel and record root of Appendices~P.8--P.9. Faithful sector resolution, applied to this history space, selects the memoryless dressing kernel used by the decorated scale-setting action.

\subsection{3.1 Information--Geometry Equivalence}

The first postulate states that spacetime geometry is the continuum expression of the capacity substrate. This is stronger than saying entanglement contributes an additional piece of stress-energy inside otherwise standard general relativity: the metric and the scalar capacity sector are two projections of one finite medium, and $S_{\mathrm{ent}}$ is not appended to an independent background geometry. In the weak field, gravitational potential is the fractional deficit of available capacity.

Two consequences of this reading should be kept distinct from the start. First, absolute $S_{\mathrm{ent}}$ is not itself ``the gravitational potential''; the observable weak-field potential comes from the fractional deficit $\delta S/S_\infty$, which is why a fixed-epoch rescaling of entropy units leaves gravity unchanged (Section~2). Second, because geometry and capacity are two descriptions of one response, the deficit is not an extra force appended to an independently existing metric. Section~10 proves that the ordinary reduced capacity functional is the Einstein constraint action in the capacity coordinate. The remaining common-parent problem is the microscopic embedding of the carrier-resolved transverse contact and the saturation boundary.

\subsection{3.2 Mass--Entropy Equivalence}

The second postulate identifies mass as the inertial reading of localized capacity commitment. At scale $\ell$,
\[
m(\ell)=\kappa_m(\ell)\,\Delta S.
\]
A particle is already a localized defect of the entanglement substrate, so $m=\kappa_m\Delta S$ asserts that the inertial content of the defect \emph{is} its entanglement content, read in mass units.

For elementary fermionic sectors the canonical defect increment is
\[
\Delta S_f=\ln 2.
\]
The one bit here is not arbitrary. An elementary fermionic exclusion is binary (the face is occupied or unoccupied), and a binary distinction carries exactly $\ln 2$ of missing entanglement. This is the simplest possible defect increment, which is why the lightest such defect, the electron, becomes the cleanest anchor for the mass--entropy map (Section~13). Composite sectors instead require their fully dressed bound-state entanglement budgets.

Two corollaries are used later. First, because mass and entanglement budget are two descriptions of the same defect, and the masses of separated defects add, capacity committed in service of one defect cannot simultaneously serve another: shared service would make the joint budget, and with it the joint mass, sub-additive. Commitment is therefore per-defect: each committed unit carries the label of the defect it serves. Second, the same bookkeeping makes the saturated early phase countable: its abundance follows from the number of committed units, with no double-counting across separated defects (Section~20).

\subsection{3.3 Many-Pasts Hypothesis}

The third postulate concerns the microscopic support of the present entanglement network. A recorded present can be compatible with many coarse-grained histories of the substrate. Many-Pasts takes those alternative pasts seriously while retaining one realized macroscopic present. Its probability theory must distinguish amplitudes for alternatives that still interfere, probabilities for recorded presents, and conditional probabilities for decoherent histories compatible with a given record.

The operational construction takes the decoherent-histories form \citep{Halliwell1994}. Appendix~H.11 derives that form for every finite retained sector from its real dynamics, record composition, and norm. A coarse history $h=(\alpha_1,\ldots,\alpha_n)$ has class operator
\[
C_h=\Pi^{(n)}_{\alpha_n}U_{n,n-1}\cdots
\Pi^{(1)}_{\alpha_1}U_{1,0},
\]
and decoherence functional
\[
\mathcal D(h,h')=\operatorname{Tr}\!\left(C_h\rho_0C_{h'}^\dagger\right).
\]
When a family decoheres, $\mathcal D(h,h')\simeq0$ for $h\ne h'$, its diagonal entries obey the ordinary probability sum rules. If $P$ denotes a final macroscopic record and $\mathcal H_P$ is a decoherent refinement of the histories ending in that record, then
\[
p(P)=\sum_{h\in\mathcal H_P}\mathcal D(h,h)
=\operatorname{Tr}(\Pi_P\rho_{\rm now}),
\qquad
p(h\mid P)=\frac{\mathcal D(h,h)}{p(P)}.
\]
The probability of the present is obtained from the common normalized measure over all records before the history distribution is conditioned on $P$. Unresolved alternatives remain combined at amplitude level; no positive probability is assigned to individual fine-grained paths that have not decohered.

This construction makes the finite operational branch standard quantum
mechanics.  Reversible coupling to a retained record gives the Kraus
operators, POVMs, conditional update, and trace-preserving reduced channel.
The norm rule gives their Born probabilities, and locality plus trace
preservation gives no-signaling.  Many-Pasts changes the history-space
interpretation without adding a collapse term or a signaling bias.

The ontological reading is record-retentive. The realized present includes both its current macroscopic configuration and the physical records that encode earlier events. Compatible pasts are distinguished to the resolution carried by those present records; they are not additional coexisting spacetimes. When a physical process erases a record distinction, the conditional history measure coarsens by summing the histories that the surviving record can no longer separate. Conditioning on a future record is defined only after that record belongs to a present configuration, so the construction supplies no future-to-past force or retrocausal update law.

For the retained one-mark sector, Appendix~H.11 derives rather than assumes
$\mathcal H_{\rm mark}\simeq\mathbb C^9$ and the ray space
$\mathbb{CP}^8$.  Reversible internal equivalence makes a branch weight a
function only of its squared norm.  Additivity under coarse-graining of
orthogonal durable records makes that function additive; positivity and
normalization then fix
\[
\boxed{p(\Pi\mid\psi)=\|\Pi\psi\|^2
=\langle\psi|\Pi|\psi\rangle.}
\]
This proof uses record completeness, nonnegativity, normalization, and finite
additivity for mutually exclusive physical records.  Gleason's theorem
\citep{Gleason1957} remains an independent consistency check rather than the
step that carries the Born derivation.  The same norm makes the finite
decoherence functional the Gram kernel of history branch vectors.  The full
continuum/Fock-space kinematics, the cosmological initial state, the histories
that actually decohere, and durable global record capacity remain open.

The postulate and the renewal theorem have separate logical roles. Many-Pasts supplies the ontology and the record-conditioned history space. The already-stated faithful full-support condition selects the memoryless dressing kernel when applied to paths at fixed admissibility marginal; ``maximum caliber'' names this history-space form and is independent of the Born measure on laboratory records. The lightest-defect functional selects the same absence of temporal memory and, within the support-to-length map, selects seven occupied channels with vanishing inter-channel correlation. A reversible dilation exports the old replaceable closure register to history. The determinant-survival action and additive-generator prescription combine the finite marked response with the Compton phase readout.

The renewal statement applies to the replaceable closure register, not to every degree of freedom carried by a charged defect. The marked internal fiber and its position degree of freedom belong to the retained system. Appendix~H.11 proves the corresponding coherence criterion: after a dilation, a spatial off-diagonal is multiplied by the overlap of the discarded records produced by its two branches. Free marked transport must therefore export no cell address or other which-path label. This is a consistency condition on the microscopic realization of the existing operational quantum postulate, not an additional founding premise.


\section{4. Relativistic Continuum Structure}

\subsection{4.1 Capacity budget and continuum symmetry}

The finite-capacity, isotropic, relational substrate admits a covariant
continuum description if its matter and geometric operators approach
one invariant cone. Relational state counting supplies no preferred
embedding coordinates and leaves that dynamical limit to the operators.

The first ingredient is a finite maximal update rate, denoted by the same constant $c$ that later appears in the transport relation $D/\tau_0=c^2$. In the present interpretation, $c$ measures the largest rate at which the substrate can propagate and reorganize information. A defect at rest spends that budget entirely on local temporal evolution. A defect in motion must spend part of the same budget on spatial transport within the surrounding network. Because the substrate is isotropic, the cost of motion depends only on the rotational scalar $v^2$ at leading order, with the temporal rate maximal at $v=0$ and vanishing when the budget is exhausted at $v=c$. These endpoint conditions alone admit many interpolating functions and so do not fix the form of the time-dilation relation. The form is fixed once the finite update speed is treated as invariant across inertial coarse descriptions: homogeneity, isotropy, and the relativity principle then select the Lorentz group rather than the Galilean one, giving the invariant interval
\[
c^2 d\tau^2 = c^2 dt^2 - d\mathbf{x}^2,
\]
and hence
\[
\frac{d\tau}{dt}=\sqrt{1-\frac{v^2}{c^2}}.
\]
The capacity-budget picture supplies the substrate interpretation of this Lorentzian kinematics: motion allocates part of the finite update budget to spatial transport, leaving the remaining fraction as proper-time evolution.

This argument fixes the continuum kinematics once a common invariant
speed exists. The charged rest gap determines \(\tau_*\); the spatial
operator determines its small-momentum coefficient. H.15 constructs
both within one finite geometry--matter model. Its raw Whitney metric
has direction-dependent limiting speeds, while the selected inverse
metric corrects the leading cone on a periodic family. A common cone
across the full geometric ensemble remains a test. The definition
\(L_*=c\tau_*\) is a causal-length conversion and does not set a bond
length, tetrahedron diameter, or face area.

The same capacity language also unifies motion-induced and gravity-induced clock slowing. In the nonlinear branch the surviving-capacity fraction is
\[
q=\frac{S_{\mathrm{ent}}}{S_\infty},
\]
so smaller $q$ means that less local update capacity remains available. Motion reduces the temporal share of the budget by consuming part of it in spatial transport; a nearby defect reduces the local budget by depleting available capacity. The two familiar time-dilation effects are therefore interpreted as two regimes of one mechanism.

The second ingredient is the relational character of the substrate. The network is not embedded in a prior physical manifold whose coordinate labels carry independent meaning; its physical content is the pattern of local capacities, defects, and neighborhood relations. Continuum coordinates are descriptive labels imposed on that relational structure, and smooth coordinate changes relabel the same underlying configuration. This is precisely why the low-energy theory must be written in generally covariant form.

The metric sector is a proposed coarse description of the relational
capacity dynamics. The Einstein action is its selected leading
gravitational continuation. A discrete preferred frame can generate
dimension-four Lorentz-violating operators with large infrared
coefficients through loops \citep{Collins2004}. The absence of an
embedding lattice in the primitive state count leaves room for an
invariant limit; causal-set sprinkling provides a separate example
of Lorentz-compatible discreteness \citep{Bombelli1987,Dowker2004}.
For the present model, the directional cochain counterexample in
H.15 shows why the spatial operator must still be checked.
The cosmological bath selects a state, while laboratory bounds
constrain the effective operators. A common cone and the absence
of unsuppressed Lorentz-violating loop terms remain required.

\subsection{4.2 Dependency Map of the Theory}

The logical flow begins with the three foundational postulates (Information--Geometry, Mass--Entropy, and Many-Pasts), with faithful full-support resolution applied to both states and histories, and runs through the static weak-field chain before reaching the conditional sectors:
\[
\begin{gathered}
\{\text{three postulates}\}\\
\rightarrow\ \text{finite-capacity substrate ontology}
\ \rightarrow\ \text{tetrahedral boundary ensemble}\\
\rightarrow\ \text{faithful resolution of states and histories}
\ \rightarrow\ \text{local replacement / history export}\\
\ \rightarrow\ \text{edge transport / loop dressing / source map}\\
\rightarrow\ \text{capacity form of the Einstein constraint}
\rightarrow\ \text{Newton / baseline no-slip / GR PPN},\\
\text{with carrier projection + thermal contact + horizon matching}\\
\rightarrow\ \{a_0,\ \text{RAR},\ \text{nonspherical fields},\ \text{lensing}\}
\ \text{in the leading EFT}.
\end{gathered}
\]
Only then come the conditional and frontier sectors (transport, clusters, cosmology, strong field, and the particle/gauge extensions), each developed as a consequence or completion of the same framework.

Two features of this map matter. First, it is a dependency graph, not an equality of closure status: the ordinary static branch is closed more tightly than the transverse, cosmological, or strong-field sectors, and Part~VIII makes that difference explicit in a closure-status table. Second, Many-Pasts appears at the top of the map because the scale-setting chain uses it: faithful full-support resolution selects the local renewal process on the history space that Many-Pasts supplies. Their combination gives a reversible present/history exchange; neither the Born history measure nor record conditioning alone selects memorylessness.


\section*{Part II. UV Coefficient Chain}
\addcontentsline{toc}{section}{Part II. UV Coefficient Chain}

Part~I fixed what the theory is about. The question now is whether the local capacity-sharing structure can actually be \emph{counted}. If the substrate has finite local capacity, the coefficients that appear in the continuum weak-field theory should not be free continuum parameters; they should descend from a finite local boundary problem. The next five sections follow that problem through: the smallest boundary cell that can carry capacity and close isotropically, the weighting that selects well-closed configurations, the cost of neighboring cells disagreeing, the local returns that dress that cost, and the continuum coefficient they leave behind. The baseline calculation is internal to the microscopic construction. Appendix~B separately propagates nearby discrete branches to the Newton scale, and Appendix~L records that this comparison is postdictive rather than historically blind.

Several of the ultraviolet choices below may look at first like independent tunings: the tetrahedral cell, the seven labels, the injective assignment, the parity doubling, the admissibility kernel, the transverse export, and the electron anchor. None is a phenomenological knob, and none varies from galaxy to galaxy. Section~5 and Appendix~B derive the finite ensemble and its entropy; Appendix~C derives the edge projection; Appendices~D and~H derive the replacement process, the reversible history export, the factorized lightest branch, and the decorated marked-transfer action; Appendix~L records the historical fork accounting. The remaining microscopic task is to embed that finite transfer vertex in a stable geometric continuum, through a specified GFT action or a capacity-decorated CDT transfer matrix in a controlled extended regime. A critical trajectory that removes the geometric cutoff is a stronger, separate test.

\section{5. Why a Tetrahedral Boundary Ensemble}

The problem is to find a finite boundary cell that can carry channel entropy, close isotropically, and hand a scalar response to the continuum. In three spatial dimensions the minimal volumetric simplex is a tetrahedron. The construction uses five ingredients:
\begin{itemize}
\item a tetrahedral volumetric cell;
\item half-integer primitive data on the two sides of each shared face;
\item positive matching after the two sides are placed in one orientation;
\item four distinguishable relational ports with single-copy channel capacity;
\item binary cell orientation.
\end{itemize}
This package is not presented as the only possible ultraviolet completion. Postulate~II assigns a half-integer primitive spin
\[
j_0=\frac12,\frac32,\frac52,\ldots
\]
to each side of a shared face. Before gluing, the pair spans
\[
V_{j_0}\otimes V_{j_0}
=
\bigoplus_{J=0}^{2j_0}V_J.
\]
After orientation transport, positive coherent matching is generated by a nonnegative mismatch operator whose unique null space is the maximal coupled sector. Its sharp limit is therefore
\[
\mathcal G_{\rm sharp}=P_{2j_0},
\]
and the transmitted face alphabet has
\[
|M|=\dim V_{2j_0}=4j_0+1
\]
states. Appendix~B derives this operator on the full tensor product, gives its finite-width spectrum, and states the Hessian tests that would falsify the primitive matching rule.

The marked sector provides a second, more selective relation. The independently constructed present/history response is a product of two spatial vectors,
\[
\mathcal H_{\rm mark}=V_1^P\otimes V_1^H
=V_0\oplus V_1\oplus V_2.
\]
The complete nonmaximal information left by maximal fusion is
\[
Q(j)
=
\left(V_j\otimes V_j\right)\ominus V_{2j}
=
\bigoplus_{J=0}^{2j-1}V_J.
\]
Because both sums are multiplicity-free,
\[
\mathcal H_{\rm mark}(s)\cong Q(j)
\quad\Longleftrightarrow\quad
j=s+\frac12.
\]
The closure response is a three-component spatial vector, so in $d=3$ it
carries $s=1$. In the selected sharp reversible vertex, unitarity makes the
complete marked carrier the retained nonmaximal fusion output, and therefore
\[
j_0=\frac32,
\qquad
V_{2j_0}=V_3,
\qquad
|M|=7.
\]
This is a derivation within the selected sharp reversible incidence vertex,
not a new minimality axiom. Unitarity fixes the retained output to the
orthogonal fusion complement, and completeness of the nine-state H.9 carrier
then makes the complement map an isometric identification. The result fixes
the carrier space; the hard-core realization, flags, phases, and detailed
routing remain selected structure of the displayed microscopic vertex. If
that vertex is replaced, $j_0=\tfrac32$ remains the branch selected by the
discrete audit of Appendix~B.5.

The four faces are relationally distinct ports. A displayed letter is the occupation of one mode in a single cell-level channel resource; the fermionic single-copy ceiling therefore forbids the same channel from being routed through two ports at once. This gives an injective assignment without antisymmetrizing away the port labels. The two global orientations remain distinct microscopic states. The resulting count is
\[
\Omega_{\mathrm{tet}}=2P(7,4)=1680,
\]
and the combinatorial sharing ceiling is
\[
g_{\mathrm{share,max}}=\ln(1680)=7.42654907240.
\]
The exact equality
\[
16=7+9
\]
has a direct meaning: in the sharp primitive pair map, seven states form the geometric link and the other nine can be retained by the marked record. Within the physical three-dimensional rotation algebra the conditional chain is
\[
d=3
\Longrightarrow
V_1^P\otimes V_1^H
\Longrightarrow
j_0=\frac32
\Longrightarrow
V_3
\Longrightarrow
7.
\]
It is not asserted as a theorem under dimensional continuation to arbitrary $\mathrm{Spin}(d)$.

The exact $K^2$ spectrum and branch audits are given in Appendix~B. The $j$-labeled tetrahedron used here coincides with the quantum tetrahedron of simplicial spin networks \citep{Barbieri1998,BaezBarrett1999}, whose discrete geometric spectra \citep{RovelliSmolin1995} arise from the same $SU(2)$ representation theory. The present construction differs in weighting these states by admissibility closure rather than by a spin-foam amplitude, and in routing them to a capacity entropy rather than to area and volume operators.


\section{6. Admissibility Closure}

\subsection{6.1 Minimal isotropic kernel}

Not every boundary configuration should count equally. The raw combinatorial ensemble is too permissive to be the complete ultraviolet input: some configurations sit close to the regular closure pattern expected of a smooth local cell, while others are badly distorted. Admissibility closure is the statement, in its mildest form, that more poorly closed configurations contribute less to the coarse ensemble. The minimal rotationally invariant measure of that distortion is a single quadratic closure-defect scalar $K^2$, and the weighting it induces is
\[
p_\eta(b)\propto e^{-\eta K^2(b)}.
\]
Normalization, isotropy, and a fixed quadratic closure moment select this maximum-entropy kernel. Higher invariants such as $K^4$ carry additional ultraviolet information and enter as subleading refinements. One distinction is fixed here and holds throughout: this capacity (admissibility) closure is soft, and its invariant $K^2$ is strictly positive on the ensemble, while the exact Gauss closure of local gauge redundancy is a different operator on a different tensor factor. The two cannot be one operator: Appendix~Q proves the impossibility and locates every quantity of the numerical spine on the gauge-inert factor.

\subsection{6.2 Closure condition and uniqueness}

The admissibility precision $\eta$ is not chosen externally; it is fixed by maximizing the normalized closure evidence. Tetrahedral closure is the vanishing of the three-component oriented-face sum, so the closure-defect space is three-dimensional, and the quadratic family on it carries a determinant weight $\eta^{3/2}$. The closure-evidence functional is therefore
\[
\mathcal F(\eta)=\ln Z(\eta)+\tfrac32\ln\eta,
\]
and its stationary point gives the closure condition
\[
\langle K^2\rangle_\eta = \frac{3}{2\eta},
\]
in which the factor $3/2$ is the determinant weight of the three independent closure components, while the discreteness and multiplicities of the spectrum stay inside the exact sum $Z(\eta)$. This is the stationary normalized-evidence point of the exact closure spectrum, and it is a maximum rather than a bare root (Appendix~B.2).

On that spectrum it is unique,
\[
\eta_* = 0.0298668443935.
\]
The closed branch is locally stiff: small fractional changes in $\eta$ produce only small fractional changes in the downstream effective sharing entropy.

\subsection{6.3 Effective sharing entropy}

The admissibility-weighted effective sharing entropy is
\[
g_{\mathrm{share,eff}} = 7.41980002357.
\]
The gap between $g_{\mathrm{share,max}}$ and $g_{\mathrm{share,eff}}$ is the difference between the raw combinatorial ceiling and the admissibility-closed effective boundary entropy that actually propagates into observable couplings.

The continuum description does not inherit the naive channel-counting ceiling; it inherits the portion of the channel space that survives after closure is imposed. The downstream couplings should therefore be read as consequences of admissibility-closed sharing, not of raw combinatorics alone.

With $\eta_*$ fixed, the effective sharing entropy carries no remaining freedom; the exact spectrum, multiplicities, and uniqueness proof are given in Appendix~B.2.


\section{7. Edge Kernel and Tree-Level Coupling}

Admissibility determines the capacity of one cell. The edge kernel determines the cost when neighboring cells differ, and that cost becomes the continuum stiffness: stronger resistance to local disagreement makes capacity deficits spread less readily. The same ultraviolet closure data fix both quantities. The geometric bridge is the tetrahedral identity
\[
\sum_{i=1}^4 \hat n_i\hat n_i^{\mathsf T}=\frac{4}{3}I_3,
\]
which implies a channel-averaged transverse fraction of $2/3$ and gives the bare edge smoothness coupling
\[
J_{\mathrm{bare}}=\frac{2}{3}\eta_*.
\]
If adjacent cells disagree strongly the edge pays a larger penalty; if they agree, the penalty is small. The factor $2/3$ is the geometric fraction that survives after averaging the four tetrahedral channel directions into the isotropic continuum limit, which is the part of the disagreement that the scalar sharing channel actually carries.

For a $z=4$ regular coarse adjacency graph, the tree-to-lattice reduction then yields
\[
J_{\mathrm{eff}}^{\mathrm{tree}}=\frac{J_{\mathrm{bare}}}{3}=\frac{2\eta_*}{9}.
\]
The division by $3$ comes from the branching geometry of the rooted $z=4$ graph. One neighboring link points back toward the source, while the remaining $z-1=3$ links carry forward transport into the tree. The net long-range transport $J_{\mathrm{eff}}^{\mathrm{tree}}$ is therefore the portion of the microscopic edge penalty that survives this local branching.

\paragraph{Origin of the horizon target.}
The horizon target
\[
\sigma_*=\frac{\pi}{g_{\mathrm{share,eff}}}
\]
is the closure-consistency value required by the horizon-normalized field convention. In the admissibility-closed boundary ensemble, one active microscopic sharing unit carries effective entropy $g_{\mathrm{share,eff}}$. In the continuum normalization used for the weak-field scalar, the occupancy variable is normalized by
\[
S=\pi Q_{\mathrm{occ}},
\]
so a coarse horizon-normalized channel with occupancy $Q_{\mathrm{occ}}=1$ carries entropy $\pi$ in the $S$-field convention. If $\sigma_*$ denotes the asymptotic conditional-independence weight seen by the rooted shell hierarchy (Appendix~B.3), consistency between the boundary entropy count and the horizon-normalized continuum field requires
\[
\sigma_*\,g_{\mathrm{share,eff}}=\pi,
\]
and therefore
\[
\sigma_*=\frac{\pi}{g_{\mathrm{share,eff}}}=0.42340665\ldots .
\]
The factor $\sigma_*$ matches the admissibility-closed microscopic entropy normalization to the horizon-normalized scalar-field convention. The rooted shell observable converges rapidly to this closure target, constraining the nonlocal correction at small shell depth. The four tetrahedral channel directions average to the isotropic tensor structure in the continuum limit, so the combinatorial data that fix admissibility also fix tree-level transport. Appendix~C gives the shell hierarchy and phase-selection checks.


\section{8. Finite-Loop Renormalization}

Tree level is not the end of the ultraviolet chain. The full lattice admits local closed-return motifs that recycle part of the transmitted information before it contributes to net coarse transport. The leading correction is organized as a local Dyson self-energy dressing,
\[
J_{\mathrm{eff}}^{(\mathrm{ren})}
=
\frac{J_{\mathrm{eff}}^{\mathrm{tree}}}{1+J_{\mathrm{eff}}^{\mathrm{tree}}\Sigma_{\mathrm{ret}}}.
\]

A purely tree-like transmission rule would let the relevant amplitude move outward once and never locally return; a real coarse graph is not that simple. Some of the transmitted information cycles back through short closed motifs before contributing to long-distance transport. The renormalized coupling is therefore the true stiffness felt by the coarse field after these local returns have been resummed.

The structure of that self-energy is not a generic loop number. The returns split into seven sector-diagonal channels and one collective mode. The seven are the face-label channels, each returning independently without mixing. The one is the permutation-symmetric combination across channels, which returns as a shared closure-singlet rather than as a channel-specific loop, and it is weighted by the same transverse projection and branch-dilution factors that define the tree edge map,
\[
\left(\frac{2}{3}\right)\left(\frac{1}{3}\right)=\frac{2}{9}.
\]
The leading local self-energy is therefore
\[
\Sigma_{\mathrm{ret}}=7+\frac{2}{9}=\frac{65}{9}.
\]
Equivalently, on the seven-channel scalar return space,
\[
R_{\mathrm{ret}}=I_7+\frac{2}{9}P_{\mathrm{sing}},
\qquad
P_{\mathrm{sing}}=|u\rangle\langle u|,
\qquad
u=\frac{1}{\sqrt7}(1,\ldots,1),
\]
with $\Sigma_{\mathrm{ret}}=\mathrm{Tr}(R_{\mathrm{ret}})$. The orthogonal six-dimensional sum-zero sector carries no net scalar charge in the coarse branch and so adds no separate scalar return. The singlet weight is fixed by the tree map, not introduced here, so no new loop parameter appears. Permutation symmetry fixes the existence of the singlet but not the placement of the suppression factors on it alone; that placement is the minimal-return-operator reading whose graph-level derivation Appendix~C.3 records as the outstanding audit. The induced uncertainty is bounded: replacing $\Sigma_{\mathrm{ret}}=65/9$ by $7$, $7+2/3$, or $8$ shifts $J_{\mathrm{eff}}^{(\mathrm{ren})}$ and $\gamma$ by at most $0.5\%$ and leaves $G$, $a_0$, and $\kappa/\gamma$ exactly unchanged, since $J_{\mathrm{eff}}^{(\mathrm{ren})}$ cancels in the source-to-stiffness ratio (Appendix~C.5).
\[
c_{\mathrm{loop}}^{(\mathrm{ren})}
\equiv
\frac{J_{\mathrm{eff}}^{(\mathrm{ren})}}{J_{\mathrm{eff}}^{\mathrm{tree}}}
=
\frac{1}{1+J_{\mathrm{eff}}^{\mathrm{tree}}\Sigma_{\mathrm{ret}}}
\approx 0.95426,
\]
and
\[
J_{\mathrm{eff}}^{(\mathrm{ren})}\approx 0.00633348.
\]
This reproduces the shell-target crossing near $J_{\mathrm{bare,cross}}\sim 0.019$ at the $0.05\%$ level.

The finite renormalization is written as an explicit local self-energy. The remaining audit task is the independent graph-level derivation of the relative diagonal and singlet weights of the same scalar-return operator, not the introduction of any new loop parameter.


\section{9. Continuum Stiffness and SI Normalization}

The last UV step reads the lattice weighting as a quantum action rather than a thermal one: the lattice quadratic form is interpreted as a Euclidean action weight,
\[
\frac{I_E}{\hbar}
=
\frac{J_{\mathrm{eff}}^{(\mathrm{ren})}}{2}
\sum_{a,i}(Q_a-Q_{a+L_*\hat n_i})^2,
\]
where the sum runs over one sublattice representative $a$ of each bipartite primitive cell and its four outgoing bonds $\hat n_i$, so each undirected nearest-neighbor edge is counted once (the convention of Appendix~C.4). The microscopic four-cell is assigned the volume
\[
\Delta V_4=\frac{L_*^4}{c}
\]
as a coarse-graining convention: the abstract tetrahedral cell complex has no space-filling regular-tetrahedron Euclidean embedding (Appendix~C.5), so cell volumes and face areas enter as normalization conventions of the coarse map, not as geometry supplied by the graph.
Up to this point the derivation has determined a dimensionless lattice weighting. The continuum EFT, however, needs a dimensionful coefficient multiplying derivatives of a field in spacetime. The Euclidean-action interpretation upgrades the lattice closure data into a continuum action density with the right units and the right covariant target.

The same tetrahedral identity used in the edge-kernel reduction then yields the continuum coefficient for the occupancy field $Q_{\mathrm{occ}}$,
\[
\gamma_Q=\frac{4\hbar c}{3L_*^2}J_{\mathrm{eff}}^{(\mathrm{ren})}.
\]
Here $L_*$ is the canonical tetrahedral spacing, with one coarse cell carrying volume $L_*^3$ up to the fixed cell-shape convention, and $J_{\mathrm{eff}}^{(\mathrm{ren})}$ is the loop-dressed edge coupling. The numerical factor $4/3$ is the isotropic projection
\[
\sum_i \hat n_i\hat n_i^T=\frac43 I_3
\]
that turns the tetrahedral edge directions into the continuum gradient tensor.

The field normalization is fixed by horizon capacity:
\[
S=\pi Q_{\mathrm{occ}}.
\]
Therefore the canonical EFT coefficient in the $\frac{\gamma}{2}(\partial S)^2$ convention is
\[
\gamma=\frac{4\hbar c}{3\pi^2L_*^2}J_{\mathrm{eff}}^{(\mathrm{ren})}.
\]
Physically, $\gamma$ is the continuum stiffness of the entanglement-capacity field. A larger $\gamma$ makes spatial gradients more costly and suppresses the capacity-deficit response to a given source; a smaller $\gamma$ allows larger variations of the field.
The faithful sector-resolution principle fixes $L_*$ without using $G$. It is nevertheless useful to define the gravitational scale induced by this length,
\[
G_*:=\frac{c^3L_*^2}{\hbar}.
\]
Then the stiffness may be written in Einstein-normalized form as
\[
\gamma=\frac{4J_{\mathrm{eff}}^{(\mathrm{ren})}}{3\pi^2}\frac{c^4}{G_*}.
\]
This is the same algebra as the familiar Planck-cell rewrite, but read in the opposite direction: the substrate cell length induces the gravitational scale rather than being chosen by first inserting the measured value of $G$. Within the Euclidean-action and cell-volume conventions stated above, the SI-normalized stiffness coefficient is fixed; because the absolute normalization of an isolated scalar functional is conventional, the invariant content of this step is the ratio $\kappa/(\gamma S_\infty)$ that the weak-field matching of Section~11 consumes. In that ratio $J_{\mathrm{eff}}^{(\mathrm{ren})}$ cancels (Appendix~C.5), so the loop-dressed coupling carries no content for the Newton normalization; its nontrivial input enters the stiffness itself and the dynamical and galactic sectors.

This completes the micro-to-continuum coefficient chain. The tetrahedral ensemble determines the effective sharing entropy; the edge kernel and loop dressing turn that entropy into a discrete stiffness; and the Euclidean matching turns the discrete stiffness into the continuum coefficient $\gamma$ of the weak-field EFT.

\paragraph{Closed UV-to-IR chain.}
The UV coefficient chain can now be summarized as
\[
\{\Omega_{\mathrm{tet}},\,K^2,\,\eta_*,\,g_{\mathrm{share,eff}},\,L_*,\,J_{\mathrm{bare}},\,J_{\mathrm{eff}}^{\mathrm{tree}},\,\Sigma_{\mathrm{ret}},\,J_{\mathrm{eff}}^{(\mathrm{ren})},\,\gamma\}
\longrightarrow
\{\kappa,\,G,\,a_0,\,g_{\mathrm{obs}}(g_{\mathrm{bar}})\}.
\]
The first bracket is the micro-to-continuum closure chain; the second collects the weak-field observables it feeds. All later weak-field coefficients come from this chain.

The ordinary branch uses the Euclidean action-kernel interpretation and the physical matching of Appendix~C.5. Fuller inhomogeneous dynamics must reproduce both prescriptions.

\section*{Part III. The Closed Static Branch: Einstein Gravity in the Capacity Variable}
\addcontentsline{toc}{section}{Part III. The Closed Static Branch: Einstein Gravity in the Capacity Variable}

\section{10. Einstein Parent Action and the Reduced Capacity Frame}

The ordinary longitudinal capacity branch has the ordinary metric action as its covariant parent; no independently varied capacity scalar is added to it:
\[
I_{0}[g,\psi]
=
\frac{c^3}{16\pi G}\int_{\mathcal M}d^4x\,\sqrt{-g}\,(R-2\Lambda)
+I_{\mathrm{GHY}}[g]
+I_{\mathrm{matter}}[g,\psi],
\]
where covariant coordinates use $x^0=ct$. When the time integral is instead written in seconds, $dx^0=c\,dt$ and the ADM prefactor is correspondingly $c^4/(16\pi G)$. Matter is coupled once, to one physical metric. The capacity functional is the static scalar-constraint reduction of $I_0$, expressed in a different field coordinate. This statement can be proved without appealing to the final Poisson equation.

\paragraph{Static scalar reduction of Einstein--Hilbert gravity.}
Use Newtonian gauge,
\[
ds^2
=-\left(1+\frac{2\Phi}{c^2}\right)(dx^0)^2
+\left(1-\frac{2\Psi}{c^2}\right)\delta_{ij}dx^idx^j,
\qquad x^0=ct,
\]
and retain the static scalar sector through quadratic order. In ADM variables \citep{ADM1962} the shift and extrinsic curvature vanish, so the Einstein--Hilbert plus Gibbons--Hawking--York action is
\[
I_{\mathrm{ADM}}^{\mathrm{static}}
=\frac{c^4}{16\pi G}\int dt\,d^3x\,N\sqrt h\,{}^{(3)}R
+I_{\mathrm{matter}}^{\mathrm{static}}
+I_\infty.
\]
Expanding $N=1+\Phi/c^2$ and $h_{ij}=(1-2\Psi/c^2)\delta_{ij}$, cancelling the reference boundary term at infinity, and using $I_{\mathrm{matter}}^{(1)}=-\int dt\,d^3x\,\rho\Phi$, gives
\[
I_{0,\mathrm{scal}}^{(2)}[\Phi,\Psi]
=
\int dt\,d^3x
\left[
\frac{1}{8\pi G}
\left((\nabla\Psi)^2-2\nabla\Phi\!\cdot\!\nabla\Psi\right)
-\rho\Phi
\right].
\]
The lapse perturbation remains a constraint variable. Its variation and the spatial-scalar variation give, respectively,
\[
\nabla^2\Psi=4\pi G\rho,
\qquad
\nabla^2(\Phi-\Psi)=0.
\]
For the isolated solution with regular interior matching and no independent harmonic boundary data, asymptotic flatness gives $\Phi=\Psi$. Eliminating $\Psi$ therefore produces
\[
\boxed{
I_{\mathrm{Newton}}[\Phi]
=
\int dt\,d^3x
\left[
-\frac{(\nabla\Phi)^2}{8\pi G}
-\rho\Phi
\right].
}
\]
Thus no independent scalar stress tensor is needed to create the linear potential, and the equality $\Phi=\Psi$ in the baseline branch is a metric constraint equation rather than an anisotropic-stress assumption. Appendix~N gives the expansion, boundary bookkeeping, and degree-of-freedom audit in full.

\paragraph{Exact reduced-action identity.}
On the renormalized static branch, write $S_{\mathrm{ent}}=S_\infty-\delta S$. Source-independent extensive terms are removed by the vacuum normalization proved for the regulated joint measure in Section~25. The capacity functional is then
\[
I_{\mathrm{cap}}^{\mathrm{static}}[\delta S;\rho]
=
\int dt\,d^3x
\left[
-\frac{\gamma}{2}(\nabla\delta S)^2
+\kappa\rho\,\delta S
\right].
\]
The field redefinition
\[
\delta S=-\frac{2S_\infty}{c^2}\Phi
\]
turns it into
\[
I_{\mathrm{cap}}^{\mathrm{static}}
=
\int dt\,d^3x
\left[
-\frac{2\gamma S_\infty^2}{c^4}(\nabla\Phi)^2
-\frac{2\kappa S_\infty}{c^2}\rho\Phi
\right].
\]
Using
\[
G=\frac{c^2\kappa}{8\pi\gamma S_\infty}
\]
gives the action-level equality
\[
\boxed{
I_{\mathrm{cap}}^{\mathrm{static}}
=Z_S I_{\mathrm{Newton}},
\qquad
Z_S\equiv\frac{2\kappa S_\infty}{c^2}.
}
\]
The field-independent factor $Z_S$ cannot affect the classical reduced equations. It can matter when the UV construction is asked to normalize fluctuations or correlation functions, but it does not represent a second determination of $G$ and it does not license adding $I_{\mathrm{cap}}$ to $I_0$.
The equality assumes the same asymptotically flat or Dirichlet boundary data on both sides. At a finite boundary the Newton surface term and its image under $\delta S=-2S_\infty\Phi/c^2$ must be included as well.

\paragraph{Scope of the background-covariant notation.}
For transport calculations the same reduced equation is packaged as
\[
I_{\mathrm{cap}}[S_{\mathrm{ent}};\chi\,|\,g_{\mathrm{ref}}]
=
\int d^4x\,\sqrt{-g_{\mathrm{ref}}}
\left[
-\frac{\gamma}{2}g_{\mathrm{ref}}^{\mu\nu}
\partial_\mu S_{\mathrm{ent}}\partial_\nu S_{\mathrm{ent}}
-\lambda S_{\mathrm{ent}}-\kappa\chi S_{\mathrm{ent}}
\right].
\]
The vertical bar is essential: $g_{\mathrm{ref}}$ and the reduced source projection $\chi$ are held fixed while $S_{\mathrm{ent}}$ is varied. This notation is useful for extending the reduced response in time, but it is not a covariant scalar--tensor parent action. In the static nonrelativistic sector $\chi\simeq\rho$; covariantly the source is the full stress tensor through $I_{\mathrm{matter}}[g,\psi]$.

\paragraph{Capacity coefficients and the source theorem.}
The UV calculation still fixes how the geometric constraint is coordinatized by the substrate variable. With
\[
\sigma_{\mathrm{def}}=\frac{\rho}{\kappa_m(L_*)},
\]
the Green-matched projection is
\[
\nabla^2\delta S
=-\frac{3L_*}{4G_{\mathrm{tet}}(0)}\,\sigma_{\mathrm{def}},
\qquad
\frac{\kappa}{\gamma}
=\frac{3L_*}{4G_{\mathrm{tet}}(0)\kappa_m(L_*)}.
\]
This is the microscopic map between defect density and the capacity coordinate on the Einstein constraint surface. It is not an extra matter coupling in the covariant parent theory.

\paragraph{The length backbone of Newton's constant.}
The weak-field normalization is
\[
G=\frac{c^2\kappa}{8\pi\gamma S_\infty}.
\]
Faithful full-support resolution and the positive marked-transfer spectrum fix
\[
L_*=-\frac32\,Z_e\lambda_e\ln\!\left(1-e^{-7g_{\mathrm{share,eff}}}\right),
\qquad
\lambda_e=\frac{\hbar}{m_ec},
\]
and therefore induces
\[
G_*=\frac{c^3L_*^2}{\hbar}
=
\frac94\,\frac{\hbar c}{m_e^2}\,
Z_e^2\ln^2\!\left(1-e^{-7g_{\mathrm{share,eff}}}\right).
\]
Substituting the source-map identities
\[
\frac{\kappa}{\gamma}=\frac{3L_*}{4G_{\mathrm{tet}}(0)\kappa_m(L_*)},
\qquad
\kappa_m(L_*)=\frac{\hbar}{cL_*\ln2},
\qquad
S_\infty^{\mathrm{cell}}=\frac{3\ln2}{32\pi G_{\mathrm{tet}}(0)},
\]
into the weak-field expression gives identically
\[
G=\frac{c^3L_*^2}{\hbar}=G_*.
\]
Thus the electron recurrence fixes $L_*$ and the candidate scale $G_*=c^3L_*^2/\hbar$. Physical $G=G_*$ follows only after the selected Einstein/horizon normalization $S_\infty=S_0$ is imposed. The stiffness and source coefficient give a consistent static-EFT representation of that matched length backbone rather than an independent numerical determination. The comparison with the measured Newton constant is therefore a test of the electron-anchored candidate together with this physical matching.

The decorated scale gives $G_*=6.6742890813\times10^{-11}\,\mathrm{m^3\,kg^{-1}\,s^{-2}}$, or $-0.073\sigma$ relative to CODATA. Because both the early entropy construction and the later residual-closing vertex were developed with the discrepancy known, Appendix~L treats this as a high-precision postdiction. The nontrivial content is the shared action that also fixes the two charged-lepton corrections.


\section{11. Capacity Variable, Bridge Law, and Variational Status}

Varying the reduced capacity-frame functional with respect to $S_{\mathrm{ent}}$ gives
\[
\gamma\Box S_{\mathrm{ent}}=\lambda+\kappa\chi.
\]
On the renormalized static, nonrelativistic branch this becomes
\[
\nabla^2\delta S=-\frac{\kappa}{\gamma}\rho.
\]

Define the surviving fractional capacity
\[
q(x)\equiv\frac{S_{\mathrm{ent}}(x)}{S_\infty}
=1-\frac{\delta S(x)}{S_\infty}.
\]
The action reduction above fixes the weak-field bridge directly,
\[
\frac{\Phi}{c^2}=-\frac{\delta S}{2S_\infty}.
\]
Equivalently,
\[
q=1+\frac{2\Phi}{c^2}+O(c^{-4}).
\]
The bounded nonlinear rule
\[
N^2=q
\]
is the continuous multiplicative completion selected by the capacity-composition rule. In a normalized static spherical vacuum exterior, $q$ equals the radial geometric scalar
\[
q=h^{ab}\partial_aR\partial_bR
=1-\frac{2GM_{\mathrm{MS}}}{c^2R},
\]
and equals \(N^2\) with the asymptotic Killing normalization.
Inside matter the clock quantity \(q_\tau=N^2\) generally differs
from the radial quantity \(q_R=(\nabla R)^2\).
Appendix~N.6 derives the exact static lapse source
\(D^2N=(4\pi G/c^4)N(\mathcal E+\mathcal S)\) and its nonlinear
capacity equation. A generic lapse remains foliation dependent,
so \(N^2=q\) alone is not a covariant constraint.
The metric-only parent treats these reduced quantities as
geometric observables without additional Cauchy data. A general
capacity observable may depend on a relational region, its state
and algebra, a clock, and boundary reference data.

Combining the static source equation with the weak-field bridge gives
\[
G=\frac{c^2\kappa}{8\pi\gamma S_\infty}.
\]
This relation uses only the invariant combination $\kappa/(\gamma S_\infty)$: a fixed-epoch rescaling of the entropy units changes $S_\infty$ and $\kappa$ together and leaves the observable potential unchanged.

\paragraph{Two coordinates on one response.}
The metric parent solves the Hamiltonian and spatial constraints for $\Phi$ and $\Psi$; the capacity frame uses $\delta S$ as a field coordinate on that reduced solution. Around a constant background, a canonical scalar stress would begin as $(\partial\delta S)^2=O(\rho^2)$ and could not be the source of the observed $O(\rho)$ potential. The action identity removes that mismatch: the linear capacity response is the reduced metric constraint itself.

\paragraph{Matter enters once.}
For the nonrelativistic static branch, the microscopic source theorem reduces the full metric source to the defect density $\rho$. Covariantly, matter enters only through $I_{\mathrm{matter}}[g,\psi]$, so the full stress tensor gravitates, including trace-free radiation, and the Bianchi identity enforces the usual conservation law. The notation $\chi\simeq\rho$ belongs only to the reduced nonrelativistic source map; an explicit universal term $S_{\mathrm{ent}}T^\mu{}_{\mu}$ is neither required nor adopted.

\paragraph{Parent-action decision.}
The preferred minimal construction for the ordinary branch is therefore metric-only:
\[
I_{\mathrm{parent}}^{\mathrm{long}}=I_0[g,\psi],
\qquad
\delta S=\delta S[g,\psi]\quad\hbox{after constraint reduction}.
\]
It propagates the two tensor polarizations of general relativity and no extra scalar. Constrained-clock and scalar--tensor alternatives remain useful control cases, but both add structure and generically add a mode; Appendix~N records why neither is selected. The galactic excess belongs to the separate carrier-resolved contact of Sections~15--16; its leading EFT is explicit, while its microscopic metric vertex remains conditional.


\section{12. Newtonian Gravity and the Point-Source Limit}

In the renormalized static weak-field sector the scalar equation reduces to
\[
\nabla^2\delta S=-\frac{\kappa}{\gamma}\rho.
\]
After the background is renormalized away and the source is taken to be nonrelativistic, the deficit field obeys an ordinary Poisson equation. Its mathematical structure is the one used in standard weak-field gravity, with $\delta S$ as the field variable.

For a point source $M$,
\[
\delta S(r)=\frac{\kappa M}{4\pi\gamma r}.
\]
Using the bridge law,
\[
\frac{\Phi}{c^2}=-\frac{\delta S}{2S_\infty},
\]
the gravitational acceleration becomes
\[
g(r)=\frac{c^2\kappa}{8\pi\gamma S_\infty}\frac{M}{r^2}
=
\frac{GM}{r^2}.
\]
Thus Newtonian gravity is recovered as the weak-field response of the entanglement-capacity medium: the sourced scalar equation and the bridge law together imply the familiar point-mass force law, with nothing further assumed.

\paragraph{Interpretation.}
A point defect produces a $1/r$ capacity deficit, and the bridge maps its gradient to the Newtonian inverse-square force. Ordinary gravity is the small-deficit, weak-curvature limit of the extended capacity strain around localized defects. Section~24 exhibits the same structure on the discrete substrate: cell-by-cell re-equilibration screens at sub-cell range, whereas a conserved capacity current with maintenance sinks obeys the massless graph-Poisson equation and produces the $1/r$ deficit. The Newtonian form therefore diagnoses the conservation law behind it.


\section{13. Electron Anchor: One-Bit Mass and Seven-Sector Length}

The electron supplies the elementary mass anchor $m_e/\ln2$. Its Compton scale also assigns the length associated with the seven-sector effective support. The decorated marked-transfer action closes the finite correction to that scale and routes the same correction through the heavier charged-lepton shells. These are two readings of one dimensionful datum, not independent measurements, so the weak-field normalization remains a single calibration.

\subsection{13.1 Why the electron is the anchor}

The mass--entropy map needs an elementary anchor because the elementary matter sector is the localized defect sector. The electron is the lightest simple charged fermionic defect, and its mass is not obscured by hadronic or QCD dressing. A single fermionic face-exclusion defect carries the canonical increment
\[
\Delta S_f=\ln 2,
\]
one bit of missing entanglement, because an excluded face is a binary occupied/unoccupied defect of the local network.

\subsection{13.2 One-bit mass anchor}

At the electron Compton scale $\ell=\lambda_e$ the mass--entropy map reads
\[
\kappa_m(\lambda_e)=\frac{m_e}{\ln 2}.
\]
Dividing the electron mass by the fixed one-bit increment fixes the mass-per-entropy conversion at the electron's own scale. Run back to the cutoff cell, the conversion is
\[
\kappa_{m,\mathrm{UV}}=\frac{\hbar}{cL_*}\frac{1}{\ln 2},
\]
with canonical running law
\[
\kappa_m(\ell)=\kappa_{m,\mathrm{UV}}\left(\frac{L_*}{\ell}\right)^{1+\alpha_{\mathrm{cl}}},
\qquad
\alpha_{\mathrm{cl}}=0
\]
in the closed branch. One bit fixes the electron-scale conversion; the running law carries it to the UV scale; the same conversion then feeds the weak-field source map. This is the only point at which the mass anchor enters gravity.

\subsection{13.3 Seven-sector length anchor}

The same electron anchors the cell length through the marked support-to-rate map. Faithful full-support resolution uniquely selects the memoryless kernel. Within the stated recurrence mass functional, fermionic exclusion and electron lightness jointly select the maximum channel count $k=7$ and vanishing inter-channel correlation $\Delta_7=0$. On the commutative renewed-state algebra, the state-weighted determinant of the likelihood multiplication operator is
\[
\Delta_{\tau_p}(\mathsf R)
=\exp[\tau_p(\ln\mathsf R)]
=e^{-g_{\mathrm{share,eff}}}.
\]
Multiplicativity across the seven factorized clouds gives
\[
r=\Delta_{\tau_p^{\otimes7}}(\mathsf R^{\otimes7})
=e^{-7g_{\mathrm{share,eff}}}.
\]
This is the record-conditioned geometric transfer rate, not the collision probability of two independently sampled blocks. The positive survival spectrum gives the baseline scale
\[
L_*^{(0)}=-\frac32\lambda_e\ln(1-r),
\]
with $3/2$ the transverse export factor of Appendix~C.5.

Appendix~H realizes this determinant transfer in a finite charged action. The construction separates the baseline recurrence from the defect-bound closure response, so the one-channel mixing projector, seven-channel survival operator, and marked-fiber determinant remain distinct. The finite occupation bound in H.2 excludes fewer-channel competitors in the specified positive-generator class; H.5 establishes the unique product minimum over seven-channel readout laws at fixed marginals and marked vertex.

This does not mean an electron is a single tetrahedral cell carrying seven simultaneous labels. The local ensemble supplies seven distinguishable dressing layers. Each layer has fermionic occupation at most one, and the recurrence mass minimum occupies all seven once. The electron is the lightest coherent one-bit defect on that selected branch and exports the transverse share of the resulting support. The Compton scale calibrates the substrate length hierarchy; the one-bit mass calibrates the source map. The two uses impose a nontrivial joint requirement on the electron's role in the gravitational normalization without duplicating a single input.

\subsection{13.4 Decorated marked-transfer vertex}

The decorated marked-transfer vertex closes the Newton and charged-lepton residuals with one marked-fiber weight and no continuously adjusted coefficient (Section~13.5). The baseline Newton normalization, before the marked vertex, was low by about $1.05\%$. Because $G_*\propto L_*^2$, the missing amplitude in the substrate length was $0.5309\%$. The baseline muon ratio required a $0.5124\%$ uplift. At the PDG tau mass the tau ratio required $0.6601\%$, which separates into the same universal uplift and a smaller $0.1470\%$ second-shell factor. Their common sign and scale motivated one small charged response with shell-dependent routing. Appendix~L records that this was an action-level postdiction.

That response also had to be additive. The tetrahedral ensemble already fixed the entropy, closure spectrum, edge projection, source map, and the micro-to-macro coefficient chain. Changing its label count, admissibility weight, or transverse export to close the residuals would move results that did not share the discrepancy. The admissible correction was therefore required to vanish in the unmarked vacuum, reuse the established closure incidence, and act only on the charged transfer graph. This criterion motivates the marked vertex below; it does not determine the answer numerically. The field content, determinant power, and routing still have to follow from the displayed action and survive the alternative-kernel audits.

The renewed closure amplitude has the exact Gaussian representation
\[
e^{-\eta_*C^2/2}
=\int\frac{d^3\xi}{(2\pi)^{3/2}}
\exp\!\left[-\frac12\xi^2+i\sqrt{\eta_*}\,\xi_aC_a\right].
\]
The amplitude-level closure incidence is therefore $\sqrt{\eta_*}$. For each occupied unordered channel pair $e=(m,m')$, the two directed scalar returns have row operator
\[
R_e=\frac27\bigl(\langle m\!\to\!m'|+\langle m'\!\to\!m|\bigr).
\]
The contraction $B_e=\sqrt{\eta_*}R_e$ obeys
\[
B_eB_e^\dagger=2\eta_*\left(\frac27\right)^2
=\frac{8\eta_*}{49}\equiv u.
\]
Its canonical unitary dilation has no-event amplitude $\sqrt{1-u}$. The seven-channel lightest branch activates all $\binom72=21$ pair records, giving
\[
Z_{\rm edge}=\left(1-\frac{8\eta_*}{49}\right)^{21/2}.
\]

The present and history Gaussian strands each have three normalized first excitations. Their ordered products form nine orthonormal marked states, so the internal trace of one hard-core marked fiber gives
\[
\boxed{\zeta_*
=9e^{-g_{\mathrm{share,eff}}}
\left(1-\frac{8\eta_*}{49}\right)^{21/2}}
=0.005123584484947.
\]
This internal trace computes a scalar vertex weight at fixed path branch; it is not a measurement that discards the defect's position label. Free transport applies the same internal contraction on every position branch, as required by Appendix~H.11.
The finite transfer graph has a universal marked alternative, one two-vertex label return, and one second-shell singlet passage. Its exact factors are
\[
Z_\mu=1+\zeta_*,
\qquad
Z_e=(1+\zeta_*)(1+7\zeta_*^2),
\qquad
Z_{\tau,2}=1+\frac27\zeta_*.
\]
With the additive-generator prescription, the dressed scale is
\[
\boxed{L_*=Z_eL_*^{(0)}}.
\]
Appendix~H writes the controlled vertex, decomposes its complete edge Hessian, and enumerates all $56{,}800$ states of the one-step routing block. H.9 also constructs a finite recurrence contraction with $E_e^{\rm evt}=-(3\hbar/2\tau_*)\ln(1-rZ_e)$; its relative energy difference is $7.37\times10^{-26}$. The additive-generator prescription supplies the exact formulas used here.

\paragraph{The anchor is a gauge choice.}
Because counting fixes only dimensionless quantities, exactly one dimensionful measurement must be supplied, and which one is a convention. The physical content of the framework is carried by anchor-invariant statements: the lepton ratios $m_\mu/m_e$ and $m_\tau/m_e$; the hierarchy
\[
\frac{m_e}{m_P}
=-\frac32Z_e\ln(1-e^{-7g_{\mathrm{share,eff}}});
\]
the abundance ratio $\Omega_c/\Omega_b$; and, within its branch assignment, $a_0/cH_0$. The horizon normalization $1/4$ enters through the C.5 matching of the stated cell convention. The electron remains the anchor because its mass is measured most precisely. Appendix~L records that the high-precision marked correction was constructed after the residuals were known, so anchor invariance does not turn the agreement into a blind prediction.

\subsection{13.5 Newton constant and charged-lepton ratios}

The marked action uses no continuously fitted coefficient. Substitution gives
\[
G_*=6.6742890813\times10^{-11}\ {\rm m^3\,kg^{-1}\,s^{-2}}
\qquad(-0.073\sigma),
\]
\[
\frac{m_\mu}{m_e}
=720\frac27Z_\mu
=206.768280237
\qquad(-0.535\sigma),
\]
\[
\frac{m_\tau}{m_e}
=720^2\left(\frac27\right)^4Z_\mu Z_{\tau,2}
=3477.343310
\qquad(-0.125\sigma).
\]
These three comparisons probe one marked-fiber weight through different graph polynomials. The scalar stiffness, source map, and weak-field bridge reproduce the same dressed $G_*$ after substitution, so they remain one consistency chain rather than an independent determination.

In fractional terms the three agreements are about $1.6\times10^{-6}$, $1.2\times10^{-8}$, and $6\times10^{-6}$; equivalently, the calibrated substrate length matches the Planck length to $8.2\times10^{-7}$. The shell integers alone, before the marked response, give $m_\mu/m_e=720\cdot\tfrac27=205.71$ and $m_\tau/m_e=720^2(\tfrac27)^4=3454.6$, already within $0.51\%$ and $0.65\%$ of the measured ratios; the marked vertex supplies the remaining fraction of a percent. A second route to Newton's constant runs through the geometric sector: the absolute angular rule and the event calibration of Appendix~C.8 give $G_0/G_*=1.030$ at the illustrative $m_4=2.35$, and exact equality requires $m_4=2.3241$. The finite-count value of $m_4$ and the interacting coefficient $G_R$ measured on the native ensemble test that route (Sections~23.1 and~24).

\paragraph{Routing-integer audit.}
The displayed action makes the coefficient of $\zeta_*^2$ in $Z_e$ equal to
the number of persistent charged labels: the second marked return sums over
the alphabet, so the factor is $1+7\zeta_*^2$ (Appendix~H.9). Holding
$\zeta_*$ fixed and temporarily treating that integer as unknown, each unit
step in $n$ moves the induced $G_*$ by $2.34$ CODATA standard deviations:
$n=6$ gives $-2.41\sigma$, $n=7$ gives $-0.073\sigma$, and $n=8$ gives
$+2.26\sigma$. Within $-200\le n\le200$, seven is the only integer within two
standard deviations. Dropping the second-order term ($n=0$) gives
$-16.4\sigma$. Appendix~B supplies two readings of this result. In the
selected sharp reversible vertex, the carrier-incidence theorem and the
three-dimensional closure vector fix the same seven internally, and the
Newton comparison checks its routing. For a different microscopic vertex,
the comparison selects the seven-state member of the stated discrete family.
Neither reading is historically blind: the target and the residual were
already known, as Appendix~L records.

Two remarks keep this honest. The resolvability is partly fortunate: $\zeta_*^2\simeq2.6\times10^{-5}$ places the rung spacing just above the CODATA uncertainty, so the integer is measurable at all; a smaller $\zeta_*$ would bury the ladder inside the error bar, and a larger one would leave $n=7$ one near-miss among many. And since the alphabet is symmetric about zero, one might expect the neutral label $m=0$ to drop out of a return sum, giving coefficient six; $n=6$ lies outside the band, so the neutral label demonstrably participates in the second marked return. With the current PDG tau mass, $2/7$ is the only member of the forty-three distinct rationals in $[0,1]$ with denominator at most eleven that lies within one experimental standard deviation. Its residual is $-0.125\sigma$; $3/11$ gives $-1.436\sigma$. This finite candidate test does not assign probabilities to microscopic models.

The same comparison has a dimensionless form. Squaring the hierarchy gives the electron's gravitational coupling,
\[
\frac{G_*m_e^2}{\hbar c}
=\frac94\,Z_e^2\ln^2\!\left(1-e^{-7g_{\mathrm{share,eff}}}\right),
\]
with $9/4$ the square of the transverse export factor and $Z_e^2$ the marked-transfer dressing. The relation also has an area reading. A horizon stores one bit in area $4\ln2\,L_P^2$, while the electron spreads its single bit over the Compton area $\lambda_e^2$. The marked factor changes the finite packing correction without altering the dominant hierarchy $e^{-14g_{\mathrm{share,eff}}}$.

\subsection{13.6 Composite sectors}

For composite hadrons the claim is weaker and different in kind. The relevant quantity is the dressed, vacuum-subtracted bound-state entropy,
\[
m_{\mathrm{hadron}}=\kappa_m(\ell_H)\,S_{\mathrm{ent},H}^{\mathrm{dressed}},
\]
with the dressed budget generated by confinement, gluonic structure, trace-anomaly dynamics, and chiral vacuum reorganization. A lattice derivation of that dressed entropy is not yet available. The present claim is limited to structural compatibility between the mass--entropy map and the standard QCD mass budget. The elementary-fermion anchor is settled in the simple sectors; the hadronic coefficients remain open.


\section{14. Baseline Metric Closure: No Slip and PPN}

For the branch that reproduces Newtonian gravity, the physical metric is already fixed by the Einstein parent action. The scalar constraint reduction of Section~10 gives
\[
\nabla^2(\Phi-\Psi)=0,
\]
so asymptotic flatness implies
\[
\Phi=\Psi.
\]
This is an action-level result. It is not inferred from the canonical stress tensor of $\delta S$, because no independently gravitating capacity scalar is present. Since the parent action of this branch is exactly Einstein--Hilbert plus minimally coupled matter, its vacuum post-Newtonian solution has
\[
\gamma_{\mathrm{PPN}}=\beta_{\mathrm{PPN}}=1
\]
and the remaining standard PPN parameters vanish, subject to the usual assumptions on the matter sector and boundary conditions. The capacity redefinition leaves these values unchanged. The carrier rule introduced below also removes localized Solar-System sources from the transverse influence functional. The leading Solar-System metric is therefore the Einstein solution in the specified carrier-resolved EFT. A nonlinear covariant audit of the record projector and the cosmic-state vector remains required for time-dependent preferred-frame terms.

\section*{Part IV. The Carrier-Resolved Galactic Branch}
\addcontentsline{toc}{section}{Part IV. The Carrier-Resolved Galactic Branch}

\section{15. Carrier-Resolved Galactic Dynamics}

The Einstein branch responds to the full stress tensor and supplies ordinary gravity at every scale. The galactic excess is a second, coarse-grained response associated with retained source records. Its leading effective completion consists of a carrier projector, a thermal contact law, and a conservative auxiliary-field action. These ingredients fix the static equations without introducing an independently coupled matter charge. Their microscopic origin in the GFT state and metric vertex remains conditional.


\paragraph{Carrier nesting.}
An active carrier $P$ is a retained, dynamically coherent source record at the resolution relevant to the long-range response. On the algebra of decohered source records, let $e_P$ be its sharp selector and $E_P$ the source-preserving conditional expectation onto its retained subalgebra. The selector belongs to that subalgebra. For an active antichain covering the modeled records,
\[
e_P^2=e_P,\qquad e_Pe_Q=0\quad(P\ne Q),\qquad
\sum_{P\in\mathfrak A}e_P=1,
\qquad \boxed{\Pi_PX=E_P(e_PX).}
\]
The conditional-expectation bimodule property and $E_P^2=E_P$ give
\[
\Pi_P^2=\Pi_P,\qquad \Pi_P\Pi_Q=0\quad(P\ne Q).
\]
Each record enters one active carrier. On a common parent record space, a finer conditional expectation $E_{\rm f}$ satisfies $E_PE_{\rm f}=E_P$. If the child selectors partition the parent and their retained algebras refine its algebra, then
\[
\sum_{A\prec P}e_A=e_P,\qquad
\Pi_PX=E_P\!\left(\sum_{A\prec P}e_AX\right)
=E_P\!\left(\sum_{A\prec P}\Pi_AX\right).
\]

Spatial sources are reconstructed after record selection. Write
\[
T_P^{\rm cg}=\mathcal C_\ell\Pi_PT,\qquad
T^{\mathfrak A}=\sum_{P\in\mathfrak A}T_P^{\rm cg},
\]
where the linear reconstruction $\mathcal C_\ell$ preserves the integrated source at the chosen coarse resolution. Its scalar realization and the parent-resolution compatibility condition are given in Appendix~N.8. Reconstructed carriers can overlap in space. Smooth weights $\chi_P\ge0$, with $\sum_P\chi_P=1$ in the response domain, interpolate their constitutive energies. The exact projection identities act on the record selectors $e_P$; the spatial weights enter the action once and carry no idempotence requirement.

The transverse functional consequently has the carrier sum
\[
\Gamma_\perp^{\rm nest}
=\sum_{P\in\mathfrak A}w_P\,
\Gamma_{\perp,w_P}[g_\pm;T^{\mathfrak A}_\pm,\chi_P,\mathcal R_P],
\qquad
T^{\mathfrak A}=\sum_{Q\in\mathfrak A}T_Q^{\rm cg}.
\]
Each carrier's share of the response is evaluated on the joint coarse source $T^{\mathfrak A}$ and weighted by its spatial interpolation profile; for a single active carrier $T^{\mathfrak A}=T_P^{\rm cg}$.
The Einstein action continues to use the unprojected $T_{\mu\nu}$. A star contributes once to its galaxy's coarse transverse source and retains its resolved stellar field in the Einstein branch. It does not also carry an independent transverse halo. This source-space nesting removes parent--descendant cross-terms before the nonlinear response is evaluated. The joint source makes the forces between active carriers reciprocal. A response computed from each carrier's own field alone would not be: the excess force of $P$ on $Q$ scales as $M_Q\sqrt{M_P}$ and that of $Q$ on $P$ as $M_P\sqrt{M_Q}$, and a Milky-Way--LMC pair would self-accelerate at about $10^2\ {\rm km\,s^{-1}\,Gyr^{-1}}$. On the joint source the excess forces cancel identically (paragraph ``Several active carriers'' below).

\paragraph{Thermal contact and the acceleration scale.}
The response cell of carrier $P$ is the tilt mode of its transverse frame, with the baryonic field as its restoring stiffness. Its inertia lever is the apparent-horizon radius multiplied by the sharing entropy of the cell,
\[
\ell=g_{\mathrm{share,eff}}\,\frac{c}{H}.
\]
For the potential $Mg\ell(1-\cos\theta)$ in a field $g$ and the kinetic term $\tfrac12M\ell^2\dot\theta^2$, small oscillations have $\omega^2=g/\ell$ for every mode mass $M$, so the mode frequency is $\omega_P=\sqrt{|\boldsymbol\nabla\Phi_{b,P}|/\ell}$. Thermalized in the apparent-horizon bath at $k_BT_H=\hbar H/2\pi$, the mode has a gap in thermal units equal to a ratio of two frequencies, in which $\hbar$ and $k_B$ cancel:
\[
x_P=\frac{\hbar\omega_P}{k_BT_H}=\frac{2\pi\omega_P}{H}=\sqrt{\frac{|\boldsymbol\nabla\Phi_{b,P}|}{a_0(z)}},
\qquad
a_0(z)=\frac{g_{\mathrm{share,eff}}}{4\pi^2}\,cH(z),
\qquad
a_0\equiv a_0(0).
\]
For a lever measured in apparent-horizon radii the factor $4\pi^2$ is the square of the $2\pi$ in the horizon temperature. The mode fixes the square root in $x_P$, the transition $x_P=1$ at $|\boldsymbol\nabla\Phi_{b,P}|=a_0$, and the redshift dependence $a_0(z)\propto H(z)$. Evaluated in the field component along a frame tilted by $|\delta|\le\pi/2$ from the source direction, its gap is $x_P\sqrt{\cos\delta}$, the director-completion gap $\widetilde x_P$ used below. The pendular frequency and the algebra of these relations are kernel checked in \path{Entropic/RAR/Pendular.lean}. The matching inputs of the acceleration scale are the identification of the cell with this mode, its thermalization in the horizon bath, and the lever. The lever coefficient follows from the thermodynamics of the cell under one further premise. The register of Appendix~B has entropy $g_{\rm share,eff}=\eta_*\langle K^2\rangle+\ln Z$ with $\eta_*\langle K^2\rangle=3/2$. Cells are renewed, so their chemical potential vanishes and the enthalpy of a cell is $g_{\rm share,eff}k_BT$; carried at the horizon arm $c/H$, it is the inertia of the tilt mode. A static frame has proper acceleration equal to the field it holds against and reads the Unruh temperature $k_BT_U=\hbar g\cos\delta/2\pi c$ of the projected baryonic field. An oriented record of entropy $s$ at that temperature has free energy $-s\,k_BT_U$, which pulls the frame toward the field with weight $s\,\hbar g/2\pi c$. The three closure components share $\eta_*\langle K^2\rangle$ equally, so the transverse pair that the tilt rotates carries exactly one nat, and $\ln Z$ carries no orientation. With $s=1$ the lever is $\ell=g_{\rm share,eff}c/H$ and $x_P^2=(2\pi)^2T_U/(g_{\rm share,eff}T_H)$. The RAR scale requires $0.85<s<1.39$, which excludes one bit and contains one nat. The premise is the Unruh reading of the projected baryonic field by the static frame, with the oriented record identified as the transverse share. The response reads the baryonic source by the nesting rule; the proper acceleration of a static observer also includes the excess field, and the premise assigns the reading to the baryonic part. The algebra from these premises to the lever and the gap is kernel checked in the same module. For a lever of $\kappa$ apparent-horizon radii, $\kappa$ is the ratio of the effective inertia of the tilt mode to its coupling to the source field. The mode is a collective excitation of the frame and carries no committed defect, so ordinary matter keeps equal inertial and gravitational mass and test-body free fall is unchanged. At $\kappa=1$ the scale is $a_0=cH_0/4\pi^2=1.66\times10^{-11}\,\mathrm{m\,s^{-2}}$, below the RAR scale by the factor $g_{\mathrm{share,eff}}$. The local frame tilt carries no occupation. The relativistic completion of the director stiffness given below, against its restoring density, gives the local tilt the frequency $c/\xi_{\theta,P}$, and for the Milky-Way coherence length of $6\,\mathrm{kpc}$ its gap is $2\pi c/(H_0\xi_\theta)\simeq5\times10^6$. The inertia of the cell is therefore of horizon scale. Near $a_0$ the mode frequency is about $H/2\pi$, a period of about $4\pi^2/H$, and in the deep regime it lies below the Hubble rate; a quasi-stationary thermal occupation of so slow a mode in an expanding background is a further open condition. The occupation is a property of the equilibrium state and requires no completed oscillation. In a stationary de Sitter bath a mode of any gap relaxes to the Gibbons--Hawking occupation at a rate set by its coupling \citep{GibbonsHawking1977}, and the stock contacts of Appendix~O.3 supply that coupling. In the deep regime $H$ drifts by order unity during one period of the mode, and the resulting departure from the stationary occupation would appear as a lag of $a_0(z)$ behind $H(z)$. At the Hubble constant of the standard linear-depletion release fit of Section~20, $H_0=68.46\,\mathrm{km\,s^{-1}\,Mpc^{-1}}$, the mode gives $a_0=1.250\times10^{-10}\,\mathrm{m\,s^{-2}}$, and at the Planck value $1.231\times10^{-10}\,\mathrm{m\,s^{-2}}$, compared with the fitted RAR scale $1.20\times10^{-10}\,\mathrm{m\,s^{-2}}$ \citep{McGaugh2016,Planck2020}. The frozen release history of Appendix~O.12, which contains no fitted dark-energy parameter, gives $H_0=68.42\,\mathrm{km\,s^{-1}\,Mpc^{-1}}$ and $a_0=1.249\times10^{-10}\,\mathrm{m\,s^{-2}}$, and the astrophysical chain without cosmological data gives $a_0=(1.27\pm0.04)\times10^{-10}\,\mathrm{m\,s^{-2}}$. Changing from the Planck normalization to this standard linear-depletion fit changes deep-regime accelerations by $0.8\%$ and moves the 22-dwarf $\chi^2$ of Section~15 by $0.1$. The archived old-disk calculations use the earlier exponential fit, $H_0=68.93885586416\,\mathrm{km\,s^{-1}\,Mpc^{-1}}$ and $a_0=1.25882886\times10^{-10}\,\mathrm{m\,s^{-2}}$. \citet{Milgrom1999} related the MOND acceleration scale to the de Sitter horizon through the Unruh temperature of accelerated observers, with a different interpolation; the tilt mode gives the exponential interpolation and the coefficient $g_{\mathrm{share,eff}}/4\pi^2$. The extrapolation $a_0(z)\propto H(z)$ is tested separately in Section~27.3. The same $\epsilon$ fixes the pre-release abundance $\Omega_c/\Omega_b=1/\epsilon$ of Section~20, so the two scales obey the parameter-free relation $a_0\,\Omega_c/\Omega_b=cH_0$.

The cell Hamiltonian of an active carrier $P$ is that of the mode,
\[
H_{{\rm cell},P}=\hbar\omega_P:N_P:=k_BT_H\,x_P:N_P:.
\]
Its equilibrium free energy and response are
\[
F_{{\rm th},P}=k_BT_H\ln(1-e^{-x_P}),
\qquad
\frac{\partial F_{{\rm th},P}}{\partial E_P}
=n_B(x_P)=\frac{1}{e^{x_P}-1}.
\]
The mode energy $E_P=k_BT_Hx_P$ follows from the tilt mode. In the deep regime $n_B\simeq1/x_P$ is the equipartition occupation of the mode, which gives $g_{\rm obs}\to\sqrt{a_0g_{\rm bar}}$, and at high acceleration the Wien suppression of the occupation returns the Newtonian field. The rank-one capacity Hessian of Appendix~N remains a thermodynamic consistency condition for one capacity variable. Its massless tangent and massive eigenvalue $A_L+A_T$ do not supply the thermal energy or an additional propagating exchange mode.

\paragraph{Conservative static action.}
Define the carrier density and its Newtonian auxiliary potential by
\[
\rho_P=\frac{u_\mu u_\nu}{c^2}T_P^{{\rm cg},\mu\nu},
\qquad
\nabla^2\Phi_{b,P}=4\pi G\rho_P,
\]
and set
\[
z_P=\frac{|\boldsymbol\nabla\Phi_{b,P}|^2}{a_0^2},
\qquad
x_P=z_P^{1/4}.
\]
The required excess function is
\[
\mathcal Q_\perp(z)
=\int_0^z\frac{ds}{e^{s^{1/4}}-1},
\qquad
\mathcal Q_\perp'(z)=n_B(z^{1/4}),
\]
so the total auxiliary function $\mathcal Q(z)=z+\mathcal Q_\perp(z)$ satisfies
\[
\mathcal Q'(z)=1+n_B(z^{1/4})
=\frac{1}{1-e^{-z^{1/4}}}.
\]
For each active carrier, introduce a total carrier potential $\Phi_P$ and the QUMOND-type action \citep{Milgrom2010}
\[
I_{Q,P}=-\int dt\,d^3x\left\{
\frac{1}{8\pi G}
\left[2\boldsymbol\nabla\Phi_P\!\cdot\!\boldsymbol\nabla\Phi_{b,P}
-a_0^2\mathcal Q(z_P)\right]
+\rho_P\Phi_P
\right\}.
\]
The functional added to the Einstein branch is the excess $I_{\perp,P}=I_{Q,P}-I_{N,P}$, where $I_{N,P}$ is the same auxiliary action with $\mathcal Q(z)=z$. Variation gives
\[
\nabla^2\Phi_P
=\boldsymbol\nabla\!\cdot\!
\left[\mathcal Q'(z_P)\boldsymbol\nabla\Phi_{b,P}\right].
\]
Writing $\phi_{\perp,P}=\Phi_P-\Phi_{b,P}$ isolates the transverse excess:
\[
\boxed{
\nabla^2\phi_{\perp,P}
=\boldsymbol\nabla\!\cdot\!
\left[n_B(x_P)\boldsymbol\nabla\Phi_{b,P}\right].
}
\]
This equation is conservative and applies to nonspherical sources. When active carriers overlap, $\Phi_{b,P}$ in the flux is replaced by the joint potential $\Psi=\sum_{Q\in\mathfrak A}\Phi_{b,Q}$ and carrier $P$'s flux is weighted by $\chi_P$; an isolated carrier is unchanged. For spherical systems the unprojected contact reduces to
\[
g_{\rm obs}
=g_{\rm bar}\left[1+n_B\!\left(\sqrt{g_{\rm bar}/a_0}\right)\right]
=\frac{g_{\rm bar}}
{1-\exp[-\sqrt{g_{\rm bar}/a_0}]}.
\]
As the force of a nonnegative retained stock, this relation is exact on profiles whose response density is nonnegative throughout the carrier and whose mass fits the retained budget, before the outer cap becomes active. High-acceleration cores can violate this admissibility condition even in weak gravity. Appendix~O.3 specifies a positive-part target normalized to the same stock mass and bounds its spherical force correction. Local positivity at one radius alone does not imply exactness there, because the enclosed mass also includes the interior projection.

For the unprojected spherical law and $a_0,g_{\rm bar}>0$,
\[
\max\!\left(g_{\rm bar},\sqrt{a_0g_{\rm bar}}+\frac{g_{\rm bar}}2\right)
\le g_{\rm obs}\le g_{\rm bar}+\sqrt{a_0g_{\rm bar}}.
\]
The response increases strictly with $g_{\rm bar}$. Its high-acceleration limit is Newtonian, and $g_{\rm obs}/\sqrt{a_0g_{\rm bar}}\to1$ in the deep regime. For every spherical profile with $M_b(r)\to M>0$, this law gives $v^4\to GMa_0$ \citep{TullyFisher1977,McGaugh2000}. These inequalities and limits are kernel checked; a finite retained-stock cap limits the physical radial range of the uncapped law. Disks and other nonspherical sources generally carry geometry-dependent curl corrections. Their predictions require direct Poisson solves using resolved baryonic maps; the pointwise spherical formula is not imposed on each pixel.

\paragraph{Derivation of the radial-acceleration relation.}
The tilt mode in the horizon bath fixes the gap $x_P$ and the scale $a_0$. With its Bose occupation read as the excess per unit baryonic field, the auxiliary action turns that occupation into the conservative field equation. The conditional stock contacts of Appendix~O.3, biased by the Bose relative entropy of the cell, relax a retained stock onto the resulting equilibrium, and they set the relaxation rate. Given these inputs, the exponential law, its bounds, the limit $v^4\to GMa_0$, and $a_0=\epsilon cH_0$ follow with no coefficient fitted to galaxies. The lever $\ell=g_{\mathrm{share,eff}}c/H$ is the ratio of the register enthalpy to the one-nat oriented record read at the Unruh temperature of the projected baryonic field, and that reading is its premise.

\paragraph{Direction-stiffness completion for resolved disks.}
The preceding local contact slaves the transverse frame to the source direction
$\mathbf e_{L,P}=\boldsymbol\nabla\Phi_{b,P}/|\boldsymbol\nabla\Phi_{b,P}|$.
At a thin-disk midplane that direction reverses across roughly two scale heights, and the local divergence acquires a narrow phantom sheet with surface density of order $n_B\Sigma_b$. The sheet multiplies the vertical self-gravity of the disk by $1+n_B$, with $n_B\simeq0.5$ at the solar radius, and produces the failed local Milky-Way benchmark below.

The selected resummation promotes the source direction to a nematic unit director $\hat{\mathbf e}_P\equiv-\hat{\mathbf e}_P$ and defines
\[
\widetilde z_P=\frac{(\hat{\mathbf e}_P\!\cdot\!\boldsymbol\nabla\Phi_{b,P})^2}{a_0^2},
\qquad
\widetilde x_P=\widetilde z_P^{1/4},
\qquad
\mathbf J_{\perp,P}=n_B(\widetilde x_P)
(\hat{\mathbf e}_P\!\cdot\!\boldsymbol\nabla\Phi_{b,P})\hat{\mathbf e}_P,
\]
with
\[
\boxed{\nabla^2\phi_{\perp,P}=\boldsymbol\nabla\!\cdot\mathbf J_{\perp,P}.}
\]
The tetrahedral kinetic term contains $Z_0X^2(\partial\hat{\mathbf e})^2$ (Appendix~N.8). Identifying its longitudinal amplitude with the carrier record potential $X_P$ gives the stiffness $Z_{\theta,P}=2c\,X_P^2$. The record potential is the potential of the carrier's response stock, $X_P\leftrightarrow\Phi_{R,P}$ (Appendix~O.3), which vanishes at the carrier vacuum. For an isolated carrier in equilibrium $\Phi_{R,P}$ is a fixed functional of $\Phi_{b,P}$. The baryonic potential then serves as a proxy, $Z_{\theta,P}=2c\,\Phi_{b,P}^2$, with the same stiff-limit solutions. Overlapping carriers require the stock potential itself (paragraph ``Several active carriers'' below). The tetrahedral term does not fix the normalization $c$. The Milky Way and satellite dispersions place it near $c\simeq0.03$ for the stock amplitude, and the resolved solutions are insensitive to it within that range. At constant stiffness the one-constant density $|\partial\hat{\mathbf e}|^2\equiv\partial_i\hat e^a\partial_i\hat e^a$ cannot be distinguished from $|\partial\hat{\mathbf e}|^2+\kappa\,s[\hat{\mathbf e}]$, where
\[
s[\hat{\mathbf e}]
=\partial_i\hat e^a\partial_a\hat e^i-(\boldsymbol\nabla\!\cdot\!\hat{\mathbf e})^2
=\partial_i\!\left[(\hat{\mathbf e}\!\cdot\!\boldsymbol\nabla)\hat e^i-\hat e^i\,\boldsymbol\nabla\!\cdot\!\hat{\mathbf e}\right].
\]
Once $Z_{\theta,P}$ follows the record potential, the $\kappa$ term integrates to the connection cross term $-\tfrac12\kappa\,\partial_iZ_{\theta,P}[(\hat{\mathbf e}\!\cdot\!\boldsymbol\nabla)\hat e^i-\hat e^i\,\boldsymbol\nabla\!\cdot\!\hat{\mathbf e}]$, one of the completion data named in Appendix~N.8. Requiring the director completion to preserve the unprojected spherical contact fixes $\kappa$. The radial field has $|\partial\hat{\mathbf r}|^2=2/r^2$ and $s[\hat{\mathbf r}]=-2/r^2$, so $|\partial\hat{\mathbf e}|^2+\kappa\,s[\hat{\mathbf e}]$ vanishes on every radial director only for $\kappa=1$:
\[
\boxed{
\mathcal S[\hat{\mathbf e}]
=|\partial\hat{\mathbf e}|^2+s[\hat{\mathbf e}]
=(\lambda_1-\lambda_2)^2+|(\hat{\mathbf e}\!\cdot\!\boldsymbol\nabla)\hat{\mathbf e}|^2\ \ge0 ,}
\]
where $\lambda_{1,2}$ are the principal splays of the director congruence. The density penalizes anisotropic splay and bend and vanishes for radial and uniform directors. Because $s$ has zero principal symbol, the director equation keeps the ellipticity of the one-constant term.

Expanding $-a_0^2\mathcal Q_\perp(\widetilde z_P)$ to second order in the tilt $\delta=\alpha-\alpha_L$ from the source direction gives the restoring density $n_B|\boldsymbol\nabla\Phi_{b,P}|^2\delta^2$. Against the stiffness $c\,X_P^2|\boldsymbol\nabla\delta|^2$ this sets the coherence length
\[
\boxed{\xi_{\theta,P}^2=\frac{c\,X_P^2}{n_B(\widetilde x_P)\,|\boldsymbol\nabla\Phi_{b,P}|^2}.}
\]
For the Milky Way $\Phi_R/\Phi_b\simeq2.8$ at the solar radius. The stock amplitude at $c=0.03$ therefore gives $\xi_\theta(R_0,0)\simeq6\,\mathrm{kpc}$, twenty times the stellar scale height, and the baryonic proxy at $c=1$ gives $12\,\mathrm{kpc}$. Either way the director cannot follow the source reversal across the disk thickness. Spherical sources carry the radial director, $\mathcal S$ vanishes, and the unprojected local contact is preserved. Its positive-stock realization obeys the admissibility condition and projection of Appendix~O.3. In a disk the near-radial director passes through the midplane without reversal. Near $z=0$, $\hat{\mathbf e}\simeq\hat{\mathbf R}+(z/R)\hat{\mathbf z}$ and $\partial_zJ_{\perp,z}\simeq n_Bg_R/R$, where the local contact has $\partial_z(n_B\partial_z\Phi_b)\simeq4\pi Gn_B\rho_b$. The ratio $g_R/(4\pi G\rho_bR)$ is $0.03$--$0.1$ in stellar disks, so the excess is a smooth, nearly round component without a phantom sheet. An isotropic scalar smoothing kernel wide enough to remove the Milky-Way sheet shifts dwarf and inner-disk SPARC accelerations by $0.03$--$0.10$ dex, which the RAR data exclude. One auxiliary action generates the Newtonian Poisson equation, the excess flux, the director equation, and a stiffness source $c\,X_P\mathcal S[\hat{\mathbf e}_P]$ in the excess equation (Appendix~N.8). The source vanishes for every spherical carrier. The director completion was formulated after the failure of the local-contact Milky-Way vertical benchmark was known. The comparisons of Section~16 are therefore development-set postdictions unless they are identified as subsequently frozen tests; the old-disk dispersions and excess shapes registered in Section~27.1 are such tests.

\paragraph{Several active carriers.}
Let the active set contain carriers whose fields overlap, such as a galaxy and a satellite that keeps an independent record. With $\Psi=\sum_P\Phi_{b,P}$ and the local contact $\mathbf J=n_B\bigl(\sqrt{|\boldsymbol\nabla\Psi|/a_0}\bigr)\boldsymbol\nabla\Psi$, two integrations by parts with $\rho_i=\nabla^2\Phi_{b,i}/4\pi G$ give the excess force on carrier $i$,
\[
F_i^\perp=-\int\rho_i\,\partial_z\phi_\perp\,d^3x
=-\frac{1}{4\pi G}\int\mathbf J\cdot\boldsymbol\nabla\partial_z\Phi_{b,i}\,d^3x .
\]
Summed over carriers the integrand becomes $n_B\,\partial_z\tfrac12|\boldsymbol\nabla\Psi|^2$, a total derivative, so
\[
\boxed{\sum_{i\in\mathfrak A}F_i^\perp=0 .}
\]
For a Milky-Way-like Plummer carrier ($6\times10^{10}M_\odot$, $b=3\,$kpc) and an LMC-like one ($3\times10^9M_\odot$, $b=1\,$kpc) $50\,$kpc apart, the joint source gives excess forces $\pm1.471\times10^{12}\,M_\odot({\rm km\,s^{-1}})^2{\rm kpc^{-1}}$ with zero net. Per-carrier responses give $+8.14$ and $-1.72\times10^{12}$, a centre-of-mass self-acceleration of $104\ {\rm km\,s^{-1}\,Gyr^{-1}}$. An independent implementation reproduced both results.

The director completion carries over with one director per carrier, a spatial interpolation weight $\chi_P$ ($\sum_P\chi_P=1$), and the joint field in $\widetilde z_P$. Translating each carrier with its support, stiffness amplitude, and director core leaves the static energy unchanged, so the forces balance on-shell. For the Milky-Way--LMC pair, the on-shell solution balances to $2\times10^{-6}$ of the mutual force. With the baryonic amplitude $Z_{\theta,P}=2c\Phi_{b,P}^2$ it has two defects:
\begin{itemize}
\item a negative excess-density layer of $-7.5\times10^9M_\odot$, two and a half times the LMC's baryons, on the support edge;
\item accelerations of 0.4--2.8\% that act on LMC-labelled but not MW-labelled matter at the same point.
\end{itemize}
A single continuous director removes the layer only by locking the satellite or by creating a compensating defect. Appendix~N.8 shows that two changes remove both defects:
\begin{itemize}
\item the nematic order melts where carrier records overlap;
\item the stiffness amplitude becomes the potential of each carrier's response stock (Appendix~O.3).
\end{itemize}
With an overlap zone half as wide as the support and normalization $c\simeq0.03$, the combined model gives the following:
\begin{itemize}
\item an interface layer of $-0.85\times10^9M_\odot$, against $-0.33\times10^9$ for a single joint field;
\item an intact satellite interior;
\item a mutual force within the finite-difference systematic of the joint-field value;
\item no label-dependent force on baryons;
\item a Milky-Way vertical potential of $1416$ and $1315\ {\rm km^2\,s^{-2}}$ for the two baryonic models of Section~16.
\end{itemize}

The joint source fixes momentum balance but not inertia. The Milky Way's reflex velocity toward the LMC, integrated along a radial infall from $250$ to $50\,$kpc at $300\ {\rm km\,s^{-1}}$, is $8.7\ {\rm km\,s^{-1}}$ when the response has no inertia of its own. The measured travel velocity is $32\pm4\ {\rm km\,s^{-1}}$ at $40$--$120\,$kpc \citep{PetersenPenarrubia2021}, and about $40$ near $100\,$kpc \citep{Chandra2024}. The response must therefore be carried by an inertial stock that moves with its carrier. Appendix~O.3 identifies that stock with the reallocated capacity and fixes its budget.

\paragraph{Empirical domain.}
The exponential law is specified for equilibrium active carriers in the regime established by rotationally supported galaxies and isolated-galaxy lensing. Dispersion-supported low-mass dwarfs provide a distinct boundary test. A GravSphere analysis of 12 nearby dwarfs with $10^4<M_{\rm bar}/M_\odot<10^{7.5}$ finds that most lie above the extrapolated RAR, trace multivalued loci, and show substantial galaxy-to-galaxy scatter; the reported tidal and nonequilibrium checks do not account for the result \citep{Julio2025}. The carrier projector could address this evidence only through a derived phase-selection prediction. The present EFT therefore makes no universal RAR claim for ultrafaint dwarfs.

Satellites and tidal dwarfs test the carrier rules directly. The support rule of Section~18 assigns a satellite its own cells only where its own baryonic field exceeds the host's. A satellite outside that condition keeps its pre-infall stock, which evolves collisionlessly through transport and stripping. A satellite inside it relaxes like an isolated dwarf. For the eight classical Milky-Way dwarfs with GravSphere central densities \citep{Read2019}, the ratio of the Milky-Way law field to the dwarf's own law field at its three-dimensional half-light radius is $5.5$ for Sextans, $2.5$ for Draco, $2.0$ for Ursa Minor and $1.7$ for Carina, which keep their stock, and $0.98$ for Sculptor, $0.38$ for Fornax, $0.20$ for Leo~II and $0.09$ for Leo~I, which relax. Scored on $\rho(150\,{\rm pc})$, this assignment gives $\chi^2=10.4$ over the eight dwarfs. Retained stock in all eight gives $17.2$, with Fornax contributing $11.8$, and relaxed stock in all eight gives $44$. The classification is unchanged for Milky-Way baryonic masses of $5$--$6\times10^{10}M_\odot$, cap ratios of $10$--$20$, projected or three-dimensional radii, and Newtonian or law fields; Sculptor lies on the boundary. Ten of the $256$ retained or relaxed assignments reach $\chi^2\le10.4$, and the best, $5.7$, also retains Leo~II, so these data support the rule at a modest level. A dwarf whose main progenitor never crossed the release gate of Section~20 keeps its allocation in the saturated source-associated phase, and its stock stays collisionless whatever its support. Across 22 Local Group dwarfs with half-light masses from the Local Volume Database, relaxed masses from the contact runs of Appendix~O.3, and gate crossing from the mean accretion history of \citet{Fakhouri2010}, this combined assignment gives $\chi^2=21.7$ at the reference gate of $17\,{\rm km\,s^{-1}}$ and $25.4$ at $14\,{\rm km\,s^{-1}}$. $\Lambda$CDM halos with supernova heating at the published coreNFW calibration give $24.0$, and NFW halos without heating give $30.0$. The two accounts describe these dwarf masses to the same degree. Over plausible accretion histories, M31 baryonic masses and the two published Tucana dispersions, the combined assignment spans $\chi^2=18.5$--$29.2$ and heating spans $23$--$25.5$, so these data do not fix the gate. Tucana and Cetus, isolated and above the gate, have the low relaxed masses the support rule predicts; Eridanus~II and Canes Venatici~II, below it, are as dense as CDM halos. The current Crater II result uses the retained history; the carrier-local stiffness and tidal-dwarf comparisons below test their stated equilibrium configurations.
\begin{itemize}
\item \textbf{Crater II.} The moving-LMC tidal calculation summarized in \emph{Empirical Results} gives a retained-stock state with $\sigma_{\rm los}=2.36\,{\rm km\,s^{-1}}$ and $R_h=0.86\,{\rm kpc}$, compared with the measured dispersion $2.7\pm0.3\,{\rm km\,s^{-1}}$ \citep{Caldwell2017}. The result is conditional on its specified progenitor, orbit, and force model.
\item \textbf{Carrier-local stiffness.} A stiffness built from the joint potential would lock satellite directors to the host. It would reduce their dynamical masses by factors of 2.4--3.9, and Fornax would fall to $7.2\ {\rm km\,s^{-1}}$. The stiffness amplitude is therefore carrier-local.
\item \textbf{Tidal dwarfs.} Descendants that inherit their parent's record sit under the parent's director and are locked. For the NGC~5291 tidal dwarfs, with the parent at $1$--$2\times10^{11}M_\odot$ and $60\,$kpc, the locked response gives dynamical-mass boosts of $1.5$--$2.0$, against the measured $1.2$--$1.5$ \citep{Lelli2015}. Treating each dwarf as its own carrier gives $3.6$--$4.9$. For a $2\times10^{11}M_\odot$ parent, $\chi^2$ over the three dwarfs is $0.5$ for inheritance and $48$ for separate carriers.
\end{itemize}

\paragraph{Retained-stock comparisons.}
The empirical companion reports LMC enclosed masses of
$5.0348\times10^{10}M_\odot$ at 30 kpc and
$8.2896\times10^{10}M_\odot$ at 50 kpc for its stated baryonic profile
and $C_R=33$. Its Crater II comparison favors the stock over the specified
NFW population, with $B=2.62\times10^5$ at 0.5-dex halo scatter and 10\%
model error; the preference changes substantially with that error model.
The same moving-LMC campaign leaves Antlia 2 unresolved: its best bound
stock remnant has dispersion $3.48\,{\rm km\,s^{-1}}$, below the measured
value near $6\,{\rm km\,s^{-1}}$. These saved runs use collisionless
stock with no subsequent contact relaxation, as the support rule assigns
to external-field-dominated satellites, and a specified approximate
force model. They do not yet implement the complete reciprocal
carrier/director dynamics or the final observational selection.

\paragraph{Status.}
The thermal contact free energy, local auxiliary action, and carrier nesting rule close the leading slowly varying static EFT once its matching data are supplied. Evaluated on the joint coarse source, the carrier sum conserves momentum between active carriers. The direction-resummed response follows from the same auxiliary action with the record-potential stiffness. Several inputs of its multi-carrier completion remain conditional: the overlap width, the normalization $c$, the tie of supports and director cores to the response stocks, and the stock dynamics. The amplitude identification $X_P\leftrightarrow\Phi_{R,P}$, the saddle-splay coefficient set by exact preservation of the unprojected spherical contact, the tilt-mode identification with its horizon thermalization and lever, active-carrier projector, and potential-zero prescription remain conditional. A microscopic GFT derivation must produce the cell Hamiltonian, projector, metric vertex, and director stiffness. The ultrafaint-dwarf phase boundary remains empirical work. Appendix~O.3 supplies a conditional selection dynamics for the response stock. Momentum-conserving stock contacts are biased by the Bose relative entropy of the response cell, in the zero-temperature limit of local detailed balance with the horizon bath. The cells are carried only by active carriers, so cluster stock evolves as collisionless matter (Section~18).

\section{16. Galactic Metric, Lensing, and Local Tests}

\paragraph{No-slip metric contact.}
The leading stationary metric response of carrier $P$ is
\[
\boxed{
h_{\mu\nu}^{\perp,P}
=-\frac{2\phi_{\perp,P}}{c^2}
\left(\bar g_{\mu\nu}+2u_\mu u_\nu\right).
}
\]
In Newtonian gauge this gives $\Delta\Phi_{\perp,P}=\Delta\Psi_{\perp,P}=\phi_{\perp,P}$. A leading closed-time-path contact that produces this mean response is
\[
\boxed{
\Gamma_\perp^{\rm QS}
=-\frac12\sum_{P\in\mathfrak A}w_P
\left\langle h_\Delta,
\mathcal E h_{\perp,P}[T^{\mathfrak A}_c]\right\rangle ,
}
\]
where $\mathcal E$ is the linearized Einstein operator and the bracket includes the invariant spacetime integral. This contact satisfies $\Gamma_\perp[g,g]=0$ and is real in the static mean. Its induced stationary source has $T^{00}_{\perp,P}=\rho_{\perp,P}c^2$ and $T^{ij}_{\perp,P}=0$, so it carries no scalar anisotropic stress. The quasistatic contact introduces no independent scalar Cauchy data and adds no pole to the vacuum graviton propagator. A canonical local scalar-gradient stress would produce anisotropic stress and is not the metric realization used here.

The linear Ward identity follows from the Bianchi identity for $\mathcal E$. Nonlinear conservation requires the record variables, reconstruction, and spatial weights to be varied with the metric. In covariant form the combined source map must obey
\[
(\mathcal C_\ell\Pi_P)[\varphi^*g,\varphi^*\psi,\varphi^*\mathcal R]
=\varphi^*(\mathcal C_\ell\Pi_P)[g,\psi,\mathcal R].
\]
The reconstruction and interpolation profiles inherit their metric and carrier dependence. A covariant realization must include these variations in its nonlinear Ward identity, which remains to be shown. The stationary force calculation includes the corresponding spatial variations.

\paragraph{Lensing response.}
No additional lensing coefficient appears. The effective density inferred by dynamics and lensing is
\[
\rho_{{\rm lens},P}
=\rho_P+\frac{1}{4\pi G}
\boldsymbol\nabla\!\cdot\mathbf J_{\perp,P}.
\]
For a deep-regime point mass, define $A=\sqrt{GMa_0}$. Then
\[
g=\frac{A}{r},
\qquad
\rho_\perp(r)=\frac{A}{4\pi Gr^2},
\]
and the projected profiles and asymptotic deflection are
\[
\Sigma_\perp(R)=\Delta\Sigma_\perp(R)=\frac{A}{4GR},
\qquad
\alpha_\infty=\frac{2\pi\sqrt{GMa_0}}{c^2}
=\frac{2\pi v_\infty^2}{c^2}.
\]
The response of an isolated carrier is bounded by its stock (Appendix~O.3). The isolated-lens calculation uses the linear-depletion benchmark for the carrier's Lagrangian baryons,
\[
M_\perp\le C_RM_b,\qquad C_R=(1-\eta_*)(1-f_b)\frac{M_h}{M_b},
\]
and the cap is reached at $r_C=r_M/\ln(1+1/C_R)$. The composition-preserving allocation gives the ensemble full-release factor $e^{-\eta_*}$ in place of $1-\eta_*$, a $0.0455\%$ budget increase; the likelihood values below refer to the displayed benchmark. The isolated-lens KiDS-1000 RAR \citep{Brouwer2021} tests this bound. For a point baryonic mass the prediction at fixed $g_{\rm bar}$ is independent of $M_b$, and the model surface density converts to $g_{\rm obs}=4G\Delta\Sigma$ as for the data.
\begin{itemize}
\item \textbf{Cap per present baryon.} $C_R=(1-\eta_*)/\epsilon=5.16$ puts the cap at the innermost lensing point. It gives $\chi^2=326$ on the seven isolation-safe points ($g_{\rm bar}\ge10^{-13}\,{\rm m\,s^{-2}}$), against $40$ uncapped, and is excluded. An independent implementation reproduced both values.
\item \textbf{Lagrangian cap with the whole budget carried.} With $M_h$ from the stellar-to-halo relation of \citet{Moster2013} at $z=0.2$, the four stellar-mass bins give $C_R=20$--$49$ and $r_C=21$--$50\,r_M$, near the virial radius. Over 31 mass-binned isolation-safe points this gives $\chi^2=86$, against $75$ uncapped and $69$ for halo masses $0.2\,$dex higher. Beyond $r_C$ this configuration is Keplerian. With an isolation radius of $4\,$Mpc the lensing signal continues the relation to $g_{\rm bar}\simeq10^{-14}\,{\rm m\,s^{-2}}$ \citep{Mistele2024}, and for $\log_{10}M_b=10.5$--$10.8$ the Keplerian exterior falls $5$--$9\sigma$ below those points.
\item \textbf{Support-limited stock.} The support rule of Section~18 carries only the stock inside the region where the lens's condensed baryons dominate the baryonic field. Hot circumgalactic gas carries no records. Where its enclosed mass exceeds the stars and cold gas the support ends, and the stock outside keeps the collisionless distribution of its assembly. For hot-gas masses of $0.1$--$0.4$ of the cosmic baryon share, isothermal to the virial radius, the support ends at $60$--$220\,$kpc in the four stellar-mass bins of \citet{Brouwer2021}. The carried stock inside relaxes to the exponential law, and the profile outside follows the halo, with no Keplerian edge.
\end{itemize}
A population forward model of the four isolated stellar-mass bins of \citet{Brouwer2021} tests the support-limited stock against $\Lambda$CDM halos. It uses 60 excess-surface-density points with their full covariance, the stellar-to-halo relation of \citet{Moster2013} with $0.15\,$dex scatter, the halo mass function, stellar-mass errors, the same hot gas in every model, and a common environment term with free amplitudes in each bin. With the halo-mass normalization and the concentration fitted as common nuisance parameters for both, the support-limited stock gives $\chi^2=140$ and truncated NFW halos give $133$. The whole-budget Keplerian configuration gives $270$ before those nuisances are added, against $252$ for NFW halos. Every model leaves $3$--$6\sigma$ residuals at $0.5$--$1.5\,$Mpc, where neighbouring structure contributes in both theories \citep{MisteleFlat2024}. The isolated-lens data therefore do not separate the support-limited stock from $\Lambda$CDM halos; the hot-gas masses of the lenses fix the support radius and are the input that sharpens the test.

Joint kinematic and weak-lensing data extend the same RAR by about $2.5$ decades in acceleration. With strict isolation and consistent stellar and gas masses, early- and late-type galaxies follow the same relation \citep{Brouwer2021,Mistele2024}. This supports the leading no-slip choice. Direct solutions for resolved nonspherical lenses remain necessary.

\paragraph{Resolved Milky-Way test.}
For the flat-inner-disk baryonic model of \citet{Wang2026MW}, two independent solvers give the local-contact value $\Delta\Phi(8.17,1.055)=1792.2\ {\rm km^2\,s^{-2}}$ against $1294.8\pm54.0$. That baryonic model also gives only $v_c(R_0)=199\ {\rm km\,s^{-1}}$ under the local response, against $237\ {\rm km\,s^{-1}}$. The director response changes the circular speed at $R_0$ by less than $0.5\%$, so it cannot supply the missing enclosed mass. The comparison below therefore continues the inner thin disk exponentially, with $h_{\rm in}=2.2$--$2.5\,\mathrm{kpc}$ and $M_b\simeq(7$--$8)\times10^{10}M_\odot$, a baryonic premise that RAR-based fits share.

With those baryons the responses compare as follows, at Wang's local stellar surface density and at the \citet{McKee2015} level. The rotation-curve $\chi^2_v$ uses 34 points and the vertical-potential $\chi^2_\Phi$ uses 43.
\begin{center}
{\small
\begin{tabular}{lcccc}
\hline
response & $\Delta\Phi(8.17,1.055)$ [${\rm km^2\,s^{-2}}$] & $v_c(8.22)$ [${\rm km\,s^{-1}}$] & $\chi^2_v$ & $\chi^2_\Phi$ \\
\hline
Newtonian baryons & $1233$ / $1135$ & $193$ & $341$ / $352$ & $183$ / $258$ \\
local contact & $1838$ / $1699$ & $236$ & $3.6$ / $5.0$ & $561$ / $357$ \\
director completion & $1416$ / $1315$ & $235.6$ / $235.0$ & $4.6$ / $6.1$ & $137.1$ / $135.9$ \\
baryonic proxy, $c=1$ & $1401$ / $1302$ & $235.5$ / $234.9$ & $4.8$ / $6.3$ & $135.7$ / $139.5$ \\
rigid radial director & $1396$ / $1297$ & $235.7$ & $4.7$ / $6.2$ & $136.0$ / $141.8$ \\
fitted NFW controls & -- & -- & -- & $135$ \\
\hline
\end{tabular}
}
\end{center}
The director completion uses the stock amplitude at $c=0.03$ (Appendix~N.8), with the Milky Way's stock the $C_R=14$ capped response of Appendix~O.3. It fits the vertical potential as well as the NFW controls, both with $\chi^2_\Phi/N\simeq3.2$, and leaves the rotation-curve fit essentially unchanged. The solution sits in the stiff limit:
\begin{itemize}
\item at $(R_0,1\,\mathrm{kpc})$ the director lies $1.0^\circ$ from radial while the source field is tilted by $14.3^\circ$;
\item varying $0.03\le c\le1$ spans $\Delta\Phi=1396$--$1416\ {\rm km^2\,s^{-2}}$, and $c=0.01$ gives $1445$;
\item the stiffness source is $-0.25\%$ of $\boldsymbol\nabla\!\cdot\mathbf J_\perp$ at $R_0$ and $-0.03\%$ of the excess mass within $30\,\mathrm{kpc}$.
\end{itemize}
The excess density at $R_0$ is $0.37$, $0.35$, and $0.30\ {\rm GeV\,cm^{-3}}$ at $|z|=0$, $1$, and $3\,\mathrm{kpc}$.

The baryonic proxy at $c=1$ gives $1401$ and $1302\ {\rm km^2\,s^{-2}}$ and $0.35$, $0.34$, and $0.29\ {\rm GeV\,cm^{-3}}$. With the proxy, rescaling $1/4\le c\le4$ spans $1397$--$1414$. For that solution two independently written solvers agree to $0.1\ {\rm km^2\,s^{-2}}$, find the same minimizer from local and radial starting fields with no negative Hessian pivots, and are unchanged when the director grid is halved. The two amplitudes differ by about 1\% in $\Delta\Phi$ and by $0.02\ {\rm GeV\,cm^{-3}}$ in the local excess density.

\paragraph{Resolved disk galaxies.}
A director solve with the stiffness source for each of the 131 SPARC galaxies passing the quality cut ($Q\ne3$, $i\ge30^\circ$, $\delta V/V<0.1$, $R<2.5\,r_M$; 1683 points) reproduces the algebraic RAR to $0.012$ dex in every $r/r_M$ bin, with rms residual $0.147$ dex for both. It stays within $0.011$ dex of the rigid radial director at every point, and the stiffness source moves no point by more than $0.001$ dex. The local two-dimensional solve departs from the algebraic relation by up to $0.025$ dex in the same bins. These solves use the baryonic proxy at $c=1$. They stay within $0.011$ dex of the rigid radial director, which is the stiff limit shared by both amplitudes.

\paragraph{Vertical dynamics of face-on disks.}
Removing the phantom sheet changes vertical dynamics, and this separates the completion from QUMOND. The DiskMass and NGC~6946 solutions below use the baryonic proxy at $c=1$, and the predictions registered in Section~27.1 keep that basis. For a face-on disk the mass-weighted dispersion $\langle\sigma_z^2\rangle=\int z\,\nu\,\partial_z\Phi\,dz/\int\nu\,dz$ follows from the solved potential at the rotation-curve mass-to-light ratio. For the DiskMass galaxies \citep{Martinsson2013a,Martinsson2013b}, with $\Upsilon_K$ from MOND rotation-curve fits \citep{Angus2015}, the director excess raises $\sigma_z^2$ above the Newtonian value by a median factor $1.20$ at the DiskMass scale heights, against $1.73$ for QUMOND. Matching the measured integrated-light dispersions of the 27 unflagged galaxies then requires stellar scale heights $h_z^*/h_{z,\rm DMS}=0.55$ (16--84\%: $0.39$--$0.77$) for the completion, $0.37$ for QUMOND, and $0.61$ for Newtonian baryons alone. As a pipeline check, Newton at the DiskMass mass-to-light ratios reproduces the measured dispersions to $3\%$.

The shortfall common to every baryons-only model traces to the integrated-light dispersion. It mixes a cold young layer with the old disk, whose scale height enters the vertical Jeans relation \citep{Aniyan2016}. A forward model of the measurement uses isothermal populations with a Milky-Way-shaped age--velocity relation, $V$-band, $K$-band, and Mg\,b light weights, the $17\ {\rm km\,s^{-1}}$ PPak instrumental profile, a single-Gaussian fit, and the $K$-band scale height. Relative to a single-population disk it gives $\Sigma_{\rm true}/\Sigma_{\rm DMS}=1.25$ (16--84\% over 66 population variants: $1.11$--$1.49$). A DiskMass quality cut keeps the 22 galaxies whose DiskMass Newtonian model reproduces its own dispersion fit within $25\%$ and whose dispersions stay above $12\ {\rm km\,s^{-1}}$. With the correction these require $h_z^*/h_{z,\rm DMS}=0.76$ ($0.66$--$0.90$) for the completion and $0.55$ ($0.46$--$0.66$) for QUMOND. The completion-to-QUMOND ratio stays within $1.32$--$1.45$ in every variant.

NGC~6946 supplies the old-disk dispersion directly, from two-component fits to high-resolution integrated light in the inner disk and planetary nebulae farther out: $\sigma_{\rm old}=65$, $65$, $32$, $21$, and $15\ {\rm km\,s^{-1}}$ at $1.5$--$9\,\mathrm{kpc}$ \citep{Aniyan2021}. With $h_z=346\,\mathrm{pc}$ from the DiskMass scaling relation, an isothermal old disk, a $15\%$ cold layer, and the measured gas, one stellar mass-to-light ratio must fit both the SPARC rotation curve and $\sigma_{\rm old}$. The completion gives $\Upsilon=0.31^{+0.09}_{-0.08}$ from $\sigma_{\rm old}$ against $0.39$ from the rotation curve, a $1.0\sigma$ difference, with $\chi^2=5.9$ for five points at the rotation-curve value. QUMOND gives $0.19^{+0.08}_{-0.04}$, a $2.6\sigma$ difference, and $\chi^2=13.7$. The completion remains acceptable for an exponential vertical profile, the TRGB distance, and $h_z$ changed by $\pm20\%$, with $\chi^2=4.9$, $4.9$, $4.5$, and $10.1$ ($p\ge0.07$). The single-component dispersions of the same galaxy give $\chi^2=154$ for the completion and $345$ for QUMOND. The DiskMass-type tension therefore disappears in this galaxy once the old-disk dispersion is used.

NGC~628 supplies a second direct case \citep{Aniyan2018}. Its old-disk dispersions are $55.4$, $50.9$, $33.8$, $22.6$, and $17.5\ {\rm km\,s^{-1}}$ at $2.6$--$12.2\,\mathrm{kpc}$, with $h_z=398\pm88\,\mathrm{pc}$ and $V_{\max}=180\pm9\ {\rm km\,s^{-1}}$ at the kinematic inclination of $8.5^\circ$. A reduced vertical model treats the old disk as an isothermal layer that carries the tracer, the cold layer and the gas as thin sheets, the excess as a round component contributing $(g_{\rm ex}/R)(\pi^2/3)h_z^2$, and QUMOND through the factor $\nu(y)$ on the disk self-gravity. On NGC~6946 it reproduces the full solutions above, with $\Upsilon=0.29$ for the completion and $0.17$ for QUMOND against $0.40$ from the rotation curve. For NGC~628 the stellar mass required by $\sigma_{\rm old}$ is $(2.6\pm0.4)\times10^{10}M_\odot$ for the completion and $(1.1\pm0.3)\times10^{10}M_\odot$ for QUMOND, and the rotation velocity requires $(2.8\pm0.5)\times10^{10}M_\odot$. The completion agrees at $0.3\sigma$; QUMOND differs by $2.9\sigma$ and fits the dispersion profile with $\chi^2=8.4$ for five points. Across variants of the scale length, gas surface density, scale height, excess axis ratio, and rotation velocity the completion stays within $1.7\sigma$ and QUMOND lies $1.3$--$3.9\sigma$ away. The old-disk test separates the completion from QUMOND. A maximal disk with a round halo has the same vertical structure as the completion, so the test leaves $\Lambda$CDM untouched.

\paragraph{Shape of the excess.}
The same solution fixes the shape of the excess. For the Milky Way the total potential has $q_\Phi=0.87$--$0.88$, $0.95$, $0.98$, and $0.99$ at $10$, $20$, $30$, and $50\,\mathrm{kpc}$. The baryonic proxy gives $0.88$--$0.89$, $0.96$, $0.98$, and $0.99$, with excess-density axis ratio $q_\rho=0.91$--$0.92$, $0.94$, $0.95$, and $0.98$. For the 53 SPARC disks with $V_{\rm flat}>150\ {\rm km\,s^{-1}}$ the radial-director limit gives $q_\Phi=0.95\pm0.05$ at $20\,\mathrm{kpc}$ and $0.98\pm0.02$ at $30\,\mathrm{kpc}$. Disks with $V_{\rm flat}=120$--$150\ {\rm km\,s^{-1}}$ are rounder, with $q_\Phi=0.99$ at $20\,\mathrm{kpc}$. Hydrodynamic $\Lambda$CDM halos of Milky-Way mass have inertia-tensor density axis ratios $c/a\simeq0.67$--$0.72$ with halo-to-halo scatter $0.1$ at these radii \citep{Prada2019,Chua2019}. An oblate NFW homoeoid with $c/a=0.67$--$0.72$ has potential axis ratio $0.85$--$0.87$ at $20\,\mathrm{kpc}$ and $0.87$--$0.89$ at $50\,\mathrm{kpc}$. In potential flattening the two predictions therefore differ by about $0.1$.

Milky-Way tracers with six-dimensional data give $q_\rho=1.00\pm0.09$ from RR Lyrae at $5$--$20\,\mathrm{kpc}$ \citep{Wegg2019} and $1.06\pm0.06$ from three streams near $20\,\mathrm{kpc}$ \citep{Palau2023}. Their weighted mean, $1.04\pm0.05$, lies $1.9\sigma$ above the predicted $0.94$. It lies $2.9$--$3.3\sigma$ above the $\Lambda$CDM halos, including their scatter, and the likelihood ratio on statistical errors is $20$--$90$ in favor of the completion. A single galaxy caps this ratio: for a prediction uncertainty of $0.02$, even a perfect measurement at the predicted value would give only $60$--$190$, because the $\Lambda$CDM prediction for one halo carries the population scatter.

A population measurement has no such cap, and \citet{Chemaly2026} provide the first extragalactic one. They fit projected tracks of nearby extragalactic streams with a halo-only axisymmetric NFW potential of free orientation and potential flattening $q$, and infer a truncated-Gaussian population. Their gold subsample of 17 streams gives $\mu_q=0.72^{+0.16}_{-0.14}$ and $\sigma_q=0.34^{+0.18}_{-0.19}$; the bronze subsample of 15 gives $0.97^{+0.19}_{-0.17}$ and $0.30^{+0.20}_{-0.16}$. Synthetic orbit-segment tracks with $1\,\mathrm{kpc}$ noise, viewed at random orientations and fitted the same way with halo mass free and the disk omitted, return $\hat\mu_q=1.02$--$1.03$ for a round halo with a Milky-Way disk and $0.86$--$0.91$ for $q_\Phi=0.87$. Both inputs have zero intrinsic scatter and still return $\hat\sigma_q=0.3$--$0.5$, so the inferred scatter carries no shape information. Bootstrapped through the same hierarchical model, the round input yields a gold-like $\hat\mu_q\le0.72$ in at most $3\%$ of 17-stream samples, and the $\Lambda$CDM-like input in $1$--$16\%$. The bronze result distinguishes neither. The combined likelihood ratio is about $4$ in favor of a $\Lambda$CDM-like halo, with a range of $3$--$14$ over the assumed per-stream likelihood width, and comes almost entirely from the gold subsample. Read as a potential flattening, $\mu_q=0.72$ corresponds to a density axis ratio near $0.4$, flatter than the simulated halos. With per-stream noise of $0.3$ in $q$, separating $0.97$ from $0.87$ at $3\sigma$ requires about 80 streams of gold quality. The Milky Way and the present stream population therefore point in opposite directions, and a disk-modeled population of order $10^2$ streams can separate the two predictions.

\paragraph{Milky-Way compilation and the tilted stellar halo.}
Nine published density axis ratios of the inner Milky-Way halo run from $0.75\pm0.03$ for $87$ streams \citep{Ibata2024} to $1.20\pm0.03$ for the Helmi streams \citep{Woudenberg2024}; the others are those of \citet{Bovy2016}, \citet{Wegg2019}, \citet{MalhanIbata2019}, \citet{Palau2023}, \citet{NibauerBonaca2025}, \citet{DillamoreSanders2025}, and \citet{Zhu2025}. A fixed-effect combination gives $\chi^2=115$ for eight degrees of freedom, so differences between methods exceed the quoted errors. A random-effects combination gives $\bar q_\rho=0.95$ ($16$--$84\%$: $0.89$--$1.01$) with an inter-study scatter of $0.14$. With that scatter marginalized, the Bayes factor for the predicted $0.93\pm0.02$ against $\Lambda$CDM halos at $0.70\pm0.10$ is $16$, between $8$ and $21$ with any one study removed, $11$ with the scatter fixed at the DerSimonian--Laird value $0.20$, and $3$--$4$ against a rounder $\Lambda$CDM prior of $0.80\pm0.10$. The largest stream sample lies $6\sigma$ below the prediction on its own errors.

The stellar debris of Gaia--Sausage--Enceladus is tilted by $25.4^\circ\pm3.2^\circ$ above the disk, with density-contour axis ratios $0.81$ and $0.73$ \citep{DillamoreSanders2025}. Test-particle integrations of a radial-orbit population with that tilt, in potentials that share the disk, bulge, and circular speed, give a tilt half-life of $0.15$--$0.25\,$Gyr in the aligned excess and in an aligned flattened NFW halo. An NFW halo tilted with the debris holds the stellar tilt at $17$--$19^\circ$. Radial orbits cross the disk at pericenter, and differential torques erase their coherent orientation in the tested aligned potentials. Under the adopted $8$--$10\,$Gyr age and long-lived-debris interpretation, the aligned Milky-Way solution fails this test.

\paragraph{Tilt supported by conserved stock.}
The retained stock supplies a continuing gravitational source whose density can carry its assembly history. Its velocity stress supports that distribution against collapse and phase mixing. Appendix~O.3 gives a conditional collisionless realization with positive phase-space mass, self-gravity, and reciprocal disk torque. The stock density's major axis is distinct from the local director orientation. The radial-director identity in N.8 allows an anisotropic stock stiffness without a transverse elastic torque on an exactly radial director; the disk's thermal forcing and finite director departures still require their own solution. A fast response field can follow this slowly evolving source without requiring the source to rearrange on a light-crossing time.

The constructed density is
\[
\rho_R(r,\hat{\mathbf n})=\rho_0(r)\left[1+A P_2(\hat{\mathbf n}\cdot\hat{\mathbf e})\right],
\qquad A=0.3,\qquad P_2(\mu)=\frac{3\mu^2-1}{2}.
\]
Its angular factor lies between $0.85$ and $1.30$ and preserves the spherical average and total stock budget. The frozen radial target uses $a_0=1.230729\times10^{-10}\,\mathrm{m\,s^{-2}}$ and a stock mass of $8.05237\times10^{11}M_\odot$. The $43^\circ$ initial axis, quadrupole amplitude, orbit weights, and construction rate $\Omega_0=0.608998\,\mathrm{rad\,Gyr^{-1}}$ specify the candidate. The angle is motivated by a potential-dependent halo inference \citep{DillamoreSanders2025}. Formation of this orbit population from an assembly history remains open.

Three live realizations use the same 1072 positive orbit totals and 8192 massless stellar tracers, with 131072 stock samples and different stock phase seeds. The stock generates its evolving field, and the disk receives the opposite torque as a rigid gyroscope. Complete $\ell=0,2,4$ multiplets, smooth radial assignment, and a symmetric pair energy define the force. Exact inversion partners enforce a centered, even-parity configuration. All three pass the preregistered descriptive retention condition: stock tilt at least $35^\circ$ and stellar tilt at least $20^\circ$ in $10$--$40\,$kpc, using moments averaged over $7.5$--$8\,$Gyr before extracting their axes. The corresponding results are
\begin{center}
\small
\begin{tabular}{lccccc}
\hline
 & \multicolumn{2}{c}{$10$--$40\,$kpc} & \multicolumn{3}{c}{Stellar, $6$--$60\,$kpc} \\
Realization & Stock tilt & Stellar tilt & Tilt & $p_*$ & $q_*$ \\
\hline
Run 1 & $40.79^\circ$ & $22.16^\circ$ & $24.38^\circ$ & $0.8487$ & $0.7905$ \\
Run 2 & $39.58^\circ$ & $21.38^\circ$ & $22.62^\circ$ & $0.8513$ & $0.7898$ \\
Run 3 & $42.36^\circ$ & $22.33^\circ$ & $24.75^\circ$ & $0.8742$ & $0.7968$ \\
\hline
\end{tabular}
\end{center}
The largest sampled relative energy error is $0.00247\%$; total angular-momentum error is below $1.30\times10^{-13}$ of the disk angular momentum. The seeds test repeatability conditional on one orbit construction and stellar sample. The lower-count controls retain smaller tilts, and their phase draws also differ. Particle-number convergence is therefore unresolved.

The stellar-shape comparison uses the position-inertia tensor inside the spherical $6$--$60\,$kpc cut of \citet{DillamoreSanders2025}; $p_*$ and $q_*$ are square-root eigenvalue ratios. Mapping the adopted density through that cut gives $p_*=0.8439\pm0.0257$ and $q_*=0.7739\pm0.0182$, with target inclination $25.39^\circ\pm3.16^\circ$. The mapping uses 1000 independent Gaussian shape draws at fixed radial profile. The three sums of squared standardized residuals in inclination, $p_*$, and $q_*$ are $0.96$, $1.62$, and $3.02$. These supplementary diagnostics omit the survey selection, finite-tracer covariance, and full orientation likelihood. They have no calibrated goodness-of-fit probability. At the final instantaneous checkpoint, separate minor-axis checks retain offsets of $1.5$--$5.6^\circ$ after azimuth alignment. Streaming-subtracted stellar dispersions in $6$--$24\,$kpc give $\beta\simeq0.879$--$0.884$, so strongly radial motion survives.

The fitted disk precession rates are $0.559$--$0.564\,\mathrm{rad\,Gyr^{-1}}$, $7.5$--$8.2\%$ below the construction rate. This compares two model rates; no observed precession rate enters that residual. A relative equilibrium at the evolved rate would require new orbit weights and torque balance. The present configurations also acquire $12$--$23^\circ$ of late halo azimuth slip relative to the disk and become rounder. Those relative changes persist under a common change of rotating frame, so exact stationarity is unestablished.

The evolved fields retain small changes in the sampled radial and even vertical observables. Relative to the frozen spherical stock plus baryons, the maximum change in $\sqrt{-Ra_R}$ is $1.35$--$1.46\%$ over the sampled $8.2$--$60\,$kpc radii. The even vertical potential at $R=8.2\,$kpc and $|z|=1.1\,$kpc changes by $-0.14$ to $+2.05\%$. The odd part, half the potential difference between the two heights, reaches $186$--$257\,\mathrm{km^2\,s^{-2}}$ in absolute value and requires a deformable-disk calculation. These are compatibility checks of the chosen configuration. They supply no new rotation-curve or vertical-equilibrium likelihood.

The weaker transverse force remains sensitive to angular resolution: adding $\ell=6$ changes it by roughly $25$--$29\%$, and adding $\ell=8$ gives further changes of $20$--$31\%$ in the inspected snapshots. The much smaller total-force changes do not establish convergence of the force supporting the tilt. Positive sample masses also leave small negative artifacts in the truncated density reconstruction, at most $1.66\times10^{-5}$ of the stock mass on the tested extended grids. These calculations establish repeatable eight-Gyr tilt retention within the specified kinetic and rigid-disk model. The coupled director response, a bending disk, force convergence, and formation of the candidate remain required for a full halo prediction. The aligned-branch failure and its registered test in Section~27.1 retain their stated scope.

\paragraph{Solar-System quadrupole and wide binaries.}
In the Solar System the Milky Way is the active carrier. The Sun and planets remain resolved descendants in the Einstein branch and do not re-enter $\Gamma_\perp$ as separate transverse sources. The Sun-generated anomalous quadrupole is therefore
\[
\boxed{Q_2^{\odot,\perp}=0}
\]
at leading order. The remaining smooth Galactic contribution is a tide of order
\[
\frac{g_{\rm MW}}{R_0}\simeq9.2\times10^{-31}\ {\rm s}^{-2},
\]
about two thousand times smaller than the current Cassini $1\sigma$ uncertainty. Cassini tracking gives $Q_2=(1.6\pm1.8)\times10^{-27}\,{\rm s}^{-2}$ and places universal nonlinear functions of the total local field in $3$--$15\sigma$ tension with galaxy fits; it also limits their Milky-Way boost at the Sun to two percent at 95\% confidence \citep{Park2026}. Carrier nesting avoids that total-field cross-term.

The avoidance fixes a minimum resolution for the parent's conditional expectation. A descendant of mass $M$ enters the parent's coarse source smoothed over a width $L$. The linear response to the Galactic field then gives, at its centre,
\[
Q_2=\frac{4}{15\sqrt{2\pi}}\,(n_\parallel-n_\perp)\frac{GM}{L^3},
\]
with $n_\perp=n_B(x_e)=0.645$ and $n_\parallel=d(n_Bg)/dg|_{g_e}=0.148$ at the Sun. The coefficient was verified analytically and numerically by an independent calculation. Cassini's $2\sigma$ interval, $-2.0$ to $+5.2\times10^{-27}\,{\rm s}^{-2}$, therefore requires:
\begin{itemize}
\item $L\ge0.049\,$pc ($10\,200\,$AU) if the Galactic director follows the local field, where the quadrupole is negative;
\item $L\ge0.024\,$pc ($4900\,$AU) if the director is locked ($n_\perp\to0$), where the quadrupole is positive.
\end{itemize}
Both lengths are comparable to the Sun's MOND radius, $0.034\,$pc. Globular clusters bound $L$ more strongly. Palomar~14 resolved in the joint source would have $\sigma=1.49\ {\rm km\,s^{-1}}$, against the measured $0.38\pm0.12$ \citep{Jordi2009}, and $0.54$ when smoothed into its parent. If Palomar~14 is a descendant of the Milky-Way carrier, $L\gtrsim25$--$50\,$pc.

Wide binaries inside the Milky-Way carrier are likewise Newtonian apart from the smooth Galactic tide. For a $1.5M_\odot$ pair on the flat Galactic rotation curve the Jacobi radius $r_J=[GM/(4\Omega^2-\kappa^2)]^{1/3}$ is about $1.6\,{\rm pc}$, and the fractional speed corrections from the tide satisfy
\[
\frac{\Delta v}{v_N}\simeq\frac12\left(\frac{s}{r_J}\right)^3\lesssim
2\times10^{-5},\ 1\times10^{-4},\ 4\times10^{-4}
\]
at separations of $10$, $20$, and $30\,{\rm kAU}$, respectively. Current Gaia analyses remain divided: one recent treatment favors Newtonian dynamics after modeling unresolved triples, while two analyses report a low-acceleration anomaly \citep{Pittordis2025,Chae2025,Chae2026}. The carrier-resolved branch predicts the Newtonian outcome.

The background vector $u^\mu$ enters the parent carrier response, but its local spatial variation is limited by the Galactic tide. Leading Solar-System PPN values and vacuum gravitational-wave propagation remain those of the Einstein branch. Time-dependent vector contractions and nonlinear preferred-frame terms remain part of the covariant projector audit.


\paragraph{Fluctuations.}
The equilibrium occupation sets the static susceptibility $\chi_0=\beta n_B(1+n_B)$, where $\beta=(k_BT_H)^{-1}$. A causal linear response with relaxation rate $\gamma>0$ and finite response bandwidth $\Omega\ge\gamma$ is
\[
\chi^R_{nE}(\omega)=-\chi_0
\frac{\gamma\Omega}{(\gamma-i\omega)(\Omega-i\omega)}.
\]
Both poles decay, and $\chi^R_{nE}(0)=-\chi_0$ preserves the thermal static contact. Quantum equilibrium fixes the symmetrized noise from this response. Appendix~N.8 derives the positive finite variance $\sigma_n^2$ and its ratio $\mathcal F=\sigma_n^2/[n_B(1+n_B)]\ge1$ to the classical population variance. For $N_{\rm eff}$ independent cells,
\[
\operatorname{Var}(\bar n)=\frac{\sigma_n^2}{N_{\rm eff}},\qquad
\frac{\sigma_g}{g}
=\sqrt{\frac{\mathcal F n_B}{N_{\rm eff}(1+n_B)}}.
\]
The sufficient cell counts for ten-percent and three-percent rms fluctuations are $N_{\rm eff}>100\mathcal F$ and $N_{\rm eff}>1112\mathcal F$. A measurement over a finite band uses the same noise spectrum with its temporal response window. The rates and spatial correlation volume determine the noise level; the stationary mean response is independent of them.

\section*{Part V. Transport, Clusters, and Cosmology}
\addcontentsline{toc}{section}{Part V. Transport, Clusters, and Cosmology}

The microstructure-to-static-weak-field chain is the most directly constrained part of the theory. The sectors of this part ask how the same finite-capacity substrate behaves when sources move. Transport asks how the capacity-strain field propagates after a source changes; the cluster sector asks how it couples to matter in distinct dynamical phases; cosmology treats the homogeneous mode; and the saturated phase treats the proposed committed carrier. These sectors share the static weak-field ontology but not its evidential status: the Einstein/capacity action equivalence, ordinary source map, and finite marked-transfer scale are the controlled baseline inside their displayed actions, while the transport, cluster, and cosmological additions carry their stated conditional or open grades.

\section{17. Causal Transport and Telegrapher Dynamics}

The static equations give settled fields. Moving sources require a causal response with the same zero-frequency limit. Relative to the substrate four-velocity $u^\mu$, the ordinary capacity deficit retains the covariant telegrapher completion
\[
\tau_0(u^\mu\nabla_\mu)^2\delta S
+u^\mu\nabla_\mu\delta S
=D h^{\mu\nu}\nabla_\mu\nabla_\nu\delta S+A\chi,
\qquad
h^{\mu\nu}=g^{\mu\nu}+u^\mu u^\nu,
\]
with $A/D=\kappa/\gamma$. The carrier-resolved transverse excess has the corresponding local-frame equation
\[
\boxed{
\tau_0\ddot\phi_{\perp,P}+\dot\phi_{\perp,P}
-D\nabla^2\phi_{\perp,P}
=-4\pi GD\,\rho_{\perp,P}[T_P^{\rm cg}],
}
\]
where
\[
\rho_{\perp,P}=\frac{1}{4\pi G}
\boldsymbol\nabla\!\cdot\mathbf J_{\perp,P},
\]
Its static limit is exactly the Poisson equation of Section~15. No separate static form factor is introduced.

Causal propagation and the minimal infrared scale choice give
\[
\frac{D}{\tau_0}=c^2,
\qquad
\tau_0^{-1}=H_0,
\qquad
D=\frac{c^2}{H_0}.
\]
For a Fourier mode, the response relative to the static solution is
\[
F(\omega,k)=
\frac{Dk^2}{Dk^2-\tau_0\omega^2-i\omega},
\qquad
F(0,k)=1\quad(k\ne0).
\]
The Laplace poles are
\[
s_\pm=\frac{-1\pm\sqrt{1-4\tau_0Dk^2}}{2\tau_0}.
\]
For $k\ne0$, both have negative real part. At $k=0$ the rates are $0$ and $-1/\tau_0$, with the zero rate representing the conserved mode. The characteristic propagation speed is $\sqrt{D/\tau_0}=c$. For galactic source frequencies $\omega\sim vk$, departures from the static branch begin at $O(v^2/c^2)\sim10^{-6}$; the linear damping correction is smaller on kiloparsec scales. The quasistatic RAR and lensing equations therefore survive without spectral suppression.

This transport law governs propagation delay, relaxation, and merger-era lag. Higher-derivative corrections at the block scale remain ultraviolet data. The cluster stock and merger phenomenology that this equation evolves is developed in Section~18.

\section{18. Cluster Stock, Conservative Contact, and Lensing}
\label{sec:clusterstock}

Cluster lensing and dynamics probe the same gravitating inventory. In the
weak-field kinetic realization of O.3, baryons and a positive retained stock
move in a common potential,
\[
\nabla^2\Phi=4\pi G(\rho_*+\rho_g+\rho_C),\qquad
\kappa=\frac{\Sigma_*+\Sigma_g+\Sigma_C}{\Sigma_{\rm crit}}.
\]
The static response prescribes a possible equilibrium of this stock.
Adding that equilibrium as a second halo would duplicate its mass.
Gas can shock while stars and weakly coupled stock continue through an
encounter. A spatial separation of gas and lensing peaks is therefore
dynamically possible; its size depends on the initial distributions,
encounter, and contact law.

\paragraph{One cohort budget.}
For a fully released closed cohort, use the formation mass convention of Appendix~O.15, in which observed baryon masses are quoted,
\[
k_C=5.333\,e^{-\eta_*}=5.1760752,\qquad C_0=k_C B_0,
\qquad f_{b,\rm coext}=\frac1{1+k_C}=0.1619151.
\]
Here $B_0$ is its ancestral baryonic mass and $C_0$ its surviving stock. The bare convention gives $e^{-\eta_*}/\epsilon=5.1642$ and $1+k_C=6.164$.
The mass convention and release history must be kept fixed across systems.
Inside an observed aperture, the retained fractions of baryons and stock
can differ. Writing them as $f_b^{\rm ret}$ and $f_C^{\rm ret}$ gives
\[
\frac{C}{B}=k_C\frac{f_C^{\rm ret}}{f_b^{\rm ret}},\qquad
\eta_{\rm ret}:=\frac{k_C B}{M_{\rm grav}-B}
=\frac{f_b^{\rm ret}}{f_C^{\rm ret}}.
\]
Thus $M_{\rm grav}=(1+k_C)B$ tests coextensiveness. It is a restrictive
special case of the transport model. For any positive inferred stock
$C=M_{\rm grav}-B$, the minimal compatible closed inventory is
\[
B_0=\max(B,C/k_C),\qquad
B_{\rm out}=B_0-B\ge0,\quad C_{\rm out}=k_CB_0-C\ge0.
\]
This construction supplies an admissible inventory. Its assembly history,
work cost, and transport time require a forward solution.

\paragraph{Why cluster stock keeps the cohort ratio.}
The selection of Appendix~O.3 acts through the occupation cells of active carriers, on the stock each carrier carries. An active carrier's records are identified with its condensed baryons, stars and cold gas, the baryons whose progenitor history remains recoverable in the sense of the release sector (Appendix~O.15). Shock-heated gas whose cooling time exceeds its dynamical time relaxes to the Gibbs state of its potential and holds no record independent of that potential. A carrier's support ends where its records stop dominating the baryonic field. For an isolated carrier with a nearly spherical baryonic distribution this is the radius at which its records stop dominating the enclosed baryonic mass. A satellite whose host's baryonic field exceeds its own at its half-light radius has no support of its own and remains under the host carrier. Stars, star clusters and tidal dwarfs formed from a host's gas inherit its record and are its descendants whatever their own field (Section~16), so the Sun and wide binaries are descendants of the Milky Way, and the Sun's own field, dominant at its radius, does not make it a carrier. Its cells act on the stock bound to the potential of its records and the stock it carries. In the seven X-COP clusters with stellar profiles, condensed baryons are 2.5--5\% of the baryons inside $R_{500}$, and the central galaxy's support ends at 47--99\,kpc, or $0.04$--$0.07\,R_{500}$. Cluster stock outside that support, or crossing it unbound, enters no cell and evolves as collisionless matter in the common potential. It keeps the cohort ratio $1+k_C=6.18$, between the X-COP medians $6.10$ at $R_{200}$ and $6.45$ at $R_{500}$. A response that acted on the cluster stock would fail in either direction. At $R_{500}$ the baryonic acceleration is $0.05$--$0.07\,a_0$ and the exponential law gives a total-to-baryon ratio of $4.2$--$5.0$. Imposed on the existing stock it removes 20--35\% of the stock from $R_{500}$; added to the cohort stock as a separate force it gives $9.4$--$10.2$. Galaxy outskirts that obey the law span the same range of $g_{\rm bar}/a_0$, at 35--50\,kpc in a Milky-Way model, so the switch cannot be a function of acceleration. The cells also bound the work available to contacts: their energy per unit carried stock tends to $V_f^2$ in the deep regime (Appendix~O.3), and the X-COP cores hold 20--130\,${\rm km^2\,s^{-2}}$ per unit cluster stock. Contact runs on the seven clusters, started from the observed profiles with the contact kernel of Appendix~O.3, keep the mass inside $R_{500}$ within 0.1\% and inside $0.1R_{500}$ within 1.5\% of the no-contact control over 10\,Gyr; with the contacts confined to the cores, the outer result follows from the support itself. The same contacts with cells throughout the cluster and the galactic work allowance leave 55--79\% inside $R_{500}$. The central galaxy's carried stock relaxes from above its galactic-law value toward it, while its support keeps seven to eleven times that value in uncarried cluster stock, consistent with the observed departure of brightest cluster galaxies and group centrals from the rotating-galaxy relation. Group and cluster stock that is not bound to the central galaxy also lies inside its support, so in X-ray groups the excess over the law begins in the core. In A262, MKW4 and RXJ1159.8+5531 it exceeds $0.1\,$dex within $5$--$13\,$kpc, while hot gas overtakes the condensed baryons only at $85$--$210\,$kpc \citep{Gastaldello2007,Vikhlinin2006}. The record boundary fixes which stock the cells act on, and the uncarried stock profile sets where the excess appears. The X-ray ellipticals on the radial acceleration relation \citep{Lelli2017RAR} test the record identification directly: over the radii used there, stars dominate their enclosed baryons, and in NGC~720 hot gas becomes the dominant baryonic component only outside about 200\,kpc \citep{Humphrey2011N720}, so excluding the hot gas from the record source lowers the outermost points by at most 0.04\,dex. Isolated ellipticals, which hold no group stock, should follow the law out to the radius where their hot gas dominates; NGC~720 lies within about $0.1\,$dex of it at $102\,$kpc in mass models extrapolated beyond the X-ray data. Resolved member galaxies and cosmological assembly remain to be included in the cluster runs.

\paragraph{Reciprocal contact with finite available work.}
The coupling of O.3 admits central pair transitions in the relative
radial velocity $w$. For a pair with reduced mass
$\mu=m_im_j/(m_i+m_j)$ and separation direction $\boldsymbol n$,
\[
\delta\boldsymbol v_i=\frac{m_j}{m_i+m_j}\delta w\boldsymbol n,
\qquad
\delta\boldsymbol v_j=-\frac{m_i}{m_i+m_j}\delta w\boldsymbol n,
\qquad
\Delta K=\frac\mu2[(w+\delta w)^2-w^2].
\]
Both linear and angular momentum are conserved by each impulse. The work
and heat registers balance $\Delta K$ event by event. A symmetric factor
$v(U)v(U-W)$ suppresses a transition when either endpoint has exhausted
its available work; the floor state permits no such transition. Local
partners include gas and stars from both progenitors. An initial cohort
label does not confer a special interaction eligibility.

Seventeen controlled runs comprise fourteen encounters and three isolated
controls. They use coextensive Plummer progenitors of mass ratio $0.3$,
the same $k_C$, a gas adiabatic index $5/3$, and a selected two-component
relative-energy target. Weak contact preserves a stock concentration near
the stars while the gas lags. Increasing the rate alters that separation,
but the strongest tested contact also changes the isolated stock median
radius by $-27\%$ and the gas median radius by $+46\%$ within six dynamical
time units. It has not selected a stable cluster equilibrium.
The zero-work run permits zero events and reproduces the no-contact
trajectory; a low-work run permits 61 events against 7114 with the larger
work reserve. The largest energy drift is below $0.007\%$ of the binding
energy, with momentum errors at roundoff. Resolution, timestep, and
contact-range controls support the qualitative mechanism. These are
selected effective encounters; their rates, target, and work reserves
have not been inferred independently from cluster data. The spatial
stress of the full depleted scalar is not evolved in this particle test.

\paragraph{Radial data and independent lensing.}
Public X-COP profiles provide twelve radial targets, seven with stellar
profiles; five overlap the MENeaCS weak-lensing sample
\citep{Ettori2019Cluster,Herbonnet2020Cluster}. With one fixed normalization,
all nominal NFW targets admit positive cumulative stock and positive
stock in the fifteen tested radial shells. The forward hydrostatic
reconstruction of A644 has negative inferred shells, so positivity depends
on a physically admissible mass reconstruction. Projected stellar masses
are approximately deprojected; unmeasured diffuse light and correlated
profile errors remain relevant.

For the five independent lensing targets, the self-consistent coextensive
$M_{500}$ ranges, in $10^{14}M_\odot$, are
\[
\begin{array}{c|cc}
 &M_{500}^{\rm coext}&M_{500}^{\rm WL}\\\hline
\mathrm{A1795}&3.96\text{--}4.23&9.3\pm2.2\\
\mathrm{A2029}&7.46\text{--}8.07&12.1\pm2.5\\
\mathrm{A2142}&9.29\text{--}10.02&9.7\pm2.3\\
\mathrm{A85}&5.20\text{--}5.43&5.7\pm2.2\\
\mathrm{ZW1215}&4.19\text{--}4.46&3.5\pm2.2
\end{array}
\]
The ranges double the nominal stellar mass and are sensitivity ranges.
A1795 and A2029 require substantial relative baryon--stock segregation
under the central lensing targets: $\eta_{\rm ret}=0.518$ and $0.704$.
A positive transport inventory can accommodate these values, but it does
not predict them. ZW1215 changes substantially with the lensing estimator;
the two estimators use overlapping data and cannot be counted as independent
confirmations. The empirical companion gives the aperture conventions
and profile-level qualifications.

\paragraph{Gas support and possible outlier mechanisms.}
Retaining the common potential fixes the local gas balance,
\[
a_g=-g-\frac{P_{\rm th}'}{\rho_g}
       -\frac{P_{\rm nt}'}{\rho_g}+\frac{f_{Cg}}{\rho_g},\qquad
M=M_{\rm th}+M_{\rm nt}+M_{\rm contact}-M_{\rm acc},
\]
where $M_{\rm contact}=r^2f_{Cg}/(G\rho_g)$ and
$M_{\rm acc}=r^2a_g/G$. Assembly segregation, acceleration, and reciprocal
contact are possible contributions. A quiet gas velocity field alone
does not determine all of these terms. A passive drag vanishes at zero
relative drift, however, and a persistent supporting force requires its
own constitutive derivation. The small measured turbulent support in
A1795 and A2029 constrains a purely turbulent explanation
\citep{Sarkar2026Cluster,XRISM2025Cluster}. These possibilities remain
conditional pending joint lensing, pressure, and transport fits.

\paragraph{A resolved merger target.}
The public CANUCS Bullet v1 reconstruction supplies a separate spatial
test \citep{Rihtarsic2026Cluster}. At $z_l=0.296$, in its
$H_0=70\,\mathrm{km\,s^{-1}\,Mpc^{-1}}$, $\Omega_m=0.3$ convention,
$\Sigma_{\rm crit}=1.8268761\times10^9M_\odot\,\mathrm{kpc}^{-2}$
for $D_{ls}/D_s=1$. Within 150 kpc of the Bullet's BCG3, the reconstructed
total mass is $8.511\times10^{13}M_\odot$ and the supplied gas mass is
$6.184\times10^{12}M_\odot$. The released subcluster halo centre lies
3.92 kpc from BCG3, while the supplied gas-density peak lies about
169 kpc away. The latter is a pixel-based gas-density diagnostic.
The non-gas map still includes stars. These are targets extracted from
an existing best-fit reconstruction, with no simulated-map likelihood.
The illustrative encounters have 24--30 kpc force softening and roughly
90 kpc map smoothing under the quoted physical scaling; they cannot test
the observed few-kiloparsec alignment. A common selected contact must
reproduce the isolated profiles and the merger morphology together.



\paragraph{Falsifiers and status.}
The cluster sector rests on the cohort budget and on the record identification of Appendix~O.3. The cosmic baryon fraction $\Omega_b/\Omega_m\simeq0.156$ lies within $0.7\sigma$ of the bare branch value $\epsilon/(1+\epsilon)=0.158$ \citep{Planck2020}. Relaxed clusters must hold total-to-baryon ratios near $1+k_C=6.18$ outside the central galaxy's support, and a population that departs from it by more than its hydrostatic and stellar-mass systematics rejects the cohort budget. Isolated ellipticals must follow the galactic law out to the radius where their hot gas outweighs their stars. Merger maps must place the lensing peaks with the stars and stock while the gas lags.

\section{19. Cosmology and the Hubble-Tension Sector}

The scalar capacity field has two cosmological roles, and the sector works only if they stay separate. Its homogeneous mode $\overline S(t)$ affects the background expansion and the sound horizon; its inhomogeneous fluctuations $s(x,t)$ still govern local weak-field gravity. The cosmological sector is the homogeneous continuation of the same medium, not an unrelated dark-energy component appended to the weak-field theory: what changes is the kinematic regime, not the ontology, as the background mode becomes dynamically relevant on horizon scales while the local branch stays encoded in the fluctuations.

The cosmological sector uses the same field split,
\[
S(x,t)=\overline S(t)+s(x,t),
\]
where $\overline S(t)$ is the homogeneous mode and $s(x,t)$ the inhomogeneous sector responsible for local weak-field dynamics. The vacuum baseline is fixed by apparent-horizon capacity,
\[
S_\infty(t)=\pi\frac{R_A(t)^2}{L_*^2}.
\]

This is the horizon-normalized representation of the same entropy field used locally. It is compatible with the cell-normalized source theorem because local observables depend on $\delta S/S_\infty$ and $\kappa/(\gamma S_\infty)$ rather than on an absolute entropy unit.

Two independent results clarify what this baseline means. First, the type~II$_1$ static-patch construction makes empty de~Sitter the maximum-entropy gravitational state and places excited semiclassical states below it \citep{ChandrasekaranEtAl2023}. Second, Bianconi obtains the same de~Sitter area scaling from a bulk rather than a horizon entropy. In the $c=1$ convention of her low-curvature Friedmann approximation, the local geometric-relative-entropy density and the causal-diamond four-volume scale as
\[
\frac{\delta s}{\delta v}\simeq \overline{\omega}_{[1]}H^2,
\qquad
V_{\rm dS}\sim H^{-4},
\]
so their product gives
\[
S_{\rm GfE}\simeq
\frac{\overline{\omega}_{[1]}}{\ell_P^4H^2}
\propto H^{-2}
\propto A_{\rm dS}.
\]
Integrating Bianconi's local volumetric information density over the observer's causal diamond gives an area-sized finite capacity \citep{Bianconi2026}. Her result does not fix the coefficient used here, replace $\ell_P$ by the independently calibrated $L_*$, or derive the homogeneous evolution. Her intrinsic geometric $k$-temperatures scale as $H^2$, whereas the Gibbons--Hawking temperature used in Section~15 scales as $H$.

Because the entanglement field couples to the trace of the stress-energy tensor, the homogeneous mode is largely dormant during radiation domination but can become active near matter--radiation equality. The conditional proposal would reduce the sound horizon and shift the CMB-inferred Hubble constant upward: the mechanism has the required sign and turns on at the required epoch. Appendix~O.3 fixes a covariant perturbation system for any supplied source-clock release history. The direct geodesic-clock evolution with radiation and baryons is implemented and fitted for both release histories in O.12. The standard linear-depletion configuration (plain DES-Dovekie, $\sum m_\nu=0.06\,$eV, the low-$\ell$ $\tau$ prior, nonlinear lensing and free Planck calibration) gives $H_0=68.46\,\mathrm{km\,s^{-1}\,Mpc^{-1}}$. The earlier exponential fit, with a free low-redshift supernova offset, $\sum m_\nu=0$, fixed Planck calibration and matched linear lensing, gives $68.9389$. This release calculation does not test the separate sound-horizon mechanism. A joint calculation tests the mode at the level the paper fixes it. The trace coupling fixes its onset at matter--radiation equality, $z_c=3400$, and the abundance tie holds $\omega_c=5.333\,\omega_b$; a transient early-energy component of the axion-like form ($n=3$, initial angle $2.83$) peaking at $z_c$ stands in for the unspecified potential, and its peak fraction $f_{\rm EDE}$ is the one added amplitude. The late universe is the direct geodesic release of Section~20. The early component and the release act at separate epochs, so the CMB spectra are the early-energy spectra at the joint acoustic angle plus the release-sector change in lensing and late integrated Sachs--Wolfe power, and the BAO and supernova distances come from the release background with the early-energy drag scale; as $f_{\rm EDE}\to0$ this construction returns the release fit to within $\Delta\chi^2=0.02$. On the likelihood of Section~20 the fit gives $f_{\rm EDE}=0.018\pm0.010$, $H_0=69.4\pm0.6\,\mathrm{km\,s^{-1}\,Mpc^{-1}}$, $n_s=0.971$ and $r_{\rm drag}=146.4\,$Mpc, with $\chi^2$ lower by $1.9$ than the release fit and by $1.0$ than $\Lambda$CDM at the same parameter count. The tie stops $\omega_c$ from absorbing a larger early energy: $f_{\rm EDE}=0.05$ raises $\chi^2$ by $13.5$, and $f_{\rm EDE}=0.07$ raises it by $34$ and drives the release history outside its physical range. The theory therefore places the Hubble constant at $69.4\pm0.6\,\mathrm{km\,s^{-1}\,Mpc^{-1}}$, below the Cepheid-calibrated value and near the tip-of-the-red-giant-branch value \citep{Freedman2025}. With the late universe held at $\Lambda$ the preferred fraction is $0.033$ at $H_0=69.46$; the release background supplies part of the same expansion history.

The local weak-field predictions are protected by the separation between $\overline S(t)$ and $s(x,t)$. This is the role of the shear-lock logic: changing the homogeneous background mode does not rewrite the local static Poisson branch that governs galactic dynamics and lensing. The two regimes therefore use one scalar medium with unchanged dimensionless galactic coefficient $\epsilon$. The acceleration scale $a_0(z)=\epsilon cH(z)$ still changes with the chosen cosmological history; a joint galaxy--cosmology comparison must use the same $H(z)$. The covariant capacity-clock equations make the corresponding dynamical separation test well posed for the full Boltzmann calculation.


\section{20. The Saturated Phase and the Cosmic Microwave Background}

The conditional lever normalization of Section~15 is applied to the early universe. The source-allocation cap is filled at formation by the map of Appendix~O.15; the recombination-era demand estimated below is a consistency condition on that cap. The galactic force law itself has an unbounded low-acceleration boost. The resulting saturated phase contains a conserved count of committed channels that redshifts as $a^{-3}$. The zero-pressure constraint theorem below is conditional on the pinned realization of that phase, and the numerical abundance inherits the transverse coefficient $\epsilon$.

\paragraph{Saturation requirement at recombination.}
Use $y\equiv g_{\mathrm{pert}}/a_0(z)$ for the acceleration ratio and $x\equiv\sqrt y$ for the thermal argument of Section~15. The linear-domain estimate $y\lesssim2\times10^{-3}$ today and $y\propto\sqrt a$ gives $y\lesssim6\times10^{-5}$ at $z\simeq1100$, hence $x\lesssim8\times10^{-3}$ and $\nu(x)\gtrsim10^2$. This exceeds the total source-allocation cap $1+1/\epsilon=6.32$. Under the adopted demand-to-allocation mapping, the estimate requires saturation for those perturbation modes. The relation $a_0(z)=\epsilon cH(z)$ keeps the excess ratio $cH/a_0=1/\epsilon$ fixed at every redshift; adding the ordinary unit response gives the total cap.

The RAR response $\nu(y)=[1-e^{-\sqrt y}]^{-1}$ exceeds $1+1/\epsilon$ whenever
\[
y<[\ln(1+\epsilon)]^2=0.02966\ldots.
\]
The number $1+1/\epsilon$ therefore bounds the adopted committed allocation per baryonic source; it cannot be a universal ceiling on the galactic acceleration ratio. The recombination estimate assumes a mapping from the unsaturated response demand to that allocation. A cross-phase microscopic calculation must establish this mapping and the recruitment dynamics.

\paragraph{Absolute capacity, source-normalized recruitment, and coherence.}
Three quantities that answer different questions must be kept separate. The availability field $q$ measures the absolute local capacity remaining for ordinary transactions; in the static normalization it obeys $q=1+2\Phi/c^2+O(c^{-4})$. It is not a universal identity with $-g_{00}$ in arbitrary FLRW coordinates. Recruitment saturation is instead defined relative to each source's cap. For a Lagrangian source element $A$ of baryonic mass $M_A$, let $M_{c,A}$ be the mass-equivalent committed allocation assigned to that source. Per-defect bookkeeping gives the cap and its normalized occupancy,
\[
M_{c,A}^{\max}=\frac{M_A}{\epsilon},
\qquad
\boxed{\sigma_A=\frac{M_{c,A}}{M_{c,A}^{\max}}
=\epsilon\frac{M_{c,A}}{M_A}\in[0,1].}
\]
The label $A$ is physical bookkeeping: allocations belonging to distinct defects may overlap spatially, but they cannot be counted twice. In the continuum $A$ becomes a Lagrangian source label, so $\sigma_A$ is source-supported rather than a local projector attached independently to every cell.

Coherence is likewise a property of the committed flow, not a microscopic Boolean register. Write $\mathcal C=1$ when the source-labelled currents combine into one hypersurface-orthogonal, single-stream flow with a single-valued renewal phase; locally this requires vanishing curl of the normalized flow, and globally it requires vanishing phase holonomy. At first shell crossing the Lagrangian map loses invertibility and several velocities occupy one Eulerian point, so $\mathcal C$ falls to zero. Loss of single-stream coherence ends the single-potential pinned description. It does not erase the Lagrangian source label or reroute the committed allocation: source-labelled multistream stock remains conserved until the separate carrier-loss condition $\mathsf E_A=1$ is met. The physical pinned branch is therefore
\[
\boxed{\sigma_A=1\ \text{for every recruited source},
\qquad \mathcal C=1,}
\]
and the subsequent sequence is
\[
\boxed{
\begin{gathered}
\text{coherent pinned committed capacity}
\\[-1mm]\downarrow\\[-1mm]
\text{source-labelled multistream committed stock}
\\[-1mm]\downarrow\\[-1mm]
\text{carrier loss and rerouting}
\end{gathered}}
\]
generally with the absolute availability still close to its vacuum value. Terminal strong-field saturation is the different condition $q_{\rm geo}=0$. Cosmological pinning saturates a small active subsector; it does not exhaust the vacuum budget.

\paragraph{The pinned reading and the zero-pressure constraint theorem.}
The \emph{pinned reading} is the statement that the saturated bath holds exactly at the source-normalized ceiling and stays there. The transfer law quantizes channel transport at one bit per tick as an upper bound; saturation is the attainment of that bound, so on the pinned reading every committed channel advances by exactly one unit per tick, and while the phase remains pinned no channel can decommit, since there is no slack for it to relax into. The committed count is then conserved and dilutes only with the expanding volume.

The unit-timelike normalization is not a further independent condition once saturation and coherence are granted. Let $\theta$ be the coherent renewal phase, increasing by one on each electron-calibrated interval $\tau_*$. The dimensionful clock $\phi=\tau_*\theta$ therefore satisfies $d\phi/d\tau=1$ along a committed worldline. Coherence identifies the flow covector as $u_\mu=-\nabla_\mu\phi$. Hence
\[
1=\frac{d\phi}{d\tau}
=u^\mu\nabla_\mu\phi
=-u^\mu u_\mu,
\qquad
\boxed{X=-\frac12\nabla_\mu\phi\nabla^\mu\phi=\frac12.}
\]
Thus the spectral tick fixes the normalization of the coherent clock; no cosmological coefficient is inserted. This proves the constraint within the pinned coherent branch. It does not derive the transition that makes $\sigma_A\to1$ and $\mathcal C\to1$.

The coarse-grained field $\phi$ consequently carries a constraint rather than a generic kinetic term. The leading action consistent with it is
\[
S_{\rm comm}=\int d^4x\,\sqrt{-g}\;\lambda\left(X-\tfrac12\right),
\]
whose variation gives $\nabla_\mu(\lambda\,\partial^\mu\phi)=0$ and, on the constraint surface,
\[
T_{\mu\nu}=\rho\,u_\mu u_\nu,\qquad \rho=E_0\,n,\qquad p=0,\qquad c_s^2=0,\qquad \rho\propto a^{-3},
\]
with $u_\mu=-\partial_\mu\phi$ and $n$ the conserved spatial density of committed cells. This belongs to the constrained-scalar class \citep{ChamseddineMukhanov2013}. Conditional on the pinned reading, this is the zero-pressure stress form of saturated committed capacity. It is an effective metric statement about the pinned capacity phase. The construction inherits caustic formation and the need for a small-scale completion. The manuscript does not derive the transition from the unconstrained weak-field branch onto $X=\tfrac12$ or the commit/decommit energy accounting. Related relativistic Milgromian dynamics provides a useful comparison \citep{Skordis2021}.

\paragraph{Carrier constraints.}
The constrained carrier supplies the conserved density needed for cosmological growth (Appendix~M). The free-fall growth clock, calibrated on the cluster radial decline, misses cosmological development by a factor of order thirty. A first-order relaxation clock cannot attain the acoustic rejection below one part in five hundred while settling over the available pre-recombination interval. Bound-type rail readings sit at the sound-speed pole and carry no perturbation; plateau approaches have $c_s^2=-1/(2n{+}1)<0$ and are gradient-unstable; generic unconstrained scalar continuations have nonzero sound speed and are not adopted as post-release degrees of freedom. The hard-cap argument in Appendix~M applies when the cell weight is the only state variable. The constraint reading supplies the additional conserved committed density. On the undepleted branch, carrier loss removes the source-label saturation condition and returns the metric-visible remainder to the ordinary source-associated capacity bookkeeping described in Appendix~O. Near depletion, the selected constitutive completion in O.3 replaces the constrained dust limit by a pressure-supported branch with a constant energy floor.

\paragraph{Abundance.}
At full saturation the commitment count gives the pre-release ratio
\[
\frac{\Omega_c^{\rm wb}}{\Omega_b}=\frac1\epsilon
=\frac{4\pi^2}{g_{\mathrm{share,eff}}}=5.320684854
\]
in the bare baryonic-mass convention. The free-nucleon and helium
conversion in Appendix~O.15 gives the formation value $5.333$, fixed by the O.15 map and held fixed in
the cosmological fit. Release reduces the actual present dust-to-baryon
ratio to about $5.2284$ for the reference history under the composition-preserving allocation of Appendix~O.3.

With the formation abundance and $\eta_*$ fixed, the assembly-and-mixing model has five cosmological parameters against six for $\Lambda$CDM. Direct geodesic perturbations with matched lensing give $\Delta\chi^2=+0.85$ against $\Lambda$CDM for Planck 2018, DESI DR2 BAO and DES-Dovekie with $\sum m_\nu=0.06\,$eV, and $-0.39$ for the configuration with massless neutrinos, fixed Planck calibration, nonlinear lensing and a free supernova offset below $z=0.1$ (Appendix~O.12). The fits are statistically tied, and the parameter count gives $\Delta{\rm AIC}=-1.15$. With the retention gate from feedback-resolved simulations, the measured kernels and the reionization redshift of the fitted optical depth, the release history has no free parameter; that frozen history leaves four cosmological parameters and gives $\Delta\chi^2=+1.07$ in the first configuration and $-0.28$ in the second. Appendix~O.12 specifies the matched likelihood and distinguishes the linear-depletion PPF integral $I_{\rm rel}=94.2116$ used by the hydrodynamic comparison. Appendix~O.12 takes the retention gate from feedback-resolved simulations. The independently measured mixing kernel, nonlinear response, and separate sound-horizon proposal remain physical tests.

The acceleration and abundance relations contain a parameter-free structural test that does not depend on the numerical value of the lever coefficient. Since $a_0=\epsilon cH_0$ and $\Omega_c/\Omega_b=1/\epsilon$,
\[
\boxed{
a_0\frac{\Omega_c}{\Omega_b}=cH_0
}.
\]
This identity uses the bare pre-release abundance. The formation mass
conversion and subsequent depletion must be applied when comparing with
the actual present dust ratio. A joint test must also use the same
$H_0$ in the galaxy and cosmology sectors.

Define the mass-weighted formation saturation by
\[
\bar\sigma:=\frac{\rho_c}{\rho_b/\epsilon}\leq1.
\]
The bare pre-release ratio is $\bar\sigma/\epsilon$. Appendix~O.15
specifies a formation map that fills the cap, $\bar\sigma=1$, under its
pinned-reading and ledger assumptions; the microscopic formation vertex
remains open.

That the committed density attains the ceiling exactly follows only within the recruitment assumptions of this branch: in Appendix~O.15 one marked recurrence exceeds the required formation weight by a factor of $2.1\times10^{29}$, so the cap fills within a small fraction of a Hubble time. In the notation above, the assumption is $\sigma_A\to1$ for every recruited Lagrangian source element. Commitment is source-driven: cells receiving no demand commit nothing, and each source's allocation is capped at $1/\epsilon$ per unit source mass. The recombination-era estimate above is a consistency condition: there the relevant modes demand $\nu\gtrsim10^2$, far above the cap, while the homogeneous background remains below it. With abundant supply and irreversible commitment in the pinned regime, each source recruits to its cap and allocations are assumed to add. The committed density is then
\[
\rho_c=\frac{1}{\epsilon}\,\rho_b
\]
pointwise at a commitment epoch. Appendix~O.6 registers a first-crossing computation that falsifies a late pointwise-tracking reading: the Silk-damped committed field never reaches a caustic. The adopted formation map supplies commitment before the acoustic epochs. Appendix~O.15 supplies the formation map and its epoch. Formation acts on conserved baryonic charge when that charge first admits a source label, between sphaleron freeze-out and the QCD crossover, and draws the energy $\rho_b/\epsilon$ from the plasma. Every mode used in the CMB and Lyman-$\alpha$ likelihoods is superhorizon then, so the formation slice imposes $\delta_c=\delta_b$, and the recombination-era demand estimate above becomes a consistency condition. The abundance remains conditional on the pinned reading and per-defect bookkeeping, and the formation vertex itself is not derived.


\paragraph{Source-clock conservation and release.}
The saturated committed-capacity phase remains source-associated while its
baryonic Lagrangian carrier is distinguishable. For each source $A$, the
phase $\theta_A$ is the renewal clock on the source worldline; $c_{A,c}$ is the
identity-free committed allocation. Shell crossing can create several
Eulerian streams, but it cannot create a which-stream record in the
one-register capacity action: exchanging identical stream assignments leaves
the occupancy history unchanged, so
\[
|e_i\rangle=|e_j\rangle,\qquad
\gamma_{ij}=1,\qquad C_A^{\rm rec}=1.
\]
Carrier loss occurs only when matter itself ceases to preserve the source
clock: gas from several progenitors becomes irreversibly mixed in one
retained common well and baryonic records no longer resolve progenitor
membership. For each elementary record block that loses its carrier,
faithful renewal fixes a complete fresh draw,
\[
\boxed{
\mathsf E_A=\mathbf1[\text{the baryonic carrier of }A
\text{ has lost progenitor identity}],\qquad f_{\rm reroute}=1.}
\]
The normalized nine-state dilation of Appendix~O.13 is isometric; its $1/3$
marked amplitude normalizes nine outputs and is not a routing fraction.
At coarse resolution, Mass--Entropy Equivalence and additive record pricing
assign released stock in proportion to the source-label information no longer
recoverable from matter. Thus partial $\epsilon_{\rm mix}$ describes a
fraction of fully renewed record blocks, rather than a partial-memory update
of any one failed block.

The enabling well depth remains a collective matter condition rather than a
microscopic capacity threshold. Atomic-cooling assembly supplies
$V_{\rm thr}\sim17\,{\rm km\,s^{-1}}$ before reionization, while photoheated
gas after reionization requires a scale nearer $30\,{\rm km\,s^{-1}}$.
The primitive action contains no local switch at either value. The transverse
scale $a_0(z)=\epsilon cH(z)$ and the source-clock release gate therefore have
different roles.

Appendix~O measures carrier loss by
\[
\epsilon_{\rm mix}(t)=1-\frac{I(A;Y_t)}{H(A)},
\]
with $A$ the Lagrangian progenitor label and $Y_t$ the retained coarse
baryonic carrier. First eligible assembly starts the mixing clock rather than
causing instantaneous release. The cosmological history is therefore
\[
\boxed{F_{\rm rel}(t)=\int^t K_{\rm mix}(t-t_f;M,z)\,dF_Z(t_f).}
\]
Appendix~O.12 fixes $F_Z$ by an absorbing first-entry excursion-set
construction. The PPF benchmark fits the pure-exponential mixing time $f=0.282706$ at the reference retention gates; the matched direct-perturbation fit with the composition-preserving allocation gives $f=0.274245$. A self-gravitating
progenitor-tagged run measures a different kernel with nonzero loss at
entry and a slower tail. Its convolved release integral reaches the
fitted target near a $14\,{\rm km\,s^{-1}}$ gate, while the reference
$17\,{\rm km\,s^{-1}}$ gate undershoots it. A retained well keeps gas that can cool and in which progenitor labels
mix, the criterion that also sets the post-reionization gate. The halo
occupation of feedback-resolved simulations measures that retention, and with the measured kernel it gives $\Omega_{\Lambda,0}=0.709\pm0.02$ against the $0.707$ selected by the data (Appendix~O.12).

Carrier loss does not create a new material component. The replaceable
closure register returns to the universal renewed state, while the
metric-visible remainder stays associated with the retained coarse carrier.
The marked history output has no independent spacetime current. The
ensemble mean of its drain fixes the vacuum potential on the capacity clock.
The constant-mode projection identifies its background normalization, while
realized release fluctuations redistribute the source-resolved and
reallocated ledgers locally. Their sum gives one metric-visible capacity
density,
\[
\boxed{\rho_C(a)=\rho_{c,0}^{\rm wb}a^{-3}
e^{-\eta_*F_{\rm rel}(a)}},
\]
together with
\[
\boxed{\rho_\Lambda(a)=\eta_*\rho_{c,0}^{\rm wb}
\int^a a'^{-3}e^{-\eta_*F_{\rm rel}(a')}\,dF_{\rm rel}(a').}
\]
The local vacuum density is $V_F(\phi)$, with its perturbation fixed by the
capacity clock. Appendix~O.2--O.3 gives the ensemble prescription, local
exchange, and stock-positivity domain.



\paragraph{Two consequences of the gate for other sectors.}
The release condition is collective and gated by retention, so a halo that
never crosses the gate never undergoes source-clock release. A source whose
gas is expelled before irreversible progenitor mixing retains its source label and its initial allocation of $1/\epsilon$ per unit Lagrangian baryonic mass, subject to the shared cosmological drain. Absence of local release does not remove that prescribed drain. That allocation remains in the saturated source-associated phase. It gravitates as committed collisionless capacity and adds no component to the response; it becomes response stock of Appendix~O.3 through the collective carrier-loss channel that gate crossing opens. For this reason Section~15 needs no additional collisionless matter species. The Local Group dwarf masses are consistent with this branch (Section~15): dwarfs whose progenitors stayed below the gate keep CDM-like stock, and dwarfs above it relax to the response law where they have their own support. The release history fixes the gate value, and the dwarf data leave it open.

\paragraph{Distinction from horizon saturation.}
The committed cosmological phase (channels advancing at capacity, a running clock) is distinct from the terminal saturation of the strong-field sector (Section~21), where capacity is exhausted and transactions cease. Whether the two termini are stages of one process is open and carries a stated consistency burden: the gravitating energy of committed capacity must coincide with the mass already accounted at infinity, with no double counting. This is recorded as an open question shared by the strong-field and cosmological sectors.


\section*{Part VI. Strong Fields, Many-Pasts, and Microstructure}
\addcontentsline{toc}{section}{Part VI. Strong Fields, Many-Pasts, and Microstructure}

The strong-field branch treats the zero-capacity boundary, and Many-Pasts supplies the record-conditioned history ontology in which the substrate's renewal process operates. The separate refresh theorem selects the memoryless dressing used in scale setting, while Section~23 asks for its microscopic dynamics. The same grading applies here: the exact spherical reduction, the operational quantum measure, and the one-tick renewal H-theorem (Appendix~P.8) are controlled results, while the boundary microphysics and the condensate embedding carry their stated conditional or open grades.

\section{21. Strong-Field Action: Spherical Closure and Its Boundary}

The bounded variable remains the natural strong-field order parameter,
\[
q=\frac{S_{\mathrm{ent}}}{S_\infty}\in[0,1].
\]
Serial composition and weak-field matching select $N^2=q$ in a static exterior. That relation is meaningful after the asymptotic Killing time has fixed the normalization of the lapse; it is not, by itself, a covariant field equation in a general foliation.

\paragraph{Why a multiplier action is not the parent.}
An exploratory ADM term $\sqrt h\,\lambda(N^2-q)$ makes the problem look variational, but if $q$ has no independent bulk dynamics or matter coupling its variation gives $\lambda=0$. The remaining constraint merely renames the lapse, which is gauge dependent, and supplies neither the capacity Poisson equation nor a new covariant relation. Giving $q$ its own kinetic term would instead add a scalar degree of freedom and reopen the fifth-force, stability, and double-counting problems. The multiplier construction is therefore a diagnostic control case, not the parent action.

\paragraph{Exact spherical reduction.}
Spherical symmetry supplies an invariant radial observable. It agrees with clock capacity in normalized static vacuum; N.6 derives their distinction inside matter. Write
\[
ds^2=h_{ab}(x)dx^adx^b+R^2(x)d\Omega^2,
\qquad a,b\in\{0,1\},
\qquad x^0=ct.
\]
After integrating the Einstein--Hilbert plus GHY action over the two-sphere and removing a two-dimensional total derivative, the bulk action is
\[
I_{\mathrm{sph}}
=\frac{c^3}{4G}\int d^2x\,\sqrt{-h}
\left[
R^2\,{}^{(2)}R+2h^{ab}\partial_aR\partial_bR+2
\right]
+I_{\mathrm{matter}}^{(2)}+I_{\partial}^{(2)}.
\]
Its vacuum equations imply conservation of the Misner--Sharp mass
\[
M_{\mathrm{MS}}
=\frac{c^2R}{2G}
\left(1-h^{ab}\partial_aR\partial_bR\right),
\qquad
\nabla_aM_{\mathrm{MS}}=0.
\]
The radial geometric invariant is therefore
\[
\boxed{
q_{\mathrm{geo}}
\equiv h^{ab}\partial_aR\partial_bR
=1-\frac{2GM_{\mathrm{MS}}}{c^2R}.
}
\]
The radial quantity $q_{\mathrm{geo}}=q_R$ is a composite geometric scalar whose vacuum profile is a first integral of the metric equations. It introduces no additional propagating mode. Its identification with the clock quantity $q_\tau=N^2$ uses the static vacuum equations and reference normalization.

For a static asymptotically flat vacuum, $M_{\mathrm{MS}}=M$ and Birkhoff's theorem gives
\[
ds^2
=-\left(1-\frac{2GM}{c^2r}\right)c^2dt^2
+\left(1-\frac{2GM}{c^2r}\right)^{-1}dr^2+r^2d\Omega^2,
\]
so
\[
q_{\mathrm{geo}}=N^2
=1-\frac{2GM}{c^2r}.
\]
This is the nonperturbative action-level realization of the bounded capacity relation in the strongest sector where an invariant local definition is currently available.

\paragraph{The $q=0$ surface.}
The equation $q_{\mathrm{geo}}=0$ locates a marginal sphere. A well-posed exterior variational problem is obtained first on a stretched timelike boundary $q=\epsilon>0$ with the usual GHY term and fixed induced data, and then by taking the null limit with the corresponding null and joint terms. The Einstein action determines that universal gravitational boundary bookkeeping. It does not determine a new substrate boundary Hamiltonian, microscopic reflectivity, or the rule that excises $q<0$.

Accordingly, the restriction
\[
\mathcal M_q=\{q_{\mathrm{geo}}>0\}
\]
is a physical postulate about the domain of the capacity EFT, not a consequence of varying the Einstein action. The exterior solution and horizon location are closed within spherical metric reduction; the claim that the classical interior is absent, and the dynamics experienced at the saturation surface, remain conditional on a microscopic boundary theory. Appendix~F and Appendix~N state this boundary between result and interpretation explicitly.

\paragraph{Status of the construction.}
The Schwarzschild exterior, its standard Hawking temperature, and ordinary exterior perturbation equations follow from the metric parent. In the factorized $j=3$ completion, the one-bit cut and coherent EPRL area fix the channel area and Immirzi normalization (Appendix~F.5); the Green-response $1/4$ identity restates the C.5 matching. Rotating and charged GR exteriors remain valid solutions of the baseline parent action, but a covariant capacity scalar that identifies their saturation surface has not yet been constructed. The universal strong-field result is therefore spherical and exterior; boundary microphysics, dynamical saturation, nonspherical capacity geometry, and any controlled departure from GR remain open.


\section{22. Many-Pasts: The History-Space Ontology}

Many-Pasts makes a conservative operational claim and a distinct ontological claim. Laboratory records obey standard quantum mechanics. The ontology assigns a conditional measure to the decoherent past histories compatible with the one realized present.

\paragraph{The operational weight.}
For a decoherent family, the joint history weight and its record-conditioned form are
\[
p(h,P)=\mathcal D(h,h),
\qquad
p(h\mid P)=\frac{\mathcal D(h,h)}{\sum_{h'\in\mathcal H_P}\mathcal D(h',h')}.
\]
Equivalently, one may define $D(h,P)=-\ln p(h,P)$ on the support of the decoherent joint measure, so that $p(h\mid P)\propto e^{-D(h,P)}$. The exponential is a reparameterization of a normalized quantum probability, not an additional classical measure placed on interfering paths. Appendix~H.11 derives the finite retained-sector norm, instruments, and history Gram kernel; Appendix~G applies them to Many-Pasts and gives the no-signaling calculation. The global cosmological state and continuum realization remain separate inputs.

\paragraph{Branch realization without many worlds.}
The weight also answers what it means for one outcome to be realized. There is no forward branching into co-real worlds and no collapse event. A definite present is a present with definite macroscopic records, and the histories the weight supports are exactly those compatible with those records; alternative outcomes correspond to alternative present records, not to coexisting branches. Probability is the measure this weighting assigns over the admissible pasts of the one realized present.

\paragraph{Record retention and the resolution of the past.}
The physical degrees of freedom that distinguish histories are present records. A durable record can be a detector state, an environmental correlation, or the exported history register of the renewal dilation. Such a record keeps distinct those past alternatives that remain distinguishable in the present record algebra. If orthogonal fine records $R_i$ are later merged into a coarser record $\bar R=\sum_{i\in I}R_i$, decoherence and additivity give
\[
p(h,\bar R)=\sum_{i\in I}p(h,R_i),
\qquad
p(h\mid\bar R)=
\frac{\sum_{i\in I}p(h,R_i)}{\sum_{h'}\sum_{i\in I}p(h',R_i)}.
\]
The history measure therefore loses resolution by ordinary quantum coarse-graining when the present loses a record. The ontology contains no separate archival copy of a distinction that has disappeared from every physical register. This does not erase consequences already carried into the current state. It identifies earlier alternatives with identical surviving records as the same present physical state.

\paragraph{No future pull.}
The class operators and unitary maps determine the joint measure before conditioning. The factor $p(h\mid P)$ updates the description of earlier alternatives after $P$ is recorded; it does not enter the Hamiltonian, the channel, or the class operator that produced $P$. A record that has not yet formed supplies no conditioning event in the realized present. Many-Pasts therefore adds no retrocausal force, final-boundary dynamics, or Born-rule bias. It gives a physical reading to retrodictive conditioning already present in standard quantum mechanics.

\paragraph{Familiar quantum examples.}
In a double-slit experiment, the alternatives through the two slits remain inside one coarse class operator until a durable which-path record decoheres them. Their cross terms are therefore retained before the record and suppressed after it. In an EPR or Bell experiment, the present is a joint record of the pair and the detectors. The joint measure gives the usual nonclassical correlations, while the local marginals remain independent of the remote setting. Measurement creates a stable record and thereby identifies the decoherent family on which conditional probabilities can be used. These are standard quantum calculations with a record-conditioned history-space reading.

\paragraph{The arrow of time.}
The same framework orients time through conditional typicality. At stationarity the maximum-caliber replacement process obeys detailed balance identically, $p(b)K_*(b,b')=p(b')K_*(b',b)=p(b)p(b')$, and the reversible dilation has an inverse. The orientation enters through the structure of the reduced dynamics: each tick consumes a fresh register and exports the previous one, so renewal is an exact relative-entropy H-theorem on the replaceable register while coherent sectors evolve unitarily (Appendix~P.8). Appendix~P.9 derives the minimal-record boundary as the root of the record order. Histories with a low-entropy past and increasing future entropy then dominate the conditioned ensemble of Appendix~G.7. The result adds no law of laboratory probability.

Bianconi supplies a useful compatibility result for that theorem. In her low-curvature matter- and radiation-dominated Friedmann approximations, the local geometric-relative-entropy density decreases as the universe dilutes while the integrated entropy grows with the expanding volume and the integrated energy approaches a constant \citep{Bianconi2026}. This demonstrates in a concrete information-geometric action that local ordering and a global entropy increase need not conflict. It does not select a low-entropy boundary, establish substrate mixing, or show that the Many-Pasts conditional measure favors ordinary histories over Boltzmann fluctuations. Appendices~P.8--P.10 supply the boundary, the record-level contraction and late-time stability within the present construction; relaxation of a given recorded subsystem requires its own calculation.

\paragraph{No external cycle ledger.}
If a proposed cosmological evolution reaches a state whose present record algebra contains no witness of an earlier macroscopic era, the descriptions ``first beginning'' and ``return to the same beginning'' are not distinct states within this ontology. Distinguishing them would require a register that survives the record-free interval or an external time parameter that counts passages. Many-Pasts supplies neither. This conditional observation does not derive a cyclic universe, a unique beginning, or a nucleation rate. It removes an otherwise hidden meta-history from the ontology.

\paragraph{Where the weight is used.}
Many-Pasts supplies the history-space realization in which the electron dressing operates, while faithful full-support resolution proves the independence of successive passes. Appendix~H shows that this condition is equivalent to maximum path entropy and uniquely selects the replacement process; the lightest-defect criterion selects the same temporal independence together with independent channel layers. The record-conditioned viewpoint also enters the proposed macroscopic arrow of time and the conditional cosmological extensions. It changes no laboratory law.

\paragraph{Local renewal and reversible history export.}
Complete local renewal has a precise quantum form. If the new cell must contain no information about the old cell even when the old cell is entangled with a reference, the one-tick channel is uniquely
\[
\mathcal E_*(\rho)=\rho_*\operatorname{Tr}\rho .
\]
This local channel need not destroy information globally. A reversible dilation transfers the old cell into an environment register while a fresh admissible register becomes the new present. Many-Pasts interprets the exported correlations as history degrees of freedom. The identification is structural: the finite dilation follows directly from the retained tensor factors, while Postulate~III supplies its history-space reading. It does not select the cosmological initial state or prove that the full continuum histories decohere. Exact recoverability forces the stored history to grow without bound: destructive reuse of history capacity contradicts the dilation, and the history already written exceeds the area-law capacity of a de~Sitter patch (Appendix~P.9).

\paragraph{Coherence under renewal.}
The replacement channel acts only on the renewed closure register. Its Heisenberg dual is
\[
\mathcal E_*^\dagger(O)=\operatorname{Tr}(\rho_*O)I,
\]
so no later observable of that register can recover an input off-diagonal. The global dilation can nevertheless retain coherent information in its complement. If two spatial branches export discarded states $|e_x\rangle$ and $|e_y\rangle$, their reduced off-diagonal is multiplied by $\langle e_y|e_x\rangle$. A cell-addressed mark record makes those states orthogonal and would destroy position coherence in one update. The allowed free vertex instead transports the marked fiber as part of the coherent defect system and leaves the discarded renewal record branch-independent. Ordinary environmental interactions may then reduce the overlap in the usual way. Appendix~H.11 gives the proof and distinguishes this condition from the weaker statement that a record merely ``moves with'' a worldline.

\paragraph{The retained quantum sector.}
Appendix~H.11 settles the local question within the selected quadratic
completion. Unitarity identifies the nine-state fusion complement with the
event response coordinates of the persistent marked fiber. On its real form,
the positive Lorentzian Hamiltonian selects an initial polar complex
structure. The full symplectic flow transports that structure exactly through
arbitrarily varying frames; squeezing relative to the instantaneous polar
basis and the induced Born--Huang term are retained rather than discarded.
The non-common frequency and moving frames generate $\mathfrak{su}(9)$ and
give $\mathcal H_{\rm mark}\simeq\mathbb C^9$. The complete fixed-one-mark
bilinear response algebra identifies normalized states exactly up to common
phase, yielding $\mathbb{CP}^8$ and its Fubini--Study geometry. Within the
H.9 Gaussian-strand realization, the fresh tangent metric also equals one
quarter of the Fisher metric, conditional on identifying its Hermite response
coordinates with the closure-score coordinates. Caticha's Hamilton--Killing
theorem is therefore a consistency and continuum guide rather than the
premise carrying the local reconstruction
\citep{CatichaHK2021,CatichaQED2025}.

\paragraph{Status.}
Finite operational quantum mechanics is derived on the retained sector.
The real positive quadratic dynamics gives the complex Hilbert bundle and
exact unitary transport; complete commuting local algebras give tensor
composition; reversible record coupling gives POVMs, instruments, conditional
updates, and reduced channels; reversible equivalence plus physical-record
additivity fixes the Born norm; and the history functional is the resulting
Gram kernel. The projected two-mark bridge generates $\mathfrak{su}(81)$ provided the multiplicative two-mark carrier is retained; the geometric gauge projection acts as the identity on the capacity factor (Q.4), so the bridge survives it (Appendix~P.8). Extension to a connected marked network additionally requires simultaneous retained marks. Renewal gives an exact one-tick H-theorem on the replaceable register, and the persistent geometry--defect sector evolves unitarily under it, so these coherent results and the thermodynamic arrow of Appendix~P.8 hold together. The global continuum/Fock realization, cosmological state, identification of the histories that decohere, the physical source of history capacity, and electromagnetic charge and dynamics retain their open or conditional grades.


\section{23. Microstructure Hamiltonian and Underlying Dynamics}

The UV closure chain has an explicit finite action on the scale-setting side and a controlled geometric witness at the semiclassical level. Appendix~H derives the replacement process, its quantum channel, the state-weighted determinant transfer, and the decorated native-cell vertex used below.

The selected factorized geometric completion is not obtained by identifying the finite capacity alphabet with the magnetic indices of the geometric spin network. The microscopic one-cell space is instead taken to factor as
\[
\boxed{
\mathcal H_{\rm cell}
=\mathcal H_{\rm geom}
\otimes\mathcal H_{\rm cap}
\otimes\mathbb C^2_{\rm orient},
}
\]
Here $\mathcal H_{\rm geom}$ carries the gauge-reduced simplicial
GFT/spin-foam data, while $\mathcal H_{\rm cap}$ carries the finite routed
capacity alphabet and its marks. Appendix~Q shows that this factorization
removes the Gauss-Casimir contradiction without changing $\Omega=1680$, the
$K^2$ spectrum, $\eta_*$, $g_{\rm share,eff}$, $\zeta_*$, $Z_e$, $L_*$, or
$G_*$. An equivariant relational-frame lock relates the two selected copies
of $V_3$. Once the lock exists, Schur's lemma fixes it up to phase. The two
factors are bookkeeping components of one substrate, and Postulate~I
restricts physical configurations to their locked sector.

On the capacity factor, Appendix~B proves that positive oriented matching has $V_3$ as its unique zero-cost transmitted sector and that the displayed homogeneous action has a unique $J=3$ minimum. Maximal fusion then gives the exact tree-independent canonical blocking ray on the selected maximal-channel branch
\[
F_N:V_3^{\otimes N}\to V_{3N},
\qquad
F_N|3,\mathbf n\rangle^{\otimes N}=|3N,\mathbf n\rangle.
\]
The microscopic representation stays fixed while blocks enter the large-spin
regime along this canonical ray. The face mismatch is the isotropic ferromagnetic
Heisenberg bond shifted by a constant. On hypercubic hosts in $d\ge3$ the
nearest-neighbor model of this class has a spontaneously ordered phase at low
positive temperature for every spin \citep{OpenAIHeisenberg2026}, which supports
local selection of the aligned branch. Exact local Gauss projection removes the
single-face bond between physical cells, so ordering of the gauge-invariant
shape variables requires multi-face, holonomy, or vertex terms
(Appendices~B and H.10). Several external constructions support the
geometric limit. Proper-EPRL amplitudes remove unwanted asymptotic sectors at
the vertex level and retain the Regge exponential
\citep{EngleProper2013,EngleVilenskyZipfel2016}. Han-type spin-foam models
provide a coupled large-spin/refinement witness to the Einstein sector; in
the controlled linearized construction, the low-energy excitations exhaust
the smooth linearized Einstein solutions and the two graviton helicities
\citep{Han2017,HanHuangZipfel2019}. Area-Regge continuum analysis gives the
same leading graviton dynamics with a first length-metric correction of order
$a^2C^2$ \citep{DittrichKogios2023}. Applying the Han-type refinement result
to the decorated proper-EPRL factor remains a coupled-embedding audit.

The scalar capacity deficit is a coordinate on the Hamiltonian-constraint
response, not an additional radiative gravitational field. Gauge-reduced
vacuum propagation remains transverse-traceless. A GFT condensate may still
realize the microscopic state and clock \citep{Oriti2016,Gielen2013}, but
Einstein propagation does not require a separate light scalar condensate
mode. Appendix~H.10 lifts any geometric rigging map through the finite
capacity factor and proves clock-test positivity and the causal boundary
value from strong-resolvent convergence.  On the selected regular
metric-Regge branch, Appendix~H.10 proves \((H_{\rm TT})\): the Einstein TT
infrared limit, together with an exact fixed-$j=3$ primitive realization whose
nonmaximal complement remains uniformly gapped.  It also proves that the
present capacity premises leave the geometric gluing tensor and measure
underdetermined. The coherent angular and count construction of
Section~23.1 specifies additional local geometric data. Its physical
constraint space and nonperturbative Einstein-phase selection remain
the tests of that candidate.  A
spin-$3$-regularized canonical host obeys the distinct condition
$(H'_{\rm can})$ on its zero fiber; identifying the two physical sectors
would require a separate theorem and is not assumed here.

On the defect side, the closure ensemble fixes the stationary marginal $p_{\eta_*}$ and maximum path entropy selects
\[
K_*(b,b')=p_{\eta_*}(b').
\]
This fixes the dimensionless history process but cannot produce seconds. The electron, already the theory's single dimensionful anchor and its lightest charged one-bit defect, supplies the clock through the positive transfer spectrum below.

Requiring the output cell to decouple from every reference system forces the replacement channel $\mathcal E_*(\rho)=\rho_*\operatorname{Tr}\rho$. The decorated vertex chooses the diagonal maximum-entropy completion
\[
\rho_*=\sum_b p_*(b)|b\rangle\langle b|,
\qquad
p_*(b)=Z^{-1}e^{-\eta_*K^2(b)},
\]
and prepares it from the amplitude $A_*(b)=\sqrt{p_*(b)}$. A reversible register permutation moves the old cell into history and the fresh amplitude into the present. For the complete $1680$-state tetrahedral register this is one native circuit layer. A direct circuit made only from disjoint face comparisons still requires the three perfect matchings of the four-face graph; the decorated action uses the complete tetrahedron as its local gate.

On the factorized seven-channel history space let
\[
|v^{(7)}\rangle=\bigotimes_{m=-3}^{3}|\sqrt{p_{\eta_*}}\rangle_m,
\qquad
P^{(7)}=|v^{(7)}\rangle\langle v^{(7)}|.
\]
On the commutative cell algebra, likelihood multiplication by $p_*(b)$ has the state-weighted determinant
\[
\Delta_{\tau_p}(\mathsf R)
=\exp\!\left(\sum_bp_b\ln p_b\right)
=e^{-g_{\mathrm{share,eff}}}.
\]
The seven-channel determinant is $r=e^{-7g_{\mathrm{share,eff}}}$. It is the almost-sure geometric transfer rate of the record-conditioned renewal history; the annealed equality probability would instead involve the collision entropy. Compressing the one-bit charged loop to its determinant line gives the positive scalar survival transfer
\[
T_{\rm surv}=1-r.
\]
The Euclidean transfer Hamiltonian
\[
H_{\rm surv}=-\frac{\hbar}{\tau_*}\ln T_{\rm surv}
\]
therefore has the exact raw gap
\[
E_{\rm raw}=-\frac{\hbar}{\tau_*}\ln(1-r).
\]
The closure amplitude has an exact Gaussian linearization. Projecting the same amplitude through the two directed scalar returns gives $u=8\eta_*/49$, and its canonical unitary dilation supplies the factor $(1-u)^{21/2}$. The nine present/history closure polarizations give the marked weight $\zeta_*$ of Section~13.4. The electron graph contributes $Z_e=(1+\zeta_*)(1+7\zeta_*^2)$. Assigning its motif weights to the charged generator as in H.9 gives the dressed transverse identification
\[
E_e=\frac32Z_eE_{\rm raw}.
\]
Identifying this lowest charged transfer gap with the electron rest energy fixes the update time,
\[
\boxed{
\tau_*=-\frac32Z_e\frac{\hbar}{m_ec^2}\ln(1-r)
}.
\]
Analytic continuation supplies the phase frequency $E_e/\hbar$. The invariant continuum speed converts the spectral cadence into the causal length $L_*=c\tau_*$. This definition does not determine the displacement made by a microscopic spatial update.

The causal length of that spectral update is
\[
\boxed{
L_*=c\tau_*=-\frac32Z_e\lambda_e\ln\!\left(1-e^{-7g_{\mathrm{share,eff}}}\right)
}.
\]
The baseline rare-event correction $-\ln(1-r)=r[1+O(r)]$ differs from $r$ only at order $10^{-23}$. The finite marked response changes the scale by $Z_e-1=0.00530828$, which is the physically relevant correction resolved by the microscopic vertex.

Retained temporal information raises the recurrence mass by $e^{I_t}$, so electron lightness independently selects $I_t=0$. Fermionic exclusion caps the occupied channel count at seven and subadditivity gives $\Delta_k\ge0$, hence the lightest resolved one-bit defect has $k=7$ and $\Delta_7=0$ at fixed per-channel marginal. The unit mixing gap, charged survival gap, and geometric spin-2 Hessian remain distinct objects. The selected transfer action, including its generator prescription, fixes the first two and the marked edge Hessian. In the ordinary geometric branch the source-coupled scalar response must be the nondynamical Hamiltonian-constraint combination; any additional scalar fluctuation of a particular condensate realization must be gapped, relaxational, gauge/constraint, or source-orthogonal. The physical radiative support is the two TT graviton modes.

The scale-setting branch combines the fixed renewal marginal, determinant-survival action, finite marked enumeration, additive-generator prescription, and electron calibration. Positivity gives the energy of the resulting transfer. The complete edge-Hessian and graph-enumeration audits are finite and reproducible. Long-time history capacity and a first-principles derivation of the carrier-resolved metric contact remain open; the ordinary semiclassical Einstein branch is supplied independently by the factorized spin-foam/GFT geometric sector.

The cell-order calculations show that the closure weight, refresh, history tilts, and conserved capacity field do not supply vacuum curvature stiffness. Under the factorized completion this is expected rather than a missing scalar mechanism: vacuum spin-2 stiffness belongs to the geometric factor, while the finite capacity sector supplies the source, records, and coefficient chain. The capacity variables do couple geometry to persistent closure failure, and Section~24 measures that local response. The CDT simulations remain an independent numerical host/interface rather than a derivation of the EPRL refinement limit. Postulate~III conditions histories within a specified state and does not select the nonperturbative gravitational vacuum.

\subsection{23.1 Coherent microscopic geometry and retained matter}
\label{aqg:main}

The local capacity ensemble admits a spatially entangled vacuum
without changing its \(1680\) probabilities. Its routing and
orientation representation has the equivariant encoding
\[
\mathcal H_{\rm cap,o}\cong\mathcal F(\mathbb C^4)
\otimes\mathbb C^{105},\qquad
\rho_{\rm cell}=\frac{I_{16}}{16}\otimes\sigma.
\]
Pairing neighboring virtual ports in
\((|10\rangle+|01\rangle)/\sqrt2\) preserves each local marginal
and gives \(\ln2\) of bipartite entanglement entropy per cut bond.
The residual entropy \(S(\sigma)\) remains local.

The absolute angular rule
\[
K_h(\theta)=e^{-\theta\ln2/(2\pi)}I_2,\qquad K_h(2\pi)=I_2/2
\]
gives both the replica trace \(2^{1-n}\) and the curvature weight
\(\exp[(\ln2/2\pi)N_h\delta_h]\).
With \(A_h=N_hA_{\rm ch}\), the bare Einstein coefficient obeys
\[
\frac{\hbar G_0}{c^3}=\frac{A_{\rm ch}}{4\ln2}.
\]
The full-turn rule is additional microscopic data. Its physical
area requires the event and volume calibration of C.8, which gives
the provisional \(G_0/G_*=1.03002\). The interacting coefficient
\(G_R\) must be measured in the same geometric and matter state.

Connected primitive patches realize the causal triangulation
weight explicitly. Integer channel bundles with area ratio
\(1:3\) give \(\alpha=7\). Four-cell chains and six-cell cycles
surject onto the required tetrahedral intertwiner spaces of
dimensions seven and nineteen. Full Gram normalization makes
their boundary maps isometric. Contracting the normalized
simplex packets then gives the selected Regge weight, including
boundary angles, on every allowed history. The patches encode a
coarse geometry on connected spin networks, and a primitive patch
need not be a literal tetrahedral subdivision.
The host measure, relabeling quotient, and volume prescription
remain explicit choices (Q.11b and J.12).

The canonical construction uses directional channel counts to
reconstruct one positive spatial metric. Normalized single and
paired reconnections implement its prescribed kinetic operator
on the microscopic switch patterns. The Einstein/ADM continuation
gives the same-coefficient Newton response, two-polarization
linear dynamics, and nonlinear Kasner evolution. Its DeWitt
measure and transported metric Weyl ordering specify the quantum
count operator. Continuous-configuration commutators quantize
the regulated classical brackets exactly; spatial and finite-count
errors remain distinct. A positive distributional rigging form
exists for the homogeneous volume--shape constraint (N.13).
The finite canonical histories have invariant signed cycles,
so they are not identified entry by entry with the positive
covariant history weight.
In the quadratic constrained response, the static source check returns
\(G/G_0=1\) only when the negative conformal direction is kept on its
rotated contour; deleting that direction gives \(4/3\), and the two
spatial TT modes alone give zero (N.14).

The retained defect also has a finite orbital realization.
On geometric cochains, \(D=d+d^\dagger\) gives
\[
\Omega(m)=\sqrt{D^2+m^2I},\qquad
\partial_m\Omega=m(D^2+m^2I)^{-1/2}.
\]
The native rest gap fixes \(m\); the mass derivative supplies the
integrated clock. A positive geometry transfer assembled from
enumerated slabs couples two boundary geometries, while harmonic
transport preserves the exact marked-sector rest gap.
An inverse fine-cochain metric intertwines the whole coarse
operator, its clock, and its evolution through the tested
refinements. Fixed-background renewal acts as the identity on
retained position and internal state and therefore preserves
interference (H.14--H.15).

The small-momentum cone is an additional test of this model.
Raw Whitney data split the directional speeds on a regular
periodic lattice; the selected inverse metric corrects the
leading cone on that family. A common relativistic limit across
changing boundaries and matter species, coherent renewal across
different geometric event records, and the interacting physical
constraint space remain to be established. The static capacity
equation in N.6 fixes the full lapse source within Einstein
gravity, while J.13 defines the separate measurement of \(G_R\).

At low curvature the common effective action should take the form
in units \(\hbar=c=1\),
\[
\Gamma[g,\psi]=\frac1{16\pi G_R}\int\sqrt{-g}\,(R-2\Lambda)
+\Gamma_m[g,\psi]+\Gamma_\partial+\sum_i c_i\mathcal O_i.
\]
The flat graviton pole leaves the higher-curvature coefficients
\(c_i\), physical boundary data, and interacting renormalized
parameters to be determined. A spacetime solution also requires
initial and boundary conditions.

\section*{Part VII. The Substrate on a Dynamical Lattice}
\addcontentsline{toc}{section}{Part VII. The Substrate on a Dynamical Lattice}

This part tests the cell ensemble numerically on a dynamical simplicial geometry, and then applies the resulting joint measure to the equilibrium vacuum and the cosmological term.

\section{24. Lattice Tests: Compatibility, Defect Response, and Transport}

The single-cell coefficient chain of Part~II can be evaluated in closed form. Its many-cell consequences require numerical tests: survival of the vacuum ensemble on fluctuating geometry, local strain around a persistent closure failure, and transport toward the $1/r$ deficit of Section~12. The calculation places the capacity ensemble with the finite collision penalty of Appendix~J on a dynamical simplicial geometry. Appendix~J gives the ensemble definitions, controls, and measured tables.

The measured ensembles in this section use their declared
capacity-decorated or face-wise host prescriptions. J.12 gives the
separate connected integer-channel host, whose vacuum weight is
realized in Q.11b. Its interacting continuum properties require
measurements with that same action and volume prescription.

That connected host links the microstructure to the lattice directly. The
connected primitive patches reproduce the causal-triangulation weight
exactly at \(\alpha=7\), with the couplings
\((\kappa_0,\Delta)=(2.202,0.326)\) fixed in J.12, so CDT at that point
is a candidate native measure of the theory. Sections~24.1--24.9 give
the face-wise and decorated-host results. Section~24.10 and
Appendix~J.14 report the completed connected-host normalization test
and distinguish it from the open matter measurement.

\subsection{24.1 The host geometry and the cell identification}

The host is a causal-dynamical-triangulations ensemble: four-dimensional triangulations of $S^1\times S^3$ weighted by the Regge action \citep{Regge1961}. CDT is used because it independently sustains an extended four-dimensional de~Sitter-like phase \citep{AmbjornJurkiewiczLoll2004,AmbjornGJL2012}. That phase is an external host datum for the substrate test, not a state selected by Many-Pasts.

The identification is immediate, and it respects a separation the theory requires. Each spatial tetrahedron of a slice hosts one cell's boundary data: its four triangles carry the labels $m=-3,\ldots,3$ of the $j_{\mathrm{eff}}=3$ sector of Section~5, each slice triangle is shared by exactly two cells, and the slice gluing supplies adjacency. Nothing else of the lattice enters. The theory's cell is its label configuration, a dimensionless unit of capacity; no size is ever assigned to it, and the granularity stays in the capacity, as Section~1.2 requires. The lattice simplices are regulator scaffolding, as they are throughout lattice gravity. The identification is nevertheless geometric in the sense that matters: the joint admissibility statistics of the labels depend on how the slice is glued, so the closure weighting can in principle tell one geometry from another, and that channel is the entire coupling between the theory and the host.

Two boundaries of the design are stated at the outset. First, the host supplies a dynamical, curvature-carrying simplicial geometry with the theory's cell structure, and nothing more is asked of it. The calculation does not assume the CDT ensemble is the theory's own vacuum, and it withholds any claim that the substrate itself derives four-dimensional emergence. At the tuned host couplings used below, increasing $N_{41}$ from about $42{,}000$ to $102{,}000$ raises both the Hausdorff estimator and the concentration of volume in an extended region, as expected when moving farther into that finite-volume phase; the order of the visible phase boundary and the joint continuum limit remain open (Appendix~J.1). Second, the weighting coupled to the host is the paper's own: the admissibility energy of Sections~5--6 at the closed point, with $\eta=\eta_*$ and injectivity enforced. It is not a spin-foam vertex amplitude. The distinction drawn in Section~29.8 was kept operational: a spin-foam amplitude was coupled first as a neighboring-theory control, and the failures exposed by that exercise supplied the control methodology used below (Appendix~J.3). The design's two-level closure structure, with geometric closure exact by construction in the host and label closure soft, is the architecture predicted by the factor separation of Appendix~Q: the exact level is carried by the geometric/gauge host and the soft level by the separate finite capacity factor.

\subsection{24.2 The coupled ensemble and its controls}

The coupled system is the joint Gibbs measure
\[
\pi(g,m)\ \propto\ \exp\!\Big(-S_{\rm Regge}(g)-\varepsilon\,(N_{41}-\bar N)^2-\beta\sum_{\text{cells }c}\big[E_c(m)-\mu\big]\Big),
\]
\[
E_c(m)=\eta_*\,K^2(m_c)+\lambda\,n_{\rm coll}(m_c),
\]
where the sum runs over the slice cells, $K^2$ is the closure invariant of Appendix~B, $n_{\rm coll}$ counts label collisions (injectivity is imposed as a penalty whose hard limit is the constraint, with the residual collision fraction reported so the softness stays visible), $\varepsilon$ pins the slice volume, and $\mu$ is the per-cell label free energy, computed by thermodynamic integration so that the closure term cannot masquerade as a shift of the bare cosmological coupling. The labels carry a uniform base measure, so at $\beta=0$ the label entropy cancels exactly and the bare host is recovered identically. The physical weighting fixed by Part~II is $\beta=1$ at $\eta_*$: once that chain is accepted there is no coupling dial left free, and the intermediate $\beta$ values serve only as a diagnostic interpolation.

Every run opens by recomputing the single-cell chain on its own tables and refuses to proceed unless $g_{\mathrm{share,eff}}=7.4198$ and $\langle K^2\rangle_{\eta_*}=3/(2\eta_*)=50.223$ are reproduced. This checks the shared single-cell spectrum and normalization; the lattice measure uses the finite collision penalty stated above. Every quoted triangulation is connected, simplicial, and correctly foliated, and comparisons use matched total volume at the same regulator. The principal control permutes energies across the $210$ sorted-label orbits, preserving the orbit-value pool and permutation symmetry. The orbit multiplicities are $1,4,6,12,24$, so this permutation changes the ordered-state density of states and one-cell thermodynamics. Comparisons with this control measure sensitivity to that combined reassignment. The run validity criteria were recorded before the corresponding data were examined (Appendix~J.3).

\subsection{24.3 What the closure sector cannot supply: vacuum stiffness}

Three cell-order calculations bound the closure sector's contribution to vacuum stiffness. The static closure weight induces only about one percent of the bare geometric coupling. The selected memoryless kernel, modeled as cell-by-cell redraws, has a closed stationary state with vacuum admissibility $0.536$ and no detectable curvature action. A separate closure-class history tilt, solved by a Doob transform on a ring, also remains short-ranged and well below the required coupling. These are calculations of specified effective models; they do not derive the refresh kernel from Postulate~III.

Two further calculations extend the exclusion beyond the label sector, to the conserved capacity field that Section~24.7 introduces as the carrier of the long-range sector. Even a field the labels cannot see might rank geometries on its own, because integrating out a conserved Gaussian field induces a purely geometric action: $\tfrac{1}{2}\ln\det'L(g)$ per channel, the spanning-tree entropy of the slice graph by the matrix-tree theorem. Computed on degree-matched proxies for smooth extended and for crumpled geometry, and on real engine slices, that entropy is nearly universal at fixed coordination: the per-cell differential is a few times $10^{-4}$, and with all seven channels the phase-tipping force is of order $10^{-3}$ against the bare $k_0\simeq2.2$. The theory's own non-Gaussianity closes the loophole tighter still. The finite budget of Postulate~I gives the field a saturation mass $m^2=2/g_{\mathrm{share,eff}}\simeq0.27$, and the massive determinant suppresses precisely the soft modes that carried the residual sensitivity, shrinking the differential by a further factor of three at the budget mass and toward zero beyond. What the field retains is a local renormalization of the host's couplings, of order $0.08$ per cell in the seven-channel count: enough to relocate the host's phase boundaries, and unable to create the phase.

Within the tested models, static weight, refresh, history tilts, and conserved fields do not supply vacuum stiffness. The refresh does couple geometry to the density of closure failure, which the defect experiment measures. The same pass verifies the locality fracture: injectivity-preserving unit shifts split the $1680$ states into $48$ components of $35$ states. This excludes that local realization of the full-entropy kernel; it does not derive the nonlocal physical operator.

\subsection{24.4 The externally hosted vacuum and the role of conditioning}

The exclusions show that the specified substrate does not dynamically select its vacuum geometry. Two questions called emergence of spacetime must therefore be separated. The first is kinematic encoding: the closure invariant $K^2$ measures the failure of the four oriented faces to close, adjacency carries proximity, and the marked-transfer dictionary of Section~13 converts entanglement increments into meters. The cells are dimensionless capacity units rather than sites of a preferred spatial grid. The second question is dynamical selection among crumpled, branched, and smooth extended geometries. Section~24.3 finds that none of the tested substrate mechanisms performs that ranking.

Many-Pasts cannot fill that dynamical gap. Appendix~G normalizes the probabilities of alternative present records before conditioning, so $p(h\mid P)$ describes histories compatible with an already specified present and does not select which $P$ occurs. In the simulations the smooth extended vacuum is supplied by the chosen Regge/CDT host phase as an external background datum. Many-Pasts may condition the compatible decoherent histories within that datum, but it neither ranks candidate geometries nor explains why the extended present is realized. The coefficients of Part~II are computed on that host state, and the medium dresses it rather than generating it. A substrate derivation of the prior state or of a probability measure over alternative vacuum geometries remains open.

Regge/CDT is a useful external host because it supplies tetrahedral cells, dynamical curvature, no rigid preferred spatial grid, and a demonstrated extended four-dimensional phase \citep{AmbjornJurkiewiczLoll2004,AmbjornGJL2012}. This choice is operational, but the completed fixed-regulator runs are not a vacuum derivation by the substrate. The host supplies the geometry on which the label theory is tested; the substrate then produces quantitative dressing of its couplings. Integrating out the labels gives $S_{\mathrm{eff}}=S_{\mathrm{host}}-\log Z_{\mathrm{label}}$. Appendix~J.6 reports the predicted extensive volume shift and its measured scaling family, including the independently predicted shuffled-control line. The corresponding curvature-sector shifts remain to be tested statistically. The hierarchy is explicit: cell-level non-closure reweights the glued tetrahedra, and integrating out the labels converts that reweighting into corrections to the host's volume- and curvature-sector couplings. This is presently a quantitative interface with an externally supplied vacuum geometry. Section~24.9 states the additional critical-surface test that would promote the interface to a continuum embedding.

\subsection{24.5 Compatibility: the weighting on dynamical geometry}

At $N_{41}\simeq100{,}000$, three eighty-slice ensembles were evolved from the same tuned host with respectively zero closure coupling, the physical weighting $\beta=1$, and $\beta=1$ after shuffling the closure energies among label orbits (Appendix~J.5). Their total-simplex counts agree within $0.6\%$, and every final triangulation remains connected, simplicial, and correctly foliated. The label sector orders exactly as the closure weight demands: the collision fraction falls from $0.654$ at $\beta=0$ (against the uniform-measure prediction $1-840/2401=0.650$) to $0.077$ under the physical weighting, while the shuffled control stays high at $0.736$. The Hausdorff values remain $3.55$--$3.59$ and the volume profiles remain extended across the matched ensembles. Repeating the physical weighting with forty rather than eighty slices gives collision fraction $0.078$ and $d_H=3.74$. These finite runs show strong admissibility ordering with stable host geometry. The shuffled comparison also changes the local thermodynamics.

\subsection{24.6 The defect experiment: geometry responds to the theory's mass}

The theory's matter is persistent closure failure: committed capacity, maintained against the refresh (Sections~3.2 and~23). The experiment inserts it by hand and repeats the comparison with two independent random seeds. In each repetition, one ensemble holds one hundred well-separated cells in maximal closure failure (all four faces at $m=0$: six collisions, $K^2=48$), while a control ensemble holds the same number in the best-closed injective configuration ($m=\{0,1,2,3\}$, the minimum $K^2=40.67$). Anchoring, protection from geometry moves, and every other update rule are identical. Every final triangulation remains connected and correctly foliated, and all one hundred marked cells survive in each ensemble. Subtracting the closed-cell control therefore isolates the physical effect of failure content. The final analysis pools the two seeds and uses across-seed scatter where replicated and integrated-autocorrelation-time-corrected standard errors otherwise.

The capacity response is measured with high statistical significance. At the first shell the mean closure energy is $1.6802\pm0.0026$ around failure pins and $1.7260\pm0.0008$ around closed pins, a difference $-0.0458$ at $16.8\sigma$. The collision fraction changes from $0.0783\pm0.0002$ to $0.0698\pm0.0009$, a $9.2\sigma$ separation. These label observables largely return by the second shell, as the short closure correlation length requires. The geometric channel appears in a different observable and farther out: at the third shell the mean coordination is lower around failure pins by $0.0398$, a $4.0\sigma$ separation from the closed-cell control after autocorrelation and seed scatter. Shell-cell counts do not separate significantly at this precision, and the first- and second-shell coordination differences are consistent with zero. The data establish a local geometric response.

Persistent closure failure measurably strains the surrounding capacity state and changes the local coordination field. The closed-cell control carries the same anchoring without the failure, isolating the failure content as the source. This finite-lattice result supplies the microscopic source-to-strain seed used by the weak-field sector. The defect-only experiment reaches three lattice steps and establishes neither continuum curvature nor a value of $G$. The separately conserved instrument of Section~24.7 tests transport beyond that short closure range.

\subsection{24.7 Reaching Newtonian range: the conservation requirement}

The strain of Section~24.6 dies within a lattice step, as the sub-cell correlation length requires. Macroscopic transport needs a different operator. If free capacity is only a per-cell budget, local maximum-entropy re-equilibration gives $(2z-L)\delta\mu=-m/\chi$ on the slice adjacency ($z=4$ faces per cell, $L$ the slice Laplacian). That operator has no small-momentum pole; on a real slice the response falls eight orders of magnitude within thirteen steps. A locally rebalanced budget therefore screens at the substrate scale.

\paragraph{Committed stock and maintenance throughput.}
Particle formation and persistence are different operations. Let $J_{c,A}^\mu$ carry the committed stock belonging to source $A$. Its balance law is
\[
\nabla_\mu J_{c,A}^\mu
=\Gamma_{{\rm form},A}-\Gamma_{{\rm rel},A}.
\]
For an already-existing stable defect both rates vanish, so its mass-equivalent committed stock remains constant. The defect nevertheless draws a maintenance throughput
\[
\mathcal P_{\rm maint}=\alpha_{\rm maint}\rho_{\rm def}.
\]
This is the rate at which the surrounding renewal bandwidth services persistence; it is not continuing conversion into more matter. Spatial free-bandwidth flux enters the defect event, while the reversible update exports the previous replaceable closure state along an internal history edge. The history register is an output of the renewal circuit, not an ordinary spacetime fluid, so no four-current $J_{\rm hist}^\mu$ is assumed. In free propagation that output carries no cell address of the persistent mark. This stock--throughput distinction removes the apparent secular growth of a static particle.

\paragraph{Augmented-graph Ward identity.}
The following result holds for the refresh-bandwidth completion. Form the directed renewal event graph whose edges are spatial bandwidth transport, ordinary temporal continuation, and the history output of a maintenance event. Let $B$ be its incidence matrix and $j_e$ the oriented edge throughput. In the absence of formation, release, or terminal absorption, each local reversible gate has one incoming active slot for every outgoing active or history slot. Its exact Kirchhoff identity is
\[
\Delta_t n_v+\sum_e B_{ve}j_e=0.
\]
At a stable defect the history-edge current is $j_v^{H}=\mathcal P_{{\rm maint},v}$; it is an outflow from the present register, not an increase of committed stock. Summing over any region cancels every internal edge and leaves only boundary transport and the explicitly counted history outputs.

On a fixed spatial slice, a local linear constitutive law has the quadratic operator
\[
L_f=B_sW B_s^T,
\]
where $B_s$ is the spatial incidence matrix and $W$ contains positive face conductances. Since $B_s^T\mathbf1=0$,
\[
\boxed{L_f\mathbf1=0.}
\]
For a connected slice the constant mode is the only exact zero mode. Locality and spatial isotropy then give
\[
\lambda_f(k)=D_qk^2+O(k^4)
\]
near $k=0$; a nonzero constant term would violate the incidence identity and represents leakage. The steady equation on a compact slice is
\[
L_f\mu_q
=\alpha_{\rm maint}
\left(\rho_{\rm def}-\overline{\rho}_{\rm def}\right),
\]
where the subtraction removes the constant mode. In an extended three-dimensional limit its pseudoinverse has
\[
G_f(r)\sim\frac{1}{4\pi D_qr},
\]
so a persistent maintenance demand produces the unscreened $1/r$ profile without tuning a scalar mass. Conservation proves the gapless transport pole within this completion. It does not prove that refresh bandwidth is the physical carrier, that defects couple only to it, or that the resulting scalar is the metric capacity coordinate.

\paragraph{Lattice symbol and the continuum theorem it requires.}
On a regular three-dimensional graph with spacing $a$ and uniform face conductance $w$,
\[
L_f(k)=2w\sum_{i=1}^3\left[1-\cos(ak_i)\right]
=wa^2k^2-\frac{wa^4}{12}\sum_ik_i^4+O(a^6k^6).
\]
A node source representing a continuum density, $M_x=\rho(x)a^3$, turns $L_f\mu=\alpha M$ into $-\nabla^2\mu=(\alpha/D)\rho$ with
\[
\boxed{D=\frac wa,}
\qquad
G_\mu(r)=\frac{1}{4\pi Dr}+O\!\left(\frac{a^2}{r^3}\right).
\]
The pseudoinverse of the $128^3$ periodic cubic Laplacian, shell-averaged and fitted to $A/r+C$ over $3\le r\le25$, gives $A=0.07692$ against $1/4\pi=0.07958$, a $3.3\%$ finite-lattice offset. On a random host, the continuum limit requires the slice Dirichlet forms
\[
\mathcal E_a[q]=\frac12\sum_{\langle ij\rangle}w_{ij}(q_i-q_j)^2=\frac12\langle q,L_aq\rangle
\]
to converge in the Mosco sense to $\tfrac12D_*\int d^3x\,\sqrt g\,g^{ij}\partial_iq\,\partial_jq$ jointly with the metric--measure limit of the slices,
\[
(\Gamma_a,d_a,\mu_a,L_a)\longrightarrow(M,g,dV,-D_*\Delta_g).
\]
The associated strong-resolvent convergence gives strong convergence of the heat semigroups at fixed positive time. Convergence of ordered eigenvalues, heat kernels and Green functions requires uniform compactness, heat-kernel bounds and inverse estimates, with the inverse taken on the complement of the constant mode. In flat three-dimensional space $-D_*\nabla^2\mu=\alpha_{\rm maint}\rho$ forces
\[
\mu(\mathbf x)=\frac{\alpha_{\rm maint}}{4\pi D_*}\int\frac{\rho(\mathbf y)}{|\mathbf x-\mathbf y|}\,d^3y,
\qquad
\delta q(r)=\frac{\chi_q\alpha_{\rm maint}}{4\pi D_*}\,\frac{M}{r},
\]
so the $1/r$ shape follows from the microscopic conserved transport operator and the junction condition below fixes only its normalization. Joint metric--measure and spectral convergence of this kind is established for random walks on critical spherical Fortuin--Kasteleyn planar maps, where all ordered eigenvalues and the heat trace converge to those of Liouville Brownian motion \citep{OpenAIFKSpectral2026}. The corresponding statement for CDT slices is the three-dimensional analogue and remains open.

\paragraph{The continuum junction condition.}
Let the renormalized static response be $\delta q=\chi_q\mu_q$, $J_q^i=-D_q\nabla^i\mu_q$, and $\mathcal P_{\rm maint}=\alpha_{\rm maint}\rho_{\rm def}$. Away from the compact zero-mode subtraction,
\[
\nabla^2q
=\frac{\chi_q\alpha_{\rm maint}}{D_q}\rho_{\rm def}.
\]
The static bridge $q=1+2\Phi/c^2$ and $\nabla^2\Phi=4\pi G_*\rho$ therefore require
\[
\boxed{
\frac{\chi_q\alpha_{\rm maint}}{D_q}
\longrightarrow\frac{8\pi G_*}{c^2}
}
\]
after renormalized-operator and lattice-spacing matching. The three quantities on the left are lattice outputs. The electron-derived $G_*$ enters only after their dimensionless regulator dependence has been removed, so the continuum comparison is target-blind rather than a fit.

In SI units the target is $8\pi G_*/c^2=1.86639459\times10^{-26}\,{\rm m\,kg^{-1}}$. In native units the condition becomes dimensionless. With $m_*=\hbar/(cL_*)=2.17644\times10^{-8}\,{\rm kg}$ and $G_*=c^3L_*^2/\hbar$, measuring length in $L_*$ and density in $m_*/L_*^3$ turns the static equation into
\[
\boxed{\widetilde\nabla^2q=8\pi\,\widetilde\rho,}
\]
so the renormalized dimensionless junction coefficient must flow to $8\pi=25.13274$. The finite-lattice shell-one ratio reported below becomes this object only after the field, source, and operator normalizations and the physical graph spacing have been extracted, and its present value supports no comparison with $8\pi$.

This transport experiment is performed on a separate forty-slice realization at $N_{41}\simeq100{,}000$ (Appendix~J.8). The field begins at the vacuum anchor $7.4198$, moves only through antisymmetric face fluxes, and is recycled through the compact zero mode; its mean remains $7.4198$ through all $4000$ sweeps. Absorption reads the persistent excess failure above the measured vacuum dressing, so the source strength is produced by the label dynamics rather than assigned by the transport law. Five fixed source values, $m=0,1,2,3,6$, are represented by thirty-two well-separated defects each. The $m=0$ value measures commitment-independent boundary dressing; the positive values test the source and response laws independently. Within-pin series use Sokal integrated-autocorrelation corrections, and final errors include across-pin and across-level scatter.

Four distinct measurements give a consistent finite-lattice source-to-field chain. First, conservation is exact. Second, after subtracting the $m=0$ dressing, the emergent maintenance charge is additive,
\[
Q-Q(0)=(0.1440\pm0.0025)m,
\qquad Q(0)=0.0075.
\]
Third, the shell-one deficit divided by that independently measured charge is constant across the positive source values,
\[
\left.\frac{\Delta f_1-\Delta f_1(0)}{Q-Q(0)}\right|_{m=1,2,3,6}
=4.9965\pm0.5244,
\]
where the uncertainty is the larger of propagated error and across-level scatter. Sources of different strength therefore couple through one finite-lattice junction ratio. The dimensionful continuum coefficient $8\pi G_*/c^2$ still requires a continuum scaling and operator-matching calculation. Fourth, the dressing-subtracted profile remains positive through eight shells and fits
\[
\Delta f(d)\propto d^{-0.62\pm0.11}.
\]
The screening diagnostic gives $d\ln(\Delta f\,d)/dd=+0.1062\pm0.0551$, statistically consistent with zero and implying the finite-range bound $\xi>4.6$ lattice steps. The exact Ward identity and the absence of detected screening therefore agree: the tested carrier is conservation-protected and long-ranged over the usable graph radius.

Two qualifications are visible in the same data. The proportional source fit has a maximum relative residual of $0.242$, driven mainly by the $m=6$ source falling below the extrapolation from weak sources; this is finite-capacity saturation, not precision linearity at arbitrary multiplicity. The exponent differs from the ideal three-dimensional massless Green-function value $-1$. The normalized profile is described by a $1/r$ Green function convolved with a finite-width source, together with the far-field offset left by the compact zero-mode subtraction,
\[
F(d)=\frac{\operatorname{erf}(d/\sqrt2\sigma)/d+C}{\operatorname{erf}(1/\sqrt2\sigma)+C},
\]
with $\sigma=1.95$ shells, $C=-0.074$, and an rms residual of $0.018$ over the eight shells. For comparison, the best unweighted normalized power law ($d^{-0.57}$) leaves $0.107$, the smeared source without offset $0.050$, and $1/d$ plus a constant $0.134$; the zero-mode offset by itself accounts for little of the flattening, and the finite source width carries most of it. Without the shell covariance this comparison is descriptive. Read literally, a $d^{-0.62}$ tail would correspond to an effective Green dimension $d_{\rm eff}\simeq2.6$, and a larger-volume study must separate the two readings through the low Laplacian spectrum, the heat trace and return probability, and the spectral dimension. These measurements establish the tested finite-lattice behavior while leaving the continuum scaling and absolute junction normalization open.

\subsection{24.8 What the finite-lattice results establish}

The fixed-lattice results establish four points. The admissibility sector orders on a dynamical host without destabilizing its matched geometry. Persistent closure failure produces a replicated local capacity response and an autocorrelation-corrected local coordination response. Conditional on refresh bandwidth as the carrier, the exact augmented-graph identity protects the long-range pole. The separate transport calculation then verifies conservation, approximately additive source strength, one lattice junction ratio across the tested source values, a power-law tail, and no detected screening. None of these mechanisms selects the vacuum geometry; CDT supplies that host state externally, while the substrate measurably dresses its $N_{41}$ coupling direction. The exact Regge volume projection of that direction is given in Section~25. The remaining CDT questions are qualitatively different: whether the joint system lies in the required extended phase with a massless TT sector, whether it also reaches an optional continuous cutoff-removal trajectory, how its lattice spacing and operators match to continuum observables, and whether the absolute junction coefficient tends to $8\pi G_*/c^2$. Section~26 records those distinct grades.

\subsection{24.9 The capacity-decorated continuum target}

The present runs compare coupled and control ensembles at fixed regulator.  Because $L_*$ is a physical microscopic length, existence of the effective theory does not require a second-order point at which the regulator spacing vanishes.  The necessary geometric condition is instead membership in an extended semiclassical phase whose long-wavelength transfer matrix contains a massless transverse--traceless sector.  A continuous critical surface is a separate, stronger universality and cutoff-removal test. If such a surface exists and its continuum limit is a four-dimensional unitary positive-energy local theory with a discrete bounded-below scaling spectrum and a physical local dilatation current, scale symmetry implies local conformal symmetry and a traceless improvement of the stress tensor with unchanged spacetime charges \citep{OpenAIScaleConformal2026}. Establishing those hypotheses at the fixed point would constrain its scalar trace modes.  For a causal triangulation $T$, let $Z_{\rm cap}[T;\mathbf g_{\rm cap}]$ be the trace over its closure labels, defect marks, renewal gates, source-labelled committed allocations, and history outputs. The joint partition function is
\[
Z(\mathbf g)=
\sum_{T\in\mathfrak T_{\rm CDT}}
\frac{1}{C_T}
e^{-S_{\rm CDT}[T;\kappa_0,\Delta,\kappa_4]}
Z_{\rm cap}[T;\mathbf g_{\rm cap}],
\]
and integrating out the capacity sector gives
\[
S_{\rm eff}[T]
=S_{\rm CDT}[T]-\ln Z_{\rm cap}[T].
\]
The theory fixes $\mathbf g_{\rm cap}$ at its microscopic values.  Phase membership must then be tested at the tuned host couplings with three observables:
\begin{enumerate}
\item a Hausdorff estimator tending to four, $d_H(N_{41})\to4$, as the volume grows;
\item a de Sitter-like volume profile, $N_3(t)\propto\cos^3(t/s_0)$ with $s_0\propto N_4^{1/4}$;
\item a gapless TT transfer mode, $\lambda_{\rm TT}(k)=1-c_{\rm TT}k^2a^2+O(k^4)$ as $k\to0$.
\end{enumerate}
The completed volume pair bears only on the trend toward the first condition and on a coarse extended-volume proxy: increasing $N_{41}$ from about $42{,}000$ to $102{,}000$ raises $d_H$ from $3.31$ to $3.56$ and the blob score from $1.46$ to $1.57$.  At about $100{,}000$ simplices, the three matched ensembles remain connected and correctly foliated and give $d_H=3.55$--$3.59$.  These results are consistent with movement farther into the finite-volume extended region, but they do not establish $d_H\to4$; no de Sitter-profile scaling fit or TT transfer-spectrum measurement has yet been performed for the decorated ensemble.  A gapped TT sector would fail the third condition even if the volume data appeared extended.

Two capacity-register numbers reported elsewhere are deliberately excluded from this phase test.  The refresh model's stationary admissibility $0.536$ and the history-channel correlation length $\xi_{\rm hist}\simeq0.57$ cells characterize the finite register dynamics, not the host's geometric correlation length or graviton spectrum.

The additive vacuum term and the finite continuum volume coupling are distinct coordinates of this approach, but the volume normalization is not an adjustable unknown. Section~25 derives the exact Regge operator
\[
\mathcal V_4=a^4(v_{41}N_{41}+v_{32}N_{32})
\]
and its corresponding source direction in the $(\kappa_4,\Delta)$ plane. On a fixed-ratio trajectory the scalar matching is $t_4=\overline v_4a^4\rho_{\rm vol,R}+o(a^4)$, with $\overline v_4=(v_{41}+\xi v_{32})/(1+\xi)$. The continuum calculation must therefore measure the limiting $\alpha$, $\xi$, lattice scale, and mixed critical eigenoperator, rather than fit a free $Z_V$. Vacuum normalization removes the complete capacity-sector bulk coefficient at every fixed regulator. Consequently a finite source-independent capacity density cannot hide in $t_4/a^4$: it would already be an extensive term before the limit and would have been included in $\mu_{\rm cap}$. If the common-parent critical limit exists, its gravitational volume coefficient still requires a matching condition; O.5 adopts $\Lambda_{\rm R}^{\rm empty\ equilibrium}=0$. The currently external CDT host does not satisfy that antecedent by assumption, so its own cosmological coupling remains a regulator input rather than a result of these runs.

If cutoff removal is additionally required, $\kappa_4$ can be tuned toward infinite volume and the remaining regulator couplings can be tested for a continuous transition. CDT supplies evidence that candidate continuous lines can occur and can support this stronger construction \citep{AmbjornSecondOrder2011,AmbjornCritical2019}. Let $\mathfrak C_*$ denote the manifold on which the geometric correlation length diverges in lattice units. The distinct symbol keeps this geometric locus separate from the coherence observable $\mathcal C$.  The cutoff-removal embedding succeeds only if the fixed capacity slice intersects that manifold,
\[
\boxed{
\mathcal P_{\rm cap}
\cap\mathfrak C_*\ne\varnothing,
\qquad
\frac{\xi_{\rm geom}}{a}\longrightarrow\infty.
}
\]
This requirement introduces no capacity fit. Failure of the fixed slice to intersect $\mathfrak C_*$ would reject the cutoff-removal embedding, but would not by itself reject a finite-$L_*$ theory whose extended phase and massless TT sector had already been demonstrated. Geometric criticality removes the ultraviolet regulator. The augmented-current Ward identity has a different job: it protects the infrared transport pole. Neither result supplies the other.

Appendix~J.9 states the phase-membership, optional critical-surface, finite-size-scaling, operator-mixing, current, and junction tests in full. A successful capacity-decorated CDT transfer matrix would provide an independent nonperturbative geometric embedding of the finite marked vertex. Until the phase and TT tests are passed, CDT remains the external host used by the simulations; the factorized EPRL/GFT construction of Section~23 and Appendices~B, H, and Q supplies the separate controlled semiclassical Einstein route.


\subsection{24.10 Registered normalization of the connected host}

The connected $\alpha=7$ history measure has been sampled at the
curvature couplings fixed in J.12. The preregistered comparison of its
homogeneous volume stiffness with its boundary normalization gives
$G_{\rm vol}/G_0=1.064$, with standard error $0.089$ in the logarithm,
and the verdict \textsc{agree} under the registered $0.80$--$1.25$ band.
This primary result uses A4, the only eligible stationary 32k run.
The reported-only stationary runs A2 at 16k and A5 at 64k give $1.060$
and $0.893$. Their de~Sitter-like profiles and stiffnesses support this
finite-volume correspondence.

The separate start-independence test remains \textsc{undetermined}.
A registered curvature-trace matter test is also \textsc{undetermined},
and exploratory scalar binding has no stable energy extraction.
Appendix~J.14 gives all run roles, uncertainties, exclusions and
records. This result is a normalization check on the selected
microscopic measure; a universal matter response and interacting
graviton sector remain open.

\section{25. Equilibrium Vacuum and Cosmological Term}

The normalization $q=1$ defines the unstrained vacuum reference. The regulated joint ensemble makes the corresponding subtraction exact, rather than leaving it as an unspecified cancellation. At fixed triangulation let
\[
Z_{\rm cap}[T]
=\mathbb E_{m\sim\mathrm{unif}}
\exp\!\left[-\beta\sum_cE_c(m)\right]
\]
and define its homogeneous bulk free energy per cell by
\[
\mu_{\rm cap}(a,\beta)
=-\lim_{N_{\rm cell}\to\infty}
\frac{1}{\beta N_{\rm cell}}\ln Z_{\rm cap}^{\rm vac}(a,N_{\rm cell}).
\]
Its dependence on the homogeneous regulator couplings is suppressed in the notation; the definition is applied along the trajectory whose continuum limit is being tested.
The vacuum-normalized capacity factor is
\[
\widehat Z_{\rm cap}[T]
:=e^{\beta\mu_{\rm cap}(a,\beta)N_{\rm cell}(T)}Z_{\rm cap}[T].
\]
By construction,
\[
\boxed{
-\lim_{N_{\rm cell}\to\infty}
\frac{1}{\beta N_{\rm cell}}
\ln\widehat Z_{\rm cap}^{\rm vac}=0.
}
\]
In the thermodynamic limit,
\[
F_{\rm cap}=\mu_{\rm cap}(a,\beta)N_{\rm cell}+o(N_{\rm cell}).
\]
The centered bulk potential
\[
\Omega_{\rm cap}=F_{\rm cap}-\mu_{\rm cap}(a,\beta)N_{\rm cell}
\]
has zero density in the homogeneous reference state. The cancellation is invariant under an arbitrary change of microscopic energy origin. If
\[
E_c\mapsto E_c+C,
\qquad
\mu_{\rm cap}\mapsto\mu_{\rm cap}+C,
\]
then
\[
\boxed{\widehat Z_{\rm cap}[T]\mapsto\widehat Z_{\rm cap}[T]}
\]
exactly. When the bulk cell count admits the usual thermodynamic reading, Gibbs--Duhem writes the same relation as $\Omega_{\rm cap}/V=-P$; an isolated equilibrium reference has $P=0$ \citep{VolovikVacuum2005}. The partition identity above is stronger for the regulated model because it does not require the conditional refresh current to serve as that thermodynamic charge. A large homogeneous source-independent zero-point term therefore cannot reappear as a capacity-sector observable; it cancels against the uniquely shifted bulk free energy, while defect and geometry-dependent differences remain.

Integrating out the normalized capacity sector gives
\[
S_{\rm eff}[T]
=S_{\rm CDT}[T]-\ln\widehat Z_{\rm cap}[T].
\]
The homogeneous extensive capacity term is absent from this action. Only the geometry- and defect-dependent free-energy difference remains to source strain, as required by the information--geometry and mass--entropy postulates. This result does not use the conditional refresh-current Ward identity. Transport conservation protects the infrared pole; vacuum normalization follows from the equilibrium partition measure. They are different conserved structures.

Appendix~J has already tested the finite-regulator algebra. Omitting the centering restores the term $+\beta\mu_{\rm cap}N_{\rm cell}=+(\beta\mu_{\rm cap}/2)N_{41}$ in the geometric action. Against the quadratic volume pin it predicts
\[
\Delta N_{41}=-\frac{\beta\mu_{\rm cap}}{4\varepsilon},
\]
and the measured uncentered-minus-centered displacement is $-62.8$ against $-62.2$. The run does not prove a continuum cosmology, but it verifies the coefficient and sign of the subtracted microscopic bulk term. One bookkeeping refinement matters here. In the standard CDT action basis,
\[
S_{\rm CDT}
=-(\kappa_0+6\Delta)N_0
+\kappa_4(N_{41}+N_{32})
+\Delta(2N_{41}+N_{32}),
\]
the restored $N_{41}$ term is the coupling displacement $\delta\Delta=+\beta\mu_{\rm cap}/2$, $\delta\kappa_4=-\beta\mu_{\rm cap}/2$, $\delta\kappa_0=-6\,\delta\Delta$, not a pure shift of $\kappa_4$. At fixed simplex ratio it contains the expected volume chemical potential; its orthogonal component renormalizes the regulator asymmetry. The measured displacement therefore verifies the predicted $N_{41}$ direction, which is the exact statement licensed by the control.

The same conclusion is independent of the centering convention. The leading exponential growth of the decorated canonical sum defines a critical surface in the full $(\kappa_4,\Delta)$ plane. An extensive change of microscopic energy origin translates the bare coupling vector and that critical surface by the same displacement displayed above. Their normal difference is invariant, and it vanishes at the infinite-volume surface. On a trajectory that holds $\Delta$ fixed after the other relevant directions have been projected out, this normal coordinate is the familiar
\[
t_4:=\kappa_4-\kappa_4^c\longrightarrow0.
\]
Thus critical tuning absorbs the additive bulk term without turning an energy-origin choice into an observable. The remaining question is the normalization of the physical volume direction, to which we now turn.

The continuum volume matching can be made explicit, with no unspecified volume factor $Z_V$. After Wick rotation, the exact Regge four-volumes of the two simplex types are \citep{AmbjornGJL2012}
\[
V_{41}=a^4v_{41}(\alpha),
\qquad
v_{41}(\alpha)=\frac{\sqrt{8\alpha-3}}{96},
\]
\[
V_{32}=a^4v_{32}(\alpha),
\qquad
v_{32}(\alpha)=\frac{\sqrt{12\alpha-7}}{96},
\]
with $\alpha>7/12$ so both Euclidean simplex types are nondegenerate. Here $a$ is the spatial regulator edge length. In the hosted reading it is scaffolding with no fixed relation to $L_*$; under the face-wise lock of Appendix~H.9a, paragraph~4a, every spatial face carries the channel area $a_{\rm ch}=4\ln2\,L_*^2$ and the edge is physical, $a^2=(16\ln2/\sqrt3)L_*^2$, $a=2.530L_*$ (Appendix~J.10). The integrated volume operator is therefore the kinematic identity
\[
\boxed{
\mathcal V_4[T]
=a^4\big[v_{41}(\alpha)N_{41}+v_{32}(\alpha)N_{32}\big].
}
\]
Let $\rho_{\rm vol,R}$ denote the coefficient of $\mathcal V_4$ in the dimensionless Euclidean effective action. If $G_{\rm R}$ is the continuum Newton coefficient and $L_G^2:=\hbar G_{\rm R}/c^3$, then
\[
\rho_{\rm vol,R}
=\frac{c^3\Lambda_{\rm R}}{8\pi\hbar G_{\rm R}}
=\frac{\Lambda_{\rm R}}{8\pi L_G^2}.
\]
On the matched weak-field branch $G_{\rm R}=G_*$ and therefore $L_G=L_*$. Keeping $L_G$ visible until that junction passes prevents the volume matching from assuming the gravitational result it is meant to join.
The two action-conjugate volume couplings are therefore
\[
\boxed{
t_{41}=a^4v_{41}(\alpha)\rho_{\rm vol,R},
\qquad
t_{32}=a^4v_{32}(\alpha)\rho_{\rm vol,R}.
}
\]
Equivalently, in the $(\kappa_4,\Delta)$ basis a pure volume-source displacement obeys
\[
\boxed{
\begin{pmatrix}\delta\kappa_4\\\delta\Delta\end{pmatrix}
=a^4\rho_{\rm vol,R}
\begin{pmatrix}2v_{32}-v_{41}\\v_{41}-v_{32}\end{pmatrix}.
}
\]
Substitution returns $\delta(\kappa_4+2\Delta)=a^4v_{41}\rho_{\rm vol,R}$ and $\delta(\kappa_4+\Delta)=a^4v_{32}\rho_{\rm vol,R}$, so this is an operator identity, not a fitted matching. At the isotropic point $\alpha=1$, $v_{41}=v_{32}=\sqrt5/96$ and the $\Delta$ component vanishes, providing a direct check.

The one-parameter formula of Section~24.9 is the fixed-ratio compression of this two-operator statement. If $\xi=N_{32}/N_{41}$ and $N_4=N_{41}+N_{32}$, then
\[
\overline v_4(\alpha,\xi)
=\frac{v_{41}(\alpha)+\xi v_{32}(\alpha)}{1+\xi},
\qquad
t_4=\overline v_4(\alpha,\xi)a^4\rho_{\rm vol,R}+o(a^4).
\]
Thus the geometric part of $Z_V$ is fixed:
\[
\boxed{Z_V^{\rm geom}(a)=\overline v_4\big(\alpha(a),\xi(a)\big).}
\]
The critical calculation still has to show that $\alpha(a)$ and $\xi(a)$ approach finite limits, set $a$ by a target-blind observable, and project out any mixed relevant eigenoperator. Those are continuum-existence and scale-setting tests; they leave no freedom in the volume normalization.

This matching does not by itself identify a CDT regulator simplex with the physical capacity cell. In the canonical reference convention used for this matching, one cell is assigned four-volume $L_*^4$ in $x^0=ct$ coordinates. This is not the physical event volume $v_{\rm ev}$ of the later C.8 paired-cell calibration. Appendix~J.10 records that this convention and the fixed face area are in tension by a factor $1.9$: a regular spatial cell of edge $2.530L_*$ has three-volume $1.91L_*^3$, so one cell swept through one tick is $1.91L_*^4$, and the event measure must be fixed by one of the two readings. The average number of regulator simplices in such a block is
\[
n_*(a)=\frac{L_*^4}{a^4\overline v_4},
\]
and its volume coupling is
\[
\boxed{
t_*:=n_*(a)t_4
=\rho_{\rm vol,R}L_*^4
=\frac{\Lambda_{\rm R}L_*^4}{8\pi L_G^2}
\xrightarrow{\,G_{\rm R}=G_*\,}
\frac{\Lambda_{\rm R}L_*^2}{8\pi}.
}
\]
The cutoff and the simplex factor cancel. CDT can therefore supply arbitrarily fine scaffolding while $L_*$ remains the finite physical capacity scale, as required by Section~24.1.

There is then no surviving capacity-sector $0/0$ ambiguity. The definition of $\mu_{\rm cap}(a,\beta)$ removes the complete thermodynamic bulk coefficient at every fixed regulator, not merely the part that diverges as $a\to0$. If a nonzero finite $\rho_{\rm cap,R}$ remained, then at any fixed $a$ the normalized vacuum action would contain
\[
\rho_{\rm cap,R}a^4
\big[v_{41}N_{41}+v_{32}N_{32}\big],
\]
which is still linear in the number of cells and therefore contradicts the defining zero bulk density of $\widehat Z_{\rm cap}^{\rm vac}$. Hence, on every regulator and on every continuum subsequence for which the volume operator has a limit,
\[
\boxed{
\rho_{\rm cap,R}^{\rm(eq)}=0,
\qquad
\Lambda_{\rm cap,R}^{\rm(eq)}=0.
}
\]
The complete equilibrium bulk term must be normalized before the regulator is removed. This order avoids inferring a continuum density from the ratio of two quantities that separately vanish.

The externally supplied CDT host has its own volume coupling. In a common geometry--capacity continuum, the complete effective action still requires a prescription for the gravitational volume coefficient. Centering fixes the normalized capacity free-energy reference, while the gravitational integration constant is fixed by a boundary or matching condition. Appendix~O.5 adopts $\rho_{\Lambda,i}=0$ for the empty equilibrium branch. This value is compatible with the centered measure, but it does not follow from the invisibility of constants in a fixed-volume probability distribution. Existence of a common-parent critical limit and its cosmological boundary data remain separate requirements.

\paragraph{A controlled nonequilibrium target.}
The external results also identify the correct next object more tightly than the phrase ``state-dependent energy'' alone. Let $C_0>0$ be the covariance of a regulated Gaussian capacity block in the centered equilibrium state and $C>0$ the covariance of another state on the same support. The dimensionless relative covariance is
\[
\mathcal G=C_0^{-1/2}CC_0^{-1/2}.
\]
The Kullback--Leibler divergence between the corresponding centered Gaussian measures is the basis-independent identity
\[
\boxed{
\Gamma_{\rm rel}(C\Vert C_0)
=D_{\rm KL}\!\left(\mathcal N(0,C)\Vert\mathcal N(0,C_0)\right)
=\frac12\operatorname{Tr}
\left[\mathcal G-I-\ln\mathcal G\right].
}
\]
If $\lambda_i>0$ are the eigenvalues of $\mathcal G$, each contribution $\lambda_i-1-\ln\lambda_i$ is nonnegative and vanishes only at $\lambda_i=1$. Thus
\[
\Gamma_{\rm rel}\ge0,
\qquad
\Gamma_{\rm rel}=0\Longleftrightarrow C=C_0,
\]
and for $\mathcal G=I+X$,
\[
\Gamma_{\rm rel}
=\frac14\operatorname{Tr}X^2
+O(X^3).
\]
Appendix~H.10 derives this functional rather than merely recognizing it: at finite regulator it is the quadratic-source Legendre transform of the equilibrium-normalized determinant of any positive dressed/reference Gaussian Hessian pair on a common physical support. The trace-log and positive mismatch forms are therefore closed as functional identities. The substrate must still derive the physical operator, its support and trace multiplicity, its dynamics and dimensions, and the conserved covariant response used in cosmology.
This introduces no new premise or fitted potential: it is the exact dimensionless relative-information functional of any positive Gaussian fluctuation sector. Multiplication by the appropriate thermodynamic scale turns it into a relative free energy. Bianconi independently promotes the same operator-convex form to a covariant geometric-relative-entropy theory,
\[
\Lambda_G
=\frac{1}{2\beta}\operatorname{Tr}
\left[G-I-\ln G\right],
\]
and obtains a low-energy Einstein limit and a local thermodynamic first law \citep{Bianconi2025,Bianconi2026}. The result therefore supplies a well-motivated target class for the homogeneous influence functional: equilibrium has zero relative energy, while a distinguishable state can carry a positive mismatch energy without restoring a source-independent vacuum constant.

Appendix~O.3 supplies a constrained covariant capacity-clock action and conserved metric stress tensor for a specified release history, thereby fixing the associated background and linear scalar exchange. The microscopic theory must still derive that clock potential from baryonic mixing and reproduce its common energy conversion and post-refresh ledger. Bianconi's $G$-field has its own derivative terms in the modified field equations, so it cannot be identified with $q$, varied as a new scalar, or imported wholesale without contradicting the metric-only decision of Section~11. Relative information nevertheless supplies a consistent local action, energy, and thermodynamic language for departures from the maximum-information vacuum. The remaining question is which departure the substrate dynamics selects.

Appendix~O develops the conditional dark-energy completion this leaves open. With the additional fixed local volume-form premise of O.1, primitive event measure makes the continuum variation unimodular and the cosmological term enters as an integration constant with $w=-1$ exactly, not as a local vacuum-energy coupling; constant shifts of the matter Lagrangian shift the constant without independently curving spacetime. The equilibrium centering fixes the microscopic free-energy reference. Appendix~O.5 separately adopts $\Lambda_i=0$ as the initial gravitational matching condition.

The source of later growth is collective loss of a baryonic source-clock carrier. Identity-free committed capacity cannot record stream assignment. When gas from several progenitors becomes irreversibly mixed in one retained bound object, matter no longer resolves progenitor membership and the still-labelled committed allocation is orphaned. Additive record pricing assigns the released stock fraction $\epsilon_{\rm mix}=1-I(A;Y)/H(A)$; faithful renewal fixes complete rerouting for each released elementary block, and Appendix~O.13 supplies the coefficient-free nine-state dilation. The closed-cohort advection--diffusion theorem makes the fully resolved position-information loss bounded and monotone; it does not by itself prove monotonicity for the coarse $Y$ used in the release law. The first cosmological zoom and the PPF and direct linear-lensing benchmarks are summarized in O.12. An independent prediction still requires a converged physical $K_{\rm mix}$ with progenitor geometry and retention rules fixed before the final direct-perturbation likelihood. The common energy conversion and post-refresh ledger remain physical premises. Once the release history is fixed, the vacuum stress retains intrinsic $w=-1$, while a separately conserved effective-fluid fit has $w_{\rm eff}<-1$ during positive release and approaches $-1$ as release ends. A retained well is a halo that keeps cooling gas in which progenitor labels mix. The Renaissance simulations follow radiative and supernova feedback at $z=8$--$15$ and find that about $10\%$ of halos of $10^7M_\odot$, half of those of $10^8M_\odot$ and all above $10^{8.5}M_\odot$ hold star-forming gas, nearly independently of redshift and region \citep{Xu2016Renaissance}. With this occupation as the retention gate and no input adjusted to the vacuum density, the measured mixing kernel gives $I_{\rm rel}=95.5\pm8$ and $\Omega_{\Lambda,0}=0.709\pm0.02$, against the likelihood target $94.2$ and $\Omega_{\Lambda,0}=0.707$. The reionization redshift and the post-reionization gate dominate the uncertainty. Gates computed without feedback lie lower and give $0.73$ from the hydrogen cooling criterion and $0.75$--$0.77$ from condensation in the archived run (Appendix~O.12). With the reionization redshift tied to the fitted optical depth, this frozen history has no free parameter and fits Planck, DESI DR2 and DES-Dovekie with four cosmological parameters at $\Delta\chi^2=+1.07$ against six-parameter $\Lambda$CDM, with $H_0=68.42\,\mathrm{km\,s^{-1}\,Mpc^{-1}}$.


\section*{Part VIII. Closure Status, Falsifiability, and Comparisons}
\addcontentsline{toc}{section}{Part VIII. Closure Status, Falsifiability, and Comparisons}

\section{26. Closure-Status Table}

The table below records each claim's premises, logical grade, and unresolved tests.

The leading word in each status uses a fixed vocabulary. \emph{Closed} means derived within the stated postulates and ensemble. \emph{Fixed} means no phenomenological freedom remains once the named branch is adopted. \emph{Conditional} means the result follows if a named reading or completion holds. \emph{Frontier} or \emph{open} means the piece is structured but incomplete. \emph{Empirical support} denotes comparison with data rather than derivation, and \emph{audit task} marks an independent check still required. A conditional premise does not demote every theorem downstream of it. Rows therefore separate exact identities and within-model lemmas from the microscopic or empirical status of their premises.

\begingroup
\scriptsize
\setlength{\tabcolsep}{3pt}
\renewcommand{\arraystretch}{1.12}
\begin{longtable}{|L{0.18\textwidth}|L{0.13\textwidth}|L{0.16\textwidth}|L{0.327\textwidth}|L{0.13\textwidth}|}
\hline
\textbf{Quantity / Claim} & \textbf{Sector} & \textbf{Status} & \textbf{Type of Support} & \textbf{Where Established} \\
\hline
\endfirsthead
\multicolumn{5}{r}{\textit{Closure-status table continued}}\\
\hline
\textbf{Quantity / Claim} & \textbf{Sector} & \textbf{Status} & \textbf{Type of Support} & \textbf{Where Established} \\
\hline
\endhead
$\Omega_{\mathrm{tet}}$, $g_{\mathrm{share,max}}$ & UV counting & Closed & exact combinatorics & Part II, App. B \\
\hline
$\eta_*$ & admissibility closure & Closed & stationary normalized closure evidence $\ln Z+\tfrac32\ln\eta$ maximized on the exact $K^2$ spectrum; the $3/2$ is the determinant weight of the three closure-defect components & Part II, App. B \\
\hline
$g_{\mathrm{share,eff}}$ & UV entropy & Closed & exact weighted evaluation & Part II, App. B \\
\hline
Memoryless refresh kernel & information / scale theorem & Closed under the foundational faithful full-support condition & applying that same condition to histories gives maximum path entropy, not a new premise: $H(B_{t+1}\mid B_t)=g_{\mathrm{share,eff}}-I(B_t;B_{t+1})$ is maximal iff $I=0$, uniquely giving $K(b,b')=p_{\eta_*}(b')$ & Part I, App. D, G--H \\
\hline
Renewal / retained marked quantum sector & quantum-channel interface & Exact channel, overlap, carrier incidence, and local Hilbert reconstruction; network extension conditional & $\mathcal E_*^\dagger(O)=\operatorname{Tr}(\rho_*O)I$ makes the renewed output coherence-blind; on the real retained $Q$ carrier the positive Hamiltonian selects the polar complex structure, while frame transport and $iS_\omega$ generate $\mathfrak{su}(9)$ and $\mathbb C^9/\mathbb{CP}^8$; the projected bridge gives $\mathfrak{su}(81)$ if the unreduced two-mark product survives the final host projection & \S3.3, \S22, App.~B.4, H.11, P.6, Q.9 \\

\hline
Decorated charged marked-transfer vertex & UV dynamics & Finite weights and routing closed; energy and geometric attachment selected & The local numerical coefficients survive coherent bond encoding. Connected normalized patches give a selected CDT weight; the joint orbital model preserves the native gap and clock. A common interacting continuum phase remains open. & \S13, \S23, App.~H.14--H.15, Q.11b \\
\hline
Gaussian GFT--geometric-relative-entropy bridge & UV / continuum & Exact Gaussian and capacity-lift identities; selected TT limit & Trace-log and covariance-mismatch identities are exact. Empty capacity lifts a supplied geometric rigging form. N.13 constructs a homogeneous form; the inhomogeneous interacting form and phase remain open. The linearized TT branch has a reflection-positive two-helicity Fock representation, and $B(0)=0$ with $\|B_1\|^2<Z_G\Delta_{\rm cap}$ is the explicit no-extra-pole test for a nonfactorized vertex. & \S25, App.~H.10, N.13 \\
\hline
Maximal-fusion / spin-foam refinement & geometric continuum & Exact maximal-channel map; selected connected completion & The pair parent has its 13-dimensional kernel. Connected chains and cycles surject onto coarse intertwiner spaces; Gram-normalized packets reproduce the selected Euclidean CDT weight. The all-edge pair parent alone has soft modes. The single-face matching operator is a ferromagnetic Heisenberg bond whose local hypercubic model orders at low positive temperature; exact Gauss projection removes that bond, so ordering of the physical shape variables needs multi-face, holonomy, or vertex terms. & App.~H.9a, H.12, Q.11b \\
\hline
Substrate cadence $\tau_*$ and causal length $L_*$ & scale setting & Cadence fixed within the decorated transfer action and electron anchor; causal length uses the invariant-speed identification & the determinant gives $r=e^{-7g_{\mathrm{share,eff}}}$, the marked vertex gives $Z_e=(1+\zeta_*)(1+7\zeta_*^2)$, and the charged gap fixes $\tau_*=-(3/2)Z_e\tau_e\ln(1-r)$; $L_*=c\tau_*$ defines the associated causal length but does not fix a graph displacement & Part I, App. D, H \\
\hline
Charged spatial propagation & matter / scale interface & Specified finite orbital model; continuum conditions open & Geometry transfer and cochain dynamics preserve the native rest gap and positive clock. The leading cone is normalized on a periodic family. Changing-boundary, species, and gauge-covariant spinor limits remain open. On the native complex a common speed requires the weights of a flat realization. & \S4, App.~H.9a, H.14--H.15 \\
\hline
$J_{\mathrm{bare}}$, $J_{\mathrm{eff}}^{\mathrm{tree}}$ & UV edge kernel & Closed & tetrahedral isotropy identity & Part II, App. C \\
\hline
$\Sigma_{\mathrm{ret}}=65/9$ & finite-loop UV & Conditional minimal return-sector completion & seven diagonal return channels plus one projected singlet are specified; an explicit microscopic return operator must derive their relative weights and exclude further motifs & Part II, App. C \\
\hline
$J_{\mathrm{eff}}^{(\mathrm{ren})}$ & finite-loop UV & Conditional on the minimal return operator & algebraic Dyson resummation once $\Sigma_{\mathrm{ret}}$ is specified & Part II, App. C \\
\hline
$\gamma$ & continuum stiffness & Conditional loop-dressed value & Euclidean normalization is exact for the one-sublattice edge convention; the numerical value inherits the minimal return operator & Part II, App. C \\
\hline
Green-matched source projection & UV source map & Closed in the canonical weak-field branch & exact defect counting + tetrahedral on-site Green function & App. C \\
\hline
$\kappa/\gamma$ & source-to-stiffness ratio & Closed in the canonical weak-field branch & $\sigma_{\mathrm{def}}=\rho/\kappa_m(L_*)$ plus $G_{\mathrm{tet}}(0)$ and tetrahedral $4/3$ projection & Part III, App. C--D \\
\hline
Weak-field action / bridge law & EFT / gravity & Closed for the ordinary longitudinal branch & Einstein--Hilbert $+$ GHY reduced in Newtonian gauge gives $I_{\rm Newton}$; $\delta S=-2S_\infty\Phi/c^2$ gives $I_{\rm cap}^{\rm static}=Z_SI_{\rm Newton}$ with no extra scalar & Part III, App. D, N \\
\hline
$G_*$ & gravitational scale & Parameter-free output of the decorated scale branch; historical postdiction & $G_*=(9/4)(\hbar c/m_e^2)Z_e^2\ln^2(1-e^{-7g_{\mathrm{share,eff}}})=6.6742890813\times10^{-11}$, or $-0.073\sigma$ relative to CODATA; the residual was known before the vertex was constructed & Part I, App. D, H, L \\
\hline
Event volume and channel area & physical scale & Conditional calibration & A Poisson event law, maximal-chain clock, paired-cell assignment, and regular mean area--volume conversion give $v_{\rm ev}=(\pi/24)m_4^4L_*^4$. A compatible quantum volume state exists; its selection is open. & App.~C.8 \\
\hline
Coherent bonds and renewal & microscopic state & Exact finite identities in a selected state & Encoding preserves the local marginal. Moving frames and background-wide replacement preserve retained coherence. Unequal path-conditioned reset counts can dephase; differing geometric event records remain unresolved. & App.~H.14, P.2, Q.11a \\
\hline
Metric response and linear constraints & geometric dynamics & Selected action; exact lattice identities and numerical tests & The Kuhn saddle gives the Einstein Hessian and same-coefficient Newton response. The staggered complex has two TT modes without corner doublers. A finite prime-field code gives an algebraic realization. & App.~H.13 \\
\hline
Nonlinear metric and ordering & quantum geometry & Exact identities within the selected ADM prescription & Six directional areas reconstruct $g$. The DeWitt measure and transported metric Weyl correction are explicit. Continuous-configuration commutators quantize the classical regulated brackets exactly. & App.~N.10--N.11 \\
\hline
Physical gravitational state space & quantum constraints & Positive homogeneous form; inhomogeneous limit open & The volume--shape spectral prescription gives a positive rigging form. Spatial and count errors are separately controlled on probes. Metricity boundaries, local constraints, and common physical-state refinement remain unresolved. & App.~N.12--N.15 \\
\hline
Connected covariant history weight & microscopic geometry & Exact within the selected normalized tensors and host counting & Boundary maps are isometric and sew to the Regge weight. Integer bundles select \(\alpha=7\). Host measure, volume treatment, and phase remain declared data. & App.~J.12, Q.11b \\
\hline
Joint orbital refinement and clock & finite coupled dynamics & Exact conditional operator identities & Isometric intertwining preserves the retained algebra, positivity, trace, polynomial dynamics, and regular clock functions. Source-family derivatives retain the moving-code and normalization terms. Locality and the continuum cone remain physical conditions. & App.~H.10, H.15 \\
\hline
Interacting Newton coefficient & covariant source response & Estimator specified; physical measurement open & Proper-length covariance gives the held-source interaction. The saved finite ensemble lacks an accepted full response and an infrared distance window. & App.~J.13 \\
\hline


Matched $G$ & weak-field gravity & Algebraic identity in the chosen normalization; not an independent determination & substituting the closed source map and $S_\infty^{\rm cell}$ into $G=c^2\kappa/(8\pi\gamma S_\infty)$ gives identically $G=c^3L_*^2/\hbar=G_*$ & Part III, App. C--D \\
\hline
Electron anchor & mass / length sector & Fixed empirical elementary anchor & one-bit fermionic defect supplies $\lambda_e$ for length setting and $m_e/\ln2$ for mass--entropy map & Part III, App. D \\
\hline
Seven-sector additivity and lightest branch & UV scale theorem & Closed entropy theorem; exact conditional mass minimum & subadditivity gives $H_k=kg_{\mathrm{share,eff}}-\Delta_k$ with $\Delta_k\ge0$; at fixed admissibility marginals, the adopted recurrence mass $m_k\propto e^{\Delta_k-kg_{\mathrm{share,eff}}}$ is minimized uniquely at the fermionic ceiling $k=7$ and $\Delta_7=0$ & App. D, H \\
\hline
Charged-lepton shell ladder & particle spectrum & Exact recurrence and overlap realizations within a selected pole-mass operator & stationary mean recurrence fixes the bare factor $720^N$ for the full-support shell register, while a common Gaussian singlet fluctuation gives the positive overlap $(2/7)^{N^2}$; closure Fisher information terminates the branch at $N=2$; identifying the resulting operator with charged-lepton pole mass remains the particle premise & App.~I.1 \\
\hline
$a_0$ & galactic EFT & Derived given the horizon-thermal tilt mode; the tilt-mode identification and the Unruh reading that fixes the lever coefficient are premises & the tilt mode of the transverse frame has frequency $\sqrt{g_{\rm bar}/\ell}$ with $\ell=g_{\mathrm{share,eff}}c/H$; at $k_BT_H=\hbar H/2\pi$ its gap is $\sqrt{g_{\rm bar}/a_0}$ with $a_0=[g_{\mathrm{share,eff}}/(2\pi)^2]cH_0$ (\path{Entropic/RAR/Pendular.lean}) & Part IV, App. N \\
\hline
Stock selection & galactic stock dynamics & Conditional effective dynamics; occupation identification, contact rate and record identification open & the Bose relative entropy of the response cell, weighted by equipartition with the excess field, biases momentum-conserving stock contacts; the zero-temperature limit of detailed balance gives descent to the budget-capped exponential law; spherical runs from NFW stock reach it to 0.01--0.05\,dex; cluster stock lies outside every record-dominated support and keeps $1+k_C$; observed-cluster runs keep $R_{500}$ within 0.1\% & Part IV, \S18, App.~O.3 \\
\hline
Carrier nesting & galactic source map & Closed within the specified leading EFT; momentum balance between active carriers closed by the joint source; interface completion conditional; microscopic projector open & active carriers partition the record algebra through sharp selectors $e_P$, with $\Pi_PX=E_P(e_PX)$ idempotent; the tower rule applies at a common parent record resolution; source-preserving reconstruction gives overlapping spatial profiles and a separate partition of unity $\chi_P$; descendants remain in the full Einstein source but do not receive separate transverse halos; the response is evaluated on $T^{\mathfrak A}=\sum_PT_P^{\rm cg}$, whose excess forces cancel identically, while per-carrier responses would self-accelerate a Milky-Way--LMC pair at $104\ {\rm km\,s^{-1}\,Gyr^{-1}}$; melting order and the stock amplitude remove the interface layer and the label-dependent forces in the tested pair & Part IV, App. N \\

Response stock and budget & galactic EFT / capacity ledger & Positive kinetic witness; transport contact conditional & Nonnegative spherical targets admit circular-orbit kinetic measures. The ensemble full-release budget is $e^{-\eta_*}(1-f_b)M_h$; the lensing benchmark uses $1-\eta_*$ and gives $\chi^2=86$ versus $75$ uncapped. A positive radial relaxation preserves angular marginals, while its physical target, energy supply, and coupled halo evolution remain open. & \S16, App.~N.8, O.3 \\
\hline
Positive source and thermal reconstruction & source map / fluctuations & Exact finite identities & A positive column-normalized map preserves mass, its adjoint fixes the probe coupling, and congruence preserves positive covariance and the common-temperature quantum thermal relation. Autonomous dynamics additionally require generator closure; hidden modes otherwise produce memory. & App.~N.8 \\
\hline

RAR law and static action & galactic EFT & Derived within the leading completion given the tilt mode; the tilt-mode identification and the Unruh reading of the lever are premises & the tilt mode in the horizon bath gives $H_{\rm cell}=\hbar\omega:N:=k_BT_Hx:N:$ and $\partial F_{\rm th}/\partial E=n_B(x)$; the QUMOND-type auxiliary action gives the conservative nonspherical field equation and the spherical exponential RAR; the lever coefficient is the register enthalpy over a one-nat oriented record at the Unruh temperature (\path{Entropic/RAR/Pendular.lean}), and that reading awaits a microscopic derivation & Part IV, App. N \\
\hline
Resolved-source Milky Way test & galactic EFT & Closed within the leading completion; amplitude and saddle-splay identifications conditional & The local contact gives $1792.2$ versus the measured $1294.8\pm54.0$. With an RAR-consistent exponential inner disk, the director completion with the stock amplitude at $c=0.03$ gives $1416$ (Wang stellar surface density) or $1315$ (McKee level), with 43-point $\chi^2_\Phi=137.1$ or $135.9$ versus $135$ for fitted NFW controls; the baryonic proxy at $c=1$ gives $1401$ or $1302$ ($135.7$ or $139.5$). The solution is the stiff radial limit, spans $20\ {\rm km^2\,s^{-2}}$ for $0.03\le c\le1$, and leaves the rotation curve and SPARC RAR unchanged at the stated precision. One auxiliary action yields the director, the excess flux, and the stiffness source; the saddle-splay coefficient is set by spherical exactness. These comparisons postdate the local-contact failure and are development-set postdictions. The flat-inner-disk radial deficit is a separate baryonic-model limitation. & Two independent director solvers (baryonic proxy); resolved SPARC check; data: \citep{Wang2026MW,McKee2015,Lelli2016}; App.~N \\
\hline
Disk vertical dynamics & galactic EFT & Empirical support; two direct old-disk cases & The excess carries no phantom sheet, so vertical forces are Newtonian plus a nearly round component. DiskMass integrated-light dispersions require $h_z^*/h_{z,\rm DMS}=0.76$ after the modeled population correction, against $0.55$ for QUMOND. The NGC~6946 old-disk dispersions and rotation curve share one mass-to-light ratio at $1.0\sigma$, against $2.6\sigma$ for QUMOND. In NGC~628 a reduced vertical model calibrated on NGC~6946 gives one stellar mass for $\sigma_{\rm old}$ and the rotation velocity at $0.3\sigma$, against $2.9\sigma$ for QUMOND. & \S16; data: \citep{Martinsson2013a,Martinsson2013b,Angus2015,Aniyan2018,Aniyan2021} \\
\hline
Shape of the excess & galactic EFT & Round aligned branch favored by some probes; live tilted-stock candidate; population test open & $q_\Phi=0.95$--$0.99$ at $20$--$50$ kpc with population scatter at most $0.05$, aligned with the baryons, against $0.85$--$0.89$ for $\Lambda$CDM halos in the same potential definition. Milky-Way six-dimensional tracers favor the completion by a likelihood ratio of $20$--$90$ on statistical errors; the extragalactic stream population favors a $\Lambda$CDM-like halo by about $4$. Nine Milky-Way measurements with an inter-study scatter of $0.14$ give a Bayes factor of $16$ for the completion, or $4$ against a rounder $\Lambda$CDM prior. The $25^\circ$ debris tilt decays within $0.25\,$Gyr in the tested aligned potentials. Three conditional live-stock realizations retain tilted stellar structure for $8\,$Gyr; formation, force convergence, and a deformable disk remain open. & \S16; data: \citep{Wegg2019,Palau2023,Chemaly2026,Ibata2024,Woudenberg2024,DillamoreSanders2025} \\
\hline
Baseline no slip / lensing & weak-field metric & Closed for the ordinary longitudinal branch & varying the unreduced Einstein scalar constraint gives $\nabla^2(\Phi-\Psi)=0$ and hence $\Phi=\Psi$ under asymptotic flatness & Part III, App. D, N \\
\hline
Transverse metric and lensing & galactic metric response & Closed within the specified quasistatic contact; nonlinear covariance open & $h_{\mu\nu}^{\perp}=-2\phi_\perp(\bar g_{\mu\nu}+2u_\mu u_\nu)/c^2$ gives no slip, one dynamical/lensing density, and no extra vacuum graviton pole; covariant variation of the record projector remains to be derived & \S16, App.~N \\
\hline
Solar-System and wide-binary limit & local gravity & Closed at leading quasistatic order in the carrier-resolved EFT; nonlinear preferred-frame audit open & the Milky Way is the active carrier, so localized Solar and binary sources remain purely Einstein; $Q_2^{\odot,\perp}=0$, wide binaries are Newtonian apart from the Galactic tide, and the vacuum tensor kinetic term is unchanged; a descendant smoothed over $L$ gives $Q_2\propto(n_\parallel-n_\perp)GM/L^3$, so Cassini requires $L\ge0.024$--$0.049\,$pc & Part III--IV, App. F, N \\
\hline
Telegrapher relation $D/\tau_0=c^2$ & transport & Closed in canonical transport branch & causal closure & Part V, App. E \\
\hline
Canonical $\tau_0^{-1}=H_0$ branch & transport & Fixed in the minimal transport closure & no-new-IR-scale choice & Part V, App. E \\
\hline
Hubble-tension mechanism & cosmology & Structurally supported extension & homogeneous trace-coupled mode & Part V, App. E \\
\hline
Equilibrium capacity-vacuum subtraction & cosmological constant / lattice dynamics & Exact regulator identity; conditional continuum theorem & the vacuum-normalized factor $\widehat Z_{\rm cap}=e^{\beta\mu_{\rm cap}N_{\rm cell}}Z_{\rm cap}$ removes the complete capacity-sector bulk coefficient at every regulator and is energy-origin invariant; the uncentered control verifies the predicted $N_{41}$ coupling displacement; the exact two-simplex Regge volume fixes $Z_V^{\rm geom}=\overline v_4$ and, on the matched $G_{\rm R}=G_*$ branch, gives $t_*=\Lambda_{\rm R}L_*^2/(8\pi)$ without setting $a=L_*$; the common-parent continuum and initial gravitational condition $\Lambda_{\rm R}^{\rm empty\ equilibrium}=0$ are separate matching requirements; continuum existence and dark-energy dynamics remain open & \S25, \S24.9, App.~J.2, J.6, J.9 \\
\hline
Cosmological term from fixed measure and collective source-clock release & dark energy & Exact finite allocation; conditional history; direct linear benchmark & Composition preservation uniquely allocates the drain and gives $a^3\bar C=B e^{-\eta_*F}$. A selected low-density action supplies a constant energy floor and regular current. Direct geodesic perturbations give $\Delta\chi^2=+0.85$ with one fewer parameter. The feedback-resolved retention occupation gives $\Omega_{\Lambda,0}=0.709\pm0.02$ against $0.707$, and the frozen history with no fitted dark-energy parameter gives $\Delta\chi^2=+1.07$ with four cosmological parameters; an independent mixing kernel and the nonlinear spatial completion remain open. & \S20, App.~O.3, O.12--O.15 \\
\hline
Gauge compatibility of the capacity spine & UV consistency & Same-factor Gauss reading excluded; separate-factor covariance closed kinematically & the certified floor $K^2\ge122/3>0$ excludes a Gauss-Casimir reading; the microscopic completion uses $\mathcal H_{\rm geom}\otimes\mathcal H_{\rm cap}\otimes\mathbb C^2_{\rm orient}$, so every capacity operator commutes strongly with the geometric gauge action; an equivariant $V_3$ frame lock preserves global rotational covariance without changing $1680$, the $K^2$ spectrum, or any downstream coefficient & \S6, \S23, App.~Q \\
\hline
Orientation doublet & UV geometry / matter interface & Persistent orientation excluded by renewal; pure-gravity decoupling closed; fermion decoupling closed in the record-oriented frame & a persistent lift would retain at least $\ln2$ of mutual information and is therefore incompatible with full-support renewal; the two $\mu=+1$ lifts share the metric and Regge phase, so the proper-gravity vertex acts as the identity on the refreshed doublet; the inversion $e\to-e$ reverses the sign of the Dirac action and no spinor map compensates, so fermion branches in separated cells coupled to refreshed lifts would lose relative phase in atom interferometers within nanoseconds; coupling fermions to the frame time-oriented by the record order makes the refreshed bit inert; parity-odd couplings require the spatial orientation of the chiral sector; that orientation is the coherent orientation of the gluing, which persists through renewal and, in the record-oriented frame, is a lift-class invariant & \S23, App.~H.9, H.9a, P.8, Q \\
Native lattice fermion & matter sector & Open: the one-spinor-per-cell operator fails; the Dirac--K\"ahler placement is relativistic and doubler-free with a dimension-five term excluded by mirror symmetry of the operator and ensemble, and its form sectors share one light cone only with flat-realization weights & on the diamond graph the spin-$\tfrac12$ lift of the native gluing carries one isotropic Dirac cone, $v=ta/\sqrt3$, with mirror-parity chirality, three nodal-loop doublers whose Weyl charges sum to zero, and a gap from a handedness-weighted term; regular cells cannot realize that graph; on every closed balanced complex the uniform zero modes have no first-order dispersion; the body-centred disphenoid complex gives sixteen zero modes in two isotropic linear families with a first-order spin splitting $|\eta|\le0.64$, above the Crab bound of $10^{-5}$; generic complexes give a quadratic touching; the Dirac--K\"ahler field has Betti-number zero modes, fits with $|c_1|/v<10^{-5}$ in cell units, and sector speeds equal to one with Voronoi--Delaunay weights; with regular-cell weights its sector speeds differ by $2$--$5\%$ on the body-centred honeycomb and by $17$--$49\%$ on random and balanced complexes & App.~H.9a \\
\hline
Record dynamics, organization, and computation & history sector & Exact algebraic core plus named conditional interpretations & renewal archives exactly; the canonical defect's source cost and record-discrimination rate are one likelihood ratio; weak fixed-marginal correlations give \(\mathcal I_H\propto D_+^2\) while the growing mode rises; the primitive face interaction generates \(\mathfrak{su}(16)\) and composes to \(\mathfrak{su}(4^N)\) on connected networks, and the gauge projection acts trivially on the capacity factor; identity-free capacity cannot record stream assignment; collective baryonic progenitor mixing supplies the conditional carrier-loss history, and the one-signed drain also inherits O.13's ledger and orientation premises & App.~P, O.10--O.13 \\
\hline
Saturated-phase committed capacity & cosmology & Conditional transition; zero-pressure constraint theorem closed within the pinned coherent branch & source-labelled saturation is $\sigma_A=1$ while absolute availability remains nonzero; one coherent phase advancing per electron-calibrated tick gives $X=1/2$ without a separate premise; the constraint action then gives $p=0$, $c_s^2=0$, and $\rho\propto a^{-3}$; this is a capacity stress form, not a particle ontology; formation is specified in O.15, and the formation vertex remains open & Part V, \S20, App.~M, O.15 \\
\hline
Pre-release committed-capacity abundance $\Omega_c/\Omega_b=1/\epsilon$ & cosmology & Conditional on the lever normalization, pinned reading, and per-defect bookkeeping & $\sigma_A=\epsilon M_{c,A}/M_A$ is defined on each Lagrangian source label; the selected formation map fills the cap, $\bar\sigma=1$, under its stated premises (O.15); tied at the formation value $5.333$, the matched linear-lensing release benchmark is competitive with $\Lambda$CDM on Planck 2018 high-$\ell$ CMB, DESI DR2, and DES-Dovekie with one fewer cosmological parameter; the joint Boltzmann evolution of the homogeneous mode remains open & Part V, \S20 \\
\hline
Formation and no late recruitment & cosmology / record dynamics & Closed for the specified transfer with the carrier-loss ledger rule; formation vertex open & formation copies the Lagrangian identity of unlabelled conserved baryonic charge, fills the cap $M_A/\epsilon$ on $\Sigma_0$, and draws $\rho_b/\epsilon$ from the plasma; thermal pairs receive no allocation; $n_\varnothing=0$ after $\Sigma_0$ and no written operation raises it; $\Delta_A$ is nondecreasing on each stream branch; the per-tick formation weight must exceed $e^{-16.1g_{\rm share,eff}}$ & App.~O.15 \\
\hline
$a_0(\Omega_c/\Omega_b)=cH_0$ & cross-sector structure & Exact conditional identity in the pre-release mass convention & multiplication cancels $\epsilon$; conversion to the later mass convention and release depletion must be included for present dust; the acceleration lever and committed-abundance premises remain separate & Part IV, Part V, \S20 \\
\hline
Post-fixation tests of $\epsilon$ & empirical support & Mixed empirical results; no combined significance assigned & $a_0$ is $+2.6\%$ from the local RAR scale at the stated reference $H_0$; the direct geodesic fit has $\Delta\chi^2=+0.85$ with one fewer cosmological parameter; cluster profiles supply positive-stock targets, with A1795 and A2029 requiring relative segregation; the redshift comparison uses the Appendix E inference and reference projection specified in \S27.3; these comparisons share theory premises and are not statistically independent & \S15, \S18, \S20, \S27.3 \\
\hline
Collective source-clock carrier loss / conditional conservation & cosmology & Exact no-stream, renewal, fractional-stock, and resolved-position monotonicity lemmas within stated premises; coarse kernel measured at low resolution; gate and labels open & shell crossing and phase drift do not release identity-free capacity; additive information pricing fixes the released stock fraction and faithful renewal fixes unit rerouting for each failed block; the covariant capacity clock fixes the exact-gradient exchange and geodesic Euler equation; a first SPH measurement gives $K_{\rm mix}=1-(1-K_0)e^{-\tau/(f_{\rm tail}\,t_{\rm dyn})}$ with $K_0\simeq0.5$ and $f_{\rm tail}\simeq2$--$6$ on three coarse partitions; a resolved zoom must confirm it and fix the gate and labels & Part V, \S20, App.~O.7--O.14 \\
\hline
Cluster cohort normalization & cluster sector & Fixed budget; transport conditional & $k_C=5.333\,e^{-\eta_*}=5.1760752$ in the O.15 mass convention fixes the closed cohort. Aperture baryons and stock can have different retained fractions; their histories require forward evolution. & \S18, App.~O.3 \\
\hline
Reciprocal finite-work contact & effective dynamics & Conservation identities and controlled encounters & Central pair transitions conserve linear and angular momentum; work and heat balance kinetic changes. The floor gates inaccessible transitions. Physical target, rate, and range remain selected inputs. & \S18, App.~O.3 \\
\hline
Observed radial clusters & empirical targets & Positive nominal NFW inventories; coextensive outliers & Twelve X-COP profiles, seven stellar profiles, and five independent lensing masses. A1795 and A2029 require relative segregation under their central lensing targets; A644 and ZW1215 require reconstruction/estimator uncertainty. & \S18 \\
\hline
Bullet lensing and gas separation & cluster dynamics & Effective mechanism demonstrated; resolved fit open & Weak-contact encounters separate gas from stock near galaxies. The released Bullet map supplies aperture masses and a 3.92 kpc halo--galaxy target, below the illustrative simulation resolution. & \S18 \\
\hline
Depleted capacity & covariant effective action & Positive resting state on a selected regular branch & A constant floor gives $\rho\to K$, $p\to-K$, and vanishing source-clock current. The tested background stays on the unchanged dust branch. Constitutive selection, null interfaces, and nonlinear evolution remain open. & App.~O.3 \\
\hline
Bounded capacity $q$, $N^2=q$ & strong field & Fixed static constitutive rule; covariant generalization open & multiplicative composition and weak-field matching fix the static lapse map; a generic lapse is foliation dependent & Part VI, App. F, N \\
\hline
ADM multiplier action & strong field & Excluded as a parent completion & with no independent $q$ dynamics, variation gives $\lambda=0$ and only renames the foliation-dependent lapse; adding dynamics introduces an extra mode & Part VI, App. F, N \\
\hline
Spherical parent reduction & strong field & Closed metric reduction; clock equality restricted to vacuum & \(q_{\rm geo}=(\nabla R)^2\) is the radial mass invariant and gives Schwarzschild in static vacuum. Clock capacity \(q_\tau=N^2\) generally differs inside matter; its lapse equation includes energy and stress. & Part VI, App.~F, N.6 \\
\hline
Capacity-exhaustion horizon $q=0$ & strong field & Geometric zero closed in spherical symmetry; domain termination conditional & $q_{\rm geo}=0$ is the marginal sphere; identifying it with substrate saturation and excluding $q_{\rm geo}<0$ requires the bounded-domain postulate and boundary theory & Part VI, App. F, N \\
\hline
Hawking temperature and exterior ringdown & strong field & GR-matching branch closed; capacity boundary condition open & Schwarzschild exterior fixes the wave operator and Euclidean temperature; $\mathcal R=0$ follows if future-horizon regularity is retained, while $\Gamma_{\partial q}$ must determine the physical substrate reflectivity & App. F \\
\hline
Bekenstein--Hawking / channel-area bridge & gravity / cut entropy & Exact bare relation under the angular rule; physical calibration conditional & Coherent bonds give $\ln2$ per cut. One angular tensor gives replica and Regge weights, hence $\hbar G_0/c^3=A_{\rm ch}/(4\ln2)$. Matching sets $G_0=G_*$; event calibration gives a provisional $1.03002G_*$. The native bulk/cut mismatch remains. & App.~C.8, F, J.11, Q.11 \\
\hline
Horizon formation and boundary microphysics & strong field & Geometric marginal-surface formation closed; saturation dynamics open & GR collapse can form $q_{\rm geo}=0$; a bounded causal $q_{\rm cap}$ evolution, its equality to $q_{\rm geo}$, the boundary action, and relaxation spectrum remain to be derived & Part VI, App. F \\
\hline
Rotating / charged stationary exteriors & strong field & Baseline metric solutions closed; capacity map open & Einstein / Einstein--Maxwell gives Kerr / Reissner--Nordstr\"om / Kerr--Newman, but no nonspherical covariant capacity scalar or domain rule has yet been derived & App. F \\
\hline
Finite retained-sector quantum reconstruction & quantum foundations & Closed for finite retained factors in the selected stable quadratic completion; network/global extension conditional & H.11 derives the exact transported $\mathbb C^9/\mathbb{CP}^8$ Hilbert bundle, tensor composition from multiplicative labels and complete commuting local algebras, the Born norm from reversible equivalence and physical-record additivity, finite instruments and no-signaling, and the finite history Gram kernel; with $\|A_R\otimes A_L\|=41/5$ a network whose bridges survive has a quasi-local infinite-volume algebra, Lieb--Robinson dynamics, and GNS representations; record completeness remains a named probability premise, while the continuum/Fock realization, cosmological state, actual decohering histories, and durable global record capacity remain open & \S3.3, \S22, App.~G, H.11 \\
\hline
No-signaling in finite operational branch & quantum foundations & Derived & reversible local record coupling gives a trace-preserving reduced channel; summing a remote instrument leaves the local marginal independent of the remote setting & Part VI, App. G, H.11 \\
\hline
Record-retention reading of Many-Pasts & quantum foundations & Closed as an ontological consequence of Postulate III and the finite record measure & the realized present includes its physical records; erasing fine record distinctions pushes the joint measure forward by ordinary coarse-graining; $p(h\mid P)$ is formed after the joint dynamics and supplies no future-to-past force & \S3.3, \S22, App. G.6 \\
\hline
Arrow of time & quantum foundations & Renewal H-theorem, orientation and record root closed; macroscopic arrow conditional on typicality; cosmological completion conditional & renewal replaces the closure register each tick; the persistent geometry--defect channel is unitary (H.11), which preserves coherence; record-writing subsystems contract relative entropy; code-preserving gates divide the closure code (576 sectors on two glued cells), so its mixing comes from renewal; the boundary is the minimal-record root, whose depth the surprisal of the present record bounds; the gravitational entropy of the root and the global cosmological state remain open & \S22, App.~G.7, P.8--P.9 \\
\hline
History capacity and record root & history sector & Non-reuse and record root closed; de~Sitter bounds obeyed by the present factor; physical source open & exact recoverability gives $\dim\mathcal H_{\rm hist}\ge\prod_nd_n$, so destructive recycling contradicts the dilation; the positive Born weight of the realized present bounds the recoverable past by its surprisal, so the record order is well founded with a minimal-record root; the written history, of order $10^{243}$ nats, exceeds the de~Sitter area capacity $3.3\times10^{122}$ by about $3\times10^{120}$, so it carries no area-law count; the N-bound and the type II$_1$ maximum constrain the retrievable present factor, every operator of the completed local action commutes with exported history, and the construction denies a finite de~Sitter Hilbert space and excludes global recurrence & App.~P.2, P.9 \\
\hline
Late-time re-entry & cosmology / history sector & Exact zero in the specified dynamics; thermal bound for number-changing completions & $[H,n_m]=0$ for every channel and release is one-signed; committed stock cannot be regenerated (O.15); records written in matter obey $P_{\rm false}(K)\le\min(e^{-K},e^{-KE_{\rm nat}/k_BT})$ with $E_{\rm nat}/k_BT_{\rm dS}\simeq3.9\times10^{39}$ & App.~P.10 \\
\hline
Microstructure Hamiltonian & UV realization & Finite marked-transfer action and capacity rigging lift closed; selected-branch Einstein theorem closed; host selection underdetermined & the native capacity vertex and marked routing are finite and audited; the selected $V_3$ phase has an exact maximal-fusion map, locked coherent phase space, and, conditional on the factorized realization, a unique equivariant frame lock; $(H_{\rm TT})$ and an exact primitive realization are closed, while a unique geometric gluing tensor is excluded as a consequence of the present capacity premises; a face-wise reading of the lock fixes the host class and bare couplings (H.9a, J.10) & Part VI, \S23, App.~B, H, Q \\
\hline
Charged-lepton spectrum & particle-sector extension & Conditional on the pole-mass identification of the shell operator (App.~I.1) and the additive-generator attachment & closure-spectrum collapse gives three shells; the baseline shell algebra is dressed by $Z_\mu=1+\zeta_*$ and $Z_{\tau,2}=1+(2/7)\zeta_*$, giving $m_\mu/m_e=206.768280237$ and $m_\tau/m_e=3477.343310$, both within $1\sigma$; the dressing was constructed after the baseline residuals were known (App.~L) & \S13, App. H--I, L \\
\hline
Abelian / weak gauge hosting & gauge sector & Coherent extension; Standard Model identification external & baseline redundancy supplies the Abelian gauge template; primitive spin data carry an $SU(2)$ action, while weak chirality, doublets, and hypercharge assignments are not derived & App. I.2 \\
\hline
Electromagnetic extension & gauge sector & Specified extension; native photon and electroweak records open & compact $U(1)$ links with an exact Gauss law and closed-graph neutrality; two transverse modes of the quadratic symbol; the global order parameter of the native $\mathbb{CP}^8$ gluing exhibits hysteresis on the $6^4$ lattice; the native leptons induce $\alpha^{-1}\simeq30$; the $SU(2)_L$ and $U(1)_Y$ stiffness targets use the diamond transfer form $t_{\rm adj}=1-2(1-k)/(d+1)$ that gives $t_8=13/14$, with targets $\alpha_2^{-1}=49.15$ and $\alpha_Y^{-1}=55.08$ at $m_P$; a common heat time and the trial $\alpha^{-1}=138.452$ fail & App. S \\
\hline
Persistent open-route color algebra & gauge sector & Conditional derivation from one open-route premise & a reference-preserving open three-route record is a reversible qutrit channel; quotienting the channel phase gives $PU(3)$ with $\mathfrak{su}(3)$ Lie algebra and eight traceless endpoint operators; phase-invariant channel-return fidelity gives the adjoint Wilson form & App. I.2 \\
\hline
Primitive color transfer $t_8=13/14$ & gauge sector & Fixed within the minimal identity-kernel completion & primitive $j_0=3/2$ recoupling gives $G=(9/10)P_{\mathbf1}+(21/20)P_{\mathbf2}$; maximum faithful throughput and route-frame isotropy give $\mathcal E=1_{\mathbf1}\oplus(13/14)I_{\mathbf8}$ & App. H.9, I.2 \\
\hline
Diamond causal color regulator & gauge sector & Conditional on the minimal product causal stack and equal primitive heat times & four-link mixed loops, six-link spatial loops, the embedded diamond volume, and electric--magnetic isotropy give $a_t/d=4\sqrt2/3$ and $g_{\rm graph,HK}^2=2\sqrt6\,g_{\rm HK,step}^2=0.24203562353$; graph-to-continuum matching remains open & App. I.2 \\
\hline
Physical QCD matter sector & gauge / particle sector & Open and outside the present premises & the global $SU(3)$ endpoint lift and central $U(1)$ identification, triality-carrying quark defects, two weak species per generation, hypercharge, a chiral lattice fermion action (H.9a tests the native one-spinor-per-cell operator), finite scheme matching, confinement, and hadron observables are not derived by the projective pure-gauge route construction & App. I.2 \\
\hline
Vacuum stiffness from the substrate & lattice dynamics & Excluded, five branches & induced curvature coupling $c_0=+0.019$ against bare $k_0\simeq2.2$; refresh stationary state geometry-blind to exponential accuracy; history-tilt ceiling twenty times below requirement; free conserved field induces near-universal spanning-tree entropy; budget saturation mass suppresses the residual soft modes & \S24.3, App.~J.4 \\
\hline
Vacuum geometry in the lattice realization & lattice dynamics & Open; externally hosted & tested substrate mechanisms do not select an extended phase; the Regge/CDT host supplies the vacuum geometry, while Many-Pasts conditions histories only after a present record is specified & \S24.4, App.~G, J.4 \\
\hline
Capacity-decorated CDT phase / continuum tests & geometric embedding & Open, with explicit tests; finite-volume host behavior measured & at the same tuned couplings, increasing $N_{41}$ from about $42{,}000$ to $102{,}000$ raises $d_H$ from $3.31$ to $3.56$ and the blob score from $1.46$ to $1.57$; this is consistent with movement farther into the extended finite-volume region but does not establish phase membership; $d_H\to4$, de Sitter-profile scaling, and especially a massless TT transfer sector remain to be measured, while intersection with a continuous critical manifold is the separate stronger cutoff-removal test; the predicted native point $(2\ln2,\simeq0)$ at $\tilde\alpha\simeq1.05$, tested with criteria fixed before data, gives seed-dependent basins and fails the bare-measure volume-pair criterion (J.10); the connected $\alpha=7$ branch has registered stiffness agreement, while its start-independence test remains undetermined (J.14) & \S24.9, App.~J.1, J.9, J.10 \\
\hline
Connected CDT normalization & geometric embedding & Registered finite-volume agreement; matter response open & the $\alpha=7$ primary ratio is $G_{\rm vol}/G_0=1.064$, with $\sigma_{\ln R}=0.089$, from the sole eligible 32k run A4; reported-only stationary 16k and 64k runs give $1.060$ and $0.893$; Q1 start independence and the separate trace test are undetermined & \S24.10, App.~J.14 \\
\hline
Ensemble fracture under local dynamics & lattice dynamics & Verified & explicit construction: $48$ components ($4!\times2$) of $35$ states each under injectivity-preserving unit shifts, conserved (ordering, parity) charge & \S24.3, App.~D.4, H.6, J.4 \\
\hline
Substrate contribution to the host volume term & lattice dynamics & Measured as a scaling family, demonstration volume & closure free energy per cell, computed from $\eta_*$ alone, displaces the equilibrium volume by the predicted amount: single point $-62.2$ predicted, $-62.8$ measured; eight coupling--volume-penalty combinations span $-10$ to $-158$ at ratios $0.96$--$1.07$ with one point at $0.89$, with the concave $\beta$-curve and $1/\varepsilon$ scaling confirmed; the shuffled control lies on its separately predicted line at ratio $1.00$; curvature-sector shifts remain to be measured with comparable statistics & \S24.4, App.~J.6 \\
\hline
Admissibility weighting on dynamical geometry & lattice dynamics & Finite-lattice compatibility measured & at $N_{41}\simeq100{,}000$, three matched eighty-slice ensembles remain connected and correctly foliated with $0.6\%$ total-volume spread; the collision fraction changes from $0.654$ without the weighting to $0.077$ with it, while the shuffled-energy control remains $0.736$; the corresponding $d_H$ values remain $3.55$--$3.59$, and a forty-slice repetition gives collision fraction $0.078$ & \S24.5, App.~J.5 \\
\hline
Local response of geometry to persistent closure failure & lattice dynamics & Demonstrated in two independent matched comparisons & for each of two random seeds at $N_{41}\simeq100{,}000$, one hundred failure defects are compared with one hundred identically anchored closed cells, with every marked cell surviving: first-shell closure-energy and collision separations are $16.8\sigma$ and $9.2\sigma$; third-shell coordination separates at $4.0\sigma$; shell-cell count and first-/second-shell coordination do not separate significantly after autocorrelation and seed corrections & \S24.6, App.~J.7 \\
\hline
Long-range propagation of the strain field & lattice dynamics & Conditional Ward theorem; finite-lattice source-to-field behavior measured; continuum junction open & on a separate forty-slice realization at $N_{41}\simeq100{,}000$, conservation is exact; $Q-Q(0)=(0.1440\pm0.0025)m$; the dressing-subtracted shell-one deficit per charge is $4.9965\pm0.5244$ across $m=1,2,3,6$; the profile is $d^{-0.62\pm0.11}$ through eight shells, a smeared $1/r$ source with far-field offset reproduces it at rms residual $0.018$, and the screening fit is consistent with zero, giving $\xi>4.6$; the $24.2\%$ maximum source-linearity residual and nonideal exponent are explicit finite-size and saturation diagnostics, while carrier identification, critical scaling, the Dirichlet-form spectral limit of the slice Laplacian, and $\chi_q\alpha_{\rm maint}/D_q\to8\pi G_*/c^2$ ($8\pi$ in native units) remain open & \S24.7--24.9, App.~J.8--J.9 \\
\hline
Formal and numerical consistency checks & validation layer & Kernel-checked mathematics and enclosures; numerical diagnostics & A Lean 4.35.0 library of 222 modules audits 5{,}038 theorems. The claim ledger records 882 statements as proved, six of them given a cited classical result, and 90 in part, with physical premises and numerical checks classified separately. Closure analysis, finite Gaussian integrals, record maps, and conditional response identities are checked. & App.~K, P.11, R.1 \\
\hline
EFT consistency checklist & field-theory audit & Supportive audit layer & ordinary longitudinal branch has no extra scalar degree of freedom; static source sign and quadratic form checked; telegrapher stability belongs to the conditional transport completion & App. D \\
\hline
\end{longtable}
\endgroup

This table is the epistemic map used for the rest of the discussion.

\paragraph{Cosmological row.}
The per-source variable $\sigma_A$ separates saturation of recruitment
from exhaustion of absolute capacity. Conditional on $\sigma_A=1$ and
one coherent tick-normalized phase, $X=1/2$ follows and the
zero-pressure constrained-capacity stress form is a theorem.
Appendix~O.15 specifies formation, its epoch, and its energy source;
the formation vertex remains open. The converted pre-release abundance is $5.333$. Direct geodesic-clock perturbations with matched lensing give $\Delta\chi^2=+0.85$ for the linear-depletion benchmark with one fewer cosmological parameter. The composition-preserving law gives exponential depletion and positive internal components on the total-stock domain. Its direct linear likelihood is competitive; an independently fixed physical release history remains open.


\section{27. Falsifiability and Observational Tests}

\paragraph{Joint empirical checks of the selected branches.}
The empirical companion carries a common-coefficient audit across local
and distant galaxies, clusters, expansion, primary CMB spectra, lensing,
stellar systems, and gravitational waves. Several tests sharpen the
remaining physical requirements. For 263 KMOS3D, KGES, and KROSS catalogue
objects, the outer-radius KMOS3D log-acceleration residual slope is
$-0.245$ dex per unit redshift with constant $a_0$ and $-0.313$ with the
common $H(z)$ law; the latter has median log-velocity residual $-0.087$ dex.
These central-value diagnostics retain a high-redshift tension and require
a resolved kinematic likelihood. In 130 SPARC galaxies with 2710 points,
profiling the stated mass, distance, and inclination priors gives a
penalized score of 11258, or 3433 with a $5\,\mathrm{km\,s^{-1}}$ velocity
floor. The median per-point $\chi^2$ in the latter calculation is 1.54.
For 100 flat-velocity galaxies the baryonic Tully--Fisher log residual has
median 0.020 dex and rms 0.067 dex. These are separate observational
reductions from the director comparison in Section~16.

The exponential release fit gives
$a_0=1.2588\times10^{-10}\,\mathrm{m\,s^{-2}}$ from its fitted $H_0$.
The bare reciprocal abundance identity acquires the specified formation
mass-convention and survival factors for the observed late-time stock;
its elementary empirical test is limited mainly by galaxy mass calibration,
at roughly twenty percent. The ACT primary-spectrum and lensing checks
give respectively $162.73$ for 135 bins and $16.21$ for ten lensing bins,
against $162.02$ and $17.86$ for the reference cosmology. Their overlapping
information is not added to the Planck fit. The predicted $E_G$ values
are $0.4375$ at $z=0.316$ and $0.3862$ at $z=0.555$, calculated from
the Weyl potential and baryon velocities. A no-slip identity alone does
not determine them. The implementation's reference-dust split makes its
ordinary total-matter density contrast unsuitable as the physical total-matter
$\sigma_8$ without reconstructing the capacity contribution.

Under the selected carrier prescription, Solar-System and wide-binary
tests contain ordinary Galactic tides and no independent binary halo.
Radial tidal estimates at 10--30 kAU for one to two solar masses span
$\Delta v/v\simeq1.1\times10^{-5}$--$6\times10^{-4}$; orientation,
selection, and contaminants are required for a Gaia likelihood. The
minimal constant-Planck-mass Einstein action with minimally coupled
scalar matter has the tensor quadratic action
\[
S_T^{(2)}=\frac{M_{\rm Pl}^2}{8}\int dt\,d^3x\,a^3
\left[\dot h_{ij}^{\rm TT}\dot h_{ij}^{\rm TT}
-\frac{\partial_kh_{ij}^{\rm TT}\partial_kh_{ij}^{\rm TT}}{a^2}\right].
\]
It has two luminal tensor modes and positive kinetic energy. Public
GWTC-4 polarization comparisons are compatible with this branch; they
do not select it over GR. The native continuum and possible curvature
operators still require matching. The finite routing and Farey tests of
the particle coefficients establish uniqueness only within their
enumerated menus and preserve the known discovery history.



At the common cosmological acceleration normalization, the reduced
old-disk solver gives NGC~628 a completion $\chi^2=2.04$ for five
dispersions and a $-0.21\sigma$ radial--vertical mass offset; its
approximate algebraic QUMOND comparison gives $8.95$ and $-2.91\sigma$.
The gas profile, inclination, scale height, and population selection
remain physical uncertainties. A round conventional halo can also give
similar vertical gravity. In the saved satellite grid, increasing the
model-error allowance from ten to thirty percent changes the Crater~II
stock/NFW evidence from $2.6\times10^5$--$2.0\times10^8$ to 94--155,
and Antlia~2 from 0.007--0.019 to 0.69--0.74. The prior and error model
therefore control much of the quoted evidence. These satellite trajectories
do not include the conservative contact.

\subsection{27.1 Static weak-field falsifiers}

The static galactic sector has seven direct checks. Rotation curves test the exponential transition and baryonic Tully--Fisher limit \citep{McGaugh2016,Lelli2016}; resolved disks test the director completion of the nonspherical equation; old-disk vertical dispersions test the absence of a phantom sheet; galaxy--galaxy lensing tests the no-slip metric; Cassini tests the carrier nesting rule; dispersion-supported dwarfs test the boundary of the equilibrium carrier phase; and stellar streams test the predicted round excess component.

Stacked weak lensing first established a low-acceleration continuation and reported an early/late-type split \citep{Brouwer2021}. A later joint analysis extends the kinematic RAR by about $2.5$ decades in acceleration and finds a common early/late relation after imposing strict isolation and using consistent stellar and gas masses \citep{Mistele2024}. Those controls remove that split. A persistent mismatch between dynamical and lensing accelerations in an isolated sample would falsify the leading no-slip contact.

The local prediction is sharper. The Milky Way is the active transverse carrier, while the Sun and wide binaries are resolved Einstein descendants. The branch predicts $Q_2^{\odot,\perp}=0$ and Newtonian wide-binary dynamics apart from the smooth Galactic tide. Cassini gives $Q_2=(1.6\pm1.8)\times10^{-27}\,{\rm s}^{-2}$; current universal total-field implementations are in $3$--$15\sigma$ tension with that result \citep{Park2026}. Wide-binary analyses remain statistically divided \citep{Pittordis2025,Chae2025,Chae2026}. A confirmed velocity boost above $1\%$ at separations of $10$--$30$ kAU, more than twenty times the largest tidal estimate, would reject carrier nesting as stated, while a significant Solar quadrupole correlated with the Galactic field would reject its local projection rule.

Systems of several carriers give three further tests:
\begin{itemize}
\item a Milky-Way reflex toward the LMC well below the $\Lambda$CDM-like value would reject an inertial response stock;
\item tidal dwarfs inside their parent's record must stay nearly Newtonian, with dynamical-mass boosts of 1.5--2 for the NGC~5291 system;
\item the residual forces on the response stocks predict small offsets between the lensing and baryonic centres of satellites.
\end{itemize}

At lower baryonic mass, the GravSphere dwarf sample of \citet{Julio2025} lies systematically above the extrapolated rotating-galaxy RAR and shows multivalued loci and large scatter. The current EFT restricts the exact exponential law to equilibrium active carriers; it must derive a phase-selection boundary before that restriction can explain the dwarf data. A universal application of the same law to dispersion-supported ultrafaint dwarfs is already disfavored by these observations.

The resolved director completion adds three galactic failure modes. With the Wang snail potentials confirmed, a local dynamical excess density at $R_0$ departing from the predicted $0.35$--$0.37\ {\rm GeV\,cm^{-3}}$ within $1\,\mathrm{kpc}$ of the plane and $0.29$--$0.30\ {\rm GeV\,cm^{-3}}$ at $|z|=3\,\mathrm{kpc}$ by more than the baryon-model uncertainty would reject its Milky-Way solution. For face-on disks it predicts the vertical dispersion of the old disk measured apart from the young cold layer. For the DiskMass galaxies at $h_{z,\rm DMS}$ and the rotation-curve mass-to-light ratio, the predicted old-disk dispersion exceeds the single-component value by a median factor $1.61$ (16--84\%: $1.34$--$1.95$), against $1.87$ for QUMOND. A median ratio of measured to predicted old-disk dispersion below $0.8$ or above $1.2$ would reject the completion; separating it from QUMOND at $3\sigma$ requires about $5\%$ precision on that median. The per-galaxy predictions are registered in machine-readable form with SHA-256 \texttt{5a72f3e1a3ae0b659963a339ae54bd92253c34e2cd76f932903a9f6ad9b74427}. For the aligned equilibrium excess it predicts $q_\Phi=0.95$--$0.99$ at $20$--$50\,\mathrm{kpc}$ with population scatter at most $0.05$ and alignment with the baryons. A disk-modeled stream population with mean $q_\Phi<0.85$ or scatter above $0.1$, or a confirmed tilt of the inner halo relative to the disk that the LMC does not produce, would reject this aligned branch. In the same potential definition typical $\Lambda$CDM halos lie at $0.85$--$0.89$, so a population mean of $0.93$ or above with small scatter would disfavor them. Nine Milky-Way measurements, combined with their inter-study scatter, favor the completion by a Bayes factor of $16$ (Section~16), the projected-track extragalactic sample mildly favors $\Lambda$CDM \citep{Chemaly2026}, and a decision requires about 80 well-measured streams. For the adopted long-lived Gaia--Sausage--Enceladus debris, Section~16 finds the aligned solution inadequate and constructs a separate tilted-stock candidate. Its three live phase realizations establish conditional retention over $8\,$Gyr. They do not alter the registered aligned-branch prediction or determine a population distribution of halo shapes; those require the stock's assembly and relaxation dynamics. A converged self-consistent realization with a live disk that holds the stellar debris tilt below $19^\circ$ after $8\,$Gyr, two standard deviations under the measured $25.4^\circ\pm3.2^\circ$, rejects the tilted-stock candidate.

The stock selection adds four tests:
\begin{itemize}
\item for a fixed baryonic target and $a_0$, relaxation reduces residuals from the exponential law with time since the last major accretion; longer outer orbits can retain larger residuals at a given age. Evolution of the target adds a history-dependent lag;
\item isolated ellipticals follow the law out to the radius where their hot gas outweighs their stars;
\item the support of an isolated system ends where its hot gas outweighs its condensed baryons, and the stacked lensing profile follows the exponential law inside that radius and the collisionless stock outside it; isolated-lens profiles binned by hot-gas mass at fixed stellar mass test the support radius, and an isolated finite stock that follows the law to projected radius $R$ needs $S\ge\tfrac\pi4\sqrt{a_0/g_{\rm bar}}-1$ stock masses per baryonic mass;
\item isolated dwarfs whose progenitors crossed the release gate lie on the law, and dwarfs whose progenitors stayed below it keep CDM-like central densities.
\end{itemize}

\subsection{27.2 Dynamical falsifiers}

The dynamical extension has two independent observational tests. \emph{Cluster-stock} tests ask whether relaxed clusters hold the cohort ratio $1+k_C=6.18$ outside the central galaxy's support, whether the excess over the galactic law in groups follows the uncarried stock profile (Section~18), and whether resolved merger maps place the lensing peaks with the stars and stock while the gas lags. \emph{Transport} tests ask whether $D/\tau_0=c^2$ and the reciprocal contacts evolve the stock correctly through a merger. Systems such as the Bullet Cluster \citep{Clowe2006} probe both. A failure of the propagation law would reject the causal completion even if the static branch survived.

\subsection{27.3 Cosmological falsifiers}

Cosmology presents a different kind of test. The question there is whether a full Boltzmann treatment allows the trace-coupled homogeneous mode to reduce the sound horizon without spoiling the CMB or structure-growth observables. If it cannot, the cosmological extension fails on its own terms. The empirical target is set by the current measurement spread: early-universe inferences near $67.4$ \citep{Planck2020}, distance-ladder determinations ranging from $\simeq70$ \citep{Freedman2025} to $73$ \citep{Riess2022}, and a tension whose proposed resolutions are reviewed in \citet{DiValentino2021}.

\paragraph{Hubble constant from galaxy dynamics.}
The baryon density alone fixes the present energy budget: the tie, the release with the feedback-resolved retention gate, and flatness give $H_0=69.3\pm2.1\,\mathrm{km\,s^{-1}\,Mpc^{-1}}$ with no fitted dark-sector parameter (Appendix~O.12), between the Planck-$\Lambda$CDM and Cepheid-calibrated values. Since $h$ varies nearly as $\omega_b^{1/2}$, a baryon density from big-bang deuterium gives the same value within its uncertainty. With the homogeneous mode at matter--radiation equality (Section~19) the theory gives $H_0=69.4\pm0.6\,\mathrm{km\,s^{-1}\,Mpc^{-1}}$ and $a_0=1.27\times10^{-10}\,\mathrm{m\,s^{-2}}$; a local value of $73$ or above with an uncertainty of $1\,\mathrm{km\,s^{-1}\,Mpc^{-1}}$ excludes the mode together with the abundance tie at $3\sigma$. The release fit of Section~20 gives $H_0=68.46\,\mathrm{km\,s^{-1}\,Mpc^{-1}}$, with the early-universe inferences and below the Cepheid-calibrated local value. The relation $a_0=\epsilon cH_0$ has no free coefficient, so a galactic measurement of $a_0$ is an independent determination of the Hubble constant, $H_0=a_0/\epsilon c$: $H_0=68.46$ requires $a_0=1.250\times10^{-10}\,\mathrm{m\,s^{-2}}$, and $73.04$ requires $1.334\times10^{-10}\,\mathrm{m\,s^{-2}}$. Only galaxies with absolute distances carry this information. With Hubble-flow distances the inferred $a_0$ scales as $H_0^{p}$, with $p$ between $1$ in the Newtonian regime and $2$ in the deep regime, and that dependence nearly cancels the predicted $a_0\propto H_0$. A fit of the exponential law to the 44 SPARC galaxies \citep{Lelli2016} with tip-of-the-red-giant-branch, Cepheid or supernova distances, with per-galaxy mass-to-light ratios, distances and inclinations marginalized under their priors, gives $a_0=(1.21\pm0.19)\times10^{-10}\,\mathrm{m\,s^{-2}}$ and $H_0=66\pm11\,\mathrm{km\,s^{-1}\,Mpc^{-1}}$ from a galaxy bootstrap. All 153 galaxies of quality 1--2 give $(1.03\pm0.11)\times10^{-10}\,\mathrm{m\,s^{-2}}$, about $2\sigma$ below $1.250$. The subsamples range from $0.97$ for gas-dominated dwarfs to $1.37$ for the massive absolute-distance galaxies, beyond their statistical errors; pressure support, beam smearing and inclination errors in dwarfs are candidate causes, and a residual dependence on galaxy type would test the universality of the law. A $2\%$ determination of $a_0$ separates $68.5$ from $73$ at $3\sigma$. At the present per-galaxy scatter it needs about $3000$ galaxies with absolute distances, with the subsample drift controlled at the same level. A common distance-scale error $k$ enters the ladder value as $k^{-1}$ and the deep-regime galactic value as $k^{-2}$, so $H_{0,\rm ladder}^2/H_{0,\rm RAR}$ is independent of $k$ under this relation. A precise $a_0$ inconsistent with $\epsilon cH_0$ for the measured $H_0$ falsifies the horizon lever.

The relation $a_0(z)=cH(z)g_{\mathrm{share,eff}}/(4\pi^2)$ predicts measurable redshift evolution. For a Planck-like background, $H(2.3)/H_0\simeq3.46$. A disk with $g_{\mathrm{bar}}\simeq2\times10^{-10}\,\mathrm{m\,s^{-2}}$ then has $g_{\mathrm{obs}}/g_{\mathrm{bar}}\simeq2.02$, compared with $1.39$ for an epoch-independent scale. Observed outer rotation curves at these redshifts indicate strong baryon dominance \citep{Genzel2017,Lang2017}, although pressure-support and stacking systematics leave the comparison unsettled. Controlled observations showing no stronger boost than matched $z=0$ systems would falsify $a_0\propto H(z)$. Hydrodynamical $\Lambda$CDM simulations also produce a characteristic acceleration that rises with redshift. In the Magneticum runs it grows by a factor of about three between $z=0$ and $z=2.3$ \citep{Mayer2022}, and EAGLE shows much weaker evolution \citep{Ludlow2017}. The direction of the trend therefore does not identify the transverse closure. Given the horizon lever and a stationary occupation of the slow mode, the closure fixes $a_0(z)/a_0(0)=H(z)/H_0$ with no free coefficient, $1.79$ at $z=1$ and $3.46$ at $z=2.3$ for the reference background. The lag of the slow mode (Section~15) is unquantified and would modify these ratios. A $3\sigma$ separation of $3.46$ from a ratio of $3.0$ requires a total uncertainty near $4.5\%$, including gas and mass calibration and the uncertainty of the alternative prediction. Systems younger than the relaxation time of the stock selection lie between the CDM relation and the exponential law, so the comparison must use relaxed samples. The electromagnetic coupling carries no horizon factor: the transfer multipliers depend only on record amplitude and dimension. The theory therefore keeps $\alpha$ fixed while $a_0$ follows $H(z)$. At $z=1.15$, where $a_0$ is predicted to exceed its local value by a factor $1.95$, quasar absorption gives $\Delta\alpha/\alpha=1.3\pm1.3_{\rm stat}\pm0.4_{\rm sys}$ parts per million \citep{Murphy2022ESPRESSO}.

The equilibrium-vacuum route requires a common-parent critical limit, a defined volume operator, and a physical gravitational boundary condition. The finite-regulator capacity centering is an exact normalization identity. A nonzero capacity bulk coefficient in that normalized reference would contradict its definition; a nonzero gravitational integration constant would instead test the separate initial matching in O.5. A common substrate can support several effective operators. The subsequent equation of state and release dynamics require covariantly conserved stress energy and a completed perturbation prescription.

A MUSE sample of $79$ star-forming galaxies at $0.33<z<1.44$ finds a radial-acceleration scale that rises with redshift \citep{Ciocan2026}; a resolved low-redshift H\,\textsc{i} sample tentatively reports the same direction \citep{Varasteanu2025}. The MUSE paper contains distinct model-dependent inferences. Its main-text dark-halo decomposition gives slope about $1.59\pm0.10$ in units of $10^{-10}\,\mathrm{m\,s^{-2}}$ and a representative $z\sim1$ scale $2.38^{+0.12}_{-0.10}$ in those units. Its Appendix E MOND-based kinematic inference instead gives intercept $1.03\pm0.05$ and slope $1.20\pm0.10$. For the reference background $H_0=67.4\,{\rm km\,s^{-1}\,Mpc^{-1}}$, $\Omega_m=0.315$, and $\Omega_\Lambda=0.685$, the uniform-grid linear projection of $a_0(z)=\epsilon cH(z)$ over $z=0.33,0.34,\ldots,1.44$ has intercept $1.002784$ and slope $1.225645$ with $g_{\rm share,eff}=7.41980002357$ and the stated background. Those coefficients agree with the Appendix E marginal intervals; they are not a fit to the individual galaxies or an agreement with the distinct main-text slope. A joint test needs the galaxy likelihood, its redshift weights and coefficient covariance, and the same fitted $H(z)$ used by the release cosmology.



\paragraph{Saturated-phase falsifiers.}
Three tests bind the saturated phase: a full Einstein--Boltzmann implementation must evolve the commitment history and constrained fluid with the declared abundance convention; the formation vertex of O.15 must be derived within its ledger and pinned-reading assumptions; and linear observables must remain compatible with the inherited dynamics. The separate sound-horizon proposal requires its own joint calculation. The subsequent release channel is more selective than collapse. The self-gravitating run produces a gate- and label-dependent kernel. The documented linear-depletion PPF benchmark requires $I_{\rm rel}\simeq94.2$ at its reference assembly inputs; the positive allocation uses the weighted integral and its own matched likelihood. An independent test must freeze the measured kernel and retention prescription and then compare their convolution with the same cosmological likelihood.

\paragraph{Dark-energy falsifiers.}
The sector of Appendix~O has a one-signed accumulation law. A reconstructed decrease of its physical $\rho_\Lambda$ with time, or $w_{\rm eff}>-1$ inferred from that same component after consistently accounting for matter transfer, would reject that law. A $w_0,w_a$ fit that assumes separately conserved CDM is a different inference and its crossing of $-1$ is not by itself this falsifier. The approximate PPF fit prefers $f\simeq0.283$ at fixed reference gates and the matched direct-perturbation fit $f=0.274245$; a measured kernel with nonzero loss at entry must be tested by its full convolution, not by equating its tail time with that $f$. Substantial pre-gate release, an independently frozen history incompatible with the matched likelihood, or an independent local determinant mode contrary to the unimodular premise would test distinct parts of the construction. The feedback-resolved occupation places the retention gate without an input adjusted to the vacuum density and gives $I_{\rm rel}=95.5\pm8$ against the target $94.2$ (O.12). A measured cold-gas retention fraction at $z\ge8$, a tighter reionization redshift and a converged mixing kernel narrow that interval, and an integral outside the narrowed interval falsifies the normalization. Retention of mixed, record-carrying gas in most halos below $10^7M_\odot$ at $z\simeq10$ would contradict the gate that the observed vacuum density requires.

\subsection{27.4 Correlated-constant falsifiers}

The inferred $L_*$, induced $G_*$, and matched weak-field $G$ are not independent legs: the matched route carries the same electron length calibration. Their agreement is a normalization audit rather than a separate falsifier. The marked action adds two independent tests because the same $\zeta_*$ enters the muon and tau through different graph polynomials. A future independent determination of $L_*$ would provide a fourth test of the shared vertex.

\subsection{27.5 Many-Pasts status}

Many-Pasts predicts no laboratory departure from the Born rule or
no-signaling, but those finite operational laws are derived within the construction.
On the selected stable retained sector, H.11 derives the complex Hilbert
bundle and exact moving-frame unitary transport, tensor composition, Born
weights, finite instruments, reduced channels, no-signaling, and the history
Gram kernel. Its remaining burdens are global: the multiplicative factors
and bridges must survive the constrained host into a network and
continuum/Fock theory; the cosmological state and the histories that decohere
must be fixed; and record completeness and the physical source of history capacity must be supplied. Renewal supplies an exact one-tick H-theorem (Appendix~P.8), the record root fixes the boundary (P.9), and late-time re-entry vanishes in the specified dynamics (P.10). Failure of the first group would remove the proposed global microscopic quantum completion; failure of the record root or of conditional typicality would remove the proposed arrow. Neither would create a
new laboratory deviation inside the finite theorem's domain.

\section{28. What the Theory Would Have to Get Wrong to Fail}

The failure modes are not all equally severe, and they are ordered here by how much of the theory each would remove.

\paragraph{Kills the core ontology.} A demonstrated failure of mass--entropy equivalence, an internal incoherence in the Many-Pasts weighting, or evidence that geometry cannot be read as the long-wavelength form of an entanglement-capacity substrate would remove the foundations on which everything else rests.

\paragraph{Kills the static capacity-response closure or its transverse completion.} A weak-field UV coefficient chain that cannot be reconciled with an independently validated microscopic derivation would break the ordinary static closure. An RAR transition shape that departs from the proposed bosonic law, old-disk vertical dispersions that require a phantom sheet, a disk-modeled stream population incompatible with the predicted round excess, persistent low-acceleration slip, unacceptable residual PPN effects, or violation of the metric Ward identity would falsify the transverse completion while leaving the ordinary Einstein/capacity equivalence intact.

\paragraph{Kills an extension only.} If the cosmological trace-coupled homogeneous mode cannot survive a full Boltzmann likelihood confrontation, the cosmology sector fails while the static weak-field branch stands. If the saturated phase fails any of its commitments (the stock-to-baryon abundance departing from the capacity ceiling, linear-regime observables departing from the inherited form, or no committed component surviving on cluster infall streams), the saturated committed-capacity reading fails in the same contained way. If the source clock is not carried by baryonic Lagrangian matter, if a tracer simulation finds insufficient or prematurely early progenitor-identity loss, or if the resulting mixing-weighted history fails a background-data fit, the dark-energy release extension fails without touching the weak-field branch. If \(q_{\rm cap}\) does not track \(q_{\rm geo}\) during collapse, or if the derived boundary response conflicts with black-hole observations, the bounded-domain interpretation fails while the spherical Einstein reduction remains. None of these touches the weak-field core.
\paragraph{Requires modification, not death.} A fuller graph calculation may revise the separate conditional stiffness self-energy, and strong-field boundary spectroscopy may change without altering the Einstein exterior. A failure of the decorated charged vertex would be more serious because the same $\zeta_*$ enters the electron scale and both heavier-lepton ratios. The positive-spectrum clock conversion survives only if another microscopic charged action replaces it.

One scale-setting commitment cuts across these tiers. The high-precision $G_*$ rests on seven-fold additivity, the record-conditioned determinant transfer, and the decorated marked vertex. A microscopic Hessian with a collective edge mode, persistent endpoint polarization, or off-diagonal charged propagation would fail the finite action and displace the correlated $G$ and lepton results. If the electron identification or mass--entropy ontology failed, the scale-setting derivation of Appendix~H would not go through.

\section{29. Comparison with Other Approaches}

The galactic excess and cosmological abundance are assigned to two phases of one medium. The comparisons below distinguish that proposal from nearby alternatives while retaining the conditional grades of Sections~15 and~20.

\subsection{29.1 Relative to \texorpdfstring{$\Lambda$}{Lambda}CDM}

The contrast with $\Lambda$CDM begins at the level of ontology. Here visible matter is interpreted as localized defects in a vacuum-capacity medium. The ordinary longitudinal response is a reduced representation of Einstein gravity; the proposed extra galactic response is carried by a transverse substrate sector. The UV entropy enters both branches, while the galactic branch additionally requires its tilt-mode lever, carrier, and metric-vertex matching conditions. Its leading lensing response is specified; cosmology remains a separate extension.

\subsection{29.2 Relative to MOND-like interpolation programs}

MOND-like programs \citep{Milgrom1983,SandersMcGaugh2002,FamaeyMcGaugh2012} usually begin from an acceleration law or interpolation function. Here the same law comes from a thermal occupation contact and a carrier-resolved auxiliary action. The source-space nesting rule removes the Solar cross-term found in universal functions of the total field, while the no-slip metric contact fixes the leading lensing response. The comparison turns on the microscopic derivation of those structures and on resolved nonspherical, dwarf, and lensing tests.

\subsection{29.3 Relative to Verlinde-style emergent gravity}

Verlinde-style emergent-gravity programs share the broad intuition that gravity may be entropic \citep{Jacobson1995,Padmanabhan2010,Verlinde2011,Verlinde2017}, but they are usually formulated through thermodynamic reasoning or horizon-inspired force laws. The present framework specifies tetrahedral counting, admissibility closure, edge coupling, return dressing, and Euclidean normalization before reaching the continuum EFT. Its correctness is an empirical question.

\subsection{29.4 Relative to TeVeS and other multi-field modified gravities}

Multi-field relativistic MOND completions such as TeVeS \citep{Bekenstein2004} introduce additional scalar and vector fields alongside the metric to obtain relativistic lensing and cosmology. Here the ordinary longitudinal response adds no field beyond the Einstein metric. The transverse sector is a carrier-conditioned metric contact with no additional vacuum graviton polarization; its leading stationary tensor structure gives no slip, while nonlinear covariance and the microscopic vertex remain open.

\subsection{29.5 Relative to AeST}

The closest modern comparator is the AeST theory of Skordis and Z\l{}o\'snik \citep{Skordis2021}, which combines the metric with a dynamical timelike vector and a scalar field and can reproduce MOND-scale galaxy phenomenology while fitting the CMB power spectrum. The present construction differs in its finite counting input and is less developed cosmologically: the joint Boltzmann treatment remains open here.

\subsection{29.6 Relative to scalar-tensor gravity}

The reduced capacity functional of Section~10 can be mistaken for the scalar part of a Brans--Dicke-type theory \citep{BransDicke1961}. The action reconstruction shows otherwise: in the ordinary static branch it is the Einstein constraint action after a field redefinition and carries no independent scalar. A scalar--tensor or disformal theory is retained only as a control case for sectors that might genuinely require an additional mode; if adopted, ordinary fifth-force, slip, and PPN constraints would apply.

\subsection{29.7 Relative to quantum-mechanical interpretations}

Because Many-Pasts occupies the role of an interpretation of quantum mechanics (Sections~3.3, 22), it should be placed against the standard options. It posits one realized present, adds no collapse term, and introduces no hidden sharp values. Its operational probability theory has the decoherent-histories form: unresolved alternatives retain amplitudes, decoherent record histories receive diagonal probabilities, and conditioning on the present occurs only after the alternative records are normalized. The new content is the ontology assigned to that conditional measure and its proposed substrate realization. For finite retained sectors, H.11 derives the Hilbert norm, tensor composition, instruments, no-signaling, and history Gram kernel. The continuum/Fock realization, cosmological state, actual decohering histories, and durable global record capacity remain open.

\subsection{29.8 Relative to CDT, spin foams, and group field theory}

The construction uses discrete tetrahedral boundary data, with primitive $j_0=3/2$ pairs whose positive oriented matching transmits the seven-state $V_3$ sector; this places it near the quantum-tetrahedron vocabulary of simplicial spin networks \citep{Barbieri1998,BaezBarrett1999,RovelliSmolin1995}. Its selection principle remains distinct from ordinary spin-foam/GFT programs: the finite capacity sector is fixed by a boundary-counting and faithful-resolution problem rather than by using a gravitational vertex amplitude as its statistical weight.

Group field theory treats spacetime as a many-body system of quantized tetrahedra, with simplicial quantum geometry as the second quantization of loop-quantum-gravity spin networks \citep{Oriti2009,Oriti2016}. Its condensate states describe homogeneous cosmology, and effective Friedmann dynamics with a quantum bounce emerge from the collective dynamics of the quanta \citep{Gielen2013,OritiSindoniWilsonEwing2016}. The present construction shares the tetrahedral quantum, the emergent continuum, and the condensate reading of the vacuum; Appendix~H.1 writes the cell ensemble as a GFT condensate of bosonic tetrahedral quanta coupled to fermionic defects. It adds a finite capacity factor with a fixed closure spectrum, a marked record register, and a release channel that sources the vacuum energy. The GFT interaction kernel and vertex amplitude of that realization remain to be selected.

An identification of the four capacity labels with the open Peter--Weyl magnetic indices is not required and is not adopted in the completed embedding. Appendix~Q instead uses
\[
\mathcal H_{\rm cell}
=\mathcal H_{\rm geom}\otimes\mathcal H_{\rm cap}\otimes\mathbb C^2_{\rm orient}.
\]
The geometric factor carries the standard gauge and intertwiner data of the simplicial amplitude; the finite capacity factor carries the injective four-port alphabet, $K^2$, the Gibbs weight, marks, and records. Their selected $V_3$ frames are related equivariantly, so geometry--capacity equivalence does not identify two distinct operators or force the capacity count to discard geometric intertwiners. Appendix~Q excludes the Gauss-index identification and proves the strong gauge invariance of the numerical spine.

This selected factorized completion supplies a direct spin-foam/GFT witness for the vacuum geometry.  The supporting external results have different jobs and are not combined into a single published theorem.  Proper-EPRL-type amplitudes isolate the gravitational Regge exponential in the semiclassical block sector \citep{EngleProper2013,EngleVilenskyZipfel2016}. The capacity matching supplies an exact canonical maximal-fusion ray $V_3^{\otimes N}\to V_{3N}$ on the selected channel, and the block-lock lemma propagates the chosen one-cell geometry--capacity lock along that ray.  Han-type semiclassical--continuum constructions independently reach Einstein gravity; in the controlled linearized construction the low-energy excitations give all smooth linearized Einstein solutions and the two graviton helicities \citep{Han2017,HanHuangZipfel2019}. The area-Regge continuum calculation gives the same leading graviton dynamics and an $O(a^2C^2)$ correction \citep{DittrichKogios2023}. Appendix~H.10 turns that selected-branch symbol into the strong-resolvent theorem $(H_{\rm TT})$ and pulls it back to an exact fixed-$j=3$ primitive realization.  It also proves that the finite capacity premises cannot determine the independent geometric gluing tensor or measure.  A derived $(H_{\rm phase})$ therefore requires an additional geometric principle before the phase test can even be posed.  Appendix~H.9a, paragraph~4a, names the one candidate already in the postulates, the face-wise reading of the relational lock, under which the host is the equilateral CDT class at predicted couplings (Appendix~J.10).  The factorized capacity decoration inherits any positive host spectrum by identity; an alternative nonfactorized vertex must preserve the same interface properties.

Section~24 supplies a numerical interface in which the capacity ensemble, extended by the finite collision penalty of Appendix~J, is coupled to a CDT host \citep{AmbjornJurkiewiczLoll2004,AmbjornGJL2012}. In the completed runs the host supplies the dynamical vacuum geometry and the paper tests how its capacity weighting dresses it. Section~24.9 and Appendix~J.9 first ask whether the fixed capacity slice lies in a stable extended phase with a massless TT transfer sector.  Intersection with a continuous geometric critical surface is the stronger optional cutoff-removal test, not the existence condition at physical $L_*$.  Success or failure of that CDT route is therefore a separate embedding test and does not alter the block-scale spin-foam witnesses.

\subsection{29.9 Relative to algebraic and information-geometric gravity}

Three continuum results address separate claims made here. The observer-dressed de~Sitter algebra establishes a maximum-entropy empty gravitational state. The horizon modular calculation of Dorau and Much equates the relative information of a coherent excitation with its Killing energy flux and, after the area input, with Einstein curvature. The geometric-relative-entropy action gives a local bulk information functional with an Einstein limit, a positive Legendre energy, cosmological thermodynamics, and de~Sitter area scaling \citep{ChandrasekaranEtAl2023,DorauMuch2026,Bianconi2025,Bianconi2026}. They support the continuum ontology but do not solve the ultraviolet problem. This paper must still derive the local bound, matter commitment, dimensional scale, and response coefficients from one finite microstructure. Bianconi's independent dynamical $G$-field belongs to a modified-gravity theory and is not imported into the metric-only longitudinal branch.

The coefficient comparison is explicit. On the matched branch $L_*^2=\hbar G_*/c^3$, the horizon normalization gives
\[
S_\infty=\frac{A_{\rm dS}}{4L_*^2}
=\frac{c^3A_{\rm dS}}{4\hbar G_*}.
\]
The observer-dressed de~Sitter construction uses this area coefficient and gives $S_{\max}-S_{\rm gen}(E)=\beta_{\rm dS}E+O(E^2)$ for a small central excitation. The horizon modular calculation uses $S_{\rm rel}=c^3\delta A/(4\hbar G_*)$ and obtains
\[
G_{\mu\nu}+\Lambda g_{\mu\nu}
=\frac{8\pi G_*}{c^4}\langle T_{\mu\nu}\rangle.
\]
The bulk route derived here gives
\[
\nabla^2\delta S=-\frac{\kappa}{\gamma}\rho,
\qquad
\frac{\Phi}{c^2}=-\frac{\delta S}{2S_\infty},
\qquad
G_* =\frac{c^2\kappa}{8\pi\gamma S_\infty},
\]
and hence $\nabla^2\Phi=4\pi G_*\rho$. The two continuum routes therefore use the same coupling. Appendix~F.5 fixes the microscopic channel-to-area normalization within the factorized coherent $j=3$ embedding. The result is a consistency matching, not independent evidence for the area law.

At Gaussian order, the relationship to Bianconi is sharper. Appendix~H.10 proves that the normalized determinant of any positive dressed GFT Hessian relative to its equilibrium Hessian on the same physical support has exactly the trace-log form used by geometric relative entropy, while its quadratic-source Legendre transform is exactly the positive covariance mismatch $\tfrac12\operatorname{Tr}'(\mathcal G-I-\ln\mathcal G)$. In the factorized completion the normalized equilibrium capacity trace contributes one to the empty vacuum amplitude, while the selected coarse geometric factor supplies the TT Einstein Hessian. This closes the Gaussian functional bridge and shows that the capacity decoration preserves the ordinary low-energy support.  The selected metric-Regge branch satisfies $(H_{\rm TT})$.  The coherent angular construction specifies a local geometric candidate, with a metric count operator and ordering. Its physical constraint space, refinement measure, and nonperturbative phase remain open.  The distinct canonical completion retains $(H'_{\rm can})$.

The cosmological scaling also agrees at the level of form. With $R_A=c/H$,
\[
S_\infty=\frac{\pi c^2}{H^2L_*^2},
\qquad
S_{\rm GfE}\simeq\frac{\overline\omega_{[1]}c^2}{\ell_P^4H^2},
\qquad
\frac{S_{\rm GfE}}{S_\infty}
\simeq\frac{\overline\omega_{[1]}L_*^2}{\pi\ell_P^4}.
\]
Bianconi therefore obtains the same $H^{-2}$ dependence, but her couplings do not determine the finite-cell coefficient used here.


\section{30. Conclusion}

One physical picture runs through the paper. The vacuum is a finite medium of
entanglement capacity; matter is that capacity tied up in stable, localized
defects; a particle's mass measures the entanglement its defect commits; and
gravity is the capacity strain the surrounding medium carries once that
commitment is made. The ordinary weak-field capacity action is the Einstein
constraint sector in the capacity variable, and the selected metric-Regge
completion has the massless two-helicity Einstein infrared limit. On galactic
scales, the same medium produces the additional long-range response usually
attributed to dark matter.

The central result is a tightly specified static weak-field construction with its remaining assumptions exposed. The tetrahedral ensemble fixes the sharing entropy; faithful renewal and the decorated marked-transfer vertex fix the electron-anchored substrate length; finite-state recurrence fixes the shell support factor; a common Gaussian scalar mode realizes the quadratic shell overlap; the edge kernel and Green-matched source map fix the ordinary response; and the separate minimal loop operator supplies the stated conditional stiffness correction. The weak-field bridge then rewrites the Einstein constraint sector in the capacity variable.

Newtonian gravity is the point-source limit of the Einstein/capacity action equivalence, and baseline lensing and PPN values follow from that parent. The carrier-resolved galactic EFT supplies the thermal Bose contact, conservative nonspherical action, source nesting rule, no-slip metric response, point-source lensing law, causal transport poles, and leading Solar-System and wide-binary limits. On the joint coarse source it conserves momentum between galaxies and satellites, and the response budget set by Lagrangian baryons passes the weak-lensing test. The radial-acceleration relation follows from the horizon-thermal tilt mode of the transverse frame. The tilt-mode identification and lever, the carrier projector, and the no-slip metric vertex remain matching premises awaiting a microscopic GFT derivation. The marked-transfer action gives $G_*$ at $-0.073\sigma$ and the muon and tau ratios at $-0.535\sigma$ and $-0.125\sigma$. The shell recurrence and overlap are explicit, while their identification with the physical pole-mass operator remains the particle premise. These are historical postdictions from one zero-continuous-fit action, not three blind predictions.

Several constructions connect these sectors. Isometric reduction preserves the retained operator algebra and its source-dependent response. Positive reconstruction preserves source mass, adjoint forces, and the quantum noise relation. Composition-preserving allocation keeps the internal capacity components nonnegative until total-stock exhaustion, and its background integrates to exponential depletion. A selected multistream allocation conserves total energy--momentum and has a smooth expanding beam realization. Positive spherical stock admits an explicit kinetic equilibrium, while radial relaxation can preserve selected angular information. A finite-energy reset also separates the present register's energy budget from its unbounded exported history. The contact that selects the stock target, the independently fixed mixing history, the physical archive supply, and the common interacting continuum remain matching requirements.

The same medium extends into sectors the paper holds more tentatively. Transport gives the field a finite propagation speed and permits relative motion between gas and retained stock in cluster mergers. The cluster sector transports a positive inertial stock under a common potential, with a conservative finite-work contact and observed radial and merger targets. The same stock reaches the galactic equilibrium through contacts biased by the Bose relative entropy of the response cell. With carrier records identified with condensed baryons, cluster stock lies outside every carrier support and keeps the cohort ratio. The separate homogeneous-mode proposal could reduce the sound horizon and raise the inferred Hubble constant; that effect has not been established by the current release fit. In the saturated early universe the medium becomes a conserved committed density that gravitates as zero-pressure committed capacity. It forms once, when conserved baryonic charge first admits a source label, and no later process regenerates it; tied at that formation value, the exponential-allocation fit with linear lensing is competitive with $\Lambda$CDM on Planck high-$\ell$ spectra, DESI, and DES-Dovekie with one fewer cosmological parameter. Direct linear perturbations and the positive-allocation likelihood are implemented, while the independently normalized release history and nonlinear spatial completion remain open. In strong fields the exact spherical reduction identifies the composite $q_{\rm geo}$ and recovers Schwarzschild, while domain termination, boundary microphysics, and the nonspherical capacity map remain open. These are extensions and frontier completions, not closed results, and the closure table marks the distinction.

The microscopic construction has also been tested on a dynamical simplicial lattice. At matched $N_{41}\simeq100{,}000$, the admissibility weight strongly orders the microscopic labels while the host geometry remains extended and nearly unchanged. Two independent comparisons between identically anchored failure defects and closed cells give local capacity separations of $16.8\sigma$ and $9.2\sigma$ and a third-shell coordination separation of $4.0\sigma$ after autocorrelation and seed corrections. A separate conserved-field calculation finds exact conservation, approximately additive emergent source charge, one junction ratio across the tested source strengths, a power-law tail through eight shells, and no detected screening. Increasing the host volume from about $42{,}000$ to $102{,}000$ raises both the Hausdorff estimator and the concentration of volume in the extended region. Extended-phase membership, a massless TT transfer sector, the optional critical-surface test, and the absolute $G$ junction remain to be measured. Section~26 records the present grades.

The normalized capacity measure is invariant under extensive shifts of microscopic energy origin, and the uncentered control tests the predicted lattice-coupling displacement. The cosmological branch adopts zero initial gravitational vacuum density, additive source-record pricing, a common record/closure energy conversion, and the post-refresh ledger. Its constrained capacity-clock action fixes the background exchange and linear scalar perturbations for any supplied release history. Common closed-cohort advection--diffusion makes fully resolved position-information loss bounded and monotone; the physical coarse kernel must be measured and checked separately. The first progenitor-tagged run and fitted-history likelihood are available, and with the retention gate from feedback-resolved simulations they give $\Omega_{\Lambda,0}=0.709\pm0.02$ against the required $0.707$. The frozen history fits the cosmological data with four parameters at $\Delta\chi^2=+1.07$. A converged kernel must be fixed before the background, growth, lensing, and interacting CMB spectra can be predicted from the measured mixing history.

Three independent continuum constructions support different parts of the proposed ontology. The type~II$_1$ de~Sitter algebra realizes an entropy-maximizing empty gravitational state. The horizon modular calculation relates an excitation's relative information to energy flux and, after the area input, to Einstein curvature. The geometric-relative-entropy action gives a local bulk information dynamics with an Einstein limit, a positive mismatch energy, a first law, and $H^{-2}$ de~Sitter entropy \citep{ChandrasekaranEtAl2023,DorauMuch2026,Bianconi2025,Bianconi2026}. Section~25 derives the same positive mismatch functional from a Gaussian covariance identity. These results support the continuum ontology but do not derive the finite cell, coefficients, cosmological state, or ultraviolet completion used here.

The geometric sector has a finite coherent realization within the
existing capacity state space. Its bonds preserve the local closure
ensemble, and one angular tensor fixes the bare cut and curvature
coefficient. Connected chain and cycle patches, with full Gram
normalization, give an exact encoded causal triangulation weight
at the integer channel ratio \(1:3\). The canonical reconnection
construction supplies a distinct selected metric operator, with
linear Einstein response, the same tree-level Newton coefficient,
and nonlinear Kasner evolution. Its signed count histories prevent
a direct entrywise identification with the positive covariant weight.

The quantum count construction fixes the DeWitt measure and metric
Weyl ordering. Its differential commutators quantize the classical
spatial-regulator brackets, while spatial and count tests resolve
separate errors. The homogeneous volume--shape constraint admits a
positive distributional rigging form; the inhomogeneous physical
state space remains unresolved. The retained matter construction
places orbital motion, the native rest gap, and a positive clock
on a finite geometry transfer. Its inverse cochain metric preserves
the complete joint operator through the tested refinements.
Selection of a common continuum phase and an interacting
\(G_R\) still requires geometric, matter, and source tests on one
ensemble.

Many-Pasts assigns a record-conditioned ontology to the decoherent histories
compatible with the present. Its operational branch changes no laboratory
prediction, and the finite formalism is reconstructed from stated operational
premises. On the real retained $Q$ carrier, the positive Lorentzian Hamiltonian
selects the initial polar complex structure. The full symplectic flow
transports it exactly through varying frames, even though a comparison with
the instantaneous polar basis contains real squeezing; the forced
Born--Huang term is kept in that full generator. The non-common frequency and
moving frames generate $\mathfrak{su}(9)$ and yield
$\mathbb C^9/\mathbb{CP}^8$. Multiplicative joint labels and two complete
commuting local matrix algebras force the tensor product, after which the
projected neighboring bridge supplies the finite entangling algebra.
Reversible internal equivalence and additive complete physical records fix
the Born norm. Reversible system--record coupling then gives POVMs, Kraus
operators, conditional updates, trace-preserving reduced channels, and
no-signaling; the finite decoherence functional is the Gram kernel of the
resulting history branch vectors.

The past is represented at the resolution retained by present physical
records; erasing a record coarsens the same joint measure, and conditioning on a record after it forms supplies no future-to-past dynamics. Renewal gives an exact one-tick relative-entropy H-theorem and orients time, while coherent sectors evolve unitarily; exact recoverability forbids recycling of history capacity and fixes a minimal-record root; and in the specified dynamics the released vacuum cannot re-enter the committed phase. Faithful
full-support resolution selects the memoryless dressing kernel. The decorated
vertex prepares its diagonal fresh closure state, reversible dilation exports
the old replaceable register into history, and the marked spectrum plus
electron anchor fixes the clock and phase frequency. The response motif,
persistent coherent marked fiber, and charged-defect shell object are
kept distinct: the event incidence identifies their coordinates but does not
turn the transient response occupation into a conserved charge. The full
continuum/Fock realization, cosmological state, actual decohering histories,
record completeness and the physical source of history capacity, electromagnetic charge and Maxwell dynamics, relativistic quantum fields, the gravitational entropy of the record root, and the formation vertex remain open. None is a new founding premise.

The particle appendix also contains a conditional projective color sector: a persistent open tetrahedral route, transported reversibly, gives the $PU(3)$ qutrit-channel group, its $\mathfrak{su}(3)$ algebra, the adjoint Wilson form, and the fixed primitive transfer $t_8=13/14$, and the same primitive fusion reproduces the nine-state representation required by the marked present/history fiber (Appendix~I.2). The global $SU(3)$ endpoint lift, chiral quark matter, continuum-scheme matching, confinement, and hadron observables remain outside the completed derivation.

In ordinary gravity, the static capacity equation includes pressure
and spatial stress through the lapse source. Clock capacity and
radial mass geometry agree in normalized static vacuum and differ
inside matter. The matched scalar normalization, chain-clock area
calibration, and native profile stiffness remain separately testable.
The selected metric-Regge branch has an explicit reflection-positive
Fock representation whose one-particle states are the two massless
helicity-two gravitons (Appendix~H.10).
The face-wise CDT campaign retains its failed bare-measure
volume-pair criterion. The connected \(\alpha=7\) construction
defines a different host; the archived held-source diagnostic
does not determine its infrared Newton coefficient. Fixed-background
renewal preserves orbital coherence, while differing geometric
event records and the physical source of history capacity remain open.
The refreshed orientation doublet is inert on the pure-gravity branch
and, with fermions coupled through the frame oriented by the record
order, in the Dirac sector. A relativistic native fermion needs a
Dirac--K\"ahler-type placement, which is doubler-free and whose form
sectors share one light cone when its weights are those of a flat
realization (Appendix~H.9a); the chiral gauging remains open.

In the galactic sector, one auxiliary action with the record-potential director stiffness, together with the stock selection functional for the stock amplitude, reproduces the Milky-Way vertical potential for an RAR-consistent inner disk and the SPARC RAR, removes the phantom sheet from disk vertical dynamics, and predicts a round excess component. The remaining theory work includes a time-dependent implementation of the already reciprocal joint-source rule for interacting carriers, a microscopic derivation of the director stiffness together with the cell Hamiltonian, carrier projector, and metric vertex, full nonlinear variation of the carrier records, and a bath calculation of the equilibration rate and noise correlation volume. Old-disk vertical dispersions, disk-modeled stream shapes, the tilt of the Gaia--Sausage--Enceladus debris, resolved nonspherical lensing maps, and the ultrafaint-dwarf phase boundary are the immediate empirical tests. The finite-loop stiffness operator still needs its microscopic return-operator audit. Cosmology needs a resolved progenitor-tagged zoom that fixes the retention gate and progenitor labels together with the mixing kernel; the documented approximate likelihood constrains the release integral near $94.2$, and a first low-resolution run reaches that value at a $14\ {\rm km\,s^{-1}}$ gate selected in the analysis. The positive-allocation likelihood is implemented; an independently measured retention gate remains to be tested. The conservative cluster contact needs physical selection and a resolved lensing-map fit, and the strong-field boundary needs its saturation dynamics and spectroscopy. In the foundations sector, the physical source of history capacity and the gravitational entropy of the record root remain open. The color branch still needs a derived chiral matter action and graph-to-continuum matching before comparison with physical QCD.

One ultraviolet count fixes the coefficients of the closed sector in advance. A measurement inconsistent with any of those linked coefficients would falsify that sector. The calculations above specify those comparisons.

\clearpage
\section*{Appendix A: Symbol Dictionary and Canonical Conventions}

Appendix A gathers the conventions used throughout the technical material that follows, fixing the units, field definitions, and couplings in one place before the denser calculations begin.

\paragraph{Plain-language terms.}
Several physical words recur throughout the paper and are collected here in plain form before the symbols:
\begin{itemize}
\item \emph{Capacity}: the entanglement support locally available in the vacuum medium.
\item \emph{Defect}: a stable, localized commitment of that capacity; coarse-grained, a particle.
\item \emph{Deficit}: capacity no longer freely available near a defect, $\delta S=S_\infty-S_{\mathrm{ent}}$.
\item \emph{Capacity strain}: the extended deficit profile whose fractional value gives the weak-field potential and whose gradient gives the gravitational field.
\item \emph{Committed capacity}: capacity locked into transferring on behalf of a defect; the conserved carrier of the saturated phase (Section~20).
\item \emph{Dressing}: the cloud an elementary defect builds by resolving the boundary sectors; its determinant transfer and marked response set the length in the decorated scale branch.
\item \emph{Saturation}: the regime in which the capacity bath has no slack left, with the available transfer channel at its ceiling.
\item \emph{Pinned reading}: the statement that the saturated phase holds exactly at that ceiling and stays there.
\item \emph{Admissibility}: the weighting that favors boundary configurations close to a regular, isotropic local cell.
\item \emph{Closure}: the condition that the four oriented face data sum to zero, so the cell closes into a regular volume; $K^2$ measures the failure of closure.
\item \emph{Many-Pasts}: the postulate that one recorded present is supported by its compatible decoherent pasts. Their operational probabilities are the diagonal decoherence-functional weights conditioned on that present; $e^{-D(h,P)}$ is only shorthand for those normalized quantum probabilities.
\end{itemize}

\subsection*{A.1 Units, signature, and entropy normalization}

All dimensional quantities are expressed in SI units unless noted otherwise. The metric signature is $(-,+,+,+)$. Covariant spacetime integrals use $x^0=ct$, so the Einstein--Hilbert coefficient is $c^3/(16\pi G)$; after writing $dx^0=c\,dt$, the ADM coefficient is $c^4/(16\pi G)$. Entropies are measured in nats, so Boltzmann's constant is absorbed into the entropy normalization. The canonical reference discretization assigns the paired primitive cell spatial volume $V_*=L_*^3$, two tetrahedral sites, and time-coordinate four-volume $\Delta V_4=L_*^4/c$. Its site density is $2/L_*^3$. In this convention $\gamma_Q=4\hbar cJ/(3L_*^2)$ and $\kappa/\gamma=3L_*/(4G_{\mathrm{tet}}(0)\kappa_m(L_*))$ are simultaneously exact (C.4--C.5). The physical event volume $v_{\rm ev}$ in length coordinates and the physical site volume are calibrated separately in C.8; changing them requires rematching lattice stiffness and source density. In the decorated electron-anchored marked-transfer branch,
\[
L_*=-\frac32\,Z_e\lambda_e\ln\!\left(1-e^{-7g_{\mathrm{share,eff}}}\right),
\qquad
\lambda_e=\frac{\hbar}{m_ec},
\qquad
Z_e=(1+\zeta_*)(1+7\zeta_*^2).
\]
The conventional Planck length $L_P=\sqrt{\hbar G/c^3}$ is used only as a comparison scale or in standard gravitational thermodynamic expressions after the gravitational scale has been identified.

These conventions matter because the argument repeatedly moves between a dimensionless ultraviolet counting problem and a dimensionful continuum EFT. The units and signature make those descriptions directly comparable.

\subsection*{A.2 Core scalar variables}

The canonical continuum variable is the vacuum-relative coarse-grained entanglement field
\[
S_{\mathrm{ent}}(x),
\]
with vacuum baseline $S_\infty$ and deficit
\[
\delta S(x)=S_\infty-S_{\mathrm{ent}}(x).
\]
For nonlinear work the bounded occupancy fraction is
\[
q(x)=\frac{S_{\mathrm{ent}}(x)}{S_\infty}=1-\frac{\delta S}{S_\infty}\in[0,1].
\]
The absolute entropy unit is fixed only after choosing a cell or horizon normalization. Under a constant rescaling of $S_{\mathrm{ent}}$, the quantities $S_\infty$ and $\kappa/\gamma$ rescale together, leaving $\delta S/S_\infty$ and $\kappa/(\gamma S_\infty)$ invariant.
The source channel is
\[
\chi(x)=-\frac{T^\mu{}_\mu}{c^2},
\]
which is the continuum trace channel of the localized defect sector and reduces to the ordinary mass density $\rho$ in the nonrelativistic static limit.

\subsection*{A.3 Couplings and derived observables}

The main-text conventions are
\begin{align}
\gamma &:\ \text{entanglement-field stiffness},\\
\kappa &:\ \text{continuum defect--entropy coupling},\\
\kappa_m(\ell) &:\ \text{mass-per-entropy map at scale }\ell,\\
\zeta_* &\equiv 9e^{-g_{\mathrm{share,eff}}}
\left(1-\frac{8\eta_*}{49}\right)^{21/2},\\
Z_e &\equiv (1+\zeta_*)(1+7\zeta_*^2),\\
L_* &\equiv -\frac32\,Z_e\lambda_e\ln\!\left(1-e^{-7g_{\mathrm{share,eff}}}\right)
\quad\text{in the decorated marked-transfer branch},\\
G_* &= \frac{c^3L_*^2}{\hbar}
=\frac94\,\frac{\hbar c}{m_e^2}Z_e^2\ln^2\!\left(1-e^{-7g_{\mathrm{share,eff}}}\right),\\
G_{\mathrm{tet}}(0) &:\ \text{tetrahedral on-site Green constant},\\
g_{\mathrm{share,max}} &= \ln(1680),\\
g_{\mathrm{share,eff}} &:\ \text{admissibility-weighted sharing entropy},\\
J_{\mathrm{bare}},\,J_{\mathrm{eff}}^{\mathrm{tree}},\,J_{\mathrm{eff}}^{(\mathrm{ren})} &:\ \text{UV edge couplings},\\
a_0 &= \frac{cH_0 g_{\mathrm{share,eff}}}{4\pi^2}.
\end{align}
The canonical weak-field bridge and Newton closure are
\[
\begin{aligned}
\frac{\Phi}{c^2}=-\frac{\delta S}{2S_\infty},
&\qquad
\frac{\kappa}{\gamma}
=
\frac{3L_*}{4G_{\mathrm{tet}}(0)\kappa_m(L_*)},\\
G&=\frac{c^2\kappa}{8\pi\gamma S_\infty}.
\end{aligned}
\]

Collected in one place, these formulas also make clear which quantities are downstream of the closure chain. The UV data determine the stiffness and source-to-stiffness ratio first; the observable weak-field constants appear after the bridge and fixed $S_\infty$ normalization are applied.

\subsection*{A.4 Notation map}

One notation set is used throughout. The effective sharing entropy is denoted $g_{\mathrm{share,eff}}$, the scalar variable is always the vacuum-relative field $S_{\mathrm{ent}}$ or its deficit $\delta S$, and the weak-field bridge is used in the single form stated above. The principal extension-sector symbols are:
\begin{align*}
\epsilon&\equiv \frac{g_{\mathrm{share,eff}}}{4\pi^2},
&\nu(x)&\equiv \frac{1}{1-e^{-x}},\\
x&\equiv \frac{E_\perp}{k_BT_H},\\
D&:\ \text{capacity diffusivity},
&\tau_0&:\ \text{transport relaxation time},\\
D/\tau_0&=c^2,\\
\sigma_*&\equiv \pi/g_{\mathrm{share,eff}},\\
a_{\mathrm{UV}}&\equiv 1/\operatorname{Var}_{\eta_*}(K^2).
\end{align*}
In the transverse EFT, $\mathcal R_P$ is the retained record algebra of carrier $P$, $\Pi_P$ its conditional-expectation record projector, $T_P^{\rm cg}=\mathcal C_\ell\Pi_PT$ its reconstructed source, $\chi_P$ its spatial interpolation weight, $\Phi_{b,P}$ its Newtonian auxiliary potential, and $\phi_{\perp,P}$ its excess potential. The variables $z_P=|\boldsymbol\nabla\Phi_{b,P}|^2/a_0^2$ and $x_P=z_P^{1/4}$ enter the thermal occupation $n_B(x_P)$ and the auxiliary function $\mathcal Q_\perp(z_P)$. The one-capacity Hessian symbols $\Delta C$, $u_C$, $A_L$, $A_T$, and $C_\times$ are retained only as static thermodynamic consistency data in Appendix~N; they are not mode frequencies.

\section*{Appendix B: UV Boundary Ensemble and Admissibility Closure}

The ultraviolet counting problem is finite. A seven-state tetrahedral ensemble, weighted by its closure defect, determines a unique admissibility-closed entropy. Sections~5--6 use this result.

\subsection*{B.1 Tetrahedral package and logical status}

The canonical ultraviolet cell is a tetrahedron, the minimal volumetric simplex in three spatial dimensions. Its four faces are distinguishable relational ports. Each side of a shared face carries a primitive half-integer representation $V_{j_0}$, and the two sides are compared only after one has been transported and orientation-reversed into the frame of the other. Appendix~B.4 shows that the positive oriented-matching operator on
\[
V_{j_0}\otimes V_{j_0}
\]
transmits the maximal coupled multiplet $V_{2j_0}$.

The value of $j_0$ has two possible grades. In the selected sharp reversible
vertex, unitarity identifies the complete marked carrier with the nonmaximal
fusion output, and the representation content fixes
\[
j_0=\frac32,
\qquad
V_{2j_0}=V_3,
\qquad
|M|=7.
\]
For a different microscopic vertex, each half-integer $j_0$ defines a
discrete zero-parameter branch, and Appendix~B.5 compares those branches
through the same downstream map.

The counting convention is then fixed by the cell ontology. A letter is the occupation of one channel in a single-copy cell resource, so one letter cannot be routed through two ports. The ports remain distinguishable, and the two global orientations are distinct. The state count and its raw entropy are therefore
\[
\Omega_{\mathrm{tet}}=2P(7,4)=1680,
\qquad
g_{\mathrm{share,max}}=\ln(1680)=7.42654907240.
\]
The closure weighting below acts on this finite ensemble. Appendix~B.4 derives the primitive pair selection and the conditional $7+9$ complement structure; Appendix~B.5 records the neighboring discrete counts and preserves their postdictive status.

\subsection*{B.2 Closure invariant, kernel, and unique fixed point}

The canonical scalar closure invariant is
\[
K^2(b)=48-\frac{1}{3}\bigl(S^2-\Sigma^2\bigr),
\qquad
S=\sum_{i=1}^4 m_i,
\qquad
\Sigma^2=\sum_{i=1}^4 m_i^2.
\]
The admissibility family is
\[
p_\eta(b)=\frac{1}{Z(\eta)}e^{-\eta K^2(b)},
\qquad
Z(\eta)=\sum_{b\in B}e^{-\eta K^2(b)}.
\]
Stationary normalized closure evidence fixes the admissibility precision
$\eta$. Classical closure would set the three-component oriented-face sum
$\mathbf c(b)=\sum_i\hat n_i m_i\in\mathbb R^3$ to zero. Its squared
magnitude is $|\mathbf c(b)|^2=(4\Sigma^2-S^2)/3$. The invariant $K^2$ used
here is instead the quantum expectation of the squared closure operator. Each
face carries the seven-state $j_{\rm eff}=3$ representation ($2j+1=7$), and
the product boundary state
$|b\rangle=\bigotimes_{i=1}^4|j,m_i\rangle_{\hat n_i}$ uses the tetrahedral
normal frame $\hat n_i\!\cdot\!\hat n_j=-\tfrac13$. The closure operator
$\widehat{\mathbf C}=\sum_i\widehat{\mathbf A}_i$ then obeys
\[
K^2(b)=\langle b|\widehat{\mathbf C}^{\,2}|b\rangle
=\sum_i j(j+1)+2\sum_{i<k}m_im_k\,\hat n_i\!\cdot\!\hat n_k
=4j(j+1)-\frac13\bigl(S^2-\Sigma^2\bigr),
\]
with $4j(j+1)=48$ for $j=3$. The decomposition
\[
K^2=|\mathbf c(b)|^2+\sum_i\bigl[j(j+1)-m_i^2\bigr]
\]
exhibits $K^2$ as the mean nonclosure plus the irreducible quantum variance of the four face operators. The expectation is taken in the unprojected product state. This implements the soft-closure convention: the evidence weight $e^{-\eta K^2}$ imposes closure statistically, whereas projection onto the gauge-invariant intertwiner subspace would give different expectation values. The seven-label alphabet and the constant $48$ are fixed together by the $j_{\rm eff}=3$ product-state boundary realization and the tetrahedral normal frame.

The evidence factor follows from the same operator. The closure operator $\widehat{\mathbf C}$ has three components, and the normalized isotropic quadratic evidence kernel on that three-dimensional defect space carries the determinant weight $\eta^{3/2}$, so the normalized closure evidence is $\eta^{3/2}Z(\eta)$ and its logarithm is
\[
\mathcal F(\eta)=\ln Z(\eta)+\tfrac32\ln\eta.
\]
With $\partial_\eta\ln Z=-\langle K^2\rangle_\eta$, the stationary condition $\mathcal F'(\eta)=0$ is the closure relation
\[
\langle K^2\rangle_\eta=\frac{3}{2\eta},
\]
so the factor $3/2$ is the determinant weight of the three independent closure components, not an equipartition rule imported onto the bounded discrete spectrum; the discreteness, multiplicities, and positive floor of the spectrum remain inside the exact finite sum $Z(\eta)$. The second derivative is $\mathcal F''(\eta)=\mathrm{Var}_\eta(K^2)-3/(2\eta^2)$. Global uniqueness follows directly from the spectral bounds. Set $x_{\min}=122/3$ and $x_{\max}=170/3$. Every stationary point lies in
\[
\mathcal I=\left[\frac{3}{2x_{\max}},\frac{3}{2x_{\min}}\right],
\]
because $\mathcal F'>0$ below this interval and $\mathcal F'<0$ above it. The variance bound for a variable in $[x_{\min},x_{\max}]$ gives
\[
\mathrm{Var}_\eta(K^2)\le\frac{(x_{\max}-x_{\min})^2}{4}=64.
\]
Throughout $\mathcal I$,
\[
\mathcal F''(\eta)\le64-\frac23x_{\min}^2
=-1038.5185\ldots<0.
\]
The derivative therefore crosses zero once. Since $\mathcal F\to-\infty$ as $\eta\to0^+$ and as $\eta\to\infty$, this point is the unique global maximum. Evaluation gives
\[
\eta_*=0.0298668443935,
\]
at which $\mathrm{Var}_{\eta_*}(K^2)=15.69$ is dwarfed by $3/(2\eta_*^2)=1681.6$, giving $\mathcal F''(\eta_*)=-1665.9<0$: $\eta_*$ is the unique global maximum of the normalized closure evidence.
Because the parity-symmetric ensemble is finite, both $Z(\eta)$ and $\langle K^2\rangle_\eta$ are exact finite sums over the spectrum. The distinct closure-defect values and their degeneracies are
\[
\begin{array}{c|cccccc}
K^2 &
\frac{122}{3} & \frac{134}{3} & \frac{142}{3} & \frac{146}{3} & \frac{152}{3} & \frac{154}{3} \\ \hline
\mathrm{mult} &
96 & 96 & 96 & 288 & 192 & 144
\end{array}
\]
\[
\begin{array}{c|ccccc}
K^2 &
\frac{158}{3} & 54 & \frac{164}{3} & \frac{166}{3} & \frac{170}{3} \\ \hline
\mathrm{mult} &
384 & 192 & 48 & 96 & 48
\end{array}
\]
with total multiplicity $1680$ as required. In particular,
\[
Z(\eta)=\sum_a n_a e^{-\eta K_a^2},
\qquad
\langle K^2\rangle_\eta=\frac{\sum_a n_a K_a^2 e^{-\eta K_a^2}}{\sum_a n_a e^{-\eta K_a^2}},
\]
where $(K_a^2,n_a)$ run over the table above. The closed-branch value $\eta_*$ is therefore the unique interior maximum of an exact finite-spectrum functional.
Writing $M(\eta)=\langle K^2\rangle_\eta$, the same finite sum gives $M'=-\operatorname{Var}_\eta(K^2)$. On the stationary window,
\[
\frac9{340}\le\eta\le\frac9{244},\qquad
\frac{d}{d\eta}(\eta M)
\ge\frac{122}{3}-\frac9{244}\,64
=\frac{7010}{183}>0.
\]
The finite weighted-variance and rational slope bounds admit exact algebraic verification. Rational exponential bounds for the tabulated spectrum bracket the unique root by
\[
0.0298668443934<\eta_*<0.0298668443936.
\]
The Gibbs derivative, existence, uniqueness, global maximality, and decimal bracket are kernel checked in \path{Entropic/ClosureRoot.lean}. The same proof gives $7.4198000220<g_{\rm share,eff}<7.4198000250$ using rational exponential and logarithm enclosures. The alphabet and normalized evidence functional are the specified model inputs.

The corresponding effective sharing entropy is
\[
g_{\mathrm{share,eff}}
=
-\sum_{b\in B}p_{\eta_*}(b)\ln p_{\eta_*}(b)
=
7.41980002357.
\]
The closed-branch moments used in the UV stiffness discussion are
\[
\begin{aligned}
\langle K^2\rangle_{\eta_*}&=50.2229154254,\\
\mathrm{Var}_{\eta_*}(K^2)&=15.6889750078,\\
a_{\mathrm{UV}}\equiv\frac{1}{\mathrm{Var}_{\eta_*}(K^2)}&=0.0637390269.
\end{aligned}
\]
These values quantify the local stiffness of the canonical closure point.
The closure-saturation product is
\[
C_{\mathrm{cl}}:=\eta_*\langle K^2\rangle_{\eta_*}=\frac32,
\qquad
C_{\mathrm{cl}}^{-1}=\frac23.
\]
The value $C_{\mathrm{cl}}=3/2$ is an identity imposed by the stationarity condition, not an independent numerical cross-check. Its reciprocal is reused as the transverse export factor in the decorated scale-setting branch discussed in Appendix~D.4.

The admissibility parameter stops being free here. The kernel introduces $\eta$, and the closure condition removes its arbitrariness again by demanding that the fluctuation scale produced by the weighting agree with the weighting itself.

\subsection*{B.3 Rooted reduction and local benchmarks}

Rooting on the shared face reduces the exact parity-symmetric ensemble to $140$ rooted microstates and $69$ rooted closure classes. The rooted classes can be labeled by $\alpha=(m_\bullet,K^2)$, so the same reduced state space supports the local evaluation, the cavity benchmark, and the later shell propagation. Let $X$ denote the root-face label and $Y_r$ the boundary record after $r$ rooted shells. The local information observable
\[
\sigma_{\mathrm{ind}}^{(r)}=\frac{H(X\mid Y_r)}{H(X)}
\]
has the principal pre-nonlocal benchmarks
\begin{align}
\sigma_{\mathrm{ind}}^{\mathrm{toy}} &= 0.44997,\\
\sigma_{\mathrm{ind}}^{\mathrm{loc}} &= 0.44708,\\
\sigma_{\mathrm{ind}}^{\mathrm{Bethe}}(J=0) &= 0.44749.
\end{align}
Here the Bethe value is the homogeneous cavity evaluation on the $69\times 69$ rooted-class interaction graph at zero transport coupling,
\[
\mu_\alpha
\propto
w_\alpha
\left(\sum_\beta U_{\alpha\beta}(0)\mu_\beta\right)^{z-1},
\qquad
\sum_\alpha \mu_\alpha=1,
\]
where $w_\alpha=n_\alpha e^{-\eta_*K_\alpha^2}$ is the rooted-class Gibbs weight, $z=4$, and $U_{\alpha\beta}(0)$ is the rooted shared-face compatibility matrix before shell transport is turned on. Thus $\sigma_{\mathrm{ind}}^{\mathrm{Bethe}}(J=0)$ is the cavity-theory benchmark of the same explicit rooted ensemble. The three benchmark values are single-instance calculations and remain provisional until they are independently rederived. When the shared seven-valued label $L$ is the sole channel, so that $X-L-Y$ is a Markov chain, data processing gives $I(X;Y)\le H(L)\le\ln7$. The finite full-support case is kernel checked. For the seven-valued root label $X$, this bound permits $H(X\mid Y)/H(X)<1/2$. The sixteen-dimensional primitive pair space follows from the fusion construction of Appendix~H.7.
The horizon target implied by the effective sharing entropy is
\[
\sigma_*=\frac{\pi}{g_{\mathrm{share,eff}}}=0.42340665.
\]
The gap between the local benchmarks and $\sigma_*$ arises from shell and loop structure. The local admissibility closure already satisfies its own condition.

The remaining ultraviolet calculation propagates the fixed local closure ensemble through transport and return structure.

\subsection*{B.4 Channel ontology and derived selection}

Four statements require separate tests.
Positive oriented matching selects one transmitted multiplet $V_{2j_0}$. The
choice of that matching rule is part of the primitive ultraviolet architecture
and has a spectral falsifier. In the selected sharp reversible vertex,
unitarity fixes the retained carrier to the fusion complement; completeness
of the H.9 carrier then fixes $j_0=\tfrac32$. If a different microscopic
vertex is chosen, $j_0=\tfrac32$ remains the surviving member of the discrete
audit in B.5.

\subsubsection*{B.4.1 Setup}

Two cells share a face. Each contributes a primitive spin-$j_0$ slot, so the two-sided space is
\[
\mathcal H_f
=
V_{j_0}\otimes V_{j_0}
=
\bigoplus_{J=0}^{2j_0}V_J.
\]
The decomposition is multiplicity-free. The face exports a classical weight alphabet to the admissibility ensemble. The microscopic questions are which coupled sector is transmitted, how the untransmitted information is retained, and why the four port assignments are injective.

\subsubsection*{B.4.2 Positive oriented matching}

\paragraph{Scope.}
Oriented gluing alone does not select $j_0$: every half-integer $j_0$ defines
a discrete member with $4j_0+1$ transmitted weights. Sections~B.4.2e--f add
the selected sharp reversible split. Its unitarity identifies the complete
marked carrier with the nonmaximal fusion output, after which the independently
fixed spin-one closure response uniquely gives $j_0=\tfrac32$. No
smallest-spin rule is used.

\subsubsection*{B.4.2a Face-sector Hessian}

\paragraph{Proposition (exact block decomposition).}
If the gauge-fixed face Hessian commutes with the diagonal $SU(2)$ action, Schur's lemma gives
\[
\boxed{
\mathcal H^{(2)}(p)
=
\bigoplus_{J=0}^{2j_0}
\left(Z_Jp^2+r_J\right)I_{V_J}.
}
\]
Thus an $SU(2)$-covariant quadratic kernel cannot mix different total-spin sectors. Its dynamical content is the finite spectrum $\{Z_J,r_J\}$.

For invariant block densities
\[
\rho_J=\sum_{M=-J}^{J}|\sigma_{JM}|^2,
\]
write the radial part of the homogeneous potential as
\[
V_{\rm eff}
=
\sum_Jr_J\rho_J
+\frac12\sum_Ju_J\rho_J^2
+\sum_{J<K}w_{JK}\rho_J\rho_K.
\]
The first instability lies in the block with the smallest $r_J$. A pure $J_*$ saddle, with $\rho_{J_*}^{(0)}=-r_{J_*}/u_{J_*}$, is locally stable against an unoccupied block $J$ when
\[
m_{J|J_*}^2
=
r_J-\frac{w_{JJ_*}}{u_{J_*}}r_{J_*}>0.
\]
These are the general tests. The minimal matching action below evaluates them explicitly.

\subsubsection*{B.4.2b Exact coherent-matching kernel}

The same geometric face has opposite outward orientations in its two incident cells. Let $g_{LR}$ transport the right frame into the left frame and define the right generator in the common orientation by
\[
\widetilde{\mathbf J}_R
=
\operatorname{Ad}_{g_{LR}}\mathbf J_R^{\vee},
\qquad J_R^{\vee a}=-(J_R^a)^T.
\]
The transpose implements the dual representation associated with reversing the outward face orientation. The generators $J_R^{\vee a}$ obey the usual $\mathfrak{su}(2)$ algebra and are unitarily equivalent to spin $j$. Perfect matching compares equal oriented data. If
\[
|j,\mathbf n\rangle
=
D^{(j)}(g_{\mathbf n})|j,j\rangle
\]
is a spin coherent state, the matched pair is
\[
|j,\mathbf n\rangle_L\otimes
|j,\mathbf n\rangle_{\widetilde R}.
\]
The highest-weight identity
\[
|j,j\rangle\otimes|j,j\rangle=|2j,2j\rangle
\]
and diagonal covariance imply
\[
|j,\mathbf n\rangle\otimes|j,\mathbf n\rangle
=
F_{2j}^{\dagger}|2j,\mathbf n\rangle,
\]
where $F_{2j}:V_j\otimes V_j\to V_{2j}$ is the maximal-channel Clebsch--Gordan coisometry.

The direction-independent positive Gram operator of exact matching is
\[
\mathcal G_{\rm sharp}
=
(4j+1)
\int_{S^2}\frac{d\Omega_{\mathbf n}}{4\pi}\,
|j,\mathbf n;j,\mathbf n\rangle
\langle j,\mathbf n;j,\mathbf n|.
\]
The coherent-state resolution of the identity on $V_{2j}$ gives
\[
\boxed{
\mathcal G_{\rm sharp}
=
F_{2j}^{\dagger}F_{2j}
=
P_{2j}.
}
\]
Thus exact oriented matching transmits the maximal coupled multiplet. At $j_0=\tfrac32$,
\[
\mathcal G_{\rm sharp}=P_3,
\qquad
\dim V_3=7.
\]
The seven-state link is the image of the pair operator on the full sixteen-state space, so the alphabet follows from the tensor product itself.

\paragraph{The contraction alternative.}
Ordinary $\varepsilon$-contraction is a different microscopic operation. It pairs dual legs and leaves an unfused alphabet of $2j_0+1$ states rather than transmitting $V_{2j_0}$. At $j_0=\tfrac32$ this gives four letters and the 48-state branch in B.5. The two operations therefore define distinct gluing ontologies with distinct closure spectra. The present theory adopts oriented fusion because the face must continue to carry closure data after gluing.

\subsubsection*{B.4.2c Finite-width completion and its falsifier}

The quantum mismatch of two perfectly aligned spin-$j_0$ slots has the irreducible floor $2j_0$. Subtracting that floor defines
\[
Q_f
=
\frac12\left[
\left(\mathbf J_L-\widetilde{\mathbf J}_R\right)^2
-2j_0I
\right].
\]
On $V_J$,
\[
Q_f\big|_{V_J}
=
\frac12\left[
(2j_0)(2j_0+1)-J(J+1)
\right]I_{V_J}.
\]
Hence $Q_f\ge0$ and its null space is exactly $V_{2j_0}$. The positive finite-width transfer
\[
T_s=e^{-sQ_f},
\qquad s>0,
\]
forms a semigroup and converges to $P_{2j_0}$ as $s\to\infty$. For $j_0=\tfrac32$,
\[
\boxed{
Q_f=6P_0+5P_1+3P_2,
\qquad
T_s=e^{-6s}P_0+e^{-5s}P_1+e^{-3s}P_2+P_3.
}
\]

A minimal homogeneous action on the complete primitive space is
\[
\Gamma_f
=
\int d^dx\left[
Z|\partial\Phi|^2
+r\,\Phi^\dagger\Phi
+\kappa\,\Phi^\dagger Q_f\Phi
+\frac{u}{2}(\Phi^\dagger\Phi)^2
\right],
\qquad
Z,\kappa,u>0.
\]
Its potential is
\[
V
=
r\rho+\frac{u}{2}\rho^2
+6\kappa\rho_0
+5\kappa\rho_1
+3\kappa\rho_2,
\qquad
\rho=\sum_{J=0}^{3}\rho_J.
\]
For $r<0$, the fusion-sector support of the homogeneous minimum is uniquely $J=3$:
\[
\rho_3=-\frac{r}{u},
\qquad
\rho_0=\rho_1=\rho_2=0.
\]
The orthogonal blocks have positive masses
\[
\boxed{
m_{0|3}^2=6\kappa,
\qquad
m_{1|3}^2=5\kappa,
\qquad
m_{2|3}^2=3\kappa.
}
\]
At fixed total density, every nonmaximal occupation raises the action, so mixed-$J$ support is excluded globally in this minimal model.

\paragraph{Theorem (maximal-fusion blocking ray).}
Let $N$ transmitted cells each carry the selected irrep $V_3$. The highest-spin irrep $V_{3N}$ occurs with multiplicity one in $V_3^{\otimes N}$. Therefore the iterated maximal-channel Clebsch--Gordan coisometries define, after one consistent phase convention, a unique map
\[
\boxed{
F_N:V_3^{\otimes N}\longrightarrow V_{3N}
}
\]
independent of the binary fusion tree. Associativity here is not an additional dynamical assumption: every binary tree that always chooses the maximal coupled irrep ends in the same multiplicity-one copy of $V_{3N}$, so two such intertwiners differ only by an overall phase.

For spin coherent states the map is exact,
\[
\boxed{
F_N\,|3,\mathbf n\rangle^{\otimes N}=|3N,\mathbf n\rangle .
}
\]
Thus a fixed microscopic $j=3$ alphabet is compatible with an arbitrarily large semiclassical block spin on the selected maximal-channel branch; the elementary representation does not run under this canonical blocking map.

\paragraph{Lemma (equivariant lock under maximal blocking).}
Let $L_3:V_3^{\rm cap}\to V_3^{\rm geom}$ be the normalized one-cell equivariant lock of Appendix~Q, and let $F_N^{\rm cap}$ and $F_N^{\rm geom}$ denote normalized maximal-channel fusion maps on the two factors.  Because $V_{3N}$ occurs with multiplicity one in $V_3^{\otimes N}$, the two normalized $SU(2)$-equivariant maps from $(V_3^{\rm cap})^{\otimes N}$ to $V_{3N}^{\rm geom}$ can differ only by a phase.  Hence there is $e^{i\phi_N}$ such that
\[
\boxed{
F_N^{\rm geom}L_3^{\otimes N}
=e^{i\phi_N}L_{3N}F_N^{\rm cap}.
}
\]
Thus, conditional on the selected one-cell lock, its propagation to every maximal block is unique and parameter free.  This is a representation-theoretic theorem; it does not by itself prove that the full interacting coarse-graining map dynamically projects arbitrary blocks onto the maximal channel.

A dynamical attraction toward the same ray is exhibited within the displayed positive all-pairs diagnostic model
\[
H_N
=\kappa\sum_{a<b}\left(9-\mathbf J_a\!\cdot\!\mathbf J_b\right)
=\frac{\kappa}{2}\left[(3N)(3N+1)-\mathbf J_{\rm tot}^2\right].
\]
It obeys
\[
\ker H_N=V_{3N},
\qquad
\Delta_N\equiv E_{3N-1}-E_{3N}=3\kappa N.
\]
The increasing gap is a diagnostic of maximal-spin phase selection; the exact fusion map above does not require an all-to-all microscopic interaction.

The same conclusion is visible directly in the face mismatch. For equal spin $j$,
\[
Q_j=j^2-\mathbf J_L\!\cdot\!\widetilde{\mathbf J}_R,
\]
and on coherent data separated by angle $\theta$,
\[
\boxed{
\langle Q_j\rangle
=j^2(1-\cos\theta)
=\frac12j^2\theta^2+O(\theta^4).
}
\]
Hence $e^{-sQ_j}$ has angular width $\Delta\theta\sim(j\sqrt{s})^{-1}$. On a maximal block $j=3N$ the accepted geometric matching therefore sharpens rather than broadens as $N$ grows.

\paragraph{Heisenberg form of the face mismatch.}
With $\mathbf J_L^2=\widetilde{\mathbf J}_R^{\,2}=j(j+1)$, the mismatch operator is the isotropic ferromagnetic Heisenberg bond shifted by a constant,
\[
Q_j=j^2-\mathbf J_L\!\cdot\!\widetilde{\mathbf J}_R,
\qquad
\mathbf J_L\!\cdot\!\widetilde{\mathbf J}_R\big|_{V_J}=\tfrac12\left[J(J+1)-2j(j+1)\right].
\]
For the primitive pair $j_0=\tfrac32$ this gives $q_J=6-\tfrac12J(J+1)$, so $J=0,1,2,3$ cost $6,5,3,0$ and the spectrum $Q_f=6P_0+5P_1+3P_2$ of B.4.2c follows with its zero on $V_3$. After fusion each block carries spin three, and the nearest-neighbor bond has the spectrum
\[
\begin{array}{c|ccccccc}
J_{ab}&0&1&2&3&4&5&6\\ \hline
9-\mathbf J_a\!\cdot\!\mathbf J_b&21&20&18&15&11&6&0
\end{array}
\]
Each positive bond selects $J_{ab}=6$, and on a connected graph the kernel of the bond sum is the maximal ray $V_{3N}$. The selected matching operator therefore belongs to the class of the nearest-neighbor isotropic quantum Heisenberg ferromagnet. For that class on $\mathbb Z^d$ with $d\ge3$ and every spin $S\in\{\tfrac12,1,\tfrac32,\ldots\}$, a translation-invariant equilibrium state with magnetization at least $S/4$ exists at every sufficiently low positive temperature and satisfies the KMS condition for the zero-field dynamics \citep{OpenAIHeisenberg2026}. In three dimensions the magnetization deficit obeys Bloch's $T^{3/2}$ law, with the coefficient fixed by the one-magnon dispersion, for every nonnegative symmetric finite-range coupling whose support generates $\mathbb Z^3$ \citep{OpenAIBloch2026}. Both $S=\tfrac32$ and $S=3$ lie in this class. On hypercubic hosts the all-pairs diagnostic above therefore has a local counterpart: the maximal ray is the zero-temperature kernel of the nearest-neighbor bond sum, and the aligned phase persists at positive temperature as a property of the ordered state.

The ordered local phase carries gapless spin waves, and the uniformly gapped complement used below comes from the hierarchical fusion parent of Appendix~H.10, paragraph~L. The theorem also concerns the global $SU(2)$ acting on unprojected face vectors. Exact local Gauss projection removes the single-face bond between physical cells (H.10, paragraph~L), so the theorem supports the gauge-fixed coherent description of the maximal branch, and ordering of the gauge-invariant shape variables is left to multi-face, holonomy, or vertex terms.

\paragraph{Corollary (additive channel area and the coarse area operator).}
The selected factorized completion of Appendix~Q locks each transmitted capacity multiplet $V_3^{\rm cap}$ to a geometric $V_3^{\rm geom}$ in the same relational orientation.  On the coherent branch the geometric face variable is the flux vector itself: a spin-$J$ coherent face carries the classical/Regge area proportional to $J$, while the kinematical LQG Casimir operator has eigenvalue proportional to $\sqrt{J(J+1)}$ \citep{PerezSpinFoam2013,HanSmallSpin2019}.  For $N$ maximally fused primitive faces,
\[
J_N=3N.
\]
Area additivity of the $N$ fine facets therefore fixes the physical coarse area observable uniquely.  If
\[
\widehat A_{\rm LQG}(J)
=8\pi\gamma_{\rm BI}L_*^2\sqrt{J(J+1)},
\]
then the renormalized operator on the maximal-fusion ray is
\[
\boxed{
\widehat A_{\rm phys}^{(N)}
=Z_A(N)\widehat A_{\rm LQG}(3N),
\qquad
Z_A(N)=\frac{3N}{\sqrt{3N(3N+1)}}
=\sqrt{\frac{3N}{3N+1}}.
}
\]
Hence
\[
\widehat A_{\rm phys}^{(N)}
=24\pi\gamma_{\rm BI}N L_*^2,
\qquad
Z_A(1)=\frac{\sqrt3}{2},
\qquad
Z_A(N)\to1.
\]
The microscopic spin remains fixed. The factor $Z_A(N)$ renormalizes the coarse observable so that one irrep $V_{3N}$ represents the additive area of $N$ coherent $V_3$ facets. It reconciles the exact small-$J$ Casimir spectrum with the linear area of the Regge phase, and approaches unity in the large-block limit. Appendix~F.5 fixes the remaining dimensionless normalization $\gamma_{\rm BI}$ from the one-bit cut entropy.

\paragraph{Controlled semiclassical-refinement witness.}
The exact blocking ray can be embedded in the established spin-foam semiclassical--continuum limit of Refs.~\citep{Han2017,HanHuangZipfel2019}. To avoid collision with the couplings elsewhere in this paper, denote their typical large-spin parameter by $\Lambda_{\rm sf}$. Their allowed family includes
\[
\Lambda_{\rm sf}(\mu)=\Lambda_0\mu^{-2+u},
\qquad 0<u<\frac25,
\qquad \mu\to0.
\]
Choose $u=1/4$ and identify the spin-foam large-spin scale only up to the fixed normalization carried by the blocked representation,
\[
N(\mu)=\left\lfloor N_0\mu^{-7/4}\right\rceil,
\qquad
J_{\rm block}=3N,
\qquad
\Lambda_{\rm sf}\propto J_{\rm block}=3N.
\]
The constant factor is asymptotically irrelevant and is absorbed into $\Lambda_0$; no equality between the symbol $\Lambda_{\rm sf}$ and the cell count is required.
One admissible regulator choice is
\[
\delta(\mu)=\delta_0\mu^{3/4},
\]
which lies inside the semiclassical--continuum window for sufficiently small $\mu$. The physically relevant refinement variable is the Regge edge length measured in the running infrared unit $\mu^{-1}$, not an identification of one capacity cell with a fixed geometric edge. With $\ell_{\rm block}\sim\sqrt{\gamma\Lambda_{\rm sf}}\,L_P$ and $L_P(G_*)=L_*$ on the matched normalization,
\[
\ell_{\rm rel}(\mu)
\equiv \mu\ell_{\rm block}
\propto\mu^{1/8}\longrightarrow0
\]
(up to the fixed dimensionless spin-foam normalization). This use of $L_*$ fixes the gravitational unit; it does not assign a literal edge length $L_*$ to the microscopic capacity cell.
Along this explicit witness trajectory,
\[
\boxed{
\Lambda_{\rm sf}^{-1}=O(\mu^{7/4}),\qquad
\Delta\theta=O(\mu^{7/4}),\qquad
\delta=O(\mu^{3/4}),\qquad
\ell_{\rm rel}^2=O(\mu^{1/4}).
}
\]
Along this trajectory, the finite-spin correction, capacity-matching width,
transverse spin-sum regulator, and leading lattice-curvature correction all
vanish while the microscopic constituent remains exactly $j=3$. In the
controlled linearized sector, Ref.~\citep{HanHuangZipfel2019} proves that the
low-energy excitations exhaust the smooth solutions of the linearized
Einstein equations: the two helicity-$2$ graviton modes. Ref.~\citep{Han2017}
gives the broader Lorentzian/Euclidean Einstein--Regge
semiclassical-continuum construction. No general nonlinear convergence
theorem is claimed here.

The microscopic test is finite. For the reduced physical Hessian define
\[
h_J
=
\frac{1}{2J+1}
\operatorname{Tr}_{V_J}
\left(P_J\mathcal H_fP_J\right).
\]
The weak pass condition is
\[
h_3<h_2,\qquad h_3<h_1,\qquad h_3<h_0.
\]
The minimal-kernel target is
\[
\boxed{
(h_0-h_3):(h_1-h_3):(h_2-h_3)=6:5:3.
}
\]
After condensation the remaining tests are $h_J+w_{J3}\rho_3>0$ for $J=0,1,2$. The construction fails if another block is lighter, inequivalent $J$ sectors mix after gauge reduction, a stable mixed-$J$ minimum survives, or the primitive transfer is not positive.

\subsubsection*{B.4.2d Primitive tensor and effective pushforward}

For a shared face $(ab)$, let $Q_{ab}$ denote the transported mismatch operator above. A concrete finite-width primitive tensor is
\[
\mathcal G_{v,s}^{(0)}
=
\mathcal P_v^{\rm inv}
\left[
\bigotimes_{(ab)\in E(v)}e^{-sQ_{ab}}
\right]
\mathcal P_v^{\rm inv},
\]
where $\mathcal P_v^{\rm inv}$ is the cellwise gauge/intertwiner projector. The zero-parameter primitive theory is defined by its sharp form,
\[
\boxed{
\mathcal G_{v,\sharp}^{(0)}
=
\mathcal P_v^{\rm inv}
\left[
\bigotimes_{(ab)\in E(v)}P_3^{(ab)}
\right]
\mathcal P_v^{\rm inv}.
}
\]
The finite-$s$ family is a positive regulator and a spectral falsifier, not an additional coupling of the sharp theory.

Let
\[
F_3:V_{3/2}\otimes V_{3/2}\longrightarrow V_3
\]
obey $F_3^\dagger F_3=P_3$ and $F_3F_3^\dagger=I_{V_3}$. On an isolated face,
\[
\boxed{
F_3T_sF_3^\dagger=I_{V_3}
}
\]
for every $s>0$. For the complete gauge-projected vertex, finite-$s$ independence additionally requires $\mathcal P_v^{\rm inv}$ to preserve the facewise coupled-$J$ decomposition, or equivalently that the spectator/intertwiner dressing be channel-blind. This condition requires a reduced-kernel or commutator audit. The sharp $P_3$ support survives whenever its overlap is nonzero; finite-width spectral ordering remains an explicit test.

With $\mathcal F_v=\bigotimes F_3$, the post-fusion geometric tensor is
\[
\mathcal G_v^{\rm eff}
=
\mathcal F_v\mathcal G_{v,\sharp}^{(0)}\mathcal F_v^\dagger.
\]
The marked-transfer vertex of Appendix~H decorates this effective object. It is not itself the primitive pair kernel.

\subsubsection*{B.4.2e Exact marked complement and conditional sharp-kernel selection}

Let a marked response be built from two spin-$s$ strands. Then
\[
\mathcal H_{\rm mark}(s)
=
V_s^P\otimes V_s^H
=
\bigoplus_{J=0}^{2s}V_J.
\]
Maximal fusion of two primitive spin-$j$ slots transmits $V_{2j}$ and leaves
\[
Q(j)
=
\left(V_j\otimes V_j\right)\ominus V_{2j}
=
\bigoplus_{J=0}^{2j-1}V_J.
\]

\paragraph{Theorem (unique complement match).}
The representations are isomorphic if and only if
\[
\boxed{
\mathcal H_{\rm mark}(s)\cong Q(j)
\quad\Longleftrightarrow\quad
j=s+\frac12.
}
\]

\paragraph{Proof.}
Both decompositions are multiplicity-free. The marked space contains one copy of each $V_J$ from $J=0$ through $2s$; the fusion complement contains one copy from $J=0$ through $2j-1$. They agree exactly when $2s=2j-1$. Equivalently,
\[
\dim Q(j)
=(2j+1)^2-(4j+1)
=4j^2,
\qquad
\dim\mathcal H_{\rm mark}(s)=(2s+1)^2,
\]
and the positive solution is $2j=2s+1$.
\hfill$\square$

The marked closure response is independently a three-component vector and
therefore has $s=1$. In the selected sharp reversible split, the complete
retained output is the fusion complement and the complete H.9 one-mark
carrier has the same nine dimensions. Hence the carrier-space identification
is isometric and
\[
\boxed{j_0=\frac32.}
\]

\paragraph{Explicit complement isometry.}
On $V_{3/2}\otimes V_{3/2}$ let
\[
P_Q=P_0+P_1+P_2.
\]
The blockwise equivariant isometry
\[
\mathcal I:Q\longrightarrow V_1^P\otimes V_1^H
\]
is defined by
\[
\mathcal I\,|(j_0j_0);J,M\rangle
=
|(1_P1_H);J,M\rangle,
\qquad J=0,1,2.
\]
The independent present and history strands have ordered first-excitation basis
\[
|a,b\rangle=|a\rangle_P\otimes|b\rangle_H,
\qquad
\langle a,b|c,d\rangle=\delta_{ac}\delta_{bd}.
\]
Let $W$ map this basis to the nine hard-core marked states,
\[
W|a,b\rangle=d_{ab}^{\dagger}|0\rangle.
\]
The exact unit Gram matrix gives
\[
W^\dagger W=I_9,
\qquad
WW^\dagger=I_{\mathcal H_d}.
\]
Define
\[
L=W\mathcal I P_Q.
\]
Then
\[
\boxed{
L^\dagger L=P_Q,
\qquad
LL^\dagger=I_{\mathcal H_d}.
}
\]
The complete sharp pair map is
\[
U_{\rm sharp}|\psi\rangle
=
F_3P_3|\psi\rangle\otimes|0\rangle_R
+|\varnothing\rangle\otimes L|\psi\rangle.
\]
The two flags are orthogonal, so
\[
\boxed{
U_{\rm sharp}^{\dagger}U_{\rm sharp}
=P_3+P_Q
=I_{16}.
}
\]
The seven-state link and nine-state marked record are therefore complementary outputs of an explicit reversible pair map.
More generally, writing the retained output as $C|\psi\rangle$, unitarity of
\[
U|\psi\rangle
=F_3P_3|\psi\rangle|0\rangle
+|\varnothing\rangle C|\psi\rangle
\]
gives
\[
P_3+C^\dagger C=I,
\qquad
\boxed{C^\dagger C=P_Q.}
\]
Because the complete H.9 carrier and $Q$ both have dimension nine, $C$ is an
isometric identification onto that carrier. This proves the carrier-space
incidence in the selected sharp vertex. It does not derive the hard-core
operator realization, the flags, or the detailed routing and feedback graph.

Since $V_0$, $V_1$, and $V_2$ occur once, Schur's lemma implies that every other equivariant minimal complement isometry differs by one phase on each block,
\[
L'
=
e^{i\phi_0}L_0
\oplus
e^{i\phi_1}L_1
\oplus
e^{i\phi_2}L_2.
\]
If no microscopic operator compares the blocks coherently, these phases are marked-field conventions. A vertex that does compare them must derive their relative phases.

\paragraph{Why the unit marked weights select the sharp limit.}
For finite $s$, the system Kraus amplitude in the canonical dilation is $T_s^{1/2}$ and the complementary amplitude is $\sqrt{I-T_s}$. Its squared norms on the three nonmaximal blocks are
\[
\ell_0^2(s)=1-e^{-6s},
\qquad
\ell_1^2(s)=1-e^{-5s},
\qquad
\ell_2^2(s)=1-e^{-3s}.
\]
The H.9 marked vertex instead has one common quadratic normalization,
\[
G_{\rm mark}=I_9=P_0+P_1+P_2.
\]
Therefore, within the minimal incidence identification in which this marked vertex is the canonical fusion environment and no compensating $J$-dependent coupling is added, exact equality of the marked weights requires
\[
\boxed{
s\to\infty,
\qquad
T_s\to P_3,
\qquad
I-T_s\to P_Q.
}
\]
A finite physical width could survive only with additional block-dependent
normalization. The sharp conclusion holds within the selected reversible
carrier incidence. For a different microscopic vertex, the equal H.9 Gram
matrix and the finite-width fusion transfer describe distinct objects and
impose no constraint on one another.

\subsubsection*{B.4.2f Conditional dimensional endpoint}

The closure defect is a spatial vector,
\[
C_a,\qquad a=1,2,3.
\]
In three spatial dimensions the vector representation of $SO(3)$ lifts to the spin-one irrep of $SU(2)$:
\[
d=3
\quad\Longrightarrow\quad
\mathbb R^3\cong V_1
\quad\Longrightarrow\quad
s=1.
\]
Within the selected sharp carrier incidence, the theorem above gives
\[
s=1
\Longrightarrow
j_0=\frac32
\Longrightarrow
V_{2j_0}=V_3
\Longrightarrow
|M|=7.
\]
The minimal three-dimensional simplex has four faces, and the present/history vector record has $3^2=9$ states. Single-copy channel capacity makes the four port assignments injective, while the two global orientations give
\[
\boxed{
\Omega=2P(7,4)=1680.
}
\]
The complete conditional integer chain is therefore
\[
\boxed{
d=3
\Longrightarrow
s=1
\Longrightarrow
j_0=\frac32
\Longrightarrow
\begin{cases}
\text{ports}=4,\\
\text{alphabet}=7,\\
\text{record}=9,\\
\Omega=1680.
\end{cases}
}
\]
For the physical $d=3$ rotation algebra one may summarize the evaluated identities as
\[
\text{ports}=d+1,
\qquad
|M|=2d+1,
\qquad
\dim\mathcal H_{\rm mark}=d^2.
\]
They are not a claimed continuation to arbitrary dimension. For $d\neq3$, the vector representation belongs to $\mathrm{Spin}(d)$ and need not be one $SU(2)$ spin. A dimension-general theorem would require repeating the complement analysis in that group.

\subsubsection*{B.4.3 Injectivity from single-copy channel capacity}

The exclusion statement concerns a single cell resource, not Pauli antisymmetry among the four port labels.

\paragraph{Lemma.}
Assume:
\begin{enumerate}
\item the letters index a single copy of the cell-level fermionic channel family, with occupations $n_m\in\{0,1\}$;
\item a face displaying $m$ is, by definition, that channel occupation routed through the face's port;
\item the four ports are distinguishable relational registers.
\end{enumerate}
Then the admissible boundary states are the injections $f:F_4\to M$, and their number before orientation doubling is $P(|M|,4)$.

\paragraph{Proof.}
If two ports displayed the same letter $m$, the one cell resource would require $n_m=2$, contrary to its single-copy ceiling. The assignment is therefore injective. The ports are distinct interfaces to distinct neighbors, so permuting letters between ports changes the routing state; the assignments are ordered. Antisymmetry acts in the occupation algebra of each channel, not by quotienting the relational port index.
\hfill$\square$

This lemma does not follow from generic fermionic statistics alone. Its physical content is that the four ports draw from one shared channel multiplet rather than four independent copies. That is the same single-copy channel ontology used by the one-bit defect and determinant sectors.

\subsubsection*{B.4.4 Uniform base measure}

\paragraph{Lemma.}
If the selected face multiplet realizes its full Shannon capacity, its unconditioned base measure is uniform.

\paragraph{Proof.}
For a probability measure $\mu$ on $M$,
\[
H(\mu)\le\ln|M|,
\]
with equality only at $\mu(m)=1/|M|$. Full use of the fixed alphabet therefore selects the uniform prior.
\hfill$\square$

This statement concerns the reference measure, not the final closure-weighted ensemble:
\[
p_\eta(b)
=
\frac{1}{Z(\eta)}
\underbrace{\mu_0(b)}_{\text{uniform base measure}}
e^{-\eta K^2(b)}.
\]
Covariance alone does not force uniformity, because the face normal permits covariant deformations such as functions of $(\mathbf J\cdot\hat n)^2$. Full-capacity saturation excludes those deformations at the prior level; the closure observable then supplies the physical nonuniformity.

\subsubsection*{B.4.5 Status ledger}

\paragraph{Derived from the primitive matching rule.}
For any fixed half-integer $j_0$, exact oriented coherent matching gives the transmitted block $V_{2j_0}$ and an alphabet of $4j_0+1$ states. The single-copy channel ontology gives injectivity; relational ports preserve ordered assignments; and the global face orientation gives the binary factor.

\paragraph{Derived in the selected sharp carrier incidence.}
The exact relation
\[
\mathcal H_{\rm mark}(s)\cong Q(j)
\quad\Longleftrightarrow\quad
j=s+\frac12
\]
combines with the three-dimensional vector response $s=1$ to fix
$j_0=\tfrac32$, seven transmitted states, nine complement states, and
$\Omega=1680$. Unitarity identifies that complement with the complete
nine-state H.9 carrier in the selected sharp reversible split. The exact unit
normalization of the marked vertex selects the sharp $P_3$ kernel within the
minimal canonical dilation.

\paragraph{Microscopic premise.}
The primitive action implements positive oriented matching. The displayed sharp tensor $\mathcal G_{v,\sharp}^{(0)}$ is the minimal completion. A more fundamental simplicial calculation must still reproduce its $J=3$ support or pass the weaker Hessian inequalities of B.4.2c.

\paragraph{Carrier incidence and selected microscopic detail.}
The representation equivalence, $L^\dagger L=P_Q$, $LL^\dagger=I_9$, and
$U_{\rm sharp}^\dagger U_{\rm sharp}=I_{16}$ are exact. They prove the
carrier-space identification in the selected sharp reversible vertex.
The hard-core realization, flags, relative block phases, routing, and
feedback graph remain selected microscopic structure rather than consequences
of the representation theorem alone.

\paragraph{Data status.}
Within the selected sharp vertex, $j_0=\tfrac32$ is internally fixed by the
carrier incidence and the audit below becomes a consistency check. For a
different microscopic vertex, the same value is selected by the postdictive
discrete audit. In neither reading does a continuous parameter choose the
alphabet.

\subsection*{B.5 Relaxation and branch audits}

The following tables propagate nearby discrete constructions through one
fixed downstream map. They are alternatives, not samples from a probability
distribution, so ``sigma'' significances do not apply. The marked sector keeps
nine response polarizations. A branch with alphabet size $n$ uses $\binom n2$
pair records, the scalar factor $2/n$, and the second-return coefficient $n$.
Because the Newton target informed the historical construction, Appendix~L
grades this as a conditional postdictive comparison.

\subsubsection*{B.5.1 Fusion and constituent-spin branches}

\[
\begin{array}{lcccc}
\hline
\text{branch} & \Omega & g_{\rm share,eff} & n\,g_{\rm share,eff} & G_*/G\\
\hline
j_0=\tfrac12\ \text{full reducible space}
&48&3.8712&15.5&7\times10^{31}\\
j_0=\tfrac32\ \text{unfused contraction}
&48&3.8712&15.5&8\times10^{31}\\
j_0=\tfrac32,\ J=2
&240&5.4785&27.4&2\times10^{21}\\
j_0=\tfrac32,\ J=3
&1680&7.4198&51.9&1.0000\\
j_0=\tfrac52,\ J=5
&15840&9.6572&106.2&7\times10^{-48}\\
j_0=\tfrac32\ \text{full reducible space}
&87360&11.2616&180.2&4\times10^{-112}\\
\hline
\end{array}
\]
The two 48-state rows have the same count but different representation content: the first mixes $V_0$ and $V_1$, while the second carries the unfused $V_{3/2}$ alphabet. The family spans about $10^{143}$ in $G_*/G$. Holding the rest of the construction fixed, the seven-state branch is the only row near the measured Newton scale.

\subsubsection*{B.5.2 Counting conventions at seven letters}

\[
\begin{array}{lcccc}
\hline
\text{convention} & \Omega & g_{\rm share,eff} & G_*/G & m_\mu/m_e\\
\hline
2P(7,4) &1680&7.4198&1.0000&206.7683\\
\text{drop orientation doubling} &840&6.7267&1.7\times10^4&207.82\\
\text{allow repeated labels} &4802&8.4483&5.5\times10^{-7}&206.09\\
\text{quotient port permutations} &70&4.2417&3.2\times10^{19}&231.01\\
\text{quotient ports and orientation} &35&3.5486&1.1\times10^{24}&256.31\\
\hline
\end{array}
\]
These rows change the counting convention while retaining the same seven-letter closure formula and downstream transport. The baseline is the only displayed convention that keeps the Newton and charged-lepton comparisons simultaneously near their measured values.

\subsubsection*{B.5.3 One continuous prior deformation}

To test one nearby nonuniform prior, introduce the parity-even, face-axis-preserving tilt
\[
w_\lambda(b)
=
\exp\!\left[-\lambda\sum_{i=1}^{4}m_i^2\right]
\]
and re-solve the closure condition at each $\lambda$. At $\lambda=0$, exact finite differences of the enumerated ensemble give
\[
\frac{dg_{\rm share,eff}}{d\lambda}=-0.278458,
\qquad
\frac{d\ln G_*}{d\lambda}=3.90135,
\qquad
\frac{d\ln(m_\mu/m_e)}{d\lambda}=1.37074\times10^{-3}.
\]
With the present CODATA uncertainty and the baseline residual retained, the one-standard-deviation Newton interval along this one deformation is approximately
\[
-5.3\times10^{-6}<\lambda<6.2\times10^{-6}.
\]
For the muon ratio, $|\Delta\lambda|\simeq1.6\times10^{-5}$ produces one current experimental standard deviation locally. This is a sensitivity scale, not a symmetric bound about zero, because the baseline prediction is not exactly centered. These statements constrain one specified deformation conditional on the frozen downstream map; they do not directly measure the microscopic prior or exhaust all nonuniform measures. Parity-odd tilts have zero linear response by the $m\mapsto-m$ symmetry and require a separate quadratic audit.

\paragraph{Provenance.}
The entropy near $7.42$ was recognized historically by inverting the measured Newton constant, and the decorated vertex was constructed after the remaining Newton and lepton residuals were known. Appendix~L retains that chronology. Once the primitive oriented-matching rule and the marked-fusion incidence identification are adopted, the sequence
\[
d=3
\Longrightarrow
V_1
\Longrightarrow
j_0=\frac32
\Longrightarrow
V_3
\Longrightarrow
7
\Longrightarrow
1680
\Longrightarrow
g_{\rm share,eff}=7.419800\ldots
\]
contains no adjustable continuous parameter. The branch tables show how sharply nearby stated constructions move the downstream observables; they do not turn the historical Newton comparison into a historically blind prediction.

\section*{Appendix C: Edge Kernel, Finite Renormalization, and Continuum Matching}

Appendix~C carries the local boundary ensemble of Appendix~B into edge transport, loop dressing, and the continuum stiffness coefficient of the weak-field EFT.

\subsection*{C.1 Channel-averaged isotropy identity and tree coupling}

Let $\hat n_i$ be the four face normals of a regular tetrahedron. The exact identity
\[
\sum_{i=1}^4 \hat n_i\hat n_i^{\mathsf T}=\frac{4}{3}I_3
\]
implies a channel-averaged transverse fraction of $2/3$. The bare edge stiffness is therefore
\[
J_{\mathrm{bare}}=\frac{2}{3}\eta_*=0.0199112296.
\]
For a rooted $z=4$ coarse adjacency graph, the tree-to-lattice map gives
\[
J_{\mathrm{eff}}^{\mathrm{tree}}=\frac{J_{\mathrm{bare}}}{z-1}=\frac{2\eta_*}{9}=0.0066370765.
\]

This is the first place where local closure data become a transport law. The tetrahedral identity fixes the isotropic projection, and the rooted branching structure determines how much of the microscopic edge penalty survives as net outward propagation on the coarse graph.

\subsection*{C.2 Horizon target and shell convergence}

The horizon-capacity target is
\[
\sigma_*=\frac{\pi}{g_{\mathrm{share,eff}}}=0.42340665.
\]
At the derived coupling the explicit shell values are
\[
\sigma_{\mathrm{ind}}^{(2)}=0.42143,
\qquad
\sigma_{\mathrm{ind}}^{(3)}=0.42166,
\qquad
\Delta_{2\to 3}=0.00023.
\]
The residual shift from the target is already small and stable by shell depth $r=2$, isolating the remaining correction to the loopy local-return sector rather than a broad nonlocal ambiguity.

The shell calculation places the tree branch near the target and assigns the residual discrepancy to local returns. No broad long-range correction is required at this order.

\subsection*{C.3 Finite-loop self-energy closure}

The minimal loopy correction is organized as a local Dyson dressing:
\[
J_{\mathrm{eff}}^{(\mathrm{ren})}
=
\frac{J_{\mathrm{eff}}^{\mathrm{tree}}}{1+J_{\mathrm{eff}}^{\mathrm{tree}}\Sigma_{\mathrm{ret}}}.
\]
The specified minimal return operator uses
\[
\Sigma_{\mathrm{ret}}=7+\frac{2}{9}=\frac{65}{9},
\]
The two terms have distinct return-channel origins. A short return motif leaves a shared face, explores a local closed loop, and re-enters the same coarse edge before contributing to net long-range transport. In the canonical label basis $m=-3,-2,\ldots,3$, the minimal operator contains seven label-diagonal returns, one for each face-label channel, contributing
\[
\mathrm{Tr}(I_7)=7.
\]
In addition to these label-preserving loops, permutation symmetry allows one collective mode shared across all channels. Writing
\[
P_{\mathrm{sing}}=|u\rangle\langle u|,
\qquad
u=\frac{1}{\sqrt 7}(1,1,\ldots,1),
\]
this shared return is rank one. Permutation symmetry permits this singlet but does not fix its weight relative to the diagonal returns or exclude longer return motifs. The minimal completion assigns it the same $2/3$ transverse projection used in the tree coupling and the rooted return factor $1/(z-1)=1/3$ on the $z=4$ graph. The collective contribution is then
\[
\mathrm{Tr}\!\left(\frac{2}{3}\frac{1}{3}P_{\mathrm{sing}}\right)=\frac{2}{9},
\]
since $\mathrm{Tr}(P_{\mathrm{sing}})=1$. Equivalently,
\[
R_{\mathrm{ret}}=I_7+\frac{2}{9}P_{\mathrm{sing}},
\qquad
\Sigma_{\mathrm{ret}}=\mathrm{Tr}(R_{\mathrm{ret}})=7+\frac{2}{9}.
\]
Thus $65/9$ is fixed inside the stated minimal return operator, not derived from the closure ensemble alone. A microscopic graph-return calculation must derive the relative diagonal and singlet weights, test longer motifs, and determine whether the operator closes on these channels. Within this conditional completion,
\[
c_{\mathrm{loop}}^{(\mathrm{ren})}
\equiv
\frac{J_{\mathrm{eff}}^{(\mathrm{ren})}}{J_{\mathrm{eff}}^{\mathrm{tree}}}
=
\frac{1}{1+J_{\mathrm{eff}}^{\mathrm{tree}}\Sigma_{\mathrm{ret}}}
\approx 0.95426,
\]
and
\[
J_{\mathrm{eff}}^{(\mathrm{ren})}\approx 0.00633348.
\]
This reproduces the shell-target crossing near $J_{\mathrm{bare,cross}}\sim 0.019$ at the stated level of agreement.

The local Dyson dressing corrects the tree branch for short motifs that recycle amplitude before it contributes to coarse transport. Summing those returns determines the renormalized coupling from the tree coupling.

\subsection*{C.4 Euclidean-action normalization and continuum stiffness}

The lattice quadratic form is interpreted canonically as a Euclidean action weight,
\[
\frac{I_E}{\hbar}
=
\frac{J_{\mathrm{eff}}^{(\mathrm{ren})}}{2}
\sum_{a,i}(Q_a-Q_{a+L_*\hat n_i})^2,
\]
where $a$ runs over one sublattice representative of each bipartite primitive cell and $i$ runs over its four outgoing bonds. Every undirected nearest-neighbor edge is counted once. This convention is essential: summing the four bonds from both sublattices would double the variation and the source normalization. The microscopic four-cell is
\[
\Delta V_4=\frac{L_*^4}{c},
\]
a coarse-graining convention rather than derived cell geometry (Section~9; the abstract cell complex has no regular-tetrahedron Euclidean embedding).
Because every edge occurs once, varying the discrete action gives $J_{\mathrm{eff}}^{(\mathrm{ren})}L_\diamond Q=s$ for a source term $-\sum_a s_aQ_a$, with $L_\diamond=4I-A$. The diagonal inverse of this same unnormalized Laplacian is the $G_{\mathrm{tet}}(0)$ used in C.5, so a point source obeys $Q(0)=sG_{\mathrm{tet}}(0)/J_{\mathrm{eff}}^{(\mathrm{ren})}$. This removes the factor-of-two ambiguity between the stiffness and source conventions.
The same tetrahedral identity then yields
\[
\gamma_Q=\frac{4\hbar c}{3L_*^2}J_{\mathrm{eff}}^{(\mathrm{ren})}
\]
for the occupancy field $Q_{\mathrm{occ}}$. With the horizon-capacity normalization
\[
S=\pi Q_{\mathrm{occ}},
\]
the canonical convention $\frac{\gamma}{2}(\partial S)^2$ gives
\[
\gamma
=
\frac{4\hbar c}{3\pi^2L_*^2}J_{\mathrm{eff}}^{(\mathrm{ren})}
=
\frac{4\hbar c}{3\pi^2L_*^2}
\frac{2\eta_*/9}{1+(2\eta_*/9)(65/9)}.
\]
Using the substrate-induced scale
\[
G_*:=\frac{c^3L_*^2}{\hbar},
\]
this is
\[
\gamma
=
\frac{4J_{\mathrm{eff}}^{(\mathrm{ren})}}{3\pi^2}\frac{c^4}{G_*}
\approx
8.556\times10^{-4}\frac{c^4}{G_*}.
\]

This step converts the dimensionless lattice weighting into the dimensionful continuum stiffness used by the weak-field action, after Euclidean normalization and faithful sector-resolution scale setting.

\subsection*{C.5 Local defect insertion and the source-side lattice constant}

The stiffness-side matching is not the only UV quantity that can be closed locally. For the canonical rigid defect insertion, excluding one of the seven admissible face labels from one face removes exactly one-seventh of the isotropically averaged local partition weight. Therefore the logarithm of the isotropically averaged partition ratio is exactly
\[
\Delta S_{\mathrm{def}}
:=
-\ln\!\Bigl\langle \frac{Z_{\mathrm{def}}}{Z_{\mathrm{vac}}}\Bigr\rangle_{\mathrm{iso}}
=
\ln\!\frac{7}{6}.
\]
This is the exact isotropic source benchmark in the canonical seven-label ensemble. The isotropically averaged defect free-energy cost differs from it only at $O(10^{-5})$ because the admissibility weighting breaks label symmetry only weakly.

The local source benchmark $\ln(7/6)$ should not be confused with the elementary fermionic one-bit anchor $\ln2$. The former is the isotropically averaged partition-ratio shift produced by removing one admissible label from the seven-label boundary ensemble. The latter is the intrinsic binary entropy of an elementary occupied/unoccupied fermionic face-exclusion defect. The source theorem uses $\ln(7/6)$ to normalize the local scalar insertion into the lattice response, while the electron anchor uses $\ln2$ to fix the mass--entropy unit of the elementary fermionic defect.

\paragraph{Exclusion sufficiency lemma.}
The reduction of the $1680$-state boundary ensemble to a single scalar source is exact at the classical source level. The canonical exclusion defect is the vacuum ensemble conditioned on the exclusion event $A$, $P_{\mathrm{def}}=P_{\mathrm{vac}}(\,\cdot\,|A)$, so its likelihood ratio is
\[
\frac{dP_{\mathrm{def}}}{dP_{\mathrm{vac}}}=\frac{\mathbf 1_A}{P_{\mathrm{vac}}(A)},
\]
a function of the exclusion indicator $X=\mathbf 1_A$ alone. The indicator is therefore a sufficient statistic for distinguishing the defect ensemble from the vacuum, and coarse-graining onto it preserves the full classical relative information,
\[
D(P_{\mathrm{def}}\Vert P_{\mathrm{vac}})=-\ln P_{\mathrm{vac}}(A)=D(P_{\mathrm{def}}^X\Vert P_{\mathrm{vac}}^X).
\]
The scope of the lemma is exactly its statement: one scalar statistic carries all the local classical source information distinguishing the canonical exclusion-defect ensemble from the vacuum. It does not by itself establish that the condensate has a single infrared mode, that orientation or defect-species structure is dynamically irrelevant, or that spatial correlations carry no further information; those are dynamical questions, and the open one is stated in Appendix~H.7. For independent exclusions the source is exactly additive, $-\ln P(\bigcap_iA_i)=\sum_i-\ln P(A_i)$, so the leading dilute source is linear in the defect number, with correlation corrections entering at pair order in the density.

To propagate that local defect into the lattice field equation one needs the on-site Green function of the tetrahedral/diamond nearest-neighbor Laplacian. The four-valent diamond graph is adopted as the coarse adjacency realization of the abstract face-sharing tetrahedral boundary complex; the complex enters combinatorially, since regular tetrahedra admit no face-to-face Euclidean tessellation (the dihedral angle $\arccos(1/3)\simeq70.53^\circ$ does not divide $360^\circ$), so the adopted graph and its Green function are the working objects while cell volumes and face areas remain conventions of the coarse map. Eliminating the two-sublattice structure gives the standard Brillouin-zone representation for the diamond lattice Green function \citep{Guttmann2010}
\[
G_{\mathrm{tet}}(0)
=
\frac{1}{(2\pi)^3}
\int_{[-\pi,\pi]^3}
\frac{4\,d^3k}
{16-\left|1+e^{ik_1}+e^{ik_2}+e^{ik_3}\right|^2}.
\]
This integral is the reproducible source-side lattice constant: it is the self-energy of a unit point insertion for the same scalar mode whose long-wavelength stiffness was matched in Appendix~C.4. Joyce's exact evaluation of the diamond-lattice Green function, in this normalization, gives \citep{Joyce1994,Guttmann2010}
\[
G_{\mathrm{tet}}(0)
=
\frac{3\Gamma(1/3)^6}{2^{14/3}\pi^4}
=
0.4482203943883814\ldots .
\]
Thus the source-side graph constant is an exact lattice invariant. Direct quadrature with endpoint extrapolation reproduces the same value.

\paragraph{Green-tensor transverse response.}
The same graph response also fixes the transverse export weight $w_\perp$, the single graph quantity that propagates downstream into the horizon channel count of Appendix~F.5 and into the global $\ln(3/2)$ closure-saturation factor of the scale-setting relation. We compute it directly from the four nearest-neighbor bond frame, with no fitting freedom.

Let $d_i$, $i=1,\ldots,4$, be the four diamond nearest-neighbor bond directions,
\[
d_1=\frac{(1,1,1)}{\sqrt3},\quad
d_2=\frac{(1,-1,-1)}{\sqrt3},\quad
d_3=\frac{(-1,1,-1)}{\sqrt3},\quad
d_4=\frac{(-1,-1,1)}{\sqrt3}.
\]
They obey
\[
\sum_{i=1}^4 d_i^a d_i^b=\frac43\delta^{ab}.
\]
Distributing the scalar on-site Green response over the tetrahedral bond frame gives
\[
{\cal G}_{\rm loc}^{ab}
=
G_{\mathrm{tet}}(0)\frac14\sum_i d_i^a d_i^b
=
\frac{G_{\mathrm{tet}}(0)}{3}\delta^{ab}.
\]
For a local horizon normal $\hat r$,
\[
P_\perp^{ab}=\delta^{ab}-\hat r^a\hat r^b,
\]
so
\[
G_\perp=P_\perp^{ab}{\cal G}_{\rm loc}^{ab}
=
\frac23G_{\mathrm{tet}}(0).
\]
Thus the transverse export weight $w_\perp=2/3$ is the transverse part of the exact local graph response.

Using the field normalization $S=\pi Q_{\mathrm{occ}}$, the rigid local defect shift is
\[
\delta Q_{\mathrm{def}}=\frac{\Delta S_{\mathrm{def}}}{\pi}
=
\frac{\ln(7/6)}{\pi}.
\]
The corresponding local source amplitude in lattice units is therefore
\[
s_{\mathrm{def}}
=
J_{\mathrm{eff}}^{(\mathrm{ren})}
\frac{\delta Q_{\mathrm{def}}}{G_{\mathrm{tet}}(0)},
\]
so that
\[
\frac{s_{\mathrm{def}}}{J_{\mathrm{eff}}^{(\mathrm{ren})}}
=
\frac{\ln(7/6)}{\pi G_{\mathrm{tet}}(0)}
=
0.109472228\ldots
\]
is a pure number fixed by the same UV lattice geometry.

The Green-function constant turns the local insertion into a continuum source theorem. Defining the defect-entropy density by
\[
\sigma_{\mathrm{def}}=\frac{\rho}{\kappa_m(L_*)},
\]
the tetrahedral projection used in the stiffness mapping gives
\[
\nabla^2\delta S
=
-\frac{3L_*}{4G_{\mathrm{tet}}(0)}\,\sigma_{\mathrm{def}}.
\]
Equating this with the weak-field source equation
\[
\nabla^2\delta S=-\frac{\kappa}{\gamma}\rho
\]
closes the canonical source-to-stiffness ratio:
\[
\boxed{
\frac{\kappa}{\gamma}
=
\frac{3L_*}{4G_{\mathrm{tet}}(0)\kappa_m(L_*)}
}.
\]
This is the source-side counterpart of the stiffness derivation. The edge-kernel calculation fixes how the scalar capacity mode resists gradients; the Green-matched defect calculation fixes how localized matter defects source that same mode.

\paragraph{The three cancellations, made explicit.}
The passage from the dimensionless lattice insertion $s_{\mathrm{def}}/J_{\mathrm{eff}}^{(\mathrm{ren})}=\ln(7/6)/(\pi G_{\mathrm{tet}}(0))$ to the source theorem turns on three reductions, none of which leaves a free constant.
\emph{(i) The factor $\pi$ cancels against the field normalization.} The insertion is written for the occupancy field, where one isotropically averaged defect shifts $\delta Q_{\mathrm{def}}=\ln(7/6)/\pi$; a fixed defect type $(\ell,f)$ instead carries $-\ln P(A_{\ell f})$, ranging over $[0.14933,\,0.15819]$ with mean $0.15415642$ against $\ln(7/6)=0.15415068$, so the source theorem uses the exact isotropic benchmark under the assumption that the coarse defect population is isotropically averaged over face and label types. Converting to the entropy field through $S=\pi Q_{\mathrm{occ}}$ multiplies by $\pi$, so the $S$-field shift is $\delta S_{\mathrm{def}}=\pi\,\delta Q_{\mathrm{def}}=\ln(7/6)$ and the explicit $\pi$ does not survive into the source theorem.
\emph{(ii) $\ln(7/6)$ is an isotropic source benchmark, not a universal per-defect cost.} For a coarse population averaged over face and label types, the exact mean partition ratio is $6/7$, so the corresponding annealed source normalization is $\ln(7/6)$. A fixed defect type instead carries $-\ln P(A_{\ell f})$. Writing the isotropically averaged source as $\sigma_{\mathrm{def}}=\rho/\kappa_m(L_*)$ folds that benchmark and the mass--entropy unit into the density. The geometric coefficient $3L_*/(4G_{\mathrm{tet}}(0))$ is fixed by the tetrahedral projection and the cell length and contains no $\ln(7/6)$.
\emph{(iii) $J_{\mathrm{eff}}^{(\mathrm{ren})}$ cancels in the ratio.} The source amplitude $s_{\mathrm{def}}$ and the stiffness $\gamma$ each carry one power of the renormalized edge coupling, so it cancels in the source-to-stiffness ratio $\kappa/\gamma$, leaving the lattice-geometric constant $3L_*/(4G_{\mathrm{tet}}(0)\kappa_m(L_*))$. The renormalized loop coupling therefore drops out of the observable normalization.

\paragraph{Physical geometry--capacity matching.}
The source theorem fixes $\kappa/\gamma$ in the chosen entropy units. With the Compton-calibrated one-bit dictionary $\kappa_m(L_*)=\hbar/(cL_*\ln2)$, the weak-field bridge gives
\[
G=\frac{c^2}{8\pi S_\infty}\frac{\kappa}{\gamma}
=\frac{c^3L_*^2}{\hbar}\,
\frac{3\ln2}{32\pi G_{\mathrm{tet}}(0)S_\infty}.
\]
Define $S_0=3\ln2/[32\pi G_{\mathrm{tet}}(0)]$. The selected canonical matching is
\[
\boxed{S_\infty^{\rm cell}=S_0=0.0461482516\ldots,\qquad G=G_* .}
\]
Equivalently it imposes $n_{\rm hor}S_\infty^{\rm cell}=1/4$ with the response number of F.5. The coefficient $1/4$ and the identification $L_P(G_*)=L_*$ provide the physical Einstein/horizon normalization. This is a matching prescription of the chosen branch. The source theorem by itself gives $G/G_*=S_0/S_\infty$.

An entropy-unit change $S'=bS$ gives $\gamma'=\gamma/b^2$, $\kappa'=\kappa/b$, and $S_\infty'=bS_\infty$. Hence $\kappa/(\gamma S_\infty)$ is invariant. Cell and horizon representatives can use different entropy units once the physical matching has been fixed. The rescaling carries that matching between representations and leaves the observable $G$ unchanged. At fixed entropy convention and source coefficients, choosing a different physical background value $S_\infty=\beta S_0$ would give $G=G_*/\beta$. The dimensionless transfer and counting results do not determine $\beta$. The Einstein/horizon matching selects $\beta=1$; a microscopic derivation of that selection would require additional geometric dynamics or boundary data.

\subsection*{C.6 Local susceptibility cross-check}

The exact local moments of the admissibility-closed ensemble supply an independent non-degeneracy check on the source-side result. From the variance in Appendix~B,
\[
a_{\mathrm{UV}}
:=
\frac{1}{\mathrm{Var}_{\eta_*}(K^2)}
=
0.0637390269,
\]
which is the local zero-mode inverse susceptibility of the closure scalar. Two features of this number matter for the source theorem in C.5. First, $a_{\mathrm{UV}}$ is finite and positive, confirming that the closed branch at $\eta_*$ has a non-degenerate zero-mode response rather than a critical singularity that would invalidate the linear Green-function matching used to derive $\kappa/\gamma$. Second, the same local susceptibility controls the stability of the closure mode propagated into the shell and loop calculations, so the residual fractional discrepancy between the local benchmark and $\sigma_*$ belongs to the loop sector rather than the local closure.

The actual closure of $\kappa/\gamma$ in the canonical branch is the Green-matched theorem in Appendix~C.5. The role of $a_{\mathrm{UV}}$ here is restricted to confirming non-degeneracy of the local mode being matched.

\subsection*{C.7 UV-to-IR payoff}

At this stage the weak-field UV coefficient chain is explicit:
\[
\Omega_{\mathrm{tet}}
\rightarrow
g_{\mathrm{share,eff}}
\rightarrow
L_*
\rightarrow
J_{\mathrm{bare}}
\rightarrow
J_{\mathrm{eff}}^{\mathrm{tree}}
\rightarrow
\Sigma_{\mathrm{ret}}
\rightarrow
J_{\mathrm{eff}}^{(\mathrm{ren})}
\rightarrow
\gamma.
\]
The same chain feeds
\[
a_0=\frac{cH_0g_{\mathrm{share,eff}}}{4\pi^2},
\]
and the decorated support-to-rate action gives the substrate-induced scale
\[
G_*=\frac{c^3L_*^2}{\hbar}
=
\frac94\,\frac{\hbar c}{m_e^2}
Z_e^2\ln^2\!\left(1-e^{-7g_{\mathrm{share,eff}}}\right).
\]
The Green-matched source theorem fixes
\[
\frac{\kappa}{\gamma}
=
\frac{3L_*}{4G_{\mathrm{tet}}(0)\kappa_m(L_*)}.
\]
The weak-field bridge then converts this source-to-stiffness ratio into the observed Newtonian normalization through the invariant combination $\kappa/(\gamma S_\infty)$. The comparison with the same $G_*$ assigned by the electron-anchored support-to-rate branch tests that this normalization propagates the substrate scale coherently; it is not a disjoint-input second derivation of $G$, since the matched route carries $L_*$ within it.

The coefficients this appendix supplies to the main weak-field chain are $J_{\mathrm{bare}}$, $J_{\mathrm{eff}}^{\mathrm{tree}}$, $\Sigma_{\mathrm{ret}}$, $J_{\mathrm{eff}}^{(\mathrm{ren})}$, $\gamma$, and the Green-matched source projection $\kappa/\gamma$. The remaining uses of $S_\infty$ belong to the fixed normalization of the weak-field bridge, not to the source sector itself.

\subsection*{C.8 Event, volume, and gravitational calibration}
\label{aqg:calibration}

The renewal cadence determines the causal length \(L_*=c\tau_*\).
A physical event density and a spatial area require additional geometric
data. Keep the canonical convention of C.4--C.5 explicit, and denote a
physical event four-volume in length coordinates by \(v_{\rm ev}\).

Consider a homogeneous Poisson event law in a four-dimensional causal
diamond, conditioned on its total count when that count is fixed.
If its proper duration is \(T\), its volume is
\(\pi(cT)^4/24\). Define the longest-chain constant \(m_4\) by
\[
L_{\rm chain}\sim m_4
\left(\frac{\pi(cT)^4}{24v_{\rm ev}}\right)^{1/4}.
\]
Identifying a longest-chain step with the cadence \(\tau_*\) gives
\[
\boxed{v_{\rm ev}=\frac{\pi}{24}m_4^4L_*^4.}
\]
This uses a spatial event law and a clock identification in addition to
temporal renewal. The Lorentz-invariance theorem for sprinklings
\citep{AQGBombelli2009} assumes that event law. Its finite-valency
obstruction also prevents identifying a covariantly selected finite
nearest-neighbor graph with the entire sprinkled causal order. The
four-port graph here is a relational coarse structure.

For a fixed total of \(N\) uniformly distributed events, a subregion of
volume fraction \(f\) has mean \(Nf\) and variance \(Nf(1-f)\).
Disjoint fractions \(f_1,f_2\) have covariance \(-Nf_1f_2\).
Fixed total count therefore permits local fluctuations. The stronger
fixed-volume-form premise used in O.1 remains separate.

Adopt one event per paired primitive cell per tick, or one half event per
tetrahedral site per tick, as the mean event assignment. The physical
site volume and regular coarse area calibration are then
\[
V_{\rm site}=\frac{v_{\rm ev}}{2L_*},\qquad
a_{\rm reg}=(6\sqrt2\,V_{\rm site})^{1/3},\qquad
A_{\rm ch}=\frac{\sqrt3}{4}a_{\rm reg}^2.
\]
The regular tetrahedron specifies the mean area--volume conversion.
It supplies no face-to-face Euclidean tessellation. Combined with the
absolute angular rule of Q.11,
\[
\boxed{\frac{G_0}{G_*}
=\frac{\sqrt3}{16\ln2}
\left(\frac{3\sqrt2\,v_{\rm ev}}{L_*^4}\right)^{2/3}
=0.1055\ldots\,m_4^{8/3}.}
\]
At the illustrative chain estimate \(m_4=2.35\),
\[
\begin{aligned}
v_{\rm ev}/L_*^4&=3.99217968,&
a_{\rm reg}/L_*&=2.56812078,\\
A_{\rm ch}/L_*^2&=2.85582456,&
G_0/G_*&=1.03002098.
\end{aligned}
\]
Equality \(G_0=G_*\) would require \(m_4\simeq2.3241\) within these
premises. The range \(m_4=2.35\pm0.10\) maps to approximately
\(0.92\)--\(1.15\) for \(G_0/G_*\); it is a calibration sensitivity
range, with substantial uncertainty in finite-count extrapolation.

\paragraph{Compatibility with a quantum pair volume.}
Restrict the additive primitive flux volume to the pair sector of H.12.
Its 13 eigenvalues span the interval
\[
1.71248L_*^3\le V_{\rm pair}\le4.53199L_*^3.
\]
The maximally mixed state has mean \(3.00949L_*^3\).
A maximum-entropy state has
\[
\rho_\lambda\propto e^{\lambda V_{\rm pair}},\qquad
\lambda=1.18686L_*^{-3},\qquad
\langle V_{\rm pair}\rangle=v_{\rm ev}/L_*=3.99218L_*^3.
\]
Its existence establishes compatibility; the mean was imposed in choosing
\(\lambda\). Dynamics must select the volume state. At primitive \(j=3\),
shape uncertainty also remains finite: the minimum RMS face-cosine
deviation from regularity is about \(0.288\).

\paragraph{The scalar normalization at a different physical volume.}
The coefficients of C.4--C.5 use the canonical reference cell and bond
displacement. Replacing that volume while retaining the same displacement
changes their physical normalization. In the tested event calibration,
keeping the bond displacement at \(L_*\) gives an inferred scalar
normalization near \(3.99G_*\). Matching the scalar and cut values would
require a displacement near \(1.97L_*\); neither the regular edge nor the
diamond bond at that volume supplies this value automatically.

On the selected Einstein continuation, the scalar coefficients instead
follow by reducing the same metric action. In tick units and in the
canonical entropy variable,
\[
\kappa_{\rm phys}=\frac1{2S_\infty},\qquad
\gamma_{\rm phys}=\frac1{16\pi G_0S_\infty^2}.
\]
For \(S_\infty=0.04614825156\) these are \(10.83465\) and \(9.06930\).
They are the continuum coefficients of that parent action. Equating them
to a native lattice source and stiffness remains a matching calculation.
Neither entropy-unit changes nor a choice of event bookkeeping can supply
that equality by themselves.

\section*{Appendix D: Weak-Field Technical Derivations, Electron Anchor, and EFT Consistency}

Appendix~D collects the weak-field derivations that are central but too dense for the main line: the bridge law, Newtonian recovery, the electron anchor, and the EFT consistency checks.

\subsection*{D.1 Bridge law from the reduced action}

The weak-field bridge follows from an action equivalence and needs no extrapolated exponential lapse. The Einstein scalar-constraint action reduces to
\[
I_{\mathrm{Newton}}[\Phi]
=\int dt\,d^3x
\left[-\frac{(\nabla\Phi)^2}{8\pi G}-\rho\Phi\right].
\]
The capacity action becomes the same functional, up to the constant $Z_S=2\kappa S_\infty/c^2$, under
\[
\delta S=-\frac{2S_\infty}{c^2}\Phi.
\]
Hence
\[
\frac{\Phi}{c^2}=-\frac{\delta S}{2S_\infty}
\]
is the unique linear field redefinition matching both the kinetic and source terms with the UV normalization. The bounded nonlinear rule is instead $N^2=q=1-\delta S/S_\infty$ in a static branch. An exponential lapse shares the first derivative at the vacuum point but is not used globally because it has no finite-capacity endpoint. The complete derivation is in Appendix~N.

\subsection*{D.2 Point source, Newton limit, and lensing}

In the renormalized static branch,
\[
\nabla^2\delta S=-\frac{\kappa}{\gamma}\rho.
\]
For a point source $M$,
\[
\delta S(r)=\frac{\kappa M}{4\pi\gamma r},
\qquad
g(r)=\frac{c^2\kappa}{8\pi\gamma S_\infty}\frac{M}{r^2}
=
\frac{GM}{r^2}.
\]
For the ordinary longitudinal branch, variation of the unreduced Einstein scalar action gives
\[
\Phi=\Psi
\]
with asymptotically flat boundary conditions. This equality follows from the spatial metric constraint, not from a canonical capacity stress tensor. The effective-halo rewrite of the additional galactic response is
\[
\rho_{\mathrm{halo}}(r)=\frac{1}{4\pi Gr^2}\frac{d}{dr}\Big[r^2(g_{\mathrm{obs}}-g_{\mathrm{bar}})\Big].
\]
Thus the same metric potential controls orbital dynamics and light bending in the baseline Einstein branch. The separate carrier-resolved contact of Section~16 specifies $\Delta\Phi_\perp=\Delta\Psi_\perp$ at leading quasistatic order and uses the same induced density for support and lensing. Its nonlinear Ward identity depends on covariant variation of the carrier records, as Appendix~N.8 states.

\subsection*{D.3 Electron anchor and composite matter}

The canonical fermionic entropy increment is
\[
\Delta S_f=\ln 2.
\]
The UV mass normalization is
\[
\kappa_{m,\mathrm{UV}}=\frac{\hbar}{cL_*}\frac{1}{\ln 2},
\]
and the running law in the closed branch is
\[
\kappa_m(\ell)=\kappa_{m,\mathrm{UV}}\left(\frac{L_*}{\ell}\right)^{1+\alpha_{\mathrm{cl}}},
\qquad
\alpha_{\mathrm{cl}}=0.
\]
At the electron Compton scale $\ell=\lambda_e$ this gives
\[
\kappa_m(\lambda_e)=\frac{m_e}{\ln 2},
\]
which is the elementary anchor used here. Composite hadrons are not reduced to a bare constituent count. Their mass budget is assigned to a dressed bound-state entropy
\[
m_{\mathrm{hadron}}=\kappa_m(\ell_H)S_{\mathrm{ent},H}^{\mathrm{dressed}},
\]
whose microscopic decomposition must include confinement, gluonic structure, trace-anomaly contributions, and chiral vacuum reorganization.

The electron is an elementary one-bit defect and provides the dimensional anchor. Hadronic inertial content instead belongs to a dressed entropy budget that includes binding and vacuum structure; a bare constituent count is insufficient.

\subsection*{D.4 Faithful sector resolution and induced \texorpdfstring{$G_*$}{G*}}

The electron anchor enters the gravitational normalization through the absolute substrate length. The exact inputs are the seven-channel entropy, the memoryless kernel selected by faithful resolution, channel factorization on the lightest branch, and the decorated marked-transfer vertex of Appendix~H. The calculation below separates the determinant recurrence, positive survival gap, and finite marked response before combining them in the physical cell length.

\paragraph{The recurrence derivation of the dictionary.}
Appendices~H.2, H.4, and H.5 show that faithful full-support resolution selects memoryless replacement and that the lightest branch has seven factorized channels. On the commutative boundary algebra define $(\mathsf Rf)(b)=p_*(b)f(b)$ and $\tau_p(f)=\sum_bp_bf(b)$. The state-weighted determinant is
\[
\Delta_{\tau_p}(\mathsf R)
=\exp[\tau_p(\ln\mathsf R)]
=\exp\!\left(\sum_bp_b\ln p_b\right)
=e^{-g_{\mathrm{share,eff}}}.
\]
Multiplicativity gives $e^{-7g_{\mathrm{share,eff}}}$ on the seven-channel product. This is the quenched geometric transfer rate of the record-conditioned likelihood operator. It is not the Born probability of a specified microstate or the collision probability of two renewed blocks.

The corresponding history statement follows from the exact recurrence theorem. For a stationary ergodic finite-alphabet process, if $R_n$ is the first return time of a length-$n$ block, then
\[
\lim_{n\to\infty}\frac{1}{n}\ln R_n=h
\qquad\text{almost surely},
\]
where $h$ is the Shannon entropy rate \citep{WynerZiv1989,OrnsteinWeiss1993}. This theorem selects Shannon entropy over raw state count and R\'enyi-2 entropy for long typical blocks. It does not by itself identify the seven distinguishable sector layers with a long temporal block, and it gives no exact finite-$n$ return probability. For the present one-layer distribution,
\[
H_1=7.4198000,
\qquad
H_2=-\ln\sum_b p_b^2=7.4125486,
\qquad
H_\infty=-\ln\max_b p_b=7.1343850.
\]
The distinctions are numerically consequential after seven layers. The refresh projector $|\sqrt p\rangle\langle\sqrt p|$ carries the participation scale $e^{H_2}$, while the determinant and almost-sure Lyapunov rate use $H_1$. They are different observables. With
\[
p_{\rm rec}\equiv e^{-7g_{\mathrm{share,eff}}},
\]
the positive survival operator has nonzero eigenvalue $1-p_{\rm rec}$. The decorated marked vertex supplies the electron factor $Z_e$. With the additive-generator prescription of H.9, the dressed electron gap is
\[
E_e=-\frac{3\hbar Z_e}{2\tau_*}\ln(1-p_{\rm rec}).
\]
Identifying the lowest charged gap with $m_ec^2$ fixes the cadence from the already-declared electron anchor,
\[
\tau_*=-\frac32Z_e\tau_e\ln(1-p_{\rm rec}),
\qquad
L_*=c\tau_*.
\]
Thus
\[
L_*=-\frac32Z_e\lambda_e\ln\!\left(1-e^{-7g_{\mathrm{share,eff}}}\right)
=\frac32Z_e\lambda_e e^{-7g_{\mathrm{share,eff}}}
\left[1+O\!\left(e^{-7g_{\mathrm{share,eff}}}\right)\right].
\]
The logarithmic survival correction is $p_{\rm rec}/2\simeq1.4\times10^{-23}$. The finite closure-response correction is $Z_e-1=0.00530828$. No second dimensional input, maximum-throughput clock, or tetrahedral-diameter equality is introduced.

\paragraph{The transfer-support dual.}
A spatial version uses the same dressed support and assigns the leading step count
\[
N_{\rm eff}=\frac{2}{3Z_e}e^{7g_{\mathrm{share,eff}}}.
\]
The positive transfer spectrum gives the temporal relation $\tau_e=N_{\rm eff}\tau_*$. Assigning the same count to space gives $\lambda_e=N_{\rm eff}L_*$ and hence $L_*=c\tau_*$. The spatial assignment is a continuum matching prescription, not a consequence of the scalar rest spectrum. A microscopic orbital transfer must reproduce it through its small-momentum dispersion; the graph bond length cannot be identified with $L_*$ without that derivation.

\paragraph{Internal locality and the native vertex.}
The selected replacement kernel reaches the full ensemble. The tested single-label move class does not: shifts preserving injectivity split each parity copy into $4!=24$ ordering sectors of $\binom{7}{4}=35$ states because two slots cannot exchange order without a collision. Its strength-weighted walk saturates at entropy $7.374$, below $g_{\mathrm{share,eff}}=7.4198$. This excludes that implementation. Appendix~H.8 gives the whole-tetrahedron replacement gate, and H.9 decorates it with the fresh-state amplitude and charged event. The native gate is carried by the separate capacity/record factor; a particular GFT condensate realization may implement it dynamically, while the ordinary geometric propagation is supplied by the controlled factorized EPRL-class branch.

The reduced electron Compton wavelength is
\[
\lambda_e=\frac{\hbar}{m_ec},
\]
and the admissibility-closed sharing entropy is
\[
g_{\mathrm{share,eff}}=7.41980002357.
\]
The closure fixed point also gives
\[
\langle K^2\rangle_{\eta_*}=\frac{3}{2\eta_*},
\qquad
\eta_*=0.0298668443935,
\]
so the closure-saturation factor is
\[
C_{\mathrm{cl}}:=\eta_*\langle K^2\rangle_{\eta_*}=\frac32,
\qquad
C_{\mathrm{cl}}^{-1}=\frac23.
\]

Separating the baseline recurrence from the finite marked response gives
\[
\ln\!\left(\frac{\lambda_e}{L_*}\right)
=7g_{\mathrm{share,eff}}-\ln\!\left(\frac32\right)-\ln Z_e
+O(e^{-7g_{\mathrm{share,eff}}}).
\]
Equivalently,
\[
\boxed{
L_*^{\rm leading}=\frac32\,Z_e\lambda_e e^{-7g_{\mathrm{share,eff}}}
}.
\]
The exact expression is $L_*=-(3/2)Z_e\lambda_e\ln(1-e^{-7g_{\mathrm{share,eff}}})$. It gives
\[
L_*=1.6162537024\times10^{-35}\ {\rm m},
\]
which differs from the CODATA Planck length by about $8.2\times10^{-7}$ in fractional terms, or $8.2\times10^{-5}$ percent.

The induced gravitational scale follows from the same algebra used in the stiffness matching:
\[
G_*:=\frac{c^3L_*^2}{\hbar}.
\]
Substituting the exact marked-transfer relation gives the closed form
\[
\boxed{
G_*=
\frac{c^3}{\hbar}
\left[-\frac32\,Z_e\lambda_e\ln\!\left(1-e^{-7g_{\mathrm{share,eff}}}\right)\right]^2
=
\frac94\,\frac{\hbar c}{m_e^2}\,
Z_e^2\ln^2\!\left(1-e^{-7g_{\mathrm{share,eff}}}\right)
}.
\]
Numerically,
\[
G_*=6.6742890813\times10^{-11}\ {\rm m^3\,kg^{-1}\,s^{-2}},
\]
or $-0.073\sigma$ relative to CODATA. The same $\zeta_*$ gives the muon and tau ratios at $-0.535\sigma$ and $-0.125\sigma$. One finite action therefore supplies a single marked response to all three observables, with its coefficient and routing fixed jointly rather than adjusted independently for the three observables. The discrepancies were already known while the vertex was being constructed, so these are high-precision postdictions. Their evidential content is the common action, the absence of a continuous fit, and the nearby kernels rejected by the finite audits.

The factor $3/2$ is fixed structurally by the tetrahedral transverse-export geometry. The tetrahedral identity gives the transverse export fraction
\[
\frac14\sum_{i=1}^4(\hat n_i\cdot\hat u)^2=\frac13,
\qquad
f_\perp=1-\frac13=\frac23,
\]
so $f_\perp^{-1}=3/2$. The admissibility fixed point returns the same value,
\[
C_{\mathrm{cl}}^{-1}
=
\left(\eta_*\langle K^2\rangle_{\eta_*}\right)^{-1}
=
\frac23,
\]
but it is built in. The condition defining $\eta_*$ is $\langle K^2\rangle_\eta=3/(2\eta)$, so $C_{\mathrm{cl}}=\eta_*\langle K^2\rangle_{\eta_*}=3/2$ holds by construction. It is not an independent cross-check. The decorated scale action reuses its reciprocal $2/3$ as the transverse export and inserts it once as a global factor, not seven times.

A discrete sensitivity audit enumerates 576 nearby formulas obtained by varying the response dimension, pair count, direction factor, determinant power, statistics, and singlet overlap. Only the decorated-vertex formula lies within one muon standard deviation. This count is not a probability distribution over theories. It shows that the landing is sensitive to the field content and graph, which is why the action and its adversarial alternatives are displayed explicitly. A shared edge mode, directed complex determinant, stable bosonic determinant, or six-state symmetric response misses the muon respectively by approximately $1.15\times10^4\sigma$, $13.8\sigma$, $56.7\sigma$, and $7.64\times10^4\sigma$. Appendix~H derives the selected entries before performing the observable trace.

The electron anchor carries three related roles. Its reduced Compton wavelength $\lambda_e$ is the non-gravitational length used in faithful sector resolution. Its status as the lightest one-bit fermionic defect identifies which elementary excitation calibrates the seven-channel dressing block. Its mass fixes the mass--entropy map through
\[
\kappa_m(\lambda_e)=\frac{m_e}{\ln2}.
\]
The same elementary defect enters the two normalization channels in distinct roles: its Compton length sets the UV cell scale, while its mass per one-bit defect entropy sets the source normalization.

The seven-sector scale-setting relation has a concrete finite-dimensional realization on a transfer-operator history space. The face label algebra is
\[
\mathcal A_7=\bigoplus_{m=-3}^{3}\mathbb C E_m,
\qquad
E_mE_n=\delta_{mn}E_m,
\qquad
\sum_{m=-3}^{3}E_m=I_7.
\]
The admissibility-closed tetrahedral state space has $1680$ oriented injective states, with stationary weight
\[
p_{\eta_*}(b)=Z^{-1}e^{-\eta_*K^2(b)}.
\]
For a single sector, the refresh kernel
\[
P_{\eta_*}(b,b')=p_{\eta_*}(b')
\]
has Perron stationary entropy $g_{\mathrm{share,eff}}$. The electron does not simultaneously occupy seven mutually exclusive face labels in one tetrahedron; the labels are sector channels in the dressing history of the one-bit defect. The appropriate support object is therefore a history space, with one admissibility-closed sector layer for each $m=-3,\ldots,3$.

Let
\[
\mathcal H_{\rm hist}
=
\bigotimes_{m=-3}^{3}\mathcal H_B^{(m)},
\qquad
\dim \mathcal H_B=1680,
\]
and let $\mathcal T_m$ act as $P_{\eta_*}$ on the $m$th factor and as the identity on the others. A representative one-pass dressing operator is
\[
\mathcal D_e^{(0)}
=
\mathcal T_{-3}\mathcal T_{-2}\cdots\mathcal T_{3}.
\]
The displayed order is only a representative of the symmetrized one-pass class. It introduces no physical ordering or extra factor of $7!$. The one-pass prescription visits each simple sector once. A history with omitted sectors is unresolved, while additional visits belong to a multi-pass dressing history rather than to repeated fermionic occupation.

A sector-resolution step should not be identified with conditioning the
1680-state ensemble on boundary states containing a given label $m$.  Such
conditioning changes the entropy and does not reproduce
$g_{\mathrm{share,eff}}$.  The sector label instead specifies which simple
face-algebra channel is being resolved while the admissibility cloud sampled
in that step remains the full closed boundary ensemble.

With the effective dimension defined by the stationary Shannon entropy of the positive history kernel, each sector contributes $g_{\mathrm{share,eff}}$, so
\[
\dim_{\rm eff}(\mathcal D_e^{(0)})
=
\exp(7g_{\mathrm{share,eff}})
=
3.60286052\times10^{22}.
\]
Only the transverse part of this support is exported into the weak-field scalar channel. The normalized export weight is
\[
w_\perp=\frac23,
\]
the same fraction fixed by the tetrahedral transverse projection and by the reciprocal closure-saturation factor. Hence
\[
\dim_{\rm eff}(\mathcal D_{e,\perp})
=
w_\perp\,\dim_{\rm eff}(\mathcal D_e^{(0)})
=
\frac23\,e^{7g_{\mathrm{share,eff}}}
=
2.40190701\times10^{22}.
\]
The decorated marked-transfer action gives the exact support-scale relation
\[
\boxed{
\frac{\lambda_e}{L_*}
=
\frac{2}{3Z_e\left[-\ln\!\left(1-e^{-7g_{\mathrm{share,eff}}}\right)\right]}
=
\frac{1}{Z_e}\dim_{\rm eff}(\mathcal D_{e,\perp})
\frac{r}{-\ln(1-r)}
}.
\]
Thus $\lambda_e/L_*=Z_e^{-1}\dim_{\rm eff}(\mathcal D_{e,\perp})[1-r/2+O(r^2)]$. Appendix~H derives the memoryless replacement, determinant recurrence, positive gap, and decorated charged vertex. Fermionic exclusion and subadditivity select $k=7$ and $\Delta_7=0$, while the marked graph fixes $Z_e$. The capacity-decorated rigging lift is exact and the selected metric-Regge branch has the Einstein TT limit.  The geometric gluing tensor and measure remain independent data not fixed by the capacity premises; the physical source of history capacity and the full source-coupled transverse/environmental spectrum also remain open.

The relation is fixed inside the displayed finite transfer action and the electron anchor. Its remaining conditionality is the broader claim that this transfer action is the charged sector of the eventual geometric GFT, not an unfixed number in the finite calculation.

This is the microscopic content of the length formula above: the elementary one-bit fermionic defect is supported over one complete transverse-exported dressing block. The matched weak-field gravitational constant then follows from the gauge-invariant bridge
\[
G=\frac{c^2}{8\pi}\frac{\kappa}{\gamma S_\infty},
\]
with the physical geometry--capacity matching and subsequent entropy-unit transformations specified in Appendix~C.5.

\subsection*{D.5 EFT consistency checklist}

The ordinary longitudinal branch is a rewriting of the Einstein constraint sector, so it introduces no independent scalar propagator to which a separate ghost or tachyon test could be applied. Its consistency checks and those of the optional transport completion must be stated separately.
\begin{itemize}
\item \emph{No extra longitudinal degree of freedom.} The Einstein parent supplies the Cauchy data and propagating tensor modes; $\delta S$ is the static scalar-constraint coordinate.
\item \emph{Correct-sign sourcing.} Positive mass produces a capacity deficit and attractive Newtonian response in the reduced action.
\item \emph{Positive static quadratic form.} The Euclidean capacity functional has positive stiffness $\gamma>0$ after the overall gravitational-sign convention is fixed.
\item \emph{Conditional causal transport.} The phenomenological telegrapher completion has finite characteristic speed when $D/\tau_0=c^2$ with $D,\tau_0>0$.
\end{itemize}
These statements do not establish a UV completion or quantize an additional scalar. They show that the controlled static representation is consistent with its Einstein parent and that the separately proposed transport equation is linearly stable in its stated parameter range.

The checklist is intentionally modest. It verifies the reduced Einstein representation and the signs of the separate transport model; it does not establish an independently quantized scalar EFT.

The one place where an explicit formula is worth recording is linear vacuum stability in the time-dependent sector. Writing a small perturbation $\delta s$ about the vacuum branch, the linearized telegrapher equation is
\[
\tau_0\ddot{\delta s}+\dot{\delta s}-D\nabla^2\delta s=0.
\]
For a plane-wave mode $e^{-i\omega t+i\mathbf{k}\cdot\mathbf{x}}$, this gives the dispersion relation
\[
\tau_0\omega^2+i\omega-Dk^2=0.
\]
With $\tau_0>0$ and $D>0$, the corresponding mode frequencies have non-growing time dependence, so the phenomenological transport vacuum is linearly stable.

\subsection*{D.6 Quadratic fluctuations and weak-field stability}

Before imposing the Einstein scalar constraint, a formal covariant quadratic representation is
\[
I_{\mathrm{formal}}^{(2)}[\delta S]
=
-\int d^4x\,\sqrt{-g}\,
\frac{\gamma}{2}\,
g^{\mu\nu}\partial_\mu\delta S\,\partial_\nu\delta S.
\]
After constraint reduction only its spatial Poisson quadratic form remains. The absence of a mass term expresses the unscreened constraint and does not prove an additional propagating boson. The galactic excess of Section~15 belongs to the separate carrier-resolved contact of Appendix~N.8.

Appendix~D provides the technical support layer for the weak-field bridge, Newton limit, electron anchor, substrate length branch, and EFT consistency audit.

\section*{Appendix E: Transport, Cosmology, and the Hubble-Tension Proposal}

Appendix~E collects phenomenological time-dependent and homogeneous extensions of the static capacity variable. These additions are not implied by the Einstein constraint reduction and must recover it in their static limit. The transport relation is fixed only in its preferred branch, while the homogeneous and perturbation sectors remain open. The release cosmology and its fit to current data are given separately in Appendix~O.12.

\subsection*{E.1 Telegrapher equation and causal closure}

The time-dependent deficit field obeys
\[
\tau_0\partial_t^2\delta S+\partial_t\delta S=D\nabla^2\delta S+A\chi,
\qquad
\frac{A}{D}=\frac{\kappa}{\gamma}.
\]
Causality requires
\[
\frac{D}{\tau_0}=c^2.
\]
In the canonical no-new-IR-scale branch,
\[
\tau_0^{-1}=H_0,
\qquad
D=\frac{c^2}{H_0}.
\]

The telegrapher form is the minimal causal completion of the static Poisson sector. It supplies propagation and relaxation while preserving the static weak-field law as its exact late-time limit.

\subsection*{E.2 Static-limit recovery for galaxies}

For a Fourier mode $k$, the telegrapher characteristic equation
\[
\tau_0 s^2+s+Dk^2=0
\]
has the roots
\[
s=-\frac{1}{2\tau_0}\pm i\omega_k,
\qquad
\omega_k\simeq ck
\]
whenever $4\tau_0Dk^2\gg 1$. Galactic wavelengths are far below the critical scale
\[
\lambda_c=\frac{4\pi c}{H_0}\approx 56\ \mathrm{Gpc},
\]
so galactic modes are deeply underdamped. Time-averaging the sourced solution over intervals large compared with $2\pi/\omega_k$ returns the static Poisson branch exactly, and the residual ponderomotive correction scales parametrically as
\[
\frac{\delta F_{\mathrm{pond}}}{F_{\mathrm{static}}}
\sim
e^{-T/(2\tau_0)}
\left(\frac{\omega_{\mathrm{orb}}}{\omega_k}\right)^2
\lesssim 10^{-6}
\]
for representative orbital speeds, with the precise value depending on the averaging interval and system scale. The estimate keeps the correction below the near-stationary weak-field branch but does not supply a universal $10^{-8}$ bound.

\subsection*{E.3 Homogeneous mode and cosmological sourcing}

The cosmological split is
\[
S(x,t)=\overline S(t)+s(x,t),
\]
with $\overline S(t)$ the homogeneous mode and $s(x,t)$ the inhomogeneous weak-field sector. The background capacity is normalized by the apparent horizon,
\[
S_\infty(t)=\pi\frac{R_A(t)^2}{L_*^2},
\qquad
R_A(t)=\frac{c}{\sqrt{H^2+k c^2/a^2}}.
\]
Because the field couples to the trace of the stress-energy tensor, the homogeneous mode is suppressed during radiation domination and turns on near matter--radiation equality.

This timing places the mode where a sound-horizon shift has the most effect: it stays small during radiation domination and becomes relevant near matter--radiation equality.

\subsection*{E.4 Sound-horizon shift and shear lock}

In the conditional cosmological proposal, the trace-sourced homogeneous mode acts as a transient early-energy contribution. It is intended to reduce the sound horizon without rewriting the local static Poisson law. Demonstrating that separation together with the committed component requires the open joint Boltzmann calculation. Section~19 reports the joint fit with the onset fixed at matter--radiation equality and the abundance tie imposed. Moving the onset to $z_c=6800$ improves $\chi^2$ by $0.7$ ($f_{\rm EDE}=0.031$, $H_0=69.75$); with the late universe at $\Lambda$, $z_c=1700$ worsens it by $4.2$. A fit with $\omega_c$ free and the late universe at $\Lambda$ reaches $f_{\rm EDE}=0.13$ and $H_0=72.3$ by raising $\omega_c$ to $0.133$, which the abundance tie excludes. Relative to the $\Lambda$CDM fit, the joint fit has $1.33$ times the abundance of $10^{10}M_\odot$ halos at $z=14.5$ and $1.17$ times at $z=10$; the release fit without the mode gives $1.15$ and $1.08$. The bright end of the ultraviolet luminosity function rises by $0.07$ dex at $z=14.5$. The free-$\omega_c$ fit has $6.4$ times the $\Lambda$CDM halo abundance at $z=14.5$. At $z=1.98$ the abundance of clusters above $1.6\times10^{14}M_\odot$ is $1.22$ times the $\Lambda$CDM value with and without the mode, since the mode has decayed long before cluster assembly.

The homogeneous mode can alter the cosmological background while leaving the coefficients of the local weak-field branch unchanged. A quantitative perturbation calculation remains open.

The transport relation and preferred branch are closed. The homogeneous mode still requires a joint Boltzmann calculation; the release cosmology of Appendix~O.12 is fitted separately.

\section*{Appendix F: Strong-Field Spherical Reduction}

The spherically reduced Einstein action gives the radial invariant
\(q_{\rm geo}=(\nabla R)^2\), the Schwarzschild exterior, and the
marginal surface \(q_{\rm geo}=0\). In normalized static vacuum
this invariant equals the clock capacity \(q_\tau=N^2\).
Inside matter the two observables generally differ, as derived
in N.6. Interpreting their common exterior zero as the physical
end of the substrate EFT requires a domain postulate and a
microscopic boundary action.

\subsection*{F.1 Bounded capacity and the unique lapse map}

The strong-field order parameter is the surviving-capacity fraction
\[
q(x)=\frac{S_{\mathrm{ent}}(x)}{S_\infty}\in[0,1].
\]
This bound follows directly from finite local channel capacity. If the vacuum channel count is finite and $S_{\mathrm{ent}}$ is the logarithmic coarse entropy of the surviving local ensemble, then no physical branch can have either negative capacity or more than the asymptotic vacuum capacity.

In a static exterior, the lapse associated with the asymptotic Killing time is determined by the local surviving capacity. Let
\[
N=f(q).
\]
The conditions are:
\[
f(1)=1,\qquad \lim_{q\to0^+}f(q)=0.
\]
The substrate-level composition axiom is that independent serial capacity losses compose multiplicatively on the lapse:
\[
f(q_1q_2)=f(q_1)f(q_2).
\]
This is the assumption that extends the linear weak-field match to a nonlinear lapse map. With continuity, the positive solutions on $(0,1]$ are $f(q)=q^\alpha$. The multiplicative classification, the positive exponent selected by the zero-capacity limit, and the weak-field normalization below are proved in \path{Entropic/BlackHole.lean}. Expanding near $q=1-\delta$ gives
\[
N=q^\alpha=1-\alpha\delta+O(\delta^2).
\]
The weak-field bridge gives
\[
N=1-\frac{\delta}{2}+O(\delta^2),
\]
so $\alpha=1/2$ and therefore
\[
N=\sqrt q,
\qquad
N^2=q.
\]
Equivalently, if one writes $N^2=F(q)$, the unique continuous multiplicative completion is $F(q)=q$. The nonlinear static lapse rule is therefore fixed by capacity composition and weak-field matching. Because a lapse depends on foliation, this is a static constitutive statement. Its equality to the invariant radial quantity is derived below in static spherical vacuum; N.6 gives the matter-dependent distinction.

\subsection*{F.2 Einstein action reduced to the capacity-adapted invariant}

Take the most general spherically symmetric line element
\[
ds^2=h_{ab}(x)dx^adx^b+R^2(x)d\Omega^2.
\]
Here $x^0=ct$, so the two-dimensional measure is $d^2x=dx^0dr$. The same action written as $dt\,dr$ acquires one additional factor of $c$.
The four-dimensional curvature decomposes as
\[
{}^{(4)}R
=
{}^{(2)}R
+\frac{2}{R^2}\left[1-(\nabla R)^2-2R\Box R\right].
\]
After the angular integral, the Einstein--Hilbert plus GHY action becomes, up to the asymptotic and corner terms retained in $I_\partial^{(2)}$,
\[
I_{\mathrm{sph}}
=
\frac{c^3}{4G}\int d^2x\,\sqrt{-h}
\left[
R^2{}^{(2)}R+2(\nabla R)^2+2
\right]
+I_{\mathrm{matter}}^{(2)}+I_\partial^{(2)}.
\]
The integration by parts that converts $-4R\Box R$ into $+4(\nabla R)^2$ is legitimate only together with the reduced GHY contribution. The boundary term is therefore part of the statement.
For $\partial\mathcal M_4=\partial\mathcal M_2\times S^2$ one has
\[
K^{(4)}=K^{(1)}+\frac{2}{R}n^a\nabla_aR.
\]
The second term cancels the surface term from the integration by parts, leaving
\[
I_\partial^{(2)}
=\varepsilon\,\frac{c^3}{2G}
\int_{\partial\mathcal M_2}dy\,
\sqrt{|\gamma_{(1)}|}\,R^2K^{(1)}
+I_{\mathrm{joint}}+I_{\mathrm{ref}},
\]
with $\varepsilon$ the standard orientation sign.

Define
\[
q_{\mathrm{geo}}\equiv(\nabla R)^2
=h^{ab}\partial_aR\partial_bR.
\]
Varying $R$ and $h^{ab}$ in vacuum gives
\[
R\,{}^{(2)}R-2\Box R=0,
\]
\[
2R\left(h_{ab}\Box R-\nabla_a\nabla_bR\right)
+h_{ab}(q_{\mathrm{geo}}-1)=0.
\]
These equations imply
\[
\nabla_aM_{\mathrm{MS}}=0,
\qquad
M_{\mathrm{MS}}
=\frac{c^2R}{2G}(1-q_{\mathrm{geo}}).
\]
Therefore
\[
q_{\mathrm{geo}}
=1-\frac{2GM_{\mathrm{MS}}}{c^2R}.
\]
The variational status is unambiguous: $q_{\mathrm{geo}}$ is a composite of the two-dimensional metric and the areal-radius dilaton, and its vacuum profile is a first integral. No scalar has been added to the Einstein degrees of freedom.

\paragraph{Audit of the multiplier form.}
If instead one writes an ADM term $\sqrt h\,\lambda(N^2-q)$ while giving $q$ no other bulk dependence, the $q$ equation sets $\lambda=0$. The constraint then only identifies an arbitrary scalar with a foliation-dependent lapse. It does not reproduce the capacity source equation and it has no invariant content away from a specified static slicing. The spherical reduction above supplies the invariant result that the multiplier form lacks.

\subsection*{F.3 Variational status of the $q_{\mathrm{geo}}=0$ boundary}

For $\delta>0$, let $\mathcal B_\delta$ be the timelike level surface $q_{\mathrm{geo}}=\delta$. The exterior Einstein problem is well posed with the standard term
\[
I_{\mathrm{GHY}}[\mathcal B_\delta]
=\varepsilon\,\frac{c^3}{8\pi G}
\int_{\mathcal B_\delta}d^3y\,\sqrt{|\gamma|}\,K
\]
and fixed induced metric on $\mathcal B_\delta$, together with the reference term at infinity and any required joints. The limit $\delta\to0^+$ is null and must be taken using the corresponding null-boundary and joint prescription; a bare timelike GHY expression cannot be evaluated directly at the null surface.

This construction fixes the universal gravitational variation. It does not make
\[
q_{\mathrm{geo}}\big|_{\partial\mathcal M_q}=0
\]
a new Euler--Lagrange boundary condition. The zero is the level set at which the spherical invariant becomes marginal. Declaring that level set to be the end of the physical substrate domain,
\[
\mathcal M_q=\{q_{\mathrm{geo}}>0\},
\]
is the bounded-capacity postulate.

Any additional functional
\[
\Gamma_{\partial q}[\sigma_{AB},\text{boundary channels}]
\]
would describe genuine substrate physics: absorption, partial reflection, relaxation, entropy, or a moving-boundary stress. Its variation may contain a boundary stress and a response conjugate to the limiting capacity, but neither its form nor its spectrum follows from Einstein--Hilbert reduction. The static exterior can be solved without inventing this term; claims about excision, infalling evolution, or finite reflectivity cannot.

\subsection*{F.4 Static spherical vacuum exterior}

In static spherical vacuum, the result follows immediately. On $\mathcal{M}_q$, the bulk equations are the vacuum Einstein equations. The unique asymptotically flat static spherical solution is the Schwarzschild exterior,
\[
ds^2
=
-\left(1-\frac{2GM}{c^2r}\right)c^2dt^2
+
\left(1-\frac{2GM}{c^2r}\right)^{-1}dr^2
+
r^2d\Omega^2,
\]
so the spherical geometric representative of the capacity variable is
\[
q_{\mathrm{geo}}(r)=N^2(r)=1-\frac{2GM}{c^2r},
\qquad
r>r_h,
\]
with
\[
r_h=\frac{2GM}{c^2}.
\]
This also agrees with the weak-field capacity deficit:
\[
\frac{\delta S(r)}{S_\infty}
=
\frac{2GM}{c^2r},
\qquad
q(r)=1-\frac{\delta S(r)}{S_\infty}.
\]

The exterior $r>r_h$ is exactly the standard Schwarzschild exterior, and $q_{\mathrm{geo}}=0$ at $r=r_h$. The geometric invariant becomes negative in the trapped region. Interpreting that sign change as exhaustion of a nonnegative substrate capacity motivates restricting the capacity EFT to $q_{\mathrm{geo}}\geq0$, but this restriction is not implied by the Einstein equations. The absence of a physical classical interior is therefore a conditional substrate claim, not an action-level theorem.

The exterior-domain result does not decide what an infalling observer experiences at the $q_{\mathrm{geo}}=0$ surface. Whether a capacity-exhaustion boundary is smooth, dissipative, anomalous, or absent requires the open functional $\Gamma_{\partial q}$, not merely the exterior Schwarzschild solution.

\paragraph{Geometric identification of the capacity variable.}
The spherical reduction supplies a normalization-independent gradient invariant on the orbit space,
\[
q_{\mathrm{geo}}=h^{ab}\partial_a R\,\partial_b R=|\nabla R|^2,
\]
which in Schwarzschild gives
\[
q_{\mathrm{geo}}=1-\frac{2GM}{c^2r}=N^2,
\]
which can be identified with the substrate fraction $q=S_{\rm ent}/S_\infty$ on the exterior. In spherical dynamical collapse with Misner--Sharp mass $M_{\rm MS}(R,t)$,
\[
q_{\mathrm{geo}}(R,t)=1-\frac{2GM_{\rm MS}(R,t)}{c^2R}.
\]
The marginal-trapped-surface condition $\theta_+\theta_-=0$ coincides with $q_{\mathrm{geo}}=0$ independently of null normalization: positive, zero, and negative $q_{\mathrm{geo}}$ label untrapped, marginal, and trapped spherical regions. Standard collapse in the metric parent can therefore produce the geometric zero. Equating that zero with substrate saturation and refusing the negative branch are the additional bounded-capacity interpretation to be tested by the transport and boundary theory.

\paragraph{Rotating and charged stationary exteriors.}
The baseline metric parent also admits Kerr, Reissner--Nordstr\"om, and Kerr--Newman exteriors when the appropriate conserved gauge sector is present. Their standard horizon entropy
\[
S=\frac{k_BA}{4L_*^2}
\]
retains its form with $A$ the appropriate horizon area, and the Hawking temperature follows from the surface gravity as usual,
\[
T_H=\frac{\hbar \kappa_{\rm sg}}{2\pi k_Bc}.
\]
A nonspherical covariant capacity scalar whose zero selects the outer horizon and whose sign defines a physical domain has not yet been derived. The absence of inner or trapped regions therefore cannot be inferred from the spherical invariant alone.

\subsection*{F.5 Horizon thermodynamics, channel area, and the Immirzi normalization}

Because the exterior geometry is unchanged, semiclassical quantities depending only on the exterior near-horizon saddle are unchanged. The Euclidean continuation is used here as an exterior-saddle calculation: the resulting periodicity depends on regularity of the near-horizon exterior geometry, not on adopting the bounded-domain interpretation past $q_{\mathrm{geo}}=0$. The Euclidean regularity argument therefore gives the standard Hawking temperature \citep{Hawking1975,GibbonsHawking1977},
\[
T_H=\frac{\hbar c^3}{8\pi GM k_B}.
\]
For the GR exterior saddle, the same calculation gives the Bekenstein--Hawking area law \citep{Bekenstein1973,Hawking1975},
\[
S_{\rm BH}=\frac{k_BA}{4L_P^2}.
\]
On the matched branch $L_P(G_*)=L_*$. The factorized $j=3$ completion
supplies a definite map between one capacity cut and physical geometric
area. The semiclassical area law provides the continuum normalization for
that map. \citet{DorauMuch2026} derive the semiclassical Einstein equations from the Araki--Uhlmann relative entropy of coherent excitations across a bifurcate Killing horizon, with this area law as the one gravitational input. The cut map supplies that input from the cell structure.

\paragraph{Projected-area form of the quarter.}
When a coarse spatial horizon is represented as the boundary of a convex body in Euclidean three-space, Cauchy's surface-area formula gives
\citep{TsukermanVeomett2016}
\[
\overline A_{\rm proj}
=\frac{1}{4\pi}\int_{S^2}A_{\rm proj}(\hat u)\,d\Omega
=\frac{A}{4}.
\]
The matched entropy law can therefore be written as
\[
\boxed{
\frac{S_{\rm BH}}{k_B}
=\frac{\overline A_{\rm proj}}{L_*^2}.}
\]
This identity gives the factor $1/4$ a geometric interpretation: it converts total surface area into mean orthogonal projection area. It does not determine the microscopic density. Assigning one nat to each projected area $L_*^2$ is equivalent to the continuum normalization already supplied by the Bekenstein--Hawking law. The formula does not by itself extend this record interpretation to a nonconvex or dynamical horizon.

\paragraph{One cut carries one capacity bit while the geometric boundary face remains.}
If the bounded-capacity domain terminates at $q_{\mathrm{geo}}=0$, the exterior cell has no capacity partner across the outward face.  The missing capacity pairing is the same elementary fermionic exclusion used by the mass--entropy postulate, so the canonical cut increment is
\[
\boxed{\Delta S_f=\ln2.}
\]
The factorized construction of Appendix~Q fixes how this statement is read.  The absent partner removes the \emph{capacity} pairing; it does not erase the geometric boundary datum.  The open boundary still carries the geometric $V_3^{\rm geom}$ representation required for spin-network/spin-foam gluing, while the corresponding capacity $V_3^{\rm cap}$ has lost its partner across the cut.  Thus a horizon puncture is a geometric $j=3$ face decorated by one binary capacity-cut defect, not an additional sevenfold thermodynamic degeneracy.

Accordingly the entropy used here is conditional on the macroscopic boundary geometry,
\[
S_{\rm hor}=S_{\rm cap}(\text{cut}\mid g_{\partial}),
\]
so the geometric spin/intertwiner labels specify the area sector and are not multiplied as an independent entropy factor.  Multiplying a separate isolated-horizon spin degeneracy by the capacity degeneracy would double count the geometry--capacity degree of freedom in this framework; such a construction would be a different microscopic theory, not an extra factor in the present one.

Q.11 realizes this conditional cut in an explicit coherent bond state
with the original local capacity marginal. Its absolute angular
transfer yields the common bare curvature and replica coefficient.
The following algebra states the matched \(G_*\) specialization;
C.8 supplies the separate physical event-volume calibration.

\paragraph{Channel-to-area normalization within the selected completion.}
For a generic LQG surface puncture the kinematical area spectrum is \citep{PerezSpinFoam2013}
\[
A_{\rm LQG}(j)
=8\pi\gamma_{\rm BI}L_*^2\sqrt{j(j+1)}.
\]
The present branch, however, is built from spin coherent faces and maximal coherent fusion.  The area variable that enters the associated Regge geometry is the coherent flux magnitude, linear in $j$; for a coarse face made from parallel small facets, the large-face area is the sum of the small $j_f$ and the coarse spin is $J=\sum_fj_f$ \citep{HanSmallSpin2019}.  Appendix~B therefore fixes the physical coarse area operator on the maximal-fusion ray to
\[
A_{\rm phys}(J)
=8\pi\gamma_{\rm BI}L_*^2J,
\]
or equivalently
\[
A_{\rm phys}(J)
=Z_A(J)A_{\rm LQG}(J),
\qquad
\boxed{Z_A(J)=\frac{J}{\sqrt{J(J+1)}}.}
\]
For $N$ primitive transmitted faces, $J=3N$ and hence
\[
Z_A(N)=\sqrt{\frac{3N}{3N+1}},
\qquad
A_N=24\pi\gamma_{\rm BI}N L_*^2.
\]
This operator renormalization is fixed by additive coherent area; it is not a new coupling.  At one primitive face $Z_A(1)=\sqrt3/2$, and $Z_A\to1$ as the block enters the ordinary large-spin Regge regime.

A patch of $N$ cut capacity channels carries
\[
\frac{S_{\rm cap}}{k_B}=N\ln2.
\]
Requiring this same geometric patch to obey the already-established Einstein/Bekenstein--Hawking normalization gives
\[
N\ln2
=\frac{A_N}{4L_*^2}
=6\pi\gamma_{\rm BI}N.
\]
Therefore
\[
\boxed{
\gamma_{\rm BI}=\frac{\ln2}{6\pi}
=0.0367726000254\ldots
}
\]
and the physical area carried by one microscopic capacity cut is
\[
\boxed{
a_{\rm ch}=\frac{A_N}{N}=4\ln2\,L_*^2,
\qquad
\rho_{\rm ch}=\frac{1}{4\ln2\,L_*^2}.
}
\]
The result is exact conditional on the selected factorized $j=3$ embedding,
the one-cell relational lock, the coherent/Regge additive-area prescription,
and the existing one-bit cut rule. It fixes the embedding normalization. The
Bekenstein--Hawking law and $L_P(G_*)=L_*$ supply the continuum normalization,
so this is not independent evidence for that law. Within the embedding, the
normalization fixes the Immirzi factor.

\paragraph{Bare and coarse area operators.}
If one equated $\ln2$ directly to the unrenormalized $j=3$ Casimir area, one would obtain $\gamma=\ln2/(4\pi\sqrt3)$.  That assignment treats a single bare puncture eigenvalue as the macroscopic area variable and fails to preserve area under the exact maximal-fusion map.  The coherent/Regge branch instead carries the renormalized operator above.  Indeed, for every $N$,
\[
Z_A(N)\,8\pi\gamma_{\rm BI}L_*^2
\sqrt{3N(3N+1)}
=4N\ln2\,L_*^2,
\]
so fine-channel additivity, maximal fusion, the spin-foam Regge area, and the horizon entropy all use one normalization.  The two numerical values are therefore not competing Immirzi predictions: the smaller value belongs to a bare operator convention that is not the coarse observable selected by this branch.

\paragraph{Independent bulk-response identity.}
The graph response supplies a separate normalization check and should not be interpreted as the puncture count.  Since
\[
{\cal G}_{\rm loc}^{ab}=\frac{G_{\mathrm{tet}}(0)}{3}\delta^{ab},
\qquad
G_\perp=\frac23G_{\mathrm{tet}}(0),
\]
the angular response number is
\[
n_{\rm hor}
=4\pi\frac{G_\perp}{\ln2}
=\frac{8\pi G_{\mathrm{tet}}(0)}{3\ln2}.
\]
It is dimensionless and is \emph{not} the literal number of $j=3$ punctures.  Using the cell-normalized capacity baseline
\[
S_\infty^{\rm cell}
=\frac{3\ln2}{32\pi G_{\mathrm{tet}}(0)}
\]
gives the response identity that C.5 imposes as the physical horizon matching
\[
\boxed{n_{\rm hor}S_\infty^{\rm cell}=\frac14.}
\]
The microscopic puncture theorem above and this bulk Green-response identity are two descriptions of the same normalized continuum coefficient, but they count different objects and must not be multiplied together.

\paragraph{Status.}
Within the selected factorized $j=3$ geometric completion and coherent/Regge area prescription, the channel-to-area normalization is therefore closed:
\[
\boxed{
\Delta S_f=\ln2,
\qquad
j_{\rm geom}=3,
\qquad
Z_A(N)=\sqrt{\frac{3N}{3N+1}},
\qquad
\gamma_{\rm BI}=\frac{\ln2}{6\pi},
\qquad
a_{\rm ch}=4\ln2\,L_*^2.
}
\]
What remains in the strong-field sector is dynamical rather than kinematical: the saturation/front evolution, boundary Hamiltonian, relaxation spectrum, reflectivity, stretched-layer corrections, and nonspherical capacity map.

\subsection*{F.6 Absorption, ringdown, and echoes}

The exterior Regge--Wheeler/Zerilli operators are unchanged. With $q_{\mathrm{geo}}=0$ at the marginal surface, the tortoise coordinate $r_*\sim r_h\ln q_{\mathrm{geo}}$ sends the surface to $r_*\to-\infty$, and the near-horizon wave equation reduces to $(c^{-2}\partial_t^2-\partial_{r_*}^2)\psi\simeq0$. If the usual GR future-horizon regularity condition is retained, it selects
\[
\psi\sim e^{-i\omega(t+r_*/c)},
\]
and therefore
\[
\mathcal R=0.
\]
The standard greybody factors and quasinormal spectrum then follow \citep{ReggeWheeler1957,Zerilli1970}. But if the substrate EFT truly terminates at the marginal surface, future-horizon regularity is a boundary choice, not a consequence of the exterior differential operator alone. The open functional $\Gamma_{\partial q}$ must determine whether that choice is correct.

If the microscopic boundary has finite reflectivity or lies on a stretched layer
\[
q=\delta>0,
\]
the region between the exterior potential barrier and that layer behaves as a cavity. A typical echo delay then scales as
\[
\Delta t_{\rm echo}
\sim
\frac{2r_h}{c}|\ln\delta|+\tau_{\rm ch},
\]
where $\tau_{\rm ch}$ is a boundary-channel relaxation time. Thus $\mathcal R=0$ is the GR-matching boundary condition, while $\mathcal R(\omega)$ is a boundary observable to be calculated rather than assumed.

\subsection*{F.7 Dynamical formation as a free-boundary problem}

Standard collapse in the metric parent can form a marginal surface $q_{\mathrm{geo}}=0$. It does not prove that a substrate capacity field saturates there or that the physical evolution terminates. That identification requires a bounded causal transport law for a capacity variable $q_{\mathrm{cap}}$, together with a constitutive relation showing $q_{\mathrm{cap}}=q_{\mathrm{geo}}$ in the regime of overlap. A candidate test system is
\[
\partial_t q_{\mathrm{cap}}+D_iJ^i
=-\Gamma(q_{\mathrm{cap}})\Sigma[T_{\mu\nu}],
\]
\[
\tau_J(\partial_t+\mathcal{L}_v)J^i+J^i
=-D(q_{\mathrm{cap}})D^i q_{\mathrm{cap}}.
\]
Here $J^i$ is the capacity flux, $\Sigma[T_{\mu\nu}]$ is a positive depletion source built from the collapsing stress-energy, $D(q_{\mathrm{cap}})$ is a bounded mobility, $\Gamma(q_{\mathrm{cap}})$ is a bounded depletion rate, and $\tau_J>0$ is a relaxation time. Eliminating $J^i$ gives a telegrapher-type equation with finite characteristic speed
\[
v_{\rm cap}\sim \sqrt{\frac{D_0}{\tau_J}}.
\]
The trial functions
\[
D(q_{\mathrm{cap}})=D_0q_{\mathrm{cap}}(1-q_{\mathrm{cap}}),
\qquad \Gamma(q_{\mathrm{cap}})=\Gamma_0q_{\mathrm{cap}}
\]
give local characteristic speed $\sqrt{D(q_{\rm cap})/\tau_J}\le\sqrt{D_0/(4\tau_J)}$ within $0\le q_{\rm cap}\le1$. For finite $\tau_J$, the flux is independent initial data. In one spatial dimension with $v=0$ and $\Sigma=0$, the exact solution
\[
q_{\rm cap}(t,x)=-b\tau_J(1-e^{-t/\tau_J}),\qquad
J(t,x)=b e^{-t/\tau_J}x,\qquad b>0,
\]
starts at $q_{\rm cap}=0$ and crosses below zero immediately. This solution and the local crossing criterion are kernel checked. An invariant interval therefore requires a compatible flux restriction or a constitutive limiter with its own invariant-domain proof. At its lower boundary $D_iJ^i\le0$ is necessary; at its upper boundary $D_iJ^i+\Gamma(1)\Sigma\ge0$ is necessary. The parabolic limit removes the independent flux, but gives a different propagation law. A causal bounded completion must retain both finite speed and the invariant-domain condition.

If $q_{\mathrm{cap}}$ first reaches zero on a two-surface and the microscopic theory licenses domain termination, that surface becomes a moving boundary
\[
\partial\mathcal{M}_q(t)=\{x\mid q_{\mathrm{cap}}(t,x)=0\}.
\]
The level-set kinematics are fixed by differentiating $q_{\mathrm{cap}}(t,X(t))=0$ along the moving surface:
\[
V_n=-\frac{\partial_t q_{\mathrm{cap}}}{|\nabla q_{\mathrm{cap}}|}
\bigg|_{q_{\mathrm{cap}}\to0^+}.
\]
In spherical symmetry this becomes
\[
\frac{dr_f}{dt}
=-\frac{\partial_t q_{\mathrm{cap}}}{\partial_r q_{\mathrm{cap}}}
\bigg|_{r=r_f(t)}.
\]
During continued infall the exterior should be Vaidya-like with a slowly varying mass parameter, settling to the Schwarzschild exterior after the front stabilizes. This defines a dynamical program whose well-posedness and invariant domain require a compatible completion: the transport coefficients, boundary action, and channel relaxation spectrum must be computed from the graph ensemble or constrained by simulation.

The dynamical system is presumed to preserve the usual covariant conservation of the combined matter-plus-capacity stress-energy, with any local matter depletion balanced by flux, boundary work, or capacity-sector stress. Showing that this conservation structure follows from a graph-derived transport action, rather than imposing it as a constitutive condition, is part of the dynamical closure work.

The concrete closure tests are correspondingly specific. A spherical collapse simulation should show formation of the first $q=0$ surface without overshoot into $q<0$. A coupled matter-plus-capacity run should approach a Vaidya exterior during accretion and a Schwarzschild exterior after settling while satisfying the combined conservation law. Exterior perturbation simulations with an absorbing boundary should reproduce standard Schwarzschild greybody factors and ringdown, while partial-reflectivity runs should produce controlled echo delays. Finally, a microscopic boundary-action calculation or a graph-ensemble Monte Carlo of saturated boundary channels should reproduce the channel-counting rule that yields $n_{\rm hor}S_\infty^{\rm cell}=1/4$. These are not new fit knobs; they are the numerical and microscopic tests that would close the dynamical and boundary sectors.

\subsection*{F.8 Weak-field boundary and closure statement}

In the weak-field Solar-System regime, the metric-only parent yields
\[
\gamma_{\mathrm{PPN}}=\beta_{\mathrm{PPN}}=1,
\]
with the remaining standard PPN coefficients vanishing under the usual assumptions. This follows because the ordinary longitudinal capacity functional is a reduced representation of Einstein gravity. In the carrier-resolved galactic EFT, localized Solar-System sources are absent from the transverse source map, so the leading local solution remains on this Einstein branch. Nonlinear preferred-frame terms still require the covariant projector audit of Appendix~N.8.

The weak-field capacity coordinate ceases to be adequate when
\[
\frac{|\Phi|}{c^2}=O(1),
\qquad
\frac{\delta S}{S_\infty}=O(1),
\]
which is the regime where the spherical invariant $q_{\mathrm{geo}}$ provides the controlled nonlinear description.

Appendix~F therefore closes two precise strong-field statements. The spherically reduced Einstein action makes $q_{\mathrm{geo}}=(\nabla R)^2=1-2GM_{\mathrm{MS}}/(c^2R)$ a composite first integral and returns the Schwarzschild exterior. Separately, the factorized coherent $j=3$ boundary construction fixes the microscopic channel area and Immirzi normalization, while the bulk Green-response product reproduces the $1/4$ imposed by the C.5 matching. The bounded capacity interpretation of $q_{\mathrm{geo}}$, exclusion of the negative branch, and physical boundary dynamics remain additional hypotheses. Rotating and charged solutions belong to the baseline metric parent, but their capacity-variable completion is also open.

\section*{Appendix G: Many-Pasts, Operational Closure, Branch Realization, and the Arrow of Time}

Appendix~G separates the three probability objects used by Many-Pasts: amplitudes for unresolved alternatives, probabilities for decoherent record histories, and conditional probabilities for histories compatible with one present record. These distinctions give the shorthand $P(H\mid P)\propto e^{-D(H,P)}$ a single meaning. The finite operational construction derived in H.11 has the standard decoherent-histories form \citep{Halliwell1994}; Many-Pasts supplies its record-conditioned ontology.

\subsection*{G.1 What is a history of the entanglement network?}

A coarse projective history $h=(\alpha_1,\ldots,\alpha_n)$ is a sequence of alternatives at ordered substrate times. The alternative $\alpha_k$ is represented by a projector $\Pi^{(k)}_{\alpha_k}$. With unitary evolution $U_{k,k-1}$ between times, its class operator is
\[
C_h=\Pi^{(n)}_{\alpha_n}U_{n,n-1}\Pi^{(n-1)}_{\alpha_{n-1}}\cdots
U_{2,1}\Pi^{(1)}_{\alpha_1}U_{1,0}.
\]
This operator retains amplitudes. It is defined before any classical probability is assigned to the individual history.

Given an initial state $\rho_0$, H.11 identifies the decoherence functional
with the Gram kernel of these branch states:
\[
\mathcal D(h,h')=\operatorname{Tr}\!\left(C_h\rho_0C_{h'}^\dagger\right).
\]
A family admits ordinary probabilities when its off-diagonal terms are negligible at the required accuracy,
\[
\mathcal D(h,h')\simeq0\qquad(h\ne h').
\]
The diagonal entries $p(h)=\mathcal D(h,h)$ are then nonnegative and additive under coarse-graining. When alternatives do not decohere, their class operators must be added before the probability is evaluated. For $C_A=\sum_{h\in A}C_h$,
\[
p(A)=\operatorname{Tr}\!\left(C_A\rho_0C_A^\dagger\right),
\]
which retains the interference terms. Many-Pasts places no classical distribution over unresolved fine-grained paths.

\subsection*{G.2 What is the present coarse configuration $P$?}

The present $P$ is a macroscopic record represented by a final projector $\Pi_P$. Let $\mathcal H_P$ be a decoherent family of histories whose final alternatives refine that record. Exhaustiveness and decoherence give
\[
p(P)=\sum_{h\in\mathcal H_P}p(h)
=\operatorname{Tr}(\Pi_P\rho_{\rm now}).
\]
The conditional Many-Pasts measure is
\[
\boxed{
p(h\mid P)=\frac{p(h)}{p(P)},
\qquad h\in\mathcal H_P.
}
\]
The common measure over all possible records is normalized first; conditioning on the realized record comes afterwards. Normalizing a new set of histories separately for each already-selected present would leave the probabilities of the alternative presents undefined.

\subsection*{G.3 What is the distance $D(H,P)$?}

For a decoherent history ending in $P$, define
\[
D(h,P)=-\ln p(h),
\]
with $D=+\infty$ when $p(h)=0$. Then
\[
p(h\mid P)=\frac{e^{-D(h,P)}}{\sum_{h'\in\mathcal H_P}e^{-D(h',P)}}.
\]
The distance notation is shorthand for the diagonal decoherence-functional weight, not a second probability law. The overlap expression $-\ln\operatorname{Tr}(\Pi_P\rho_{H\to\mathrm{now}})$ is recovered when $H$ already denotes a decohered preparation history and only the final record remains unresolved.

\subsection*{G.4 The derived finite Born branch}

A laboratory setting $x$ is represented by the quantum instrument
$\{\mathcal M_{a\mid x}\}_a$ obtained by reversible coupling to a retained
record and then retaining or forgetting its alternatives (H.11). Each map is
completely positive and trace non-increasing, while
$\sum_a\mathcal M_{a\mid x}$ is trace preserving. For an initial state $\rho$,
\[
p(a\mid x)=\operatorname{Tr}\!\left[\mathcal M_{a\mid x}(\rho)\right].
\]
If $C_{h,a\mid x}$ refines the histories ending in record $a$, then
\[
p(h\mid a,x)=
\frac{\operatorname{Tr}(C_{h,a\mid x}\rho C_{h,a\mid x}^\dagger)}{p(a\mid x)}
\]
for a decoherent refinement. The record marginal is the Born probability
because the branch weight is the derived squared Hilbert norm. Reversible
internal equivalence makes the weight depend only on that norm, and finite
additivity under coarse-graining of orthogonal physical records fixes the
dependence to be linear. The remaining premises are ordinary probability
normalization and nonnegativity, together with record completeness for the
decompositions used. Gleason's theorem and generalized-record results remain
independent consistency checks \citep{Gleason1957,Hartle2016}. The global
state, continuum/Fock realization, and which cosmological histories actually
decohere are not fixed by this finite theorem.

\subsection*{G.5 No-signaling}

Let Alice and Bob act locally on $\rho_{AB}$ with instruments $\{\mathcal M^A_{a\mid x}\}_a$ and $\{\mathcal N^B_{b\mid y}\}_b$. Their joint record probability is
\[
p(a,b\mid x,y)=\operatorname{Tr}\!\left[
(\mathcal M^A_{a\mid x}\otimes\mathcal N^B_{b\mid y})(\rho_{AB})
\right].
\]
Summing over Bob's record gives
\begin{align}
p(a\mid x,y)
&=\sum_b p(a,b\mid x,y)\\
&=\operatorname{Tr}\!\left[
(\mathcal M^A_{a\mid x}\otimes\mathcal N^B_y)(\rho_{AB})
\right]\\
&=\operatorname{Tr}\!\left[\mathcal M^A_{a\mid x}(\rho_A)\right]
=p(a\mid x),
\end{align}
where $\mathcal N^B_y=\sum_b\mathcal N^B_{b\mid y}$ is trace preserving. Alice's marginal is independent of $y$, and the same calculation applies to Bob. Global conditioning on a joint present record does not create a controllable signaling channel because the unconditioned local marginals remain those of standard quantum mechanics.

\subsection*{G.6 Branch realization}

The postulate contains one realized present with its records. Its multiplicity lies among the admissible \emph{pasts} of that present, with no forward branching into co-real macroscopic worlds and no collapse event selecting among them. An outcome is specified by the records that obtain in the present, and the Born weights give their statistics. This account of branch realization adds neither many worlds nor a collapse dynamics.

The present record algebra fixes the resolution at which past alternatives are physically distinguished. Let $\{R_i\}_{i\in I}$ be mutually orthogonal fine records that a later physical channel maps to one coarse record $\bar R=\sum_{i\in I}R_i$. For a decoherent joint refinement,
\[
p(h,\bar R)=\sum_{i\in I}p(h,R_i),
\qquad
p(h\mid\bar R)=
\frac{\sum_{i\in I}p(h,R_i)}{\sum_{h'}\sum_{i\in I}p(h',R_i)}.
\]
This is the push-forward of the same joint quantum measure. Record loss needs no additional rule and does not alter the earlier unitary dynamics. It reduces the distinctions represented in the realized present. The history register of H.8 is one microscopic carrier of such distinctions; environmental and macroscopic records are others. The paper does not assume an inaccessible second archive after every physical carrier of a distinction has been erased.

The conditioning also has no forward dynamical role. The class operator $C_h$, the state, and the quantum instruments determine $p(h,P)$ before the ratio $p(h\mid P)$ is formed. A later record permits retrodictive conditioning once it exists; the conditional probability never appears as a force in the preceding evolution. Many-Pasts therefore preserves the causal and no-signaling content of the finite channel structure derived in H.11.

This record-retention reading removes one apparent cosmological alternative under a precise condition. If a state has a trivial macroscopic record algebra with respect to an earlier era, two descriptions that differ only by whether that era occurred before the record-free state are represented by the same present state. A physically meaningful cycle count would require a record surviving the interval. The framework neither supplies such a meta-register nor derives that a record-free cosmological state is reached. The result is an ontological equivalence and makes no claim about cosmic recurrence.

\subsection*{G.7 Arrow of time from conditional typicality}

The arrow of time rests on a typicality statement. Let $h=\{M_t\}_{t_i\le t\le t_0}$ be a decoherent macrohistory conditioned on present records $M_{t_0}$. A coarse Markov description would assign
\[
P(h\mid M_{t_0})\propto
\mu_i(M_{t_i})
\prod_{t_i\le t<t_0}T(M_{t+\Delta t}\mid M_t),
\]
where $\mu_i$ is a boundary measure and $T$ is the conditional macro-transition probability obtained only after summing the microscopic transitions compatible with each pair of macrostates. Multiplying a separate factor $e^{S(M_t)}$ at every time would generally count the same microscopic multiplicity twice. A neutral transition law is also insufficient: conditional counting without a low-entropy boundary condition is dominated by high-entropy pasts and Boltzmann-fluctuation histories. The required theorem must show that $\mu_i$ and $T$ exponentially suppress those histories strongly enough for ordinary entropy-increasing histories to dominate. The boundary measure has a forced support. Many-Pasts posits no meta-time and no unrecorded archive: absent a surviving record, ``a prior era occurred and all traces vanished'' and ``no prior era occurred'' are not distinct physical states, so the fundamental boundary is the minimal-record state, and a heavily recorded state cannot be initial because its record content constitutes a past (Appendix~P.4). $\mu_i$ is therefore defined on the first physically distinguishable record-bearing state, and the arrow is conditional typicality from that anchor; what the quotient removes is the freedom to place the boundary at an arbitrary high-entropy state. Appendix~P.8 supplies the substrate transition law on finite connected networks; the gravitational entropy of the boundary state remains open.

The boundary condition can be stated directly. A \emph{Substrate Past Hypothesis} would require
\[
\operatorname{supp}\rho_0\subseteq\mathcal H_{M_{\rm low}},
\qquad
t_0-t_i\ll \min(t_{\rm relax},t_{\rm rec}),
\]
where $\mathcal H_{M_{\rm low}}$ is a low-entropy macro-subspace and the elapsed time is short compared with equilibration and recurrence. For recorded macrovariables with mixing coarse dynamics, standard large-deviation counting favors entropy growth away from that boundary. Appendix~P.8 supplies the microscopic part: renewal resets the closure register each tick, record-writing subsystems contract relative entropy, and coherent sectors evolve unitarily. Appendix~P.9 derives the support of $\rho_0$ as the minimal-record root of the record order. The low gravitational entropy of that root is not derived. The saturated phase of Section~20 supplies the capacity-sector half of the hypothesis by construction, since the pinned state is a near-zero-entropy configuration of the capacity sector. For the canonical defect class, the visibility--source contrapositive of Appendix~P.3 supplies the matter-sector clause: a state with no elementary massive defects carries no defect records, so matter-sector simplicity at the anchor is a named conditional. The quantitative companion is the well-foundedness bound of Appendix~P.4, which prices a false durable record of content $K$ at $e^{\ln N_{\rm trials}-K}$ with $\ln N_{\rm trials}\sim5.6\times10^{2}$ nats, at the conditional grade it inherits from Section~20. Naming these premises isolates the remaining arrow-of-time problem from the already-closed operational probability branch.

The faithful-history-resolution theorem of Appendix~H does not supply the missing boundary. Its stationary kernel satisfies
\[
p(b)K_*(b,b')=p(b)p(b')=p(b')K_*(b',b),
\]
so it is exactly detailed-balanced and time-reversal symmetric. The quantum dilation is globally reversible as well. Local export distinguishes a present register from a history register. Forward in time each tick consumes a fresh register and exports the previous one, so the reduced map on the present is a channel and contracts relative entropy (P.8); the time-reversed description draws its inputs from correlated history and has no such property. The orientation therefore comes from the export structure together with the minimal-record root (P.9). Appendix~P.10 treats late-time fluctuations: re-entry into the committed phase vanishes in the specified dynamics, and records written in matter obey $P_{\rm false}(K)\le\min(e^{-K},e^{-KE_{\rm nat}/k_BT})$.

\subsection*{G.8 Memoryless dressing and the connection to $L_*$}

The dressing kernel's temporal reading has a spatial twin. For the canonical elementary defect, the rate at which the exported history becomes distinguishable from vacuum history and the scalar source insertion of Appendix~C.5 are the same likelihood ratio, with the one-bit defect statistic sufficient; the identity, its isotropy benchmark $\ln(7/6)$, and its exact decomposition are established in Appendix~P.3. Postulate~III and the weak-field source theorem, separable modules until that identity, consume one shared constant.

The history measure does not imply memorylessness merely because it is written as an exponential. The foundational faithful full-support principle selects it when applied to paths at the fixed admissibility marginal; maximum caliber is only the conventional name for that application. For any stationary dressing kernel with one-time marginal $p_{\eta_*}$,
\[
H(B_{t+1}\mid B_t)
=g_{\mathrm{share,eff}}-I(B_t;B_{t+1})
\le g_{\mathrm{share,eff}}.
\]
Equality holds only when consecutive configurations are independent, which fixes $K(b,b')=p_{\eta_*}(b')$ on the support. More generally, the entropy rate of any stationary process with this one-time marginal is bounded by $g_{\mathrm{share,eff}}$, with equality only when the present is independent of its complete past. The quantum replacement channel transfers the old local information into a dilation register rather than destroying it. Many-Pasts gives that register its history-space interpretation. The finite marked event is a separate statement from this renewal theorem and is supplied by the decorated vertex in H.9.

\subsection*{G.9 Relation to many-worlds, collapse, hidden variables, and decoherence}

Many-Pasts is a form of history-space realism built on the decoherent-histories formalism. It posits one realized record-bearing present and no stochastic collapse term. Its probabilities are the diagonal entries of the finite history Gram kernel, conditioned on that present only after the alternative records have been normalized. The distinctive claim lies in the ontology assigned to this conditional measure and in its proposed substrate realization. On finite retained sectors, H.11 derives the Born norm, record instruments, and no-signaling from the real quadratic carrier, reversible record coupling, and locality. The continuum/Fock realization, cosmological initial state, actual decohering histories, and durable global record capacity remain open.


\subsection*{G.10 Familiar quantum examples: double-slit, EPR/Bell, and measurement}

In a double-slit experiment the two path alternatives remain combined in one class operator while no which-path record exists. The probability of a detection event therefore includes their interference term. A durable which-path record defines a decoherent refinement, after which the path histories admit separate conditional probabilities and the fringe term is suppressed.

In an EPR or Bell experiment the record is joint, so its history measure carries the standard nonclassical correlations. Summing over the remote record invokes a trace-preserving local map and returns a marginal independent of the remote setting, as shown in G.5. Conditioning on the observed joint record does not alter that prior no-signaling marginal.

Measurement creates a durable record and identifies a decoherent family. Operationally the calculation is ordinary quantum mechanics. The additional claim is that the realized present is supported by the conditional ensemble of compatible past histories.

These examples introduce no new laboratory predictions. They show where positive history probabilities are legitimate and where amplitudes must remain combined.

\section*{Appendix H: Microscopic Realization and Coarse-Graining}

This appendix develops a microscopic realization of the scalar stiffness and defect ontology. The weak-field coefficient chain does not depend on this realization. Section~H.9a gives the complete ordinary-gravity derivation in one sequence; the other sections establish the finite action, prove its component results, identify countermodels, and test the geometric embedding.

\subsection*{H.1 GFT condensate realization and coarse-graining}

The candidate microscopic realization is a GFT condensate of bosonic
tetrahedral quanta \citep{Oriti2016,Gielen2013,OritiSindoniWilsonEwing2016} coupled to fermionic defects. Denote the corresponding
fields by $\phi(g_1,\ldots,g_4)$ and $\psi$. In the condensate regime, write
the coarse field as
\[
\sigma(x)=\sqrt{n(x)}\,e^{i\theta(x)}.
\]
The hydrodynamic identity
\[
|\nabla_\mu \sigma|^2
=
\frac{(\nabla_\mu n)^2}{4n}
+n(\nabla_\mu\theta)^2
\]
shows that if
\[
S_{\mathrm{ent}}(x)=S_0+\alpha \ln\!\frac{n(x)}{n_{\mathrm{bg}}},
\]
then the coarse action contains a positive scalar stiffness
\[
\gamma \sim \frac{Z_\sigma n_{\mathrm{bg}}}{2\alpha^2}>0.
\]
The coarse source channel arises from fermionic face exclusion. A localized condensate defect appears macroscopically as matter, and the surrounding reduction of available occupancy becomes the long-wavelength EFT field. The continuum kinetic term and source channel can therefore arise from the proposed substrate realization. Deriving every inhomogeneous continuum coefficient from the underlying kernel remains open.

This condensate argument supplies a microscopic realization compatible with the EFT ontology and sign choices. The coefficient derivation remains the one given earlier.

The condensate picture addresses the emergence of continuum geometry. The remainder of this appendix derives the defect dynamics behind faithful sector resolution from the admissibility ensemble of Part~II together with the Many-Pasts postulate.

\subsection*{H.2 The dressing Hamiltonian and the lightest seven-channel branch}

The faithful sector-resolution relation of Appendix~D.4 was written there as a support-scale identity on the history space $\mathcal H_{\rm hist}=\bigotimes_{m=-3}^{3}\mathcal H_B^{(m)}$, using the refresh kernel $P_{\eta_*}(b,b')=p_{\eta_*}(b')$ on each sector layer. Its dynamical content is carried by an explicit defect Hamiltonian, from which the choice of that kernel can be derived.

The seven face-label channels of the admissibility-closed ensemble define the orthogonal projectors $E_m$, $m=-3,\ldots,3$, of the face algebra $\mathcal A_7$. By Postulate~II the elementary defect is fermionic, so each channel carries an occupation number $n_m\in\{0,1\}$ with fermionic creation and annihilation operators $c_m^\dagger,c_m$ and $n_m=c_m^\dagger c_m$. When channel $m$ is occupied it is dressed by a cloud state described by a density operator $\rho_m$ on the boundary ensemble $\mathcal H_B$. The defect Hamiltonian is
\[
H
=
\underbrace{\sum_{m}\bigl(\varepsilon_0\,n_m-n_m\,F_m(\rho_m)\bigr)}_{\text{single-channel terms}}
+
\underbrace{\sum_{m<m'}V_{mm'}(\rho_m,\rho_{m'})}_{\text{inter-channel coupling}},
\]
with $\varepsilon_0$ the bare cost to occupy a channel, $F_m$ the free energy released by dressing channel $m$ with its cloud, and $V_{mm'}$ the residual coupling between the clouds of distinct channels. Two properties of the ground state of $H$ supply the two ingredients of the sector-resolution relation.

\paragraph{One-pass occupation and the exponent seven.}
A channel is energetically occupied when its dressing gain exceeds the fixed bare cost, $F_m(\rho_m)>\varepsilon_0$. Symmetry makes this inequality the same for all seven channels, so an all-bound configuration is a possible symmetric ground-state branch when the inequality holds. The static Hamiltonian alone cannot identify that branch with the lightest charged particle, because its ordering depends on the undetermined balance between $\varepsilon_0$ and the dressing free energy.

The adopted recurrence mass functional supplies the missing ordering. Let $k=\sum_m n_m$ be the number of resolved channels. Fermionic exclusion gives $0\le k\le7$, and with the admissibility-closed marginal fixed on every occupied layer the support relation generalizes to
\[
\frac{\lambda_k}{L_*}=\frac23\exp\!\left(kg_{\mathrm{share,eff}}-\Delta_k\right),
\qquad
m_k=\frac{3\hbar}{2cL_*}\exp\!\left(\Delta_k-kg_{\mathrm{share,eff}}\right),
\]
where $\Delta_k\ge0$ is the total correlation among the occupied layers. For each $k$, the minimum occurs at $\Delta_k=0$. Between those minima,
\[
\frac{m_{k+1}}{m_k}=e^{-g_{\mathrm{share,eff}}}\simeq5.99\times10^{-4},
\]
so every missing channel raises the minimum mass by $e^{g_{\mathrm{share,eff}}}\simeq1.67\times10^3$. The unique lightest resolved one-bit defect in this support-to-length branch therefore has $k=7$ and $\Delta_7=0$. The factor seven and the product cloud are selected together, without separate one-pass and factorization postulates. This theorem holds at fixed marginals $p_{\eta_*}$. Optimizing the cloud over a different marginal would define a different ultraviolet ensemble.

\paragraph{Remark: finite marked corrections and channel occupation.}
Write $H=g_{\mathrm{share,eff}}$ and $f(x)=-\ln(1-x)$. At a common proper-time cadence, a resolved $k$-channel branch in the additive-generator class has
\[
E_k=\frac{3\hbar}{2\tau_*}A_k f(e^{-S_k}),
\qquad S_k\le kH.
\]
Here $S_k$ is the joint readout entropy and $A_k$ is its positive motif sum. Retaining the unit unmarked motif gives $A_k\ge1$. The selected seven-channel product has $E_7=(3\hbar/2\tau_*)Z_e f(e^{-7H})$. Monotonicity and convexity of $f$, with $f(0)=0$, give
\[
\frac{E_k}{E_7}
\ge\frac{f(e^{-kH})}{Z_e f(e^{-7H})}
\ge\frac{e^{(7-k)H}}{Z_e},\qquad 1\le k<7.
\]
The second inequality follows from $f(cx)\ge c f(x)$ for $c\ge1$ and $cx<1$. Its weakest bound is $e^H/Z_e=1659.8886\ldots$ at $k=6$. Thus the finite marked correction preserves the occupation ordering even for correlated lower-channel clouds, provided their energy uses this same survival prescription and their motif sum retains the unit term. Correlations among seven occupied channels are treated at fixed vertex in H.5.

\paragraph{Factorized cloud and the additive support.}
The dressing dimension multiplies across channels precisely when the joint cloud state is a product, $\rho=\bigotimes_m\rho_m$; the seven contributions $g_{\mathrm{share,eff}}$ then add. The inter-channel term $V_{mm'}$ determines whether the static ground state of $H$ has this form. At fixed marginals, the recurrence mass functional independently selects the seven-channel product as the lightest branch. Each selected channel samples its full ensemble, and the minimizing readout has no inter-layer correlation.

The static Hamiltonian identifies the channel and cloud variables but does not order all of their branches. The next subsections prove the entropy ceiling, the unique memoryless kernel that reaches it, and the joint $k=7$, $\Delta_7=0$ minimum of the adopted recurrence mass functional.

\subsection*{H.3 Slot coupling, layer factorization, and the additivity theorem}

Two distinct correlation structures appear in the closure data, and separating them is essential: conflating them leads to the false conclusion that the cloud cannot factorize.

\paragraph{The closure invariant is a pure pair coupling.}
Expanding $S^2=\bigl(\sum_i m_i\bigr)^2=\Sigma^2+2\sum_{i<j}m_im_j$ in $K^2(b)=48-\tfrac13(S^2-\Sigma^2)$ cancels the self-terms exactly and leaves the identity
\[
K^2(b)=48-\frac23\sum_{i<j}m_im_j.
\]
The admissibility weight $e^{-\eta_*K^2(b)}$ therefore contains only cross-terms between distinct face slots of a single boundary state; it does not factorize over those four slots. This is the exact origin of the residual correlation between face slots on the closed ensemble,
\[
I(\text{slot}_0;\text{slot}_1)=0.1545~\text{nats},
\qquad
\frac{I}{H(\text{slot})}=0.079,
\]
so the four faces of one tetrahedron are about eight percent correlated, a structural feature of the closure invariant.

\paragraph{Slot coupling does not obstruct layer factorization.}
The factorization the support relation requires is over the seven \emph{sector layers} $b^{(-3)},\ldots,b^{(3)}$ of $\mathcal H_{\rm hist}$, each layer being a full boundary state drawn from the entire $1680$-state ensemble. The slot coupling $-\tfrac23 m_im_j$ lives \emph{inside} a single layer's boundary state: it relates the four faces of that one tetrahedron and never couples layer $m$ to layer $m'$. The two structures act on different objects:
\begin{center}
\small
\begin{tabular}{p{0.26\textwidth}p{0.30\textwidth}p{0.30\textwidth}}
\hline
Structure & What it couples & Role \\
\hline
slot coupling $-\tfrac23 m_im_j$ & the four faces \emph{within} one boundary state $b$ & the eight-percent mutual information; lives inside each $\mathcal H_B^{(m)}$ \\
layer coupling $V_{mm'}$ & the seven sector layers $b^{(m)}$ of $\mathcal H_{\rm hist}$ & controls factorization; acts \emph{between} the factors \\
\hline
\end{tabular}
\normalsize
\end{center}
The substrate's intrinsic correlation, the most natural candidate obstruction to factorization, therefore acts at the wrong level to obstruct it: it is internal to a layer, not between layers.

\paragraph{Additivity of the decoherent preparation overlap.}
Appendix~G.3 defines $D(h,P)=-\ln p(h)$ from the diagonal decoherence-functional probability. For the special case in which $H$ already denotes a decohered preparation and only the final record remains unresolved, this reduces to $D(H,P)=-\ln\operatorname{Tr}(\Pi_P\rho_{H\to\mathrm{now}})$. If that preparation state and the resolution projector factorize over channels, $\rho=\bigotimes_m\rho_m$ and $\Pi_P=\bigotimes_m\Pi_m$, the trace factorizes and the logarithm converts the product into a sum,
\[
\operatorname{Tr}(\Pi_P\rho)=\prod_m\operatorname{Tr}(\Pi_m\rho_m)
\quad\Longrightarrow\quad
D=\sum_m\bigl[-\ln\operatorname{Tr}(\Pi_m\rho_m)\bigr]=\sum_m D_m .
\]
The overlap distance is therefore additive in this special factorized preparation. The effective-support statement used downstream follows independently from Shannon subadditivity: a product cloud with seven equal marginals has joint entropy $7g_{\mathrm{share,eff}}$ and perplexity $e^{7g_{\mathrm{share,eff}}}$. Whether the electron's dressing is that product is decided by the two conditions in the next subsections.

\paragraph{The ceiling and the two conditions.}
Subadditivity bounds the support from above. With each channel marginal fixed to the admissibility weight, $H(B_m)=g_{\mathrm{share,eff}}$, the joint entropy of a single readout obeys
\[
H\bigl(B_{-3},\ldots,B_3\bigr)\le\sum_{m=-3}^{3}H(B_m)=7g_{\mathrm{share,eff}},
\]
and the deficit
\[
\Delta=\sum_{m}H(B_m)-H\bigl(B_{-3},\ldots,B_3\bigr)\ge 0
\]
measures the total correlation among the seven channels. The full support $\dim_{\rm eff}=e^{7g_{\mathrm{share,eff}}}$ is reached precisely when each channel carries its full entropy $g_{\mathrm{share,eff}}$ and the channels are mutually independent, $\Delta=0$. The first is a condition on each layer in substrate time; the second is a condition on the seven layers at a single readout. The next two subsections establish the unique faithful-resolution kernel and the minimal-mass factorized branch, while separating those results from the still-open microscopic implementation.

\subsection*{H.4 Maximum caliber: complete renewal is uniquely selected}

The closure calculation fixes the one-time marginal $p_{\eta_*}$. Applying faithful full-support resolution to entire paths maximizes their entropy rate, a condition conventionally called maximum caliber \citep{Jaynes1980}. This is the temporal form of the existing foundational requirement. For an arbitrary stationary process, its entropy rate is
\[
h_\mu
:=
\lim_{n\to\infty}H(B_0\mid B_{-1},\ldots,B_{-n})
\le H(B_0)=g_{\mathrm{share,eff}}.
\]
Equality holds only when $B_0$ is independent of its complete past. Stationarity then makes the process independent and identically distributed, with every conditional distribution equal to $p_{\eta_*}$. Maximum path entropy therefore selects complete renewal uniquely. Finite capacity alone does not imply renewal. The selection follows from faithful full-support resolution, whose temporal formulation is maximum caliber.

The one-parameter family below illustrates the cost of retained memory:
\[
K_a(b,b')=a\,\delta(b,b')+(1-a)\,p_{\eta_*}(b'),
\qquad a\in[0,1],
\]
which repeats the current state with probability $a$ and otherwise redraws from the stationary weight. Every member has $p_{\eta_*}$ as its stationary distribution and the same single-time marginal, so the equilibrium ensemble cannot say which one governs the dressing. The per-channel conditional entropy $H(b'\mid b)=\sum_b p_{\eta_*}(b)\,H\bigl(K_a(b,\cdot)\bigr)$, evaluated on the exact $1680$-state ensemble, nonetheless slides with the memory $a$:
\begin{center}
\small
\begin{tabular}{ccc}
\hline
$a$ (memory) & per-channel $H(b'\mid b)$ & status \\
\hline
$0.0$ (refresh) & $7.41980=g_{\mathrm{share,eff}}$ & full entropy \\
$0.1$ & $6.99953$ & reduced \\
$0.3$ & $5.80153$ & reduced \\
$0.5$ & $4.40051$ & reduced \\
$0.9$ & $1.06642$ & reduced \\
\hline
\end{tabular}
\normalsize
\end{center}
Only the memoryless endpoint $a=0$, the replacement kernel $K_*(b,b')=p_{\eta_*}(b')$ of Appendix~D.4, returns the full $g_{\mathrm{share,eff}}$. For any stationary Markov kernel with that marginal,
\[
H(B_{t+1}\mid B_t)=H(B_{t+1})-I(B_t;B_{t+1})\le g_{\mathrm{share,eff}},
\]
and equality holds if and only if $I(B_t;B_{t+1})=0$. The Markov corollary is therefore
\[
\boxed{K_*(b,b')=p_{\eta_*}(b')}.
\]
Single-label local moves conserve slot ordering, fracturing each parity copy into twenty-four sectors of thirty-five states, so no such local kernel reaches the full ensemble at any parameter value (Appendix~D.4). The native-vertex construction of H.8 implements the selected nonlocal-in-label-space replacement as one operation on the complete local tetrahedral register.

The lightest-defect branch selects the same endpoint independently within the adopted support map. If a per-pass temporal mutual information $I_t$ reduces the fresh entropy from $g_{\mathrm{share,eff}}$ to $g_{\mathrm{share,eff}}-I_t$, then at fixed $L_*$
\[
\lambda(I_t)=\lambda(0)e^{-I_t},
\qquad
m(I_t)=m(0)e^{I_t}.
\]
Every retained temporal correlation raises the recurrence mass, so electron lightness again selects $I_t=0$. Faithful full-support resolution selects the vacuum history process; mass minimization independently selects the same process in the defect sector.

A finite irreducible detailed-balance generator with stationary weight $p_{\eta_*}$ approaches the replacement kernel. Transpositions of occupied labels preserve the unordered occupied set, so an irreducible generator must also connect different occupied sets. An exact dimensionless example is
\[
Q(b,b')=p_{\eta_*}(b')-\delta_{bb'},
\qquad
e^{\tau Q}(b,b')=e^{-\tau}\delta_{bb'}+(1-e^{-\tau})p_{\eta_*}(b').
\]
Its off-diagonal rates obey detailed balance, its rows sum to zero, and its nonstationary eigenvalues are $-1$. Thus $e^{\tau Q}$ converges to replacement as $\tau\to\infty$. Faithful full-support resolution selects this endpoint and leaves the approach to it unspecified. H.8 derives the quantum replacement channel and reversible update; H.9 supplies the fresh amplitude, marked event, and update interaction. The physical source of history capacity (P.9) and selection of a stable interacting geometric phase remain open; Q.11b supplies a connected finite encoding of a selected covariant host.

\subsection*{H.5 Independent channels: the electron as the lightest defect}

The second condition is that the occupied channels factorize at one readout, so that the deficit $\Delta_k$ of H.3 vanishes. Appendix~H.2 already showed that electron lightness selects both the maximum occupation $k=7$ and $\Delta_7=0$ once the support-to-length dictionary is adopted. The calculation here isolates the factorization part of that joint minimum.

Keep the deficit explicit in the baseline support relation. With each channel at its full entropy, the joint entropy is $7g_{\mathrm{share,eff}}-\Delta$, and
\[
\frac{\lambda}{L_*^{(0)}}=\frac23\,e^{\,7g_{\mathrm{share,eff}}-\Delta}.
\]
At a fixed substrate scale $L_*$, a one-bit charged defect whose dressing carries correlation $\Delta$ has spatial support $\lambda\propto e^{-\Delta}$, and through $m=\hbar/(c\lambda)$ a mass
\[
m\propto e^{\Delta}.
\]
Correlation therefore contracts the support and raises the recurrence mass. Because $\Delta\ge0$, the lightest one-bit charged defect is the product dressing $\Delta=0$; correlated dressings describe heavier branches within the same recurrence/support map. Independence is forced by mass minimization in the theory's own definition of the electron anchor.

\paragraph{Correlated clouds at fixed marked vertex.}
The finite correction preserves the product minimum within the specified vertex. Hold the seven marginals at $p_{\eta_*}$, all occupations at $n_m=1$, and the internal routing of H.9 fixed. The marked action then contains only the constants $g_{\mathrm{share,eff}}$, $\eta_*$, and the occupied pair projectors. Its auxiliary records retain their specified contractions. A change in the joint cloud readout law $P$ changes its entropy, while the internal motif operator is unchanged. Consequently
\[
Z_e[P]=(1+\zeta_*)(1+7\zeta_*^2)=Z_e.
\]
With $H=g_{\mathrm{share,eff}}$, the fixed-marginal KL identity gives
\[
\Delta=7H-H(P)=D\!\left(P\middle\|p_{\eta_*}^{\otimes7}\right)\ge0,
\qquad r(P)=e^{-H(P)}=r e^\Delta.
\]
For the selected additive generator,
\[
\boxed{\frac{E(P)}{E(p_{\eta_*}^{\otimes7})}
=\frac{f(r e^\Delta)}{f(r)}\ge e^\Delta,
\qquad f(x)=-\ln(1-x).}
\]
Equality with the minimum requires $\Delta=0$, hence $P=p_{\eta_*}^{\otimes7}$. This establishes the unique product minimum over the fixed-marginal readout laws of the same marked vertex. A model that retains only $Z_\mu$ has the weaker lower bound $e^\Delta/(1+7\zeta_*^2)$ and leaves $\Delta<\ln(1+7\zeta_*^2)=1.83741\times10^{-4}$ nats unexcluded. Such a model changes the specified electron return block. A coupling between the cloud correlations and the auxiliary records would likewise require a different vertex.

Faithful full support fixes memoryless refresh. The occupation bound in H.2 and the fixed-vertex result above select $k=7$ and $\Delta_7=0$ within the marked-transfer prescription. H.6 gives the baseline survival scale and H.9 its finite dressing:
\[
L_*=-\frac32\,Z_e\lambda_e\ln\!\left(1-e^{-7g_{\mathrm{share,eff}}}\right),
\qquad
G_*=\frac{c^3L_*^2}{\hbar},
\]
with no gravitational quantity used as input.

The decorated action realizes the nonlocal replacement, fresh state, and charged marked event. In the factorized completion those operations live on the finite capacity/record factor and need not be re-derived from the geometric magnetic indices. A specific many-body GFT realization must still determine how the decoration couples to inhomogeneous matter and transverse response; it need not re-select memorylessness, the clock conversion, or $\Delta=0$.

\subsection*{H.6 Mass from the positive survival transfer operator}

Faithful full-support resolution fixes the replacement process on history space but not a dimensional duration. The duration comes from the paper's existing electron anchor once the already-declared effective-support event is represented in the transfer space. The electron supplies the dimensional anchor; the determinant-survival and marked-generator prescriptions specify its dynamical readout.

\paragraph{The marked seven-channel transfer.}
On the factorized lightest branch, likelihood multiplication on each renewed cloud has determinant $e^{-g_{\mathrm{share,eff}}}$. The product determinant is
\[
r:=\Delta_{\tau_p^{\otimes7}}(\mathsf R^{\otimes7})
=e^{-7g_{\mathrm{share,eff}}}.
\]
This is the scalar charged determinant line of the decorated transfer, not the probability of an arbitrarily chosen seven-layer microstate. The selected one-bit survival action assigns the positive no-loop transfer
\[
T_{\rm surv}^{(0)}=1-r,
\]
equivalently the one-dimensional Grassmann determinant $\int d\bar c\,dc\,e^{-\bar c(1-r)c}=1-r$. One may therefore define
\[
H_{\rm surv}^{(0)}=-\frac{\hbar}{\tau_*^{(0)}}\ln T_{\rm surv}^{(0)},
\]
with exact raw gap
\[
\boxed{E_{\rm raw}^{(0)}=-\frac{\hbar}{\tau_*^{(0)}}\ln(1-r)}.
\]
The unit mixing gap on the renewed probability space is a different operator and is not identified with this exponentially small charged gap.

\paragraph{Electron calibration and exact cell length.}
Before the finite marked response is included, the transverse export gives $H_e^{(0)}=(3/2)H_{\rm surv}^{(0)}$. The baseline electron calibration is
\[
m_ec^2
=-\frac{3\hbar}{2\tau_*^{(0)}}\ln(1-r).
\]
It fixes
\[
\boxed{
\tau_*^{(0)}=-\frac32\tau_e\ln(1-r),
\qquad
L_*^{(0)}=c\tau_*^{(0)}=-\frac32\lambda_e\ln(1-r)
}.
\]
For $r=e^{-7g_{\mathrm{share,eff}}}$,
\[
L_*^{(0)}=-\frac32\lambda_e\ln\!\left(1-e^{-7g_{\mathrm{share,eff}}}\right)
=\frac32\lambda_e e^{-7g_{\mathrm{share,eff}}}
\left[1+\frac12e^{-7g_{\mathrm{share,eff}}}+O(e^{-14g_{\mathrm{share,eff}}})\right].
\]
The relative logarithmic correction to the leading expression is $1.4\times10^{-23}$. The baseline induced scale is
\[
G_*^{(0)}
=\frac94\frac{\hbar c}{m_e^2}
\left[\ln\!\left(1-e^{-7g_{\mathrm{share,eff}}}\right)\right]^2,
\]
whose leading form is proportional to $e^{-14g_{\mathrm{share,eff}}}$. H.9 derives the finite response and gives $\tau_*=Z_e\tau_*^{(0)}$, $L_*=Z_eL_*^{(0)}$, and $G_*=Z_e^2G_*^{(0)}$.

\paragraph{Why this supplies the phase bridge.}
A positive Euclidean transfer operator defines an energy spectrum. Analytic continuation supplies the Lorentzian phase frequency of the same charged mode. An independent identification of an integer renewal count with a condensate winding is unnecessary. H.9 realizes the determinant recurrence and marked response in the decorated charged vertex.

The temporal recurrence gives $\tau_e=N_{\rm eff}\tau_*$. Defining the associated causal length by $L_*=c\tau_*$ then gives $\lambda_e=N_{\rm eff}L_*$. This is the scale convention used in the static graph calculation. It does not identify one graph displacement with $L_*$. The physical displacement per update must follow from the coupled spatial transfer, and a geometric tetrahedron may carry a different length fixed by its embedding.

\paragraph{Internal locality and the native vertex.}
The tested single-label move cannot implement replacement. Injectivity-preserving shifts fracture the $1680$ states into
\[
48\ \text{components}=4!\ \text{orderings}\times2\ \text{orientations},
\qquad
35\ \text{states per component}.
\]
H.8 constructs a one-layer reversible dilation when the complete tetrahedral register is the native local vertex. This establishes circuit depth one, not the physical duration; the duration is fixed spectrally by the electron equation above. H.9 decorates that gate with the fresh amplitude and marked charged event.

\subsection*{H.7 Conditional dressing functionals and the condensate-spectrum lemma}

The dressing Hamiltonian named in Section~23,
\[
H=\sum_m\big(\varepsilon_0\,n_m-n_m F_m(\rho_m)\big)+\sum_{m<m'}V_{mm'}(\rho_m,\rho_{m'}),
\]
carries three named functionals $\varepsilon_0,F_m,V$. A proposed mean-field condensate gives a conditional reconstruction of these functionals. The required condensate has not yet been derived from a specified GFT action.

\paragraph{Mean-field reduction.}
In a GFT condensate the boundary-data order parameter is $\langle\hat\varphi(b)\rangle=\sigma(b)$. The admissibility calculation fixes a normalized modulus profile $p_{\eta_*}(b)$, so a candidate condensate may be written
\[
|\sigma(b)|^2=\bar n\,p_{\eta_*}(b).
\]
This fixes neither the overall density $\bar n$ nor any of the phases. A general 1680-component condensate has a common phase and relative-phase directions, and their physical status is determined by the kinetic and interaction kernels of the specified GFT action. Likewise, $\eta_*^{-1}$ may be used as a formal closure-temperature parameter, but the identity
\[
\eta_*\langle K^2\rangle_{\eta_*}=\frac32
\]
is the stationarity condition generated by the three-component determinant weight, not a literal equipartition theorem for the bounded discrete spectrum. Conditional on a condensate realization, $T_*=1/\eta_*=33.48$ may nevertheless be read as the dimensionless substrate closure temperature, with $\langle K^2\rangle_{\eta_*}=\tfrac32 T_*=50.22$; this is an interpretation of the same stationary point and introduces no new number.

\paragraph{The refresh projector.}
The memoryless transition matrix is $K_*(b,b')=p_{\eta_*}(b')$. Let $D_p=\operatorname{diag}(p_{\eta_*})$ and define the detailed-balance symmetrization
\[
\widetilde K_*=D_p^{1/2}K_*D_p^{-1/2}.
\]
With $v_b=\sqrt{p_{\eta_*}(b)}$, one obtains
\[
\widetilde K_*=|v\rangle\langle v|\equiv P,
\qquad
\widetilde L_{\rm mix}=I-\widetilde K_*=P_\perp.
\]
The classical replacement process has one stationary direction in the weighted probability space and removes every orthogonal probability mode after one discrete update. The vector $|v\rangle$ belongs to the symmetrized classical transfer representation; it is not automatically the quantum state of the cell. It does not follow that $P_\perp$ is a Hermitian mass matrix in a closed Lorentzian GFT. The Markov update, a quantum density operator, and the coherent fluctuation Hessian are different objects.

\paragraph{Conditional refresh--Hessian bridge.}
The minimal reversible bridge uses the positive refresh Dirichlet form as the internal Euclidean quadratic cost,
\[
\mathcal K_{\rm int}^{(2)}=\Delta_{\rm ref}P_\perp,
\qquad \Delta_{\rm ref}>0.
\]
A fixed-$j$ Euclidean effective action realizing this bridge is
\[
\Gamma_E[\sigma]=\int d^4x_E\left[
Z\,\partial_\mu\sigma^\dagger\partial^\mu\sigma
-\mu^2\sigma^\dagger\sigma
+\frac{\lambda}{2}(\sigma^\dagger\sigma)^2
+\Delta_{\rm ref}\sigma^\dagger P_\perp\sigma
\right].
\]
It has $\sigma_0=\sqrt{\bar n}\,e^{i\theta_0}v$ with $\bar n=\mu^2/\lambda$. With canonically normalized real fluctuations, its Euclidean quadratic operator contains a radial singlet with $m_h^2=2\mu^2/Z$, a common phase with $m_\pi^2=0$ when number $\mathrm U(1)$ is exact, and $1679$ complex relative modes with $m_\perp^2=\Delta_{\rm ref}/Z$. Lorentzian propagation requires the usual analytic continuation and the global $(-,+,+,+)$ convention of Appendix~A. The exclusion source couples at linear order to the density fluctuation $h$ because $\delta(\sigma^\dagger\sigma)=2\sqrt{\bar n}\,h+\cdots$. In the constant-phase realization displayed here, the common phase is derivatively coupled and has no linear zero-frequency overlap with that static density source. This is not a structural decoupling theorem for a relational GFT condensate. If $\sigma=\sqrt n\,e^{i\theta}$ and the background carries relational phase current $\partial_\chi\theta_0\ne0$, the Madelung term contains
\[
n(\partial_\chi\theta)^2
\supset
2(\partial_\chi\theta_0)\,\delta n\,\partial_\chi\delta\theta,
\]
so density and phase can mix. The complete coherent Hessian and defect/source vertex must determine whether a phase mode is gapped, remains light but source-orthogonal, or mediates an additional response. Simplicial interactions may break number $\mathrm U(1)$ and gap it.

An open-system completion may instead implement renewal as the quantum replacement channel
\[
\mathcal E_*(\rho)=\rho_*\operatorname{Tr}\rho.
\]
The classical closure ensemble fixes the diagonal probabilities of $\rho_*$ in the boundary basis. H.9 adopts the diagonal maximum-entropy completion and constructs its reversible update. A dissipative representation with Lindblad generator $\dot\rho=\Gamma_{\rm ref}[\mathcal E_*(\rho)-\rho]$ gives every traceless input perturbation a relaxation gap. A geometric condensate embedding must determine how this transfer description appears in its coherent fluctuation spectrum.

\paragraph{Selection, mixing, and marked survival are distinct operators.}
The combinatorial Hamiltonian selects the admissibility profile. The mixing generator $\widetilde L_{\rm mix}=P_\perp$ controls decay toward that profile. The exponentially small electron scale comes from a third object: the scalar determinant transfer of the seven-channel likelihood operator,
\[
r:=\Delta_{\tau_p^{\otimes7}}(\mathsf R^{\otimes7})
=e^{-7g_{\mathrm{share,eff}}},
\]
whose one-bit survival transfer is $T_{\rm surv}^{(0)}=1-r$. Over a baseline proper-time step $\tau_*^{(0)}$ its exact spectral gap is
\[
\mathrm{gap}(H_{\rm surv}^{(0)})
=-\frac{1}{\tau_*^{(0)}}\ln(1-r),
\qquad
m_ec^2=\frac32\hbar\,\mathrm{gap}(H_{\rm surv}^{(0)}).
\]
Solving gives $\tau_*^{(0)}=-(3/2)\tau_e\ln(1-r)$ and $L_*^{(0)}=c\tau_*^{(0)}$. H.9 enumerates $Z_e$ and uses its additive-generator prescription to obtain $L_*=Z_eL_*^{(0)}$. The unit mixing gap and charged survival gap are not identified.

\paragraph{$\varepsilon_0$, the bare cost.}
$\varepsilon_0$ is the condensate chemical potential $\mu_0=\partial E/\partial N$, the energy to commit one cell against the mean field, of order $E_*=\hbar c/L_*$. Energies in $H$ are measured in units of $E_*$, and $T_*=1/\eta_*$ is dimensionless. The rest mass is read from the positive marked-survival operator $H_{\rm surv}$, not from the static spectrum of $H$ or the unit mixing gap of $L_{\rm mix}$.

\paragraph{$F_m$, the dressing free energy.}
A bound channel carries a cloud $\rho_m$ over the boundary ensemble. The proposed released free energy is $F_m(\rho_m)=T_*S[\rho_m]$. At fixed mean closure its maximum is the Gibbs state $p_{\eta_*}$ with entropy $g_{\mathrm{share,eff}}$. On the selected lightest branch the seven clouds carry total entropy $7g_{\mathrm{share,eff}}$. The faithful state-weighted determinant of their likelihood operator is exactly $e^{-7g_{\mathrm{share,eff}}}$. This is a quenched multiplicative transfer rate, not a raw single-draw Gibbs probability.

\paragraph{Two energies, kept distinct.}
The separation cost and rest energy are different quantities. The separation cost is the entropic free energy required to dissolve a defect into the substrate; in this construction it is large and linear in $g_{\mathrm{share,eff}}$. The rest energy is $m_ec^2=\hbar\Gamma_e$, with $\Gamma_e$ set by the rare joint recurrence and therefore exponential in $-7g_{\mathrm{share,eff}}$. Reading the separation cost as the mass would place the electron near tens of $E_*$ and remove the required suppression.

\paragraph{$V$, the inter-channel coupling.}
The permutation-invariant singlet
\[
|s\rangle=\frac{1}{\sqrt7}\sum_m|m\rangle
\]
is the natural channel through which a scalar capacity fluctuation couples uniformly to the seven sectors. Within the minimal constant-phase refresh--Hessian completion displayed above, the radial density fluctuation is the only linearly source-coupled singlet at static order and the relative sector is gapped. Integrating out that density mode gives
\[
V_{mm'}=v\,\langle m|s\rangle\langle s|m'\rangle=\frac{v}{7},
\]
a rank-one outer product. The per-incidence overlap is $1/7$; the separate orientation doubling gives the $2/7$ used in Appendix~I.1. The implication is
\[
\text{one source-coupled density singlet}\quad\Longrightarrow\quad \operatorname{rank}(V)=1.
\]
This result is derived inside the minimal bridge, not from the closure ensemble alone. A microscopic interaction that mixes $P$ and $P_\perp$, closes the relative gap, or sources the common phase would add exchange channels and deform the uniform structure.

\paragraph{Positive-transfer phase and the exchanged mode.}
The projector calculation separates the static source from the transfer clock. The static exclusion source couples to the radial density singlet. Independently, the positive Euclidean survival operator defines the electron energy, and ordinary analytic continuation gives the Lorentzian phase frequency $\omega_e=E_e/\hbar$ of that same charged mode. The decorated vertex supplies the finite charged routing, and the additive-generator prescription specifies its energy readout. The remaining transverse question concerns the microscopic origin of the carrier source projector and metric contact: the interaction must preserve the established projection and must not create an additional sourced scalar pole beyond the Einstein constraint response.

\paragraph{The remaining microscopic lemma.}
The effective projector model narrows the missing calculation without requiring geometry--capacity mixing to vanish.  For a specified fixed-$j$ simplicial GFT action admitting $\sigma_b=\sqrt{\bar n\,p_{\eta_*}(b)}e^{i\theta}$, compute its complete Bogoliubov or Schwinger--Keldysh kernel.  After the within-tick evolution $\mathcal U^\dagger$, exact renewal gives
\[
\mathcal T^\dagger=P\mathcal U^\dagger
=\begin{pmatrix}A&B\\0&0\end{pmatrix}
\quad\text{on }P\mathcal K\oplus P_\perp\mathcal K,
\]
so the replaceable sector cannot carry an autonomous nonzero transfer pole even when $B\ne0$; H.9a proves the spectral statement.  The remaining calculation is instead to show that the retained block $A$ has only the Einstein constraint and transverse--traceless gravitational support at low energy, while any additional persistent relative mode is gapped, relaxational, gauge/constraint, or source-orthogonal.  The decorated transfer vertex of H.9 fixes the charged recurrence and its edge Hessian. A geometric GFT must still determine whether its retained condensate realization produces a coherent mass, a relaxation gap, or both, and verify that no additional source-coupled light field survives. The closure weight fixes $P$ exactly, while the refresh--Hessian bridge for the geometric condensate remains a selected effective completion.

\subsection*{H.8 Faithful history resolution and reversible Many-Pasts renewal}

Four distinct statements fall under ``refresh'': the classical history process, the quantum channel, its reversible implementation, and the conversion of a positive charged spectrum into proper time. They are treated separately below. Section~H.9 supplies the charged marked vertex. Maximum-caliber terminology and the clock construction add no foundational premise.

\paragraph{Finite-record entropy identity.}
For a joint law $P_n$ of $n$ records with fixed marginal $p_*$ on each record,
\[
D_{\rm KL}(P_n\Vert p_*^{\otimes n})
=nH(p_*)-H(P_n)\ge0.
\]
Equality holds precisely at $P_n=p_*^{\otimes n}$. Faithful resolution, defined here as maximal joint entropy subject to the fixed marginals, therefore selects independent layer records. Applied to every finite temporal block it selects independent successive records as well. The stationary entropy-rate statement below gives the same endpoint.

\paragraph{Maximum path entropy.}
Let $\mathcal B$ be the finite admissible boundary space and let $p_*(b)=p_{\eta_*}(b)$ be its fixed stationary marginal. For any stationary process $\{B_n\}$ on $\mathcal B$,
\[
h_\mu
=
\lim_{N\to\infty}\frac1N H(B_1,\ldots,B_N)
=
H(B_0\mid B_{-1},B_{-2},\ldots)
\le H(p_*).
\]
Equality holds exactly when $B_0$ is independent of its complete past. Applying the same statement after every time translation factorizes every finite joint distribution,
\[
P(B_1=b_1,\ldots,B_N=b_N)=\prod_{n=1}^N p_*(b_n).
\]
The unique maximum-caliber history process is therefore i.i.d. For a Markov representation its transition kernel is
\[
\boxed{K_*(b,b')=p_*(b')}.
\]
Here ``maximum caliber'' is terminology for faithful full-support resolution on history space \citep{Jaynes1980}. It is not a sixth theory-defining input. The foundational requirement already says that an admissible resolution carries the full available entropy compatible with its fixed data; applying that same requirement to a path gives the equality case above. Finite capacity alone would not imply renewal, but the paper's stronger faithful full-support premise does.

The selected process carries no thermodynamic orientation. It obeys
\[
p_*(b)K_*(b,b')=p_*(b)p_*(b')=p_*(b')K_*(b',b),
\]
so its stationary path measure is exactly invariant under time reversal. Faithful history resolution derives local renewal, not a low-entropy past and not a physical cadence.

\paragraph{Proper time from the positive marked-transfer spectrum.}
The renewal theorem fixes the discrete process but says nothing about seconds per update. Proper time comes from the determinant survival operator of H.6,
\[
r=\Delta_{\tau_p^{\otimes7}}(\mathsf R^{\otimes7})=e^{-7g_{\mathrm{share,eff}}},
\qquad
T_{\rm surv}^{(0)}=1-r.
\]
Defining the baseline transfer over $\tau_*^{(0)}$ gives
\[
H_{\rm surv}^{(0)}:=-\frac{\hbar}{\tau_*^{(0)}}\ln T_{\rm surv}^{(0)},
\qquad
E_{\rm raw}^{(0)}=-\frac{\hbar}{\tau_*^{(0)}}\ln(1-r)
\]
on its nonzero support. The baseline transverse identification $m_ec^2=(3/2)E_{\rm raw}^{(0)}$ therefore yields
\[
\boxed{\tau_*^{(0)}=-\frac32\tau_e\ln(1-r)},
\qquad
\boxed{L_*^{(0)}=c\tau_*^{(0)}=-\frac32\lambda_e\ln(1-r)}.
\]
The leading form is $L_*^{(0)}=(3/2)\lambda_e r[1+O(r)]$. H.9 supplies the finite marked factor $Z_e$ and hence the physical $L_*=Z_eL_*^{(0)}$. No update diameter, throughput maximization, or geometric conversion factor enters. The complete tetrahedron being a native gate is a locality and circuit-depth statement only; its physical embedding size is a separate consistency question.

A positive Euclidean transfer spectrum also supplies, by analytic continuation, the Lorentzian phase frequency. The clock and phase use the same charged spectral mode and require no independent winding postulate. H.9 derives the recurrence as the state-weighted determinant and implements its finite marked dressing.

\paragraph{Quantum replacement theorem.}
Let $\mathcal H_A$ be the replaceable closure-register Hilbert space and $\mathcal E:A\to A'$ a trace-preserving quantum channel. It does not include the retained position and internal fiber of a persistent charged defect. Complete local renewal means that the output contains no information about the input, including information visible through an arbitrary reference $R$:
\[
(\operatorname{id}_R\otimes\mathcal E)(\rho_{RA})
=
\rho_R\otimes\rho_*
\qquad
\text{for every }\rho_{RA}.
\]
Taking product inputs first shows that every normalized input has the same output $\rho_*$. Linearity then gives the unique channel on all operators,
\[
\boxed{\mathcal E_*(X)=\rho_*\operatorname{Tr}X}.
\]
Every traceless perturbation is annihilated after one update. As a superoperator, $\mathcal E_*$ has one stationary direction and zero on the traceless operator subspace. The closure ensemble fixes $\langle b|\rho_*|b\rangle=p_*(b)$ in the boundary basis. H.9 selects the diagonal maximum-entropy completion for the decorated transfer action; a different coherent GFT state would define a different microscopic completion and must be tested against that vertex. H.11 proves that this statement is compatible with particle interference only when the renewed output excludes the persistent marked-position fiber and the discarded free-renewal record is branch-independent.

\paragraph{Reversible history export.}
Local replacement is compatible with global unitarity by Stinespring dilation \citep{Stinespring1955}. Introduce the old present register $A$, a fresh register $F$, its purifier $R$, and a history register $H$. For
\[
|\Psi_*\rangle_{FR}
=
\sum_{b\in\mathcal B}\sqrt{p_*(b)}\,|b\rangle_F|b\rangle_R
\]
in the diagonal closure-state realization, a register permutation can act as
\[
|\psi\rangle_A|\Psi_*\rangle_{FR}|0\rangle_H
\longmapsto
|\Psi_*\rangle_{AR}|0\rangle_F|\psi\rangle_H.
\]
Tracing out $R,F,H$ leaves $\rho_*$ on the new closure present and removes every dependence on the old replaceable state. The inverse unitary recovers that state from $H$, so no global information has been destroyed. Postulate~III supplies the interpretation of the exported register as history structure; the dilation theorem alone calls it an environment. For an unmeasured freely propagating defect, the state exported by this map may not encode the cell address of the mark. The marked fiber is transported in the retained defect system, and only a genuine interaction may correlate it with a discarded environmental record. Repeating the construction indefinitely requires fresh and history capacity or a controlled recycling mechanism, which the present paper has not derived.

\paragraph{The tetrahedral circuit.}
For
\[
\mathcal B
=
\{(s;m_1,m_2,m_3,m_4):s=\pm,\ m_i\in\{-3,\ldots,3\},\ m_i\ne m_j\},
\qquad |\mathcal B|=1680,
\]
the register permutation above is one circuit layer if the complete tetrahedron is a native local gate and the fresh admissible register has already been prepared. Neither injectivity nor the collective function $K^2(m_1,m_2,m_3,m_4)$ then requires a sequence of face repairs; the gate replaces one admissible four-face object by another.

A depth-one layer of disjoint two-face gates cannot test all six injectivity relations. The direct pair-comparison construction uses the three perfect matchings
\[
\{(1,2),(3,4)\},
\qquad
\{(1,3),(2,4)\},
\qquad
\{(1,4),(2,3)\},
\]
and therefore has depth three. This proves the stated depth for the direct comparator and excludes depth one in a pairwise architecture without a jointly coupled ancilla. It is not a lower-bound proof for every possible ancilla-assisted state-preparation algorithm. H.9 takes the complete tetrahedron as the native decorated gate; a geometric GFT embedding must realize that locality and sustain the prepared state dynamically.

\paragraph{Consequences and remaining gap.}
The foundational faithful full-support condition closes the classical replacement kernel when applied to histories. Detailed-balance symmetrization then closes $P=|\sqrt p\rangle\langle\sqrt p|$ and $P_\perp=I-P$ in probability space. Reference decoupling closes the quantum replacement-channel form, and the native tetrahedral architecture supplies a one-layer reversible dilation. H.9 adds the marked recurrence, after which positivity and the electron anchor close the clock, length, and phase-frequency conversion. These results remove an arbitrary transition matrix, geometric clock factor, and separate winding assumption while showing how local memorylessness can coexist with a globally retained past.

They do not by themselves close the coherent GFT Hessian. The exclusion sufficiency lemma of C.5 proves that one scalar statistic contains all local classical information distinguishing the canonical defect from the vacuum. The replacement process removes orthogonal probability memory. A microscopic spectrum must still show that no additional conserved or coherent light field survives and that the source couples only to the capacity singlet. The bridge from $P_\perp$ to a Euclidean mass or Schwinger--Keldysh relaxation pole therefore retains its conditional grade. The stationary reset carries no orientation; the macroscopic arrow comes from the export structure and the record root (G.7, P.8--P.9).


\subsection*{H.9 Decorated simplicial marked-transfer vertex}

The renewed marginal, lightest seven-channel branch, determinant recurrence,
and reversible present/history exchange determine a finite charged vertex with
a fully specified marked sector. In the selected factorized completion of
Section~23 and Appendix~Q, this vertex acts on the capacity/record factor; an
EPRL-class simplicial factor carries the vacuum geometry. The equivariant
$V_3$ frame lock relates the two factors without identifying their magnetic
indices. On the native tetrahedron, the vertex realizes faithful full-support
history resolution, and the established one-bit fermionic defect supplies the
hard-core mark. The field space and routing are physical choices: the
countervertices below fail the finite audits.

\paragraph{The marked fiber from primitive fusion.}
The nine-state marked space also has an independent geometric derivation. A shared face begins with two primitive fermionic slots of spin $j_0=3/2$. Maximum-capacity fusion selects the unique maximal channel,
\[
V_{3/2}\otimes V_{3/2}
=V_3\oplus\underbrace{(V_0\oplus V_1\oplus V_2)}_{Q},
\qquad \dim Q=9.
\]
The selected $V_3$ is the seven-state effective bulk link used by the closure ensemble. The nonmaximal information is the representation $Q=V_0\oplus V_1\oplus V_2$. The marked response was obtained from a different construction: the closure response is a vector $V_1$, and one response excitation on each of the present and history strands gives
\[
\mathcal H_{\rm mark}=V_1^P\otimes V_1^H
=V_0\oplus V_1\oplus V_2.
\]
Thus there is an $SU(2)$ intertwiner
\[
\mathcal I:Q\longrightarrow\mathcal H_{\rm mark},
\qquad
\mathcal I\,|(j_0j_0);J,M\rangle
=|(1_P1_H);J,M\rangle,
\quad J=0,1,2,
\]
with $\mathcal I^\dagger\mathcal I=P_Q$ and $\mathcal I\mathcal I^\dagger=I_9$. Each irrep occurs once, so the intertwiner is unique up to one phase on each $J$ block.

The match is more selective than the dimension count. For a response strand of spin $s$,
\[
V_s\otimes V_s=\bigoplus_{J=0}^{2s}V_J,
\]
whereas the complement of maximal fusion for two primitive spin-$j$ slots is
\[
Q(j)=\left(V_j\otimes V_j\right)\ominus V_{2j}
=\bigoplus_{J=0}^{2j-1}V_J.
\]
The two representations agree if and only if $j=s+1/2$. The independently
fixed closure-vector value $s=1$ therefore selects $j=3/2$. Appendix~B proves
the representation theorem and shows that unitarity of the selected sharp
reversible split fixes its retained output to $Q$. Since the complete H.9
one-mark carrier also has dimension nine, the complement map is an isometric
carrier-space identification. This proves the incidence used here without
deriving every microscopic detail: the equal weights, flags, relative block
phases, hard-core realization, and routing remain selected structure of the
explicit vertex below.

\paragraph{Fresh-state amplitude from the closure action.}
The Gaussian response uses an auxiliary vector whose norm includes the quantum variance in Appendix~B.2. Define
\[
\mathbf c(b)=\sum_i m_i\hat n_i,
\qquad
\mathbf C(b)=\sqrt{\frac{K^2(b)}{|\mathbf c(b)|^2}}\,\mathbf c(b).
\]
Injectivity ensures $|\mathbf c(b)|^2=(4\Sigma^2-S^2)/3>0$, so the definition is regular on the full ensemble. It gives $C_aC_a=K^2$ and transforms as a vector under a common frame rotation. The coordinates $C_a$ are a selected auxiliary realization of the scalar closure weight; their definition does not identify three commuting measurements of the quantum closure operator. The diagonal renewed state has amplitude
\[
A_*(b)=Z^{-1/2}\exp\!\left[-\frac{\eta_*}{2}K^2(b)\right],
\qquad
|A_*(b)|^2=p_*(b).
\]
For a standard three-component Gaussian auxiliary $\xi_a$,
\[
\exp\!\left[-\frac{\eta_*}{2}C^2\right]
=\int\frac{d^3\xi}{(2\pi)^{3/2}}
\exp\!\left[-\frac12\xi^2+i\sqrt{\eta_*}\,\xi_aC_a\right].
\]
The amplitude-level closure incidence is therefore $\sqrt{\eta_*}$. The selected charged decoration reuses this auxiliary strand with the same incidence. This fixes the response coupling within the chosen vertex. Faithful full-support resolution is being applied on a marked record space with no further joint constraint. Adding a defect-only invariant would define a nonminimal countervertex; it is allowed as a different theory and is one of the falsifiers of the construction, not an extra premise needed by the selected action.

The reversible update has independent present and history Gaussian strands. Their first normalized Hermite excitations obey
\[
\langle a|c\rangle_P=\delta_{ac},
\qquad
\langle b|d\rangle_H=\delta_{bd}.
\]
The ordered products $|a,b\rangle=|a\rangle_P\otimes|b\rangle_H$ form nine orthonormal states,
\[
\langle a,b|c,d\rangle=\delta_{ac}\delta_{bd}.
\]
They carry $\mathbb R^3\otimes\mathbb R^3=\mathbf1\oplus\mathbf3\oplus\mathbf5$. Introduce hard-core marked states $d_{ab}^\dagger|0\rangle$ and the stranded vertex
\[
V_Q=\sum_{a,b=1}^3d_{ab}^\dagger\xi_a^P\xi_b^H+\mathrm{h.c.}
\]
Its reduced Gram matrix is $I_3\otimes I_3=I_9$. Equal marked weights follow from this two-strand contraction, not from rotational symmetry or the number nine alone. Trace and epsilon contractions would split the three rotation sectors and falsify the vertex.

\paragraph{Pair contraction and canonical dilation.}
For each unordered occupied pair $e=(m,m')$, the two directed scalar returns are orthogonal alternatives. Define
\[
R_e=q\bigl(\langle m\!\to\!m'|+\langle m'\!\to\!m|\bigr),
\qquad q=\frac27,
\]
and $B_e=\sqrt{\eta_*}R_e$. Then
\[
B_eB_e^\dagger=2\eta_*q^2
=\frac{8\eta_*}{49}
\equiv u,
\qquad 0<u<1.
\]
The canonical unitary dilation of this contraction is
\[
U_e=
\begin{pmatrix}
\sqrt{I-B_e^\dagger B_e}&-B_e^\dagger\\
B_e&\sqrt{I-B_eB_e^\dagger}
\end{pmatrix}.
\]
On the occupied scalar-event support its no-event amplitude is
\[
\alpha=\sqrt{1-u}
=\sqrt{1-\frac{8\eta_*}{49}}.
\]
This identifies the origin of every entry: $\eta_*$ comes from the same closure amplitude that prepares the vacuum, $2q^2$ is the norm of the directed return pair, and the square root follows from reversible dilation.

For an auxiliary Grassmann representation, introduce
$\Gamma_e=(\gamma_e^P,\gamma_e^H)^T$ and
$J=\left(\begin{smallmatrix}0&1\\-1&0\end{smallmatrix}\right)$. With
$n_m=c_m^\dagger c_m$ and $n_d=\sum_{ab}d_{ab}^\dagger d_{ab}$, the marked action is
\[
\boxed{
\begin{aligned}
S_{\rm mark}
={}&g_{\mathrm{share,eff}}n_d
+\Lambda_dn_d(n_d-1)\\
&+\frac12\sum_{m<m'}
\Gamma_{mm'}^T
\left[1+n_dn_mn_{m'}(\alpha-1)\right]J
\Gamma_{mm'},\\
&\hspace{9em}\Lambda_d\to+\infty.
\end{aligned}
}
\]
The quadratic coefficient is $\alpha=\sqrt{1-u}$, not the probability $1-u$. The Grassmann integral returns a Pfaffian amplitude, so inserting $1-u$ in each $2\times2$ Majorana block would count the no-event probability twice and produce $(1-u)^{21}$ instead of the required amplitude $(1-u)^{21/2}$.

The Majorana variables represent the antisymmetric present/history amplitude; they are not additional propagating particles. When the marked fiber or either incident channel is absent, the reference block is $J$ and has Pfaffian one. On the established $n_m=1$ seven-channel branch, all $\binom72=21$ blocks become $\alpha J$, so
\[
Z_{\rm edge}=\alpha^{21}
=\left(1-\frac{8\eta_*}{49}\right)^{21/2}.
\]
These event records do not belong to the physical composite space $\Lambda^2\mathbb C^7$. The natural rank-one singlet lift on that composite space has the exact $6+15$ split derived in the audit; it does not act on the separately stranded internal present/history response records used here.

This distinction is why all 21 pair records survive even though the scalar source couples through one rank-one singlet. The singlet projector acts on physical channel amplitudes. The 21 Majorana blocks label mutually distinguishable internal present/history incidences, so quotienting them by the physical singlet would identify different incidences and violate faithful resolution. Their edge labels are internal to the marked response and carry no spatial cell address under the free vertex.

\paragraph{Complete decorated vertex.}
Let $\mathcal G_v$ denote the geometric fixed-spin gluing tensor and $W_*$ the native update that prepares $A_*(b)$ and exports the old replaceable closure register. The position and nine-state marked fiber belong to the retained defect system, not to a cell-addressed environment chain. Under free transport $W_*$ must therefore be translation-covariant and must leave the discarded renewal output independent of the defect's spatial branch. This is the no-which-path condition derived in H.11; it restricts the realization of the already-required quantum vertex and adds no foundational postulate. The incidence $V_Q$ above identifies the nine marked basis vectors with the ordered Gaussian response products. The transfer space also contains the unmarked vacuum. Its projector is
\[
\Pi_Q=|0\rangle\langle0|+\sum_{a,b=1}^3|a,b\rangle\langle a,b|.
\]
For $e=(m,m')$, define
\[
P_e=n_dn_mn_{m'},
\qquad
\widetilde U_e=(I-P_e)+P_eU_e.
\]
The complete one-step decorated transfer is
\[
\boxed{
\mathcal A_v
=\mathcal G_vW_*\Pi_Q
\exp\!\left[-g_{\mathrm{share,eff}}n_d-\Lambda_dn_d(n_d-1)\right]
\prod_{m<m'}\widetilde U_{mm'}.
}
\]
The empty state therefore has unit weight, while each of the nine one-mark states has fugacity $e^{-g_{\mathrm{share,eff}}}$. The factors $\widetilde U_e$ act on distinct edge-history strands and commute. The Grassmann action above is the Pfaffian trace representation of these controlled pair factors; it is not an additional operator multiplied into $\mathcal A_v$. Faithful full-support resolution on this internal response space, which has no further joint edge constraint, selects the tensor product of its one-edge marginals. Taking its finite internal trace at fixed position branch gives
\[
\boxed{
\zeta_*
=9e^{-g_{\mathrm{share,eff}}}
\left(1-\frac{8\eta_*}{49}\right)^{21/2}
=0.005123584484947.
}
\]

\paragraph{Complete edge-Hessian audit.}
On the 21-dimensional edge-label space, two distinct edges either share one endpoint or are disjoint. Let $A_1$ and $A_0$ be their adjacency matrices. The Johnson edge algebra has spectra
\[
\operatorname{spec}(A_1)=\{-2^{\times14},3^{\times6},10^{\times1}\},
\]
\[
\operatorname{spec}(A_0)=\{-4^{\times6},1^{\times14},10^{\times1}\}.
\]
A generic permutation-covariant edge kernel can therefore carry three eigenvalues. The conditioned vertex instead gives
\[
\mathcal K_{\rm edge}^{(2)}=I_{21}\otimes\alpha J,
\qquad
\operatorname{spec}_{\rm edge}=\{\alpha^{\times21}\},
\]
with zero off-diagonal edge norm. Coherently averaging the induced edge permutation matrix over all $7!=5040$ channel relabelings gives $\mathbf1\mathbf1^T/21$, of rank one. That operation would collapse the result to a shared singlet. It is not the charged gluing used here: $m=-3,\ldots,3$ are physical spin weights, $K^2$ is not invariant under arbitrary relabeling, and the defect Hamiltonian contains diagonal occupations and density interactions rather than hopping terms $c_m^\dagger c_{m'}$.

\paragraph{Primitive routing enumeration.}
Write $w=e^{-g_{\mathrm{share,eff}}}\alpha^{21}$, so $\zeta_*=9w$. The one-step transfer block contains three finite state spaces. The universal sector has one empty state and nine marked states, hence
\[
Z_\mu=1+9w=1+\zeta_*.
\]
The electron feedback sector has one empty state and states $|m;I_L,I_R\rangle$, with seven persistent charged labels and two independently renewed nine-state endpoints. Its dimension is
\[
1+7\times9\times9=568,
\]
and its trace is
\[
Z_{e,\mathrm{return}}=1+7\times81w^2
=1+7\zeta_*^2.
\]
The second-shell sector has one empty state and nine marked states with one scalar passage $q=2/7$, giving
\[
Z_{\tau,2}=1+9qw=1+\frac27\zeta_*.
\]
The direct product contains $10\times568\times10=56{,}800$ configurations. Exhaustive enumeration reproduces the analytic product exactly. The finite routing factors are therefore
\[
Z_\mu=1+\zeta_*,
\qquad
Z_e=(1+\zeta_*)(1+7\zeta_*^2),
\qquad
Z_{\tau,2}=1+\frac27\zeta_*.
\]
One native update applies each controlled gate once. Hard-core unit charge restricts the marked return to vacuum plus one return motif; the charged label persists diagonally while the response endpoints renew; and the shell grading terminates after $N=2$. These restrictions implement structures already present elsewhere in the paper: one native renewal pass, the one-bit fermionic anchor, the physical spin-weight labels, and the three-shell closure ladder. They add no separate foundational premise. Relaxing them gives explicit controls. A shared endpoint polarization changes the electron factor to $1+7\zeta_*^2/9$ and moves $G$ to $-14.6\sigma$; off-diagonal charged propagation gives $1+49\zeta_*^2$ and moves it to $+98.1\sigma$.

\paragraph{Charged generator and proper-time matching.}
The internal enumeration gives positive dimensionless motif weights $w_i$. For the electron their values and multiplicities are
\[
\begin{array}{c|rrrr}
w_i&1&w&w^2&w^3\\ \hline
\text{multiplicity}&1&9&567&5103
\end{array},
\qquad
\sum_iw_i=(1+9w)(1+567w^2)=Z_e.
\]
The dynamical completion assigns these weights additively to the charged Euclidean generator. Write $r=e^{-7g_{\mathrm{share,eff}}}$, $a=-(3/2)\ln(1-r)>0$, and let $n_e$ project onto the retained electron sector. For dimensionless marked activity $s$, the selected transfer obeys
\[
T(s+t)=T(s)T(t),\qquad T(0)=I,\qquad T(1)=e^{-a n_e}.
\]
Continuity and real positivity on the occupied scalar block fix $T(s)=e^{-as n_e}$, as checked by the Lean theorem \texttt{AppH1.scalar\_semigroup}. Assigning activity $w_i$ to each motif gives
\[
\boxed{
T_e=e^{-aZ_e n_e},\qquad
H_e=-\frac{\hbar}{\tau_*}\ln T_e
=-\frac{3\hbar Z_e}{2\tau_*}\ln(1-r)n_e.}
\]
The vacuum eigenvalue is one and the charged eigenvalue lies strictly between zero and one. This specifies the positive effective charged transfer after the internal records have been eliminated. The equality of motif weights and generator activities is a constitutive rule of the selected vertex. A derivation from a native microscopic evolution would have to reproduce this matching. Multiplying a transfer by a partition trace would instead add $-(\hbar/\tau_*)\ln Z_e$ to its energy; that operation is absent from the present prescription.

\paragraph{A finite recurrence completion.}
The same motif weights admit an explicit quadratic contraction that decorates a single charged recurrence. Let $M=5680$ be the electron motif count, let $v_i=\sqrt{w_i}$, and introduce auxiliary Grassmann variables for those $M$ alternatives and one charged determinant variable. Choose the dimensionless quadratic kernel
\[
\mathcal K_{\rm evt}=
\begin{pmatrix}
1&-\sqrt r\,v^\dagger\\
-\sqrt r\,v&I_M
\end{pmatrix},
\qquad v^\dagger v=\sum_iw_i=Z_e.
\]
The auxiliary variables represent internal recurrence alternatives. Their uncoupled determinant is one. Eliminating them gives the charged Schur complement and normalized Grassmann determinant
\[
K_{\rm charged}=1-rv^\dagger v=1-rZ_e,
\qquad \det\mathcal K_{\rm evt}=1-rZ_e.
\]
The spectrum of $\mathcal K_{\rm evt}$ consists of $1\pm\sqrt{rZ_e}$ and $M-1$ unit eigenvalues, so it is positive when $rZ_e<1$. With $n_e$ the charged projector, the positive effective transfer and its occupied energy are
\[
T_{\rm evt}=(I-n_e)+(1-rZ_e)^{3/2}n_e,
\qquad
E_e^{\rm evt}=-\frac{3\hbar}{2\tau_*}\ln(1-rZ_e).
\]
The electron anchor would then give
\[
L_*^{\rm evt}=-\frac32\lambda_e\ln(1-rZ_e).
\]
The coupling $\sqrt r\,v_i$ specifies how the motifs attach to the charged determinant line. It defines this alternative finite action without a fitted coefficient; the isolated internal trace does not select that attachment. Its difference from the additive-generator energy is
\[
\frac{E_e^{\rm evt}}{E_e}-1
=\frac{f(rZ_e)}{Z_e f(r)}-1
=\frac{Z_e-1}{2}r+O(r^2)
=7.36676\times10^{-26}.
\]
At fixed $m_e$ the fractional length change is the same, and the induced $G_*$ changes by $1.47335\times10^{-25}$ relatively. The displayed numerical precision is unchanged. For a fixed-vertex correlated cloud, $r\mapsto re^\Delta$ gives $E_e^{\rm evt}(P)/E_e^{\rm evt}(P_*)\ge e^\Delta$ by the same convexity argument as H.5. The fixed-marginal Shannon bound $H(P)\ge H$ ensures $Z_e r(P)\le Z_e e^{-H}<1$ throughout this readout class.

The coefficient chain below uses the additive-generator prescription $E_e=(3\hbar/2\tau_*)Z_e f(r)$. The finite recurrence contraction supplies a second fully specified action. Selection between their exact spectra requires a physical rule for attaching the marked response to the charged evolution.

The independent-factor completion gives another check with the same weights. Its positive transfer
\[
T_{\rm ind}=\prod_i(1-rw_i)^{(3/2)n_e}
\]
has occupied energy $-(3\hbar/2\tau_*)\sum_i\ln(1-rw_i)$. Relative to $H_e$, its fractional difference is $-7.32\times10^{-26}$ at the stated ensemble. This estimate applies to this specified alternative. Positivity alone supplies no universal bound on other generator prescriptions.

\paragraph{Results and scope.}
The marked vertex gives $G_*$, $m_\mu/m_e$, and $m_\tau/m_e$ at $-0.073\sigma$, $-0.535\sigma$, and $-0.125\sigma$, fractional agreements of about $1.6\times10^{-6}$, $1.2\times10^{-8}$, and $6\times10^{-6}$, with no continuously adjusted coefficient; the expressions and values are displayed below. The comparisons use CODATA 2022 for $m_e$, $G$, and $m_\mu/m_e$, with the exact SI definitions of $h$ and $c$. The PDG tau mass is $1776.93\pm0.09$ MeV; the selected ratio gives $m_\tau c^2=1776.9188$ MeV and a residual $-0.125\sigma$ \citep{CODATA2022Astra,PDG2026Astra}. These residuals use experimental standard uncertainties and do not include a probability distribution over alternative microscopic prescriptions.


The physical scale and charged-lepton ratios are
\[
L_*=Z_e\left[-\frac32\lambda_e\ln(1-e^{-7g_{\mathrm{share,eff}}})\right],
\qquad
G_*=\frac{c^3L_*^2}{\hbar},
\]
\[
\frac{m_\mu}{m_e}=720\frac27Z_\mu,
\qquad
\frac{m_\tau}{m_e}=720^2\left(\frac27\right)^4Z_\mu Z_{\tau,2}.
\]
They evaluate to $G_*=6.6742890813\times10^{-11}\,\mathrm{m^3\,kg^{-1}\,s^{-2}}$, $m_\mu/m_e=206.768280237$, and $m_\tau/m_e=3477.343310$. The residuals are $-0.073\sigma$, $-0.535\sigma$, and $-0.125\sigma$.

The finite marked weights and charged transfer prescriptions are specified; the numerical chain uses the displayed additive-generator prescription. Its mass minimum holds at fixed marginals and marked vertex, while its ordinary gravitational normalization uses the physical matching of C.5. The alternative kernels
remain useful falsifiers: a collective edge mode, a physical pair-composite
interpretation, persistent endpoint polarization, additional defect-only
coupling, or off-diagonal charged propagator changes the answer. The geometric
branch is also more specific than a ``stable condensate'' requirement.
Appendix~B supplies exact $J=3$ phase selection and maximal-fusion blocking;
Appendix~Q supplies a separate gauge-covariant capacity factor with an
equivariant frame lock; and EPRL/Regge constructions supply controlled
two-helicity block witnesses. H.9a gives the full ordinary-gravity chain.
H.10 proves the massless-TT infrared theorem on the selected metric-Regge
branch, constructs an exact fixed-$j=3$ perfect-action realization, closes the
factorized capacity audits, and proves the clock-test theorem. It also shows
that the capacity premises do not fix the geometric gluing tensor or measure.
That geometric selection problem, the physical source of history capacity, and the microscopic origin of the carrier-resolved transverse contact remain open; Appendix~P.8 gives the exact one-tick H-theorem of renewal and the domain of the thermodynamic arrow.

\paragraph{Reproduction.}
The supplementary scripts reproduce the ensemble invariants, state-weighted determinant, $48\times35$ locality fracture, Gaussian closure amplitude, unitary pair dilation, nine-state Gram matrix, 21-block edge Hessian, and 56,800-state routing trace. They also verify the dressed scale and dressed charged-lepton ratios of Appendix~I.1. Appendix~R contains the self-contained numerical-spine script; the complete Hessian and graph enumerations remain in the accompanying audit files.

\subsection*{H.9a From the finite cell to the Einstein action}

\paragraph{Purpose and logical structure.}
The ordinary-gravity claim is not that the $1680$-state capacity code secretly contains the Einstein action.  Geometry and capacity are distinct operator factors in the selected realization, and each step below has its own logical grade.  The derivation is
\[
\boxed{
\begin{gathered}
\text{finite capacity cell}
\longrightarrow
\text{selected relational quantum geometry}
\longrightarrow
\text{metric Regge branch}\\
\longrightarrow
\text{Einstein gravity in the infrared}.
\end{gathered}}
\]
The representation theory and maximal-fusion ray are exact once the factorized geometric completion is selected.  Metricity, block-scale Lorentz invariance, and membership in the appropriate semiclassical phase are the named geometric premises.  The capacity action fixes the channel, record dynamics, source, and dimensional scale; it does not by itself choose the complete ultraviolet geometric propagator.

\paragraph{1. Relational frame geometry.}
In the selected spatial-frame completion, each tetrahedral cell $c$ carries a local frame $R_c$, while only relative transport between adjacent cells is observable:
\[
g_{cc'}=R_cU_{cc'}R_{c'}^{-1}\in SU(2).
\]
Here $U_{cc\prime}$ is genuine link transport expressed in reference frames, with $U_{c\prime c}=U_{cc\prime}^{-1}$. A change of local frame, $R_c\mapsto h_cR_c$, acts by
\[
g_{cc'}\mapsto h_cg_{cc'}h_{c'}^{-1}.
\]
A loop based at $c$ has holonomy $R_cU_{\rm loop}R_c^{-1}$ and can carry curvature. Thus the selected relational description has the standard local $SU(2)$ gauge redundancy.  If $\mathbf J_{cf}$ is the flux through face $f$ in the frame of cell $c$, gauge invariance imposes the Gauss constraint
\[
\boxed{\mathbf G_c=\sum_{f\subset\partial c}\mathbf J_{cf}=0.}
\]
The geometric state space of a primitive tetrahedron with four spin-three faces is therefore
\[
\boxed{
\mathcal H_c^{\rm geom}
=\operatorname{Inv}_{SU(2)}\!\left(V_3^{\otimes4}\right).}
\]
Since
\[
V_3\otimes V_3=\bigoplus_{J=0}^{6}V_J,
\]
each intermediate $J=0,\ldots,6$ can be paired with the same $J$ in the other pair to form a singlet.  Hence
\[
\boxed{\dim\operatorname{Inv}_{SU(2)}(V_3^{\otimes4})=7.}
\]
This seven-dimensional intertwiner space is geometric.  It must not be identified with the independent seven-state capacity carrier.  The cell space remains
\[
\boxed{
\mathcal H_{\rm cell}
=\mathcal H_{\rm geom}\otimes\mathcal H_{\rm cap}
\otimes\mathbb C_{\rm orient}^{2},}
\]
with the precise factorization and gauge action given in Appendix~Q.4.

Classically one writes the face flux as
\[
\mathbf E_f=A_f\mathbf n_f,
\qquad
\sum_fA_f\mathbf n_f=0.
\]
For nondegenerate outward normals, the closure data reconstruct a convex tetrahedron up to translation.  The dictionary is therefore
\[
\boxed{
j_f\leftrightarrow A_f,\qquad
\iota_c\leftrightarrow\text{tetrahedral shape},\qquad
g_{cc'}\leftrightarrow\text{relative connection}.}
\]
This is the first gravity step: the selected geometric factor carries the usual holonomy--flux data of relational three-geometry.  It is a selected completion of the finite substrate, not a consequence of the $1680$ count alone.

\paragraph{2. Carrier selection and the exact coarse ray.}
The primitive positive mismatch operator acts in the capacity fusion factor:
\[
Q_f=6P_0+5P_1+3P_2,
\qquad
\ker Q_f=V_3^{\rm cap}.
\]
Consequently
\[
F_3e^{-sQ_f}F_3^\dagger=I_{V_3^{\rm cap}}
\qquad(s>0).
\]
The equivariant lock of Appendix~Q.5 then pairs this selected carrier with a geometric copy,
\[
L_3:V_3^{\rm cap}\longrightarrow V_3^{\rm geom},
\qquad
L_3U_{\rm cap}(g)=U_{\rm geom}(g)L_3.
\]
Thus $Q_f$ selects the capacity carrier and the lock selects the matching geometric representation.  Neither operation supplies curvature:
\[
\boxed{
\begin{gathered}
Q_f\text{ and }L_3\text{ select and relate the face data;}\\
\text{the geometric host supplies their propagation and curvature.}
\end{gathered}}
\]

For a block of $N$ aligned primitive faces, the highest-spin irrep occurs once in $V_3^{\otimes N}$.  There is therefore a unique normalized equivariant fusion map up to phase,
\[
\boxed{F_N:V_3^{\otimes N}\longrightarrow V_{3N},
\qquad J_N=3N.}
\]
On spin coherent states,
\[
F_N|3,\mathbf n\rangle^{\otimes N}=|3N,\mathbf n\rangle,
\]
and coherent transport fuses exactly:
\[
\boxed{
\langle3,\mathbf n|D^3(u)|3,\mathbf n'\rangle^N
=\langle3N,\mathbf n|D^{3N}(u)|3N,\mathbf n'\rangle.}
\]
The semiclassical parameter is therefore the block size, not a running microscopic spin:
\[
\boxed{j_{\rm primitive}=3\ \text{fixed},
\qquad N\to\infty,
\qquad J=3N\to\infty.}
\]
The fusion identity is exact on the maximal coherent ray.  It is not a proof that every interacting refinement trajectory flows onto that ray; the inter-cell dynamical refinement remains a host-theory test.

\paragraph{2a. The locked ray has an exact collective phase space.}
The number and phase of a pre-glued spin-three face mode form the collective
phase space of the coarse spin. For coherent amplitude
\[
\alpha=\sqrt{\bar N}\,e^{-i\theta}
\]
and transported endpoint normals $\mathbf n_L,\mathbf n_R$, the Berry one-form is
\[
\vartheta_{\rm B}
=\bar N\,d\theta
+3\bar N\bigl[\mathcal A(\mathbf n_L)-\mathcal A(\mathbf n_R)\bigr].
\]
With
\[
J=3\bar N,
\qquad
\xi=\frac{\theta}{3},
\]
this becomes
\[
\boxed{
\vartheta_{\rm B}
=J\,d\xi
+J\bigl[\mathcal A(\mathbf n_L)-\mathcal A(\mathbf n_R)\bigr],}
\]
the twisted-geometry symplectic potential on $T^*SU(2)$.  The identification is exact on the locked ray.  Because every allowed spin is $J=3N$, the twist fiber obeys
\[
\boxed{
(J,\xi)\in3\mathbb N\times
\mathbb R/(2\pi/3)\mathbb Z
\simeq
(N,\theta)\in\mathbb N\times S^1.}
\]
The resulting $\mathbb Z_3$ quotient is structural, not a coordinate artifact.

The same ray has an exact overcomplete family.  With
\[
c_N(\alpha)=e^{-|\alpha|^2/2}\frac{\alpha^N}{\sqrt{N!}},
\qquad
|\alpha,\mathbf n\rangle\!\rangle
:=\sum_{N=0}^{\infty}c_N(\alpha)\sqrt{6N+1}\,|3N,\mathbf n\rangle,
\]
Glauber orthogonality and the spin-coherent resolution give
\[
\boxed{
\int\frac{d^2\alpha}{\pi}
\int\frac{d\Omega(\mathbf n)}{4\pi}
|\alpha,\mathbf n\rangle\!\rangle
\langle\!\langle\alpha,\mathbf n|
=\bigoplus_{N=0}^{\infty}I_{V_{3N}}.}
\]
Thus the collective number, twist, and normals provide a genuine coherent-state path-integral chart, with relative uncertainties $\Delta J/\bar J=O(\bar J^{-1/2})$.  This construction requires a physical phase: the selected GFT field must be real or its global $U(1)$ phase must be broken.  It is stated for a pre-glued link because the twist is relational between its two ends.

\paragraph{2b. Lorentzianization is applied after fusion.}
The representation fixes the ordering. On the lowest $SU(2)$ block $V_j$ of
the principal-series representation $(\rho,k)=(\gamma j,j)$, rotational
covariance and $\mathbf K\!\cdot\!\mathbf L=\rho k$ give
\[
\boxed{
P_jK_iP_j=\frac{\gamma j}{j+1}L_i.}
\]
Applying the Lorentzian $Y$ map separately to primitive $j=3$ facets therefore gives the coefficient $3\gamma/4$.  Fusing first to $J=3N$ and applying $Y$ only at block scale gives instead
\[
\frac{\gamma J}{J+1}
=\gamma\frac{3N}{3N+1}
\longrightarrow\gamma.
\]
The two orderings differ by $25\%$ at primitive spin.  The area normalization below was fixed from the large-$J$ law with $J=3N$, so its consistent implementation is
\[
\boxed{\text{exact }SU(2)\text{ fusion first, block-scale }Y_{\gamma_*}\text{ second}.}
\]
Facet-first Lorentzianization would require a different Immirzi normalization and would define a different theory.

\paragraph{2c. The orientation doublet is renewed, geometrically inert, and inert for fermions in the record-oriented frame.}
The binary orientation factor cannot be a hidden persistent register.  If a uniformly populated lift label $\sigma\in\{+,-\}$ survived from one renewal step to the next, then the renewed boundary data would obey
\[
I(B_t;B_{t+1})\ge H(\sigma)=\ln2,
\]
contradicting the full-support memoryless kernel proved in H.4.  The doublet is therefore refreshed rather than transported as an additional geometric degree of freedom.

On the selected pure-gravity branch this causes no loss of geometric information.  The two lifts $(\nu,\omega)=(+,+)$ and $(-,-)$ have the same Regge-sector label $\mu=\nu\omega=+1$, determine the same metric, and carry the same $e^{+iS_R}$ asymptotic phase.  The proper vertex may consequently be extended as
\[
\boxed{\mathcal G_v=A^{\rm prop}_{v,\mu=+}\otimes I_{\rm orient},}
\]
so that gravity depends only on the metric-equivalence class and the refreshed bit is inert.

That conclusion does not automatically extend to fermions, which couple to a tetrad rather than to the metric alone.  A matter completion must verify
\[
\boxed{
\mathcal A_{\rm matter}[e,\psi]
=\mathcal A_{\rm matter}[e',\psi']
\quad\text{under the simultaneous lift flip at fixed }g_{\mu\nu},}
\]
with the corresponding spinor transformation included. The discrete frame flips at fixed $g_{\mu\nu}$ decide the question. With spinor representatives $S$ obeying $S^{-1}\gamma^aS=\Lambda^a{}_b\gamma^b$ and the physical conjugate $\bar\psi'=\psi'^\dagger\gamma^0$, the parity flip $P$ ($S=\gamma^0$) leaves the kinetic and mass terms invariant with the volume factor $\sqrt{-g}$, the time flip $T$ ($S=\gamma^5\gamma^0$) leaves them invariant with the signed volume form $\det e$, and the inversion $e\to-e$ ($S=\gamma^5$) reverses the sign of both terms with either volume factor. A compensating map would have to satisfy both $S^\dagger\gamma^0S=\gamma^0$ and $S^\dagger\gamma^0\gamma^aS=-\gamma^0\gamma^a$ for every $a$. At $a=0$ the latter condition would give $S^\dagger S=-I$, which is impossible. The finite complex-matrix obstruction and the parity/time-flip bilinears are kernel checked. For the inversion representative,
\[
\boxed{S_D[-e,\gamma^5\psi]=-S_D[e,\psi].}
\]
Bivector data fix a 4-simplex up to inversion \citep{BarrettEtAl2009}, so the lift pair that shares both the bivectors and the metric is related by $e\to-e$; reading $(\nu,\omega)$ as orientation and inversion makes the flip $T$. Both reverse the time orientation of the frame, and the Dirac sector distinguishes at least one of them. The statement assumes the physical conjugate and Grassmann fields; in a Euclidean integral with $\bar\psi$ independent, the map $\psi\to\gamma^5\psi$, $\bar\psi\to\bar\psi\gamma^5$ compensates the inversion as a chiral rotation, and the sign reappears in the anomalous measure. The inversion reverses $H$, so a branch evolves by $e^{\mp iH\tau_*/\hbar}$ according to its lift. Renewal is cell-local, so two branches of a superposition in separated cells draw independent lifts, and their relative phase diffuses at the rate $(mc^2/\hbar)^2\tau_*$: $3\times10^{-2}\,{\rm s^{-1}}$ for an electron, $8\times10^{8}\,{\rm s^{-1}}$ for a rubidium atom, and $7\times10^{13}\,{\rm s^{-1}}$ for a $25{,}000\,$u molecule. Rubidium wave packets separated by half a metre stay coherent for about $1\,$s \citep{Kovachy2015}, and such molecules for milliseconds \citep{Fein2019}. Coupling fermions to the time-oriented representative of the frame class, with $e^0$ future-directed along the record order of Appendix~P.8--P.9, restores the equality. With $\sigma={\rm sgn}(e^0\!\cdot n_{\rm rec})$,
\[
\boxed{\mathcal A_{\rm matter}=S_D[\sigma e,\psi]}
\]
is invariant under both the inversion and $T$, and the refreshed bit stays inert for fermions. Pseudoscalar and axial couplings, including Holst- and Nieh--Yan-type terms, change sign under $P$; they describe ordinary parity violation and require a spatial orientation, which the chiral weak sector must supply and the record order does not. None of the pure-geometric results below depends on these conditions.

\paragraph{2d. The native lattice fermion.}
On the diamond graph, the coarse adjacency realization of Appendix~C.5, each cell has four neighbors across faces with normals $\hat n_j$, and a spinor crossing face $j$ carries the lift of the native gluing $R_j$, $U_j=\exp(-i\pi\,\hat n_j\cdot\mathbf J)$, and the one-particle hopping operator is
\[
H(\mathbf k)=\begin{pmatrix}0&Q(\mathbf k)\\Q(\mathbf k)^\dagger&0\end{pmatrix},\qquad Q(\mathbf k)=t\sum_{j=0}^{3}U_j\,e^{i\mathbf k\cdot\mathbf d_j},
\]
with $\mathbf d_j$ the bond vectors along the face normals. For spin $\tfrac12$, $U_j=-i\hat n_j\cdot\boldsymbol\sigma$. The four normals sum to zero and $\sum_j\hat n_j\hat n_j^{T}=\tfrac43I$, so $Q(0)=0$ and
\[
\boxed{H\simeq v\,\tau_x\otimes\boldsymbol\sigma\cdot\mathbf k,\qquad v=\frac{ta}{\sqrt3},}
\]
an isotropic massless Dirac fermion at $\mathbf k=0$. Its right- and left-handed components are the symmetric and antisymmetric combinations of a cell and its mirror neighbor. Spin $\tfrac32$ gives no state at $\mathbf k=0$ and a quadratic touching at the W points, and spin $0$ gives only nodal lines, so the native gluing selects spin $\tfrac12$ for a relativistic fermion. The remaining zero modes form three closed nodal loops between the X and W points, each of length $3.33\,(2\pi/a)$ with sublattice winding $+1$. A Zeeman field that breaks time reversal splits the cone into Weyl nodes of charge $\pm1$ and each loop into a further pair, and the charges sum to zero, as the Nielsen--Ninomiya theorem requires \citep{NielsenNinomiya1981}. The amplitude through two gluings, $U_l^\dagger U_j=\hat n_l\cdot\hat n_j+i(\hat n_l\times\hat n_j)\cdot\boldsymbol\sigma$, contains the spin--orbit structure of the Fu--Kane--Mele diamond model \citep{FuKaneMele2007} and leaves states at the Dirac energy. A second-neighbor term weighted by the handedness $h_c$ of the cell,
\[
H_W=r\,h_c\sum_{j<l}\big[1-\cos\mathbf k\cdot(\mathbf d_j-\mathbf d_l)\big],\qquad h_A=-h_B,
\]
vanishes at $\mathbf k=0$, opens a gap of $0.34\,t$ over the three loops for $r=0.1$--$0.3\,t$, and leaves the cone unchanged. The graph therefore carries one doubler-free Dirac fermion, and the lattice sublattice symmetry is broken in the manner of the Wilson term. The cone also splits at first order in the cell size,
\[
E_\pm=vk\Big[1\mp\frac{ak}{4}\frac{|\boldsymbol\kappa|}{k^3}\Big],\qquad \boldsymbol\kappa=\big(k_x(k_y^2-k_z^2),\,k_y(k_z^2-k_x^2),\,k_z(k_x^2-k_y^2)\big),
\]
with $\boldsymbol\kappa$ the cubic Dresselhaus vector, so the splitting vanishes along $\langle100\rangle$ and $\langle111\rangle$ and is largest along $\langle110\rangle$.

\paragraph{2e. A persistent spatial orientation from the gluing.}
The sign $h_c$ is the handedness of a cell's labelled face normals, $\hat n_1\cdot(\hat n_2\times\hat n_3)$, measured against a spatial orientation, and a chiral coupling needs the same orientation. Three facts supply it. First, the frame enters through its lift class $\{e,\Lambda e\}$, with $\Lambda$ the inversion or the time flip, and the record order of Appendix~P.8--P.9 selects the member with $e^0$ future-directed. That representative is the same for both members of the class, so its spatial orientation is a class invariant that the refreshed bit cannot change: for the inversion it is fixed by the 4-orientation, which the inversion preserves, and for the time flip it is the spatial orientation the two members share. Second, the 4-orientation and the spatial orientation are coherent orientations of the complex, and on a connected orientable manifold the face identifications fix them up to one global sign. Spinors require an orientable spin manifold, so the host is orientable wherever fermions propagate. The identifications belong to the geometric factor, which renewal does not replace, and a Pachner move replaces a ball by a ball with the same boundary, onto which the orientation extends uniquely. The coherent orientation therefore persists. It carries one bit for each connected component, while the memoryless kernel of H.4 constrains registers cell by cell, so the two statements are compatible. Third, the persistence is required. On a periodic diamond graph of $7^3$ cells with $r=0.2\,t$, the coherent orientation, with either global sign, leaves only the four states of the cone below $|E|=0.5\,t$. Refreshing the orientation of $1\%$ of the cells places $22$ further states in that window, bound to the refreshed cells, and refreshing every cell gives $152$ extended states, which restore the doublers. The doubler-free fermion therefore exists exactly when the orientation is carried by the gluing.

The orientation defines which hand a chiral coupling selects. The coupling itself remains open: the weak group must act on the mirror-odd combination alone, or interactions must decouple the mirror-even partners, and a local bilinear operator can do neither \citep{NielsenNinomiya1981}. Exactly gauge-invariant Ginsparg--Wilson constructions \citep{Luscher1999} and symmetric mass generation \citep{WangYou2022} are the candidate routes.

\paragraph{2f. The native operator on closed complexes.}
Regular cells cannot realize the diamond graph. In a closed triangulated three-manifold $F=2T$ and $\chi=0$ give $E=V+T$, so the mean number of cells around an edge is $6T/(V+T)<6$ and some edge is surrounded by a cycle of at most five cells. The shortest cycle of the diamond graph has six. Mirror gluing that preserves the port labels makes the complex balanced: the face colors are the colors of a proper four-coloring of the vertices, each color class of faces is a perfect matching between the two orientation classes of cells, and the cells around each edge alternate two colors in a closed cycle. Summing displacements around these cycles gives the same mean displacement $\mathbf D$ across faces of every color. The four uniform states, $\chi_A$ on every cell of one class and $\chi_B$ on every cell of the other, are exact zero modes on every balanced complex because $\sum_jU_j=0$, and their first-order momentum term is $\big(\sum_j\hat n_j\cdot\boldsymbol\sigma\big)\otimes\mathbf D=0$. Two-step and handedness-weighted terms built from color sequences leave it zero.

The nearest periodic case is the tetragonal-disphenoid honeycomb on the body-centred cubic lattice. It is balanced, with four or six cells around each edge and a mean of $36/7$ against the flat equilateral value $2\pi/\arccos\tfrac13=5.104$. Twelve further zero modes at $\mathbf k=0$ couple to the uniform states, and the sixteen modes form two isotropic linear families with speeds $t\ell/\sqrt5$ and $4t\ell/(5\sqrt3)$ for cube edge $\ell$. Their relative splitting reaches $\pm\ell k/(5\sqrt3)$ along $\langle111\rangle$ and vanishes along $\langle100\rangle$ and $\langle110\rangle$. Written as $E^2=p^2+\eta\,p^3/M_P$, with $\ell$ fixed by the edge $a=2.530L_*$, the splitting reaches $|\eta|=0.64$, and on the diamond graph $|\eta|$ reaches $0.60$. The Crab Nebula bounds this coefficient for leptons at $10^{-5}$ \citep{Maccione2007}.

Generic balanced complexes have no extra zero modes. On barycentric subdivisions of periodic Poisson--Delaunay triangulations, with mean edge degree $5.09$, and on the subdivided body-centred honeycomb, only the four uniform zero modes exist, and a twisted boundary condition gives $E(2k)/E(k)=4.00$. The touching is quadratic, $E\simeq(0.85$--$0.91)\,t\,(kb)^2$ for mean bond length $b$. The one-spinor-per-cell operator with the native gluings therefore yields a relativistic fermion on the diamond graph and on the ordered body-centred complex, with a Planck-suppressed splitting above the Crab bound in both, and a nonrelativistic touching on generic complexes. The Dirac--K\"ahler construction places fields on all simplices and is defined on any simplicial complex \citep{BecherJoos1982,Rabin1982}; on a closed complex its zero modes are the harmonic cochains, counted by the Betti numbers. It is the standard candidate for the native fermion, and Lorentz invariance requires equal speeds in its separate form sectors. The orientation of paragraph 2e belongs to the complex and is available to any such construction.

\paragraph{2g. The Dirac--K\"ahler fermion on the complex.}
The Dirac--K\"ahler field places one amplitude on every vertex, edge, face, and cell, and its Hamiltonian is $D=d+d^\dagger$, with $d^\dagger$ the adjoint of the coboundary under diagonal Hodge weights $\star_p$. On every complex tested its zero modes at $\mathbf k=0$ number $1,3,3,1$ by degree, the Betti numbers of the three-torus, so the field carries the eight components of two Dirac tastes and no lattice doublers. Reality gives $D(-\mathbf k)=\overline{D(\mathbf k)}$. This relates the unordered spectra at opposite momenta and leaves an odd, imaginary antisymmetric part in the effective small-$k$ block. That part splits degenerate transverse polarizations linearly in $k$, the lattice form of optical activity, and is the dimension-five term. Inversion maps $\mathbf k\to-\mathbf k$ and acts as $-1$ on the vector-like sectors, so it removes the odd part on centrosymmetric complexes. The operator $D=d+d^\dagger$ with diagonal Hodge stars uses no orientation data, and the host ensemble weights a complex and its mirror image equally: the CDT action depends only on simplex counts, and the closure code is parity-doubled. The ensemble-averaged term therefore vanishes. On periodic Delaunay complexes the odd block is below $10^{-17}$ for a centrosymmetric complex. A random chiral realization has an odd block growing as $k^3$, and its mirror image has exactly the opposite coefficient. Per realization the isotropic coefficient has random sign, and its rms falls as $N^{-1/2}$ with the number of cells, reaching $1.4\times10^{-3}$ per unit $ka$ at $128$ cells. Averaged over the volume of a TeV wavelength this gives $|\eta|\sim10^{-26}$, against the Crab bound of $10^{-5}$. Fits of $v(k)=v+c_1k+c_2k^2$ on the tested complexes give $|c_1|/v<10^{-5}$ in cell units. The one-spinor-per-cell operator of 2f uses the cell handedness $h_c$ and is chiral, which is why it carries the splitting. A finite Immirzi parameter in the geometric factor couples to fermions through the parity-odd Holst term; it produces an axial contact interaction and no propagation term. At small $k$ the spectrum splits into three sectors: the scalar sector (vertex forms and exact edge forms) with speed $v_0$, two transverse sectors (coexact edge forms and exact face forms) with speed $v_T$, and the cell sector (cell forms and coexact face forms) with speed $v_3$.

With the Voronoi--Delaunay weights of a flat realization ($\star_0$ the Voronoi volume, $\star_1$ and $\star_2$ the ratios of dual to primal measures, $\star_3$ the inverse cell volume) all three speeds equal one in every direction, with corrections of order $k^2$, on the body-centred honeycomb and on periodic Poisson--Delaunay triangulations, as in random-lattice field theory \citep{ChristFriedbergLee1982}. With the weights of the regular cell, $\star_0\propto n_T(v)$ and $\star_1\propto\deg(e)$ with $\star_2$ and $\star_3$ constant, the speeds separate. On the body-centred honeycomb $v_0^2:v_T^2:v_3^2=\tfrac56:\tfrac45:\tfrac34$, isotropic in every direction. On Poisson--Delaunay triangulations the ratios self-average to $v_T/v_0=0.826$ and $v_3/v_0=0.814$, and on balanced barycentric subdivisions they fall to $0.64$--$0.74$ and $0.51$--$0.58$. Constant rescalings of each $\star_p$ equalize the speeds on one complex and leave them unequal on another: the rescaling that equalizes the honeycomb gives $0.848$ and $0.859$ on Poisson--Delaunay triangulations. Every diagonal star built from regular-cell data reduces to such a rescaling. Speed differences of this size within one fermion multiplet are excluded by many orders of magnitude \citep{KosteleckyRussell2011}.

The Dirac--K\"ahler field is relativistic and free of doublers, and its dimension-five term vanishes on centrosymmetric complexes and in the mirror-symmetric ensemble. Its eight components share one light cone with the weights of a flat realization of the complex. The regular-cell weights that the face-wise lock assigns give each form sector its own light cone. A photon carried by edge forms on the same complex moves with $v_T$, so one condition sets the common light cone of matter and radiation, and radiative corrections spread any residual difference across species \citep{Collins2004}. A principle that gives matter fields the weights of the flat realization, together with an account of the two tastes, completes the native fermion.

\paragraph{3. One channel fixes one area increment on the matched branch.}
The physical area on the coherent maximal-fusion ray is linear in the additive flux,
\[
A(J)=8\pi\gamma_*L_*^2J.
\]
A boundary cut removes one capacity partner and contributes
\[
S_N/k_B=N\ln2.
\]
Matching this microscopic cut to the Bekenstein--Hawking normalization on the Einstein branch gives
\[
N\ln2
=\frac{A_N}{4L_*^2}
=\frac{8\pi\gamma_*L_*^2(3N)}{4L_*^2},
\]
and hence
\[
\boxed{
\gamma_*=\frac{\ln2}{6\pi},
\qquad
A_N=4\ln2\,L_*^2N,
\qquad
a_{\rm ch}=4\ln2\,L_*^2.}
\]
This is a normalization match, not an independent derivation of the black-hole area law; Appendix~Q.11 gives the finite-$N$ area-operator statement and that caveat in full.

The dimensional scale is not left free here.  Appendix~D.4 and Sections~H.6--H.9 use the electron recurrence and the marked-transfer factor to obtain
\[
\boxed{
L_*=-\frac32Z_e\lambda_e
\ln\!\left(1-e^{-7g_{\mathrm{share,eff}}}\right),
\qquad
G_*=\frac{c^3L_*^2}{\hbar}.}
\]
Thus the area relation fixes the geometric normalization in units of $L_*$, while the independent charged-transfer calculation fixes the numerical length and therefore the numerical Newton scale.  The two statements should not be conflated.

\paragraph{4. Metricity is a selected physical subspace.}
Gauge-closed tetrahedra alone span a twisted-geometry space: adjacent cells can agree on a face area while assigning different intrinsic shapes to that face.  The ordinary one-metric branch additionally imposes shape matching,
\[
\boxed{
q_{ab}^{(c)}\big|_f=q_{ab}^{(c')}\big|_f.}
\]
On this support, neighboring tetrahedra glue into a single piecewise-flat metric geometry.  Therefore
\[
\boxed{
\text{Gauss closure}+\text{shape matching}
\Longrightarrow\text{Regge geometry}.}
\]
Shape matching is not generated by $Q_f$.  It is part of the selected metric completion and must be preserved by the geometric vertex and by its coupling to the capacity decoration.

\paragraph{4a. Equal face areas force disphenoid cells; the face-wise lock forces regular cells.}
Two facts narrow the selected subspace.  First, every primitive face carries $j=3$, so on the coherent ray every face has the same area $a_{\rm ch}=4\ln2\,L_*^2$ (paragraph~3).  A tetrahedron whose four faces have equal area is a disphenoid: opposite edges are equal and the four faces are congruent.  Primitive cells are therefore disphenoids exactly, with no selection, and the only residual shape freedom is the shape of the congruent face triangle, two parameters per cell.  Second, the capacity-closure polynomial of Appendix~Q.7, $|\sum_ib_iv_i|^2=4\sum_ib_i^2-(\sum_ib_i)^2$, is written on the vertex directions $v_i$ of the regular tetrahedron, and Appendix~Q.4 implements Postulate~I by the equivariant lock on the two relational frames.  If that lock is read face-wise, so that capacity port $i$ and geometric face $f_i$ are one face of one substrate and $\mathbf n_{f_i}=R_cv_i$ for a single rotation $R_c$, then the four normals form a regular tetrahedron and the cell is regular.  Every face is then an equilateral triangle of area $a_{\rm ch}$, adjacent cells agree on face shape identically, and
\[
\boxed{\text{face-wise lock}\Longrightarrow\text{shape matching holds kinematically}.}
\]
Twisted configurations are excluded rather than selected, and the native spatial geometry is the equilateral-triangulation class with fixed edge
\[
\boxed{a^2=\frac{16\ln2}{\sqrt3}L_*^2=6.403L_*^2,\qquad a=2.530L_*.}
\]
With time-like edges set by the renewal cadence and equal by homogeneity, the spacetime building blocks are the two causal-dynamical-triangulation simplex types with one asymmetry $\tilde\alpha=\ell_t^2/a^2$: the native geometric host is the CDT class with Regge weight $(\ln2/2\pi)\sum_hN_h\Theta_h$ and $G=G_*$.  The first fact is a theorem.  The second is a reading of Postulate~I that the present text does not derive; without it the cells are disphenoids with a two-parameter shape freedom and metricity returns to the selected status of paragraph~4.  Appendix~J.10 draws the consequences: the CDT bare couplings become predictions rather than tuned inputs, and the phase of the native point becomes a definite question.

Capacity transport cannot replace this reading.  Every capacity observable has the form $I_{\rm geom}\otimes O$ (Appendix~Q.4), while shape mismatch is a geometric datum, so no condition built from bandwidth transport, face labels, or face holonomies can enforce shape matching.

\paragraph{5. Local composition fixes the Regge form, not its coefficient.}
In a piecewise-flat metric geometry, curvature is concentrated on triangular hinges $h$ with area $A_h$ and deficit or boost angle $\Theta_h$.  Let the leading local hinge amplitude be a continuous phase $W(A,\Theta)\in U(1)$.  Invariance under subdivision of a hinge and composition of successive curvature wedges requires
\[
W(A_1+A_2,\Theta)=W(A_1,\Theta)W(A_2,\Theta),
\]
\[
W(A,\Theta_1+\Theta_2)=W(A,\Theta_1)W(A,\Theta_2).
\]
The two continuous Cauchy equations imply
\[
\boxed{W(A,\Theta)=e^{i\alpha A\Theta}}
\]
for a real constant $\alpha$.  Local composition therefore fixes the bilinear Regge form,
\[
S_{\rm geom}=\hbar\alpha\sum_hA_h\Theta_h,
\]
but it does not determine $\alpha$.  Matching this branch to the already-fixed Newton normalization sets
\[
\boxed{
\alpha=\frac{c^3}{8\pi\hbar G_*}
=\frac{1}{8\pi L_*^2}.}
\]
The resulting action and its channel form are
\[
\boxed{
S_R=\frac{c^3}{8\pi G_*}\sum_hA_h\Theta_h,}
\]
\[
\boxed{
\frac{S_R}{\hbar}
=\frac{\ln2}{2\pi}\sum_hN_h\Theta_h
=\gamma_*\sum_hJ_h\Theta_h.}
\]
The last equation is the central microscopic normalization identity.  Each cut channel supplies $\ln2$ of boundary entropy, and bending its associated area through $\Theta$ supplies the phase $(\ln2/2\pi)\Theta$ on the matched Einstein branch.  Proper-EPRL asymptotics provide a controlled witness of the same Regge phase; they do not turn the composition argument into an independent derivation of its coefficient.

\paragraph{6. Smooth metric limit.}
For a sequence of Regge geometries approaching a smooth metric,
\[
2\sum_hA_h\Theta_h
=\int d^4x\sqrt{-g}\,R
+O\!\left(a^2\int d^4x\sqrt{-g}\,\mathcal R^2\right),
\]
where $a$ is the coarse lattice spacing and $\mathcal R^2$ denotes the appropriate curvature-squared combination.  Hence
\[
S_R
=\frac{c^3}{16\pi G_*}
\int d^4x\sqrt{-g}\,R
+S_{\rm higher}.
\]
Allowing the volume term gives
\[
\boxed{
S_{\rm grav}
=\frac{c^3}{16\pi G_*}
\int d^4x\sqrt{-g}\,(R-2\Lambda)
+S_{\rm higher}.}
\]
The fixed-measure argument of Appendix~O determines how $\Lambda$ enters the ordinary branch; it is logically separate from the local curvature derivation above.

The stronger infrared uniqueness statement is conditional.  If the selected phase has one metric, locality at leading derivative order, diffeomorphism covariance, block-scale local Lorentz invariance, and only two massless helicity-$2$ modes, then the two-derivative nonlinear action is Einstein--Hilbert plus a cosmological term.  The cited Han-type and proper-EPRL results witness the Regge/Einstein and two-helicity sectors in their controlled regimes.  They are not one theorem for the full decorated action.  In particular,
\[
\boxed{
\begin{gathered}
\text{local Lorentz invariance at block scale is selected and witnessed,}\\
\text{not derived from the primitive }SU(2)\text{ cell}.
\end{gathered}}
\]
The ordinary infrared logic and the capacity-side compatibility theorem do not require a proof of every spin-foam refinement.  A claim that a particular proper or causal EPRL vertex is the faithful block image does, however, retain the open inter-cell fusion and host-convergence tests of H.10.

The uniqueness statement stops at two derivatives.  Gauge-invariant, parity-even quasilocal operators suppressed by powers of $kL_{\rm block}$ leave the infrared coefficient chain unchanged, and the finite $1680$ count does not fix their Wilson coefficients.  The corresponding TT kernel has the schematic expansion
\[
K_{\rm eff}(k)=Z_Gk^2P_{\rm TT}+c_4L_{\rm block}^2k^4+\cdots,
\]
so the controlled Einstein regime is $kL_{\rm block}\ll1$ and the first ordinary-branch ultraviolet deviations are expected at the block scale defined in H.10, not automatically at one primitive length.  In four dimensions, Lovelock uniqueness and massless spin-two self-coupling select Einstein--Hilbert only after locality, one-metric covariance, and the two-helicity Lorentzian phase are supplied \citep{Lovelock1971,Deser1970}.  Therefore
\[
\boxed{
S_{\rm higher}=\sum_i c_i\,\mathcal O_i,
\qquad c_i\text{ not fixed by the present postulates}.}
\]
Pure-gravity curvature-squared terms are field-redefinition or topological at this order, while matter moves the corresponding freedom into contact operators.  Predictions at the holonomy-turnover or other Planckian scales are consequently conditional on the stated minimal completion $c_i=0$ beyond the lock-generated tower; the infrared Einstein universality class is not.

\paragraph{7. The transfer defect becomes a metric source.}
Let a stable localized defect have transfer energy
\[
E_d=m_dc^2.
\]
Over a proper-time interval $\Delta\tau$, its Lorentzian transfer phase is
\[
\mathcal A_d(\Delta\tau)
=\exp\!\left(-\frac{iE_d\Delta\tau}{\hbar}\right).
\]
Composition along a worldline therefore gives
\[
\boxed{
S_d[z,g]=-m_dc^2\int d\tau=-m_dc\int ds.}
\]
At leading derivative order, this is also the unique local reparametrization-invariant scalar action built from one metric and the worldline tangent.  Its metric variation produces
\[
\boxed{
T_d^{\mu\nu}(x)
=\frac{m_dc}{\sqrt{-g}}
\int d\tau\,u^\mu u^\nu
\delta^{(4)}\!\bigl(x-z(\tau)\bigr).}
\]
The ordinary action is consequently
\[
\boxed{
S_{\rm ord}
=\frac{c^3}{16\pi G_*}
\int d^4x\sqrt{-g}\,(R-2\Lambda)
-\sum_dm_dc^2\int d\tau_d
+\cdots,}
\]
and variation of the same metric gives
\[
\boxed{
G_{\mu\nu}+\Lambda g_{\mu\nu}
=\frac{8\pi G_*}{c^4}T_{\mu\nu}.}
\]

For propagating charged leptons, the field-theory completion of the same
ordinary branch uses the shell-grade multiplet
$\Psi=(\psi_0,\psi_1,\psi_2)^T$ and the positive mass operator of
Appendix~I.1:
\[
\boxed{
S_{g\Psi}=\frac{c^3}{16\pi G_*}
\int d^4x\sqrt{-g}\,(R-2\Lambda)
+\int d^4x\sqrt{-g}\,
\bar\Psi\!\left(i\hbar c\,\gamma^ae_a{}^\mu D_\mu
-c^2\widehat M\right)\!\Psi .}
\]
Here $D_\mu$ is the spin- and gauge-covariant derivative defined by the same
tetrad and charged connection.  The kinetic operator is diagonal in shell
grade at this order. With the global $(-,+,+,+)$ signature and the displayed
$i\gamma^\mu D_\mu$ convention, the Clifford algebra is
$\{\gamma^\mu,\gamma^\nu\}=-2g^{\mu\nu}I$. In a local
inertial frame, a mass eigenstate obeys
\[
(\gamma^\mu p_\mu-m_Nc)\psi_N=0,
\qquad
\boxed{E_N^2=c^2\mathbf p^2+m_N^2c^4},
\]
where the second equation follows by multiplying by
$\gamma^\nu p_\nu+m_Nc$ and using the Clifford relation.  Metric variation
of this same action gives the Dirac stress tensor in the Einstein equation
above.  The eigenvalue $m_N$ therefore enters both the inertial dispersion
relation and the stress tensor obtained by metric variation.  The localized WKB limit is the worldline
action $S_d$.

This gives a single leading-order dynamics for the displayed lepton spectrum,
relativistic propagation, and ordinary gravity on the selected Lorentzian
branch.  Its covariant Dirac kinetic term is an infrared completion.  The
native decorated vertex must still generate that term, reproduce its speed
$c$, couple through the record-oriented frame of paragraph~2c, and preserve the precision
mass ratios under its radiative corrections.

Appendix~O's fixed covariant event measure supplies a conditional microscopic reading of the proper-time composition.  For a defect worldtube of invariant four-volume $\mathcal V_4=A_\perp L$, the premise $\mathcal V_4=v_*N_{\rm evt}$ gives, at fixed transverse profile,
\[
N_{\rm evt}\propto L=c\tau.
\]
Equal transfer phase per event then produces
\[
\prod_{n=1}^{N_{\rm evt}}e^{-i\delta\varphi_d}
=\exp\!\left(-\frac{i}{\hbar}m_dc^2\tau\right).
\]
For a Poisson covariant count,
\[
\mathbb E[N_{\rm evt}]=\frac{A_\perp L}{v_*},
\qquad
\operatorname{Var}N_{\rm evt}=\mathbb E[N_{\rm evt}],
\qquad
\frac{\delta\tau}{\tau}=O(N_{\rm evt}^{-1/2}).
\]
This explains how microscopic event cadence can count invariant proper time, conditional on the fixed-measure premise.  Covariance and one-metric coupling fix the worldline form; the event measure is its proposed microscopic realization, not an additional proof of Lorentz invariance.  Variation with respect to the path gives $u^\nu\nabla_\nu u^\mu=0$, while variation of the same functional with respect to the metric gives $T_d^{\mu\nu}$.  One functional therefore fixes both inertial and gravitational mass.

\paragraph{8. Newtonian limit and the weak equivalence principle.}
For
\[
ds^2=-\left(1+\frac{2\Phi}{c^2}\right)c^2dt^2
+\left(1-\frac{2\Psi}{c^2}\right)d\mathbf x^2,
\]
a slowly moving defect obeys
\[
d\tau
=dt\left[1+\frac{\Phi}{c^2}
-\frac{v^2}{2c^2}+O(c^{-4})\right].
\]
Therefore
\[
S_d
=\int dt\left[-m_dc^2+\frac12m_dv^2-m_d\Phi+\cdots\right].
\]
The same transfer coefficient multiplies the inertial and gravitational terms:
\[
\boxed{m_{\rm inertial}=m_{\rm gravitational}=E_d/c^2.}
\]
There is no independent gravitational charge to tune.  The linearized $00$ equation yields
\[
\boxed{\nabla^2\Phi=4\pi G_*\rho,}
\]
and ordinary matter with negligible anisotropic stress gives
\[
\boxed{\Phi=\Psi}
\]
at leading order.  These are the Newtonian force law and baseline no-slip lensing of Appendix~D, obtained here as consequences of the common metric action.

The three-dimensional commutant
$\operatorname{End}_{SU(2)}(V_0\oplus V_1\oplus V_2)\simeq\mathbb C^3$
does not introduce three gravitational charges.  That space describes the internal marked record fiber; the source tensor structure is fixed instead by the covariant worldline functional.  Its $V_0\oplus V_1\oplus V_2$ channel content can dress the response, but cannot replace the single coefficient $m_d$ in $T_d^{\mu\nu}$.

\paragraph{9. Renewal cannot add a second ordinary gravitational pole.}
Let $\Pi_*$ be the stationary projector on the renewed capacity operator space and define
\[
P=I_{\rm geom}\otimes\Pi_*,
\qquad Q=I-P.
\]
For an arbitrary bounded within-tick geometry--capacity evolution $\mathcal U^\dagger$, the renewal step is
\[
\mathcal T^\dagger=P\mathcal U^\dagger.
\]
On $P\mathcal K\oplus Q\mathcal K$ it has the exact block form
\[
\boxed{
\mathcal T^\dagger
=\begin{pmatrix}A&B\\0&0\end{pmatrix},}
\qquad
(\mathcal T^\dagger)^n
=\begin{pmatrix}A^n&A^{n-1}B\\0&0\end{pmatrix}.
\]
Consequently
\[
\boxed{
\operatorname{spec}(\mathcal T^\dagger)\setminus\{0\}
=\operatorname{spec}(A)\setminus\{0\}.}
\]
Within-tick capacity fluctuations may dress the retained evolution through $B$, but a replaceable capacity mode has no autonomous persistence after renewal.  The result holds without requiring geometry--capacity mixing to vanish.  A conservation-protected collective variable already contained in $A$ is not excluded; that is why the infrared galactic response is graded separately from ordinary vacuum gravity.

\paragraph{10. Closure status.}
The complete chain is:
\[
\boxed{
\begin{aligned}
\text{selected relational cells}
&\Rightarrow SU(2)\text{ flux geometry},\\
j_{\rm primitive}=3
&\Rightarrow J=3N,\\
N\ln2
&\Rightarrow A_N=4\ln2\,L_*^2N,\\
\text{shape matching}
&\Rightarrow\text{piecewise-flat metric},\\
\text{local composition}+G_*
&\Rightarrow S_R=\frac{c^3}{8\pi G_*}\sum_hA_h\Theta_h,\\
\text{smooth Lorentzian two-helicity phase}
&\Rightarrow S_{\rm EH},\\
E_d+\text{proper-time composition}
&\Rightarrow S_d=-m_dc^2\int d\tau,\\
\delta(S_{\rm EH}+S_d)=0
&\Rightarrow G_{\mu\nu}+\Lambda g_{\mu\nu}
=\frac{8\pi G_*}{c^4}T_{\mu\nu},\\
\text{renewal triangularity}
&\Rightarrow\text{no additional replaceable capacity pole}.
\end{aligned}}
\]
Exact results in this chain include the representation decompositions, the
selected equivariant lock, maximal-channel fusion, the locked coherent phase
space, the channel--area algebra after continuum normalization, the transfer
scale, and renewal triangularity. The selected completion supplies the
geometric factor, shape-matched one-metric support, and ultraviolet propagator.
The one-metric requirement and local composition fix the leading
Regge/Einstein universality class; the channel area and $G_*$ fix its
normalization. Proper-EPRL, Han-type, and area-Regge results independently
witness that the selected coarse data realize this class. H.10 then proves
the TT infrared limit and gives an exact fixed-$j=3$ perfect-action pullback.
The coherent bond and angular rules of Q.11 provide one explicit set of
additional geometric data, with the pair interaction in H.12 and metric
count dynamics in N.10--N.12. Selection of those data from the capacity
premises alone remains unproved; their physical phase and constrained
state space are separate tests.

\subsection*{H.10 Gaussian GFT--geometric-relative-entropy bridge and reduced embedding audits}

At any finite regulator, positive dressed and reference Hessians on a common
gauge-reduced support determine a unique normalized Gaussian determinant and
positive relative-covariance functional. This is an algebraic result. In the
factorized completion, the finite capacity sector fixes the equilibrium
projector, records, defects, and coefficient chain. The geometric
spin-foam/GFT sector supplies the gauge-reduced spin-2 geometry. The selected
maximal-fusion ray and the Han-type refinement results of Appendix~B and
Section~23 provide a controlled semiclassical witness. Proper-EPRL one-sector
asymptotics supply a separate compatibility check. The exact operator
identification for the relative-information functional, the coupled
geometry--capacity audits below, and the all-state nonperturbative Lorentzian
continuation remain conditional.

\paragraph{A. Physical Hessians and relative covariance.}
Let $\mathcal H'_{\rm phys}$ be a finite-dimensional regulated real fluctuation space after gauge fixing and removal of the same exact zero modes in the equilibrium and dressed backgrounds. For a complex GFT field this means its real Nambu form, with the conjugate components included in the physical trace. On that common support, take
\[
S_0^{(2)}[\varphi]
=\frac12\langle\varphi,\mathbb K_0\varphi\rangle,
\qquad
S^{(2)}[\varphi]
=\frac12\langle\varphi,\mathbb K\varphi\rangle,
\qquad
\mathbb K_0>0,
\quad
\mathbb K>0.
\]
Whitening by the equilibrium Hessian defines
\[
\Theta
=\mathbb K_0^{-1/2}\mathbb K\mathbb K_0^{-1/2}>0,
\qquad
\mathcal G
=\Theta^{-1}
=\mathbb K_0^{1/2}\mathbb K^{-1}\mathbb K_0^{1/2}>0.
\]
Equivalently, with $C_0=\mathbb K_0^{-1}$ and $C=\mathbb K^{-1}$,
\[
\mathcal G=C_0^{-1/2}CC_0^{-1/2}.
\]
Every prime below means this identical gauge and zero-mode prescription. Equality of the two physical supports is a hypothesis to be tested in the microscopic spectrum; it is not supplied by the determinant identity.

\paragraph{B. Finite-regulator determinant theorem.}
The equilibrium-normalized Gaussian functional is
\[
\boxed{
\widehat Z_{\rm G}
=
\frac{\displaystyle\int_{\mathcal H'_{\rm phys}}D\varphi\,
e^{-\frac12\langle\varphi,\mathbb K\varphi\rangle}}
{\displaystyle\int_{\mathcal H'_{\rm phys}}D\varphi\,
e^{-\frac12\langle\varphi,\mathbb K_0\varphi\rangle}}
=(\det{}'\Theta)^{-1/2}
=(\det{}'\mathcal G)^{1/2}.
}
\]
Consequently,
\[
\boxed{
2\ln\widehat Z_{\rm G}
=-\operatorname{Tr}'\ln\Theta
=\operatorname{Tr}'\ln\mathcal G.
}
\]
The proof is just the finite Gaussian integral
\[
Z[\mathbb K]
=(2\pi)^{N_{\rm phys}/2}(\det{}'\mathbb K)^{-1/2}.
\]
Taking the ratio cancels the measure normalization, while multiplicativity of the determinant under the congruence defining $\Theta$ gives the displayed result. The proof requires neither $[\mathbb K_0,\mathbb K]=0$ nor simultaneous diagonalization. Positivity defines the real Gaussian measure and the principal operator logarithm. The finite-dimensional Lebesgue integral, orthogonal change of variables, determinant normalization, and trace-log identity are kernel checked in \path{Entropic/AppH10}. The physical support and the choice of Hessian remain inputs.
The real/Nambu convention fixes the exponent $-1/2$; a complex-coordinate calculation gives the same result only after its conjugate variables and Jacobian are counted consistently in $N_{\rm phys}$ and $\operatorname{Tr}'$.

\paragraph{Sources under Gaussian elimination.}
On the same finite physical support, split the variables into retained $x$ and eliminated $y$ with action
\[
S=\tfrac12x^{\mathsf T}Ax+x^{\mathsf T}By
+\tfrac12y^{\mathsf T}Dy-j_x^{\mathsf T}x-j_y^{\mathsf T}y,
\qquad D>0.
\]
Completing the square in $y$ gives the exact effective action
\[
\begin{aligned}
S_{\rm eff}={}&\tfrac12x^{\mathsf T}(A-BD^{-1}B^{\mathsf T})x
-(j_x-BD^{-1}j_y)^{\mathsf T}x\\
&-\tfrac12j_y^{\mathsf T}D^{-1}j_y
+\tfrac12\ln\det D+\text{constant}.
\end{aligned}
\]
The source shift, source contact, and determinant all contribute when their arguments vary. Retaining them makes successive eliminations agree with joint elimination. Dropping them preserves only part of the quadratic operator and can change its physical response. The completed-square and minimum identities, together with the finite positive Gaussian integral and its measure normalization, have kernel-checked proofs. Source-dependent physical supports and continuum limits require their stated additional conditions.

\paragraph{C. Exact quadratic-source Legendre transform and Gaussian KL divergence.}
Write the whitened dressed Hessian as
\[
\Theta=I-U>0,
\qquad
U=U^\dagger,
\]
and define the normalized connected functional
\[
W[U]
=\ln\widehat Z_{\rm G}[U]
=-\frac12\operatorname{Tr}'\ln(I-U).
\]
Its matrix derivative is
\[
2\frac{\delta W}{\delta U}
=(I-U)^{-1}
=\mathcal G.
\]
The convex conjugate with the corresponding quadratic-source normalization is
\[
\Gamma_{\rm rel}[\mathcal G]
=\sup_{\substack{U=U^\dagger\\ I-U>0}}
\left\{
\frac12\operatorname{Tr}'(\mathcal G U)-W[U]
\right\}.
\]
Stationarity gives $U=I-\mathcal G^{-1}$. Substitution therefore yields the exact identity
\[
\boxed{
\Gamma_{\rm rel}[\mathcal G]
=\frac12\operatorname{Tr}'
\left[\mathcal G-I-\ln\mathcal G\right].
}
\]
If $g_i>0$ are the eigenvalues of $\mathcal G$, then $g_i-1-\ln g_i\ge0$, with equality only at $g_i=1$. Hence
\[
\Gamma_{\rm rel}\ge0,
\qquad
\Gamma_{\rm rel}=0
\Longleftrightarrow
\mathcal G=I
\Longleftrightarrow
C=C_0
\]
on the common support. In probabilistic language this is precisely
\[
\Gamma_{\rm rel}
=D_{\rm KL}\!\left(
\mathcal N(0,C)\,\Vert\,\mathcal N(0,C_0)
\right).
\]
Thus the positive mismatch functional in Sections~25 and~29.9 is the convex conjugate of the normalized Gaussian trace logarithm.

\paragraph{D. Conditional relation to geometric relative entropy.}
Bianconi's geometric-relative-entropy construction uses a positive metric-response operator $G_B=\Theta_B^{-1}$ and the same operator-convex mismatch form \citep{Bianconi2026}. If a specified decorated GFT gluing tensor derives the operator identification
\[
\Theta_B=\Theta,
\qquad
G_B=\mathcal G,
\]
then the Bianconi trace-log and Legendre functionals are the Gaussian effective functionals of the dressed GFT covariance relative to equilibrium. Under that still-conditional identification, $G_B$ can be read in this theory as a composite covariance or susceptibility ratio, slaved to the condensate, metric, matter, and relational boundary data. The functional identity proved above is exact; the microscopic gluing and operator map are not. Bianconi's independent derivative dynamics are not imported into the metric-only ordinary branch, and this paper does not attribute the composite reading to her construction itself.

\paragraph{E. Euclidean sign and continuation.}
The Euclidean one-loop contribution obtained by integrating out the Gaussian fluctuation is
\[
\boxed{
\Delta\Gamma_{E,{\rm eff}}
=-\ln\widehat Z_{\rm G}
=\frac12\operatorname{Tr}'\ln\Theta
=-\frac12\operatorname{Tr}'\ln\mathcal G.
}
\]
This sign is not the sign of the positive Legendre/KL mismatch. The logarithm of a multiplicity or partition ratio and the Euclidean effective free-energy contribution carry opposite signs, while the convex conjugate measures the nonnegative cost of changing the covariance. A common Lorentzian parent must fix the analytic continuation and the sign of its retarded influence kernel. The determinant theorem and positive Legendre identity do not depend on that unresolved continuation.

\paragraph{F. Conditional local derivative expansion.}
Suppose the regulated operator is quasilocal in a geometric cell basis,
\[
\Theta_{vv'}=\delta_{vv'}\Theta_v+\mathcal R_{vv'},
\]
with $\mathcal R$ derivative-suppressed on fields varying slowly compared with $L_*$. Only under that assumption may the exact trace be organized as
\[
2\ln\widehat Z_{\rm G}
=\nu_*
\int d^4x\,\sqrt{|g|}\,
\big[-\operatorname{tr}_{\rm phys}\ln\Theta(x)\big]
+\mathcal O\!\left(
L_*\nabla\Theta,
L_*^2\nabla^2\Theta
\right).
\]
Here $\operatorname{tr}_{\rm phys}$ is the trace over the local gauge-reduced geometric fiber and $\nu_*$ is its continuum block density. Neither is fixed by the $1680$-state capacity probability distribution. In the factorized completion they follow from the geometric gluing and blocking calculation, while the equilibrium capacity factor contributes its rank-one stationary transfer projector in probability space, and matter/record insertions couple through the equivariant lock. The finite determinant is exact before locality is assumed; applying it as a local continuum density remains conditional on the quasilocal derivative expansion.

\paragraph{G. Low-curvature and static Einstein corollary.}
If the covariant operator supplied by the embedding has the small-response form
\[
\Theta=I-\widetilde U,
\qquad
\widetilde U
=\beta\widetilde R-\alpha\widetilde M
+\mathcal O(R^2,RM,M^2,\nabla^2),
\]
where $\widetilde R$ and $\widetilde M$ are the dimensionless curvature and matter insertions on the physical fiber, then
\[
-\operatorname{Tr}'\ln(I-\widetilde U)
=\operatorname{Tr}'\widetilde U
+\frac12\operatorname{Tr}'\widetilde U^2+\cdots.
\]
For the controlled geometric branch the EPRL/Regge refinement results of Appendix~B supply the Einstein--Hilbert spin-2 support independently of this trace-log expansion. The extra question here is whether the relative-covariance operator itself inherits the same covariant fiber trace, Ward identity, and normalization when capacity defects dress that geometry. The already-controlled longitudinal reduction fixes its static target:
\[
I_{\rm cap}^{\rm static}
=\int dt\,d^3x
\left[
-\frac{\gamma}{2}(\nabla\delta S)^2
+\kappa\rho\,\delta S
\right],
\qquad
\delta S=-\frac{2S_\infty}{c^2}\Phi,
\]
\[
G=\frac{c^2\kappa}{8\pi\gamma S_\infty},
\qquad
\boxed{I_{\rm cap}^{\rm static}=Z_S I_{\rm Newton}.}
\]
The ordinary longitudinal capacity response is not a propagating scalar pole. Under the capacity--metric dictionary it is the Hamiltonian-constraint combination of the Einstein metric, so its static Green function may be unscreened without adding independent Cauchy data. A microscopic condensate realization therefore passes the ordinary-gravity audit when every scalar fluctuation beyond that constraint response is gapped, relaxational, pure gauge/constraint, or exactly source-orthogonal in the physical retarded kernel. The radiative gauge-reduced support must remain transverse-traceless. This spectrum statement follows from the geometric completion, not from the trace-log identity by itself.

\paragraph{H. Marked-record support-preservation lemma.}
The pair contraction and the seven-channel recurrence are different quantities. From H.9,
\[
u=\frac{8\eta_*}{49}
=0.00487621949\ldots,
\qquad
\alpha=\sqrt{1-u}
=0.9975589108\ldots,
\]
\[
\alpha^{21}=0.9499693634\ldots,
\qquad
r=e^{-7g_{\rm share,eff}}
=2.77557\times10^{-23}.
\]
The positivity of the marked edge factor follows from $0<u<1$, not from the exponential smallness of $r$. On the established seven-channel one-mark branch, integrating out the 21 internal present/history Majorana blocks multiplies the remaining bosonic geometric transfer by the strictly positive scalar $\alpha^{21}$. Taking the branch-independent internal trace of the nine-state hard-core marked fiber gives
\[
\boxed{
\zeta_*
=9e^{-g_{\rm share,eff}}\alpha^{21}
=0.00512358448\ldots>0.
}
\]
Therefore the integrated marked-record sector, by itself, neither removes nor creates a bosonic geometric mode. This is a decoration-side support lemma. On the factorized vacuum branch the geometric TT support is fixed by the controlled spin-foam continuum; for non-equilibrium defect backgrounds one must still verify that dressing does not introduce new gauge-reduced zero modes, alter the constraints, or create an additional sourced pole.

\paragraph{I. Exact internal-fiber identity and limited ADM reading.}
The primitive fusion identity of H.9 is
\[
V_{3/2}\otimes V_{3/2}
=V_3\oplus V_0\oplus V_1\oplus V_2.
\]
Removing the selected maximal channel gives
\[
Q
=(V_{3/2}\otimes V_{3/2})\ominus V_3
=V_0\oplus V_1\oplus V_2,
\qquad
\dim Q=1+3+5=9.
\]
Independently, the marked response fiber obeys
\[
\mathcal H_{\rm mark}
=V_1^P\otimes V_1^H
=V_0\oplus V_1\oplus V_2.
\]
The intertwiner already constructed in H.9 therefore proves equivalence of the two internal fibers, with one copy of each irrep. Under spatial rotations, the ADM shift and symmetric spatial metric decompose as
\[
h_{0i}\in V_1,
\qquad
h_{ij}\in V_0\oplus V_2.
\]
The nine marked states have the spatial-rotation content of the three shift
components and six spatial-metric components. The lapse $h_{00}$ remains a
separate scalar constraint multiplier. This $SU(2)$ identity does not make
the marked fiber the metric. It supplies the blockwise representation target
for the equivariant geometry--capacity lock. Inter-cell covariance, the lapse
constraint, the diffeomorphism Ward identity, and TT propagation remain in
the geometric factor; the marked factor can only dress those responses.

\paragraph{J. Equilibrium capacity normalization and physical support.}
The stationary distribution has two related but distinct representations. In the detailed-balance/symmetrized probability space,
\[
|v\rangle
=\sum_{b\in\mathcal B}\sqrt{p_{\eta_*}(b)}\,|b\rangle,
\qquad
P_{\rm stat}=|v\rangle\langle v|
\]
is the rank-one stationary projector. The actual renewed quantum register is instead the diagonal density operator already fixed in H.8,
\[
\rho_*=\sum_b p_{\eta_*}(b)|b\rangle\langle b|,
\qquad
\operatorname{Tr}\rho_*=1,
\]
and the replacement channel is $\mathcal E_*(X)=\rho_*\operatorname{Tr}X$.
Here $P_{\rm stat}$ is rank one in transfer/probability space, whereas
$\rho_*$ has the full support required by faithful renewal.

After the complete homogeneous capacity free energy is centered as in Section~25, the normalized homogeneous equilibrium reference contributes unit capacity weight.  Thus for that reference state,
\[
\boxed{
\mathcal A_{\rm full}^{\rm ref}
=\mathcal A_{\rm geom}^{\rm ref}
}
\]
up to the already-divided gauge/measure normalization.  For a generic empty geometry $T$, however,
\[
\mathcal A_{\rm full}[T]
=\mathcal A_{\rm geom}[T]\,\widehat Z_{\rm cap}[T],
\]
where centering removes the source-independent homogeneous extensive term.
Geometry-dependent finite contributions need not equal unity. The capacity
Hilbert space is also not one-dimensional. Instead, the renewed local closure
register has one stationary superoperator direction: $\mathcal E_*$
annihilates every traceless perturbation of that replaceable register. The
renewed register therefore cannot furnish an additional coherent pole. Any
gapless capacity pole must lie in a retained or conservation-protected
collective sector, such as the transport carrier tested in Section~24.
Persistent marked fibers must satisfy the source-overlap conditions above.

Schematically, on the controlled linearized witness branch, the combined continuum results imply a physical Hessian of the form
\[
\mathbb K_{\rm phys}(k)
\sim Z_G k^2P_{\rm TT}
+O(a^2k^4)+O(\Lambda_{\rm sf}^{-1}k^2),
\]
where the separate terms summarize the Einstein spin-2 limit and the known derivative/large-spin suppressions rather than quote one theorem for a single decorated proper-EPRL model.
The first explicit length-metric correction in the area-Regge continuum is $O(a^2C^2)$ \citep{DittrichKogios2023}. The independent area-metric RG analysis of Ref.~\citep{BorissovaEtAl2025} supports the decoupling of nonmetric shape-mismatch modes in the infrared, but it is not used here as a theorem about the full EPRL flow. Its failure to make parity generic is instead handled by the exact microscopic orientation exchange imposed in Appendix~Q.

\paragraph{K. The selected geometric hosts and why their specification matters.}
The capacity construction fixes the $j=3$ carrier, its maximal-fusion ray, and the infrared normalization; it does not fix a unique ultraviolet geometric propagator.  Two distinct selected completions are therefore kept separate.

\emph{Covariant fixed-facet host.}
Every microscopic geometric face is a spin-three facet with its own relational incidence identity and its own dual loop.  A coarse face is a partition into $N$ such facets.  Refinement creates new interfaces; it never replaces one boundary cycle by many indistinguishable copies carrying a common holonomy variable.  Schematically,
\[
\boxed{
f_{\rm coarse}=\bigsqcup_{a=1}^{N}f_a,
\qquad j(f_a)=3,
\qquad g_{f_a}\text{ retained separately}.}
\]
This distinction answers the known stack-localization objection.  Repeating the same face on one cycle produces
\[
\omega_h\sim[\lambda\tau_k(g_h)]^{p_h},
\qquad
\operatorname{Re}S_h
=p_h\ln\frac{|\tau_k(g_h)|}{d_k^2}\le0,
\]
with equality only at flat holonomy.  It therefore localizes onto flat $SU(2)$ connections as $p_h\to\infty$, even at fixed $k$ \citep{HanStack2026}.  A coarse-label host does not evade the problem merely by restricting $J$ to multiples of three:
\[
\sum_{J\in3\mathbb N}d_J\chi_J(\phi)
=\frac13\sum_{r=0}^{2}
\sum_Jd_J\chi_J(\phi)e^{2\pi irJ/3},
\]
which is supported on flat or $\mathbb Z_3$-twisted-flat class angles.  By contrast, a relational facet partition has weight $\prod_aw(g_{f_a})$ on distinct loops and contains no power of a single $\tau_k(g)$, so that specific mechanism does not directly apply.  This is not a proof that every RG flow avoids a topological phase; the non-topological clause in the convergence condition below remains essential.

Han's many-small-spin construction supplies only the single-vertex half of the desired blocking statement: a generalized vertex with parallel fixed-spin facets has the same large total spin $J=\sum_a j_a=3N$ and the corresponding Euclidean Regge critical data \citep{HanSmallSpin2019}.  It does not show that many primitive four-simplices coarse-grain into one block vertex.  The grades are therefore
\[
\boxed{
\begin{aligned}
&F_N|3,\mathbf n\rangle^{\otimes N}=|3N,\mathbf n\rangle
&&\text{exact kinematics},\\
&\text{parallel facets within one vertex}
&&\text{controlled Euclidean witness},\\
&\text{inter-cell vertex fusion}
&&\text{open dynamics}.
\end{aligned}}
\]
The parallel-normal restriction is exactly the maximal channel $V_{3N}\subset V_3^{\otimes N}$.  In the amplitude realization
\[
\mathcal A_{\rm fine}=\mathcal A_NF_N,
\qquad
\|Q_N\psi\|_{\rm phys}=0,
\]
so the nonmaximal geometric complement is null rather than a second set of propagating face modes.  The marked $V_0\oplus V_1\oplus V_2$ records remain entirely in $\mathcal H_{\rm cap}$, consistent with the factorization of Appendix~Q.  Proper and causal EPRL amplitudes are then block-scale witnesses on $V_{3N}$, not primitive definitions; their face dimension $6N+1$ is the fused image of the primitive $d_3=7$ carrier.  Their grade remains conditional on the open inter-cell line above.

\emph{Canonical compression freezes the area spectrum.}
Let $q(\Psi)=\sum_I\|C_I\Psi\|^2$ be a closable positive master form and let $\mathcal H_*$ be the closure of spin networks whose edge spins lie in $3\mathbb N$.  Restricting and closing the form defines a positive self-adjoint compression $\mathbf M_*$ with
\[
\langle\Psi,\mathbf M_*\Psi\rangle
=\sum_I\|C_I\Psi\|^2,
\qquad
\ker\mathbf M_*=\ker\mathbf M\cap\mathcal H_*.
\]
This is compression in form sense, not the sum of separately projected constraints:
\[
\boxed{
\langle\Psi,P\mathbf MP\Psi\rangle
-\sum_I\|PC_IP\Psi\|^2
=\sum_I\|(I-P)C_IP\Psi\|^2\ge0.}
\]
For the standard fundamental-holonomy regularization, $C_I^\dagger C_I$ changes any edge spin by at most two units.  Projection back to $3\mathbb N$ therefore retains only the zero shift.  Hence
\[
\boxed{
[\mathbf M_*,\widehat A_e]=0
\quad\text{for every edge }e,}
\]
and its symbol on the locked coherent ray is independent of the conjugate twist $\xi$.  Areas are frozen: the compressed standard master constraint is excluded as a geometric host, rather than listed as an admissible candidate.

The lock itself selects the regularization.  Replace the fundamental trace by the normalized spin-three trace
\[
\operatorname{Tr}_{1/2}\!\bigl(h[h^{-1},V]\bigr)
\longmapsto
\frac{c_3}{3\cdot4\cdot7}
\operatorname{Tr}_{3}\!\bigl(h[h^{-1},V]\bigr),
\]
with $c_3$ fixed by the same classical limit.  Locked transitions $0,\pm3$ then survive, and $\mathbf M_*^{(3)}$ carries $\Delta J=0,\pm3,\pm6$ with the nontrivial symbol
\[
\mathbf M_*^{(3)}\supset
\frac{2}{9}\bigl(1-\cos3\xi\bigr)
=\xi^2-\frac34\xi^4+O(\xi^6).
\]
Thus
\[
\langle\mathbf M_*^{(3)}\rangle
=\mathbf M[J,\xi,\mathbf n]
+O(\xi^4)+O(\bar N^{-1}),
\]
and the canonical alternative is specifically the spin-three-regularized master constraint, not the compression of the standard one.  Its fixed holonomy scale and $2\pi/3$ periodicity are the structural $\mathbb Z_3$ of the locked phase space.  Apparent higher-representation zeros at $\xi+2\pi/3$ are the same locked state, not additional kernels.  This canonical completion and the covariant fixed-facet completion are both admissible, but their equivalence is not assumed.

\paragraph{L. Exact gapped embedding and perfect-action pullback.}
A one-scale nearest-neighbor alignment does not justify discarding nonmaximal block channels.  On a connected primitive graph $G$,
\[
H_G=\kappa\sum_{(ab)\in E(G)}
\bigl(9-\mathbf J_a\!\cdot\!\mathbf J_b\bigr)
\]
has the maximal-spin ray as its kernel, but on the one-magnon sector
\[
H_G=3\kappa L_G,
\qquad
\Delta_G=3\kappa\lambda_2(G)\longrightarrow0
\]
for growing local blocks.  Nearest-neighbor alignment therefore leaves a gapless spin-wave complement.  On a periodic cube of linear size $L$, $\lambda_2(L_G)=2[1-\cos(2\pi/L)]$, so
\[
\frac{\Delta_L}{\kappa}=6\left(1-\cos\frac{2\pi}{L}\right)\simeq\frac{12\pi^2}{L^2},
\]
equal to $1.757$, $0.457$, $0.115$, and $0.0289$ at $L=8,16,32,64$. The same local model orders: on $\mathbb Z^d$ with $d\ge3$ its magnetized positive-temperature state exists for every spin (Appendix~B and \citep{OpenAIHeisenberg2026}). Local Heisenberg alignment thus supplies phase ordering with Goldstone spin waves, and the hierarchical parent below supplies the uniform gap.

\emph{Gauss-projection lemma.} The magnetization above is an order parameter for the global $SU(2)$ acting on unprojected face vectors, while the geometric $SU(2)$ is a local frame redundancy. For one geometric tetrahedron the physical space is $\operatorname{Inv}_{SU(2)}(V_3^{\otimes4})$, of dimension seven, with projector $P_c$. Each face flux $J_{cf}^a$ is a spin-one tensor operator, and the Wigner--Eckart theorem gives
\[
\boxed{P_cJ_{cf}^aP_c=0.}
\]
For two neighboring cells and every fixed holonomy $R(U)$,
\[
(P_c\otimes P_{c'})\,\mathbf J_{cf}\!\cdot\!R(U)\mathbf J_{c'f}\,(P_c\otimes P_{c'})=0,
\]
the $V_3$ form of the vanishing transported bilinear between gauge-projected spin-$\tfrac32$ slots in Appendix~P.8. Softening the projection by the Gauss penalty $H_0=\lambda(\mathbf G_c^2+\mathbf G_{c'}^2)$ and perturbing with $V=-g\,\mathbf J_{cf}\!\cdot\!R\mathbf J_{c'f}$ leaves the same result at second order. Each flux carries its cell from total spin zero to total spin one, so every intermediate state costs $4\lambda$, and rotational invariance gives $P_cJ_{cf}^aJ_{cf}^bP_c=\tfrac13j(j+1)\delta^{ab}P_c=4\delta^{ab}P_c$ for $j=3$. Hence $PV^2P=48g^2P$ and
\[
\boxed{H_{\rm eff}^{(2)}=-\frac{12g^2}{\lambda}\,I,}
\]
a constant that leaves the seven intertwiner states degenerate. Single-face Heisenberg alignment with exact local Gauss invariance therefore produces no physical shape ordering. The seven shape states carry their own spectrum: in the recoupling basis $J_{12}=0,\ldots,6$,
\[
P_c(\mathbf J_1\!\cdot\!\mathbf J_2)P_c=\tfrac12\left[J_{12}(J_{12}+1)-24\right],
\qquad
\operatorname{spec}=\{-12,-11,-9,-6,-2,3,9\}.
\]
Physical geometric interactions act on these gauge-invariant Casimir and shape variables through several faces, holonomy dynamics, or the spin-foam vertex.

The tree fusion map supplies a uniformly gapped parent instead.  At each node $v$ joining spins $j_L,j_R$, define
\[
Q_v=j_Lj_R-\mathbf J_L\!\cdot\!\mathbf J_R,
\qquad
q_J=\frac12\bigl[(j_L+j_R)(j_L+j_R+1)-J(J+1)\bigr].
\]
Nested-block Casimirs commute, so
\[
\boxed{
H_T=\kappa\sum_{v\in T}Q_v,
\qquad
\ker H_T=V_{3N},
\qquad
\operatorname{gap}H_T\ge6\kappa.}
\]
The kernel is independent of the fusion tree by multiplicity one; only the decoupled excited spectrum remembers the tree.

Let $F_N:V_3^{\otimes N}\to V_{3N}$ be the normalized coisometry, $F_NF_N^\dagger=I$, and $Q_N=I-F_N^\dagger F_N$.  For any positive self-adjoint coarse operator $K_N^{\rm coarse}$, define
\[
\boxed{
K_N^{\rm prim}
=F_N^\dagger K_N^{\rm coarse}F_N
+\Delta_NQ_N,
\qquad \Delta_N\ge6\kappa.}
\]
On the physical ray,
\[
F_NK_N^{\rm prim}F_N^\dagger=K_N^{\rm coarse},
\qquad
F_N(K_N^{\rm prim}-z)^{-1}F_N^\dagger
=(K_N^{\rm coarse}-z)^{-1}.
\]
Moreover,
\[
\operatorname{spec}K_N^{\rm prim}
=\operatorname{spec}K_N^{\rm coarse}\cup\{\Delta_N\}_{Q_N}.
\]
The construction also preserves coupled matter--geometry dynamics.  At a
finite regulator let $\mathcal H_d$ carry the persistent defect variables,
set
\[
\widetilde F_N=F_N\otimes I_d,
\qquad
\widetilde Q_N=I-\widetilde F_N^\dagger\widetilde F_N,
\]
and let $K_{N,g d}^{\rm coarse}$ be a positive self-adjoint operator on
$V_{3N}\otimes\mathcal H_d$.  It may contain metric-dependent defect
transport and backreaction, so no geometry--matter factorization is needed.
The joint pullback
\[
\boxed{
K_{N,g d}^{\rm prim}
=\widetilde F_N^\dagger K_{N,g d}^{\rm coarse}\widetilde F_N
+\Delta_N\widetilde Q_N}
\]
satisfies
\[
\widetilde F_N(K_{N,g d}^{\rm prim}-z)^{-1}\widetilde F_N^\dagger
=(K_{N,g d}^{\rm coarse}-z)^{-1}.
\]
For a deparametrized self-adjoint Hamiltonian the same block identity gives
\[
\widetilde F_Ne^{-itH_{N,g d}^{\rm prim}/\hbar}
\widetilde F_N^\dagger
=e^{-itH_{N,g d}^{\rm coarse}/\hbar}.
\]
The interacting spectrum, propagation, backreaction, and their correlations
are therefore retained together on the embedded physical sector.

The flat complement $\Delta_NQ_N$ may be replaced by the hierarchical $H_T$ above; its excited eigenvalues then spread while its kernel and uniform lower gap are unchanged.  Poles, residues, the $i0$ prescription, and the zero fiber therefore pull back exactly.  Associativity of maximal fusion,
\[
F_{N+M}=F_{(N,M)}\circ(F_N\otimes F_M)
\quad\text{up to phase},
\]
makes $I_N:=F_N^\dagger$ a projective family.  For $\mathcal A_{\rm fine}:=\mathcal A_NF_N$ one has $\mathcal A_{\rm fine}\circ I_N=\mathcal A_N$, which is cylindrical consistency at the amplitude level. This construction proves that a positive primitive operator can carry a chosen healthy coarse Einstein sector. It does not show that the strictly local primitive weights dynamically generate that perfect operator. In four dimensions the corresponding perfect action is generically nonlocal \citep{DittrichKaminskiSteinhaus2014}; here the nonlocality is confined to the first reliable semiclassical block,
\[
\boxed{
L_{\rm block}:=\sqrt{A_0}
=\sqrt{a_{\rm ch}N_0}
=\sqrt{4\ln2\,N_0}\,L_*,}
\]
where $N_0$ is the smallest facet number for which $J_0=3N_0$ supports the required proper/causal Lorentzian asymptotics at the stated numerical tolerance.  This definition follows from the physical coarse area $A_0=a_{\rm ch}N_0$; $\sqrt{3N_0}L_*$ would instead confuse the spin label with an area in units of $L_*^2$.

The present calculation does not derive $N_0$.  For numerical planning it uses the provisional working bracket $J_0\sim10$--$30$, hence $N_0\sim4$--$10$ and
\[
\boxed{L_{\rm block}\sim3.3\text{--}5.3\,L_*.}
\]
Existing analytic and single-vertex numerical asymptotic studies motivate this test range \citep{EngleVilenskyZipfel2016,Gozzini2021}, but do not establish it for the decorated vertex. A direct calculation must replace the bracket with a measured tolerance bound. Provisionally, $kL_*\ll0.2$--$0.3$ is sufficient for $kL_{\rm block}\ll1$, so ordinary-branch ultraviolet corrections should appear at a few $L_*$ rather than exactly at $L_*$. This is an estimated locality range, not a closed prediction.

\paragraph{Selected metric-Regge host theorem.}
The perfect-action construction above permits a stronger statement than an
existence witness.  Choose the covariant host to be a regular, shape-matched
metric-Regge block sequence on the ordinary branch.  Four-dimensional
linearized Regge calculus has propagating gauge-invariant lattice curvature
degrees of freedom \citep{Hoehn2015}, and its weak-field limit recovers the
standard continuum theory \citep{RocekWilliams1981}.  On the regular
centrally subdivided lattice, the modern area-Regge calculation gives the
gauge-reduced length-metric symbol
\citep{DittrichKogios2023}
\[
\boxed{
\mathbb K_n^{\rm TT}(k)
=Z_Gk^2P_{\rm TT}
+O(a_n^2k^4),
\qquad Z_G>0,}
\]
where the metric/shape-matched branch is the length-Regge subspace and does
not retain the additional area-metric modes.  Here
\(a_n=L_{\rm block}/L_{{\rm obs},n}\to0\) is an infrared scaling parameter;
the physical microscopic cutoff \(L_*\) is not sent to zero.

This symbol also completes the operator statement.  For
\(z\notin\mathbb R\), identify the regular-lattice TT spaces with their
Fourier interpolation in the continuum TT Hilbert space and write
\[
R_n(k;z)=[\mathbb K_n^{\rm TT}(k)-z]^{-1}.
\]
For every fixed physical momentum,
\[
R_n(k;z)\longrightarrow
[Z_Gk^2P_{\rm TT}-z]^{-1},
\qquad
\lVert R_n(z)\rVert\le\frac1{|\operatorname{Im}z|}.
\]
Dominated convergence first gives resolvent convergence on TT Schwartz wave
packets of compact Fourier support, and the uniform bound extends it by
density:
\[
\boxed{
\mathbb K_n^{\rm TT}
\xrightarrow{\rm s.r.}
Z_G(-\Delta)P_{\rm TT}.}
\]
Equivalently, positive functional calculus gives for the deparametrized
Hamiltonians
\[
H_n=c\sqrt{\mathbb K_n^{\rm TT}/Z_G}
\xrightarrow{\rm s.r.}
c|\mathbf k|P_{\rm TT}.
\]
Thus \(\inf\sigma(H_{\rm TT})=0\), genuine TT states occur arbitrarily close
to zero energy, and the retarded kernel has two massless helicity-\(2\) poles
with positive residue.  This is the relevant masslessness criterion; a
nonzero pointwise density at the spectral threshold is not required.

The exact primitive pullback carries the result to fixed microscopic
spin:
\[
\mathbb K_n^{\rm prim}
=F_n^\dagger\mathbb K_n^{\rm TT}F_n+H_{T,n},
\qquad
H_{T,n}F_n^\dagger=0,
\qquad
H_{T,n}|_{Q_n}\ge6\kappa,
\]
\[
\boxed{
F_n(\mathbb K_n^{\rm prim}-z)^{-1}F_n^\dagger
=(\mathbb K_n^{\rm TT}-z)^{-1}.}
\]
The physical maximal-fusion ray therefore inherits the Einstein resolvent
exactly, while every nonmaximal fusion mode remains uniformly gapped.  Denote
this closed branch theorem by
\[
\boxed{
(H_{\rm TT})\quad
\text{metric-Regge branch}\Longrightarrow
\text{massless two-helicity Einstein TT limit}.}
\]
Together with the perfect-action pullback, it also proves that a compatible
fixed-\(j=3\) primitive realization exists.

\paragraph{Free-graviton Fock representation of the TT branch.}
The limit operator also fixes the Hilbert space of the linearized branch. At nonzero momentum write $P_{ij}=\delta_{ij}-\hat k_i\hat k_j$ and
\[
P^{\rm TT}_{ij,kl}=\tfrac12\left(P_{ik}P_{jl}+P_{il}P_{jk}-P_{ij}P_{kl}\right),
\qquad
\operatorname{Tr}P_{\rm TT}=2 .
\]
The one-particle space and Hamiltonian are
\[
\mathfrak h_1=L^2(\mathbb R^3,d^3k)\otimes\mathbb C^2,
\qquad
h_1=c|\mathbf k|,
\]
with the two-dimensional factor spanned by the TT polarizations, and second quantization gives
\[
\boxed{
\mathcal F_{\rm grav}=\Gamma_s(\mathfrak h_1),
\qquad
H_{\rm grav}=d\Gamma(c|\mathbf k|).}
\]
The Euclidean TT action $\tfrac12Z_G\int[c^{-2}(\partial_\tau h)^2+(\nabla h)^2]$ has covariance
\[
C_E(\omega,\mathbf k)=\frac{c^2}{Z_G}\,\frac{P_{\rm TT}}{\omega^2+c^2|\mathbf k|^2},
\qquad
C_E(\tau,\mathbf k)=\frac{c}{2Z_G|\mathbf k|}\,e^{-c|\mathbf k||\tau|}P_{\rm TT}.
\]
For a real TT test tensor $f$ supported at $\tau>0$ and its time reflection $\Theta f$,
\[
\langle\Theta f,C_Ef\rangle
=\int\frac{c\,d^3k}{2Z_G|\mathbf k|}
\left\|\int_0^\infty e^{-c|\mathbf k|\tau}P_{\rm TT}\hat f(\tau,\mathbf k)\,d\tau\right\|^2\ge0,
\]
so the gauge-reduced Gaussian TT theory is reflection positive, and Osterwalder--Schrader reconstruction returns $\mathcal F_{\rm grav}$ with its vacuum and $H_{\rm grav}$. On the selected metric-Regge branch,
\[
\boxed{
(H_{\rm TT})+\text{positive Gaussian transfer}
\Longrightarrow
\text{two-helicity massless graviton Fock representation}.}
\]
The statement is linear and concerns the TT sector; the continuum realization of the full nine-state, history, and matter theory is a separate problem. A microscopic version at finite regulator requires the spacetime symbol
\[
\mathbb K_a^{\rm TT}(\omega,\mathbf k)\longrightarrow Z_G\left(\frac{\omega^2}{c^2}+|\mathbf k|^2\right)P_{\rm TT}
\]
together with reflection positivity of the discrete transfer matrix, so that the temporal and spatial coefficients converge to one light cone. The continuum construction of the two-dimensional $O(3)$ model, which proves convergence of the lattice Schwinger functions, Euclidean invariance, reflection positivity, and clustering before reconstruction \citep{OpenAIO3Continuum2026}, is the template for that sequence. For the Gaussian TT sector the characteristic functionals $\exp[-\tfrac12\langle f,C_af\rangle]$ converge when $\mathbb K_a^{\rm TT}$ converges in the strong-resolvent sense and has a uniform positive lower bound on a common invariant TT subspace. Restricting Fourier support away from $\mathbf k=0$ supplies the infrared domain; the common lower bound controls $\mathbb K_a^{-1}$.

This is not yet a phase theorem for the native theory.  The Hessian result is
restricted to the selected shape-matched branch, and the perfect operator is
constructed from the healthy coarse operator.  Neither controls which
configurations dominate a nonperturbative sum.  Once a definite local
geometric tensor and measure have been supplied, the remaining covariant host
condition would be
\[
\boxed{
(H_{\rm phase})\quad
\mu_{\rm native}\xrightarrow{\rm IR}
\mu_{\rm metric\text{-}Regge/Einstein},}
\]
meaning that the resulting measure must concentrate at long distance on
shape-matched, non-topological metric configurations with a gapless
positive-residue TT sector, rather than on a BF-like, degenerate, bifurcated,
or otherwise non-Einstein phase.  The underdetermination theorem below shows
that the present premises do not yet define \(\mu_{\rm native}\), so this
condition is not presently a well-posed derived calculation.  An
unrestricted Lorentzian EPRL sum and uniqueness of the selected ultraviolet
propagator are not claimed.

\paragraph{Capacity-side closure on the displayed factorized branch.}
The capacity audits listed above are consequences of the finite decoration and the host identities above, not additional spectral assumptions.  After the internal marked fiber is traced,
\[
\boxed{
G_{\rm dec}=\zeta_*G_{\rm geom+particle},
\qquad \zeta_*=9e^{-g_{\rm share,eff}}\alpha^{21}>0.}
\]
Therefore the decoration changes no denominator: its poles are those of geometry plus the worldline particle, and the signs of all residues are preserved.  It introduces neither a ghost nor a tachyon.  Renewal triangularity separately removes autonomous replaceable-register poles, while the tensor lift below preserves the positive physical inner product and causal boundary value.

There are two equivalent bookkeeping pictures.  In the dilated picture the complete update is unitary, with
\[
V^\dagger V
=P+\frac19\sum_{a=1}^{9}W_a^\dagger W_aQ
=I,
\]
so Hermiticity and positivity are explicit. In the traced Euclidean picture,
the corresponding requirement is an Osterwalder--Schrader audit of the
complete transfer kernel, including its marked subblocks. Vacuum renewal is
counted per invariant event and the defect cadence is proper time. The
capacity sector contains no vector $n^\mu$ from which to construct a leading
preferred-frame curvature operator. Refresh-first ordering also gives exact
trigger independence:
\[
(\operatorname{id}_R\otimes\mathcal E_*)(\rho_{R,\rm ren})
=\rho_R\otimes\rho_*,
\qquad
\langle K^2\rangle_{\rm cond}=\langle K^2\rangle_{\eta_*}.
\]
Thus the factorized ordinary branch closes the capacity pole, residue, positivity, causality, preferred-frame, and refresh-order audits.  A different nonfactorized ultraviolet realization must reproduce these identities.

\paragraph{Local Einstein-phase stability under finite capacity dressing.}
The exact factorized theorem above is stronger than what is needed for a
general gauge-compatible decoration.  Suppose an undecorated host already
has the gauge-reduced infrared kernel
\[
K_g(k)=Z_Gk^2P_{\rm TT}+O(k^4L_*^2),
\qquad Z_G>0,
\]
with no additional massless gravitational pole.  If the capacity sector has
no conservation-protected vacuum pole, its static geometric correction is
local or quasilocal, and its TT correction is bounded by the host stiffness,
then the Ward identity forbids a TT mass term:
\[
\delta K_{\rm TT}(0)=0,
\qquad
\delta K_{\rm TT}(k)
=\delta Z_Gk^2P_{\rm TT}+O(k^4L_*^2).
\]
Consequently, for \(\lvert\delta Z_G\rvert<Z_G\),
\[
\boxed{
K_{\rm full}(k)
=Z_G^{\rm eff}k^2P_{\rm TT}+O(k^4L_*^2),
\qquad
Z_G^{\rm eff}=Z_G+\delta Z_G>0,}
\]
and no new gravitational zero mode appears.  Thus a finite, pole-free,
gauge-compatible, perturbatively bounded decoration is locally stable around
an already established Einstein phase.  This is not a claim that a positive
\(Z_{\rm cap}[T]\) leaves every saddle unchanged, that it cannot renormalize a
curvature coefficient, or that it is phase-neutral globally.  The microscopic
normalization \(G_*\) fixes the target renormalized coefficient; a coupled
realization must still pass the normalization audit below.  On the displayed
factorized vacuum branch the stronger identity \(G_{\rm dec}=\zeta_*G_{\rm
geom+particle}\) establishes these hypotheses exactly.  For a nonfactorized
embedding they are explicit acceptance conditions.

\paragraph{Schur-complement form of the stability condition.}
The bound on the capacity correction has an explicit quadratic form. Write the coupled gauge-reduced kernel as
\[
\mathbb K(k)=
\begin{pmatrix}
Z_Gk^2P_{\rm TT}&B(k)\\
B^\dagger(k)&M_{\rm cap}(k)
\end{pmatrix},
\qquad
M_{\rm cap}(k)\ge\Delta_{\rm cap}I\ \text{near }k=0 .
\]
The primitive complement of Appendix~B has $\Delta_{\rm cap}=m_{2|3}^2=3\kappa$, and the hierarchical parent of paragraph~L has $\Delta_{\rm cap}\ge6\kappa$. Eliminating the capacity block gives
\[
K_{\rm eff}(k)=Z_Gk^2P_{\rm TT}-B(k)M_{\rm cap}^{-1}(k)B^\dagger(k).
\]
A nonzero $B(0)$ produces a $k^0$ TT term, which the Ward identity forbids, so $B(0)=0$. With $B(k)=kB_1+O(k^2)$,
\[
K_{\rm eff}(k)=k^2\left[Z_GP_{\rm TT}-B_1M_{\rm cap}^{-1}(0)B_1^\dagger\right]+O(k^3),
\qquad
\left\|B_1M_{\rm cap}^{-1}(0)B_1^\dagger\right\|\le\frac{\|B_1\|^2}{\Delta_{\rm cap}} .
\]
A sufficient condition for the absence of an additional massless gravitational pole is therefore
\[
\boxed{
B(0)=0,
\qquad
\|B_1\|^2<Z_G\Delta_{\rm cap}.}
\]
Under it the capacity modes renormalize $Z_G$ and leave the TT pole structure unchanged. A proposed nonfactorized vertex is tested by computing $B(0)$ and $B_1$ from its gauge-reduced Hessian and comparing $\|B_1\|^2$ with $Z_G\Delta_{\rm cap}$. The factorized branch has $B=0$, consistent with $G_{\rm dec}=\zeta_*G_{\rm geom+particle}$.

\paragraph{Geometric-host underdetermination theorem.}
The preceding stability result removes the capacity factor from the host
question; it does not make the host follow from the capacity premises.  Write
the complete factorized vertex as
\[
\boxed{
\mathcal A_v[\mathcal G]=\mathcal G_v\,\mathcal D_v,}
\qquad
\mathcal D_v
=W_*\Pi_Q
e^{-g_{\rm share,eff}n_d-\Lambda_dn_d(n_d-1)}
\prod_{m<m'}\widetilde U_{mm'}.
\]
The finite factor \(\mathcal D_v\) fixes the \(V_3\) carrier, stationary
measure, renewal projector, marked transfer, and dimensional coefficient
chain.  The geometric gluing tensor \(\mathcal G_v\) supplies propagation and
curvature.  If \(\mathcal G_v\) is replaced by another gauge-covariant tensor
\(\mathcal G'_v\) with the same \(V_3\) boundary representation and
equivariant-lock interface, every closed capacity invariant is unchanged.
This follows because the gauge action factorizes and all capacity operators
commute with its geometric factor.  In addition, the exact vacuum lift
gives
\[
\boxed{
\mathcal H_{\rm phys}^{\rm full,vac}
\cong\mathcal H_{\rm phys}^{g},}
\]
so on the displayed branch the full vacuum phase is the phase of the chosen
geometric tensor.

The manuscript already contains a concrete separation.  The sharp primitive
object
\[
\mathcal G_{v,\sharp}^{(0)}
=\mathcal P_v^{\rm inv}
\left[\bigotimes_{(ab)\in E(v)}P_3^{(ab)}\right]
\mathcal P_v^{\rm inv}
\]
enforces the representation and intertwiner support but contains no
curvature phase.  The selected Einstein completion instead carries
\[
\mathcal G_v^{E}\sim
\exp\!\left[
\frac{i}{\hbar}\frac{c^3}{8\pi G_*}
\sum_hA_h\Theta_h
\right].
\]
Both admit the same finite capacity decoration.  Requiring the geometric
factor to admit an EPRL/Regge branch is therefore a host acceptance condition,
not a consequence of the capacity construction.

Even the existence of the Regge saddle does not determine phase dominance.
At finite regulator, a local positive gauge-invariant factor
\[
\exp\!\left[
\lambda\sum_h f(L_*^2\mathcal R_h)
\right],
\qquad
f(x)=0\ \text{near }x=0,
\]
can leave the smooth saddle, its perturbative Hessian, \(G_*\), and the
two-helicity branch unchanged while exponentially favoring rough
configurations at large volume.  Hence
\[
\boxed{
\text{existence of an Einstein saddle}
\not\Longrightarrow
\text{Einstein-phase dominance}.}
\]
It follows that the proposition
\[
\boxed{
\text{the present finite-capacity premises uniquely determine }
\mathcal G_v}
\]
is excluded.  The perfect-action pullback proves that an
Einstein-compatible fixed-\(j=3\) completion exists; it does not remove this
selection freedom.  A hosted minimal completion may select that operator and
inherit \((H_{\rm TT})\) exactly. The coherent completion in Q.11 supplies
an absolute angular tensor and correlated vacuum; H.12 and N.10--N.12
supply a pair parent, metric operator, and ordering. These selected data
do not remove the underdetermination theorem. Their local definition
makes further tests possible, while a full refinement measure and
positive physical state space remain necessary for
\((H_{\rm phase})\).

\paragraph{Area-law tensor networks and the host.}
A unique gapped ground state of a local finite-dimensional Hamiltonian obeys an entanglement area law, and on the square lattice it admits a projected entangled-pair approximation of polynomial bond dimension \citep{OpenAIAreaLaw2026,OpenAIPEPS2026}. A parent Hamiltonian $\sum_f\kappa Q_f+H_{\rm closure}+H_{\rm geom}+H_{\rm lock}$ with such a ground state would determine a local tensor network and hence a gluing tensor. The geometric sector violates those hypotheses in three ways. The local alignment model has the $(6N+1)$-dimensional ground multiplet $V_{3N}$ and the gap $\Delta_L\simeq12\pi^2\kappa/L^2\to0$. The hierarchical parent $H_T$ restores a uniform gap, keeps the multiplet, and is hierarchical in place of finite-range. The gravitational sector is gapless by construction through its TT branch. The area-law and tensor-network theorems therefore apply to gapped capacity subblocks and leave the underdetermination of $\mathcal G_v$ unchanged.

\paragraph{What a nonfactorized geometry--capacity embedding must preserve.}
The external geometric witness leaves the following interface checks for any
alternative fully coupled vertex. They are already satisfied by construction
on the displayed factorized branch, but a proposed nonfactorized realization
is acceptable only if all conditions hold on the same gauge-reduced background:
\begin{enumerate}
\item \emph{Spectrum/source overlap:} the capacity decoration introduces no additional source-coupled zero pole; beyond the Einstein scalar constraint, scalar fluctuations are gapped, relaxational, pure gauge/constraint, or exactly source-orthogonal.  In condensate realizations with a nonzero relational-time phase gradient, density--phase mixing must be checked in the full Hessian rather than assumed to vanish from a zero-gradient truncation.
\item \emph{Constraint and Ward structure:} coupling the capacity/record factor leaves the geometric constraint algebra and diffeomorphism Ward identity intact, including in defect backgrounds.
\item \emph{Normalization:} the static source, bulk curvature, and conditional cut use the same renormalized coefficient. Recovering the matched value $G_*=c^3L_*^2/\hbar$ from the physical event and channel calibration is an additional test.
\item \emph{Transfer/physical-inner-product structure:} for a Euclidean/GFT transfer realization, the complete decorated transfer, and not merely its finite marked subblocks, must satisfy the required positivity/reflection property; self-adjointness alone is not sufficient.  For a directly Lorentzian covariant realization, the corresponding requirement is positivity of the rigging form together with consistent causal gluing, rather than Osterwalder--Schrader positivity.
\item \emph{Lorentzian continuation:} the retarded poles have the correct causal prescription and the capacity dressing does not create ghosts, tachyons, or preferred-frame radiative modes.
\item \emph{Geometric observable support:} the dressed covariance or Hessian $T_a$ must approach the geometric von Neumann algebra $M_{\rm geom}$. One universal hyperreflexivity constant $C$ holds for all unital von Neumann algebras \citep{OpenAIKadison2026}, so
\[
\operatorname{dist}(T_a,M_{\rm geom})\le C\sup_{e\in\mathcal P(M_{\rm geom}')}\|(1-e)T_ae\|,
\]
and vanishing off-block mixing under every commutant projection, $\sup_e\|(1-e)T_ae\|\to0$, gives $T_a\to M_{\rm geom}$ in norm. On the factorized branch the capacity trace is a positive scalar dressing and the condition holds trivially.
\end{enumerate}
These are interface-preservation tests, not a demand that the finite capacity factor rederive the Einstein graviton sector.  Failure would reject the alternative realization, not reopen the exact factorized capacity theorem.

\paragraph{M. Exact rigging-map factorization and the Lorentzian endpoint.}
The controlled spin-foam result is stronger than a purely kinematic embedding and weaker than an exact all-state covariant projector. Refs.~\citep{Han2017,HanHuangZipfel2019} give a coupled large-spin/refinement regime in which Regge solutions approach the Einstein sector; in the controlled linearized construction the low-energy excitations give all smooth linearized Einstein solutions and hence exactly the two graviton helicities. This establishes the identity and existence of the ordinary semiclassical gravitational branch.

The unmodified EPRL edge map is not idempotent, so the formal sum over
arbitrary foams does not define a projector as written; this obstruction is
independent of the detailed spin-foam model
\citep{PerezSpinFoam2013,ThiemannZipfel2014}. The Riemannian ``Warsaw''
modification yields an idempotent edge projector, but no corresponding general
Lorentzian theorem is available. Ref.~\citep{BrunoEtAl2026} derives a
rigging-map structure only after assuming the appropriate distributional
continuum limit. The raw Lorentzian EPRL state sum is therefore not treated
as an exact physical projector here.

\paragraph{Exact capacity-decorated lift theorem.}
The capacity contribution to the rigging problem can be evaluated exactly
when the mixed renewal state is represented on its operator space. Let
\[
\omega_*(A)=\operatorname{Tr}(\rho_*A),
\qquad
\rho_*=\sum_bp_{\eta_*}(b)|b\rangle\langle b|
\]
be the faithful stationary state of the renewed quantum register and let $\mathcal K_{\rm cap}$ be the GNS/Hilbert--Schmidt completion of the capacity observable algebra with inner product
\[
\langle A,B\rangle_{\omega_*}
=\operatorname{Tr}(\rho_*A^\dagger B).
\]
The Heisenberg renewal map is
\[
\Pi_*(A)
\equiv\mathcal E_*^\dagger(A)
=\omega_*(A)I.
\]
It obeys
\[
\boxed{
\Pi_*^2=\Pi_*,
\qquad
\Pi_*^\dagger=\Pi_*,
\qquad
\operatorname{Ran}\Pi_*=\mathbb C I
}
\]
with the adjoint taken in the $\omega_*$ inner product: indeed
$\langle A,\Pi_*B\rangle_{\omega_*}
=\overline{\omega_*(A)}\omega_*(B)
=\langle\Pi_*A,B\rangle_{\omega_*}$.
The rank-one projector $|\sqrt p\rangle\langle\sqrt p|$ used in the detailed-balance probability representation is the commutative-coordinate image of this same one-dimensional stationary transfer direction; neither object says that the quantum state $\rho_*$ is pure.

Let
\[
\eta_g:D_g\longrightarrow D_g^*
\]
be any positive geometric rigging map.  On the empty, unmarked vacuum sector define
\[
\boxed{
\eta_{\rm full}(\phi\otimes A)[\psi\otimes B]
=\eta_g(\phi)[\psi]\,
\langle A,\Pi_*B\rangle_{\omega_*}.
}
\]
Sesquilinearity, the constraint-annihilation property, reality, and positivity are inherited from $\eta_g$ and the orthogonal projector $\Pi_*$.  Quotienting its null space and completing gives
\[
\boxed{
\mathcal H_{\rm phys}^{\rm full,vac}
\cong
\mathcal H_{\rm phys}^{g}\otimes\operatorname{Ran}\Pi_*
\cong
\mathcal H_{\rm phys}^{g}.
}
\]
This theorem is exact and regulator independent for the equilibrium capacity decoration.  It says that renewal adds no vacuum constraint solution or physical pole once a geometric rigging map exists.

Renewal is also compatible with an arbitrary coupled dynamics of the
persistent sector.  Let $S$ denote the joint geometry--defect system and $R$
the replaceable closure register.  In the Schr\"odinger picture,
$\mathcal E_*(X)=\rho_*\operatorname{Tr}X$, so every joint state satisfies
\[
\boxed{
(\operatorname{id}_S\otimes\mathcal E_*)(\rho_{SR})
=\operatorname{Tr}_R(\rho_{SR})\otimes\rho_* .}
\]
For any joint geometry--defect unitary $U_S$, a renewal step followed by the
physical evolution gives
\[
(U_S\otimes I_R)(\operatorname{id}_S\otimes\mathcal E_*)(\rho_{SR})
(U_S^\dagger\otimes I_R)
=U_S\operatorname{Tr}_R(\rho_{SR})U_S^\dagger\otimes\rho_* .
\]
The replacement changes no reduced state or internal correlation of the
persistent geometry--defect system.  A microscopic dilation must additionally
satisfy the no-which-path condition: its discarded output may not acquire a
cell address correlated with a defect trajectory.  Thus renewal and coupled
propagation coexist exactly once the persistent defect dynamics is supplied;
the renewal identity does not derive that dynamics.

\paragraph{Exact finite-regulator group averaging.}
There is also an exact regulated construction.  Parametrized GFT makes the reparametrization constraint explicit; for a self-adjoint regulated geometric Hamiltonian $H_\Lambda$, introduce an auxiliary clock pair $(\chi,p_\chi)$ and
\[
C_\Lambda=p_\chi+H_\Lambda.
\]
Then refined algebraic quantization gives
\[
\boxed{
\eta_\Lambda
=2\pi\delta(C_\Lambda)
=\int_{\mathbb R}d\tau\,e^{i\tau C_\Lambda}
}
\]
as a distributional group average, with the physical inner product obtained by inserting $\delta(C_\Lambda)$ \citep{CalcinariGielen2025}.  At any finite spin/mode/particle-number regulator a Hermitian truncation is self-adjoint, so this construction is exact even with interactions retained inside that finite matrix.  Ref.~\citep{CalcinariGielen2025} proves the clock-neutral/deparametrized equivalence explicitly for the free GFT; the interacting cutoff removal is not supplied by that paper and is not imported here as a theorem.

The same exact mechanism is available canonically: the Lorentzian LQG master-constraint programme constructs a positive symmetric master constraint with a Friedrichs self-adjoint extension, and its direct-integral zero spectral fiber defines the physical Hilbert space \citep{HanMa2006,HanThiemann2010}.  This is an exact nonperturbative canonical witness for how the geometric rigging map can exist.  It is not a proof that the raw Lorentzian EPRL foam sum equals that canonical projector.

\paragraph{Clock-test-space convergence theorem.}
Strong-resolvent convergence suffices on a clock test space.  Let $H_\Lambda$ be self-adjoint on the regulated GFT Hilbert space and define the deparametrized constraint
\[
C_\Lambda=p_\chi+H_\Lambda
\quad\text{on}\quad
L^2(\mathbb R_\chi)\otimes\mathcal H_{\rm GFT}.
\]
For product test vectors $\Phi=f\otimes\phi$ and $\Psi=g\otimes\psi$ with $f,g\in\mathcal S(\mathbb R)$, group averaging is exactly
\[
\boxed{
\eta_\Lambda(f\otimes\phi)[g\otimes\psi]
=2\pi\langle\phi,h_{fg}(H_\Lambda)\psi\rangle,}
\qquad
h_{fg}(\lambda)
=\overline{\widehat f(-\lambda)}\,\widehat g(-\lambda).
\]
Indeed, $p_\chi$ and $H_\Lambda$ act on separate tensor factors, and the $\tau$ integral gives $2\pi\delta(p+\lambda)$ in clock-momentum space.  Since $f$ and $g$ are Schwartz, $h_{fg}\in C_0(\mathbb R)$.  Therefore
\[
\boxed{
H_\Lambda\xrightarrow{\rm s.r.}H
\quad\Longrightarrow\quad
\eta_\Lambda(\Phi)[\Psi]\longrightarrow\eta(\Phi)[\Psi].}
\]
The limiting form is positive,
\[
\eta(f\otimes\phi)[f\otimes\phi]
=2\pi\int_{\mathbb R}|\widehat f(-\lambda)|^2
\,d\mu_{\phi\phi}(\lambda)\ge0,
\]
and it annihilates $C$ on the test space.

The same Schwartz smearing supplies the causal boundary value.  The spectral density of $C_\Lambda$ on these vectors is the convolution
\[
\rho_\Lambda(E)
=\int
\overline{\widehat f(E-\lambda)}
\widehat g(E-\lambda)
\,d\mu_\Lambda(\lambda).
\]
It is uniformly bounded, Lipschitz in $E$, and integrable because the spectral measure is finite and the clock kernel is Schwartz.  Strong-resolvent convergence at fixed $\varepsilon>0$ followed by the Plemelj limit therefore gives
\[
\boxed{
G_F=(C-i0)^{-1}
=\operatorname{PV}C^{-1}+i\pi\delta(C),
\qquad
\operatorname{Im}G_F(\Phi,\Phi)
=\pi\rho_{\Phi\Phi}(0)
=\frac12\eta(\Phi)[\Phi]\ge0.}
\]
Thus density domination, the pointwise-$\rho(0)$ condition, and limiting absorption hold on the Schwartz clock test space and need not be assumed.  Only strong-resolvent convergence remains for this particular all-state rigging-map route.  It is a sufficient mathematical construction, not a requirement that the physical cutoff be removed.

The theorem applies to deparametrized evolution. The equation $p_\chi+H=0$ does not project onto $\ker H$, and $\chi$ must be a physical relational matter clock. If $\chi$ is the usual massless-scalar GFT clock, the low-energy spectrum contains two TT gravitational helicities and the clock scalar. It contains no additional capacity-gravity scalar. A condensate Bogoliubov mode identified with the clock must satisfy $c_s^2\to1$ as $k\to0$. A master constraint $\mathbf M\ge0$ defines a separate zero-fiber problem and cannot replace $H$ in the clock equation. The masslessness test for the selected branch is
\[
\boxed{
\begin{aligned}
(H_{\rm TT})\quad
&H_n\xrightarrow{\rm s.r.}H_{\rm Ein},
\qquad
\inf\sigma(H_{\rm Ein}|_{\rm TT})=0,\\[-1mm]
&H_{\rm Ein}|_{\rm TT}
=c|\mathbf k|+O(k^3L_{\rm block}^2),
\quad\text{with two positive-residue TT helicities}.
\end{aligned}}
\]
The selected metric-Regge theorem above establishes \((H_{\rm TT})\), and the
perfect-action identity pulls it back to an exact fixed-\(j=3\) realization.
It does not establish the native-measure condition \((H_{\rm phase})\).  The
distinct canonical completion retains its own open test,
\[
\boxed{
(H'_{\rm can})\quad
\mathbf M_\Lambda^{(3)}\xrightarrow{\rm s.r.}\mathbf M^{(3)}
\text{ with a non-topological zero fiber carrying TT Einstein support}.}
\]
Equivalence between the covariant and canonical completions is not assumed.
Fixed-complex $q$-deformed Lorentzian EPRL models supply finite
nonzero-$\Lambda$ regulators \citep{FairbairnMeusburger2012}; they do not
establish the native phase condition or regulate the empty \(\Lambda=0\)
branch universally.

The status is therefore
\[
\boxed{
\begin{gathered}
\text{Gaussian determinant bridge: exact},\\
\text{quadratic-source Legendre/KL mismatch: exact},\\
\text{marked support preservation: exact on the established branch},\\
\text{maximal-fusion blocking: exact},\\
\text{renewal triangularity: exact},\\
\text{selected-branch Einstein theorem }(H_{\rm TT})\text{: closed},\\
\text{Einstein-compatible fixed-}j=3\text{ completion: exists},\\
\text{capacity premises uniquely fix }\mathcal G_v\text{: excluded},\\
\text{derived native phase }(H_{\rm phase})\text{: undefined until }\mathcal G_v\text{ is fixed},\\
\text{capacity-decorated rigging-map lift: exact},\\
\text{finite-regulator group averaging: exact for self-adjoint truncations},\\
\text{clock-test positivity and causal boundary value: exact},\\
\text{canonical }(H'_{\rm can})\text{: open and not required by the selected covariant host}.
\end{gathered}
}
\]
The vacuum rigging-map problem contains no unresolved \emph{capacity} step.
The selected metric-Regge branch satisfies the Einstein TT infrared theorem
\((H_{\rm TT})\), and it has a compatible fixed-\(j=3\) primitive realization.
The host itself remains underdetermined: the current premises leave
\(\mathcal G_v\) and its geometric measure as selected data. A new geometric
principle must fix both before \((H_{\rm phase})\) becomes a well-posed derived
theorem about the nonperturbative measure. A canonical--covariant
correspondence would be a further result. The raw EPRL sum does not supply that
projector.

\paragraph{Candidate selection principle.}
The face-wise interpretation of the relational lock (Appendix~H.9a, paragraph~4a) supplies the one candidate already present in the postulates for the additional geometric principle named above.  It does not choose $\mathcal G_v$ among gauge-covariant tensors; it restricts the configuration space on which any such tensor acts to equilateral causal triangulations with fixed edge $a=2.530L_*$, on which the Regge weight is $(\ln2/2\pi)\sum_hN_h\Theta_h$ with no free coefficient.  Under this interpretation $\mu_{\rm native}$ is the CDT measure at couplings fixed by the theory, and $(H_{\rm phase})$ becomes the well-posed numerical question of Appendix~J.10: whether that ensemble, with the capacity weighting switched on, is extended with a gapless TT sector at its predicted point.  This interpretation remains a premise.

\subsection*{H.11 Coherence under renewal and reconstruction of the retained marked Hilbert sector}

\paragraph{The question.}
The renewed register contains no information about its input, whereas quantum interference is carried by off-diagonal operators. Coherence can survive only if the interfering degrees of freedom remain in the retained system and the discarded complement carries no path record. The calculation below first states this division exactly. It then asks whether the retained nine-state mark must be assigned a complex Hilbert space or whether that structure follows from the displayed fusion and Lorentzian action. It adds no fourth foundational postulate.

\paragraph{Exact renewed-output theorem.}
For
\[
\mathcal E_*(X)=\rho_*\operatorname{Tr}X,
\]
the Heisenberg dual is fixed by
\[
\operatorname{Tr}\!\left[\mathcal E_*(X)O\right]
=\operatorname{Tr}\!\left[X\mathcal E_*^\dagger(O)\right].
\]
Substitution gives
\[
\boxed{
\mathcal E_*^\dagger(O)
=\operatorname{Tr}(\rho_*O)I.
}
\]
Every later observable on the renewed output therefore pulls back to a multiple of the identity, and every traceless input operator is annihilated. Equivalently, for any reference $R$,
\[
(\operatorname{id}_R\otimes\mathcal E_*)(\rho_{RA})
=\rho_R\otimes\rho_*.
\]
The replacement channel is entanglement-breaking on the renewed register. A global dilation can still preserve the old state in its complementary output, as the history register in H.8 does. The theorem excludes coherence from the replaced output only; it does not exclude global coherence.

\paragraph{Complementary-record coherence theorem.}
Let $X$ denote the defect position and $D$ its retained internal marked fiber. For two freely transported branches, write the relevant part of a one-step isometry as
\[
V\bigl(|x\rangle_X|q\rangle_D|0\rangle_E\bigr)
=|F(x)\rangle_XU_D|q\rangle_D|e_x\rangle_E,
\]
and similarly for $y$. The environment $E$ denotes only what is discarded when the defect state is read. Directly tracing $E$ gives
\[
|x\rangle\langle y|
\longmapsto
\underbrace{\langle e_y|e_x\rangle}_{\gamma_{xy}}
|F(x)\rangle\langle F(y)|
\otimes U_D|q\rangle\langle q|U_D^\dagger.
\]
Thus
\[
\boxed{\mathcal C_{xy}(1)=\gamma_{xy}\mathcal C_{xy}(0),
\qquad \gamma_{xy}=\langle e_y|e_x\rangle.}
\]
For independent record factors over $n$ updates,
\[
\mathcal C_{xy}(n)
=\left(\prod_{k=1}^{n}\gamma_{xy}^{(k)}\right)\mathcal C_{xy}(0).
\]
The equation separates the two possible microscopic readings. If the vertex leaves a cell-addressed record, $|e_x\rangle=|x\rangle_E$ and $|e_y\rangle=|y\rangle_E$, then $\gamma_{xy}=0$ for $x\ne y$: position coherence dies in one update. If the persistent internal fiber remains inside $X\otimes D$ and the discarded free-renewal output is branch-independent, $|e_x\rangle=|e_y\rangle$, then $\gamma_{xy}=1$ and renewal produces no intrinsic position decoherence. Intermediate overlaps produce ordinary partial decoherence.

Calling a record ``worldline-attached'' is not by itself sufficient. A spatially localized record that moves along two different worldlines is still a which-path record if it is discarded. In this paper the phrase means the stronger structure: the marked response fiber is retained as part of the coherently transported defect, while the discarded renewal register contains no branch address. Genuine interactions with detectors or an environment may, and ordinarily do, make the discarded states distinguishable.

\paragraph{What $Q\perp V_3$ does and does not prove.}
Primitive fusion gives the orthogonal decomposition
\[
V_{3/2}\otimes V_{3/2}=V_3\oplus Q,
\qquad Q=V_0\oplus V_1\oplus V_2,
\]
and hence
\[
P_QP_3=0,
\qquad \operatorname{Tr}P_Q=9,
\qquad \operatorname{Tr}P_3=7.
\]
If a discarded record actually resolves these two sectors, the same overlap theorem sets every $Q$--$V_3$ off-diagonal to zero. This gives exact marked/unmarked sector dephasing under a sector-resolving export. It is compatible with charge-sector superselection, but it is not yet a derivation of the full electric-charge superselection rule. That stronger statement requires the mark to be identified with the generator of a local $U(1)$ gauge symmetry and requires the physical observable algebra to preserve its Gauss-law sectors. Orthogonality of two representation subspaces alone does not force the environment to measure them.

\paragraph{Forced vertex constraint, not a new premise.}
The complete microscopic Hilbert space at a marked vertex must therefore be read as
\[
\mathcal H_v
=\mathcal H_{\rm ren}
\otimes\mathcal H_X
\otimes\bigl(\mathbb C|0\rangle\oplus Q\bigr),
\]
where $\mathcal H_{\rm ren}$ is replaced by $\rho_*$, while $\mathcal H_X\otimes Q$ is transported coherently. In the free sector the dilation of $W_*$ must satisfy
\[
\boxed{
\langle e_y|e_x\rangle=1
\quad\text{for unmeasured position branches }x,y
}
\]
after the canonical translation identifying their background renewal states. The exact equality defines the minimal ideal free vertex; electron-interference experiments impose a finite lower bound on the product of overlaps over their actual coherence time \citep{Tonomura1989}. A cell-fixed alternative is not another interpretation of the same action. It is an experimentally excluded countervertex.

Sections~3.3 and~22 already require the operational branch to reproduce ordinary quantum interference. The overlap theorem converts that existing requirement into an explicit microscopic test of $W_*$, so the condition adds no ontological premise. Electron interference would falsify an action that violates it. The finite calculation verifies the dual-channel identities, the $Q$--$V_3$ projector orthogonality, and the zero, unit, and intermediate record-overlap cases to machine precision.

\paragraph{Fast/slow architecture.}
The result separates the fast and slow roles within renewal. The fast closure register supplies the objective replacement process and its exported history correlations. The persistent marked fiber carries internal phase and position coherence. The fast output cannot feed old off-diagonal data back into present-cell observables because $\mathcal E_*^\dagger(O)$ is scalar; the slow coherent sector can nevertheless evolve unitarily because it is not the argument of $\mathcal E_*$. Renewal and interference are therefore compatible provided the complete vertex respects the tensor split above.

\paragraph{Retained orbital propagation.}
The overlap theorem fixes the coherence requirement on discarded
records, while the charged spectrum fixes the rest cadence.
H.14 realizes background-wide replacement as the identity on
retained position and internal state. H.15 joins a cochain orbital
operator to an enumerated geometry transfer and preserves its
native rest gap and clock through explicit refinements.
The general gauge-covariant spinor theory and a common invariant
speed across geometries and species remain continuum conditions.
The causal length \(L_*=c\tau_*\) fixes no independent graph
displacement.

\paragraph{Complementary outputs of primitive fusion.}
The two sides of the primitive split have different infrared roles,
\[
V_{3/2}\otimes V_{3/2}
=\underbrace{V_3}_{\text{matched geometric channel}}
\oplus
\underbrace{Q}_{\text{retained mismatch record}}.
\]
Appendix~H.10 proves $(H_{\rm TT})$ on the maximal $V_3$ ray: two
positive-residue, massless helicity-$2$ modes with $E=c|\mathbf k|$, together
with an exact fixed-$j=3$ pullback whose nonmaximal gravitational complement
is gapped.  The reconstruction below concerns the orthogonal $Q$ output.  It
does not replace the graviton theorem; it explains why the retained record
acquires quantum state geometry.  On the factorized vacuum branch the finite
capacity/record decoration adds no graviton pole.

\paragraph{Three nine-state objects and their relation.}
Three objects that share the same local nine coordinates must not be
identified.  First, the H.9 field $d_{ab}$ is a transient response motif: the
incidence vertex creates and removes it, and its occupation number is not a
conserved charge.  Second, the persistent marked defect carries a coherent
internal fiber from one renewal event to the next.  Third, the marked shells
used in the charged-lepton construction are persistent excitations of that
defect.  The incidence theorem below identifies the \emph{response
coordinates} at one event with the fusion complement $Q$; it does not make
the transient occupation number the persistent object.  The Hilbert space
belongs to the coherently transported internal fiber.  At each renewal event
the diagonal persistent label is re-expressed through the response endpoints,
while the free vertex passes their relative phases to the next event without
exporting a cell address.  The lepton shells are structured states carried by
that same persistent defect.  This is the ontology map used below:
\[
\boxed{
\begin{gathered}
\text{event response coordinates}
\longleftrightarrow Q
\longleftrightarrow
\text{persistent coherent fiber},\\
n_d\ \text{is not the persistent charge}
\end{gathered}}
\]

\paragraph{Real carrier space and the incidence theorem.}
The local reconstruction begins without complex scalar multiplication as a
quantum-kinematic premise.  Because $Q=V_0\oplus V_1\oplus V_2$ contains only
integer-spin irreducible representations, it has the invariant real form
\[
\boxed{
Q_{\mathbb R}
=V_{0,\mathbb R}\oplus V_{1,\mathbb R}\oplus V_{2,\mathbb R}
\simeq\mathbb R^9.}
\]
The complex notation used in the parent GFT action is a convenient
complexification; H.11 restricts its quadratic part to $Q_{\mathbb R}$ before
constructing any complex structure.

The sharp reversible split of B.4 has a transmitted map $F_3P_3$ and a
retained map $C$.  Unitarity gives
\[
P_3+C^TC=I_{16},
\qquad
\boxed{C^TC=P_Q.}
\]
In the selected H.9 incidence vertex the retained carrier is the complete
nine-state one-mark space, so $\dim Q_{\mathbb R}=\dim\mathcal H_{d,\mathbb R}
=9$ and $C$ is an isometric identification.  Rotational covariance fixes it
blockwise up to the multiplicity-one freedom.  Denote the resulting real
isometry by $L$:
\[
L^TL=P_Q,
\qquad
LL^T=I_9,
\qquad
\boxed{d=L\chi,\quad \chi=P_Q\Phi_{\mathbb R}=L^Td.}
\]
This proves the carrier-space incidence within the selected sharp reversible
vertex.  It does not derive the hard-core implementation, flags, or detailed
H.9 routing from fusion; those remain selected parts of the marked vertex.

\paragraph{Quadratic inheritance on the real form.}
Restrict Q.9 to real $Q$ fluctuations about the selected $V_3$ background.
At fixed spatial momentum, its positive quadratic operator is
\[
\mathcal H_Q(k)
=Zk^2I_Q+M_c^2I_Q+\kappa(6P_0+5P_1+3P_2).
\]
The condensate may shift the common term $M_c^2$ but cannot remove the
relative splitting.  Since $\delta\chi=L^T\delta d$, the chain rule gives
\[
\delta^2\Gamma_f
=\delta d^T L\mathcal H_QL^T\delta d,
\qquad
\boxed{\mathcal H_D=L\mathcal H_QL^T.}
\]
An extra one-mark quadratic term would be a new defect-only invariant and
would define a different, nonminimal action.  The hard-core constraint does
not alter the fixed-one-mark quadratic sector.

For a local relative frame, $L=L(g(x))$, allow the retained subspace
to vary as $P_Q(x)=L(x)^TL(x)$, with $P_3(x)=I-P_Q(x)$. Define
\[
A_\mu=L\partial_\mu L^T\in\mathfrak{so}(9),
\qquad
D_\mu=\partial_\mu+A_\mu,
\qquad
B_\mu=P_3\partial_\mu L^T .
\]
Since $L^TL=P_Q$ and $P_QP_3=0$,
\[
\partial_\mu(L^Td)=L^TD_\mu d+B_\mu d,
\]
and the two terms are orthogonal.  The inherited kinetic term is therefore
\[
\boxed{
Z\,\partial_\mu\chi^T\partial^\mu\chi
=Z(D_\mu d)^TD^\mu d
+Z\,d^TB_\mu^TB^\mu d .}
\]
The second term is the finite-dimensional Born--Huang correction. If $P_Q$
is fixed, $P_3\partial_\mu L^T=0$ and only the tangent connection remains.
Moving $P_Q$ gives a nonzero term precisely along directions that mix $Q$ with $V_3$.  Its presence is forced by the
primitive-invariance obstruction: the full geometric generators cannot
preserve the nine-dimensional complement.  For spatial gradients it is a
positive semidefinite induced scalar matrix; in the Lorentzian Hamiltonian
the temporal and spatial contributions enter with the corresponding metric
signs.  The complete symmetric Hessian, including this term, is the one used
in the exact transport theorem below.  The mismatch contribution has lower
gap $3\kappa$, so leakage into the matched sector is perturbatively suppressed
when frame gradients are small compared with that gap, but the term is not
discarded in the exact statement.

For a relative rotation of one primitive spin-$3/2$ factor, let $J_R^a=I\otimes J^a$, $\mathcal B_a=P_3J_R^aP_Q$, and $\omega_i^a$ denote its spatial frame gradients; the exact leakage sum gives the following isotropic contraction:
\[
\begin{gathered}
\sum_a\mathcal B_a^\dagger\mathcal B_a=\frac{21}{20}P_2,
\qquad
\sum_i\omega_i^a\omega_i^b=\frac{\omega_{\rm fr}^2}{3}\delta^{ab},
\\[3pt]
\boxed{\sum_{i,a,b}\omega_i^a\omega_i^b
\mathcal B_a^\dagger\mathcal B_b
=\frac7{20}\omega_{\rm fr}^2P_2.}
\end{gathered}
\]
This correction acts on the internal $V_2$ block, whose identification with a galactic response mode remains part of the open source--metric calculation.

\paragraph{Positive Lorentzian dynamics selects the initial polarization.}
At a fixed frame and fixed spatial momentum, define the positive frequency
operator on the real marked carrier,
\[
\boxed{
\boldsymbol\Omega
=\left[
k^2I+\frac{M_c^2}{Z}I
+\frac{\kappa}{Z}\,Q_f\big|_{Q_{\mathbb R}}
\right]^{1/2}.}
\]
Its eigenfrequencies are
\[
\omega_J(k)
=\sqrt{k^2+M_c^2/Z+(\kappa/Z)q_J},
\qquad
(q_0,q_1,q_2)=(6,5,3).
\]
The stability premise is $\boldsymbol\Omega^2>0$: the complement Hessian is
positive in the selected phase.  Its mismatch part is bounded below by
$3\kappa/Z$, so this premise holds on the retained branch used here.

Start only with the real canonical phase space
\[
v=(q,p)\in W=Q_{\mathbb R}\oplus Q_{\mathbb R}^*
\simeq\mathbb R^{18},
\qquad
\omega_{\rm can}=\sum_{a=1}^9dq_a\wedge dp_a.
\]
Hamilton's equations have the real generator
\[
\boxed{
K_0=
\begin{pmatrix}
0&Z^{-1}I\\
-Z\boldsymbol\Omega^2&0
\end{pmatrix}.}
\]
The phase part of its polar decomposition is
\[
\boxed{
J_0=-K_0(-K_0^2)^{-1/2}
=
\begin{pmatrix}
0&-(Z\boldsymbol\Omega)^{-1}\\
Z\boldsymbol\Omega&0
\end{pmatrix}.}
\]
It obeys $J_0^2=-I_{18}$, and
\[
g_0(u,v)=\omega_{\rm can}(u,J_0v)
\]
is positive because
\[
g_0(v,v)
=Zq^T\boldsymbol\Omega q
+\frac1Zp^T\boldsymbol\Omega^{-1}p>0
\qquad(v\ne0).
\]
The corresponding complex amplitude is
\[
\boxed{
z_0=\frac1{\sqrt2}\left[
(Z\boldsymbol\Omega)^{1/2}q
+i(Z\boldsymbol\Omega)^{-1/2}p
\right].}
\]
At a fixed frame the real equations give
\[
\boxed{i\hbar\dot z_0=\hbar\boldsymbol\Omega z_0.}
\]
Thus both $i$ and the fixed-frame Schr\"odinger equation follow from the
positive real Lorentzian dynamics rather than being inserted as kinematic
postulates.

\paragraph{Exact moving-frame transport.}
A changing frame generally mixes the unequal-frequency blocks.  Consequently
the instantaneous polar structure
$J_{\rm inst}(t)=-K_{\rm inst}(t)[-K_{\rm inst}(t)^2]^{-1/2}$ is not, in
general, preserved by the tangent connection alone.  In an instantaneous
particle basis this appears as a Bogoliubov or squeezing term.  Exact Hilbert
transport nevertheless follows from the full real dynamics, without an
adiabatic approximation.

Let $K_{\rm full}(t)\in\mathfrak{sp}(W,\omega_{\rm can})$ be the complete
first-order generator along the frame path, including $A_t$, the mismatch
Hessian, and the Born--Huang contribution above.  Let its fundamental matrix
$S(t,t_0)$ solve
\[
\dot S(t,t_0)=K_{\rm full}(t)S(t,t_0),
\qquad S(t_0,t_0)=I.
\]
Hamiltonian evolution makes $S$ symplectic.  Transport the initial polar
structure by
\[
\boxed{J(t)=S(t,t_0)J_0S(t,t_0)^{-1}.}
\]
Then, exactly,
\[
\boxed{
J(t)^2=-I,
\qquad
\dot J=[K_{\rm full},J],
\qquad
S(t_2,t_0)=S(t_2,t_1)S(t_1,t_0).}
\]
Compatibility and positivity are preserved as well.  Defining
\[
g_t(u,v)=\omega_{\rm can}(u,J(t)v),
\qquad
h_t(u,v)=g_t(u,v)+i\omega_{\rm can}(u,v),
\]
symplecticity gives
\[
\boxed{
h_t(Su,Sv)=h_{t_0}(u,v),
\qquad
g_t(v,v)=g_{t_0}(S^{-1}v,S^{-1}v)>0.}
\]
The propagator is therefore complex linear from $(W,J_0)$ to $(W,J(t))$
and unitary between their Hermitian products.  This is an exact
nine-dimensional Hilbert bundle for arbitrary smooth frame speed.  Squeezing
has not been declared absent: it is the non-complex-linear change of
coordinates between $J(t)$ and $J_{\rm inst}(t)$.  It does not obstruct the
exact transported polarization.

Choose any smooth $h_t$-orthonormal complex frame
$E_t:\mathbb C^9\to(W,J(t))$.  The coordinate propagator
\[
U(t,t_0)=E_t^{-1}S(t,t_0)E_{t_0}
\]
is unitary, so
\[
\boxed{
i\hbar\,\partial_t\psi=H_{\rm eff}(t)\psi,
\qquad
H_{\rm eff}=i\hbar\,\dot U U^\dagger
=H_{\rm eff}^\dagger.}
\]
In the parallel frame $E_t=SE_{t_0}$ this equation is the covariant statement
$D_t\psi=0$; another unitary frame restores the dynamical and geometric
phases.  Different choices are Hilbert-bundle frames. The propagator is a basis-independent map between its endpoint fibers; when the polarization closes, its holonomy changes by conjugation under a change of endpoint frame.  The
instantaneous formula $H=\hbar\boldsymbol\Omega-i\hbar A_t$ is recovered only
when the chosen instantaneous polarization is preserved; otherwise its exact
completion includes the squeezing connection encoded by $J(t)$.

\paragraph{Endpoint polarization and operational controls.}
The transported structure $J(t)=SJ_0S^{-1}$ can depend on the full control path. Two paths ending at the same geometric configuration need not return the same polarization. A transformation represents a unitary control on one fixed polar fiber only when its real symplectic map satisfies $SJ_0S^{-1}=J_0$. General symplectic controls instead connect different polar fibers. The operational realization below assumes access to the stated complex-linear controls on a common retained fiber. The bundle isometry alone does not establish that every physical frame motion supplies such a control.

\paragraph{The one-mark unitary algebra.}
At any fixed event, choose the co-moving polar fiber
\[
\boldsymbol\Omega_{\rm ref}
=\omega_0P_0+\omega_1P_1+\omega_2P_2
\]
and remove its common phase,
\[
\boxed{
S_\omega
=\boldsymbol\Omega_{\rm ref}
-\frac{\operatorname{Tr}\boldsymbol\Omega_{\rm ref}}9I.}
\]
Because the three frequencies are unequal, $S_\omega$ is a nonzero real
symmetric traceless operator.  In the selected complex-linear control realization, the frame generators act as
$\mathfrak{so}(9)$, and the orbit of $S_\omega$ spans
$\operatorname{Sym}_0(9)$.  Hence
\[
\boxed{
\operatorname{Lie}\!\left(\mathfrak{so}(9),iS_\omega\right)
=\mathfrak{su}(9),}
\]
with $36+44=80$ generators.  Restoring the common frequency gives
$\mathfrak u(9)$.  Exact symplectic transport carries this complete algebra
unitarily from fiber to fiber.  Thus
\[
\boxed{\mathcal H_{\rm mark}(t)\simeq\mathbb C^9}
\]
throughout the stable fixed-one-mark quadratic trajectory.  Algebraic
controllability states that these generators span the full unitary algebra;
it is distinct from the later dynamical question of which gate measure the
completed GFT actually samples.

\paragraph{Why the physical states are rays.}
Write $z$ in one fixed polar fiber and normalize $z^\dagger z=1$.
The hard-core marked algebra contains
the matrix units
\[
d_a^\dagger d_b=|a\rangle\langle b|,
\]
and $\mathfrak{su}(9)\oplus\mathbb RI$ spans every Hermitian $9\times9$
fixed-one-mark bilinear response operator.  For such an operator $T$, the response
$f_T(z)=z^\dagger Tz$ is unchanged by $z\mapsto e^{i\theta}z$.  Conversely,
if normalized $z$ and $w$ give the same value of $f_T$ for every traceless
Hermitian $T$, then
\[
\operatorname{Tr}\!\left[T(zz^\dagger-ww^\dagger)\right]=0
\quad\text{for all }T.
\]
Thus $zz^\dagger-ww^\dagger=cI$.  Equal unit traces give $c=0$, and equality
of the rank-one projectors gives $w=e^{i\theta}z$.  The complete
fixed-one-mark bilinear response algebra therefore separates normalized
states exactly modulo common phase:
\[
\boxed{\mathcal S_{\rm mark}=S^{17}/U(1)=\mathbb{CP}^{8}.}
\]
Projectivization is an operational consequence of the available responses,
not an asserted ray rule.  It concerns the normalized fixed-one-mark sector
and does not assert a superselection rule between the vacuum and one-mark
sectors.

\paragraph{Fubini--Study geometry and the factor of two.}
The quotient of the Hermitian metric gives
\[
ds_{\rm FS}^2=dz^\dagger dz-|z^\dagger dz|^2.
\]
With $z_a=\sqrt{P_a}\,e^{i\phi_a}$ and $\sum_aP_a=1$,
\[
\boxed{
ds_{\rm FS}^2
=\frac14\sum_a\frac{dP_a^2}{P_a}
+\sum_aP_a\left(d\phi_a-\sum_bP_bd\phi_b\right)^2.}
\]
The probability block is one quarter of the Fisher metric. In the canonical normalization used to define $z=(Q+iP^{\rm can})/\sqrt2$, the real variables obey $Q_a=\sqrt{2P_a}\cos\phi_a$ and $P_a^{\rm can}=\sqrt{2P_a}\sin\phi_a$. Therefore
\[
dQ_a\wedge dP_a^{\rm can}=dP_a\wedge d\phi_a.
\]
The K\"ahler form associated with $ds_{\rm FS}^2$ is half this canonical form after projective reduction. With $\Phi_a=\hbar\phi_a$ and the metric scaled by $\hbar$,
\[
\boxed{
\Omega_{\rm K}=\frac12\sum_a dP_a\wedge d\Phi_a,
\qquad
ds^2=\sum_a\left[
\frac{\hbar}{4P_a}dP_a^2+\frac{P_a}{\hbar}d\Phi_a^2
\right]}
\]
on horizontal tangent vectors satisfying $\sum_aP_a d\phi_a=0$. The physical canonical two-form in action units is $2\Omega_{\rm K}$. Keeping these normalizations separate preserves the Fisher factor and the canonical oscillator convention.

\paragraph{Conditional Gaussian-strand Fisher check.}
This check belongs to the H.9 Gaussian-strand response realization; it is
not needed for the Hilbert or projective-space theorem above. Its additional
identification is that the three response coordinates carried by the first
Hermite modes are the closure-score coordinates. The one-strand bridge is
explicit:
\[
\int d^3\xi\,\xi_a
e^{-\xi^2/2+i\sqrt{\eta_*}\,\xi\cdot C}
=i\sqrt{\eta_*}\,C_a
\int d^3\xi\,
e^{-\xi^2/2+i\sqrt{\eta_*}\,\xi\cdot C}.
\]
Thus two first-Hermite strands produce the amplitude direction
$C_a^PC_b^H$. Within this response realization, tetrahedral symmetry makes the second moment isotropic. The identity $C^2=K^2$ and the closure stationarity equation then give
\[
\langle C_aC_c\rangle=\frac{1}{2\eta_*}\delta_{ac},
\]
and the normalized marked response is
\[
r_{ab}=2\eta_*C_a^PC_b^H,
\qquad
\langle r_{ab}r_{cd}\rangle=\delta_{ac}\delta_{bd}.
\]
Perturbing the square-root probability amplitude by
$A_\epsilon=A_0(1+\epsilon^{ab}r_{ab})+O(\epsilon^2)$ gives statistical score
$s_{ab}=2r_{ab}$.  Hence
\[
I^{\rm Fisher}_{ab,cd}=4\delta_{ac}\delta_{bd},
\qquad
g^{\rm mark}_{ab,cd}=\delta_{ac}\delta_{bd},
\]
and therefore
\[
\boxed{g_{\rm mark}=\frac14I_{\rm Fisher}=g_{\rm FS}}
\]
at the fresh marked tangent point, conditional on that coordinate
identification. The definition in H.9 and the variance decomposition in Appendix~B.2 distinguish these response coordinates from
the ordinary mean-closure vector $\mathbf c(b)$, whose squared norm omits the
irreducible quantum-variance term in $K^2$. Extending the equality over all
of $\mathbb{CP}^8$ would additionally require the corresponding finite
response family away from the fresh state.

\paragraph{Tensor composition from independent record factors.}
The microscopic multiplicative-label premise says that two simultaneously
retained marks have all $9\times9=81$ ordered joint labels before the final
host projection.  Applying the positive-frequency construction to their free
joint quadratic dynamics gives a complex carrier
\[
\mathcal H_{12}\simeq\mathbb C^{81}.
\]
Let $\mathcal A_1$ and $\mathcal A_2$ be the complete response algebras of the
two marks.  Independence before the bridge means
\[
\mathcal A_1\simeq M_9(\mathbb C),
\qquad
\mathcal A_2\simeq M_9(\mathbb C),
\qquad
[\mathcal A_1,\mathcal A_2]=0,
\]
with faithful actions on $\mathcal H_{12}$.  A faithful representation of
$M_9(\mathbb C)$ on $\mathbb C^{81}$ is unitarily equivalent to
$M_9(\mathbb C)\otimes I_m$, and $9m=81$ fixes $m=9$.  Hence
\[
\mathcal A_1=M_9(\mathbb C)\otimes I_9,
\qquad
\mathcal A_1'=I_9\otimes M_9(\mathbb C).
\]
Commutativity gives $\mathcal A_2\subseteq\mathcal A_1'$.  Both algebras have
complex dimension $81$, so the inclusion is equality.  Therefore
\[
\boxed{
\mathcal H_{12}\cong\mathbb C^9\otimes\mathbb C^9,
\qquad
\mathcal A_2=\mathcal A_1'.}
\]
The tensor product is thus fixed by the multiplicative joint labels and their
two complete commuting local algebras; it is not an additional quantum
postulate.  The free joint dynamics fixes $J_{12}$ first.  The projected
bridge is then an interaction on that fixed complex structure rather than an
ingredient used to manufacture the factorization.

For the primitive pairs $(1,2)$ and $(3,4)$ this gives
\[
\boxed{\mathcal H_{2,{\rm retained}}=Q_{12}\otimes Q_{34}.}
\]
Conditional on the two retained factors and their bridge surviving the final
host projection, the same physical bridge used in P.6 acts on the neighboring
complements.
For primitive pairs $(1,2)$ and $(3,4)$, take
\[
H=2J_z=\operatorname{diag}(3,1,-1,-3).
\]
\[
\begin{aligned}
H_{QQ}
&=(P_Q^{12}\otimes P_Q^{34})(H_2\otimes H_3)
(P_Q^{12}\otimes P_Q^{34})\\
&=A_R\otimes A_L,
\end{aligned}
\]
where $A_R=P_Q(I\otimes H)P_Q|_Q$ and
$A_L=P_Q(H\otimes I)P_Q|_Q$.  Exact rational spin-$3/2$ evaluation gives
\[
\operatorname{rank}A_R=\operatorname{rank}A_L=8,
\quad
\operatorname{Tr}A_R=\operatorname{Tr}A_L=0,
\quad
\operatorname{Tr}A_R^2=\operatorname{Tr}A_L^2=38.
\]
Thus the bridge is nonzero, nonlocal, and cannot be written as a sum of
one-mark Hamiltonians.  With the two local algebras, the tensor-product lemma
of P.6 yields
\[
\boxed{
\operatorname{Lie}\!\left(
\mathfrak{su}(9)_1,\mathfrak{su}(9)_2,
iA_R\otimes A_L\right)=\mathfrak{su}(81).}
\]
Induction along a connected marked network gives
\[
\boxed{
\mathcal H_N=(\mathbb C^9)^{\otimes N},
\qquad
\mathfrak g_N=\mathfrak{su}(9^N),}
\]
provided simultaneous retained marks exist and every spanning-tree bridge
survives the complete gauge/intertwiner projection of the geometric host.

\paragraph{Infinite-volume network algebra.}
Exact evaluation of the projected bridge gives
\[
\operatorname{spec}A_R=\left\{0,\ \pm2,\ \pm\frac{\sqrt{205}}{5},\ \pm\left(1+\frac{2\sqrt{15}}{5}\right),\ \pm\left(\frac{2\sqrt{15}}{5}-1\right)\right\},
\]
whose squares sum to $38$, so
\[
\|A_R\|=\frac{\sqrt{205}}{5}=2.863564,
\qquad
\boxed{\|A_R\otimes A_L\|=\frac{41}{5}.}
\]
The bridge is a uniformly bounded finite-range interaction. On any locally finite host $\Gamma$ of bounded degree, set $\mathcal A_F=\bigotimes_{x\in F}M_9(\mathbb C)$ for finite $F$ and form the quasi-local algebra
\[
\mathcal A=\overline{\bigcup_{F\Subset\Gamma}\mathcal A_F}.
\]
Lieb--Robinson bounds for bounded finite-range interactions give the infinite-volume dynamics
\[
\tau_t(A)=\lim_{\Lambda\uparrow\Gamma}e^{itH_\Lambda}Ae^{-itH_\Lambda}
\]
in norm for every local observable, and every state $\omega$ on $\mathcal A$ has a GNS representation $(\pi_\omega,\mathcal H_\omega,\Omega_\omega)$. The same $C^*$-dynamical framework carries the spontaneous-magnetization theorem for the quantum Heisenberg ferromagnet \citep{OpenAIHeisenberg2026}. Under the bridge-survival condition above, the retained network therefore has a rigorous infinite-volume quantum algebra and Hilbert-space representations. The open questions are the physical state, its continuum limit, whether its excitations organize into Fock particles, and the physical supply of durable history capacity.

\paragraph{Born weights from reversible equivalence and physical records.}
The Hilbert norm fixes the probability rule without using Gleason's
theorem as the derivation.  Let $v\in\mathcal H\simeq\mathbb C^d$ be an
unnormalized physical branch vector and let $W(v)\ge0$ be its weight.  Three
probability requirements are used explicitly:
\begin{enumerate}
\item an exhaustive normalized record has total weight one;
\item mutually exclusive durable records have finitely additive weights;
\item a reversible internal relabeling cannot change a branch weight merely
because it rotates the branch within the complete response space.
\end{enumerate}
The third statement is $W(Uv)=W(v)$ for the reversible $SU(d)$ equivalence
group.  It is a consistency requirement on branch weights, not a claim that
the uncontrolled GFT dynamics samples every unitary with equal frequency.
Because $SU(d)$ is transitive on equal-norm vectors,
\[
W(v)=F(\|v\|^2)
\]
for some nonnegative function $F$.  Now take two orthogonal alternatives
$v\perp w$ whose distinguishing records are physically available.  Forgetting
that distinction replaces them by the coarse alternative $v+w$, while finite
record additivity gives
\[
W(v+w)=W(v)+W(w).
\]
Since $\|v+w\|^2=\|v\|^2+\|w\|^2$,
\[
F(x+y)=F(x)+F(y).
\]
Nonnegativity makes $F$ monotone on $[0,1]$; a monotone additive function on
that interval is linear.  Normalization $F(1)=1$ fixes its coefficient:
\[
\boxed{W(v)=\|v\|^2.}
\]
For a complete record projector $\Pi_a$ and a normalized state $|\psi\rangle$,
\[
\boxed{
p(a\mid\psi)=\|\Pi_a\psi\|^2
=\langle\psi|\Pi_a|\psi\rangle,}
\]
and a rank-one record gives $p(a\mid\psi)=|\langle a|\psi\rangle|^2$.
The named scope condition is \emph{record completeness}: every orthogonal
decomposition used in this argument must correspond to an exhaustive set of
physically recordable alternatives.  Ordinary nonnegativity, normalization,
and finite additivity remain probability premises.  Gleason's theorem is an
independent consistency check on the resulting projector measure, and
its $d>2$ restriction is not inherited by
this argument.

If a preparation record $r$ is forgotten, the same coarse-graining gives
\[
\rho=\sum_r p_r|\psi_r\rangle\langle\psi_r|,
\qquad
\boxed{p(a)=\operatorname{Tr}(\rho\Pi_a).}
\]
For a retained composite, applying the pure-state rule to
$\Pi_a\otimes I$ gives the reduced state by the partial trace.  Mixed states
and reduced states are therefore bookkeeping for unresolved physical
preparation or subsystem records, not additional kinematic postulates.

\paragraph{Finite measurements and channels from reversible record coupling.}
Let $S$ be a retained system and $R$ a retained record factor, initially in
$|0\rangle_R$.  Tensor composition has already been derived.  For a
reversible interaction $U$ and orthogonal final records $|a\rangle_R$, define
\[
\boxed{M_a={}_R\!\langle a|U|0\rangle_R.}
\]
Then
\[
U(|\psi\rangle|0\rangle)
=\sum_aM_a|\psi\rangle|a\rangle,
\qquad
\boxed{\sum_aM_a^\dagger M_a=I.}
\]
The derived norm rule gives
\[
\boxed{
E_a=M_a^\dagger M_a,
\qquad
p(a)=\operatorname{Tr}(\rho E_a),
\qquad
\sum_aE_a=I.}
\]
Thus the finite POVM is the system-side image of a reversible
system--record interaction.  Conditioning on the retained record gives
\[
\rho\longmapsto
\frac{M_a\rho M_a^\dagger}
{\operatorname{Tr}(M_a\rho M_a^\dagger)},
\]
whereas forgetting it gives the trace-preserving channel
\[
\boxed{\rho\longmapsto\sum_aM_a\rho M_a^\dagger.}
\]
Extra unobserved labels $\mu$ give a general finite instrument
$\mathcal M_a(\rho)=\sum_\mu M_{a\mu}\rho M_{a\mu}^\dagger$.  Conversely,
any finite Kraus family satisfying
$\sum_{a,\mu}M_{a\mu}^\dagger M_{a\mu}=I$ defines the isometry
\[
V|\psi\rangle
=\sum_{a,\mu}M_{a\mu}|\psi\rangle|a,\mu\rangle,
\]
which extends to a unitary on a sufficiently large retained space.  Finite
quantum instruments are therefore exactly reversible record dynamics viewed
after conditioning or forgetting.

No-signaling follows from the same construction.  For a local
trace-preserving channel on $B$,
\[
\mathcal E_B(\rho)=\sum_\mu
(I\otimes M_\mu)\rho(I\otimes M_\mu^\dagger),
\qquad
\sum_\mu M_\mu^\dagger M_\mu=I,
\]
every $A$-local observable obeys
\[
\operatorname{Tr}[(O_A\otimes I)\mathcal E_B(\rho)]
=\operatorname{Tr}[(O_A\otimes I)\rho].
\]
Locality plus trace preservation therefore leaves the remote reduced state
unchanged.

\paragraph{The finite decoherence functional as a Gram kernel.}
For an initial state $\rho_0$ and a history class operator $C_h$, define the
branch vector $|\Psi_h\rangle=C_h|\Psi_0\rangle$ for a pure preparation.  The
canonical Hermitian overlap derived above gives
\[
\boxed{
\mathcal D(h,h')
=\langle\Psi_{h'}|\Psi_h\rangle
=\operatorname{Tr}(C_h\rho_0C_{h'}^\dagger).}
\]
Its diagonal is the Born weight
$p(h)=\mathcal D(h,h)=\|C_h\Psi_0\|^2$.  For a coarse event
$A$, amplitudes combine first,
\[
C_A=\sum_{h\in A}C_h,
\qquad
p(A)=\left\|\sum_{h\in A}|\Psi_h\rangle\right\|^2.
\]
The off-diagonal entries of $\mathcal D$ are precisely the interference cross
terms.  Durable records make distinct branch vectors orthogonal; only then
does the expression reduce to $p(A)=\sum_{h\in A}p(h)$.  The mathematical
form of the decoherence functional is therefore derived for every finite
retained system.  What remains open is its global physical realization: the
continuum/Fock limit, the cosmological initial state $\rho_0$, which full
histories actually decohere, the physical source of durable history capacity (P.9), and the gravitational entropy of the minimal-record boundary.

\paragraph{Projective phase is not electric charge.}
The ray construction supplies the kinematic principal bundle
$S^{17}\to\mathbb{CP}^8$ and its local representative connection
\[
\mathcal A_\mu=-iz^\dagger\partial_\mu z,
\qquad
\mathcal A_\mu\mapsto\mathcal A_\mu+\partial_\mu\theta.
\]
This does not identify the bundle with electromagnetism.  In fact, the free
decorated block preserves the response-mark number $n_d$, but the incidence
vertex does not:
\[
[n_d,V_Q]
=\sum_{ab}\left(d_{ab}^\dagger\xi_a^P\xi_b^H-\mathrm{h.c.}\right)\ne0.
\]
The temporary H.9 response mark is therefore not the conserved electric
charge.  Any electromagnetic generator must belong to the persistent charged
defect and must still produce Gauss law, a gapless Maxwell field, and the
observed coupling.

\paragraph{Relation to Hamilton--Killing reconstruction and remaining scope.}
Caticha's reconstruction supplies an independent general theorem connecting
Fisher geometry, a compatible K\"ahler structure, complex amplitudes, and
linear unitary evolution \citep{CatichaHK2021,CatichaQED2025}.  Here the
finite one-mark complex structure, ray space, K\"ahler geometry, and full
unitary algebra have instead been obtained from the displayed fusion and
quadratic action.  Reversible equivalence and record additivity then fix the
Born norm; tensor composition, instruments, no-signaling, and the finite
decoherence functional follow in sequence.  Caticha's result is a consistency
check and a guide for the continuum extension.  What remains is to show that
the final constrained GFT host preserves the marked factors and bridges,
construct indefinite record capacity or a controlled continuum/Fock limit,
derive the global cosmological state and its decohering histories, and derive
electromagnetic charge and Maxwell dynamics.  None of those requirements
weakens the finite theorem or the exact renewal and overlap results.

\paragraph{Closure statement.}
\[
\boxed{
\begin{gathered}
\text{renewed-output coherence blindness: exact},\\
\text{record-overlap multiplier: exact},\\
Q\perp V_3:\ \text{exact},\\
\text{cell-addressed marked export: excluded by interference},\\
\text{branch-independent retained marked fiber: required vertex condition},\\
\mathbb C^9/\mathbb{CP}^8\text{ one-mark reconstruction: exact in the selected quadratic completion},\\
\text{arbitrary-speed Hilbert transport: exact under the full symplectic flow},\\
\text{instantaneous-polarization squeezing and Born--Huang term: retained exactly},\\
g_{\rm mark}=I_{\rm Fisher}/4\text{: conditional Gaussian-strand check},\\
\text{two-factor tensor product: exact from multiplicative labels and commuting full algebras},\\
\mathfrak{su}(81)\text{ two-mark closure: exact for the unreduced product before final host projection},\\
\mathfrak{su}(9^N)\text{ network extension: conditional on surviving projected bridges},\\
\text{quasi-local algebra, }\tau_t\text{, and GNS representations: exact given surviving bridges},\\
\text{Born weights, instruments, no-signaling, and finite history kernel: derived},\\
\text{continuum/Fock realization, cosmological state, history capacity, and QED: open}.
\end{gathered}
}
\]

\subsection*{H.12 Primitive pair interaction and collective shape spectrum}
\label{aqg:parent}

An explicit gauge-invariant pair interaction realizes the selected
four-bundle fusion sector. Hamiltonian formulas use \(\hbar=1\);
the numerical geometric tests also set \(c=L_*=1\). On
\(\mathcal I_3=\operatorname{Inv}(V_3^{\otimes4})\), of dimension seven,
define
\[
X=\mathbf J_1\cdot\mathbf J_2,\quad
Y=\mathbf J_1\cdot\mathbf J_3,\quad Z=-12I-X-Y,
\qquad (x,y,z)=(X+4I,Y+4I,Z+4I).
\]
Set \(u=(x,y,z)\), \(B=\sum_ru_r^2\), and \(Q_V=i[X,Y]\).
For two cells write
\[
\mathcal D_s=\sum_ru_r\otimes u_r,\quad
\mathcal E_s=\sum_r(u_r^2\otimes u_r+u_r\otimes u_r^2),\quad
\mathcal T_s=\sum_ru_r^2\otimes u_r^2.
\]
The dimensionless parent is
\[
\begin{split}
\widehat H_{\rm pair}=\frac1{13^4}\big[&
3\,13^4I+13^2(B\otimes I+I\otimes B-8\mathcal D_s)\\
&+52\mathcal E_s+B\otimes B-2\mathcal T_s
-144Q_V\otimes Q_V\big].
\end{split}
\]
For the raw fusion map \(T:\mathcal I_3\otimes\mathcal I_3\to\mathcal I_6\),
\[
T_{K,kl}=169\sqrt{(2k+1)(2l+1)(2K+1)}
\left\{\begin{matrix}3&3&k\\3&3&l\\6&6&K\end{matrix}\right\}^{\!2},
\qquad F=(TT^\dagger)^{-1/2}T.
\]
The raw map has nontrivial Gram eigenvalues; the polar normalization gives
\(FF^\dagger=I_{13}\). Exact rational nullspaces and real-root counting give
\[
\widehat H_{\rm pair}\ge0,\qquad
\ker\widehat H_{\rm pair}=\operatorname{im}F^\dagger,\qquad
\dim\ker\widehat H_{\rm pair}=13.
\]
There are 36 positive eigenvalues and a gap \(1.93012418\) in parent
units. The common annihilator family first appears at total shape degree
four: its dimension is four, and the stacked rank is 36.
The swap-odd annihilators contain two-cell correlations. The physical
scale multiplying \(\widehat H_{\rm pair}\) is independent of this
dimensionless kernel calculation.

Extend the interaction by the identity on capacity and marks. It then
preserves the original capacity spectrum and weights. For a supplied
coarse coupled Hamiltonian \(H_c\),
\[
H_{\rm fine}=\widetilde F^\dagger H_c\widetilde F
+J_p\widehat H_{\rm pair}\otimes I_{\rm matter},
\qquad \widetilde F=F\otimes I_{\rm matter},
\]
has the exact retained evolution
\(\widetilde F e^{-itH_{\rm fine}}\widetilde F^\dagger=e^{-itH_c}\).
The polynomial specifies the discarded-sector interaction; dynamics
inside the retained sector are supplied by \(H_c\).
Four-bundle fusion doubles aligned face flux. Its regular volume scales
as \(2^{3/2}\), whereas two spatially additive cells have twice the volume.
The additive primitive volume in C.8 keeps these observables distinct.

\paragraph{Connected pairs.}
For one common port identification there is an invertible matrix
\[
C_{kr}=\frac{(-1)^k}{\sqrt{2k+1}}\binom6r
\langle3,3;3,-3|k,0\rangle
\langle3,3-r;3,r-3|k,0\rangle
\]
such that \(v(z)=C(1,z,\ldots,z^6)^T\) and
\(\ker H_{12}=\operatorname{span}_z v(z)\otimes v(z)\).
In the rational similarity basis,
\(\det[\operatorname{diag}(2k+1)^{-1/2}C]=-375/7553546\).
Transforming each factor by \(C^{-1}\), an edge constraint makes a
coefficient depend only on the sum of its two indices. On a connected
graph, moving index units along edges relates every configuration with the
same total. Hence
\[
\dim\ker\sum_{\langle ij\rangle}H_{ij}=6n+1.
\]
The two-, three-, and four-site gaps are \(1.930124\), \(0.965202\),
and \(0.565412\). Twisted port identifications define a different problem.

The normalized tangent to \(v(z)\) produces a variational matrix
\(a_sL_\Gamma\), with \(a_s=0.96528324\). On an \(n\)-site path,
\[
\Delta_n\le2a_s[1-\cos(\pi/n)]\sim a_s\pi^2/n^2.
\]
The tangent subspace has nonzero leakage, so this is an upper bound on
the gap. It excludes using this all-edge parent alone as the desired
relativistic propagation Hamiltonian. The coherent completion keeps the
parent inside disjoint primitive pairs and supplies inter-block metric
dynamics through the angular contraction and its ADM continuation.

\subsection*{H.13 Metric response and local linear constraints}
\label{aqg:linear}

The angular normalization can be tested on an infinite four-dimensional
Kuhn triangulation, with 24 simplices and 15 edge classes per translation
cell. This is a metric regulator. Each physical capacity event can span
several regulator cells.

On a closed six-simplex complex, the angular contraction equals the
Einstein--Regge weight at the same \(G_0\) for both equilateral and
asymmetric lengths. Independent finite differences of the action and
the Schl\"afli-reduced variation agree to relative errors
\(5.0\times10^{-11}\) and \(3.9\times10^{-11}\).
This checks the nonlinear metric variation on a finite complex.

At a flat simplex, area and dihedral-angle derivatives obey the
Schl\"afli identity, giving the Hessian of \(-\sum_hA_h\delta_h\) as
\(H_\sigma=J_A^TJ_\theta\). Bloch assembly gives
\[
H_\ell(p)=\sum_\sigma P_\sigma^\dagger H_\sigma P_\sigma,\qquad
G_{e\mu}(p)=\frac{d_{e\mu}}{\ell_e}(e^{ip\cdot d_e}-1),\qquad
H_\ell G=0.
\]
The response of an edge to a plane-wave metric is
\[
\delta\ell_e=\frac{d_e^\mu d_e^\nu h_{\mu\nu}}{2\ell_e}
e^{ip\cdot d_e/2}\operatorname{sinc}(p\cdot d_e/2).
\]
Eliminating the five nonmetric length directions by their saddle
equations leaves the ten-component metric form
\[
\begin{split}
h^\dagger H_{\rm eff}h=\frac{V_4}{4}\big[
&p^2h_{\mu\nu}^*h_{\mu\nu}-p^2|\operatorname{tr}h|^2
-2|(ph)_\nu|^2\\
&+2\operatorname{Re}\{(\operatorname{tr}h)^*
p_\mu h_{\mu\nu}p_\nu\}\big]+o(p^2).
\end{split}
\]
The eliminated Euclidean block has one negative direction. The Schur
complement is therefore a saddle reduction. Multiplying by
\(b/A_{\rm ch}=1/(8\pi G_0)\) gives the Einstein quadratic action.
The Lorentzian equation kernel has four gauge directions off shell and
six directions on the null cone, leaving two on-shell polarizations.

Two modestly sheared backgrounds give relative tensor errors
\(6.5\times10^{-5}\) and \(1.1\times10^{-4}\) at dimensionless
momentum \(0.025\). With source coupling \(V_4T^{\mu\nu}h_{\mu\nu}/2\),
the static response approaches \(h_{44}=-8\pi G_0\rho/k^2\).
The two extrapolated Newton coefficients are \(1.030020976\) and
\(1.030020970\), against \(G_0=1.030020983\). A scalar finite-element
probe has the same leading cone. H.15 supplies a finite retained
orbital model and its clock; a common covariant spinor limit across
the geometric ensemble remains to be derived.

\paragraph{Staggered constraint complex.}
On a periodic cubic regulator let \(D_i^+=T_i-I\),
\(D_i^-=I-T_i^T\), and \(\Delta=\sum_iD_i^-D_i^+\).
Order the metric components as \(xx,yy,zz,xy,xz,yz\), with diagonals at
vertices, off-diagonals at plaquette centers, and vector parameters at
edge centers. Define
\[
(G\xi)_{ii}=2D_i^-\xi_i,\qquad
(G\xi)_{ij}=D_i^+\xi_j+D_j^+\xi_i,
\]
\[
\begin{aligned}
W&=\operatorname{diag}(1,1,1,2,2,2),&
\mathsf T&=(I,I,I,0,0,0),\\
\mathsf R&=-G^TW/2,&
\mathsf d&=(D_x^-,D_y^-,D_z^-),\qquad
\mathsf B=\mathsf d\mathsf R .
\end{aligned}
\]
The scalar constraint and kinetic matrix are
\[
\mathsf C=\mathsf B-\Delta\mathsf T,\qquad
\mathsf K=W^{-1}-\tfrac12\mathsf T^T\mathsf T.
\]
The spatial quadratic form is
\[
\mathsf L=-W\Delta+\mathsf T^T\Delta\mathsf T
-2\mathsf R^T\mathsf R-\mathsf T^T\mathsf B-\mathsf B^T\mathsf T .
\]
Block identities are implicit in this notation. Exact integer identities
\[
\mathsf CG=0,\qquad \mathsf LG=0,\qquad
\mathsf C(2\mathsf K)=-\mathsf dG^T
\]
preserve \(\mathsf Cq=0\) and \(G^Tp=0\) under canonical quadratic
shears. Replacing the trace coefficient \(1/2\) by \(\lambda\) leaves
the exact residual
\[
\mathsf C(2\mathsf K_\lambda)+\mathsf dG^T
=2(2\lambda-1)\Delta\mathsf T,
\]
fixing \(1/2\) within this ansatz. The staggered-complex identities and this coefficient are kernel checked.

The TT reduction has two equal frequencies with
\(\widehat k^2=4\sum_i\sin^2(k_i/2)\). There are no additional zeros
at Brillouin-zone corners. The physical Hamiltonian is
\[
H_{\rm lin}=\frac{16\pi G_0}{V_B}p^T\mathsf Kp+
\frac{V_B}{64\pi G_0a_B^2}q^T\mathsf Lq,
\]
so the continuum TT speed is one in tick units.

Over the prime field of order 13, the same integer complex gives commuting
Weyl constraints with ranks \(3(V-1)\) and \(V-1\).
There are \(2(V-1)+6\) logical coordinates; the six homogeneous metric
moduli are separated from propagating polarizations. Finite Clifford
shears preserve the code exactly. This algebraic example does not fix
the real-time normalization or establish a gravitational phase; its
prime dimension is independent of the 13-dimensional pair intertwiner.

\paragraph{Finite-count mobility and the local canonical map.}
A direct replacement \(\mathsf K\mapsto J(q)\mathsf KJ(q)\), with
\(J_{AA}=(1-q_A^2)^{1/4}\), violates the undeformed linear constraints
at cubic order. At amplitude \(0.16\), the tested scalar and vector
constraint velocities are \(0.00297\) and \(0.00245\).
A TT projection removes them but couples distant sites.
The local canonical map
\[
Q_A=F(q_A),\quad F'(q)=(1-q^2)^{-1/4},\qquad
P_A=(1-q_A^2)^{1/4}p_A
\]
instead gives exact interior constraints
\(\mathsf CF(q)=0\) and \(G^T(Jp)=0\).
The transformed Hamiltonian is still the free quadratic theory.
Here \(F(1)=\tfrac12B(1/2,3/4)=1.19814023\), the scalar measure is
\(dq/J(q)\), and the first finite-count interval gap converges to
\(\pi^2/[8F(1)^2]=0.85939823\).
Nonlinear gravitational interactions use the metric reconstruction and
curvature potential of N.10--N.12.

\subsection*{H.14 Native proper-time sources and coherent renewal}
\label{aqg:nativeclock}

The retained rest generator of H.9 supplies a prescribed proper-time
source. Let \(P_e\) project onto the retained electron occupation and
\[
\epsilon_e=-\frac32Z_e\log(1-e^{-7g_{\rm share,eff}})
=4.18545879558\times10^{-23}.
\]
The proper-time attachment \(s=\tau_\gamma/\tau_*\) gives
\[
T_e(s)=e^{-\epsilon_e sP_e},\qquad
U_e[\gamma]=e^{-i\epsilon_e P_e\tau_\gamma/\tau_*}.
\]
Durations add along joined curves, so the same rest energy fixes their
composed phase and stationary lapse coupling. This attachment is a
selected interaction rule. An energy eigenstate carries a phase;
readable timing requires relative phases or a further clock observable.

Two intrinsically transported normal geodesics have been constructed
on every member of the 1,920-slab \(\alpha=7\) family of Q.11b.
Their durations range approximately from \(6.54\) to \(6.99\) ticks.
The native logarithmic transfer-sewing error is
\(1.72\times10^{-16}\), and joined-slab durations add to rounding
accuracy. A local Wick endpoint check for fixed curve segments does
not establish reflection positivity of the full marked path rule.

Geometry-dependent duration is necessary for a connected source
response. One fixed timelike edge per slice gives length
\(Ka\sqrt\alpha\), independent of the geometry, so its rest-mass factor
cancels from the normalized two-source ratio.
For the varying prescribed lines,
\[
Z(m_1,m_2)=\sum_{\mathcal T}p_0(\mathcal T)
e^{-m_1\ell_1(\mathcal T)-m_2\ell_2(\mathcal T)},\qquad
\partial_{m_1}\partial_{m_2}\log Z
=\operatorname{Cov}_{m_1,m_2}(\ell_1,\ell_2).
\]
The finite free-line covariance is \(-0.00171521593106\).
At the native electron gap, the leading connected logarithm is
\(-3.00472647\times10^{-48}\), evaluated through derivatives and
logarithms. These nonstatic trajectories do not measure a held-source
force. J.13 states the required source preparation and limits.

\paragraph{Background renewal and the retained path.}
A correct fresh capacity marginal is insufficient to preserve
interference. For a specified SWAP-reset dilation with bath slots in
\(\rho_*\), branches that trigger \(n\) and \(m\) resets acquire the
off-diagonal multiplier
\[
\gamma_{nm}=\operatorname{Tr}\rho_*^{|n-m|+1}.
\]
One extra reset gives \(0.0006036303\) for the native capacity state.
Exported bath records distinguish the reset schedules.
This is a counterexample for that dilation; it does not apply to
every renewal implementation.

The compatible free realization renews background cells independently
of the retained particle's unmeasured position:
\[
\mathcal E=\operatorname{id}_{X,Q}\otimes\mathcal E_{\rm background},
\qquad
\operatorname{Tr}_{\rm background}\mathcal E(\rho)
=\operatorname{Tr}_{\rm background}\rho.
\]
Proper-time phase accumulation triggers no extra erasure.
The retained-density error is \(5.73\times10^{-17}\);
spacelike-disjoint reset orderings agree to \(3.63\times10^{-17}\).
Local event schedules suffice on a fixed geometry.
The moving graph frame of Q.11a handles the encoded bond basis.
Coherence between geometries with different exported event records
and full gravitational refoliation remain additional conditions.

\paragraph{A positive mass-conjugate clock.}
For a supplied half-transfer \(A\) and positive middle-boundary duration
operator \(N\), choose
\[
K(m)=A^\dagger e^{-mN}A,\qquad
H(m)=-\frac1{\Delta t}\log K(m).
\]
On fixed support, \(C=-\partial_mK\succeq0\), and
\[
\partial_mH=\frac1{\Delta t}\int_0^\infty
(K+sI)^{-1}C(K+sI)^{-1}\,ds\succeq0.
\]
If \(N_{\min}I\preceq N\preceq N_{\max}I\), then
\[
\frac{N_{\min}}{\Delta t}I\preceq\partial_mH
\preceq\frac{N_{\max}}{\Delta t}I.
\]
The theorem follows from the Fr\'echet derivative of the logarithm
and \(\int_0^\infty(K+sI)^{-1}K(K+sI)^{-1}ds=I\).
It gives a positive integrated response
\(iU_t^\dagger\partial_mU_t\).
A seven-state matrix test and its 169-dimensional connected-code lift
check the identity. Entrywise positive source weights alone do not
guarantee this operator inequality; a finite counterexample has a
negative clock eigenvalue.
The duration operator must still be supplied. Replacing a duration
distribution by its mean loses cumulants, while replacing coherent
intermediate amplitudes by orthogonal history labels changes the transfer.

\subsection*{H.15 Joint geometry, orbital motion, and clock}
\label{aqg:orbital}

The retained space has the split
\[
\mathcal H_{\rm ren}\otimes\mathcal H_X
\otimes(\mathbb C|0\rangle\oplus Q),\qquad \dim Q=9.
\]
Use cochains on vertices, edges, faces, and tetrahedra for
\(\mathcal H_X\), as in the Dirac--K\"ahler construction
\citep{BecherJoos1982}; H.9a, paragraph~2g, gives the native
relativistic and doubler-free conditions for this placement. The internal \(Q\) fiber remains intact.
With incidence coboundary \(d_p\) and positive cochain metric \(M_p\),
\[
d_p^*=M_p^{-1}d_p^\dagger M_{p+1},\qquad
D_g=d+d^*,\qquad
\Omega_g(m)=\sqrt{D_g^2+m^2I}.
\]
The last expression selects a positive-energy orbital completion.
Hamiltonian formulas use \(\hbar=c=L_*=1\), and the native rest
input is \(m=\epsilon_e\).
Harmonic modes have \(D_g=0\) and energy \(m\).
The mass derivative is
\[
\partial_m\Omega_g=m(D_g^2+m^2I)^{-1/2}.
\]
It equals one at rest and lies in
\([m/\sqrt{\|D_g\|^2+m^2},1]\) at finite cutoff.
For a supplied positive lapse and mass-independent self-adjoint
geometric coupling,
\[
H(m)=\widehat H_g+
\bigoplus_g N_g^{1/2}\Omega_g(m)N_g^{1/2},\qquad
\partial_mH\succeq0.
\]
Thus \(U_t=e^{-itH}\) is unitary and
\[
\mathcal T(t)=iU_t^\dagger\partial_mU_t
=\int_0^tU_s^\dagger(\partial_mH)U_s\,ds
\]
is a positive clock response with the usual composition law.
Positive Euclidean operators can still have signed matrix elements.

When the coarse symbol is \(E=\sqrt{h^{ij}p_ip_j+m^2}\),
\(\dot x^i=Nh^{ij}p_j/E\) and \(d\tau/dt=Nm/E\).
The Legendre transform gives \(-m\int d\tau\), agreeing with H.14
without an additional gravitational charge.
The assumed spatial symbol is a physical requirement, tested below.
The clock has no uniform strictly positive lower bound as the momentum
cutoff grows. The square root also retains both signs of the
Dirac--K\"ahler spectrum as orbital components; particle/antiparticle
assignment and reduction to one physical fermion species remain open.

\paragraph{One finite geometric transfer.}
Let \(g_0\) be the five-tetrahedron boundary of a four-simplex.
Replace one tetrahedron by four regular tetrahedra to obtain \(g_1\).
These are different physical metrics, with cochain dimensions 30 and 44.
The four boundary-pair families at \(\alpha=7\) contain
\(1920,1992,1992,2064\) causal slabs.
Their angular weights, with one common normalization, give
\[
A=\begin{pmatrix}
0.1169508311&0.3422031803\\
0.3422031803&1
\end{pmatrix},\qquad K_g=A^\dagger A.
\]
The smaller eigenvalue of \(A\) is negative.
The eigenvalues of \(K_g\) are approximately
\(1.85597649\times10^{-8}\) and \(1.24788351\).
Define
\[
H_g=-\frac{\log K_g}{\Delta t},\qquad
\Delta t=2a\sqrt{7+3/8}=13.94846613.
\]
This synchronization is twice a chosen boundary-normal duration;
moving paths have their own elapsed times.
The logarithm is evaluated from the eigenvalues of \(A\) before
squaring to retain numerical precision.

Raw stellar interpolation has maximum singular value \(1.48324\).
Using it as geometry transport would lower the marked geometric
ground energy by \(0.131992\), overwhelming the native rest gap.
Normalize the harmonic and orthogonal excited sectors separately by
rectangular polar factors. With harmonic projectors \(P_0,P_1\), the
selected map \(R\) obeys
\[
R^\dagger R=I,\qquad RP_0=P_1R,\qquad
\widehat H_g=
\begin{pmatrix}
(H_g)_{00}I&(H_g)_{01}R^\dagger\\
(H_g)_{10}R&(H_g)_{11}I
\end{pmatrix}.
\]
It need not intertwine excited operators on different physical metrics.
Its lower bound is the empty-sector ground energy \(E_{g,0}\).
At unit reference lapse, \(\Omega_g\ge mI\), and transported harmonic
modes attain the joint bound. Consequently,
\[
\boxed{E_{\rm marked,0}(m)-E_{\rm vacuum,0}=m.}
\]
The native tiny gap is protected algebraically.
At control mass \(0.23\), the gap error is \(4.45\times10^{-16}\);
the vacuum transfer agrees with the enumerated \(K_g\) to the same
accuracy. The clock derivative and real-time unitarity checks give
\(2.84\times10^{-10}\) and \(1.76\times10^{-15}\).
A packet moves while coherently changing geometry in this finite model.
These tests do not determine its long-distance gravitational self-energy.

\paragraph{Microscopic metric preserving the entire operator.}
For two descriptions of the same supplied coarse geometry, let
interpolation \(P\) and retraction \(R\) satisfy
\[
RP=I,\qquad d_fP=Pd_c,\qquad Rd_f=d_cR.
\]
Set \(Q_f=I-PR\) and choose any positive auxiliary metric \(M_a\):
\[
\boxed{M_f=R^\dagger M_cR+Q_f^\dagger M_aQ_f.}
\]
Every fine vector is uniquely \(Px+y\) with \(Ry=0\), and
\[
\|Px+y\|_f^2=x^\dagger M_cx+y^\dagger M_ay.
\]
The coarse image and complement are orthogonal and invariant under
\(d_f\) and its metric adjoint. In orthonormal coordinates the
encoding \(J\) therefore satisfies
\[
J^\dagger J=I,\qquad D_fJ=JD_c.
\]
Functional calculus preserves the full orbital dispersion, clock,
and evolution:
\[
\Omega_fJ=J\Omega_c,\qquad
(\partial_m\Omega_f)J=J(\partial_m\Omega_c),\qquad
e^{-it\Omega_f}J=Je^{-it\Omega_c}.
\]
This inverse construction selects a microscopic inner product for a
supplied operator. Any unwanted coarse dispersion or species is
preserved with it. The generally off-diagonal metric defines an
orbital code, without asserting equality to an unmodified fine Regge
metric.

Two stellar refinements give dimensions \(30\to44\to58\).
Direct and sequential interpolations agree exactly in the calculation;
Dirac and clock intertwining errors are \(5.99\times10^{-15}\) and
\(3.85\times10^{-15}\). The two harmonic modes remain, with no extra
zero modes in the tested complement.
Encoding each of the two physical geometries and transporting their
coupling gives a joint \(74\to102\) dimensional lift, with \(Q\)
understood in both spaces. Generator, clock, and evolution errors are
\(1.83\times10^{-13}\), \(7.79\times10^{-15}\), and
\(3.11\times10^{-14}\).
Two transported external test-source operators agree to
\(1.83\times10^{-13}\); they are not a calibrated held-mass force.
The complement does not mix with the prepared code and is not required
to lie above every excited code state.

\paragraph{Isometric reduction and source response.}
Let $J:\mathcal H_c\to\mathcal H_f$ be an isometry, $J^\dagger J=I$, and suppose that self-adjoint operators obey $H_fJ=JH_c$. The retained projector and operator lift are
\[
P=JJ^\dagger,\qquad \iota(A)=JAJ^\dagger,
\qquad P^2=P=P^\dagger.
\]
Direct multiplication gives
\[
\iota(A)\iota(B)=\iota(AB),\quad
J^\dagger\iota(A)J=A,\quad
\operatorname{Tr}_f\iota(A)=\operatorname{Tr}_c A,\quad
A\ge0\ \Longrightarrow\ \iota(A)\ge0.
\]
The lift sends the retained identity to $P$. Orthogonality of $P$ fixes the physical norm. Intertwining gives $H_f^nJ=JH_c^n$ for every nonnegative integer $n$; finite-dimensional spectral calculus then gives the corresponding evolution and massive square root whenever their defining functions are regular on the spectrum. Differentiation in mass preserves the clock identity when $J$ is mass independent; a varying $J$ requires the connection terms below. Composed refinements obey the same identities. The inverse cochain metric above supplies one realization of these hypotheses.

Physical sources require this relation for a source-dependent family. For a differentiable parameter $s$,
\[
\partial_s(J^\dagger H_fJ)
=J^\dagger(\partial_sH_f)J
+(\partial_sJ^\dagger)H_fJ+J^\dagger H_f(\partial_sJ).
\]
The last two terms account for motion of the retained subspace. A source calculation must also retain the state, discarded-sector determinant, and normalization derivatives present in its generating functional. Appendix~H.10 gives the Gaussian elimination with these sources retained; N.8 gives the corresponding positive reconstruction and adjoint force map. These identities preserve a supplied coupled operator. Its locality, common causal cone, and physical source vertex are separate conditions.

\paragraph{Common light cone and species.}
The full Whitney metric reproduces a common low-momentum speed on
flat embedded controls. Giving every tetrahedron its intrinsic regular
metric instead leaves, on the periodic BCC test,
\[
v_0=\sqrt{5/6},\qquad v_T=\sqrt{4/5},\qquad v_3=\sqrt{3/4}.
\]
Including off-diagonal integrals therefore does not remove this split,
in agreement with the native-complex result of H.9a, paragraph~2g.
A periodic inverse normalization using the exterior-form metric of a
chosen coarse geometry, fixed by its intrinsic volume before the
speed measurement, restores the leading common cone in both tested
blocks. The irregular block retains a subleading transverse splitting,
approximately linear in the tested wave number and about
\(9\times10^{-5}\) relatively at the smallest diagonal momentum.
This periodic correction has not been assembled over arbitrary
changing CDT boundaries. Its compatibility with the joint finite
model, the species interpretation, and the interacting gravitational
phase remain physical tests.

\section*{Appendix I: Mass and Gauge Extensions}

Appendix~I develops possible particle and gauge extensions of the entanglement logic. The closed weak-field derivation does not depend on them.

\subsection*{I.1 Charged-lepton spectrum from the shell algebra}

The electron is the ground-state one-bit charged defect, the $N=0$ configuration whose seven-channel dressing sets the substrate length (Appendix~D.4). The muon and tau are heavier shell excitations of the same defect. Each adds one radial entanglement shell, realized microscopically as the exclusion of one face label from the seven-state alphabet of Appendix~B. The closure family has three nondegenerate shell grades. Identifying charged-lepton shells exclusively with its closure-response modes gives the selected three-generation branch. Their baseline masses follow from the shell recurrence load and singlet log-overlap prescription, while the decorated marked-transfer vertex supplies the finite response. Because the ratios carry no dimensionful input, they test the same marked action that enters $G_*$. 

With that picture fixed, the spectrum is written as a shell expansion,
\[
\log\frac{m_N}{m_e}=B_0 N+A_0 N^2,
\qquad
N=0,1,2,
\]
with the electron the $N=0$ ground state and the muon and tau the $N=1,2$ radial entanglement-shell excitations of the same core structure.

\paragraph{Three-generation termination from closure-spectrum collapse.}
Combinatorial injectivity alone allows up to $N=3$ shells, since the reduced alphabet must retain at least four distinct labels to support an injective face assignment ($7-N\geq 4$). The sharper structural bound comes from the admissibility-closure machinery. For the reduced ensemble at shell $N$, the closure invariant $K^2$ inherits a discrete spectrum whose dispersion gives the admissibility kernel $e^{-\eta K^2}$ real weight. For $N=0,1,2$ the $K^2$ spectrum has multiple distinct values, and the closure equation $\langle K^2\rangle_\eta=3/(2\eta)$ has a non-degenerate solution.

At $N=3$ the spectrum collapses. With only four labels left, every injective assignment to the four faces uses all of them, and $K^2$ depends on the labels only through the permutation-invariant sums $S$ and $\Sigma^2$, so every microstate of the reduced four-label ensemble carries the same value of $K^2$. With a single spectral value $K_0^2$ the closure equation $\langle K^2\rangle_\eta=3/(2\eta)$ still fixes $\eta=3/(2K_0^2)$ formally; what fails at $N=3$ is selection, not solvability. The kernel $e^{-\eta K^2}$ acts trivially when every microstate carries the same $K^2$, so it weights nothing and leaves no admissibility dispersion to define a branch. Shell $N=3$ therefore does not define a non-degenerate admissibility-closed branch in the same sense as the earlier shells: the cutoff rests on the collapse of the $K^2$ dispersion.

The closure-fluctuation statement can be made exact. For $p_{N,\eta}(b)=Z_N^{-1}e^{-\eta K_N^2(b)}$, define
\[
|s_N\rangle=\sum_b\sqrt{p_{N,\eta}(b)}
\bigl(K_N^2(b)-\langle K_N^2\rangle\bigr)|b\rangle.
\]
The Fisher information for $\eta$ and the fluctuation norm coincide:
\[
I_N(\eta)=\sum_bp_{N,\eta}(b)
\left(\partial_\eta\ln p_{N,\eta}(b)\right)^2
=\operatorname{Var}_{N,\eta}(K_N^2)=\langle s_N|s_N\rangle.
\]
Every four-label alphabet has $I_3=0$ and $|s_3\rangle=0$. Every five- and six-label alphabet has positive dispersion. The selected particle identification assigns an additional lepton shell only to a nonzero closure-response mode and identifies the equivalent label realizations within each grade. Under that identification the allowed grades are $N=0,1,2$. The theorem excludes the fourth closure mode; excluding particles built from other operators is part of the particle-identification premise.

\paragraph{Linear coefficient from the first-shell reduced-alphabet entropy.}
The linear coefficient is the entropy cost of adding one shell excitation. The selected shell encoding deletes one magnetic label and retains injective assignments on the other six. Its oriented state count is
\[
\Omega_1=2P(6,4)=720.
\]
Let $p_b>0$ be the stationary weight of each shell state and let $T_b$ be the first return time to $b$ under memoryless renewal. Conditional on $b$, the return time is geometric,
\[
P(T_b=t\mid b)=p_b(1-p_b)^{t-1},
\qquad
\mathbb E[T_b\mid b]=\frac1{p_b}.
\]
A stationary choice $B\sim p$ therefore gives the exact finite-state recurrence identity \citep{Kac1947}
\[
\boxed{
\mathbb E_{B\sim p}\!\left[\mathbb E(T_B\mid B)\right]
=\sum_b p_b\frac1{p_b}=\Omega_1=720.}
\]
The geometric mean of the same conditional mean return times is instead
\[
\exp\!\left(\mathbb E_B\ln\mathbb E[T_B\mid B]\right)
=\exp\!\left(-\sum_b p_b\ln p_b\right)=e^{H_1}.
\]
The mass functional reads the shell load from the stationary arithmetic return time, so $B_0=\ln\Omega_1$. The electron determinant uses the distinct state-weighted logarithmic likelihood. With the closure parameter solved separately for each reduced six-label alphabet, its Shannon entropy satisfies
\[
\ln720-H_1\simeq0.00508\text{--}0.00782
\]
over the seven excluded labels. Substituting $e^{H_1}$ for $720$ would lower a one-shell mass factor by about $0.5\%$--$0.8\%$. The two factors are therefore distinct recurrence statistics. Reading the arithmetic recurrence load as an inertial pole-mass factor is the particle-sector matching premise of the shell model.

\paragraph{Why the same $720$ at every shell.}
Each added shell uses a copy of the same six-label encoding. The $N$-shell register has Cartesian support of size $720^N$, and the stationary arithmetic return identity above equals $720^N$ for every full-support distribution on that register; statistical independence is not needed for this support statement. The generation test separately studies the reduced alphabet with $7-N$ labels. These are two specified operations in the shell model.

The scalar singlet coupling supplies a common channel but does not by itself construct the six-label basis. Indeed, for $Q=I-|u\rangle\langle u|$,
\[
\langle a|Q|b\rangle=\delta_{ab}-\frac17,
\]
so all seven projected label vectors remain nonzero. The chosen shell encoding consists of an isometry from a six-dimensional shell-label register to $u^\perp$, together with injectivity imposed in that register's orthonormal label basis. The same register is used for each shell. This encoding yields $2P(6,4)=720$; rank-one singlet coupling alone fixes neither that basis nor its injectivity rule.

\paragraph{Quadratic coefficient from singlet-projection shell algebra.}
The quadratic coefficient follows from three premises already in the framework: (i) mass arises from scalar capacity response (Postulate II); (ii) the scalar EFT response is quadratic in the total source; (iii) the coarse scalar branch projects onto the permutation-invariant singlet of the seven-state face alphabet.

Let
\[
\mathcal H_7=\mathrm{span}\{|m\rangle:m=-3,\ldots,3\},
\qquad
|u\rangle=\frac{1}{\sqrt 7}\sum_{m=-3}^{3}|m\rangle,
\qquad
P_{\rm sing}=|u\rangle\langle u|.
\]
A shell excitation that marks the face state $|a_r\rangle$ contributes to a coherent shell source
\[
|J_N\rangle=\sum_{r=1}^{N}|a_r\rangle.
\]
The coarse scalar field $\delta S$ is permutation-symmetric on face labels and therefore couples only through $P_{\rm sing}$. The scalar-channel response is quadratic in the total source, so the relevant object is the source norm
\[
\langle J_N|P_{\rm sing}|J_N\rangle
=
\sum_{r,s=1}^{N}\langle a_r|P_{\rm sing}|a_s\rangle.
\]
Because $\langle a|P_{\rm sing}|b\rangle=1/7$ for any $a,b$, this reduces to $N^2/7$ independent of which specific face labels are excluded. Including the binary orientation degeneracy of the boundary ensemble, the orientation-summed scalar transfer weight per ordered shell incidence is
\[
\mathcal M_{rs}=2\,\langle a_r|P_{\rm sing}|a_s\rangle=\frac{2}{7}.
\]
The ordered-pair count decomposes as $N^2=N+2\binom{N}{2}$: $N$ self-incidences along the diagonal and $2\binom{N}{2}$ directed cross-incidences. Two distinct steps enter here, and only the first is the quadratic EFT response. The quadratic singlet response fixes that there are $N^2$ ordered incidences, each carrying the per-incidence overlap $\mathcal M_{rs}=2/7$; it does not by itself make the mass multiplicative. The passage from these $N^2$ additive incidences to a multiplicative factor is the log-overlap mass map: the mass is the exponential of a dressing free energy whose pairwise term is $\sum_{r,s}\ln\mathcal M_{rs}=N^2\ln(2/7)$, the same log-overlap form that carries the gravitational dressing (Appendix~H.3) and stated explicitly below. With that map, the $N$-shell dressing is
\[
\prod_{r,s=1}^{N}\mathcal M_{rs}=\left(\frac{2}{7}\right)^{N^2},
\]
so
\[
A_0=\ln\frac{2}{7}=-1.253\ldots,
\]
compared with the empirical fit $A_0=-1.255$ (agreement at $0.15\%$).

The quadratic scalar response fixes the ordered-bilinear structure under the three premises just stated: the source norm $\langle J_N|P_{\rm sing}|J_N\rangle$ runs over both indices independently because the integrated-out scalar contribution is quadratic in $J_N$.

\paragraph{Positive realization of the coherent overlap.}
The $N^2$ law is the characteristic function of one common scalar fluctuation. Let $X$ be a centered Gaussian variable with variance
\[
v=-2\ln q,
\qquad q=\frac27,
\]
and let an $N$-shell coherent source acquire phase $NX$. Its averaged overlap is
\[
\boxed{
\mathbb E\!\left[e^{iNX}\right]
=e^{-N^2v/2}=q^{N^2}.}
\]
This gives an explicit positive overlap kernel with exactly the shell factor used above. The shared fluctuation is essential: independent shell phases would give $q^N$. The quadratic exponent therefore records coherent coupling of all shells to the same permutation-singlet scalar mode.

\paragraph{Physical mass readout.}
On the selected grades let $\widehat N$ have eigenvalues $0,1,2$. The log-overlap prescription defines the positive mass operator
\[
\widehat M=m_e\exp\!\left[\widehat N\ln720+
\widehat N^2\ln(2/7)+\ln Z_{\widehat N}\right],
\]
where $Z_0=1$, $Z_1=1+\zeta_*$, and $Z_2=(1+\zeta_*)(1+2\zeta_*/7)$. Appendix~H.9a places this operator in the covariant Einstein--Dirac action of the selected ordinary branch. Its eigenvalues are then the free pole masses and source the same metric through the Dirac stress tensor. This specifies their physical readout within the shell model. Deriving the mass and kinetic operators together from a coupled native defect Hamiltonian remains a separate matching calculation.

\paragraph{Mass ladder and numerical accuracy.}
Combining the two coefficients,
\[
\boxed{
\frac{m_N^{(0)}}{m_e}=720^N\left(\frac{2}{7}\right)^{N^2},
\qquad N=0,1,2.
}
\]
Equivalently,
\[
\log\frac{m_N}{m_{N-1}}=\ln 720+(2N-1)\ln\frac{2}{7},
\]
so each added shell contributes $\ln 720$ plus an odd number of new scalar-overlap incidences ($1,3,5,\ldots$, summing to $N^2$). The second difference is fixed at
\[
\Delta^2\log m_N=2\ln\frac{2}{7}=-2.506\ldots,
\]
compared with the empirical $-2.509$ ($0.15\%$). The marked vertex of H.9 then supplies
\[
Z_\mu=1+\zeta_*,
\qquad
Z_{\tau,2}=1+\frac27\zeta_*.
\]
Thus
\[
\frac{m_\mu}{m_e}
=720\frac27Z_\mu
=206.768280237
\qquad(-0.535\sigma),
\]
\[
\frac{m_\tau}{m_e}
=720^2\left(\frac27\right)^4Z_\mu Z_{\tau,2}
=3477.343310
\qquad(-0.125\sigma).
\]
\paragraph{The dressing map, explicitly.}
The three factors follow from one construction. Write the $N$-th charged lepton as a coherent source raised on the electron ground dressing by exciting $N$ shells, $\mathcal S_N=\sum_{k=1}^{N}s_k$, where each $s_k$ excludes one face label from the seven-state alphabet. Its dressing free energy is a log-overlap functional of the kind used for the gravitational dressing in Appendix~H.3, and it separates into a per-shell term and a pairwise term. Each shell carries the bare reduced-alphabet support term $\ln\Omega_1=\ln 6!=\ln 720$, the permutation count of the six surviving labels, which builds the linear factor $720^N$. Each ordered pair of shells contributes a singlet log-overlap: a single label projected on the closure singlet $u=\tfrac1{\sqrt7}(1,\dots,1)$ of Section~8 has weight $|\langle u|m\rangle|^2=1/7$, and orientation summation doubles it to $2/7$. The $N$ shells form $N^2$ ordered pairs, so the pairwise term is $N^2\ln(2/7)$, and exponentiating the free energy returns
\[
\frac{m_N^{(0)}}{m_e}=\Omega_1^{\,N}\left(\frac{2}{7}\right)^{N^2}=720^N\left(\frac{2}{7}\right)^{N^2}.
\]
The coefficients are fixed: $720=6!$, the $N^2$ pair count is combinatorial, and $2/7$ is the orientation-summed singlet projection behind $\Omega_{\rm tet}=2\,P(7,4)=1680$. The marked factors are finite traces of the same decorated vertex used in the electron scale.

\paragraph{The support exponent, tested against the spectrum.}
The ladder also tests the exponent in the support-to-length identification of Appendix~D.4, which reads the physical ratio as the first power of the entropy support, $\lambda/L_*\propto e^{H}$. Letting that power float, $\lambda/L_*\propto e^{\alpha H}$, turns each lepton ratio into a measurement of $\alpha$, since $m_N/m_e=\bigl[720^N(2/7)^{N^2}\bigr]^{\alpha}$:
\[
\alpha_N=\frac{\ln(m_N/m_e)_{\rm PDG}}{\ln\bigl[720^N(2/7)^{N^2}\bigr]},
\qquad
\alpha_\mu=1.0010,
\qquad
\alpha_\tau=1.0008.
\]
The leading shell values give $\alpha_\mu=1.0010$ and $\alpha_\tau=1.0008$, while square-root and quadratic support maps miss by orders of magnitude. The decorated action does not float this exponent. It keeps the first-power support map and computes the finite response factors separately.

\paragraph{Status.}
The selected shell model fixes the baseline coefficients and marked response. The shell integers alone give $m_\mu/m_e=205.71$ and $m_\tau/m_e=3454.6$, within $0.51\%$ and $0.65\%$ of the measured ratios; the marked response raises them to $206.768280237$ and $3477.343310$, at $-0.535\sigma$ and $-0.125\sigma$ (Section~13.5 and Appendix~H.9). Appendix~L records the provenance of the marked response. Its three-generation interpretation uses the closure-mode identification, six-label encoding, and log-overlap mass functional stated above. The finite counts and traces are exact within those choices. A native charged-defect Hamiltonian must reproduce their physical mass and species assignments.

The electron remains the elementary weak-field anchor; the heavier charged leptons are correlated outputs of the same marked shell algebra, not closure-defining inputs. Composite hadrons remain part of the dressed bound-state entropy program. The neutrino and quark sectors are not treated here. Extending the closure-spectrum-collapse argument to those defect classes remains a separate task.

\subsection*{I.2 Gauge structure and the conditional projective color sector}

The framework's primary target remains the gravitational and dark sector. The Abelian and weak sectors below are substrate-natural hosting arguments whose physical Standard Model identification remains external. Color has a further conditional construction. A colored defect is identified with one persistent open tetrahedral-routing leg, transported reversibly until singlet closure. This identification fixes the projective qutrit transport algebra and the adjoint form of the local pure-gauge action. A minimal kinetic completion then fixes its primitive adjoint transfer. This does not derive the global $SU(3)$ endpoint lift, chiral quarks, hypercharge assignments, quark masses, confinement, or the finite matching to a continuum QCD scheme.

\paragraph{Baseline-redundancy template.}
The same logic that organizes the gravity sector extends to any conserved-charge sector. For a conserved charge $Q$, introduce an entropy-like potential $S_Q(x)$ and require that physical observables depend only on differences of that potential, not on its absolute baseline. Localizing that redundancy requires a compensating connection. In the Abelian case,
\[
D_\mu S_Q=\partial_\mu S_Q-qA_\mu,
\qquad
S_Q\to S_Q+\alpha(x),
\qquad
A_\mu\to A_\mu+\frac{1}{q}\partial_\mu\alpha,
\]
allowing a Maxwell kinetic term for $A_\mu$ once its dynamics is supplied. Appendix~S gives an explicit charged-link construction and separates that additional physical input from the gauge redundancy. The non-Abelian generalization, for multiplet-valued entropic potentials transforming under a compact Lie group $G$, gives Yang--Mills covariant derivatives and field strengths in the standard form.

The template tells us what form a gauge interaction takes once a conserved-charge sector with baseline redundancy is present. It does not select which specific groups are physically realized.

\paragraph{Substrate-natural components.}
Two ingredients of Standard Model gauge structure can be represented within the substrate kinematics.

\emph{$U(1)_Y$.} The baseline-redundancy template permits an Abelian conserved-charge sector of the hypercharge type. The associated entropy potential is conserved-charge in character; localizing its baseline redundancy permits a connection but does not determine its kinetic term or coupling. The argument is independent of the gravity-sector derivation, applying the same template to a separate conserved-charge potential.

\emph{$SU(2)_L$.} The half-integer face data carry an $SU(2)$-representation structure. On the marked-fusion incidence branch, the three-dimensional vector response and the complement theorem select the primitive spin $j_0=\tfrac32$ (Section~5 and Appendix~B); without that identification, the same spin is the surviving discrete branch. The static $K^2$ ensemble breaks this $SU(2)$ to its $U(1)$ Cartan through the magnetic-quantum-number projection, but the full $SU(2)$ remains present in the quantum face data discarded by that projection. Identifying this representation-theoretic $SU(2)$ with the internal weak group $SU(2)_L$, including chirality and doublet assignments, remains external to the present derivation.

\paragraph{Three tetrahedral routes and the closed-cell obstruction.}
Four faces admit exactly three opposite-face pairings,
\[
A=(12\,|\,34),\qquad
B=(13\,|\,24),\qquad
C=(14\,|\,23).
\]
Tetrahedral permutations act transitively on this set, so no route is preferred. On a closed effective cell define
\[
R_A=m_1m_2+m_3m_4,
\quad
R_B=m_1m_3+m_2m_4,
\quad
R_C=m_1m_4+m_2m_3.
\]
Their sum is fixed by the closure invariant,
\[
R_A+R_B+R_C
=\frac12\left[\left(\sum_i m_i\right)^2-\sum_i m_i^2\right]
=72-\frac32K^2.
\]
The symmetric route is therefore the scalar closure mode already read by gravity. The exact $1680$-state sum at $\eta_*$ gives the connected covariance
\[
C_R=
\begin{pmatrix}
27.45571280&-7.84449077&-7.84449077\\
-7.84449077&27.45571280&-7.84449077\\
-7.84449077&-7.84449077&27.45571280
\end{pmatrix}.
\]
With
\[
P_{\mathbf1}=\frac13\mathbf1\mathbf1^T,
\qquad
P_{\mathbf2}=I-P_{\mathbf1},
\]
this becomes
\[
\boxed{C_R=11.76673126\,P_{\mathbf1}+35.30020358\,P_{\mathbf2}}.
\]
The singlet eigenvalue is fixed analytically,
\[
\lambda_{\mathbf1}
=\frac13\operatorname{Var}(R_A+R_B+R_C)
=\frac34\operatorname{Var}(K^2).
\]
A closed cell thus supplies $\mathbf3_{\rm route}=\mathbf1_{\rm closure}\oplus\mathbf2_{\rm imbalance}$. It cannot furnish an irreducible color triplet as a closed scalar observable.

\paragraph{Primitive route geometry.}
The route carrier belongs to the primitive spin layer, before neighboring slots fuse into the effective $j_{\rm eff}=3$ link. Couple the four primitive $j_0=3/2$ slots according to $A$, $B$, or $C$, taking each pair through its maximal $J=3$ channel and closing the two intermediate spins to a singlet. Standard $SU(2)$ recoupling gives
\[
\langle A|B\rangle
=\langle A|C\rangle
=\langle B|C\rangle
=-\frac1{20}.
\]
The route Gram matrix is therefore
\[
G=
\begin{pmatrix}
1&-1/20&-1/20\\
-1/20&1&-1/20\\
-1/20&-1/20&1
\end{pmatrix}
=\frac9{10}P_{\mathbf1}+\frac{21}{20}P_{\mathbf2}.
\]
Its eigenvalue ratio is
\[
\boxed{\frac{\lambda_{\mathbf2}}{\lambda_{\mathbf1}}=\frac76},
\]
the same exact ratio that supplies the isotropically averaged one-label-exclusion source benchmark $\ln(7/6)$ in Appendix~C.5. This equality follows before any color datum is introduced. It is a cross-sector recoupling identity, not a determination of the strong coupling.

The layer ordering is fixed by the construction. Two primitive shared-face slots first form $V_{3/2}\otimes V_{3/2}$, maximum-capacity fusion selects $V_3$, and the four resulting effective links carry the $K^2$ closure ensemble. The color route therefore resides in an open primitive record that survives between fusion events. Appendix~H.9 gives the complementary result that the nonmaximal primitive fusion channels $V_0\oplus V_1\oplus V_2$ are exactly the nine-state representation already required by the present/history marked fiber.

\paragraph{Persistent open routing and the gauge algebra.}
The conditional color identification is precise: a colored fermionic defect carries one unclosed tetrahedral-routing leg, and transport preserves its correlations with an arbitrary reference until the leg closes into a singlet. The three unresolved route alternatives span the qutrit
\[
\mathcal H_{\rm route}\simeq\mathbb C^3.
\]
A reversible quantum channel between equal finite-dimensional systems is unitary conjugation,
\[
\mathcal E_U(X)=UXU^\dagger,
\qquad U\in U(3).
\]
The channel is unchanged under $U\mapsto e^{i\theta}U$, so its physical transport group is
\[
PU(3)=U(3)/U(1)\simeq SU(3)/Z_3,
\]
with Lie algebra $\mathfrak{su}(3)$. This fixes the local connection algebra and its adjoint representation. It does not choose the global lift of the projective channel to endpoint amplitudes. A fundamental $SU(3)$ endpoint and $Z_3$ triality require such a lift together with an identification of the central $U(1)$ phase; those are part of the open matter-sector completion.

The endpoint observable algebra provides the same result in a complementary form. A route-resolving endpoint supplies $P_A=|A\rangle\langle A|$; tetrahedral permutations generate $P_B$ and $P_C$. If $S_{ab}$ carries route $b$ to route $a$, then
\[
P_aS_{ab}P_b=|a\rangle\langle b|\equiv R_{ab}.
\]
The nine rank-one operators span $M_3(\mathbb C)=\mathbf1\oplus\mathbf8$ because the three route vectors are linearly independent. Their products obey $R_{ab}R_{cd}=G_{bc}R_{ad}$. Choosing an orthonormal basis of the same span gives matrix units. The identity is the common phase direction and the traceless Hermitian operators are the eight connection directions. The route endpoint is physically resolved, while its basis in transit is gauge covariant. These endpoint operators connect the discrete $S_3$ triple to a continuous gauge algebra; $S_3$ alone would generate only $\mathbb C\oplus M_2(\mathbb C)$.

Choosing route bases independently at the cells gives the link transformation
\[
U_{yx}\longmapsto V_yU_{yx}V_x^\dagger.
\]
Around a loop based at $x$, $U_p=\prod_{\ell\in p}U_\ell$ therefore transforms by conjugation, $U_p\mapsto V_xU_pV_x^\dagger$. This is the local gauge redundancy of the transported route frame.

\paragraph{Adjoint Wilson form from channel-return fidelity.}
The transported object fixed above is a channel on route density operators, so its return cost must be unchanged by $U_p\mapsto e^{i\theta}U_p$. The difference of two normalized route projectors is traceless and Hermitian; on that irreducible adjoint space, the positive frame-invariant quadratic form is unique up to an overall scale. The scale is absorbed into $\kappa_{\rm open}$. For a normalized route state $q$, set $P_q=|q\rangle\langle q|$ and define the resulting Hilbert--Schmidt channel-return failure
\[
\Delta_p(P_q)
=\left\|P_q-U_pP_qU_p^\dagger\right\|_{\rm HS}^2
=2\left(1-|\langle q|U_p|q\rangle|^2\right).
\]
For an unbiased qutrit, the Haar second moment gives
\[
\overline{|\langle q|U_p|q\rangle|^2}
=\frac{\operatorname{Tr}(U_pU_p^\dagger)+|\operatorname{Tr}U_p|^2}{3(3+1)}
=\frac{3+|\operatorname{Tr}U_p|^2}{12}.
\]
Using $\chi_{\mathbf8}(U)=|\operatorname{Tr}_{\mathbf3}U|^2-1$,
\[
\overline{\Delta_p}
=\frac32\left(1-\frac{|\operatorname{Tr}_{\mathbf3}U_p|^2}{9}\right)
=\frac43\left(1-\frac18\chi_{\mathbf8}(U_p)\right).
\]
A cost $\kappa_{\rm open}$ per unit return failure therefore gives the adjoint Wilson form \citep{Wilson1974},
\[
\boxed{
S_{\rm gauge}
=\beta_{\rm adj}\sum_p
\left(1-\frac18\operatorname{Tr}_{\mathbf8}\operatorname{Ad}U_p\right),
\qquad
\beta_{\rm adj}=\frac43\kappa_{\rm open}.
}
\]
This action is well defined on $PU(3)$ and is blind to the center, exactly as the channel construction requires. A fundamental-trace Wilson term would require the unresolved $SU(3)$ endpoint lift. The group argument alone does not fix $\kappa_{\rm open}$.

\paragraph{Minimal primitive kinetic normalization.}
The primitive recoupling matrix fixes the missing local coefficient after one named microscopic choice: geometric gluing uses the identity pairing on the primitive boundary Hilbert space before an independent kinetic operator is added. In the nonorthogonal route basis, the identity pairing has matrix $G$. A direct branch $K=sG$ admits a trace-preserving reversible dilation only when $K^\dagger K\le I$. Maximum faithful throughput saturates the largest eigenvalue, so
\[
s=\frac{20}{21},
\qquad
K=\frac{20}{21}G
=P_{\mathbf2}+\frac67P_{\mathbf1},
\qquad
L=\sqrt{I-K^\dagger K}
=\frac{\sqrt{13}}7P_{\mathbf1}.
\]
Together, $K$ and $L$ define an isometry and therefore admit a reversible unitary dilation. Only the symmetric closure component enters the complementary branch; the primitive imbalance doublet has unit direct amplitude.

Transport through the vacuum carries no preferred absolute route frame. Replacing $P_{\mathbf1}$ by an arbitrary rank-one projector $P=|n\rangle\langle n|$ and averaging the equivalent frames gives
\[
K_P=I-\frac17P,
\qquad
L_P=\frac{\sqrt{13}}7P.
\]
The qutrit Haar identities
\[
\int dP\,P=\frac{I}{3},
\qquad
\int dP\,PXP=\frac{X+\operatorname{Tr}(X)I}{12}
\]
then yield the covariant channel
\[
\begin{aligned}
\mathcal E(X)
&=\int dP\,\bigl(K_PXK_P^\dagger+L_PXL_P^\dagger\bigr)\\
&=\frac{13}{14}X+\frac1{42}\operatorname{Tr}(X)I
=1_{\mathbf1}\oplus\frac{13}{14}I_{\mathbf8}.
\end{aligned}
\]
Thus
\[
\boxed{t_8=\frac{13}{14}},
\qquad
-\ln t_8=\ln\frac{14}{13}.
\]
For generators normalized by $\operatorname{Tr}(T^aT^b)=\delta^{ab}/2$, one primitive transfer interval $\tau_{\rm r}$ has the isotropic adjoint generator
\[
\mathcal L(X)
=-\frac{\ln(14/13)}{3\tau_{\rm r}}
\sum_{a=1}^{8}[T^a,[T^a,X]].
\]
In the heat-kernel coordinate $t_R=\exp[-g_{\rm HK}^2C_R/2]$ with $C_A=3$,
\[
\boxed{g_{\rm HK,step}^2=\frac23\ln\frac{14}{13}=0.04940531477.}
\]
This value was obtained from the primitive geometry without a QCD target. It is conditional on the identity kinetic pairing, maximum faithful throughput, and route-frame isotropy.

\paragraph{Causal blocking on the diamond regulator.}
The paper already uses the four-valent diamond graph as the embedded adjacency realization of the tetrahedral bond frame. Its minimal product causal completion is $D_{\rm diamond}\times\mathbb Z$, with one temporal successor at each spatial site and no diagonal shortcuts. Let $d$ be the spatial bond length, $a_t$ the temporal spacing, $r=a_t/d$, and let $\rho$ be the ratio of temporal to spatial primitive heat times. The shortest mixed time--space loops have four links; the shortest spatial loops are diamond hexagons with six links. Their per-site oriented area tensors are
\[
M_{\rm E}^{ab}=\frac23r^2d^4\delta^{ab},
\qquad
M_{\rm M}^{ab}=\frac{32}{9}d^4\delta^{ab}.
\]
The diamond primitive translations $a_1=d(n_1-n_2)$, $a_2=d(n_1-n_3)$, $a_3=d(n_1-n_4)$ have determinant $16d^3/(3\sqrt3)$ and contain two sites, fixing the embedded volume per site to
\[
V_s=\frac{8}{3\sqrt3}d^3.
\]
Matching the electric and magnetic area sums to one isotropic continuum action gives
\[
g_{\rm E}^2=\frac{8\sqrt3}{3r}(1+\rho)t_s,
\qquad
g_{\rm M}^2=\frac{3\sqrt3}{2}rt_s,
\]
where $t_s$ is the primitive spatial heat time. The condition $g_{\rm E}^2=g_{\rm M}^2$ fixes
\[
\frac{a_t}{d}=\frac43\sqrt{1+\rho},
\qquad
g_{\rm iso}^2=2\sqrt{3(1+\rho)}\,t_s.
\]
The equal-transfer completion $\rho=1$ treats spatial route transfer and present/history transfer as uses of the same primitive instrument. It gives
\[
\boxed{\frac{a_t}{d}=\frac{4\sqrt2}{3}},
\qquad
\boxed{\frac{g_{\rm iso}^2}{g_{\rm HK,step}^2}=2\sqrt6}.
\]
The corresponding projective pure-gauge measure is
\[
Z_{\rm gauge}
=\int[dU]\,
\prod_{p_{\rm E}}K_{4t_s}(U_{p_{\rm E}})
\prod_{p_{\rm M}}K_{6t_s}(U_{p_{\rm M}}),
\]
and its graph heat-kernel coordinate is
\[
\boxed{
g_{\rm graph,HK}^2
=2\sqrt6\,g_{\rm HK,step}^2
=\frac{4\sqrt6}{3}\ln\frac{14}{13}
=0.24203562353.
}
\]
Here $U_\ell\in PU(3)$ and $K_t$ is the heat kernel for the projective group; locally it has the same $\mathfrak{su}(3)$ Casimir normalization used above.
This is a graph heat-kernel coordinate, not an $\overline{\rm MS}$ coupling. Conversion to a Wilson or continuum scheme requires the corresponding class integral and finite matching; neither is supplied by the blocking identity. The dimensionless result also does not identify the embedded bond length $d$ with the electron-calibrated $L_*$; that geometric dictionary belongs to the full GFT embedding.

\paragraph{Color and generation factorize.}
The shell and open-route spaces are independent in the minimal colored-defect extension,
\[
\mathcal H_{\rm defect}
=\mathcal H_{\rm shell}\otimes\mathbb C^3_{\rm route},
\qquad
K_{\rm total}^2=K_{\rm shell}^2\otimes I_3,
\qquad
T_{\rm color}^a=I_{\rm shell}\otimes T^a.
\]
It follows that $[K_{\rm total}^2,T_{\rm color}^a]=0$. The shell cutoff is also alphabet-independent: every one- and two-label exclusion leaves a nondegenerate closure branch, while every three-label exclusion makes $K^2$ constant. A colored realization would therefore carry three route components at each of the same three shell depths. The factor of three is gauge multiplicity and is not an additive $\ln3$ contribution to an isolated particle mass; physical states require their color-flux dressing and singlet constraints.

\paragraph{Status of the color result.}
The persistent-open-route identification is the conditional physical premise. Given it, reference-preserving reversible transport fixes the $PU(3)$ channel group, $\mathfrak{su}(3)$ connection algebra, and eight adjoint directions; phase-invariant channel-return fidelity fixes the adjoint Wilson form. The identity primitive kinetic pairing, maximum-throughput rule, and route-frame isotropy fix $t_8=13/14$. The embedded-diamond, equal-transfer causal completion fixes the dimensionless blocking factor $2\sqrt6$ and places the microscopic projective pure-gauge branch in a weak-coupling regime.

Physical QCD requires further structure that this derivation does not contain: a global $SU(3)$ endpoint lift and central $U(1)$ identification, triality-carrying chiral quark defect classes, two weak species per generation, hypercharge assignments, a lattice fermion action, finite graph-to-$\overline{\rm MS}$ matching, and nonperturbative scale setting. The projective pure-gauge action does not prove confinement or predict a hadron mass. The result is a conditionally closed microscopic local pure-gauge sector, not a derivation of the Standard Model matter sector.

\paragraph{Hadronic matter and the mass--entropy bridge.}
The mass--entropy bridge $m=\kappa_m(\ell)\,\Delta S$ applies to fermionic defects regardless of their color label. For charged leptons, the defect is color-singlet and the marked shell action of Appendix~I.1 gives both measured ratios within one standard deviation. For hadrons, the dressed entropy is dominated by the QCD-internal contributions of confinement-scale gluonic flux, trace-anomaly structure, chiral vacuum reorganization, and quark binding. The mass--entropy bridge then organizes the full dressed bound-state entropy budget,
\[
S_{\mathrm{ent,H}}^{\mathrm{dressed}}=S_{\mathrm{defect}}+S_{\mathrm{bind}}+S_{\mathrm{conf}}+S_{\chi\mathrm{SB}},
\]
without requiring the framework to re-derive QCD. The structural claim is compatibility: the dressed entropy budget is the quantity QCD calculates, and the mass--entropy bridge converts it to inertial mass through the same running $\kappa_m(\ell)$ used for elementary sectors.

\paragraph{Scope summary.}
The gauge sectors have distinct grades. The baseline-redundancy template supplies an Abelian hosting mechanism; Appendix~S constructs its compact electromagnetic extension, exact Gauss constraint, quadratic photon sector, and conditional induced flux energy. The native fidelity gluing exhibits hysteresis on the $6^4$ lattice. Running fixes the required electroweak stiffnesses in the diamond transfer form (S.9--S.10). Native weak-isospin and hypercharge records, finite-capacity realization, and coupling matching remain open. Primitive spin data supply an $SU(2)$ representation, while its identification with the chiral weak group remains external. The persistent-open-route premise gives a conditional microscopic $PU(3)$ projective pure-gauge sector with a fixed minimal adjoint transfer channel and causal regulator. Its global $SU(3)$ lift, the Standard Model quotient, chiral representations, matter multiplets, and renormalized couplings have not been derived. The gravitational and charged-lepton results do not depend on this color extension.

The sectors in this appendix are structurally linked to the same entanglement logic but are not part of the controlled ordinary Einstein/capacity branch.

\section*{Appendix J: Lattice Microstructure Calculations}

Appendix~J records the technical content of Section~24: the host ensemble and the checks that its geometry remains connected, simplicial, foliated, thermalized, and volume-matched; the coupled measure and its declared modeling forks; the control methodology; the exact cell-order calculations; and the measured compatibility, defect, and capacity-transport results. It then states what a continuum critical-surface calculation would have to establish. Where Appendix~K recomputes the coefficient chain from its definitions, this appendix couples that chain to a dynamical substrate and reports what the finite-lattice system does. The report, analyzer outputs, run configurations, verification scripts, and raw tables are maintained with the manuscript materials at \texttt{jaigp.org}; the report is generated from the archived CSV and log files rather than transcribed by hand.

\subsection*{J.1 Host ensemble and engine discipline}

The host is a standard causal-dynamical-triangulations ensemble \citep{AmbjornJurkiewiczLoll2004,AmbjornGJL2012}: triangulations of $S^1\times S^3$ sampled by local Pachner-class moves under the Regge action, with the $(4,1)$-simplex number pinned by a quadratic penalty. Early demonstration calculations used eighty spatial slices near $(k_0,\Delta)=(2.2,0.6)$. The largest-volume calculations use the directly mapped point $(2.0,0.4)$ and auto-tune $k_4$ separately for the $40{,}000$ and $100{,}000$ volumes and for the forty- and eighty-slice realizations; the final values are $0.7560$, $0.7608$, and $0.7700$, respectively. The action is tracked incrementally and checked against full recomputation throughout every run, and every quoted final triangulation is verified to be a connected simplicial manifold with the required foliation.

The coupling scan maps finite-volume behavior, not Einstein-phase membership. A $(k_0,\Delta)$ grid with $k_4$ auto-tuned to pseudo-criticality at every point (the volume-pin residual serving as the feedback error, as in standard CDT practice where $k_4$ is retuned per coupling point) finds an extended region at $k_0=1.0$--$2.0$, with slice-volume coherence peaking at $(k_0,\Delta)=(2.0,0.4)$ and a dead row by $k_0=3.5$: a visible boundary. A fixed $k_4=0.9$ obstructs the scan, because the volume pin then dominates the sampling at every untuned cell. Before the largest-volume data were examined, the expected finite-volume trend at the selected couplings was stated explicitly: both the Hausdorff estimator and the concentration of volume in an extended blob should rise with volume. They do. The $N_{41}\simeq42{,}112$ ensemble gives $d_H=3.31$, blob score $1.46$, and $23$ active slices; the $N_{41}\simeq102{,}002$ ensemble gives $d_H=3.56$, blob score $1.57$, and $24$ active slices. This establishes an extended finite-volume compatibility region. It does not distinguish the de~Sitter phase from a nearby bifurcation phase, determine the order of the visible boundary, establish a continuum trajectory, or demonstrate a massless TT sector.

Validation of the cell identification in Section~24.1 requires a foliation filter on the move set. Unfiltered moves preserve the four-dimensional manifold conditions but can fold a spatial slice, leaving a few percent of slice triangles in violation of the required closed-$3$-manifold condition. The engine therefore rejects every proposal that would break the foliation. Its local test is complete because the link of an affected triangle necessarily passes through a newly created pentachoron, and it preserves detailed balance in the same way as the other locality filters. Every ensemble reported here uses this engine and passes the global audit: each slice triangle belongs to exactly two spatial tetrahedra, and every pentachoron is correctly threaded between its slices.

No result below infers more from that comparison than finite-volume host suitability. In particular, the point $(2.0,0.4)$ is used only as a compatibility test for switching on the capacity weighting; it is not used to prove \((H_{\rm phase})\).  Distinguishing the de~Sitter phase from a nearby extended bifurcation phase would require the standard alternating-slice curvature and singular-vertex order parameters, but that classification is not an antecedent of the analytic local-stability result in H.10. Every substrate claim compares ensembles at the same regulator: coupled against bare, failure against closed control, or one source strength against another. Establishing a continuous geometric transition and taking the decorated common-parent continuum limit are the separate tasks of J.9.

\subsection*{J.2 The coupled measure and its fidelity anchors}

The coupled measure is the joint Gibbs distribution of Section~24.2, with the per-face label sum normalized so that each label carries base measure $1/7$. Under that base measure two properties hold exactly: at $\beta=0$ the label entropy cancels and the geometry marginal is identically the bare host, so the uncoupled ensemble reproduces the host by construction; and uniform label assignment at geometry moves is exactly detailed-balanced, so no proposal correction is needed for label births. The centering constant $\mu$ is the per-cell label free energy,
\[
\mu(\beta)=-\frac{1}{\beta N_{\rm cells}}\,\ln\,\mathbb{E}_{m\sim{\rm unif}}\!\left[e^{-\beta\sum_c E_c}\right]
=\frac{1}{\beta}\int_0^\beta\!\langle \bar E\rangle_s\,ds,
\]
computed by thermodynamic integration on the thermalized base configuration and logged with each ensemble, so the extensive part of the closure term is neutralized rather than allowed to masquerade as a host coupling. Total volume is nevertheless checked directly and matched among comparisons rather than assumed from the centering. Equivalently, the normalized factor is $\widehat Z_{\rm cap}=e^{\beta\mu N_{\rm cells}}Z_{\rm cap}$. It is exactly invariant under a common energy-origin shift $E_c\mapsto E_c+C$, $\mu\mapsto\mu+C$, which is the finite-regulator identity used in Section~25. Because $N_{\rm cells}=N_{41}/2$, the omitted centering is exactly an $N_{41}$ source. In the standard $(\kappa_4,\Delta)$ action basis this means $\delta\kappa_4=-\beta\mu/2$, $\delta\Delta=+\beta\mu/2$, and $\delta\kappa_0=-6\,\delta\Delta$ (the $N_0$ coefficient is unchanged), not a pure $\kappa_4$ displacement; Section~25 separates its physical volume component from the regulator-asymmetry component.

Every run opens by rebuilding the single-cell tables and checking the Part~II anchors ($\Omega_{\rm tet}=1680$, $\eta_*=0.0298668$, $\langle K^2\rangle_{\eta_*}=50.223$, $g_{\mathrm{share,eff}}=7.4198$), and refuses to sample if any fails.

Three modeling forks are declared. Injectivity is imposed as a penalty $\lambda$ per colliding pair ($\lambda=3$ in the quoted runs) because a hard seven-color constraint at the maximal conflict degree does not guarantee an ergodic single-face heat bath; the hard constraint is the $\lambda\to\infty$ limit, the residual collision fraction is reported in every table so the softness stays visible, and the physical $\beta=1$ weighting is read at the ordered end. The closure sum is evaluated in the unoriented reading, in which the two cells sharing a face see the same label; the oriented alternative, in which they see opposite signs, requires oriented-slice bookkeeping through all moves and is recorded as an open fork. Parity contributes exactly $\ln 2$ per cell, extensive at pinned volume, and drops out of the coupled dynamics.

\subsection*{J.3 Controls and the reading discipline}

Four validity conditions precede any physical reading. Every quoted triangulation must remain a connected simplicial manifold with the required foliation; its action and volume time series must be stationary over the measured window; compared ensembles must have matched total-simplex number, because volume differences can masquerade as dimension differences; and the Markov chain must continue moving, because a frozen chain does not sample the measure. The interpretation of each possible outcome, including which outcomes would license no claim at all, was written before the corresponding data were examined.

The main control shuffles the closure energies. The cell has $210$ sorted-label orbits; the control reassigns their energy entries among those orbits, preserving the orbit-value pool and permutation symmetry, and is evolved at the same $\beta=1$ coupling. The orbits have unequal multiplicities, so the reassignment changes the ordered-state energy histogram and one-cell free energy. Differences from this control combine the changed closure assignment with those thermodynamic changes. A control that preserves the physical energy histogram would restrict its permutations to orbits of equal multiplicity.

An EPRL calculation exposed the need for this control. A vertex amplitude \citep{EPRL2008} at frozen boundary spin $j=3$, computed with the \texttt{sl2cfoam-next} library \citep{Gozzini2021}, was coupled to the same host. An analysis specified before reading the data found five misspecifications: off-equilibrium centering shifted the bare cosmological coupling; the label heat bath used the wrong exponent; label births lay outside the acceptance ratio; the result depended at order unity on an arbitrary slot convention; and sampling required a positivization of the sign-alternating amplitude. Every effect vanished at $\beta=0$ and grew with $\beta$, so an uncoupled comparison could not reveal them. An entry-shuffled vertex tensor reproduced the apparent geometric steer at matched coupling and identified it as sampler machinery. The sampler was rebuilt as an exact joint Gibbs measure, and the closure calculations include the same shuffled-energy control. This episode is reported as methodology. The frozen-spin truncation and positivization already limit its physics scope, and the EPRL intertwiner space is distinct from the seven face labels of the present ensemble (Section~29.8).

\subsection*{J.4 Exact results at cell order}

Six calculations, independent of any Monte Carlo run, fix what the substrate can and cannot do for the vacuum before the coupled measurements are read.

\paragraph{Induced curvature coupling.} The geometry marginal of the coupled measure weights a slice by the free energy of its labels, and the leading geometric dependence enters through the coordination $q$ of slice edges, the discrete curvature variable. A transfer-matrix evaluation of the label free energy on the ring of cells around an edge gives, at the physical $\beta=1$ weighting, an induced coupling $c_0=+0.019$ per slice edge, against the bare $k_0\simeq2.2$: about one percent. The label correlation length is about half a lattice step, and the beyond-quadratic dependence on $q$ is negligible at physical coordinations. The closure weighting alone therefore cannot supply the host's geometric stiffness, and its predicted effect on the coupled sweep is a steer too small to see at single-seed resolution, a prediction the calculation of J.5 confirms.

\paragraph{The refresh at cell order.} Modeling the selected memoryless kernel as whole-boundary redraws arriving at unit rate per cell, the single-cell stationary state is exact: the probability that $k$ faces retain their last self-drawn values is $k!(4-k)!/120$, vacuum admissibility is $0.536$, and $\langle K^2\rangle=48.74$ against the ensemble value $50.22$. The state is geometry-blind to exponential accuracy and induces no vacuum curvature action. It does couple geometry updates to the local density of closure failure when their rates depend on that density. This calculation tests the effective refresh model, not a derivation of its microscopic operator.

\paragraph{The history-weighting ceiling.} Two computations bound the strongest vacuum role a Many-Pasts tilt over closure-class observables could play. Amplifying the closure link interaction to $128$ times the physical $\beta=1$ weighting saturates the link-agreement ratio near $0.55$ with correlation length $1.7$ links and a curvature-sector signal near $0.1$: a structural ceiling, set by the ground-state degeneracy of the closure-plus-injectivity interaction, twenty times below the bare coupling. Independently, the exact Doob-transform solution of a closure-maintaining history tilt (paths weighted by time-integrated closure quality) on a ring of six seven-state faces drives the nearest-neighbor label correlation from $-0.12$ through zero to only $+0.13$ across the full tilt range, with longer-range correlations near $0.03$ throughout: no long-range order at any tilt strength. Together with the two paragraphs above this closes the label sector: static weight, refresh, and closure-class history weighting each fail, quantitatively, to supply vacuum geometric stiffness at cell order.

\paragraph{The free conserved field, integrated out.} A conserved Gaussian field back-coupled to the geometry induces, on integration, the geometric action $\tfrac{1}{2}\ln\det'L(g)$ per channel and slice: the spanning-tree entropy of the slice graph, by the matrix-tree theorem. Computed on degree-matched graph classes, the entropy per cell is $0.607$--$0.613$ on diamond-cubic lattices (the smooth extended proxy), $0.609$ on random $4$-regular graphs (the crumpled proxy), and $0.56$--$0.57$ on real engine slices: nearly universal at fixed coordination, with the genuinely nonlocal soft-mode part diluting as $\ln N/N$ per cell. The per-cell differential between extended and crumpled geometry is a few times $10^{-4}$; with all seven channels the phase-tipping force is of order $10^{-3}$ against $k_0\simeq2.2$.

\paragraph{The budget-bounded field.} The finite budget of Postulate~I makes the field non-Gaussian, with saturation curvature $m^2=2/g_{\mathrm{share,eff}}\simeq0.27$, and the massive determinant differential between the same proxies shrinks from $1.2\times10^{-3}$ per cell at zero mass to $4\times10^{-4}$ at the budget mass and toward zero beyond: the theory's own non-Gaussianity suppresses the modes where the sensitivity lived. What survives, measured on eight engine slices, is a local effective-action shift of scale $\simeq0.08$ per cell in the seven-channel count.  This number is not itself a displacement in \(\Delta\), or a norm in the \((k_0,\Delta)\) plane; obtaining that displacement requires projection onto the CDT conjugate operators.  The result shows perturbative dressing of the host. The two containers no branch covers are named in Section~24.3: long-range field kernels and non-equilibrium geometry--field co-evolution, both unspecified by the theory.

\paragraph{The fracture construction.} Under unit label shifts preserving injectivity, the $1680$-state move graph has $48$ connected components of $35$ states, labeled by slot ordering and parity. This proves that the tested move class cannot equilibrate the full admissibility ensemble. More general microscopic dynamics are not excluded.

\subsection*{J.5 Compatibility at matched volume}

The closure comparison starts from common tuned geometries at $N_{41}\simeq100{,}000$. Three eighty-slice ensembles use respectively the uncoupled value $\beta=0$, the physical weighting $\beta=1$, and $\beta=1$ after the closure energies have been shuffled among label orbits. The physical weighting is also repeated with forty slices. In every case the final triangulation is connected, simplicial, correctly foliated, and stationary over the measurement window. Across the three eighty-slice ensembles the total-simplex spread is $0.6\%$.

\begin{center}
{\small
\begin{tabular}{lrrrrr}
\hline
ensemble & $N_4$ & $N_{41}$ & collision fraction & $d_H$ & centering $\mu$ \\
\hline
$\beta=0$ (80-slice control) & $323{,}514$ & $100{,}216$ & $0.654$ & $3.55$ & $3.9913$ \\
$\beta=1$ (80-slice physical) & $323{,}232$ & $99{,}998$ & $0.077$ & $3.59$ & $2.4990$ \\
$\beta=1$ (80-slice shuffled) & $325{,}346$ & $99{,}944$ & $0.736$ & $3.59$ & $3.1420$ \\
$\beta=1$ (40-slice repeat) & $316{,}238$ & $100{,}020$ & $0.078$ & $3.74$ & $2.4983$ \\
\hline
\end{tabular}
}
\end{center}

The uncoupled ensemble remains on the exact uniform-measure benchmark $1-840/2401=0.650$ within sampling error. The physical weighting suppresses collisions by nearly an order of magnitude, and the shuffled control remains high despite carrying the same energy-value pool. At the same time the matched eighty-slice $d_H$ values occupy the narrow interval $3.55$--$3.59$, their extended volume profiles persist, and the forty-slice repetition reproduces the ordered collision fraction. Thus the microscopic closure structure is strongly active while the host geometry is stable at this resolution. A $20{,}000$-volume run reaches the same conclusion with a larger finite-volume caveat.

\subsection*{J.6 The Regge bridge: a predicted volume-sector dressing}

The bridge is an ordinary statistical-mechanical statement made quantitative. Integrating the labels out at fixed geometry defines an effective action for the host,
\[
e^{-S_{\rm eff}[g]}=e^{-S_{\rm host}[g]}\,Z_{\rm label}[g],
\qquad
S_{\rm eff}[g]=S_{\rm host}[g]-\log Z_{\rm label}[g],
\]
with $S_{\rm host}$ the Regge action plus the volume pin. The label free energy $-\log Z_{\rm label}[g]$ is a function of the geometry alone. Its extensive part is $\beta\mu(\beta)\,N_{\rm cells}$, with $\mu$ the per-cell label free energy of J.2; its geometry-dependent remainder is the induced curvature-sector action priced by the ring expansion of J.4. The centered ensemble subtracts the extensive part by construction; left in, it is the leading coefficient of the effective action, and the host's response to it can be predicted before it is measured.

The centering constant $\mu$ of J.2 therefore gives a direct prediction. In the centered ensemble, the subtraction of $\mu$ removes the extensive label free energy from the geometry marginal. If that subtraction is omitted, the geometry sees a residual positive linear contribution $+\beta\mu N_{\rm cells}$ to the effective action. Here $N_{41}$ denotes the total number of $(4,1)$ and $(1,4)$ pentachora, so each spatial cell, shared by one pentachoron of each orientation, is counted once by the pair and $N_{\rm cells}=N_{41}/2$. The positive linear tilt competes with the quadratic volume pin $\varepsilon(N_{41}-\bar N)^2$, and the shift follows in two lines. The volume-dependent part of the effective action is
\[
S_{\rm vol}(N_{41})
=\varepsilon\,(N_{41}-\bar N)^2+\beta\mu\,N_{\rm cells}
=\varepsilon\,(N_{41}-\bar N)^2+\frac{\beta\mu}{2}\,N_{41},
\]
and setting $dS_{\rm vol}/dN_{41}=2\varepsilon(N_{41}-\bar N)+\beta\mu/2=0$ gives the equilibrium displacement
\[
\Delta N_{41}=N_{41}^{\rm eq}-\bar N=-\frac{\beta\,\mu}{4\,\varepsilon}.
\]
Here $\mu$ is computed by annealed thermodynamic integration of the label ensemble on the base geometry, before any coupled dynamics run; the zero-lattice reference (the exact single-cell free energy from direct enumeration of the $7^4$ label states, computable from $\eta_*$ and the injectivity penalty alone) is $2.477$ at this point, within half a percent of the lattice value, the difference being the shared-face correlations. At the demonstration point ($\beta=1$, $\mu=2.488$, $\varepsilon=0.01$) the prediction is $\Delta N_{41}=-2.488/0.04=-62.2$. Three ensembles from a common equilibrated base at $N_{41}=2000$ measure it:

\begin{center}
{\small
\begin{tabular}{lcc}
\hline
ensemble & mean $N_{41}$ & shortfall \\
\hline
bare ($\beta=0$) & $1993.2$ & $-6.8$ \\
centered ($\beta=1$) & $1969.8$ & $-30.2$ \\
uncentered ($\beta=1$) & $1907.0$ & $-93.0$ \\
\hline
\end{tabular}
}
\end{center}

The uncentered-minus-centered difference is $-62.8$ against the predicted $-62.2$. The subtraction isolates the $\mu$-term, since the two $\beta=1$ ensembles share every other piece of machinery; the centered ensemble's own $-30$ is the second-order non-neutrality of the thermodynamic-integration centering. It cancels in the difference and is kept out of other comparisons by matching their volumes directly. A number computed from the ultraviolet specification thus appears, at the predicted magnitude, as an extensive $N_{41}$ contribution to the dynamical host action. This is a dressing of the host, not a derivation of its continuum limit. In the standard CDT basis it is a correlated $(\kappa_4,\Delta)$ displacement; only projection with the two-simplex Regge volume operator identifies its physical volume component. The measurement is at demonstration volume and uses one seed. The corresponding curvature-sector measurement remains open.

The demonstration point is one member of a predicted family whose functional form was recorded before sampling. At fixed geometry-side settings the displacement obeys $\Delta N_{41}(\beta,\varepsilon)=-\beta\mu(\beta)/(4\varepsilon)$, and each factor is separately testable. In $\beta$ the displacement curve is convex: $\beta\mu(\beta)=\int_0^\beta\langle\bar E\rangle_s\,ds$ is concave as the Gibbs mean falls, and $\Delta N_{41}=-\beta\mu/(4\varepsilon)$ reverses that curvature, with local slope $d(\Delta N_{41})/d\beta=-\langle\bar E\rangle_\beta/(4\varepsilon)$, an Ehrenfest-type relation in which the same calculation supplies both sides, since $\langle\bar E\rangle_\beta$ is measurable in the ensemble whose displacement it predicts. In $\varepsilon$ the law is exactly inverse at fixed $\beta$; departures from $1/\varepsilon$ measure the curvature of the background free energy that the matched-pair subtraction cancels at first order, so the $\varepsilon$ family doubles as a control on the derivation's one assumption. The shuffled table carries its own computable free energy: at $\beta=1$ the single-cell values are $\mu_{\rm shuf}=3.17$ against $\mu=2.48$, so the shuffled ensemble is predicted to displace by $-\beta\mu_{\rm shuf}(\beta)/(4\varepsilon)$, a nonzero displacement $28\%$ away from the physical value. The measured number must track the table from which it was computed, so the shuffled calculation supplies a second quantitative prediction for its altered local spectrum.

The family has been measured at demonstration volume, with every prediction written to disk before sampling began. Eight of nine planned coupling--volume-penalty combinations produced readable data; the ninth failed at launch and contributes no datum. Across a fifteen-fold range of predicted displacements, from $-10.5$ to $-157.6$, the measured-to-predicted ratios are $0.96$--$1.07$ with one point at $0.89$. Along $\beta$ the measured displacements $-43$, $-78$, $-121$ track the predicted $-43$, $-76$, $-125$, including the departure from a straight line that the free-energy integral fixes. Halving $\varepsilon$ doubles the displacement to within four percent. The shuffled pair lands on its own line, displaced thirty percent from the physical-table results and at ratio $1.00$ to the prediction computed from its shuffled table. A repetition of this separate bridge calculation at the largest volume has not yet been performed.

\subsection*{J.7 The pinned-defect experiment}

Two independent matched comparisons resume from thermalized $N_{41}\simeq100{,}000$ ensembles with the physical $\beta=1$ weighting and run $3000$ further sweeps. For each random seed, one ensemble pins one hundred cells at maximal closure failure (all four faces $m=0$: six collisions, $K^2=48$), while the control pins one hundred cells at the best-closed injective configuration ($m=\{0,1,2,3\}$, $K^2=40.67$). Pinned faces are excluded from the label heat bath, and the carrier cells are protected from removal by the geometry moves identically. Every final triangulation remains connected and correctly foliated, and every ensemble retains all one hundred pins. The final analyzer pools the two seeds. When an observable--shell cell is present in both seeds its uncertainty is set by across-seed scatter; otherwise it uses the Sokal-windowed integrated-autocorrelation standard error. The largest measured integrated autocorrelation time is $349.5$ measurement intervals.

\begin{center}
{\small
\begin{tabular}{llcccc}
\hline
condition & shell & mean $E$ & collision fraction & mean coordination & cells per shell \\
\hline
failure   & 1 & $1.6802\pm0.0026$ & $0.0698\pm0.0009$ & $5.8452\pm0.0024$ & $4$ \\
          & 2 & $1.7299\pm0.0006$ & $0.0783\pm0.0002$ & $6.0869\pm0.0281$ & $7.7289\pm0.2255$ \\
          & 3 & $1.7269\pm0.0002$ & $0.0786\pm0.0001$ & $6.1986\pm0.0097$ & $12.7960\pm0.4518$ \\
\hline
closed control & 1 & $1.7260\pm0.0008$ & $0.0783\pm0.0002$ & $5.8429\pm0.0146$ & $4$ \\
          & 2 & $1.7288\pm0.0003$ & $0.0788\pm0.0001$ & $6.0816\pm0.0204$ & $7.8633\pm0.2143$ \\
          & 3 & $1.7277\pm0.0002$ & $0.0787\pm0.0001$ & $6.2384\pm0.0022$ & $12.5920\pm0.4737$ \\
\hline
\end{tabular}
}
\end{center}

The failure-minus-closed-control differences make the reading unambiguous. At shell one, the mean closure energy changes by $-0.0458$ ($16.8\sigma$) and the collision fraction by $-0.0085$ ($9.2\sigma$). At shell two their separations are $1.7\sigma$ and $2.1\sigma$, respectively. At shell three the mean coordination changes by $-0.0398$ ($4.0\sigma$), while its first- and second-shell separations are only $0.2\sigma$. Shell-cell counts separate by at most $0.4\sigma$ anywhere in the measured range. Thus persistent failure produces a replicated, short-range capacity disturbance and a statistically resolved local coordination response. A shell-volume signal seen in a smaller run does not survive the larger-volume autocorrelation and seed treatment.

\subsection*{J.8 The conservation law and the transport instrument}

The dichotomy of Section~24.7 in operator form. If free capacity is a per-cell budget re-equilibrated by maximum entropy (link shares $c_l\ge0$ with $\sum_{l\ni x}c_l=C-m_x$), stationarity gives one multiplier per cell, and the linearized constraint reads $(2z\,\mathbb{1}-L)\,\delta\mu=-m/\chi$ with $z=4$ and $L$ the slice Laplacian. The operator has no small-momentum pole; on a measured $169$-cell slice its Green function falls eight orders of magnitude within thirteen steps.

The transport instrument tests a different object. Its pinned commitment level $m_x$ is a fixed source label, while $Q_x$ is the maintenance throughput the source draws. On a closed slice the simulated free field obeys
\[
f_x(t+dt)-f_x(t)
=-\sum_{\text{faces }xy}J_{xy}-Q_x+\overline Q,
\qquad J_{xy}=-J_{yx},
\qquad \sum_x f_x={\rm const},
\]
where $\overline Q$ is the uniform return required by the compact zero mode. This is steady service, not $f\to m$ stock conversion. Any local flux law linearizes to $J\propto-D\nabla f$ and the steady state solves $L\,\delta f=(Q/D)(\delta_{\rm source}-{\rm uniform\ return})$: the massless graph-Poisson equation, whose Green function on a three-dimensional slice falls as $1/r$. Section~24.7 derives the corresponding incidence-matrix Ward identity and identifies the continuum junction coefficient still to be measured.

The instrument implements the conservation reading with the sink left to emerge. A field on the slice cells is initialized at the vacuum anchor $g_{\mathrm{share,eff}}=7.4198$ and transported by antisymmetric face fluxes; the code asserts conservation at machine precision every sweep and aborts on any violation. The pin levels are fixed source labels, so the conserved quantity during this test is the free-field sum $\sum_x f_x$. Absorption is proportional to a cell's \emph{persistent excess} closure failure above the measured vacuum baseline. A pinned defect therefore draws maintenance only through the failure generated by the ordinary label dynamics: its flux is measured, not set. Pins are placed at the five source values $m=0,1,2,3,6$. The zero value measures commitment-independent boundary dressing; the positive values test source additivity and the deficit per unit emergent charge. At large distance the calculation tests whether the response follows a graph-Coulomb power law and whether a nonzero screening mass is detected. This implementation freezes geometry while labels and capacity remain dynamical, isolating the transport law; conservative transport through geometry moves is the specified extension.

On small validation volumes the measured deficit increases with commitment, recovers monotonically with distance, and carries an emergent flux that rises with source value while the total field remains exactly conserved. The highest-commitment pin draws less flux than the linear extrapolation because it depletes its own neighborhood, a self-screening the model produces without being told.

Two failed definitions are retained as diagnostics. Absorption on total failure makes the vacuum itself a sink and produces the predicted four-to-seven-step screening. Absorption on instantaneous excess rectifies the vacuum's mostly-zero collision process and also creates a screening mass. Both failures follow quantitatively from their own update laws. The adopted definition therefore uses a $200$-sweep persistence average, equilibrates the field before measurement, and subtracts the same-slice far-field reference to remove the compact zero mode. These analysis choices were fixed before the largest-volume data were generated.

The transport calculation uses an ordered forty-slice realization at $N_{41}\simeq100{,}000$. It places thirty-two pins at each of the five source values, measures through shell eight, and runs $4000$ sweeps. The mean field remains $7.4198$ throughout. Per-pin time series are corrected with Sokal integrated autocorrelation times, and the level errors use across-pin scatter; the profile and junction errors additionally use the larger of propagated and across-level scatter. The finite-lattice results are
\[
Q-Q(0)=(0.1440\pm0.0025)m,
\qquad Q(0)=0.0075,
\]
\[
\left\langle\frac{\Delta f_1-\Delta f_1(0)}{Q-Q(0)}\right\rangle_{m>0}
=4.9965\pm0.5244,
\]
and the normalized dressing-subtracted shell profile is
\[
\begin{array}{c|cccccccc}
d&1&2&3&4&5&6&7&8\\ \hline
\Delta f(d)/\Delta f(1)&
1.000&0.822&0.675&0.532&0.411&0.303&0.212&0.144.
\end{array}
\]
The power-law fit gives $-0.62\pm0.11$. The screening derivative is $+0.1062\pm0.0551$, statistically consistent with zero; the corresponding finite-range lower bound is $\xi>4.6$ steps. The source fit is not perfectly linear: its maximum relative residual is $0.242$, dominated by visible saturation at $m=6$. The tail is also shallower than the ideal $d^{-1}$. A $1/r$ Green function smeared over a source of width $\sigma=1.95$ shells, with the far-field offset of the zero-mode subtraction, reproduces the normalized profile at rms residual $0.018$ (Section~24.7), so the finite source extent accounts for most of the measured exponent at this volume; the continuum scaling study must test convergence. Subject to those explicit diagnostics, the finite-lattice calculation demonstrates conservation, approximate source additivity, one junction ratio across the tested source values, power-law propagation, and no detected screening.

The continuum test uses the slice adjacency and conductance matrices directly. Across volumes it measures the low eigenvalues $\lambda_j(L_f)$, the heat trace $\operatorname{Tr}e^{-tL_f}$ and return probability, the spectral dimension, the properly normalized Green amplitude, and the renormalized junction coefficient against its native-unit target $8\pi$. Together these test the joint metric--measure--spectral convergence of Section~24.7 and decide whether the present tail is a finite-size effect or an anomalous spectral dimension.

\subsection*{J.9 Capacity-decorated CDT: phase membership and the continuum test}

These measurements establish the finite-lattice source-to-field behavior listed above. A microscopic CDT completion must still promote the capacity data to part of the joint transfer state.  Because $L_*$ is a physical microscopic length, existence of the effective theory does not require a second-order point at which $a\to0$.  Its necessary geometric test is instead membership in an extended semiclassical phase whose transfer matrix contains a massless transverse--traceless sector and whose large-volume observables approach four-dimensional metric behavior.  A continuous critical surface is a stronger universality and cutoff-artifact-removal test, not the existence condition for the finite-$L_*$ theory.  For each triangulation $T$, the minimally sufficient configuration bundle is
\[
\mathfrak X_{\rm cap}(T)
=\left\{
b_c,\ f_c,\ c_{A,c},\ \theta_A,\ H_{A,c;\ell}
\right\}.
\]
Here $b_c$ is the existing closure state, $f_c$ is transported free renewal bandwidth, and $c_{A,c}$ is the mass-equivalent committed allocation carrying the physical source label $A$. The variables $\theta_A$ and $H_{A,c;\ell}$ are respectively the source/worldline renewal phase and the internal history output at renewal depth $\ell$. The index $\ell$ labels circuit records, not a fifth spacetime coordinate. The source-supported observables are
\[
\begin{aligned}
M_{c,A}^{\rm lat}&=\sum_c c_{A,c},\\
\sigma_A&=\frac{M_{c,A}^{\rm lat}}{M_{c,A}^{\max}},\\
\mathcal C&=\text{single-stream integrability of the combined committed flow}.
\end{aligned}
\]
There is no local per-channel ``saturated'' state and no microscopic coherence bit. After vacuum normalization the coarse availability observable is obtained from $f_c$, while $\sigma_A$ is computed against the independently fixed per-source ceiling. The transfer kernel must route formation and release through the existing marked defect vertex, so $c_{A,c}$ is not an externally prescribed matter source. For a stable mark those stock-changing rates vanish; its ordinary closure dynamics instead determine the maintenance draw from $f_c$. Terminal exhaustion is omitted from this bulk transfer matrix because it belongs to the separate $q_{\rm geo}=0$ boundary theory.

The ultraviolet configuration space is the disjoint union over causal triangulations,
\[
\mathfrak X_{\rm UV}
=\bigsqcup_{T\in\mathfrak T_{\rm CDT}}
\mathfrak X_{\rm cap}(T),
\]
with joint partition function
\[
Z(\mathbf g)
=\sum_T\frac{1}{C_T}
e^{-S_{\rm CDT}[T;\kappa_0,\Delta,\kappa_4]}
Z_{\rm cap}[T;\mathbf g_{\rm cap}].
\]
The microscopic theory fixes $\mathbf g_{\rm cap}$; the physical capacity slice is therefore
\[
\mathcal P_{\rm cap}
=\{(\kappa_0,\Delta,\kappa_4;\mathbf g_{\rm cap}^{*})\}.
\]
Integrating out the capacity variables gives $S_{\rm eff}=S_{\rm CDT}-\ln Z_{\rm cap}$, making their displacement of the host critical manifold a measurable output.

The measured capacity backreaction is perturbative in the host action, but
the quoted \(0.08\) per cell is not yet a coordinate displacement in the CDT
coupling plane.  To obtain one, the geometry-dependent term must be projected
onto the operators conjugate to the CDT couplings:
\[
-\log Z_{\rm cap}[T]
\simeq aN_0+bN_{41}+cN_{32}+\cdots ,
\]
and then matched to
\[
S_{\rm CDT}
=-(\kappa_0+6\Delta)N_0
+\kappa_4(N_{41}+N_{32})
+\Delta(2N_{41}+N_{32}).
\]
Only that projection defines
\((\delta\kappa_0,\delta\Delta,\delta\kappa_4)\) or a measured boundary
shift.  Thus \(0.08\not\equiv\delta\Delta\), and it cannot be compared
directly with a plotted distance to a phase boundary.  H.10 needs no such
comparison: it proves local stability of an already existing Einstein phase
under a finite, pole-free, gauge-compatible correction smaller than the host
TT stiffness.  Phase membership itself must still come from the geometric
host.

Increasing the number of simplices at a generic coupling point is not by itself evidence of either phase membership or a continuum limit.  The existing CDT runs establish only that strong capacity ordering is compatible with an extended finite-volume geometry at the sampled regulator; they are not used to prove \((H_{\rm phase})\).  The finite-$L_*$ route must instead show a stable extended phase, four-dimensional large-volume metric behavior, and a massless TT transfer sector.  If one additionally seeks cutoff removal, $\kappa_4$ must approach its critical value so the four-volume can diverge and a remaining relevant direction $t$ must approach a continuous geometric transition at which
\[
\frac{\xi_{\rm geom}}{a}\sim |t|^{-\nu}\longrightarrow\infty.
\]
Four-dimensional CDT contains candidate continuous transition lines that can support this construction \citep{AmbjornSecondOrder2011,AmbjornCritical2019}. In the enlarged decorated theory define
\[
\mathfrak C_*
=\left\{\mathbf g:
\kappa_4=\kappa_4^{c},\quad
\xi_{\rm geom}/a=\infty
\right\}.
\]
That stronger cutoff-removal route succeeds only if
\[
\boxed{\mathcal P_{\rm cap}\cap\mathfrak C_*\ne\varnothing.}
\]
If a capacity coupling is irrelevant, it shifts the location of the critical surface without adding a tuning. If it is relevant, its already-fixed microscopic value must lie on the surface. If the geometry and capacity sectors form a new joint fixed point, its exponents become predictions of the decorated theory. Failure of the fixed capacity slice to intersect $\mathfrak C_*$ rejects the cutoff-removal embedding, but does not by itself reject a finite-$L_*$ theory already demonstrated to lie inside the required semiclassical phase.

The volume-operator normalization itself is already fixed kinematically. With Euclidean asymmetry $\alpha$,
\[
\mathcal V_4
=a^4\left[
\frac{\sqrt{8\alpha-3}}{96}N_{41}
+\frac{\sqrt{12\alpha-7}}{96}N_{32}
\right].
\]
After the mixed non-volume directions have been projected out, the two independently read source coefficients must satisfy
\[
\boxed{
\frac{t_{41}}{a^4v_{41}}
=\frac{t_{32}}{a^4v_{32}}
=\rho_{\rm vol,R}.
}
\]
Equivalently, the target-blind double ratio
\[
\boxed{
\mathcal R_V(a)
:=\frac{t_{41}v_{32}}{t_{32}v_{41}}
\longrightarrow1
}
\]
must converge before either coefficient is converted to $\Lambda_{\rm R}$. This equality uses two measurements and has no adjustable relative normalization. At fixed $\xi=N_{32}/N_{41}$ it reduces to $Z_V^{\rm geom}=(v_{41}+\xi v_{32})/(1+\xi)$. The scaling trajectory and its operator projection remain to be measured.

The required finite-size program is:
\begin{enumerate}
\item auto-tune $\kappa_4$ at each $(\kappa_0,\Delta)$, repeat the phase scan across increasing $N_4$, measure $\alpha(a)$ and $\xi(a)$, set $a$ from a target-blind geometric observable, and verify the two-source equality $t_{41}/(a^4v_{41})=t_{32}/(a^4v_{32})$ after critical-eigenoperator projection;
\item identify a continuous boundary using susceptibility peaks, Binder-cumulant crossings, histogram shape, transfer-matrix gaps, and pseudo-critical drift;
\item diagonalize the connected covariance matrix of the action-conjugate geometric and capacity operators,
\[
\mathsf C_{ij}=\langle\mathcal O_i\mathcal O_j\rangle_c,
\qquad
\mathcal O_i=-\frac{\partial\ln w}{\partial g_i},
\]
to identify the relevant mixed eigenoperator and the renormalized $q$ observable;
\item verify the augmented-event incidence identity configuration by configuration and show that any screening mass vanishes with increasing volume;
\item measure $D_q$, $\chi_q$, and $\alpha_{\rm maint}$ separately, including their lattice-spacing factors, before forming $\chi_q\alpha_{\rm maint}/D_q$;
\item apply the electron-derived $L_*$ and $G_*$ only after the dimensionless continuum extrapolation and test the junction condition $\chi_q\alpha_{\rm maint}/D_q=8\pi G_*/c^2$.
\end{enumerate}

The phase and continuum tests do different work.  Phase membership supplies the long-wavelength geometric branch at physical $L_*$; a critical surface would additionally remove the geometric cutoff and test universality.  The incidence Ward identity separately protects the hydrodynamic zero mode of the proposed bandwidth carrier. None supplies the others, and conservation alone does not establish the metric bridge, source purity, causal response, or junction normalization. At fixed lattice spacing, the larger host has the expected stronger extended-region indicators, and the transport calculation demonstrates conservation, approximate source additivity, a source-independent lattice junction ratio, power-law propagation, and no detected screening over the measured range. The massless TT transfer sector, asymptotic four-dimensional phase behavior, visible boundary's order, optional decorated critical trajectory, target-blind operator and lattice-spacing matching, and absolute junction to $8\pi G_*/c^2$ remain open. No continuum claim is assigned to the completed finite-volume results.

\subsection*{J.10 Predicted bare couplings and the native point}

The face-wise lock of Appendix~H.9a, paragraph~4a, defines a registered equilateral CDT branch with spatial edge $a=2.530L_*$ and matched $G=G_*=L_*^2$ ($\hbar=c=1$). Its bare couplings follow within that branch. The coarse-channel construction of Q.11a uses a different relation between physical cuts and regulator faces.  Writing the Euclideanized Regge action in the standard form
\[
S_E=-(\kappa_0+6\Delta)N_0+\kappa_4(N_{41}+N_{32})+\Delta(2N_{41}+N_{32})
\]
and using the exact four-dimensional action of \citep{AmbjornJurkiewiczLoll2001} with $k=a^2/(8\pi G_*)$ and the counting identity $N_2^{\rm TL}=2N_0+N_{41}+2N_{32}$ that follows from their Dehn--Sommerville relations, the equilateral point $\tilde\alpha=1$ gives
\[
\boxed{\kappa_0=\frac{\sqrt3\,a^2}{8G_*}=2\ln2=1.3863,\qquad\Delta=0,}
\]
independently of the bare cosmological term.  Along $\tilde\alpha$ with $\lambda=0$ (Appendix~O.5), $\kappa_0$ rises slowly, $1.39$, $1.47$, $1.93$, $2.33$ at $\tilde\alpha=1,2,5,8$, and $\Delta$ rises from $0$ through $0.11$ and $0.26$ to $0.36$.  The implementation reproduces the published coefficients of the Lorentzian and Euclidean actions at $\alpha=\pm1$ to three decimals, and its inversion of the standard de~Sitter point $(\kappa_0,\Delta)=(2.0,0.4)$ returns $\tilde\alpha=0.71$, $a^2/G=26$, and a bare cosmological coefficient $\lambda\simeq86$, in the direction and magnitude of the published $\tilde\alpha(\Delta)$ curve at $\kappa_0=2.2$ \citep{AmbjornGJL2012}.  The inversion exposes a structural fact: at the standard de~Sitter point the positive $\Delta$ is almost entirely induced by the bare cosmological term, $\Delta_\lambda=\lambda[\sqrt{8\tilde\alpha-3}-\sqrt{12\tilde\alpha-7}]/96$, while the geometric asymmetry contributes about $0.01$.  The lattice spacing is a second cross-check: $a=2.53L_*$ is $2.53$ Planck lengths on the matched branch, against the roughly two Planck lengths that semiclassical de~Sitter fits assign to the CDT lattice spacing \citep{AmbjornGJL2008dS}.

\paragraph{The time asymmetry.}
A $(4,1)$ simplex places its apex above the base circumcenter at Lorentzian proper-time height $h=a\sqrt{\alpha+3/8}$, so the number of renewal ticks per slice step is $n=h/L_*$.  Two kinematic floors precede any dynamics: $n>R/L_*=1.55$ for time-like edges to be time-like, and $n>2.48$ for the Euclideanized $(3,2)$ simplices to be nondegenerate ($\tilde\alpha>7/12$).  One tick per slice is impossible; the slice separation is at least three ticks.  The reference assignment of one physical event to a simplex is an additional convention. Assigning each such simplex four-volume $L_*^4$ gives $a^4v_{41}(\tilde\alpha)=L_*^4$ at $\tilde\alpha=1.06$ and $a^4v_{32}(\tilde\alpha)=L_*^4$ at $\tilde\alpha=1.04$, hence
\[
\boxed{\tilde\alpha\simeq1.05,\qquad n=3.0\ \text{ticks per slice},}
\]
with the two simplex volumes agreeing to two percent and $n=3$ coinciding with the depth-three tetrahedral update of Appendix~H.8.  The alternative reading, one regular cell of three-volume $a^3/(6\sqrt2)=1.91L_*^3$ swept through one tick, gives $v_*=1.91L_*^4$ and $\tilde\alpha=2.3$--$2.9$.  Either way $\tilde\alpha$ is of order one.  The native curve crosses the published bifurcation--de~Sitter line only at $\tilde\alpha\simeq7.1$, seven ticks per slice and volume $3.1L_*^4$ per simplex. Neither reference assignment above selects that value. C.8 calibrates an event per primitive pair, and J.11 keeps the event-per-simplex conversion explicit; neither identifies a triangulation simplex with that physical pair automatically.

\paragraph{The capacity asymmetry term.}
The uncentered capacity free energy $\beta\mu N_{\rm cells}=\beta\mu N_{41}/2$ is an $N_{41}$-only term and therefore the coupling displacement $(\delta\kappa_0,\delta\Delta,\delta\kappa_4)=(-6,+1,-1)\,\beta\mu/2$, with $\delta\Delta=+1.24$ (Appendix~J.2).  It cannot supply the de~Sitter asymmetry, for two independent reasons.  The closure energy $K^2=48-(S^2-\Sigma^2)/3$ carries a conventional origin, and $E_c\mapsto E_c+C$ leaves $\widehat Z_{\rm cap}$ exactly invariant, so no origin-independent value of the extensive coefficient exists; per-cell centering fixes the chosen capacity reference; the gravitational initial condition $\Lambda_i=0$ is supplied separately in O.5.  And an $N_{41}$-only term has no $N_0$ part, so $\delta\kappa_0=-6\delta\Delta$: the uncentered value sends $\kappa_0$ to $-6.1$, and even the $\delta\Delta\simeq+0.35$ that the bifurcation line would require sends it to $-0.7$.  The same arithmetic is why the standard de~Sitter point needs $a^2/G\simeq26$ rather than $6.4$.

\paragraph{The native point and its test.}
The registered face-wise branch therefore has coupling point
\[
\boxed{(\kappa_0,\Delta)_{\rm native}=(2\ln2,\ \simeq0.01)}
\]
under the $L_*^4$ convention, or $(1.5$--$1.6,\ 0.13$--$0.17)$ under the swept-cell convention.  On the published bare-measure phase diagram \citep{AmbjornNewPhase2017} the first lies within the scatter of the second-order $B$--$C_b$ line at this $\kappa_0$ and the second inside the bifurcation phase; neither is the de~Sitter phase.  Three readings are open.  The face-wise lock may be wrong, returning metricity to the selected status of Appendix~H.9a, paragraph~4.  The published diagram may not be the right map: it is drawn for the bare $1/C_T$ measure, whereas the theory's measure carries $e^{-\eta_*K^2}$ with injectivity on every face, a strong non-extensive constraint (Appendix~J.5 records that it moves the slice collision fraction from $0.65$ to $0.08$) whose quadratic curvature dressing $c_0\simeq0.019$ does not bound its effect on phase boundaries.  Or the native point is critical, sitting on CDT's second-order line without tuning.  The second reading retains the face-wise and volume premises and asks whether their decorated measure changes the phase.  Its direct test is a run of the weighted ensemble at the predicted couplings, with the phase diagnostics of Appendix~J.1.  The bare-measure localization uses the tuned grid at $N_{41}=10^4$, eighty slices, $\kappa_4$ auto-tuned, at $\kappa_0\in\{1.0,1.386,1.8\}\times\Delta\in\{0,0.1\}$ and at $(1.55,0.15)$, read against the same engine's grid at $\Delta\ge0.2$.  Extended: adjacent-slice correlation, blob score, and $d_H$ comparable to the engine's reference point $(2.0,0.4)$ at the same volume ($0.48$, $2.78$, $2.56$), with maximum slice fraction below $0.05$.  Time-collapsed: maximum slice fraction of order one with a single active slice.  A pass is followed by the volume-pair confirmation of Appendix~J.1.  This bare calculation locates the predicted couplings in the bare phase map; the decorated run determines whether the capacity measure moves the relevant boundary.  Either outcome is a finite-volume phase diagnostic at the grade of Appendix~J.1, not a continuum or massless-TT statement.

\paragraph{Bare-measure result at \(N_{41}=10^4\).}
The native point gives different basins across seeds and fails the registered volume-pair criterion at forty thousand simplices. The reference point passes under the same protocol. The registered bare grid was run with the engine of Appendix~J.1 (eighty slices, four hundred sweeps, $\kappa_4$ auto-tuned over three bursts, one seed per point); every triangulation passed the manifold and foliation audits.  No point collapsed.

\begin{center}
\small
\begin{tabular}{cc|ccccc|c}
\hline
$\kappa_0$ & $\Delta$ & $\kappa_4$ & $d_H$ & blob & active & slicecorr & maxfrac \\
\hline
$1.000$ & $0.00$ & $0.680$ & $2.79$ & $3.99$ & $15$ & $0.29$ & $0.050$ \\
$1.386$ & $0.00$ & $0.752$ & $2.70$ & $5.05$ & $9$ & $0.17$ & $0.063$ \\
$1.800$ & $0.00$ & $0.808$ & $2.58$ & $6.61$ & $5$ & $0.03$ & $0.083$ \\
$1.000$ & $0.10$ & $0.720$ & $2.67$ & $3.15$ & $13$ & $0.17$ & $0.039$ \\
$1.386$ & $0.10$ & $0.780$ & $2.69$ & $4.24$ & $9$ & $0.30$ & $0.053$ \\
$1.800$ & $0.10$ & $0.788$ & $2.65$ & $2.57$ & $18$ & $0.00$ & $0.032$ \\
$1.550$ & $0.15$ & $0.756$ & $2.65$ & $2.99$ & $21$ & $-0.01$ & $0.037$ \\
\hline
$2.000$ & $0.40$ & $0.752$ & $2.56$ & $2.78$ & $12$ & $0.48$ & $0.035$ \\
\hline
\end{tabular}
\end{center}

The last row is the engine's reference point at the same volume.  The native point $(2\ln2,0)$ is extended: nine active slices, a maximum slice fraction of $0.063$ against the order-one value of a time-collapsed geometry, and $d_H=2.70$, near the top of the grid (the $(1.0,0)$ point gives $2.79$).  Its volume profile is more concentrated than the reference (blob $5.0$ against $2.8$, fewer active slices, adjacent-slice correlation $0.17$ against $0.48$).  Along $\Delta=0$ the profile concentrates monotonically with $\kappa_0$ (active slices $15\to9\to5$, blob $4.0\to5.0\to6.6$), so the native point lies midway along the concentration direction.  The reading-B point $(1.55,0.15)$ is indistinguishable from the engine's $\Delta=0.2$ row.

\paragraph{Volume ladder and collapse control.}
\begin{center}
\small
\begin{tabular}{lc|ccccc}
\hline
point & $N_{41}$ & $d_H$ & blob & active & slicecorr & maxfrac \\
\hline
native $(2\ln2,0)$, seed 0 & $5{,}000$ & $2.42$ & $2.94$ & $14$ & $0.20$ & $0.037$ \\
native, seed 0 & $10{,}000$ & $2.70$ & $5.05$ & $9$ & $0.17$ & $0.063$ \\
native, seed 0 & $20{,}000$ & $2.71$ & $5.02$ & $13$ & $-0.06$ & $0.063$ \\
native, seed 0 & $40{,}000$ & $2.48$ & $4.72$ & $13$ & $0.04$ & $0.059$ \\
native, seed 1 & $20{,}000$ & $1.69$ & $8.53$ & $5$ & $-0.09$ & $0.107$ \\
native, seed 1 & $40{,}000$ & $2.88$ & $9.42$ & $6$ & $-0.04$ & $0.118$ \\
control $(1.386,-0.3)$ & $10{,}000$ & $2.56$ & $12.3$ & $7$ & $-0.09$ & $0.154$ \\
control $(2.5,-0.3)$ & $10{,}000$ & $2.43$ & $10.7$ & $10$ & $-0.09$ & $0.134$ \\
reference $(2.0,0.4)$ & $10{,}000$ & $2.56$ & $2.78$ & $12$ & $0.48$ & $0.035$ \\
reference $(2.0,0.4)$ & $40{,}000$ & $3.00$ & $1.77$ & $18$ & $0.20$ & $0.022$ \\
\hline
\end{tabular}
\end{center}

The control fixes what collapse looks like in this engine at this volume: at $\Delta=-0.3$ the maximum slice fraction is $0.13$--$0.15$ and the blob score $11$--$12$, not the order-one values of the published time-collapsed phase, so the published collapse signature does not transfer to this engine's conventions any more than the published de~Sitter boundary does.  Against that calibration the first seed at the native point is a stable extended geometry, unchanged from ten to twenty thousand simplices in $d_H$, blob, and maximum slice fraction, with the blob rising from five to ten thousand and then flat; the registered criterion is met on the first rung and flat on the second, and the maximum slice fraction did not fall.  The second seed at twenty thousand simplices lands elsewhere: $d_H=1.69$, blob $8.5$, five active slices, maximum slice fraction $0.107$, much closer to the $\Delta=-0.3$ control.  The two seeds at the same couplings sample different basins. Longer runs from both starts test whether these basins persist at the native couplings.  The forty-thousand-simplex rung of the first seed does not resolve it in the extended direction: the point stays uncollapsed (maximum slice fraction $0.059$, thirteen active slices, pin residual $0.012$), but the blob score is flat to falling ($5.02\to4.72$) and $d_H$ falls from $2.71$ to $2.48$, the opposite of the reference point's rise with volume in Appendix~J.1.  The registered volume-pair criterion therefore fails at this rung.  The confound that the grid protocol uses four hundred sweeps whereas the rising $d_H$ of Appendix~J.1 came from four-thousand-sweep runs was tested directly: the reference point $(2.0,0.4)$ at forty thousand simplices under the same four-hundred-sweep protocol gives $d_H=3.00$, up from $2.56$ at ten thousand, with the profile spreading (blob $1.77$, eighteen active slices, maximum slice fraction $0.022$).  The protocol reproduces the volume trend at the reference point.  A second seed at the native point at forty thousand simplices lands in the concentrated basin ($d_H=2.88$, blob $9.4$, six active slices, maximum slice fraction $0.118$), as it did at twenty thousand.  The registered bare-measure volume-pair test therefore fails at the predicted couplings while passing at the reference under identical conditions.   The first reading (the face-wise lock is wrong, returning metricity to the selected status of Appendix~H.9a), the second (the decorated measure moves the phase boundary), and the third (a critical, seed-unstable boundary point) remain open.  The weighted run tests the second reading; a production study with several seeds is required to separate the remaining boundary behavior.

\paragraph{The shape sector does not change this.}
A frozen-area control tests whether the remaining cell shapes can carry the missing metric modes. Area-Regge calculus permits two fluctuating shape parameters for each disphenoid cell and massive shape-mismatch modes \citep{DittrichKogios2023}; the fixed-edge CDT ensemble restricts this continuous geometry.  A linearized area-Regge computation on a periodic Kuhn-subdivided hypercubic lattice (generic flat background; flat deficits, the Schl\"afli identity, and Hessian symmetry verified to $10^{-6}$) reproduces the known structure for the full theory: ten massless modes per vertex, four of them lattice diffeomorphisms and six with $\lambda\propto k^2$, and forty massive modes with gap $|\lambda|\simeq0.5$--$0.8$ in lattice units.  With the spatial areas frozen at $a_{\rm ch}$, as the fixed $j=3$ capacity faces require, the continuous sector retains for spatial momenta only two $k^2$ modes of transverse-shift type and for momenta along the slicing direction none: the purely spatial transverse-traceless polarizations are frozen with the spatial metric.  On this frozen-area branch the shape modes cannot carry spatial TT gravitons, and the remaining block geometry must come from connectivity. The directional-count construction of Q.11a allows coarse areas to vary and has a different continuous metric sector, tested in H.13 and N.10.  The same modes do not fix the time asymmetry either: their one-loop effective potential for $\tilde\alpha$ is dominated by the measure and by the arbitrary background deformation (opposite slopes under flat-in-area and scale-free measures, slope differing threefold between backgrounds), so $\tilde\alpha$ remains a modulus at this order.  The measured mismatch gap gives $m a\simeq0.75$; with $a=2.53L_*$, the mismatch modes decouple beyond $L_{\rm block}=a/(ma)\simeq3.4L_*$, inside the provisional bracket.

\subsection*{J.11 Volume-profile normalization and the channel cut}
\label{aqg:newton}

Use \(\hbar=c=1\) for the volume-profile comparison. The conversion of a geometric lattice into physical clock events can
be separated from the measured volume-profile stiffness.
Let \(\widetilde s_0\) be the dimensionless de~Sitter profile width,
\(k_1\) its kinetic coefficient, and \(\tau_{\rm slice}\) the physical
slice separation. With the same volume-count convention in the profile
and kinetic fit,
\[
\frac{G}{\tau_{\rm slice}^2}
=\frac{2\pi}{9}\frac{\widetilde s_0^4}{k_1},\qquad
\boxed{\frac{G}{L_*^2}
=\frac{\widetilde s_0^2}{3\sqrt6\,k_1}
\sqrt{\frac{n_{\rm ev}v_{\rm ev}}{L_*^4}}.}
\]
Here \(n_{\rm ev}\) is the physical event count per counted \((4,1)\)
simplex, with the \((3,2)\)-simplex volume share included consistently.
Lattice spacing, simplex shape, and the intermediate time and volume
conventions cancel. The event assignment remains a physical input.
The simulations contain geometric simplices and capacity variables;
they contain no independently identified clock-event register from
which to measure \(n_{\rm ev}\).

For \(v_{\rm ev}\) near \(4L_*^4\), the available profile fits would
require roughly \(0.046\)--\(0.145\) events per counted simplex to give
\(G=G_*\), or about seven to twenty-two simplices per event. An edge
near \(L_*\) can realize that conversion, but the absence of an
embedding lattice does not determine an edge length. The choice
\(n_{\rm ev}=\eta_*\) has no derivation: \(\eta_*\) is the closure
stiffness and, under the common energy conversion of O.13, the released
share of committed capacity. It is not an event probability.
Counting causal links gives definition-dependent channel densities and
does not resolve the calibration.

The angular construction offers a stronger comparison, because bulk and
cut measurements use the same bare action. For the registered native
pilot, the convention-free ratio is
\[
\frac{G_{\rm bulk}}{G_{\rm cut}}
\simeq0.076503\,\frac{\widetilde s_0^{4/3}}{k_1}.
\]
At \(\widetilde s_0=0.7494807394\) and \(k_1=0.03467387683\),
this gives \(1.502071459\).
Equality would require \(k_1=0.0520826408\).
With the illustrative event calibration \(m_4=2.35\), the same
comparison reads
\[
G_{\rm cut}=1.030021G_*,\qquad
G_{\rm bulk}=1.547165G_*.
\]
Changing \(m_4\) or relabeling clock events rescales both and leaves
their ratio unchanged. The spread of values at different bare
couplings is not an uncertainty interval for this run.
Finite-volume effects, the kinetic fit, and renormalization require
tests within one common ensemble before interpreting this mismatch.
The boundary normalization cannot independently replace the bulk
coupling of that ensemble.

The face-wise branch of J.10 and the multi-channel coarse cuts of Q.11a
thus define distinct geometric implementations. Freezing all primitive
spatial areas removes the continuous spatial TT modes, as the J.10
control shows. Allowing directional channel counts to vary restores a
metric response and gives the H.13 and N.10 tests, without identifying
one binary channel with each regulator triangle.
The registered native campaign and its acceptance criteria remain
tests of the face-wise and event-per-simplex conventions under which they were registered.
Passing a time calibration under those conventions does not also
establish the required Newton normalization.

\subsection*{J.12 Connected geometry and the covariant host dictionary}
\label{aqg:host}

The connected tensors of Q.11b realize the selected causal
triangulation weight with normalized microscopic boundary maps.
The event and area calibration of C.8 gives, in geometric units,
\[
\frac{a^2}{G_0}=\frac{16\ln2}{\sqrt3}=6.40301938055.
\]
The illustrative chain coefficient \(m_4=2.35\) gives
\(a=2.56812077553L_*\) and \(G_0/G_*=1.03002098319\).
Changing this coefficient rescales \(a^2\) and \(G_0\) together;
it leaves the dimensionless curvature coupling unchanged. It does
not identify independently chosen volume ensembles.

Three regulator prescriptions enter the comparisons. The face-wise
matching of J.10 sets \(G_0=G_*\), \(a\simeq2.530L_*\), and its
near-unit asymmetry through the specified event assignment.
The continuous geometric interpolation with the paired calibration
uses \(\alpha=11.312\). The connected integer-channel construction
uses \(\alpha=7\). For the standard Euclidean count action,
\[
S_E=-(\kappa_0+6\Delta)N_0
 +(\kappa_4+2\Delta)N_{41}+(\kappa_4+\Delta)N_{32},
\]
the two paired-calibration curvature maps are
\[
\begin{array}{c|rrr}
\alpha&\kappa_0&\Delta&\kappa_4^{\rm curv}\\ \hline
11.312&2.70257004&0.43691236&-0.42026496\\
7&2.20205311516&0.326138328034&-0.304774619960
\end{array}
\]
The curvature contribution \(\kappa_4^{\rm curv}\) is not the tuned
critical cosmological coupling. Fugacities, a physical-volume window,
or a fixed count specify additional ensemble data.

The simplex four-volumes are
\[
V_{41}^{L,E}=\frac{a^4}{96}\sqrt{8\alpha\mathbin{\pm}3},
\qquad
V_{32}^{L,E}=\frac{a^4}{96}\sqrt{12\alpha\mathbin{\pm}7},
\]
with the plus sign Lorentzian. For
\(r=N_{32}^{\rm slab}/N_{41}^{\rm slab}\) and a trapezoidal
spatial-volume convention, the volume-equivalent duration satisfies
\[
\frac{c\tau_{\rm vol}^{L,E}}{L_*}
=\frac{a/L_*}{4\sqrt2}
\left[\sqrt{8\alpha\mathbin{\pm}3}
 +r\sqrt{12\alpha\mathbin{\pm}7}\right].
\]
At \(\alpha=11.312\),
\[
\frac{c\tau_{\rm vol}^{L}}{L_*}
=4.38972153+5.42399632r.
\]
This volume-to-duration conversion depends on composition. A host
interval is not generally one native tick. The volume clock,
an intrinsic geodesic duration, and the native matter clock of H.15
are separately defined observables.

The normalized tensor construction permits direct vacuum sampling
of its CDT weight without summing the primitive indices afresh.
The graph counting, automorphism quotient, boundary normalization,
and volume prescription must still match Q.11b. A sampling move
sequence adds no physical history multiplicity. Existing
capacity-decorated, face-wise, or differently constrained campaigns
do not measure the \(\alpha=7\) ensemble merely by sharing its
simplex types. The 1,920-slab sewing check is a restricted finite
family, not a phase or continuum measurement.

\subsection*{J.13 Proper-length probes and the interacting Newton test}
\label{aqg:force}

For mass-independent source preparations with Euclidean proper
lengths \(L_1,L_2\), define in units \(\hbar=c=1\)
\[
W_i=e^{-m_iL_i},\qquad
R=\frac{\langle W_1W_2\rangle}
        {\langle W_1\rangle\langle W_2\rangle}.
\]
The ratio removes the isolated-source contributions, and
\[
\left.\partial_{m_1}\partial_{m_2}\log R\right|_0
=\operatorname{Cov}(L_1,L_2).
\]
For a stationary long-duration preparation at held physical
separation \(r\), its interaction coefficient is
\[
Q(r,T)=-\frac{\operatorname{Cov}(L_1,L_2)}{T},
\qquad Q(r,T)\longrightarrow-\frac{G_R}{r}
\]
in an Einstein infrared regime. Fixing one length on every
geometry would remove the covariance. A duration defined by an
ensemble clock average preserves this response.
For the supplied Einstein Gaussian propagator, the finite-duration
benchmark is
\[
Q(r,T)=-\frac{2G}{\pi r}
\left[\arctan(T/r)-\frac{r}{2T}\log(1+T^2/r^2)\right].
\]
It checks normalization and the large-\(T\) limit for that propagator;
it does not determine the interacting triangulation propagator.

Holding separation requires intrinsic source transport. Signed
simplex edge lengths determine the Lorentzian metric on barycentric
tangent vectors, and face isometries transport the first source
and its frame. In the induced slice metric a corridor calculation
minimizes the physical segment lengths through face contacts.
If the endpoint separation is \(D\),
\[
\dot D=A_D+e\cdot v,\qquad
v_{\min}=-A_De,\qquad U_2=U_\perp-A_De,
\]
where \(e\) is the unit spatial endpoint gradient and \(A_D\)
includes the changing geometry and first source. The minimum-norm
held velocity is timelike when \(-N^2+A_D^2<0\).
Adaptive step reduction tests the actual connecting segments.
Nonunique contacts require a branch prescription.

An explicit path gives an upper bound on distance. A continuous
finite-element slice field, independently rescaled to satisfy
\(|\nabla f|\leq1\), gives the global lower bound
\[
f(q)-f(p)\leq d(p,q).
\]
The field can be refined without changing the physical triangulation.
In the four-interval diagnostic, all 403 sampled initial, endpoint,
and midpoint locations pass the separation criterion. This is a
sampled feasibility check; it supplies neither a continuous
supremum bound nor whole-ensemble source coverage.

Let \(\mu_g\) be a normalized initial-mark distribution on each
geometry and leave final marks unobserved. The sourced partition
function is
\[
Z_{m_1m_2}=\sum_gP_0(g)\int\mu_g(dX_0)\,
 e^{-m_1L_1-m_2L_2}.
\]
It has \(Z_{00}=1\) when the source exists on all of the support.
Deterministic transport is a normalized delta kernel; its coordinate
Jacobian adds no geometry weight. Rejecting geometries on which
the source fails would change \(P_0\). A response is therefore
withheld when its requested ensemble coverage is incomplete.
Initial-mark mixing also contributes a within-mark covariance
that must be separated from geometric response. This external
preparation is distinct from a derived native matter path measure.

The saved diagnostic uses forty inputs from the \(\alpha=11.312\)
selected action, with an approximately two-percent hard Lorentzian
four-volume window and no extra volume fugacity. Its primary
separations are \(r/a=0.3,0.5,0.7\), with durations of one, two,
and four host intervals. The \(0.9a\) source separation remains
unresolved. The archived data supply no accepted complete
ensemble force curve. The action chain is short and drifting,
with roughly five effective samples in the saved diagnostic,
and these sub-edge separations provide no window
\(a\ll r\ll L_{\rm box}\).

An independent exact two-state check verifies the estimator's
normalization, force sign, and suppression of incomplete coverage.
An infrared \(G_R\) still requires equilibrated ensembles, a
resolved distance window, duration stability, and a calibrated
physical source. The negative covariance of the freely transported
native lines in H.14 is not a held-separation force measurement.

\subsection*{J.14 Registered connected-host normalization and matter checkpoint}
\label{aqg:alpha7result}

The connected $\alpha=7$ construction fixes the dimensionless curvature
normalization and the two curvature couplings,
\[
 \frac{G_0}{a^2}=\frac{\sqrt3}{16\ln2}=0.1561763194,\qquad
 (\kappa_0,\Delta)=(2.20205311516,0.326138328034).
\]
The associated geometry weight is sampled with the capacity gate off.
The cosmological fugacity and volume feedback remain ensemble controls.
The registration of 17 October 2026 and its amendment were frozen before
sampling at these couplings. The amendment made $N_{41}=32{,}000$ the
primary volume and retained the 16k and 64k runs as reported-only checks.
The target counts $N_{41}$ simplices; the total number of four-simplices is
$N_4=N_{41}+N_{32}$.

\paragraph{Registered observables.}
Q1 tests whether the de~Sitter-like blob is independent of its starting
configuration. Q2 compares the Newton coefficient inferred from the
homogeneous volume stiffness with the boundary coefficient $G_0$:
\[
 R\equiv\frac{G_{\rm vol}}{G_0}
 \simeq 0.076503\,\frac{\widetilde s_0^{4/3}}{k_1}.
\]
The width $\widetilde s_0$ and bias-corrected local stiffness $k_1$ are
measured; their conversion is fixed by the registered dictionary. The
ratio is independent of the provisional chain constant $m_4$. No
renormalized stiffness or source response was tuned to make $R=1$.
The written prior was $R\simeq1.4$, with disagreement considered more likely.

The final window contains 200,000 sweeps per run. Blob eligibility requires
a relative profile-fit error at most $0.05$, width $2.5\leq w\leq K/8$,
and every registered half-window stationarity statistic to satisfy
$|z|\leq3$. Here $K=80$ is the number of time slices. Q2 uses eligible
32k runs only. Its registered agreement rule places the central geometric
mean in $0.80\leq\overline R\leq1.25$; it does not require the entire
uncertainty interval to lie inside that band.

\begin{center}
\small
\begin{tabular}{lrrrll}
\hline
Run & $N_{41}$ target & $R$ & $\sigma_{\ln R}$ & Blob eligibility & Role\\
\hline
A3 & 32,000 & 1.198 & 0.059 & Drifting & Primary pool\\
A4 & 32,000 & 1.064 & 0.089 & Pass & Primary pool\\
A6 & 32,000 & 1.027 & 0.062 & Drifting & Primary pool\\
A1 & 16,000 & 1.164 & 0.061 & Drifting & Reported only\\
A2 & 16,000 & 1.060 & 0.065 & Pass & Reported only\\
A5 & 64,000 & 0.893 & 0.054 & Pass & Reported only\\
\hline
\end{tabular}
\end{center}

\paragraph{Normalization result.}
A4 is the sole eligible primary run. It gives
\[
 \overline R=1.064,\qquad \sigma_{\ln\overline R}=0.089,
 \qquad \text{Q2: AGREE}.
\]
The approximate two-standard-error interval is $0.89$--$1.27$.
Its relative profile-fit error is $0.013$. A2 and A5 also have stationary
de~Sitter-like profiles, with fit errors $0.015$ and $0.012$, and their
ratios lie inside the same agreement band. They support the volume
comparison without changing the primary verdict. Each eligible volume
has one stationary run. The change from A4 to A5 is $1.7$ standard errors
in $\ln R$; a large-volume limit has not been estimated from these data.
The specified microscopic history measure thus has a finite-volume
stiffness compatible with its own boundary normalization under the
registered conversion.

\paragraph{Phase and estimator checks.}
Q1 remains \textsc{undetermined}. A3 fails stationarity through $\xi$
at $|z|=3.3$, and A6 through $\widetilde s_0$ at $|z|=3.3$.
A3 and A4 also differ in largest-slice fraction by $6.1$ standard errors.
The existence of stationary blobs therefore coexists with an unresolved
test of start independence. A post hoc inverse-covariance calculation,
after removing the active small-sample correction floor, gives
$R=1.179,1.022,0.973$ for A2, A4 and A5. These are exploratory estimator
checks; the local-stiffness values above determine the registered result.

\paragraph{Matter-response status.}
A separately registered curvature-trace test used two new streams at
fixed $\kappa_4$, with 64 saved geometries and two source marks per
geometry in each stream. Its dimensionless diagnostic was
$\gamma_{\rm trace}=1.82\pm1.04$, with approximate 95\% interval
$[-0.65,4.28]$, against the target $G_0/a^2=0.1562$.
Both streams also exceeded the registered maximum blob width of 10,
with widths $11.17$ and $10.61$. The verdict is
\textsc{undetermined}; no matter Newton coefficient was reported.
This separate fixed-fugacity gate failure leaves Q2 unchanged.

Subsequent quenched-scalar measurements on those 128 saved geometries
are development analyses. They use both saved sources and six bare
masses. The inferred binding slopes depend on distance range, binning
and scalar discretization, while fixed nongravitating controls also
produce apparent short-range slopes. A conditional-variance diagnostic
does not restore a stable binding interpretation. These measurements
supply no accepted binding energy or gravitational null result.
Further matter-response production is deferred pending a calibrated
finite-geometry observable and an independent power assessment.

The completed evidence concerns the microscopic realization, stationary
finite-volume profiles and the registered homogeneous stiffness ratio.
Equality with a matter-source Newton coefficient, a common matter cone,
and the interacting massless transverse--traceless sector remain open.
The illustrative $G_0/G_*=1.030$ at $m_4=2.35$ retains its separate
event-volume and chain-clock premises.

\paragraph{Evidence record.}
The pre-data registration, amendment, six configurations, final verdicts,
analysis code and the matter-development readouts are copied with their
original bytes into the companion CDT evidence directory
\nolinkurl{evidence/20261020_alpha7_checkpoint}. Its manifest records
source paths, lengths and SHA-256 digests. The raw simulation archives
remain at their recorded locations. This packet supports the numerical
claims without rerunning or altering the campaigns \citep{Alpha7Checkpoint2026}.

\section*{Appendix K: Numerical and Consistency Checks}

Appendix~K recomputes the numerical consequences of the stated derivations and collects the substrate-scale ledger. These checks support reproducibility but do not establish the ontology independently. Section~24 and Appendix~J perform a separate task by coupling the same coefficient chain to a dynamical substrate. The complete self-contained reproduction code appears in terminal Appendix~R.

The numbers recomputed here, and where each enters the main text, are: the admissibility entropy $g_{\mathrm{share,eff}}$ (Section~13, Appendices~B and~D.4); the substrate length $L_*$ (Sections~2 and~13, Appendix~D.4); the induced scale $G_*$ (Sections~10 and~13, Appendices~D.4 and~K); the galactic scale $a_0$ (Section~15, Appendix~C.7); the diamond-lattice Green constant $G_{\mathrm{tet}}(0)$ (Appendix~C.5); the edge-kernel chain $J_{\mathrm{bare}}\to\gamma$ and the return sum $\Sigma_{\mathrm{ret}}$ (Appendices~C.3--C.4); the charged-lepton ladder with its support-exponent cross-check (Section~13, Appendix~I.1); and the primitive route, adjoint-channel, and causal-blocking identities of the conditional color branch (Appendix~I.2).

\subsection*{K.1 Cross-sector numerical checks}

The cross-check program includes:
\begin{itemize}
\item the one-bit fermionic defect check $\Delta S_f=\ln 2$;
\item the rooted-shell convergence check $\sigma_{\mathrm{ind}}^{(2)}\simeq \sigma_{\mathrm{ind}}^{(3)}$;
\item the UV closed-branch moments $\langle K^2\rangle_{\eta_*}$, $\mathrm{Var}_{\eta_*}(K^2)$, and $a_{\mathrm{UV}}$;
\item the Gaussian closure-amplitude identity, nine-state response Gram matrix, unitary pair dilation, 21-block edge Hessian, and 56,800-state routing enumeration;
\item the closed routing covariance, primitive route Gram spectrum, marked-fusion representation match, $t_8=13/14$ channel, and $2\sqrt6$ diamond blocking factor;
\item cross-sector consistency among the electron anchor, the substrate length $L_*$, the induced scale $G_*$, the Green-matched Newton closure, and the galactic scale $a_0$.
\end{itemize}
These checks independently reproduce the coefficient chain in every sector where closure is claimed. The derivations remain the source of the coefficients.

\subsection*{K.2 Reproducibility ledger for the substrate scale}

The substrate-length calculation is short enough to record as a numerical ledger. Enumerating the $1680$ oriented injective tetrahedral states with labels $m=-3,\ldots,3$, weighting them by $e^{-\eta K^2}$, and solving
\[
\langle K^2\rangle_\eta=\frac{3}{2\eta}
\]
gives
\[
\eta_*=0.02986684439352237,
\qquad
g_{\mathrm{share,eff}}=7.419800023570903.
\]
The one-pass seven-sector history support and its transverse export are then
\[
e^{7g_{\mathrm{share,eff}}}
=
3.602860521062804\times10^{22},
\]
\[
\frac23 e^{7g_{\mathrm{share,eff}}}
=
2.401907014041869\times10^{22}.
\]
Using $\lambda_e=\hbar/(m_ec)=3.86159267435238\times10^{-13}\ {\rm m}$ and $r=e^{-7g_{\mathrm{share,eff}}}$ first gives the baseline
\[
L_*^{(0)}=-\frac32\lambda_e\ln(1-r)
=
1.60771947114394\times10^{-35}\ {\rm m},
\]
while the marked vertex gives
\[
\zeta_*=0.005123584484947,
\qquad
Z_e=1.005308283809514.
\]
Therefore
\[
L_*=Z_eL_*^{(0)}
=1.61625370238286\times10^{-35}\ {\rm m},
\]
and
\[
G_*=\frac{c^3L_*^2}{\hbar}
=
6.67428908131019\times10^{-11}\ {\rm m^3\,kg^{-1}\,s^{-2}}.
\]

The source-side Green constant is obtained independently from the diamond-lattice integral in Appendix~C.5. Uniform-grid quadrature with endpoint extrapolation gives
\[
G_{\mathrm{tet}}(0)=0.448220394388\ldots ,
\]
confirming the exact Joyce value
\[
G_{\mathrm{tet}}(0)
=
\frac{3\Gamma(1/3)^6}{2^{14/3}\pi^4}.
\]
This is the value used in the source theorem. These numbers are not additional inputs; they are the numerical evaluation of the finite spectrum, the seven-sector history support, and the graph Green function already defined in the derivation.

\subsection*{K.3 Conditional color and record-retention audits}

The color audit first recomputes the weighted route covariance directly from the same $1680$ states used for admissibility. It returns
\[
C_R=11.7667312558\,P_{\mathbf1}
+35.3002035759\,P_{\mathbf2},
\]
and verifies
\[
\lambda_{\mathbf1}
=\frac34\operatorname{Var}(K^2)
\]
to numerical precision. The primitive recoupling matrix has spectrum
\[
\operatorname{spec}G=\left\{\frac9{10},\frac{21}{20},\frac{21}{20}\right\},
\qquad
\frac{\lambda_{\mathbf2}}{\lambda_{\mathbf1}}=\frac76.
\]
After maximum-throughput normalization, the direct and complementary branches satisfy
\[
K=P_{\mathbf2}+\frac67P_{\mathbf1},
\qquad
L=\frac{\sqrt{13}}7P_{\mathbf1},
\qquad
K^\dagger K+L^\dagger L=I.
\]
The phase-invariant loop cost obeys the representation identity
\[
\overline{\Delta_p}
=\frac32\left(1-\frac{|\operatorname{Tr}_{\mathbf3}U_p|^2}{9}\right)
=\frac43\left(1-\frac18\operatorname{Tr}_{\mathbf8}\operatorname{Ad}U_p\right),
\]
and is unchanged by multiplying any lift of $U_p$ by a common phase. This checks that the derived action is well defined on the projective channel group without assuming the still-open fundamental endpoint lift.
The Haar identities of Appendix~I.2 then give a unital trace-preserving channel with singlet eigenvalue one and eight equal traceless eigenvalues $13/14$. The heat-kernel coordinate and equal-transfer diamond blocking evaluate to
\[
g_{\rm HK,step}^2
=\frac23\ln\frac{14}{13}
=0.04940531477,
\]
\[
2\sqrt6=4.89897948557,
\qquad
g_{\rm graph,HK}^2
=\frac{4\sqrt6}{3}\ln\frac{14}{13}
=0.24203562353.
\]
These computations audit the conditional projective pure-gauge construction. They are not a comparison with a continuum QCD coupling.

The generation cutoff has also been checked without choosing a preferred excluded alphabet. Exhaustive enumeration over all reduced alphabets gives nonzero closure variance for every one-label exclusion (all seven cases, variance range $9.95556$--$16.17778$) and every two-label exclusion (all 21 cases, range $1.24444$--$26.02667$). Every three-label exclusion (all 35 cases) gives zero variance. The cutoff therefore follows from the remaining alphabet size and the permutation symmetry of $K^2$ alone, for every choice of excluded labels.

The record-retention formula has a separate normalization check. For any orthogonal set of fine records merged into $\bar R$, summing $p(h,R_i)$ over both $h$ and $i$ gives $p(\bar R)$, and the conditional distribution in Section~22 therefore sums to one. Record erasure changes the accessible partition of the same joint measure; it adds no transition weight and cannot alter the unconditioned Born marginals.

\section*{Appendix L: Fork accounting for the scale chain}

This appendix records the construction choices behind the induced value of $G_*$ and quantifies the surrounding discrete-choice fork space. The comparison shows how much freedom was available in the scale chain.

\paragraph{Provenance.} By the author's recollection, a target entropy near $7.42$ was identified early in the framework's development by inverting the measured Newton constant, with the ensemble constructed to deliver it jointly with other structural constraints; the documentary record does not reach that genesis layer. The realized ensemble was fixed and motivated in print by information-independence arguments before the induced-$G_*$ chain existed. Its original electron scale and the baseline charged-lepton ladder landed at discrepancies of roughly one percent and one-half percent. Those residuals were public parts of the construction before the decorated marked vertex was developed.

\paragraph{Why the residuals suggested one missing object.}
The amplitudes can be compared on the same scale. The original Newton value was low by about $1.05\%$, while the length entering it was low by half that amount because $G_*\propto L_*^2$. The three required multiplicative uplifts were
\[
\begin{array}{c|c|c}
\text{baseline quantity} & \text{required uplift} & \text{structural reading}\\
\hline
L_*^{(0)}\ \text{inferred from }G & 0.5309\% & \text{electron feedback}\\
m_\mu^{(0)}/m_e & 0.5124\% & \text{universal marked response}\\
m_\tau^{(0)}/m_e & 0.6601\% & \text{universal response plus second shell}
\end{array}
\]
Dividing the tau uplift by the muon uplift leaves a further $0.1470\%$ second-shell contribution. The pattern is not three identical residuals, and it was not used to impose a universal rescaling. It is more specific: two observables require nearly the same charged-transfer amplitude, while the third requires that common amplitude followed by one smaller shell passage. This was the motivation for seeking one marked response with different graph traces.

\paragraph{Why the repair had to be additive.}
By that stage the tetrahedral construction was already doing several independent jobs. Its label space and admissibility weight fixed the closure spectrum and sharing entropy; its edge projection fed the ordinary weak-field coefficient; its transverse export entered the scale, horizon normalization, and galactic branch; and its shell algebra produced the baseline charged-lepton ladder. Removing or retuning any of those ingredients to close one residual would propagate through the rest of the microscopic and macroscopic chain. A viable repair therefore had to preserve the unmarked tetrahedral ensemble and its ordinary transfer operator.

This requirement narrowed the search to a defect-bound additive sector: it must be invisible in the unmarked vacuum, reuse the closure response already present in the renewed cell, activate only on occupied charged-channel pairs, and route through the existing electron and shell graph. These conditions explain the form of the calculation in Appendix~H. They do not by themselves prove its nine-state fiber, 21 pair records, Majorana determinant, or routing polynomials; those are the independent action-level results and audits. Showing this part of the development is appropriate because the residuals were already known. It makes clear which features were motivated by the discrepancy, which were protected by the existing theory, and which had to be derived rather than chosen.

The vertex in Appendix~H was consequently found with the discrepancies known. It contains no continuously adjusted coefficient: $g_{\mathrm{share,eff}}$, $\eta_*$, $2/7$, the nine-state trace, 21 pair records, and the three routing polynomials are fixed by the displayed action. That fact makes the result an action-level postdiction rather than a fit, but it does not make it a blind prediction. Its added evidential content is that one field space and one routing graph account for $G$, $m_\mu/m_e$, and $m_\tau/m_e$ simultaneously, while the nearby shared-mode and off-diagonal alternatives fail. The cluster budget, committed abundance, redshift dependence of $a_0$, and low-acceleration rotation-curve regime are empirical tests of the original ensemble; none entered either the early entropy calibration or the marked residual calculation.

\paragraph{The audit.} The exact substrate length is proportional to $-Z_e\ln(1-e^{-7g_{\mathrm{share,eff}}})$, with leading behavior $L_*\propto Z_ee^{-7g_{\mathrm{share,eff}}}$, so the calibration is meaningful only relative to the space of constructions that could have been written instead. The original grammar contains six label counts, five multiplicity rules, five closure coefficients, four support exponents, and five export prefactors, for 3000 candidate constructions. Those counts predate the marked vertex and must be recomputed with its $Z_e(g,\eta)$ dependence before they are used quantitatively.

The marked-sector audit separately enumerates 576 nearby formulas obtained by varying the response dimension, pair count, direction factor, determinant power, statistics, and overlap. Only the decorated-vertex formula lies within one muon standard deviation, and that exclusivity must be read carefully. The one-sigma muon window on $\zeta_*$ is about four parts per million relative, far finer than the spacing of any discrete formula menu, so at most one candidate could ever land inside it, and the failure of the other $575$ carries little information on its own. The menu tests only qualitative brittleness: a shared edge register, directed complex determinant, stable bosonic determinant, or six-state symmetric response gives a visibly different result, and no probability measure is assigned to the menu. The quantitative audit is the one-integer scan of Section~13.5, which holds $\zeta_*$ fixed and asks the measured Newton constant to select the routing integer: $n=7$ is the unique survivor over $|n|\le200$, its neighbors fail at $2.3$--$2.4\sigma$, and the current PDG tau comparison selects $2/7$ within the stated finite rational set. The dark-energy sector carries the opposite provenance: its accumulation law, equation-of-state shape, and branch windows were fixed before the distance-data comparison of Appendix~O.12 was run, so that comparison is a test rather than a calibration.

The galactic scale depends only linearly on $g_{\mathrm{share,eff}}$, so it tests the coupling form more strongly than the exact ensemble choice. The corrected $G_*$ and charged-lepton ratios test the shared marked vertex through three different graph factors, with the postdiction status stated above.

\paragraph{Fork accounting for the acceleration normalization.}
The coefficient $\epsilon=g_{\mathrm{share,eff}}/N$ enters both $a_0=\epsilon cH_0$
and $\Omega_c/\Omega_b=1/\epsilon$. The entropy is calibrated on the Newton
constant; the loading of one sharing entropy into the tilt-mode lever is
the premise of Section~15, and the horizon temperature supplies $N=(2\pi)^2$
for a lever measured in apparent-horizon radii. The earliest surviving
manuscript contains both relations with $N=(2\pi)^2$, so the record does not
fix which comparison selected $N$. Three facts bound the evidential
content. First, the RAR scale with its mass-to-light systematic and either
current value of $H_0$ requires $33.7\le N\le54.8$. Of the standard compact
volumes $\pi$, $2\pi$, $4\pi$, $\pi^2$, $2\pi^2$, $(2\pi)^2$, $8\pi^2$,
$16\pi^2$, $(2\pi)^3$, and the unit three- and four-balls, only $(2\pi)^2$
lies in that window; the menu is assembled after the comparison and carries no probability measure, and the conditional probability below does not depend on it. Second, at $N=(2\pi)^2$ the alternative abundance maps
$\Omega_m/\Omega_b=N/g$, $\Omega_c/\Omega_b=N/g+1$, and
$\Omega_c/\Omega_b=(1-\eta_*)N/g$, evaluated in the full-release linear-depletion benchmark, give pulls of $-16\sigma$, $+15\sigma$, and
$-3.1\sigma$; the release model selects $\Omega_c/\Omega_b=N/g$, since the
CMB samples the committed stock before release. Third, a normalization drawn
log-uniformly across the $a_0$ window reproduces the abundance within one
standard deviation with probability $0.05$. If the acceleration scale selected
$N$, the abundance is a test at that level, and with a fourfold penalty for the
map the Bayes factor lies between $8$ and $22$ for prior ranges of the ratio
from $[3,10]$ to $[1,30]$. If the abundance selected $N$, the acceleration
scale tests it only to its $20\%$ systematic, and the Bayes factor is of order
unity. Eleven numbers native to the construction lie inside the $a_0$ window,
among them $122/3$, $\sqrt{1680}$, $42$, and $48$; only $(2\pi)^2$ reproduces
the abundance within one standard deviation, and $122/3$ lies at $+1.8\sigma$.
The weight of the abundance comparison therefore rests on measuring the
lever in apparent-horizon radii with coefficient one sharing entropy.
Measured in apparent-horizon diameters, circumferences, or reduced radii
$c/2\pi H$, the same coefficient gives $N=2\pi^2$, $2\pi$, or $(2\pi)^3$,
all outside the window, so the menu above applies to the lever unit. Section~15 obtains the coefficient as the
register enthalpy over a one-nat oriented record read at the Unruh temperature, with the arm $c/H$; the
Unruh reading and the arm remain premises. A
derivation of that reading from the cell action, or an independent measurement of $\epsilon$ in resolved cluster maps
(Section~18), would remove the dependence on the order of comparison.

\section*{Appendix M: Growth and acoustic response of saturated capacity}

The microwave background requires a gravitating component that carries the growing mode of the baryon perturbations while rejecting their acoustic oscillation. The free-fall and first-order response clocks below fail the joint growth and acoustic requirements. The constrained carrier of Section~20 supplies a conserved density for the growing mode.

\paragraph{Growth-limited kernels.} An auxiliary phase-weight model supplies a diagnostic of a proposed free-fall response clock; its weights do not replace the transported mass inventory of Section~18. For development dynamics $dW/dt=(W_{\max}-W)/\tau$ with $\tau=\beta\,t_{\rm ff}(\rho_{\rm local})$, the measured cluster radial decline of the source weight fixes $\beta\simeq15.6$, while saturation of the cosmological bath by recombination requires $\beta\lesssim0.56$: the window is empty by a factor of order thirty. The radial ratio between $R_{500}$ and $R_{200}$, a $\beta$-independent prediction of the free-fall clock, agrees with the data at the several-percent level, so the exclusion is of the absolute clock, not of the radial structure.

\paragraph{Lagged response.} A first-order lag transmits a fraction $T=[1+(\omega\tau)^2]^{-1/2}$ of an oscillation at frequency $\omega$. The transmitted component acts as additional effective baryon loading $(1/\epsilon-1)\,T$ in the acoustic driving; Einstein--Boltzmann computation shows a $0.3\%$ temperature-spectrum tolerance bounds $T\le2\times10^{-3}$, requiring $\beta\gtrsim37$ against the development requirement $\beta\lesssim0.56$: empty by a factor of order sixty-six. These bounds exclude the displayed first-order clock. Higher-order kernels require separate growth, transient-response, and acoustic-transfer tests across the relevant modes.

\paragraph{Sound-speed classification.} For $L=f(X)$, perturbations carry $c_s^2=f_X/(f_X+2Xf_{XX})$. Generic unconstrained continuations away from the saturation surface have nonzero propagation speed, while plateau approaches $f=f_{\max}-A/X^n$ give $c_s^2=-1/(2n{+}1)$ and are gradient-unstable. Both natural cap readings (stationarity of the canonical momentum and of the energy density) impose $f_X+2Xf_{XX}=0$, the pole of $c_s^2$, and carry no healthy propagating perturbation. The only locus with $c_s^2=0$ and a healthy density response is the multiplier constraint surface $X=1/2$ used in Section~20. Appendix~O shows that carrier loss does not continue this field onto any of the unconstrained scalar branches: the source-label constraint is removed, the metric-visible remainder returns to the ordinary source-associated capacity bookkeeping, and only the homogeneous marked-history projection changes the cosmological term.

\paragraph{Relaxation plus a hard cap cannot carry the secular component.} Within dynamics consisting of relaxation toward a demand together with a hard cap, with the cell weight as the only state variable, no coupling of the weight to the demand carries a secular component of the modulation while rejecting its oscillation: a pinned weight retains no memory of the modulation, and any instantaneous coupling of bounded periodic inputs is itself bounded and periodic. A secular component requires an additional conserved integrating variable. The conserved spatial density of committed cells is that variable, and it is native to the saturated phase.

\section*{Appendix N: Action Reconstruction, Parent-Theory Audit, and Transverse-Sector Target}

This appendix asks whether the ordinary capacity functional is a second scalar action or the gravitational constraint action written in a different variable. The static longitudinal calculation answers that question. The final subsections give the specified leading transverse EFT and isolate its remaining microscopic and nonlinear-covariance requirements.

\subsection*{N.1 One parent metric, one matter coupling}

The minimal covariant parent of the ordinary branch is
\[
I_0[g,\psi]
=I_{\mathrm{EH}}[g]+I_{\mathrm{GHY}}[g]
+I_{\mathrm{matter}}[g,\psi],
\]
\[
I_{\mathrm{EH}}
=\frac{c^3}{16\pi G}
\int_{\mathcal M}d^4x\,\sqrt{-g}\,(R-2\Lambda).
\]
Here and throughout the covariant formulas, $x^0=ct$. Thus $d^4x=dx^0d^3x$ and the Einstein--Hilbert coefficient is $c^3/(16\pi G)$. After the split $dx^0=c\,dt$, the ADM action written with $t$ in seconds carries $c^4/(16\pi G)$. This convention makes the final Newton action an ordinary $\int dt\,L$ action.
All matter, including radiation, couples to the one physical metric. Its stress tensor is
\[
T_{\mu\nu}
=-\frac{2c}{\sqrt{-g}}
\frac{\delta I_{\mathrm{matter}}}{\delta g^{\mu\nu}},
\]
and the equations are
\[
G_{\mu\nu}+\Lambda g_{\mu\nu}
=\frac{8\pi G}{c^4}T_{\mu\nu}.
\]
The Bianchi identity and the matter equations give $\nabla_\mu T^{\mu\nu}=0$. No explicit coupling $S_{\mathrm{ent}}T^\mu{}_\mu$ is present. Such a term would vanish as a source for classical radiation, would have a nontrivial metric variation through $\delta T^\mu{}_\mu/\delta g^{\alpha\beta}$, and would duplicate the matter response already carried by the metric equation.

The capacity variable enters after gauge choice and constraint reduction:
\[
\delta S=\delta S[g,\psi]\big|_{\mathrm{reduced}}.
\]
This is analogous to using the Newtonian potential as a reduced gravitational variable. A reduced variable can have a useful functional without being a new covariant matter field.

\subsection*{N.2 Direct reduction of the Einstein action}

\paragraph{ADM form and boundary terms.}
With the GHY term included, the bulk gravitational action is
\[
I_{\mathrm{ADM}}
=\frac{c^4}{16\pi G}
\int dt\,d^3x\,N\sqrt h
\left({}^{(3)}R+K_{ij}K^{ij}-K^2-2\Lambda\right)
+I_\infty.
\]
This line uses $t=x^0/c$ and is therefore consistent with the covariant $c^3$ normalization above.
For a static scalar perturbation in Newtonian gauge,
\[
N=\sqrt{1+\frac{2\Phi}{c^2}}
=1+\frac{\Phi}{c^2}+O(c^{-4}),
\qquad
N^i=0,
\qquad
h_{ij}=\left(1-\frac{2\Psi}{c^2}\right)\delta_{ij},
\qquad
K_{ij}=0.
\]
The asymptotic term $I_\infty$ subtracts the flat reference contribution and cancels the total divergence generated when the spatial curvature is integrated by parts.

Define $\phi=\Phi/c^2$, $\psi_N=\Psi/c^2$, and write $h_{ij}=e^{2\zeta}\delta_{ij}$, with
\[
\zeta=\frac12\ln(1-2\psi_N)
=-\psi_N-\psi_N^2+O(\psi_N^3).
\]
In three dimensions,
\[
{}^{(3)}R
=e^{-2\zeta}
\left[-4\nabla^2\zeta-2(\nabla\zeta)^2\right].
\]
Before integration by parts, the expansion through second order is
\[
N\sqrt h\,{}^{(3)}R
=4\nabla^2\psi_N
+4\phi\nabla^2\psi_N
+4\psi_N\nabla^2\psi_N
+6(\nabla\psi_N)^2
+O(3).
\]
The first term is a boundary term. Integrating the next two Laplacians by parts and combining all surface pieces with $I_\infty$ gives
\[
N\sqrt h\,{}^{(3)}R
=\frac{2}{c^4}
\left[
(\nabla\Psi)^2-2\nabla\Phi\cdot\nabla\Psi
\right]
+\nabla_iB^i+O(c^{-6}).
\]
The boundary integral of $B^i$ is removed by the reference-subtracted variational prescription. For nonrelativistic matter with $\rho$ the rest-mass density, expansion of $-mc^2\int d\tau$ gives
\[
I_{\mathrm{matter}}^{(1)}
=-\int dt\,d^3x\,\rho\Phi.
\]
Consequently,
\[
\boxed{
I_{0,\mathrm{scal}}^{(2)}[\Phi,\Psi]
=
\int dt\,d^3x
\left\{
\frac{(\nabla\Psi)^2-2\nabla\Phi\cdot\nabla\Psi}{8\pi G}
-\rho\Phi
\right\}.
}
\]

\paragraph{Constraint equations.}
Variation with respect to the lapse potential gives
\[
\delta_\Phi I_{0,\mathrm{scal}}^{(2)}=0
\quad\Longrightarrow\quad
\nabla^2\Psi=4\pi G\rho.
\]
Variation with respect to the scalar part of the spatial metric gives
\[
\delta_\Psi I_{0,\mathrm{scal}}^{(2)}=0
\quad\Longrightarrow\quad
\nabla^2(\Phi-\Psi)=0.
\]
The elimination order is essential: $\Phi$ and $\Psi$ are varied independently, both constraint equations are obtained, and only then is their solution substituted back. Imposing $\Phi=\Psi$ in the two-potential action before variation would erase one constraint and is not an equivalent reduction.
For asymptotically flat isolated sources, the harmonic difference vanishes:
\[
\Phi=\Psi.
\]
Substituting this solution back into the two-potential action yields
\[
\boxed{
I_{\mathrm{Newton}}[\Phi]
=
\int dt\,d^3x
\left[
-\frac{(\nabla\Phi)^2}{8\pi G}
-\rho\Phi
\right].
}
\]
The sign of the reduced kinetic term is not a ghost signal. $\Phi$ is a nondynamical constraint potential, not a propagating scalar oscillator. In the full covariant parent, the lapse and shift impose constraints and the propagating gravitational content is the two tensor polarizations.

\subsection*{N.3 Capacity field redefinition and coefficient identity}

On the renormalized static branch, the capacity functional is
\[
I_{\mathrm{cap}}^{\mathrm{static}}
=
\int dt\,d^3x
\left[
-\frac{\gamma}{2}(\nabla\delta S)^2
+\kappa\rho\,\delta S
\right].
\]
The most general linear local field redefinition compatible with a constant vacuum background is $\delta S=A\Phi$. Matching the signs of the source and the weak-field clock shift fixes $A<0$. The bounded relation $q=1-\delta S/S_\infty=N^2$ at first order fixes
\[
A=-\frac{2S_\infty}{c^2}.
\]
Substitution gives
\[
I_{\mathrm{cap}}^{\mathrm{static}}
=
\int dt\,d^3x
\left[
-\frac{2\gamma S_\infty^2}{c^4}(\nabla\Phi)^2
-\frac{2\kappa S_\infty}{c^2}\rho\Phi
\right].
\]
The equality
\[
G=\frac{c^2\kappa}{8\pi\gamma S_\infty}
\]
then makes the kinetic coefficient and source coefficient agree simultaneously:
\[
\boxed{
I_{\mathrm{cap}}^{\mathrm{static}}
=
\frac{2\kappa S_\infty}{c^2}
I_{\mathrm{Newton}}.
}
\]
The off-shell identity proves that the two expressions are the same classical functional, up to a field-independent normalization, throughout the reduced static field space. Agreement of their Poisson equations alone would not establish this equivalence.
The statement uses matched Dirichlet or asymptotically flat boundary conditions. With a finite boundary, equality also requires mapping the Newton boundary functional through the same linear field redefinition.

\paragraph{Normalization freedom.}
The observable static equation depends only on
\[
\frac{\kappa}{\gamma S_\infty}.
\]
A change of entropy units
\[
\delta S\to a\,\delta S,\qquad
S_\infty\to aS_\infty,\qquad
\gamma\to\frac{\gamma}{a^2},\qquad
\kappa\to\frac{\kappa}{a}
\]
leaves both the action and the observable bridge invariant. Multiplying the entire reduced action by a further constant also leaves its classical equation invariant. The UV calculation can fix these conventions for fluctuation normalization and correlation functions, but the classical Newton constant is fixed only once through the invariant combination. A convenient unit-normalized representative has $Z_S=1$, equivalently
\[
\kappa=\frac{c^2}{2S_\infty},
\qquad
\gamma=\frac{c^4}{16\pi GS_\infty^2},
\]
but the manuscript's canonical UV convention need not be changed to use the equivalence.

\paragraph{Why canonical scalar stress cannot do this job.}
If $\delta S=O(\rho)$ about a constant background, then
\[
T_{\mu\nu}^{(\delta S)}
\sim\partial_\mu\delta S\,\partial_\nu\delta S
=O(\rho^2).
\]
The bridged potential is $O(\rho)$. It therefore cannot be attributed to the canonical stress of a second scalar. In the reconstruction above, the linear response comes from the metric constraint and the apparent scalar functional is its reduced representation.

\subsection*{N.4 Parent-action options and degree-of-freedom audit}

\begingroup
\small
\setlength{\tabcolsep}{4pt}
\renewcommand{\arraystretch}{1.15}
\begin{longtable}{|L{0.18\textwidth}|L{0.20\textwidth}|L{0.13\textwidth}|L{0.35\textwidth}|}
\hline
\textbf{Route} & \textbf{Variables / action} & \textbf{Modes} & \textbf{Audit result} \\
\hline
\endfirsthead
\hline
\textbf{Route} & \textbf{Variables / action} & \textbf{Modes} & \textbf{Audit result} \\
\hline
\endhead
Metric-only / composite capacity
&
$I_0[g,\psi]$; $\delta S$ defined after weak-field constraint reduction and $q_{\rm geo}=(\nabla R)^2$ in spherical symmetry
&
two tensors
&
Passes the weak-field action identity, couples radiation correctly, gives baseline no slip and GR PPN, and passes the spherical exterior test. This is the preferred ordinary branch. A generic local covariant definition of capacity outside the controlled reductions remains open.
\\
\hline
Algebraic ADM multiplier
&
$I_0+\int\sqrt h\,\lambda(N^2-q)$
&
no new mode only because the construction is empty
&
$\delta q$ gives $\lambda=0$ when $q$ has no other bulk term. The relation only renames a foliation-dependent lapse and does not yield the capacity source equation. Excluded as a parent completion.
\\
\hline
Clock-constrained auxiliary capacity
&
$I_0+\int\sqrt{-g}\,\Lambda(qX+1)$, $X=g^{\mu\nu}\partial_\mu\tau\partial_\nu\tau$
&
trivial if $q$ is purely algebraic; extra scalar or zero-pressure mode once dynamics is supplied
&
The $q$ equation again drives $\Lambda$ to zero unless additional $q$ dependence is introduced. Kinetic or higher-derivative terms generically introduce a preferred foliation, mimetic zero-pressure behavior, or a scalar stability problem. Retain only as a control model with a full Hamiltonian analysis.
\\
\hline
Scalar--tensor / disformal
&
$I[g,q]+I_{\rm matter}[\widetilde g,\psi]$, $\widetilde g_{\mu\nu}=C(q)g_{\mu\nu}+D(q)u_\mu u_\nu$
&
two tensors plus at least one scalar
&
Universal matter coupling can be arranged, but fifth forces, gravitational slip, PPN shifts, and radiative stability become new obligations. It is not needed for the ordinary branch and is not selected by the present evidence.
\\
\hline
\end{longtable}
\endgroup

The metric-only route is minimal in a precise sense: it adds no propagating field, uses the full stress tensor, and reproduces both controlled reductions. This does not prove that every substrate sector is metric-only. It proves that the ordinary longitudinal capacity mode must not be counted again as an independent force.

\subsection*{N.5 Exact spherical reduction and the Misner--Sharp first integral}

Let
\[
ds^2=h_{ab}(x)dx^adx^b+R^2(x)d\Omega^2.
\]
We use $x^0=ct$, so $d^2x=dx^0dr$. If the orbit-space integral is instead written as $dt\,dr$, its prefactor is $c^4/(4G)$.
The curvature identity is
\[
{}^{(4)}R
=
{}^{(2)}R
+\frac{2}{R^2}
\left[1-(\nabla R)^2-2R\Box R\right].
\]
After the angular integral,
\[
I_{\mathrm{EH}}
=\frac{c^3}{4G}
\int d^2x\,\sqrt{-h}
\left[
R^2{}^{(2)}R+2\{1-(\nabla R)^2\}-4R\Box R
\right].
\]
Integrating the last term by parts together with the reduced GHY term gives
\[
\boxed{
I_{\mathrm{sph}}
=\frac{c^3}{4G}
\int d^2x\,\sqrt{-h}
\left[
R^2{}^{(2)}R+2(\nabla R)^2+2
\right]
+I_{\mathrm{matter}}^{(2)}+I_\partial^{(2)}.
}
\]
The reduced boundary term is not left implicit. For a product boundary $\partial\mathcal M_4=\partial\mathcal M_2\times S^2$,
\[
K^{(4)}=K^{(1)}+\frac{2}{R}n^a\nabla_aR.
\]
The second term cancels the boundary contribution generated when $-4R\Box R$ is integrated by parts. With the standard orientation sign $\varepsilon$, the surviving term is
\[
I_\partial^{(2)}
=\varepsilon\,\frac{c^3}{2G}
\int_{\partial\mathcal M_2}dy\,
\sqrt{|\gamma_{(1)}|}\,R^2K^{(1)}
+I_{\mathrm{joint}}+I_{\mathrm{ref}}.
\]
The vacuum Euler--Lagrange equations are
\[
R\,{}^{(2)}R-2\Box R=0,
\]
\[
2R\left(h_{ab}\Box R-\nabla_a\nabla_bR\right)
+h_{ab}\left[(\nabla R)^2-1\right]=0.
\]
Define
\[
q_{\mathrm{geo}}\equiv(\nabla R)^2,
\qquad
M_{\mathrm{MS}}
\equiv\frac{c^2R}{2G}(1-q_{\mathrm{geo}}).
\]
Contracting and differentiating the metric equation, with the $R$ equation used to remove the curvature term, gives
\[
\nabla_aM_{\mathrm{MS}}=0
\]
in vacuum. Hence
\[
q_{\mathrm{geo}}
=1-\frac{2GM}{c^2R}.
\]
In the static areal gauge $R=r$, asymptotic flatness fixes the remaining normalization and returns
\[
ds^2
=-\left(1-\frac{2GM}{c^2r}\right)c^2dt^2
+\left(1-\frac{2GM}{c^2r}\right)^{-1}dr^2+r^2d\Omega^2,
\]
\[
q_{\mathrm{geo}}=N^2.
\]
Thus the radial $q_{\mathrm{geo}}$ is a composite spherical scalar with a vacuum first integral. Its equality to the clock quantity uses the normalized static vacuum solution. The result adds no auxiliary scalar and is nonperturbative within this symmetry sector.

\subsection*{N.6 Clock capacity, radial geometry, and the nonlinear lapse}
\label{aqg:nonlinearcapacity}

A static clock measures the norm of the timelike Killing field with
its normalization fixed at the reference boundary. Write
\[
ds^2=-N^2c^2dt^2+h_{ij}dx^idx^j,\qquad
q_\tau=N^2,\qquad s_\tau=S_\infty(1-N^2).
\]
The reference normalization \(N\to1\) is part of this observable.
In spherical symmetry the radial metric supplies a separate scalar,
\[
q_R=(\nabla R)^2=1-\frac{2Gm(R)}{c^2R}.
\]
For a connected static vacuum exterior with the stated normalization,
\(q_\tau=q_R\). Matter generally separates them. With mass density
\(\rho\), radial pressure \(p_r\), and areal coordinate \(r\),
\[
\frac{d}{dr}\log\frac{q_\tau}{q_R}
=\frac{8\pi Gr}{c^2q_R}\left(\rho+\frac{p_r}{c^2}\right).
\]
Inside a supported spherical shell, \(q_R=1\) while the clock lapse can
be a constant smaller than one. The interior is locally flat and
force free, yet its clock is redshifted relative to the exterior
reference. A local curvature scalar cannot encode this boundary
comparison.

Let \(n^\mu\) be the unit static normal, let
\(\mathcal E=T_{\mu\nu}n^\mu n^\nu\) be the energy density, and put
\(\mathcal S=h^{ij}T_{ij}\). At zero cosmological constant the static
Einstein equation gives
\[
D^2N=\frac{4\pi G}{c^4}N(\mathcal E+\mathcal S).
\]
Consequently
\[
D^2q_\tau=\frac{8\pi G}{c^4}q_\tau(\mathcal E+\mathcal S)
 +\frac{|Dq_\tau|^2}{2q_\tau},
\]
and the exact static capacity equation on \(N>0\) is
\[
\boxed{\;
D^2s_\tau=
-\frac{8\pi GS_\infty}{c^4}
 \left(1-\frac{s_\tau}{S_\infty}\right)(\mathcal E+\mathcal S)
-\frac{|Ds_\tau|^2}{2(S_\infty-s_\tau)}.\;}
\]
The spatial metric, matter equations, and equilibrium conditions
complete this system. The lapse equation alone does not determine them.
For nonrelativistic matter and a weak field,
\[
\nabla^2s_\tau=-\frac{8\pi GS_\infty}{c^2}\rho,\qquad
s_\tau=-\frac{2S_\infty}{c^2}\Phi.
\]
The corresponding physical coefficients are
\[
\kappa_{\rm phys}=\frac{c^2}{2S_\infty},\qquad
\gamma_{\rm phys}=\frac{c^4}{16\pi GS_\infty^2}.
\]
These identities express the capacity normalization for a supplied
Einstein coefficient; they do not independently determine \(G\).

In a dynamical ADM slicing, the lapse remains a slicing variable.
For \(c=1\), \(K_{ij}=-\tfrac12\mathcal L_nh_{ij}\), and \(q_N=N^2\),
the trace evolution equation gives
\[
\begin{split}
D^2q_N={}&2q_N\left[K_{ij}K^{ij}
 +4\pi G(\mathcal E+\mathcal S)-\Lambda\right]\\
&-2N(\partial_t-\mathcal L_\beta)K
 +\frac{|Dq_N|^2}{2q_N}.
\end{split}
\]
Even maximal slicing leaves a choice of boundary and slicing data.
This equation is therefore a dynamical lapse relation, not a general
foliation-independent definition of clock capacity.

A general capacity observable can depend on a relational region,
its algebra, the quantum state, a clock, and reference data:
\[
q_{\rm cap}(x)=
\mathcal Q[g_{\mu\nu},\rho_{\rm q};
 \mathcal D_x,\mathcal A(\mathcal D_x),\mathcal C,\mathcal B].
\]
Such dependence can encode the shell redshift without new local
Cauchy data. Its static clock restriction must reproduce \(q_\tau\);
\(q_R\) remains a distinct radial observable. The two restrictions
agree in their common static vacuum domain.

Relative entropy constrains a possible informational realization.
For a smooth perturbation of a faithful reference state,
\[
D(\rho_\epsilon\Vert\rho_0)
=\Delta\langle K_0\rangle-\Delta S=O(\epsilon^2).
\]
A differentiable function of this quantity alone has no generic
linear gravitational response at equilibrium. An informational
construction must retain appropriate modular-energy, geometric, or
reference data. This first-variation restriction does not invalidate
the exact hard-exclusion identities of Appendix~P, whose state change
is outside that smooth full-support perturbation.

Bianconi's geometric-relative-entropy action supplies a volumetric
comparison: a trace logarithm has an Einstein--Hilbert low-curvature
term, and its Legendre transform gives the convex mismatch functional
of Section~25 \citep{Bianconi2026}. Its independently dynamical response
operator does not supply the composite observable above. The ordinary
capacity branch requires a causal or explicitly quasilocal observable,
relational reference data, the full stress-tensor response, and
compatibility with the metric Ward identity without additional
propagating gravitational data. A bounded continuation near saturation
also needs a constitutive rule; relative entropy itself is unbounded.

\subsection*{N.7 The $q_{\mathrm{geo}}=0$ variational boundary}

The exterior action should be defined first on a stretched timelike surface
\[
\mathcal B_\delta=\{q_{\mathrm{geo}}=\delta\},
\qquad \delta>0.
\]
For Dirichlet metric data, include
\[
I_{\mathrm{GHY}}[\mathcal B_\delta]
=\varepsilon\,\frac{c^3}{8\pi G}
\int_{\mathcal B_\delta}d^3y\,\sqrt{|\gamma|}\,K,
\]
the reference term at infinity, and the appropriate joints. Only after the variation is defined should the null limit $\delta\to0^+$ be taken, using the null-boundary form and a fixed normalization of the generators. This supplies a well-posed exterior Einstein variational principle.

Three logically different claims must not be conflated:
\begin{enumerate}
\item The reduced metric equations produce a marginal sphere where $q_{\mathrm{geo}}=0$. This is derived.
\item The substrate capacity fraction equals $q_{\mathrm{geo}}$ up to that surface. This is a constitutive identification, strongly matched in static spherical vacuum.
\item The physical spacetime ends there and does not realize $q_{\mathrm{geo}}<0$. This is the bounded-domain postulate and requires new boundary dynamics.
\end{enumerate}

The standard gravitational boundary term fixes neither a channel Hamiltonian nor a reflectivity. The missing microscopic functional may be denoted
\[
\Gamma_{\partial q}
[\sigma_{AB},\zeta_{\partial};\rho_{\partial}],
\]
where $\zeta_{\partial}$ are boundary channel variables and $\rho_{\partial}$ is their state. It must determine:
\begin{itemize}
\item whether the GR ingoing condition $\mathcal R(\omega)=0$ is recovered;
\item whether a stretched layer $q_{\mathrm{geo}}=\epsilon$ is dynamically selected;
\item the relaxation spectrum and any echo kernel;
\item the shape variation and boundary stress balance;
\item consistency of the derived $a_{\rm ch}=4\ln2\,L_*^2$ area normalization under any future microscopic boundary Hamiltonian (the kinematic count itself is fixed in F.5).
\end{itemize}
Until this functional is derived, the Schwarzschild exterior is a result, while interior excision and boundary phenomenology are hypotheses.

\subsection*{N.8 Carrier-resolved transverse EFT and metric contact}

This subsection gives a self-contained action audit of the galactic branch. The leading EFT contains four distinct ingredients: a source-space nesting rule, a thermal occupation contact, a conservative auxiliary-field action, and a no-slip metric contact. The first three determine which stress enters the galactic response and how its static potential is computed. The metric contact determines lensing. The retained record algebra, thermal matching, and derivative order specify the leading EFT. The constructions below establish its source accounting and linear response; nonlinear record evolution and the values of its response rates require further matching.


\paragraph{Carrier projector and nesting.}
Let $e_P$ be a central sharp selector in the decohered record algebra and let $E_P$ be a conditional expectation onto a retained subalgebra containing $e_P$. The expectation preserves the source functional $\mathcal I$, so $\mathcal I(E_PX)=\mathcal I(X)$ on its sector. Set
\[
\Pi_PX=E_P(e_PX)=e_PE_PX.
\]
For disjoint record selectors, the bimodule property gives
\[
\Pi_P^2X=e_PE_P(e_PE_PX)=e_PE_PX,\qquad
\Pi_P\Pi_QX=E_P(e_Pe_QE_QX)=0\quad(P\ne Q).
\]
An active antichain partitions the modeled records, $\sum_Pe_P=1$, and therefore preserves their total source:
\[
\sum_P\mathcal I(\Pi_PX)=\sum_P\mathcal I(e_PX)=\mathcal I(X).
\]
These statements have a finite realization as blockwise conditional averages with positive reference weights. Each record belongs to exactly one block, and averaging a block twice leaves its first average unchanged.

For nested algebras on a common parent record space, the conditional expectations use the same reference measure and obey $E_PE_{\rm f}=E_P$. Child selectors obey $\sum_{A\prec P}e_A=e_P$. If their retained blocks refine the parent's retained blocks, $E_{\rm f}$ can be assembled from the child expectations, giving
\[
\Pi_PX=E_P\!\left(\sum_{A\prec P}\Pi_AX\right).
\]
The tower identity requires this refinement condition in addition to the antichain rule.

\paragraph{Spatial reconstruction and overlap.}
Write $T_P^{\rm cg}=\mathcal C_\ell\Pi_PT$ for the reconstructed source. In the static scalar sector a positive normalized kernel realizes the map:
\[
\rho_P(\mathbf x)=\int d^3y\,K_\ell(\mathbf x,\mathbf y)
(\Pi_P\rho)(\mathbf y),\qquad
K_\ell\ge0,\qquad \int d^3x\,K_\ell(\mathbf x,\mathbf y)=1.
\]
It follows directly that $\int\rho_P\,d^3x=\mathcal I(\Pi_P\rho)$ and that the sum of reconstructed masses equals the modeled mass. At a common coarse resolution the parent source is assembled by the preceding tower identity. Changing the reconstruction scale requires an explicit resolution map with the same source normalization. Tensor reconstruction is defined relative to the carrier clock and must transform covariantly with the record variables in the metric variation below.

The interpolation weights in the constitutive action are separate spatial functions. Positive profiles $s_P(\mathbf x)$ define
\[
\chi_P=\frac{s_P}{\sum_Qs_Q},\qquad
\sum_P\chi_P=1
\]
on their response domain. Two profiles can have $\chi_P=\chi_Q=1/2$ while their record selectors remain orthogonal. The reconstructed density already contains its full carrier normalization; the spatial weight multiplies the constitutive energy. The smooth interface profiles specified below implement this partition of unity.

The Einstein action uses the full $T_{\mu\nu}$. The transverse functional uses the joint reconstructed source and sums over active carriers:
\[
\Gamma_\perp^{\rm nest}
=\sum_{P\in\mathfrak A}w_P
\Gamma_{\perp,w_P}[g_\pm;T^{\mathfrak A}_\pm,\chi_P,\mathcal R_P],
\qquad T^{\mathfrak A}=\sum_{Q\in\mathfrak A}T_Q^{\rm cg}.
\]
Thus source assignment, spatial overlap, and nonlinear response enter as distinct operations. Translation-covariant reconstruction and interpolation profiles preserve the common-translation symmetry of the static action. On-shell force balance follows when all explicit source, support, stiffness, and core variations are included.

\paragraph{Positive reconstruction and its adjoint.}
At finite resolution use cell masses $m_\alpha$, with volume factors included in conversion to densities. A reconstruction matrix $C$ obeying
\[
C_{i\alpha}\ge0,\qquad \sum_iC_{i\alpha}=1
\]
maps positive fine sources to positive coarse sources and preserves their total mass. A coarse probe $f$ couples through
\[
f^{\mathsf T}Cm=(C^{\mathsf T}f)^{\mathsf T}m.
\]
Thus the same map fixes source deposition and force interpolation. Composing normalized positive reconstructions preserves these statements, including overlapping spatial profiles. The record projector, reconstruction, and constitutive weights remain distinct: $\Pi_P$ is idempotent, $C$ preserves source mass, and $\sum_P\chi_P=1$ weights the constitutive energy. This provides a common finite source map for the nesting identity and reciprocal variation.

\paragraph{Tetrahedral kinetic normalization.}
Let $q_i(x)$, $i=1,\ldots,4$, denote fluctuations on the four equivalent tetrahedral face directions $\mathbf n_i$, normalized by
\[
\mathbf n_i\!\cdot\!\mathbf n_j=
\begin{cases}
1,&i=j,\\
-1/3,&i\ne j,
\end{cases}
\qquad
\sum_i\mathbf n_i=0.
\]
For the normalization argument use the Euclidean quadratic term
\[
\Gamma_{E,\mathrm{kin}}=\frac{Z_0}{2}\int d^4x_E\sum_i
\partial_\mu q_i\partial_\mu q_i.
\]
The $q_i$ are dimensionless occupation fluctuations, and $Z_0$ has units of action per length squared. Lorentzian response is obtained by analytic continuation to the signature of Appendix~A.
The four-dimensional permutation representation splits into a uniform singlet and a three-dimensional sum-zero sector. On the latter define
\[
A^a=\frac{\sqrt3}{2}\sum_i n_i^{\,a}q_i,
\qquad
\sum_iq_i=0.
\]
Using $\mathbf n_i\cdot\mathbf n_j=4\delta_{ij}/3-1/3$ gives the exact norm identity
\[
\mathbf A^2=\sum_iq_i^2,
\qquad
\Gamma_{E,\mathrm{kin}}=\frac{Z_0}{2}\int d^4x_E\,
\partial_\mu\mathbf A\cdot\partial_\mu\mathbf A.
\]
The same result follows from Schur's lemma: the tetrahedral three-dimensional representation is irreducible, so a symmetry-invariant quadratic operator is proportional to the identity on it. A slowly varying baryonic field selects $\mathbf e_L=\nabla\Phi_{\rm bar}/|\nabla\Phi_{\rm bar}|$ and permits
\[
\mathbf A=X\mathbf e_L+\mathbf Y,
\qquad
\mathbf e_L\cdot\mathbf Y=0.
\]
For a locally fixed $\mathbf e_L$, the longitudinal mode and both transverse components inherit the same coefficient,
\[
Z_L=Z_T=Z_0.
\]
Writing $Y_1=R\cos\vartheta$ and $Y_2=R\sin\vartheta$ gives
\[
(\partial Y_1)^2+(\partial Y_2)^2
=(\partial R)^2+R^2(\partial\vartheta)^2,
\]
so the radial occupation mode that mixes with $X$ also carries $Z_0$.
The real pair $(Y_1,Y_2)$ has the single polar angle $\vartheta$. The response cell of Section~15 is the tilt mode of this frame, and for a lever in apparent-horizon radii its factor $(2\pi)^2$ comes from the horizon temperature (paragraph ``Thermal contact'' below).
Spatial variation of $\mathbf e_L$ produces derivative connection terms suppressed by the curvature scale of the source direction. Higher-derivative and background-dependent operators can split the coefficients; the equality is a leading long-wavelength result.

Allowing the frame to vary makes the leading omitted operator explicit. Write
\[
\mathbf A=X\hat{\mathbf e}+Y_a\hat{\mathbf e}_a,
\qquad \hat{\mathbf e}^{,2}=1.
\]
Then
\[
(\partial\mathbf A)^2=(\partial X)^2+(\partial Y_a)^2
+X^2(\partial\hat{\mathbf e})^2
+Y_aY_b\,\partial\hat{\mathbf e}_a\!\cdot\!\partial\hat{\mathbf e}_b
+\text{connection cross terms}.
\]
Thus the same $Z_0$ that fixes $Z_L=Z_T$ supplies a direction stiffness $Z_\theta=Z_0X^2$, plus a nonnegative occupation-weighted frame term. The selected infrared identification maps the longitudinal amplitude to the carrier record potential, the potential of the carrier's response stock, $X_P\leftrightarrow\Phi_{R,P}$ (Appendix~O.3), with $\Phi_{R,P}\to0$ at the capacity vacuum. For an isolated carrier in equilibrium $\Phi_{R,P}$ is a fixed functional of $\Phi_{b,P}$, which serves as a proxy. This fixes the displayed stiffness profile but leaves the potential zero for an embedded carrier, the occupation-weighted frame term, and the connection cross terms as named completion data.

\paragraph{One-capacity-invariant potential.}
After the relative sector of Appendix~H is gapped, define the longitudinal and transverse occupations $N_L=\sigma^\dagger P_L\sigma$ and $N_T=\sigma^\dagger P_T\sigma$. The minimal infrared premise is that matter sees one total-capacity invariant,
\[
\Delta C=w_LN_L+w_TN_T-N_*.
\]
This premise is stronger than tetrahedral symmetry. The general local quartic potential contains
\[
V_{\rm gen}
=\mu_LN_L+\mu_TN_T
+\frac{\lambda_L}{2}N_L^2
+\lambda_{LT}N_LN_T
+\frac{\lambda_T}{2}N_T^2+\cdots .
\]
Its quartic coupling matrix has rank at most one when
\[
\lambda_{LT}^2=\lambda_L\lambda_T.
\]
It has rank one only if the matrix is nonzero. A stable one-invariant potential further requires the matrix to be positive semidefinite, so $\lambda_L,\lambda_T\ge0$ and not both vanish. Under those conditions $\lambda_L=u_Cw_L^2$, $\lambda_T=u_Cw_T^2$, and $\lambda_{LT}=u_Cw_Lw_T$ with $u_C>0$ and suitable signed weights. The linear terms must align with the same weights. A thermodynamic capacity-pressure variable makes the stable sign explicit through
\[
V_{\rm eff}(\Delta C)=\sup_{\lambda_C}
\left[\lambda_C\Delta C-\frac{\chi_C}{2}\lambda_C^2\right].
\]
The stationary pressure $\lambda_C=\Delta C/\chi_C$ gives $V_{\rm eff}=(\Delta C)^2/(2\chi_C)$. A path-integral Hubbard--Stratonovich representation requires the contour and Lorentzian/Euclidean sign to be chosen consistently; the equation above is the thermodynamic Legendre form. More generally $V_{\rm eff}=F(\Delta C)$ with $F'(0)=0$ and $F''(0)>0$.

Near a transverse background, $N_L=N_{L0}+\zeta_LX+\cdots$ and $N_T=R^2=Y_1^2+Y_2^2$, so
\[
\Delta C=\alpha_C+g_CX+\beta_C R^2.
\]
For fields $\varphi_i=(X,R)$ on the stationary surface $\Delta C=0$,
\[
H_{ij}=\left.\frac{\partial^2V_{\rm eff}}
{\partial\varphi_i\partial\varphi_j}\right|_{\Delta C=0}
=F''(0)\,\partial_i\Delta C\,\partial_j\Delta C.
\]
For $F^{\prime\prime}(0)>0$, this outer product has rank one wherever $\nabla\Delta C\ne0$ and rank zero where that gradient vanishes. Explicitly,
\[
H=F''(0)
\begin{pmatrix}
g_C^2 & 2g_C\beta_C R_0\\
2g_C\beta_C R_0 & 4\beta_C^2R_0^2
\end{pmatrix},
\qquad
C_\times^2=A_LA_T.
\]
For $A_T>0$, minimizing at fixed $\delta X$ gives $\delta R_*=-C_\times\delta X/A_T$ and $|C_\times|/A_T=\sqrt{A_L/A_T}$; the signed ratio $C_\times/A_T=\sqrt{A_L/A_T}$ holds when $C_\times\ge0$. These are static thermodynamic identities. The Hessian has a tangent zero eigenvalue and an orthogonal eigenvalue $A_L+A_T$. It cannot supply two independent mode energies, and it is not used as the propagator that generates the RAR. A metric vertex orthogonal to the soft tangent also cannot obtain the observed response by exchanging that tangent mode. The rank-one result is therefore retained only as a one-capacity consistency check.

This result is protected at the effective level when all local corrections remain functions of the same invariant $\Delta C$. To state the decoupling condition explicitly, define $D=w_TN_L-w_LN_T$ and
\[
V(C,D)=F(C)+\frac{\Delta_D}{2}D^2+\eta_DCD^2+J_DD+\cdots,
\qquad C\equiv\Delta C.
\]
The stationary solution is $D=-J_D/(\Delta_D+2\eta_DC)+\cdots$. For large $\Delta_D$, integrating out $D$ changes $F(C)$ only through constants and powers of $\Delta_D^{-1}$. This is an effective decoupling statement, not proof that the bare vertex is rank one. When $A_T>0$, a standalone $m_X^2X^2/2$ with $m_X^2>0$ outside $F(\Delta C)$ raises the Hessian rank; a Yukawa interpretation also requires the kinetic operator; the $X^2$ curvature tied to the entries of the rank-one outer product does not.

The susceptibility is not fixed by the normalized 1680-state internal distribution. It depends on fluctuations of total condensate occupation,
\[
\chi_C=\langle \widehat C^2\rangle-\langle\widehat C\rangle^2,
\qquad
\widehat C=w_L\widehat N_L+w_T\widehat N_T,
\]
in a specified GFT state. In an energy-normalized convention the corresponding susceptibility includes $\beta_H\equiv1/(k_BT_H)$. For a single mode $\widehat C=c_0\widehat N$, a thermal state gives $\chi_C=c_0^2\bar N(1+\bar N)$, while a coherent condensate gives $\chi_C=c_0^2\bar N$. The closure distribution fixes neither state nor $\bar N$. The susceptibility controls the massive direction and nonlinear corrections but cancels from the rank-one ratio.

\paragraph{Thermal contact.}
The temporal part of the frame term $Z_0X^2(\partial\hat{\mathbf e})^2$ gives the local director tilt the frequency $c/\xi_\theta$, far above the horizon rate for every galactic coherence length, so that tilt carries no occupation (Section~15). The response cell is the horizon-scale tilt mode of the frame, with the source field as its restoring stiffness. With the inertia lever $\ell=g_{\mathrm{share,eff}}c/H$ the tilt mode has $\omega_P=\sqrt{|\boldsymbol\nabla\Phi_{b,P}|/\ell}$, and thermalized in the apparent-horizon bath at $k_BT_H=\hbar H/2\pi$ it has
\[
x_P=\frac{2\pi\omega_P}{H}=\sqrt{\frac{|\boldsymbol\nabla\Phi_{b,P}|}{a_0}},
\qquad
a_0=\frac{g_{\mathrm{share,eff}}}{4\pi^2}\,cH_0,
\qquad
H_{{\rm cell},P}=\hbar\omega_P:N_P:=k_BT_Hx_P:N_P:.
\]
The corresponding equilibrium free energy is
\[
F_{{\rm th},P}=k_BT_H\ln(1-e^{-x_P}),
\qquad
\frac{\partial F_{{\rm th},P}}{\partial E_P}
=\frac{1}{e^{x_P}-1}=n_B(x_P).
\]
The tilt mode fixes $E_P/(k_BT_H)=x_P$ and $a_0$ together. The identification of the cell with this mode, its horizon thermalization, and the lever are the matching inputs of this effective contact. The lever coefficient $g_{\mathrm{share,eff}}$ is the ratio of the register enthalpy to a one-nat oriented record read at the Unruh temperature of the projected baryonic field (Section~15); that reading and the other inputs each require a microscopic GFT derivation.

\paragraph{Static auxiliary action.}
Project the density and solve its Newtonian auxiliary field,
\[
\rho_P=\frac{u_\mu u_\nu}{c^2}T_P^{{\rm cg},\mu\nu},
\qquad
\nabla^2\Phi_{b,P}=4\pi G\rho_P.
\]
With $z_P=|\boldsymbol\nabla\Phi_{b,P}|^2/a_0^2$ and $x_P=z_P^{1/4}$, define
\[
\mathcal Q_\perp(z)=\int_0^z\frac{ds}{e^{s^{1/4}}-1},
\qquad
\mathcal Q(z)=z+\mathcal Q_\perp(z).
\]
Then $\mathcal Q'(z)=1+n_B(z^{1/4})$. The carrier action is
\[
I_{Q,P}=-\int dt\,d^3x\left\{
\frac{2\boldsymbol\nabla\Phi_P\cdot\boldsymbol\nabla\Phi_{b,P}
-a_0^2\mathcal Q(z_P)}{8\pi G}
+\rho_P\Phi_P
\right\},
\]
and the transverse contribution is $I_{\perp,P}=I_{Q,P}-I_{N,P}$, with $I_{N,P}$ obtained by setting $\mathcal Q(z)=z$. Its two Euler--Lagrange equations are
\[
\nabla^2\Phi_{b,P}=4\pi G\rho_P,
\qquad
\nabla^2\Phi_P
=\boldsymbol\nabla\!\cdot\!
\left[\mathcal Q'(z_P)\boldsymbol\nabla\Phi_{b,P}\right].
\]
Consequently, for $\phi_{\perp,P}=\Phi_P-\Phi_{b,P}$,
\[
\boxed{
\nabla^2\phi_{\perp,P}
=\boldsymbol\nabla\!\cdot\!
\left[n_B(x_P)\boldsymbol\nabla\Phi_{b,P}\right].
}
\]
This is the minimal conservative nonspherical completion in the QUMOND class \citep{Milgrom2010}. Spherical sources recover the exponential RAR. Nonspherical sources require the displayed Poisson solve and generally have geometry-dependent curl corrections.

\paragraph{Direction resummation and director action.}
At a thin-disk midplane the source direction $\mathbf e_L$ reverses through $\pi$ across roughly two vertical scale heights, and the local divergence contains a narrow phantom sheet with surface density of order $n_B\Sigma_b$. The resummation replaces the slaved direction by a nematic director $\hat{\mathbf e}_P$ and uses
\[
\widetilde z_P=\frac{(\hat{\mathbf e}_P\!\cdot\!\boldsymbol\nabla\Phi_{b,P})^2}{a_0^2},
\qquad
\widetilde x_P=\widetilde z_P^{1/4},
\]
\[
\boxed{
\nabla^2\phi_{\perp,P}
=\boldsymbol\nabla\!\cdot\!\left[
n_B(\widetilde x_P)
(\hat{\mathbf e}_P\!\cdot\!\boldsymbol\nabla\Phi_{b,P})
\hat{\mathbf e}_P\right].}
\]
Write $\mathcal Q_\perp=\mathcal Q-z$, $Z_{\theta,P}=2c\,X_P^2$ with $X_P$ the record potential, and take the director density $\mathcal S[\hat{\mathbf e}]=|\partial\hat{\mathbf e}|^2+s[\hat{\mathbf e}]$ of Section~15. The auxiliary action is
\[
\begin{aligned}
I_{{\rm dir},P}=-\int dt\,d^3x\Bigg\{
&\frac{1}{8\pi G}\Big[
2\boldsymbol\nabla\Phi_P\!\cdot\!\boldsymbol\nabla\Phi_{b,P}
-|\boldsymbol\nabla\Phi_{b,P}|^2
-a_0^2\mathcal Q_\perp(\widetilde z_P)\\
&\qquad
+\frac{Z_{\theta,P}}{2}\,\mathcal S[\hat{\mathbf e}_P]
+\lambda_P(\hat{\mathbf e}_P^2-1)
\Big]+\rho_P\Phi_P\Bigg\}.
\end{aligned}
\]
The saddle-splay part integrates by parts to
\[
\int\frac{Z_{\theta,P}}{2}\,s[\hat{\mathbf e}_P]
=-\frac12\int\partial_iZ_{\theta,P}\left[(\hat{\mathbf e}_P\!\cdot\!\boldsymbol\nabla)\hat e_P^i-\hat e_P^i\,\boldsymbol\nabla\!\cdot\!\hat{\mathbf e}_P\right]+\text{boundary},
\]
so it acts only through gradients of the stiffness. At constant $Z_\theta$ the action reduces to the one-constant tetrahedral term of the kinetic normalization above.

Variation with respect to $\Phi_P$ gives $\nabla^2\Phi_{b,P}=4\pi G\rho_P$. Variation with respect to $\Phi_{b,P}$ collects $2\boldsymbol\nabla\!\cdot\mathbf J_{\perp,P}$ from $\mathcal Q_\perp(\widetilde z_P)$. With the baryonic proxy, $X_P=\Phi_{b,P}$, the stiffness contributes $2c\,X_P\mathcal S$ through $\partial Z_{\theta,P}/\partial X_P=4c\,X_P$, and the variation gives the boxed equation below directly. With the stock amplitude, $X_P=\Phi_{R,P}$ is sourced by the stock, $\nabla^2\Phi_{R,P}=4\pi G\rho_{R,P}$, and is held fixed in this auxiliary action, whose $\Phi_{b,P}$ variation gives $\nabla^2\phi_{\perp,P}=\boldsymbol\nabla\!\cdot\mathbf J_{\perp,P}$. The stiffness then enters through the stock's selection functional of O.3. In its Gaussian limit the stock descends
\[
F_R[\rho_R]=\frac{1}{8\pi G}\int|\boldsymbol\nabla\Phi_R-\mathbf J_\perp|^2+\frac{c}{8\pi G}\int\Phi_R^2\,\mathcal S[\hat{\mathbf e}],
\]
and $\rho_R\mapsto\Phi_R$ is invertible, so stationarity in $\rho_R$ is stationarity in $\Phi_R$:
\[
\nabla^2\Phi_R=\boldsymbol\nabla\!\cdot\mathbf J_\perp+c\,\Phi_R\,\mathcal S[\hat{\mathbf e}].
\]
In the static limit the stock potential is $\phi_\perp$, and both readings give
\[
\boxed{
\nabla^2\phi_{\perp,P}
=\boldsymbol\nabla\!\cdot\mathbf J_{\perp,P}
+c\,X_P\,\mathcal S[\hat{\mathbf e}_P].}
\]
With the stock amplitude the added term is stock density $c\Phi_R\mathcal S/4\pi G$, which gravitates and lenses as ordinary matter. The equation holds in the Gaussian limit of the selection functional where the positivity and budget constraints are inactive, and it depends on the potential zero of $\Phi_R$, which is completion data.

Director variation gives
\[
\begin{aligned}
\left(\mathbf 1-\hat{\mathbf e}_P\hat{\mathbf e}_P^{\mathsf T}\right)\Big\{&
\boldsymbol\nabla\!\cdot\!\left(Z_{\theta,P}\boldsymbol\nabla\hat{\mathbf e}_P\right)
-(\boldsymbol\nabla\!\cdot\!\hat{\mathbf e}_P)\boldsymbol\nabla Z_{\theta,P}
+(\boldsymbol\nabla\hat{\mathbf e}_P)\!\cdot\!\boldsymbol\nabla Z_{\theta,P}\\
&+2n_B(\widetilde x_P)
(\hat{\mathbf e}_P\!\cdot\!\boldsymbol\nabla\Phi_{b,P})
\boldsymbol\nabla\Phi_{b,P}
\Big\}=0,
\end{aligned}
\]
with $[(\boldsymbol\nabla\hat{\mathbf e})\!\cdot\!\boldsymbol\nabla Z]_a=\partial_a\hat e^i\,\partial_iZ$. The second and third terms come from the saddle-splay part. For a radial director $n=\hat{\mathbf r}$, the elastic bracket obeys the stronger identity
\[
\boldsymbol\nabla\!\cdot(Z\boldsymbol\nabla n)
-(\boldsymbol\nabla\!\cdot n)\boldsymbol\nabla Z
+(\boldsymbol\nabla n)^{\mathsf T}\boldsymbol\nabla Z
=-2n\left(\frac{Z}{r^2}+\frac{n\cdot\boldsymbol\nabla Z}{r}\right).
\]
Its transverse projection vanishes for arbitrary angular dependence of $Z$, and $\mathcal S[n]=0$ away from the core. The thermal term is also radial for a spherical source, giving the local law as an exact solution of the unprojected auxiliary equations. For a disk or another nonspherical thermal source, its transverse forcing still has to be solved. The positive-stock interpretation additionally obeys Appendix~O.3. Replacing $z_P$ by $\widetilde z_P$ only in the total $\mathcal Q$ action would give the excess flux $\mathcal Q'(\widetilde z_P)(\hat{\mathbf e}\!\cdot\!\nabla\Phi_b)\hat{\mathbf e}-\nabla\Phi_b=\mathbf J_\perp-(\mathbf 1-\hat{\mathbf e}\hat{\mathbf e}^{\mathsf T})\nabla\Phi_b$, whose transverse part survives as $n_B\to0$ and spoils the Newtonian limit. The explicit Newtonian term and $\mathcal Q_\perp$ keep the excess flux proportional to $n_B$.

For an axisymmetric carrier write $\hat{\mathbf e}=\cos\alpha\,\hat{\mathbf R}+\sin\alpha\,\hat{\mathbf z}$ with $\alpha=\arctan(z/R)+\beta$. With $\ell=-\sin\alpha\,\partial_R\alpha+\cos\alpha\,\partial_z\alpha$ and $b=\cos\alpha\,\partial_R\alpha+\sin\alpha\,\partial_z\alpha$,
\[
\mathcal S=\left(\ell-\frac{\cos\alpha}{R}\right)^2+b^2 ,
\]
the squared difference of the meridional and azimuthal splays plus the squared bend. The numerical solutions minimize the static energy directly in the smooth variable $\beta$, which vanishes on the axis and midplane and at the isolated boundary, using bilinear Gauss elements and an exact sparse Hessian. The linearized coherence length is
\[
\boxed{
\xi_{\theta,P}^2=
\frac{c\,X_P^2}
{n_B(\widetilde x_P)|\boldsymbol\nabla\Phi_{b,P}|^2},}
\]
with $\xi_\theta(R_0,0)\simeq6\,\mathrm{kpc}$ for the Milky-Way stock at $c=0.03$ and $12\,\mathrm{kpc}$ for the baryonic proxy at $c=1$. Section~16 gives the resolved solutions. They lie in the stiff limit. With the stock amplitude, $0.03\le c\le1$ moves the Milky-Way vertical potential by $20\ {\rm km^2\,s^{-2}}$. With the proxy, $1/4\le c\le4$ moves it by $17$. The potential zero of an embedded carrier, the occupation-weighted frame term, and the remaining connection cross terms stay completion data. A radial director costs no stiffness about any center. In satellite and binary configurations the thermal term sets the core positions, and in the two-carrier solution below each core settles within $20$--$240\,$pc of its body's centre.

\paragraph{Stock variation of the stiffness.}
For the stock amplitude, $Z_\theta=2c\Phi_R^2$, the isolated-carrier stiffness energy is
\[
E_\theta=\frac{c}{8\pi G}\int\Phi_R^2\mathcal S[n]\,d^3x,
\qquad \nabla^2\Phi_R=4\pi G\rho_R.
\]
At fixed stationary director and with the same potential-zero convention, its variation with respect to stock density is
\[
U_\theta(y)=\frac{\delta E_\theta}{\delta\rho_R(y)}
=-\frac{c}{4\pi}\int\frac{\Phi_R(x)\mathcal S[n(x)]}{|x-y|}\,d^3x,
\qquad \nabla^2U_\theta=c\Phi_R\mathcal S[n].
\]
This term enters the stock's mechanical force. Its Laplacian $c\Phi_R\mathcal S$ is the stiffness source in the stationarity condition of the stock selection functional above. In a multi-carrier configuration the support weights and their variations enter as well. The identity fixes the stiffness contribution to the coupled force; a complete transport action must also specify the remaining stock and bath terms.

\paragraph{Several active carriers: director solutions.}
For overlapping carriers the joint coarse potential $\Psi=\sum_P\Phi_{b,P}$ replaces $\Phi_{b,P}$ in the thermal term, each carrier keeps its own director, and spatial interpolation weights $\chi_P\ge0$ with $\sum_P\chi_P=1$ weight their constitutive energies. Eliminating the auxiliary potentials gives the static energy
\[
\boxed{
U=-\frac{1}{8\pi G}\int|\boldsymbol\nabla\Psi|^2
-\frac{a_0^2}{8\pi G}\int\sum_P\chi_P\,\mathcal Q_\perp(\widetilde z_P)
+\frac{1}{8\pi G}\int\sum_P\chi_P\,\frac{Z_{\theta,P}}{2}\,\mathcal S[\hat{\mathbf e}_P],}
\]
with $\widetilde z_P=(\hat{\mathbf e}_P\!\cdot\!\boldsymbol\nabla\Psi)^2/a_0^2$ and excess flux $\mathbf J_\perp=\sum_P\chi_P\,n_B(\widetilde x_P)(\hat{\mathbf e}_P\!\cdot\!\boldsymbol\nabla\Psi)\hat{\mathbf e}_P$. Each director minimizes $U$. Because every director is stationary, $F_P=-dU/dz_P$ collects only explicit dependences:
\begin{itemize}
\item the flux force $-(4\pi G)^{-1}\int\mathbf J_\perp\!\cdot\!\boldsymbol\nabla\partial_z\Phi_{b,P}$ of Section~15;
\item the support variation $\partial\chi_P/\partial z_P$;
\item the variation of the stiffness amplitude;
\item the force on each director's hedgehog core.
\end{itemize}
Translating every carrier together with its support, amplitude and core leaves $U$ invariant, so $\sum_PF_P=0$ on-shell.

The Milky-Way--LMC pair tests this directly. The carriers are Plummer spheres of $6\times10^{10}$ and $3\times10^9M_\odot$ ($b=3$ and $1\,$kpc) $50\,$kpc apart, with $\chi_L=\tfrac12[1-\tanh((r_L-r_s)/w)]$ around the LMC and baryonic amplitude $Z_{\theta,P}=2c\Phi_{b,P}^2$. Both directors were minimized by damped Newton iteration on a stretched axisymmetric grid out to $5000\,$kpc, with the core heights free. Forces were computed from the explicit derivatives, referenced to the joint isotropic field, and from finite differences of the minimized $U$. For $r_s=12\,$kpc:
\begin{center}
{\small
\begin{tabular}{lcccc}
\hline
case & $F_{\rm MW}/F_L$ & support force & negative excess & label-dependent\\
 & [$10^{12}$] & on LMC [$10^{12}$] & near LMC & acceleration\\
\hline
$c=1$, $w=1.5\,$kpc & $+1.745/-1.744$ & $+0.160$ ($9.2\%$) & $-7.5\times10^9M_\odot$ & $1.6\%$ / $2.3\%$\\
$c=1$, $w=4\,$kpc & $+1.731/-1.745$ & $+0.128$ ($7.3\%$) & $-3.4\times10^9M_\odot$ & $2.4\%$ / $2.8\%$\\
$c=4$, $w=1.5\,$kpc & $+1.726/-1.725$ & $+0.197$ ($11.4\%$) & $-8.4\times10^9M_\odot$ & $0.4\%$ / $1.0\%$\\
joint field, $c\to0$ & $+1.779/-1.779$ & $0$ & $-0.35\times10^9M_\odot$ & $0$\\
\hline
\end{tabular}
}
\end{center}
Forces are in $M_\odot({\rm km\,s^{-1}})^2{\rm kpc^{-1}}$. The support force is positive away from the Milky Way, with its fraction of the mutual pull in brackets. The label-dependent acceleration is quoted at the LMC centre and the inner support edge, as a fraction of the local gravitational acceleration.

An independent implementation, written from the specification alone, reproduced the $c=1$, $w=1.5$ row, the core positions, the zero-stiffness limit and the LMC support monopole of $1.07\times10^{10}M_\odot$. It found a net force of $2\times10^{-6}$ of $|F|$, falling as the domain grows.

Four results follow:
\begin{itemize}
\item \textbf{Cores.} Each response core is pulled toward the other carrier with about $4\times10^{11}$, a quarter of the mutual pull. Free cores settle $20\,$pc (MW) and $150$--$240\,$pc (LMC) from the body centres. Pinned cores pass that force to their bodies and change the totals by $0.2\%$. Evaluating the explicit-derivative forces with pinned cores but without the core force understates each body force by about a quarter.
\item \textbf{Support force.} The support variation pushes the satellite away from its host with 7--11\% of the mutual pull. The net pull ends up 2--3\% below the joint-field value.
\item \textbf{Interface layer.} The layer lies at $r_L=10$--$15\,$kpc and carries one to three times the LMC's baryonic mass. A joint field has only the saddle feature between the bodies.
\item \textbf{Label dependence.} $\Phi_{b,P}$ depends on which baryons belong to $P$, so LMC-labelled and MW-labelled matter at the same point accelerate differently.
\end{itemize}

\paragraph{A single continuous director.}
A continuous nematic field can carry a satellite hedgehog only with a compensating defect: a hyperbolic hedgehog on the axis, or a $-\tfrac12$ Saturn-ring disclination. With $Z=2c[\chi_{\rm MW}\Phi_{b,{\rm MW}}^2+\chi_L\Phi_{b,L}^2]$ and defect positions minimized, the outcome depends on the support radius:
\begin{itemize}
\item At $r_s=12\,$kpc the ground state has no LMC hedgehog. The director crosses the satellite tilted by $42^\circ$ from its local field, the LMC excess inside $8\,$kpc falls by $46\%$, and the layer is $-2.2\times10^9M_\odot$.
\item At $r_s=20\,$kpc the ground state has the hedgehog plus a hyperbolic defect at $r_L\simeq13\,$kpc. The layer returns at $-5.5\times10^9M_\odot$, and the label-dependent acceleration reaches $44\%$ of the LMC's gravity at its centre.
\end{itemize}

\paragraph{Interface flux.}
The layer has a simple cause. Gauss's law ties the net excess inside any surface around the satellite to the flux through it. Inside, the satellite's response is nearly isotropic and its flux leaves radially. Outside, the host director is near-radial about the host and passes only the flux component along itself. The mismatch in normal flux is absorbed as surface charge on the interface. Moving the mismatch into a defect only relocates the charge. Melting the order in a thin shell does not remove it either, because the blocked flux is set by the host director outside the shell. The layer disappears only where both directors give way to the isotropic response, over a zone wide enough to carry the lateral flux.

\paragraph{Interface order and stock amplitude.}
Two changes remove both defects. The first lets the nematic order melt where reconstructed carrier profiles overlap, in the manner of a Landau--de Gennes order parameter:
\[
\boxed{
\begin{aligned}
&q=(\chi_{\rm host}-\chi_{\rm sat})^2,\qquad
\mathsf M_P=(1-q)\mathbf 1+q\,\hat{\mathbf e}_P\hat{\mathbf e}_P^{\mathsf T},\\
&\widetilde z_P=\frac{\boldsymbol\nabla\Psi\cdot\mathsf M_P\cdot\boldsymbol\nabla\Psi}{a_0^2},\qquad
\mathbf J_\perp=\sum_P\chi_P\,n_B(\widetilde x_P)\,\mathsf M_P\boldsymbol\nabla\Psi .
\end{aligned}}
\]
The stiffness density is weighted by $q^2$. A pure carrier has $q=1$ and is unchanged, while equal overlap gives $q=0$ and the isotropic joint response. The thermal term in the director equation becomes $2\chi_Pq\,n_B(\widetilde x_P)(\hat{\mathbf e}_P\!\cdot\!\boldsymbol\nabla\Psi)\boldsymbol\nabla\Psi$.

The second takes the stiffness amplitude from the carrier's response stock (Appendix~O.3):
\[
\boxed{Z_{\theta,P}=2c\,\Phi_{R,P}^2,\qquad \nabla^2\Phi_{R,P}=4\pi G\rho_{R,P},}
\]
with the supports and director cores also tied to the stocks. Baryons then couple to the carrier labels only through $\Psi$ and fall universally, and every label-dependent force acts on the stocks. Varying $\Phi_{R,P}$ places the stiffness source $c\,\chi_Pq^2\Phi_{R,P}\mathcal S[\hat{\mathbf e}_P]$ in the stock's equilibrium. For an isolated carrier in equilibrium the stock is a fixed functional of the baryons, so the same amplitude serves the single-carrier solutions.

The tests model each stock as the capped RAR excess with the Lagrangian budget of Appendix~O.3 ($C_R=14$ for the Milky Way, $33$ for the LMC). For these stocks $\Phi_R/\Phi_b\simeq2.8$ at $R_0$ and about $17$ at the LMC support edge. On a $500\,$kpc grid, finite-difference forces carry a systematic of about 5\% (from shifting both bodies together). With the joint-field reference at $-0.33\times10^9M_\odot$ of negative excess and $10.66\times10^9M_\odot$ of LMC excess inside $8\,$kpc, the results are:
\begin{center}
{\small
\begin{tabular}{lcccc}
\hline
amplitude, $c$, order & $r_s/w$ & negative excess & LMC excess $<8\,$kpc & $F_L/F_L^{\rm joint}-1$\\
 & [kpc] & [$10^9M_\odot$] & [$10^9M_\odot$] & \\
\hline
baryonic, $1$, no melt & $12/1.5$ & $-7.47$ & $11.3$ & $-2.4\%$\\
baryonic, $1$, melt & $12/1.5$ & $-7.42$ & $11.4$ & \\
baryonic, $1$, melt & $12/4$ & $-3.42$ & $10.8$ & \\
baryonic, $1$, no melt & $16/8$ & $-4.71$ & $10.5$ & $-11\%$\\
baryonic, $1$, melt & $16/8$ & $-0.86$ & $10.8$ & $-2\%$\\
stock, $1$, melt & $16/8$ & $-0.82$ & $10.6$ & $-12\%$\\
stock, $0.03$, melt & $16/8$ & $-0.85$ & $10.7$ & $-1.2\%$\\
stock, $0.01$, melt & $16/8$ & $-0.78$ & $10.8$ & $-1.4\%$\\
\hline
\end{tabular}
}
\end{center}
Melting in a narrow shell leaves the layer unchanged, as the flux argument predicts. An overlap zone half as wide as the support reduces it to about 2.6 times the joint-field value.

The stock amplitude does not touch the layer. At $c=1$ its deep potentials bring back a repulsive collective force of 12\%, which is gone for $c\le0.03$. The stock accelerations are small:
\begin{itemize}
\item with the stock amplitude at $c=0.03$ they are $3.7$ and $5.0\ ({\rm km\,s^{-1}})^2{\rm kpc^{-1}}$ at the LMC centre and inner support edge, $0.6$ and $0.85\%$ of the Milky Way's pull;
\item with the baryonic amplitude the corresponding accelerations on baryons are $12.9$ and $13.7$;
\item the LMC's self-gravity at $8$--$16\,$kpc is $460$--$970\ ({\rm km\,s^{-1}})^2{\rm kpc^{-1}}$, so these forces displace the stock from its baryons by a small fraction of its radius.
\end{itemize}

Within $40\,$kpc of the LMC, $4$--$9\times10^9M_\odot$ less excess remains than for a joint field. The deficit is smooth and reappears farther out. It comes from the Milky-Way director staying near-radial about the Milky Way, tilted by $7$--$11^\circ$.

The normalization is bounded from both sides. The Milky-Way benchmark with the stock amplitude (Wang~B baryons) gives:
\begin{center}
{\small
\begin{tabular}{lcccc}
\hline
$c$ & $1$ & $0.03$ & $0.01$ & $0.003$\\
\hline
$\Delta\Phi(8.17,1.055)$ [${\rm km^2\,s^{-2}}$] & $1396.5$ & $1416.4$ & $1444.9$ & $1500.1$\\
\hline
\end{tabular}
}
\end{center}
The stiff limit therefore holds for $c\gtrsim0.03$. At $c=0.03$ the two baryonic models give $1416$ and $1315\ {\rm km^2\,s^{-2}}$ with $\chi^2_\Phi=137.1$ and $135.9$ (Section~16).

With the stock amplitude, a satellite's stiffness depends on its own stock, which lies between two bounds:
\begin{itemize}
\item the stock relaxed to its equilibrium in the host field, giving $\Phi_R/\Phi_b=9$--$16$ at the half-light radius;
\item the stock carried from infall with a Lagrangian cap of $300$--$3000$, giving $28$--$47$ for Fornax and Sculptor and $190$--$600$ for Crater~II and Antlia~2.
\end{itemize}
In internally dominated dwarfs the stiff director is radial, which is also the field direction in the interior. A rigid radial director changes the virial dispersion of Fornax and Sculptor by $-4$ to $+3\%$. In external-field-dominated dwarfs the director lies along the host field and the stiffness decides whether it follows the internal tilts. Linearizing about the host field gives the flux multiplier $m(\mathbf k)=a_\parallel\cos^2\theta_k+[1+n_ef(k)]\sin^2\theta_k$, with $a_\parallel=1+n_e+\tfrac{x_e}{2}n_B'(x_e)$ and director transfer $f(k)=1/(1+k^2\xi_\theta^2)$. Then $\sigma^2=\langle m\rangle\sigma_N^2$. Relative to the isotropic response this gives:
\begin{center}
{\small
\begin{tabular}{lcccccc}
\hline
 & \multicolumn{3}{c}{Crater II} & \multicolumn{3}{c}{Antlia 2}\\
$c$ & relaxed & cap $300$ & cap $3000$ & relaxed & cap $300$ & cap $3000$\\
\hline
$1$ & $-1\%$ & $-23\%$ & $-32\%$ & $0$ & $-14\%$ & $-25\%$\\
$0.1$ & $0$ & $-9\%$ & $-17\%$ & $0$ & $-3\%$ & $-10\%$\\
$0.03$ & $0$ & $-4\%$ & $-9\%$ & $0$ & $-1\%$ & $-5\%$\\
$0.01$ & $0$ & $-2\%$ & $-5\%$ & $0$ & $0$ & $-2\%$\\
\hline
\end{tabular}
}
\end{center}
A fully locked director costs half of $\sigma$. If satellite stocks relax after infall the completion is harmless for any $c\le1$. If they keep their Lagrangian stock, $c\lesssim0.03$, and the Milky Way requires $c\gtrsim0.03$, so the two constraints meet at $c\simeq0.03$.

The multi-carrier completion has three conditional parts:
\begin{itemize}
\item the overlap width, the normalization and the rule tying supports and cores to the stocks, which were set by these tests and are not yet derived;
\item the rigid-blend estimate for a dwarf embedded in its host, $9$--$16\,M_b$ of negative excess, which suggests that dwarf interfaces also need melting and has not been solved with free directors;
\item the melted-interface and stock-amplitude runs, which pinned the cores at the body centres and have not yet been reproduced independently.
\end{itemize}

\paragraph{Closed-time-path metric response.}
The open-system definition remains useful for deriving the leading contact from a microscopic state. Integrating out transverse substrate variables $\xi^A$ in their physical state $\rho_\perp$ gives
\[
e^{\,i\Gamma_\perp[g_+,g_-]/\hbar}
=
\operatorname{Tr}_{\perp}
\left(
U_{\perp}[g_+]\,\rho_\perp\,
U_{\perp}^{\dagger}[g_-]
\right).
\]
The total effective functional is
\[
\Gamma_{\mathrm{eff}}
=I_0[g_+,\psi_+]-I_0[g_-,\psi_-]
+\Gamma_\perp[g_+,g_-;\psi_+,\psi_-].
\]
The physical metric equation is obtained by varying the difference field and then setting the two histories equal. Unitarity requires
\[
\Gamma_\perp[g,g]=0,
\qquad
\Gamma_\perp[g_+,g_-]^*
=-\Gamma_\perp[g_-,g_+].
\]
Diffeomorphism invariance of the influence functional must yield the Ward identity that makes the induced response covariantly conserved.

The specified leading stationary response is
\[
h_{\mu\nu}^{\perp,P}
=-\frac{2\phi_{\perp,P}}{c^2}
\left(\bar g_{\mu\nu}+2u_\mu u_\nu\right),
\qquad
\Delta\Phi_{\perp,P}=\Delta\Psi_{\perp,P}.
\]
It is generated by the quasistatic contact
\[
\boxed{
\Gamma_\perp^{\rm QS}
=-\frac12\sum_{P\in\mathfrak A}w_P
\left\langle h_\Delta,
\mathcal E h_{\perp,P}[T^{\mathfrak A}_c]\right\rangle .
}
\]
This functional vanishes when the two histories coincide and is real in the stationary mean. It gives $T^{00}_{\perp,P}=\rho_{\perp,P}c^2$ and $T^{ij}_{\perp,P}=0$, hence zero scalar anisotropic stress. It also leaves the vacuum graviton propagator unchanged. The causal material response described below supplies damped transport poles without adding a vacuum metric polarization.

At linear order, the Ward identity follows from the Bianchi identity for $\mathcal E$. Nonlinear conservation requires covariant variation of the carrier records, reconstruction, and interpolation profiles:
\[
(\mathcal C_\ell\Pi_P)[\varphi^*g,\varphi^*\psi,\varphi^*\mathcal R]
=\varphi^*(\mathcal C_\ell\Pi_P)[g,\psi,\mathcal R].
\]
The metric variation acts on $e_P$, $E_P$, $\mathcal C_\ell$, and $\chi_P$ through their carrier and geometric arguments. A specified covariant source map fixes these local contributions. The full nonlinear dynamics requires the same variables to be included in the action.

\paragraph{Quadratic data and response convention.}
Here $g_+$ and $g_-$ label the two CTP branches. Let $h_{\mu\nu}^{(1,2)}$ be their metric perturbations and define the average and difference fields $h_c=(h^{(1)}+h^{(2)})/2$ and $h_\Delta=h^{(1)}-h^{(2)}$. With $x^0=ct$, the linear metric coupling is
\[
\delta I_\perp
=\frac{1}{2c}\int d^4x\,\sqrt{-\bar g}\,
h_{\mu\nu}\widehat T_\perp^{\mu\nu}.
\]
Define the densitized source
\[
J_{\mu\nu}\equiv\frac{\sqrt{-\bar g}}{2c}h_{\mu\nu}.
\]
Let $J^c=(J^{(1)}+J^{(2)})/2$ and $J^\Delta=J^{(1)}-J^{(2)}$. Then the universal quadratic form in this source convention is
\[
\Gamma_\perp^{(2)}
=
\int d^4x\,d^4y\,
J^\Delta_{\mu\nu}(x)
\widehat\Pi_R^{\mu\nu\alpha\beta}(x,y)
J^c_{\alpha\beta}(y)
+\frac{i}{2}\int d^4x\,d^4y\,
J^\Delta_{\mu\nu}(x)
\widehat{\mathcal N}^{\mu\nu\alpha\beta}(x,y)
J^\Delta_{\alpha\beta}(y).
\]
The retarded kernel is
\[
\boxed{
\widehat\Pi_R^{\mu\nu\alpha\beta}(x,y)
=+\frac{i}{\hbar}\theta(x^0-y^0)
\left\langle
\left[
\widehat T_\perp^{\mu\nu}(x),
\widehat T_\perp^{\alpha\beta}(y)
\right]
\right\rangle_{\rho_\perp}
+\widehat\Pi_{\mathrm{local}}^{\mu\nu\alpha\beta}(x,y).
}
\]
The sign follows from the stated $+\tfrac12h_{\mu\nu}T^{\mu\nu}$ action coupling and the Kubo susceptibility convention. The local or \emph{seagull} term comes from the explicit metric variation of $\widehat T_\perp^{\mu\nu}$ and of the covariant measure. It is required, together with the commutator, for the diffeomorphism Ward identity. Define $\Delta\widehat T_\perp^{\mu\nu}=\widehat T_\perp^{\mu\nu}-\langle\widehat T_\perp^{\mu\nu}\rangle_{\rho_\perp}$. The noise kernel is the correspondingly normalized symmetrized correlator,
\[
\widehat{\mathcal N}^{\mu\nu\alpha\beta}(x,y)
=\frac{1}{2\hbar}
\left\langle
\left\{
\Delta\widehat T_\perp^{\mu\nu}(x),
\Delta\widehat T_\perp^{\alpha\beta}(y)
\right\}
\right\rangle_{\rho_\perp}.
\]

A microscopic calculation begins with the transverse Hessian and metric vertex,
\[
\mathcal K_{AB}(x,y)
=\frac{\delta^2 I_{\perp}}
{\delta\xi^A(x)\delta\xi^B(y)},
\qquad
V_A^{\mu\nu}(x;y)
=\frac{\delta^2 I}
{\delta\xi^A(x)\delta g_{\mu\nu}(y)}.
\]
In the physical state,
\[
G_{R}^{AB}(x,y)
=+\frac{i}{\hbar}\theta(x^0-y^0)
\operatorname{Tr}\!\left(
\rho_\perp[\xi^A(x),\xi^B(y)]
\right),
\]
and schematically
\[
\widehat\Pi_R^{\mu\nu\alpha\beta}
=V_A^{\mu\nu}G_R^{AB}V_B^{\alpha\beta}
+\widehat\Pi_{\mathrm{contact}}^{\mu\nu\alpha\beta}.
\]
This correlator/Hessian calculation is the microscopic audit that must reproduce the specified contact, its projector, and its local Ward term.

\paragraph{Lensing and point-source limit.}
The no-slip contact gives one effective density,
\[
\rho_{{\rm lens},P}
=\rho_P+\frac{1}{4\pi G}
\boldsymbol\nabla\!\cdot\mathbf J_{\perp,P}.
\]
For a point mass in the deep regime, $A=\sqrt{GMa_0}$ and
\[
g=\frac{A}{r},
\qquad
\rho_\perp=\frac{A}{4\pi Gr^2},
\qquad
\Sigma_\perp=\Delta\Sigma_\perp=\frac{A}{4GR},
\qquad
\alpha_\infty=\frac{2\pi A}{c^2}.
\]
This fixes the leading lensing normalization. Resolved disks and lenses require direct solutions of the nonspherical auxiliary equation.

\paragraph{Causal response and stability.}
The transverse transport equation is
\[
\tau_0\ddot\phi_{\perp,P}+\dot\phi_{\perp,P}
-D\nabla^2\phi_{\perp,P}
=-4\pi GD\rho_{\perp,P},
\qquad
\frac{D}{\tau_0}=c^2,
\qquad
\tau_0^{-1}=H_0.
\]
The Fourier response relative to the static solution and its Laplace poles are
\[
F(\omega,k)=\frac{Dk^2}{Dk^2-\tau_0\omega^2-i\omega},
\qquad
F(0,k)=1\quad(k\ne0),
\]
\[
s_\pm=\frac{-1\pm\sqrt{1-4\tau_0Dk^2}}{2\tau_0}.
\]
For $k>0$ both poles have negative real part, and the equation has characteristic speed $c$. At $k=0$ the poles are $0$ and $-1/\tau_0$. The forced quasistatic response has corrections beginning at $O(v^2/c^2)$; applying this limit also requires sufficiently small homogeneous transients. Underdamped transients decay on the time scale $2\tau_0$, so causal propagation alone does not establish rapid equilibration. On galactic wavelengths these source-free modes oscillate at approximately $ck$ despite their long damping time. The persistent tilted field of Section~16 is sourced by mechanically supported stock; slow wave damping alone would not preserve its density quadrupole.

\paragraph{Local nesting test.}
For the Solar System, the Milky Way is the active carrier and the Sun is a resolved Einstein descendant. The transverse functional never evaluates a nonlinear response on $\boldsymbol g_\odot+\boldsymbol g_{\rm MW}$. Thus
\[
Q_2^{\odot,\perp}=0,
\]
while the smooth parent field contributes only a tide of order $g_{\rm MW}/R_0\simeq9.2\times10^{-31}\,{\rm s}^{-2}$. The same rule makes wide binaries Newtonian apart from that tide. This is a source-space projection; a spatial cutoff alone would leave the coarse Solar cross-term and would double-count descendant halos. The parent's conditional expectation still contains the Sun smoothed over its resolution $L$. That gives $Q_2=[4/15\sqrt{2\pi}](n_\parallel-n_\perp)GM_\odot/L^3$, so Cassini requires $L\ge0.049\,$pc for a following Galactic director and $L\ge0.024\,$pc for a locked one (Section~16).


\paragraph{Noise and dissipation.}
The thermal free energy fixes the local static slope $\chi_0=-\partial n_B/\partial E=\beta n_B(1+n_B)$. Let $\delta n$ denote the coarse occupation response about that mean. A stable linear realization with two relaxation poles is
\[
\left[\partial_t^2+(\gamma+\Omega)\partial_t+\gamma\Omega\right]\delta n
=-\chi_0\gamma\Omega\,\delta E+\xi,
\qquad \Omega\ge\gamma>0.
\]
Here $\xi$ is the bath force in this normalization. With Fourier convention $e^{-i\omega t}$, the retarded susceptibility is
\[
\boxed{\chi^R_{nE}(\omega)=-\chi_0
\frac{\gamma\Omega}{(\gamma-i\omega)(\Omega-i\omega)}.}
\]
For $\Omega>\gamma$ its time-domain kernel is
\[
\chi^R_{nE}(t)=-\chi_0\frac{\gamma\Omega}{\Omega-\gamma}
\left(e^{-\gamma t}-e^{-\Omega t}\right)\theta(t),
\]
with continuous limit $-\chi_0\gamma^2t e^{-\gamma t}\theta(t)$ when the rates coincide. The kernel vanishes before the perturbation and has poles at $-i\gamma$ and $-i\Omega$. These are local material relaxation poles. Spatial transport uses the causal operator specified above.

For positive frequency the dissipative part is
\[
\boxed{D(\omega)=-\operatorname{Im}\chi^R_{nE}(\omega)
=\chi_0\frac{\gamma\Omega(\gamma+\Omega)\omega}
{(\gamma^2+\omega^2)(\Omega^2+\omega^2)}>0.}
\]
For all $\gamma,\Omega>0$, including coincident rates, it obeys the kernel-checked static dispersion sum rule
\[
\frac{2}{\pi}\int_0^\infty\frac{D(\omega)}{\omega}\,d\omega=\chi_0.
\]
The positive spectral density therefore admits a Gaussian oscillator-bath realization. Fixing its temperature determines the noise paired with this response. For the symmetrized occupation spectrum, without the $1/\hbar$ normalization used for the metric influence kernel above, the equilibrium relation is \citep{Clerk2010}
\[
\boxed{\mathcal N_{nn}(\omega)
=\hbar\coth\!\left(\frac{\beta\hbar\omega}{2}\right)D(\omega).}
\]
The product is even and nonnegative. Since $D(\omega)=O(|\omega|^{-3})$ at high frequency, the equal-time variance
\[
\sigma_n^2=\int_{-\infty}^{\infty}\frac{d\omega}{2\pi}\,
\mathcal N_{nn}(\omega)
\]
is finite. The two rates and the static slope completely specify this linear model. A microscopic matching can determine the rates and test the selected spectral shape.

The inequality $x\coth x\ge1$ and the dispersion sum rule give
\[
\sigma_n^2\ge\frac{2}{\pi\beta}\int_0^\infty\frac{D(\omega)}{\omega}\,d\omega
=\frac{\chi_0}{\beta}=n_B(1+n_B).
\]
The classical population variance is recovered when the response spectrum lies in the thermal regime. At finite bandwidth the variance follows from the spectral integral. This Gaussian theory describes a coarse response coordinate calibrated to the Bose mean and static slope. Its use for force fluctuations requires small excursions about the stationary response; higher occupation cumulants require a nonlinear completion.

At $|\omega|\ll\Omega$ the mean response reduces to the single-relaxation form $-\chi_0\gamma/(\gamma-i\omega)$. In the additional classical regime $\beta\hbar|\omega|\ll1$, its noise becomes
\[
\mathcal N_{nn}(\omega)\simeq
\frac{2\gamma n_B(1+n_B)}{\gamma^2+\omega^2}.
\]
For the assumed horizon temperature, $\beta\hbar=2\pi/H$, so the classical frequency condition is $|\omega|\ll H/(2\pi)$. The full quantum spectrum applies outside this interval. Over a stationary epoch the bath and response parameters are held fixed; their cosmological evolution requires an adiabatic approximation or the corresponding two-time kernels.

For independent equivalent cells define $\mathcal F=\sigma_n^2/[n_B(1+n_B)]$. Linearizing $g=g_b(1+n_B)$ gives
\[
\frac{\sigma_g^2}{g^2}
=\frac{\mathcal F n_B}{N_{\rm eff}(1+n_B)}.
\]
A normalized temporal measurement window $W(0)=1$ replaces $\sigma_n^2$ by
\[
\sigma_{n,W}^2=\int\frac{d\omega}{2\pi}\,
|W(\omega)|^2\mathcal N_{nn}(\omega).
\]
Spatial cell correlations enter through their covariance matrix and determine the effective independent count. Predictions for gravitational fluctuations also require the response from occupation to the measured field. The model supplies a causal, positive quantum noise spectrum while leaving these rate and averaging scales as effective matching data.

\paragraph{Response time, noise, and the sign of tracking.}
For the two-pole response, define the mean delay
$\bar\tau=\gamma^{-1}+\Omega^{-1}$ and
$\mathcal F=\sigma_n^2/(\chi_0/\beta)$. Quantum thermal pairing at
$\beta\hbar=2\pi/H$ implies
\[
\mathcal F\,H\bar\tau\ge4.
\]
The derivation uses the zero-point lower bound on the noise integral
and the logarithmic-mean inequality for the two positive rates.
For the illustrative matter-era occupation at $x=1$, three-percent
tracking requires $H\bar\tau\lesssim0.0253$, hence
$\mathcal F\gtrsim158$. A three-percent force rms then requires
$N_{\rm eff}\gtrsim6.47\times10^4$ under independent-cell averaging.
Spatial correlations, observational windows, and the physical coupling
determine the actual force noise. This inequality constrains the selected
two-pole contact and is not an observational exclusion.

Passive lag also has a definite sign. If an equilibrium target $q(t)$
declines and $\dot n=\Gamma(q-n)$ with $\Gamma>0$ and $n(t_i)\ge q(t_i)$,
then $n(t)\ge q(t)$; positive relaxation cascades retain the inequality.
At fixed baryons, lag behind declining $H$ therefore cannot suppress an
initially equilibrated high-redshift response. Underoccupied assembly
states would require their own preparation history. Relabeling the same
stock between channels leaves its total stress unchanged. In the tested
adiabatic mass-loss model, the surviving mass ratio from $z=2.3$ is
$0.989398$, giving an inferred $a_0$ ratio $m^8=0.9183$, compared with
$H_0/H(2.3)=0.2972$. Neither effect selects the required galactic tracking.


\paragraph{Coarse response, noise, and memory.}
For a fixed reconstruction $C$ and adjoint source coupling, the response and symmetrized noise transform together:
\[
\chi_c^R=C\chi_f^RC^\dagger,\qquad
\mathcal N_c=C\mathcal N_fC^\dagger.
\]
Positive fine covariance gives positive coarse covariance. If the fine system has one temperature and
$\mathcal N_f=\hbar\coth(\beta\hbar\omega/2)D_f$, then
$D_c=CD_fC^\dagger$ gives the same quantum relation after reconstruction. This includes spatially correlated cells. Independent-cell counting is an additional covariance assumption.

An autonomous retained evolution requires a further identity. If the fine generator is $L$, exact closure on $Cx$ requires $CL=L_cC$. Eliminating a coupled hidden block generally retains memory: for
$\dot x=Ax+By$, $\dot y=Ex+Dy$,
\[
\dot x(t)=Ax(t)+Be^{Dt}y(0)
+\int_0^t Be^{D(t-s)}E x(s)\,ds.
\]
The initial hidden state and the memory kernel must be retained or controlled by a stated approximation. The local relaxation model above is a selected response realization. Positive reconstruction and the thermal relation alone do not make it Markovian.

\paragraph{Closure audit.}
The leading EFT fixes the thermal contact, local nonspherical action, carrier nesting, momentum balance between active carriers through the joint coarse source, no-slip metric response, point-source lensing law, stable causal transport, Solar quadrupole, wide-binary limit, and a causal finite-bandwidth quantum occupation response. The director equations follow from one auxiliary action; they reproduce the resolved Milky-Way and SPARC tests and remove the phantom sheet from disk vertical dynamics. The other conditional inputs are the amplitude identification $X\leftrightarrow\Phi_R$ with its baryonic proxy for isolated carriers, the interface order, the saddle-splay coefficient of the director stiffness, the tilt-mode identification and lever, active-carrier projector, potential-zero prescription, no-slip tensor vertex, and canonical transport parameters. The nonlinear carrier variation, response-rate matching, and spatial noise correlations determine the dynamical range of this leading EFT. Its microscopic embedding provides a further test of the selected constitutive terms. Old-disk vertical dispersions, disk-modeled stream shapes, resolved nonspherical lensing maps, and the ultrafaint-dwarf phase boundary remain observational tests.

\subsection*{N.9 Action-level outcome}

The action problem divides into four parts:
\[
\boxed{
\begin{aligned}
\text{ordinary longitudinal capacity}
&=\text{reduced Einstein constraint action},\\
\text{spherical nonlinear capacity}
&=\text{composite Misner--Sharp invariant},\\
\text{galactic transverse response}
&=\text{specified carrier-resolved leading EFT},\\
&\quad\text{with explicit record and response matching;}\\
&\quad\text{microscopic GFT derivation still required},\\
\text{capacity-exhaustion boundary}
&=\text{open microscopic boundary functional}.
\end{aligned}
}
\]
This division solves the parent--descendant double-counting problem and closes the ordinary weak-field parent action. The transverse sector is explicit in the slowly varying limit: the active-carrier projector removes parent--descendant cross-terms, the joint coarse source makes the forces between active carriers reciprocal, the thermal contact and local auxiliary action give the nonspherical RAR, and the metric contact gives no-slip lensing without changing the vacuum graviton pole. The director system is a variational completion at thin-disk direction reversals, with its saddle-splay coefficient set by spherical exactness. The record algebra, reconstruction, thermal matching, and metric contact specify this effective realization. Its microscopic embedding tests the cell Hamiltonian, phase structure, tensor vertex, director stiffness, and local Ward term. The full nonlinear variation of carrier records, the variational dynamics of the positive-stock projection in Appendix~O.3, and the response-rate and correlation-volume matching remain open. The boundary-channel functional at $q_{\mathrm{geo}}=0$ is separate work.

\subsection*{N.10 Nonlinear geometry from directional channel counts}
\label{aqg:metric}

The coarse area observables of Q.11a determine a spatial metric without
changing the four-port primitive cell. Measure cuts normal to
\(e_x,e_y,e_z,(e_x+e_y)/\sqrt2,(e_x+e_z)/\sqrt2,(e_y+e_z)/\sqrt2\).
Write their area ratios as \(r_A=A_A/A_0=1+q_A>0\). The cofactor matrix
and metric are
\[
\mathcal C_{ii}=r_i^2,\qquad
\mathcal C_{ij}=r_{ij,+}^2-\frac{r_i^2+r_j^2}{2},\qquad
g=\sqrt{\det\mathcal C}\,\mathcal C^{-1}.
\]
The metric domain is \(\mathcal C>0\), and
\(\sqrt{\det g}=(\det\mathcal C)^{1/4}\).
These are six coarse directional measurements, each containing many
primitive channels. Holding every directional area fixed would remove
the metric response tested in H.13.

Let \(E_a\) be an orthonormal basis of symmetric \(3\times3\) matrices,
with unit diagonal elements and off-diagonal elements \(1/\sqrt2\).
For \(g_a=\operatorname{tr}(E_ag)\), define
\[
J_{aA}=\frac{\partial g_a}{\partial q_A},\qquad
D^{ab}=
\frac{\operatorname{tr}(E_agE_bg)
-\tfrac12\operatorname{tr}(gE_a)\operatorname{tr}(gE_b)}
{\sqrt{\det g}},\qquad
K=J^{-1}DJ^{-T}.
\]
The differential used to evaluate \(J\) is
\[
dg=\sqrt{\det\mathcal C}\left[
\tfrac12\operatorname{tr}(\mathcal C^{-1}d\mathcal C)\mathcal C^{-1}
-\mathcal C^{-1}d\mathcal C\,\mathcal C^{-1}\right].
\]
On the selected ADM continuation, in units \(\hbar=c=1\),
\[
H_{\rm kin}=\frac{16\pi G_0}{V_{\rm ref}}p^TK(q)p,\qquad
V_{\rm curv}=-\frac{V_{\rm ref}\sqrt{\det g}}{16\pi G_0}R[g].
\]
Here \(V_{\rm ref}\) is the coordinate volume of a coarse block.
The kinetic signature is five positive directions and one negative
direction. The Hamiltonian constraint is retained; this indefinite form
is not an unconstrained positive oscillator Hamiltonian.

\paragraph{Homogeneous nonlinear check.}
The vacuum Kasner solution with exponents
\((p_1,p_2,p_3)=(-1/3,2/3,2/3)\) has
\[
g_{ii}=(t/t_0)^{2p_i},\qquad
r_i=(t/t_0)^{1-p_i},\qquad
r_{ij,+}=\sqrt{\frac{r_i^2+r_j^2}{2}}.
\]
Its diagonal canonical momentum density is
\(\pi^{ii}=(p_i-1)/(16\pi G_0t_0g_{ii})\); the count momentum follows
from the canonical Jacobian, with the reference volume included when
using block-integrated momenta. Integrating from \(t=90\) to \(110\)
with \(t_0=100\) gives relative metric error \(2.45\times10^{-11}\)
and normalized Hamiltonian-constraint residual
\(1.47\times10^{-16}\). Freezing \(K\) at the identity instead gives a
\(7.95\%\) velocity error. This test resolves the nonlinear homogeneous
kinetic response. Spatial curvature requires the separate tests below.

\paragraph{A coupled continuum solution.}
One Einstein--Dirac action supplies a working common metric for the
capacity clock, a self-gravitating source, and a quantum probe. In its
static scalar reduction let \(s\) denote the canonical capacity deficit:
\[
N(s)=\sqrt{1-s/S_\infty},\qquad
ds^2=-N^2dt^2+W^2d\mathbf x^2.
\]
For the barotropic fluid
\(p=K_{\rm fluid}\rho_0^2\), \(\mathcal E=\rho_0+p\), with
\(K_{\rm fluid}=4\) and central enthalpy \(0.3\), the TOV solution on the
matched \(G=G_*\) branch has
\[
R=1.78920,\quad M=0.306425,\quad
2M/R=0.342527,\quad N(0)=0.600690,\quad s(0)=0.0294966.
\]
Lengths use electron Compton units \(L_*/\epsilon_e\); \(M\) is the
geometric mass in those length units. The fluid parameters define a
demonstration source. The capacity and proper-volume readings agree between areal and
isotropic coordinates to \(6\times10^{-10}\).
After the spinor rescaling \(\psi\mapsto W^{3/2}\psi\), the Hermitian
static Dirac Hamiltonian is
\[
H_D=\beta mN+\frac12\{\boldsymbol\alpha\cdot\mathbf p,N/W\}.
\]
This form is the static isotropic specialization of the covariant Dirac
Hamiltonian \citep{AQGObukhov2011}. Varying this same operator with
respect to \(s\) fixes its capacity source; the finite-difference check
agrees to \(2.4\times10^{-9}\).
Three spatial resolutions give converging wave-packet delay and bending;
the last refinement changes them by \(0.50\%\) and \(0.047\%\).
Norm and energy errors are below \(4\times10^{-13}\) in the finest run.
The probe's additional metric backreaction is neglected.
The continuum action supplies the coupling in this example.
The finite orbital construction in H.15 preserves the native gap and
clock. Its full Dirac limit, species structure, and dynamical source
backreaction must still reproduce this continuum example.

\subsection*{N.11 Configuration measure and quantum ordering}
\label{aqg:ordering}

The six-area change of variables fixes a kinematic measure once the ADM
kinetic form and lapse convention have been chosen. Direct determinants give
\[
\det D=-\tfrac12\det g,\qquad
|\det J|=64\sqrt2\,\frac{\prod_A r_A}{\det\mathcal C},\qquad
\sqrt{|\det K^{-1}|}
=128\,\frac{\prod_A r_A}{(\det\mathcal C)^{5/4}}.
\]
Drop the constant 128 and write
\[
\mu(q)=\frac{\prod_A(1+q_A)}{(\det\mathcal C)^{5/4}},\qquad
\phi\in L^2(\mu\,dq),\qquad \psi=\sqrt\mu\,\phi.
\]
The last expression is the half-density in the orthonormal count basis.
For a microscopic pattern \(x\) in the normalized symmetric count orbit,
\[
\psi_{\rm pattern}(x)=
\frac{\psi(n(x))}{\sqrt{\prod_A\binom{M_A}{n_A(x)}}}.
\]
This configuration measure leaves the capacity Gibbs state unchanged.
A physical inner product on constrained states requires a further
construction.

Put \(c_K=16\pi G_0/V_{\rm ref}\). The scalar Laplace--Beltrami ordering is
\[
H_{\rm LB}=-c_K\mu^{-1}\partial_A(\mu K^{AB}\partial_B).
\]
Coordinate covariance allows additional scalar quantum terms.
The completion used here Weyl-orders the kinetic symbol in metric
components and transports that operator unitarily into area coordinates.
In flat \(dg_a\) half-density coordinates,
\[
\operatorname{Op}_W(D^{ab}\pi_a\pi_b)
=-\partial_a(D^{ab}\partial_b)
-\tfrac14\partial_a\partial_bD^{ab}.
\]
The metric configuration measure is \(\mu_g\propto(\det g)^{-1/2}\).
For \(l_a=\partial_a\ln\mu_g\), the scalar term generated by transporting
its Laplace--Beltrami operator to flat measure is
\[
U=\tfrac12\partial_a(D^{ab}l_b)+\tfrac14l_aD^{ab}l_b
=\frac{15}{32\sqrt{\det g}},\qquad
\partial_a\partial_bD^{ab}=\frac{27}{8\sqrt{\det g}}.
\]
Both coefficients follow by exact evaluation at \(g=I\),
\(SL(3)\) congruence invariance, and homogeneity under \(g\mapsto sg\).
Independent checks away from the identity agree within
\(5\times10^{-15}\). Thus, in the count half-density representation,
\[
\boxed{H_{W,\rm count}=H_{{\rm LB},\rm count}
-c_K\frac{21}{16\sqrt{\det g}}.}
\]
The correction carries \(\hbar^2\) when units are restored. It fixes
the ordering selected above without changing the classical principal
symbol or the bare coefficient \(G_0\). Weyl ordering after an arbitrary
nonlinear canonical change would generally give a different operator.

\paragraph{Exact differential-operator identity.}
At a fixed finite spatial regulator, keep metric configurations continuous.
Suppose
\[
H[N]=\sum_iN_iT_i(g_i,\pi_i)+V[N](g),\qquad
D[\xi]=\sum_i\pi_i\cdot(\mathcal L_\xi g)_i,
\]
where each \(T_i\) is local and quadratic in momenta,
\(V[N]\) is coordinate-only and linear in the externally specified lapse,
and the spatial Lie-derivative stencil is linear in \(g\).
For the metric Weyl prescription,
\[
\boxed{[\widehat C_1,\widehat C_2]
=i\,\operatorname{Op}_W(\{C_1,C_2\}),}
\]
for \((C_1,C_2)=(H[N],H[M]),(D[\xi],H[N]),(D[\xi],D[\zeta])\).
Kinetic terms at different cells use disjoint variables; at the same
cell they are proportional and cancel in the first commutator.
Higher odd Moyal terms in a kinetic--potential bracket vanish because
the kinetic symbol has no third momentum derivative.
The diffeomorphism symbol is a quadratic polynomial in the full
canonical variables, so its third total derivatives vanish.

These are identities on smooth functions compactly supported in the
interior of the positive metric domain. They do not choose a
self-adjoint extension at its boundary. They also retain every
classical spatial-regulator defect: if
\(\{H[N],H[M]\}=D[\beta]+E_a[N,M]\), the right side contains
\(\operatorname{Op}_W(E_a)\). A finite-difference replacement of the
configuration derivatives introduces a separate error.
Exact symbolic two-site tests with nonconstant kinetic coefficients
give zero residual in all three identities; omitting the Weyl scalar
term gives a nonzero control residual.

\subsection*{N.12 Nonlinear spatial constraints and finite-count tests}
\label{aqg:constraints}

Nonlinear Regge discretizations generally break the exact vertex
displacement symmetry of the flat linearized theory
\citep{AQGDittrich2010}. The spatial and configuration limits must
therefore be tested separately.

For a spatially varying sector with two transverse Killing directions, use
\[
g_{ij}(x)=
\begin{pmatrix}a&0&0\\0&b&d\\0&d&c\end{pmatrix},
\qquad h=\begin{pmatrix}b&d\\d&c\end{pmatrix},
\qquad B=h^{-1}h'.
\]
The canonical momentum of \(d\) is \(p_d=2\pi^{yz}\).
The four configuration variables and two first-class continuum
constraints retain two local configuration degrees of freedom.
The diffeomorphism generator and integrated curvature are
\[
D[\xi]=\int dx\,[p_a(\xi a'+2a\xi')
+p_b\xi b'+p_c\xi c'+p_d\xi d'],
\]
\[
\int dx\,N\sqrt gR=
\int dx\left[
\frac{N\sqrt{\det h}}{4\sqrt a}
\{(\operatorname{tr}B)^2-\operatorname{tr}(B^2)\}
+\frac{N'\sqrt{\det h}}{\sqrt a}\operatorname{tr}B\right].
\]
Generic smooth periodic data, away from the constraint surface, test
the full Hamiltonian vector fields against
\[
\begin{aligned}
\{H[N],H[M]\}&=D[(NM'-MN')/a],\\
\{D[\xi],H[N]\}&=H[\xi N'],\\
\{D[\xi],D[\zeta]\}&=D[\xi\zeta'-\zeta\xi'].
\end{aligned}
\]
At 128 spatial sites the relative defects are
\begin{center}
\begin{tabular}{lrr}
\hline
Bracket & second-order stencil & fourth-order stencil\\
\hline
\(H,H\) & \(1.256\times10^{-3}\) & \(1.607\times10^{-5}\)\\
\(D,H\) & \(4.104\times10^{-3}\) & \(8.375\times10^{-5}\)\\
\(D,D\) & \(3.979\times10^{-3}\) & \(2.179\times10^{-5}\)\\
\hline
\end{tabular}
\end{center}
The last measured fourth-order convergence rates are \(3.964\),
\(3.982\), and \(3.989\). This checks the interacting spatial
constraint limit in the stated symmetry sector. It does not establish
exact closure at finite spacing or uniform convergence on quantum states.

\paragraph{Finite-area Hamiltonian commutator.}
For count configurations \(x,y\) differing only at spatial site \(i\),
the off-diagonal matrix element obeys the exact finite identity
\[
[\widehat H[N],\widehat H[M]]_{xy}
=T_{i,xy}\left[
N_i\{V[M]_y-V[M]_x\}
-M_i\{V[N]_y-V[N]_x\}\right].
\]
The corresponding continuous-configuration current is
\[
v_i=2c_KK_i(M_i\nabla_iV[N]-N_i\nabla_iV[M]),\qquad
\widehat D[v]=-i(v\cdot\partial+\tfrac12\operatorname{div}v),
\]
with \([\widehat H[N],\widehat H[M]]=i\widehat D[v]\).
Three spatial sites, all six metric coordinates per site, and the
nonlinear curvature potential test this limit on a smooth complex
half-density. At count spacings \(0.02,0.01,0.005,0.0025\), corresponding
to mean counts \(100,200,400,800\), the relative current errors are
\[
0.0104584,\quad0.00261967,\quad0.000655234,\quad0.000163828.
\]
The count convergence is second order. On the same three-cell
spatial regulator, comparison with the actual spatial Lie derivative
gives a separate \(8.5389\%\) error. The smallest count-current error
is therefore not the full constraint-algebra error. The target-current
derivative cross-check is \(7.2\times10^{-12}\).
The local Weyl correction cancels from this particular commutator;
the ordering tests in N.11 are independent checks of that term.

The full six-coordinate graph lift can also be evaluated without
projecting away the microscopic patterns. With two switches per
coordinate it has 729 count states, 4096 patterns, and 21 single or
paired move directions. The microscopic/count intertwining error is
\(6.4\times10^{-17}\), Hermiticity is exact to stored precision, and
the unitary norm error is \(2.2\times10^{-16}\).
The smallest cofactor eigenvalue in that test is \(0.847\).
These results establish a finite kinematic realization of the selected
operator on its tested metric domain.

\paragraph{Physical state space.}
A finite reflecting count box exposes the next obstruction.
The one-cell homogeneous Weyl-ordered operator in the 729-state test
has 725 positive and four negative eigenvalues, no zero eigenvalue
at tolerance \(10^{-8}\), and a closest eigenvalue of
\(-68.09835\) in the test units. Its eigenvector residual is
\(1.2\times10^{-9}\). The constant null state of the Laplace--Beltrami
control is removed by the Weyl term.
This boundary-dependent example neither rules out physical states in
a refined theory nor supplies one. Subtracting an eigenvalue would
change the constraint.

An exact undeformed coordinate transformation also cannot act on a
fixed finite grid of commuting coordinate operators: the diagonal of
\([D,Q]\) vanishes in their joint eigenbasis, whereas a nonzero
coordinate velocity is diagonal. Equivalently, every derivation of
the finite algebra of functions vanishes on its idempotents.
A continuum limit, a deformed algebra, or a different observable
representation is required.

The homogeneous spectral construction in N.13 supplies a positive
rigging form on a different, unbounded metric domain. It leaves this
finite reflecting-box result intact. The inhomogeneous completion
needs compatible spatial and count refinement, control of
\(\det\mathcal C=0\), common domains, and a physical inner product
for all local constraints. A possible coarse window
\(1\ll\bar N\ll\lambda^2/A_{\rm ch}\) separates area resolution
from a physical wavelength \(\lambda\); uniform physical-state
convergence across it remains open. The exact orbital refinements
of H.15 preserve a supplied joint kinematic operator and do not
solve this gravitational constraint problem.

\subsection*{N.13 Homogeneous volume, shape, and a positive rigging form}
\label{aqg:rigging}

The homogeneous metric domain admits an explicit distributional
constraint construction. Separate volume and unimodular shape by
\[
g=e^{2\alpha_{\rm vol}}h,\qquad \det h=1,\qquad
\alpha_{\rm vol}=\frac16\log\det g.
\]
The DeWitt line element and its absolute density become
\[
ds_{\rm DW}^2=e^{3\alpha_{\rm vol}}
\left[-24\,d\alpha_{\rm vol}^2+
 {\rm tr}(h^{-1}dh\,h^{-1}dh)\right],\qquad
d\mu\propto e^{9\alpha_{\rm vol}}d\alpha_{\rm vol}\,
 d{\rm vol}_{\rm sh}.
\]
The shape space is \(SL(3,\mathbb R)/SO(3)\).
For the homogeneous Weyl constraint of N.11, set
\(\phi=e^{-3\alpha_{\rm vol}}\chi\). Its action is
\[
\widehat C\phi=\frac{c_K}{24}e^{-6\alpha_{\rm vol}}
\left(\partial_{\alpha_{\rm vol}}^2-24\Delta_{\rm sh}
 -\frac{81}{2}\right)\chi.
\]
Transport to the flat volume density with
\(u=e^{9\alpha_{\rm vol}/2}\phi\), and let
\(B=e^{3\alpha_{\rm vol}/2}\). The congruence
\[
\widetilde C=\frac{24}{c_K}B C_{\rm flat}B
=\partial_{\alpha_{\rm vol}}^2-24\Delta_{\rm sh}-\frac{81}{2}
\]
reduces the spectral problem to a wave equation on shape space.

Choose a self-adjoint spectral realization of
\(A=-\Delta_{\rm sh}\). If \(\lambda>27/16\), put
\(\omega_\lambda=\sqrt{24\lambda-81/2}\).
On a suitable spectral test space away from the threshold,
\[
\widetilde\eta(f,g)=
\sum_{s=\pm}\int_{\lambda>27/16}
\frac{d\nu(\lambda,\sigma)}{2\omega_\lambda}\,
\overline{\widehat f(s\omega_\lambda,\lambda,\sigma)}
\widehat g(s\omega_\lambda,\lambda,\sigma)
\]
is positive semidefinite and annihilates \(\widetilde C\) on either
argument. Here \(\sigma\) labels spectral multiplicity and the hat
includes the Fourier transform in \(\alpha_{\rm vol}\).
The original constraint has the induced form
\[
\eta_C(f,g)=\widetilde\eta(Bf,Bg),
\]
on a compatible domain. Quotienting its null vectors and completing
gives a positive physical Hilbert space for this homogeneous
prescription. Constraint rescaling and domain choices enter such a
construction \citep{AQGLouko2011}; this form does not establish their
uniqueness.

An explicit Iwasawa generalized shape mode has
\(\lambda=5/2\), hence
\(\omega_\lambda=\sqrt{39/2}\). It is distributional and requires
spectral smearing for physical wave packets. Substitution into the
original area-coordinate Weyl stencil gives residuals
\[
0.0586840,\quad0.0146319,\quad0.00365554,\quad
0.000913731,\quad0.000228424
\]
at steps \(0.02,0.01,0.005,0.0025,0.00125\), respectively.
The second-order decrease tests the coordinate transport and ordering.
No auxiliary matter clock enters this construction:
\(\alpha_{\rm vol}\) is a geometric variable. The local constraints,
metricity boundaries, and common spatial refinement of an inhomogeneous
theory remain outside this one-constraint result.

\subsection*{N.14 Finite histories and the constrained conformal response}
\label{aqg:histories}

A finite Hermitian count Hamiltonian defines an ordered-jump expansion
of \(e^{-\beta H}\), with each jump carrying \(-H_{yx}\).
Diagonal rephasing leaves the product of jump amplitudes around a
closed cycle invariant. The 729-state operator contains a triangle
with product \(-0.254534167364\) in the specified test units.
Its count histories therefore cannot all be made positive by
diagonal rephasing. The positive bond curvature contraction of Q.11
does not remove these Hamiltonian signs.

The negative scale direction can be isolated algebraically. With
\(v=\sqrt{\det g}\) and the count-coordinate scale vector
\(r=\mathbf1+q\), the kinetic matrix splits as
\[
K_{\rm conf}=-\frac{rr^T}{6v},\qquad
K_{\rm sh}=K+\frac{rr^T}{6v}.
\]
The shape matrix is positive semidefinite of rank five. Its tested
finite stencil still contains nine negative triangles. Replacing
signed transitions by their absolute values changes the operator
and introduces a spurious ultraviolet mode in the tested construction.

Static gravitational response also requires the constrained metric
components. For a Gaussian contour \(h=Ry\), both the source and
observable rotate with the field. The covariance is
\[
R|\lambda|^{-1}R^T=K^{-1}.
\]
The transpose implements the analytic bilinear form. In the static
source check, this prescription gives \(G/G_0=1\). Deleting the
negative direction gives \(4/3\); using its absolute eigenvalue with
an unrotated source gives \(5/3\); retaining only the two spatial TT
wave modes gives zero. These are tests of the quadratic constrained
response, with \(G_0\) already present in its action.

A conditional nonlinear integration is available at zero
cosmological constant. Write \(g_{\mu\nu}=\Omega^2\widehat g_{\mu\nu}\)
and include the appropriate boundary terms. In units
\(\hbar=c=1\),
\[
S_E=-c_R\int\sqrt{\widehat g}\,
\left[\Omega^2\widehat R+6(\widehat\nabla\Omega)^2\right],
\qquad c_R=\frac1{16\pi G_0}.
\]
A prescribed proper-length source contributes
\(\int\sqrt{\widehat g}\,j\Omega\).
Fix the mean scale, put \(\Omega=1+i\chi\) with real zero-mean
\(\chi\), and let \(P\) remove the constant mode. Define
\[
A=P(-6\widehat\Delta+\widehat R)P,\qquad
\eta=P(j-2c_R\widehat R).
\]
For \(A>0\) and the flat contour measure \(D\Omega\), Gaussian
integration gives, up to shape-independent normalization and the
common fixed-mode contour phase,
\[
\begin{split}
S_{\rm eff}={}&-c_R\int\sqrt{\widehat g}\,\widehat R
 +\int\sqrt{\widehat g}\,j\\
&+\frac1{4c_R}\langle\eta,A^{-1}\eta\rangle
 +\frac12\log\det A .
\end{split}
\]
The inverse and determinant act on the nonconstant subspace.
A cosmological or fixed-four-volume term brings an \(\Omega^4\)
dependence. Positivity of \(A\) restricts the shape domain, and the
flat contour measure has not been derived from the finite count
histories. A negative Gaussian expectation of a proposed
\((1+i\chi)^8\) insertion excludes that insertion at the tested
variance; it does not exclude every possible complete measure.

The coordinate density in a transfer kernel must be treated
consistently. If \(T\) is the orthonormal-basis matrix and \(\mu_x\)
the coordinate density, then
\[
K_{xy}=\frac{T_{xy}}{\sqrt{\mu_x\mu_y}},\qquad
\prod_x\mu_x\prod_{\langle xy\rangle}K_{xy}
=\prod_{\langle xy\rangle}T_{xy}
\]
on a closed history, with one measure factor for each occurrence of
a vertex. Inserting the coordinate density again changes the history
weight. This cancellation fixes the finite kernel convention but
does not derive the continuum contour measure.

\subsection*{N.15 An auxiliary proper-clock transfer}
\label{aqg:auxclock}

A supplied positive bounded lapse operator
\(N=\sqrt{I-s/S_\infty}\) permits a proper-clock construction on a
common operator domain. Let \(H_0\) be the Hermitian count generator,
and define
\[
H_N=N^{1/2}H_0N^{1/2},\qquad
\mathbb H=H_N+NP_\tau+jB_V,\qquad
B_V=N^{1/2}(V/V_{\rm ref})N^{1/2}.
\]
Here \(P_\tau\) is an ideal auxiliary clock momentum. It is distinct
from the finite retained marked register and the native mass
derivative clock of H.14--H.15. At fixed auxiliary momentum \(p\),
\[
U_t(p,j)=e^{-it(H_N+pN+jB_V)}.
\]
At zero sources the integrated clock is
\[
\mathcal T_t=iU_t^\dagger\left.\partial_pU_t\right|_0
=\int_0^tU_s^\dagger N U_s\,ds,\qquad
\mathcal T_{t+s}=\mathcal T_s+U_s^\dagger\mathcal T_tU_s.
\]

For a fixed reference energy \(E\), the equation
\[
(H_N-iN\partial_\tau)\Psi=E\Psi
\]
conserves the flux norm \(\Psi^\dagger N\Psi\).
The half-density \(\chi=N^{1/2}\Psi\) obeys
\[
i\partial_\tau\chi=(H_0-EN^{-1})\chi,\qquad
V_E(\tau)=e^{-i\tau(H_0-EN^{-1})}.
\]
On a self-adjoint domain \(V_E\) is unitary, while
\[
A_E=N^{-1/2}V_EN^{1/2},\qquad A_E^\dagger N A_E=N
\]
is the corresponding scalar transfer. A faithful renewal isometry
\(W_*\) can be joined to the graph-dependent frame of Q.11a:
\[
D_{\rm out}(I\otimes W_*)(V_E\otimes I_{\rm cap})D_{\rm in}^\dagger.
\]
This is an isometry when the indicated domains and frame maps match.
Subdividing a numerical step introduces no additional physical reset
or particle phase. A turning point of spatial volume does not stop
this clock. At a lapse zero the inverse-lapse construction needs a
new domain treatment, even if the reference-time unitary still exists.

Multiple local clocks require the normal-deformation algebra. The
commutator of two normal deformations has the prescribed spatial
diffeomorphism term; it need not vanish. After subtracting that term,
any residual regulator defect remains an obstruction.
Exponentiating the regulated generators cannot remove it. This
single-clock transfer establishes unitary evolution on its stated
domain, without resolving many-clock refoliation or the full physical
constraint space.

\section*{Appendix O: The Cosmological Term from Fixed Event Measure and Marked Source-Clock Release}

The capacity-clock action fixes the local exchange for a supplied ensemble history. A composition-preserving allocation gives the background relation $\omega_{\Lambda,0}=\eta_*\omega_c^{\rm wb}I_{\eta,\rm rel}$, with the weighted integral defined in O.3. The closure precision and formation abundance are fixed; the assembly and mixing history sets the remaining normalization. O.12 reports the matched direct-perturbation likelihood for the composition-preserving allocation and the approximate PPF likelihood for the linear-depletion benchmark; each numerical result keeps its own model definition.

The dark-energy sector remains conditional. Fixed primitive event measure
gives the continuum reduction, and a coefficient-free dilation into the
nine-state marked register gives the accumulation law. A caustic can create
multiple clock branches, but identity-free committed capacity cannot record
which stream is which. Release begins only when irreversible baryonic mixing
removes the distinguishable Lagrangian carrier of a source clock. Faithful
renewal then fixes complete rerouting for each failed elementary record block.
Additive record pricing turns progenitor mutual-information loss into the
released-stock fraction. A closed-cohort advection--diffusion theorem makes
its fully resolved position version monotone; the coarse physical observable
requires a separate check. A constrained capacity-clock action fixes the
covariant exchange for a supplied release history.
The cosmological normalization follows from the supplied assembly and
mixing history. The approximate PPF benchmark gives $I_{\rm rel}=94.2116$
and $\Omega_{\Lambda,0}\simeq0.70666$. The first self-gravitating run
supplies a label-dependent mixing kernel. With the retention gate taken from
the halo occupation of feedback-resolved simulations, it gives
$\Omega_{\Lambda,0}=0.709\pm0.02$ against the required $0.707$. O.12 reports these results, and O.14 states
the remaining premises and systematics.

\subsection*{O.1 Fixed primitive event measure and the unimodular reduction}

The microscopic history is a sequence of admissibility-closed network
configurations, one per substrate tick (Appendix~G). The sharp primitive
face map resolves each update reversibly into a transmitted geometric
sector and a marked complementary record,
\[
V_{3/2}\otimes V_{3/2}\;\longrightarrow\;V_3\oplus\left(V_1^P\otimes V_1^H\right),
\qquad 16\to 7+9,
\]
supplying a concrete primitive event register. A mean physical measure
\(v_{\rm ev}\) gives the coarse volume \(\mathcal V_4=v_{\rm ev}N\).
C.8 separates this measure from the canonical reference volume and
states a paired-cell event assignment. Conditioning a point process on
fixed total \(N\) still permits local count fluctuations; it does not
by itself fix a local volume form. The following continuum premise is
stronger. The equations use the matched \(G=G_*\) branch; a different
physical normalization replaces that coefficient throughout.

\paragraph{Continuum premise.}
\emph{The primitive event count supplies a fixed, nondynamical volume
form: continuum geometry resides in connectivity and relational channel
data at fixed microscopic event measure, and the metric determinant is
not an independent local degree of freedom.} Under this premise the
continuum variation is unimodular, $g_{\mu\nu}\delta g^{\mu\nu}=0$, and
variation of the Einstein--Hilbert and matter actions gives the
trace-free equation
\[
\boxed{\;
R_{\mu\nu}-\tfrac14 Rg_{\mu\nu}
=8\pi G_*\left(T_{\mu\nu}-\tfrac14 Tg_{\mu\nu}\right).\;}
\]
If the matter sector is covariantly conserved, the contracted Bianchi
identity gives $\nabla_\nu(R+8\pi G_*T)=0$, hence $R+8\pi G_*T=4\Lambda$
for a spacetime constant, and
\[
G_{\mu\nu}+\Lambda g_{\mu\nu}=8\pi G_*T_{\mu\nu},
\qquad
T^{(\Lambda)}_{\mu\nu}=-\rho_\Lambda g_{\mu\nu},
\qquad
\boxed{\,w_\Lambda=-1\,}
\]
exactly, with $\rho_\Lambda=\Lambda/8\pi G_*$. The cosmological term is
an integration constant conjugate to the primitive history volume, not a
local vacuum-energy coupling \citep{Unruh1989,HenneauxTeitelboim1989}.

\paragraph{Vacuum-shift decoupling.}
A constant shift $\mathcal L_m\mapsto\mathcal L_m-\rho_{\rm vac}$ leaves
the traceless source invariant and contributes only an overall factor
$e^{-i\rho_{\rm vac}\mathcal V_4}$ to a fixed-event history sector,
which cancels from normalized history weights. It shifts the
integration constant, $\Lambda\mapsto\Lambda-8\pi G_*\rho_{\rm vac}$,
but cannot independently curve spacetime. This decoupling is classical
and holds at the level of normalized history weights; its stability
under quantization requires the fixed-measure property to survive
coarse-graining, which is not established here
\citep{Weinberg1989}. The equilibrium normalization of Section~25
supplies the complementary statement that the equilibrium bulk
contribution is exactly zero; the present reduction explains why
residual constant shifts are gravitationally sequestered rather than
merely subtracted.

\paragraph{Falsifier of the reduction.}
If the coarse-grained microscopic Hessian of the condensate embedding
contains an independently propagating determinant (conformal-volume)
mode, the unimodular reduction does not follow and ordinary
vacuum-energy sensitivity returns. This is a definite target for the
fixed-$N_4$ simulation program: hold the event count fixed, vary
connectivity and edge geometry, and test whether a scalar volume mode
survives in the effective spectrum.



\subsection*{O.2 Volume projection and the homogeneous-clock prescription}

For a coarse history slab of $N$ equal-volume primitive events, let $\mathbf r$ be a release profile and define
\[
P_0=\frac{\mathbf1\mathbf1^T}{N},\qquad
\mathbf r=\bar r\mathbf1+\mathbf r_\perp,\qquad
\sum_e r_{\perp e}=0.
\]
With $\mathcal V_4=v_*N$ and $r_e=\eta_*q_e$, the volume-integrated contribution is
\[
v_*\sum_e r_e=\eta_*\bar q\,\mathcal V_4.
\]
Equal event measure selects this constant-mode component uniquely. Its conversion into a homogeneous vacuum potential is the constitutive choice of the cosmological branch.

Let $\phi$ be the capacity clock and let $\mathbb E_{\rm bg}$ denote the statistically homogeneous release ensemble specified by the initial statistics and the astrophysical release model. Define
\[
q=\Gamma_{\rm rel}\rho_u,\qquad
\bar q(\phi)=\mathbb E_{\rm bg}[q](\phi).
\]
The background history fixes a one-variable potential,
\[
\boxed{V_{F,\phi}(\phi)=\eta_*\bar q(\phi),\qquad
\rho_\Lambda(x)=V_F(\phi(x)).}
\]
The initial vacuum density fixes the integration constant. The local clock evaluates this potential at each event. The function $V_F$ is held fixed in the field equations and their perturbations, so a realized change in a remote release rate supplies no instantaneous update to the local vacuum term. The expectation is determined by the chosen background dynamics; local exchange is governed by the covariant action of O.3.

A spatial average on a clock slice estimates $\bar q$:
\[
\langle q\rangle_{\Sigma,V}=\frac1V\int_{\Sigma,V}q\,dV,\qquad
\mathbb E_{\rm bg}[\langle q\rangle_{\Sigma,V}]=\bar q.
\]
For connected equal-clock correlation $C_q(\mathbf x,\mathbf y)$, its sampling variance is
\[
\operatorname{Var}(\langle q\rangle_{\Sigma,V})
=\frac1{V^2}\int_{\Sigma,V}dV_x\int_{\Sigma,V}dV_y\,
C_q(\mathbf x,\mathbf y).
\]
An integrable spatial correlation gives a variance decreasing as $V^{-1}$ at large volume. The ensemble law uses $\bar q$ independently of the finite-volume sampling fluctuation. The excursion-set and mixing calculation supplies this statistical history through $F_{\rm rel}$.

On the homogeneous background, $\phi=t$ and $\dot{\bar\rho}_\Lambda=\eta_*\bar q$. For scalar perturbations,
\[
\delta\rho_\Lambda=V_{F,\phi}\,\delta\phi,\qquad
\delta\rho_\Lambda-
\frac{\dot{\bar\rho}_\Lambda}{\dot{\bar\phi}}\delta\phi=0.
\]
Vacuum homogeneity is therefore a statement on uniform-clock hypersurfaces. The interaction still changes capacity growth and the metric potentials, whose lensing and integrated Sachs--Wolfe signals follow from the coupled perturbation equations \citep{Martinelli2019}. The homogeneous drain acts per unit proper volume and obeys the local stock budget in O.3. Specifying the release ensemble and its initial normalization fixes this effective branch without introducing a new local vacuum degree of freedom.

\subsection*{O.3 Source-associated capacity after carrier loss and the accumulation law}

While a baryonic progenitor label remains recoverable, the source-associated
allocation may occupy the saturated constraint surface of Section~20. Carrier
loss removes that source-label constraint. It does not create a new material
species. The renewal theorem sends the replaceable closure register to the
universal state $\rho_*$, and the marked complement is an internal history
output rather than an ordinary spacetime fluid. The selected infrared ledger assigns no independent zeroth-derivative metric stress to the orthogonal history output. Absence of an additional propagating pole alone does not establish this contact-term assignment; the coupled metric vertex must reproduce it:
\[
\boxed{
T^{\mu\nu}_{H,\perp}\big|_{k^0}=0,\qquad
\delta T^{\mu\nu}_{H,\perp}
=O[(kL_{\rm block})^2]}
\]
within the local derivative expansion. Any such higher-gradient influence is
a correction to the retained metric/capacity response, not a new conserved
component.

The metric-visible remainder is instead associated with the retained coarse
baryonic carrier $Y$. At scales larger than that carrier, covariance and the
one-metric worldtube theorem give the leading monopole stress
\[
\boxed{
T_C^{\mu\nu}=\rho_C U^\mu U^\nu+\Delta T_C^{\mu\nu},\qquad
\frac{|\Delta T_C|}{\rho_C}
=O\!\left(\frac{\sigma_Y^2}{c^2},k^2R_Y^2\right).}
\]
This is an effective source-locked capacity stress, not a particle ontology.
The zero-pressure leading form follows from the absence of an independent
local scalar pressure operator at this order; finite carrier size and
velocity dispersion supply the displayed corrections.

\paragraph{Covariant capacity-clock completion.}
The source-associated capacity and its homogeneous vacuum ledger follow from one constrained scalar action,
\[
\boxed{
S_{C\Lambda}=\int d^4x\sqrt{-g}\,
\left[\lambda\left(X-\tfrac12\right)-V_F(\phi)\right],
\qquad
X=-\frac12\nabla_\mu\phi\nabla^\mu\phi.}
\]
It extends the pinned action of Section~20 by a clock potential and belongs to the same constrained-scalar class \citep{ChamseddineMukhanov2013,ChamseddineMukhanovVikman2014}. Variation gives
\[
X=\frac12,
\qquad
\nabla_\mu\!\left(\lambda\nabla^\mu\phi\right)=V_{F,\phi},
\qquad
T_{C\Lambda}^{\mu\nu}=\lambda u^\mu u^\nu-V_Fg^{\mu\nu},
\]
where $u_\mu=-\nabla_\mu\phi$. The unit-gradient constraint makes this flow geodesic, $u^\mu\nabla_\mu u_\nu=0$. With $\rho_C=\lambda$ and $\rho_\Lambda=V_F$,
\[
\nabla_\mu T_C^{\mu\nu}=\nabla^\nu V_F,
\qquad
\nabla_\mu T_\Lambda^{\mu\nu}=-\nabla^\nu V_F.
\]
\paragraph{Positive internal allocation.}
Write $C=\rho_u+R$, where $\rho_u$ is source-resolved capacity and $R$ is reallocated capacity. A local release event transfers $q=\Gamma_{\rm rel}\rho_u\ge0$ between them. The ensemble potential supplies a total drain $Q=\dot V_F\ge0$ per unit proper volume. Requiring this drain to preserve the internal composition fixes its allocation uniquely. For $C>0$, the conditions
\[
d_u+d_R=Q,\qquad Rd_u=\rho_u d_R
\]
give $d_u=\rho_u Q/C$ and $d_R=RQ/C$. The covariant currents are
\[
\begin{aligned}
Q_u^\nu&=-q u^\nu+\frac{\rho_u}{C}\nabla^\nu V_F,\\
Q_R^\nu&=q u^\nu+\frac{R}{C}\nabla^\nu V_F,\\
Q_\Lambda^\nu&=-\nabla^\nu V_F,
\qquad Q_u^\nu+Q_R^\nu+Q_\Lambda^\nu=0.
\end{aligned}
\]
On the common clock, $\nabla^\nu V_F=-Q u^\nu$, so
\[
\begin{aligned}
\dot\rho_u+\theta\rho_u&=-q-\rho_u Q/C,\\
\dot R+\theta R&=q-RQ/C,\\
\dot C+\theta C&=-Q.
\end{aligned}
\]
At either internal boundary the release and fractional drain point into the nonnegative domain, provided total capacity remains positive. For $F=R/C$,
\[
\dot F=\Gamma_{\rm rel}(1-F),\qquad 0\le F\le1.
\]
Thus $F$ measures the reallocated fraction of surviving capacity. A positive finite update, in comoving cell masses, is
\[
\rho_u^+=s(1-p)\rho_u,\qquad R^+=s(R+p\rho_u),
\qquad C^+=sC,\qquad 0\le p,s\le1.
\]
The lost mass $(1-s)C$ supplies the drain. Conservation and positivity follow algebraically for every such step. Composition preservation is the constitutive premise selecting this allocation; conservation alone allows other allocations.

On the homogeneous ensemble background, $Q=\eta_*\bar q$ and $F=F_{\rm rel}$ give
\[
a^3\bar C=B e^{-\eta_*F},\quad
\bar\rho_u=(1-F)\bar C,\quad \bar R=F\bar C,
\quad \bar q=B a^{-3}e^{-\eta_*F}\dot F,
\]
where $B=\rho_{c,0}^{\rm wb}$ and the initial fraction is zero. Hence
\[
V_F(\phi)=\rho_{\Lambda,i}+\eta_*B
\int_{t_i}^{\phi}a(t')^{-3}e^{-\eta_*F(t')}\,dF(t').
\]
The ensemble evolution and expansion specify this potential before variation. Realized local release changes $\rho_u$ and $R$ but does not change the prescribed $V_F$ at distant events. In an inhomogeneous element $Q$ remains the prescribed clock drain; replacing it by $\eta_*q(x)$ would define a different vacuum model. The construction assumes a common capacity flow. After shell crossing, one clock current cannot be parallel to every stream. The stream allocation below supplies the accompanying momentum exchange on its positive-stock domain.

Vacuum homogeneity fixes the associated momentum exchange. For $T_\Lambda^{\mu\nu}=-Vg^{\mu\nu}$, conservation assigns the capacity the transverse force density
\[
f_C^\mu=h^\mu{}_{\nu}\nabla^\nu V,\qquad
h^\mu{}_{\nu}=\delta^\mu{}_{\nu}+u^\mu u_\nu.
\]
The choice $V=V_F(\phi)$ makes $f_C^\mu=0$ and preserves geodesic flow. A vacuum growth rate tied to the realized local density generally generates spatial gradients on a clock slice and hence momentum transfer. A density-dependent shutoff at local stock exhaustion therefore requires a spatially varying vacuum and a compatible momentum equation, or a revised exchange sector. The volume-singlet branch retains the uniform-clock drain and its stated local positivity domain.

\paragraph{Exchange between distinct streams.}
\phantomsection\label{sec:multistreamexchange}
The prescribed vacuum can exchange with several streams without requiring
their velocities to coincide. Let $n^\mu$ be the future unit clock vector,
$\nabla^\mu V=-Qn^\mu$, and let
$T_a^{\mu\nu}=\rho_a u_a^\mu u_a^\nu$ describe stream $a$.
For nonnegative weights $w_a$ with $\sum_a w_a=1$, choose
\[
\nabla_\mu T_a^{\mu\nu}=-Qw_a n^\nu.
\]
With $\gamma_a=-u_a\cdot n$, its longitudinal and transverse parts are
\[
\begin{aligned}
\nabla_\mu(\rho_a u_a^\mu)&=-Qw_a\gamma_a,\\
\rho_a u_a^\mu\nabla_\mu u_a^\nu
&=Qw_a(\gamma_a u_a^\nu-n^\nu).
\end{aligned}
\]
The force is orthogonal to $u_a$, and the summed exchange cancels
$\nabla_\mu T_\Lambda^{\mu\nu}=Qn^\nu$ exactly. A moving stream loses
rest energy and receives the corresponding force; it need not remain
geodesic during exchange. At fixed spatial momentum its speed increases
as its rest mass decreases. This is an energy drain.

One positive-domain constitutive choice uses the energy available above
the null limit. Relative to $n$, write
\[
e_a=T_a^{\mu\nu}n_\mu n_\nu,\qquad
j_a^\mu=-h_n{}^\mu{}_{\alpha}T_a^{\alpha\nu}n_\nu,\qquad
A_a=e_a-\sqrt{j_a\cdot j_a},\qquad
w_a=\frac{A_a}{\sum_b A_b}.
\]
For a timelike dust stream, $A_a=\rho_a/(1+|\mathbf v_a|)>0$.
The common-flow limit gives $w_a=\rho_a/\sum_b\rho_b$, recovering the
composition-preserving allocation. Conservation permits other weights;
this choice specifies the extension.

For a source substep at fixed proper volume, let $(E_a,\mathbf p_a)$ be
the stream cell energies and momenta, and put
$A=\sum_a(E_a-|\mathbf p_a|)$. A vacuum-energy increment
$0\le\Delta E_\Lambda<A$ has the exact update
\[
\begin{aligned}
s&=1-\frac{\Delta E_\Lambda}{A},&
E_a^+&=|\mathbf p_a|+s(E_a-|\mathbf p_a|),\\
\mathbf p_a^+&=\mathbf p_a,&
E_\Lambda^+&=E_\Lambda+\Delta E_\Lambda.
\end{aligned}
\]
Every stream stays future timelike, and total energy and spatial momentum
are conserved. Its infinitesimal limit is the displayed exchange law.
Internal unresolved and retained labels use their actual rest-mass survival
factor $\sigma_a=m_a^+/m_a$ in the positive update above; $\sigma_a=s$
only in the cold limit. A common clock-time release rate corresponds to
$\Gamma_a=\gamma_a\Gamma_{\rm clock}$ per stream proper time.
The reserve condition is essential: a local shutoff at exhaustion would
require the spatial-vacuum or exchange completion already specified above.

The same source admits an expanding spacetime example with distinct
velocities. For
$ds^2=-dt^2+a_x^2dx^2+a_y^2(dy^2+dz^2)$, take equal counterstreams,
total number density $N$, particle energy $E$, opposite momenta of magnitude
$p$, stock energy density $C=NE$, and longitudinal pressure $P=Np^2/E$.
In units $c=8\pi G=1$, their equations include
\[
\begin{aligned}
\theta&=H_x+2H_y,&
\dot H_x&=\tfrac12(b+C+2V+P)-\theta H_x,\\
\dot H_y&=\tfrac12(b+C+2V-P)-\theta H_y,&
\dot b&=-\theta b,\qquad \dot N=-\theta N,\\
\dot p&=-H_xp,&
\dot E&=-H_xp^2/E-Q/N,\qquad \dot V=Q.
\end{aligned}
\]
The Einstein constraint
$\mathcal C=2H_xH_y+H_y^2-b-C-V$ obeys
$\dot{\mathcal C}=-2\theta\mathcal C$. Smooth initial data with
$N>0$, $E>|p|$ and $\mathcal C=0$ therefore have a positive-duration
solution on this branch. Controls use the native $\eta_*$ and $\epsilon$,
$b(0)=\epsilon$, $C(0)=1$, $V(0)=F(0)=0$,
$H_x(0)=H_y(0)=\sqrt{(1+\epsilon)/3}$, and unit initial particle mass,
with $p(0)=0$, $0.2$ or $0.7$.
The prescribed $Q(t)$ follows an independent homogeneous existence example
with the same initial densities and $\Gamma_{\rm clock}=H_{\rm ref}$.
The beam solutions remain positive through dimensionless
$t=20$. Their maximum relative constraint residual is below
$9\times10^{-12}$; the zero-momentum control reproduces the reference
FLRW solution to $1.4\times10^{-12}$.
These are smooth expanding examples; bound-galaxy formation solutions remain to be constructed.
The projected stock target, director and support forces, and the depleted
pressure branch still require their common dynamical realization.

\paragraph{The reallocated stock carries the galactic excess.}
Inside galaxies $R$ is the retained stock. Its equilibrium agrees with the transverse response of Section~15 on admissible profiles. On galactic timescales where the cosmological exchange is negligible, each active carrier carries a positive stock obeying:
\[
\boxed{\partial_t\rho_{R,P}+\boldsymbol\nabla\!\cdot(\rho_{R,P}\mathbf v_{R,P})=0,\qquad \rho_{R,P}\ge0,}
\]
Its relaxation target must have the same mass and nonnegative density. For an isolated spherical carrier, write the capped, budget-matched target as $\varrho_{{\rm tar},P}=\chi_P\boldsymbol\nabla\!\cdot\mathbf J_\perp/(4\pi G)$ on its finite carrier domain $\mathcal D_P$, with $\int_{\mathcal D_P}\varrho_{{\rm tar},P}\,d^3x=M_{R,P}>0$. The selected static admissibility prescription is
\[
\boxed{\rho^{\rm adm}_{R,P}
=\frac{M_{R,P}[\varrho_{{\rm tar},P}]_+}
{\int_{\mathcal D_P}[\varrho_{{\rm tar},P}]_+\,d^3x},
\qquad [u]_+=\max(0,u).}
\]
The denominator is positive for a target of positive total mass. This prescription preserves the stock mass and reduces to the original target when that target is nonnegative. Clipping alone would add the absolute mass of the negative region. The projection specifies an equilibrium target. The kinetic and radial-transport constructions below supply positive realizations; their physical selection and common variational coupling remain completion data. In nonspherical systems it applies to the stock contribution after the polarization split below, whose full coupled dynamics must be supplied. The existing auxiliary-action and reciprocity results do not establish that projected dynamics. The ADM mass of an isolated carrier remains $M_b+M_R$.

\paragraph{Spherical admissibility and the size of the correction.}
Before the outer cap, put $s=d\ln M_b/d\ln r$ and $x=\sqrt{GM_b/(a_0r^2)}$. The contact target has enclosed mass $m_{\rm tar}=M_b/(e^x-1)$, hence
\[
\boxed{\frac{\varrho_{\rm tar}}{\rho_b}
=\frac{1}{e^x-1}
\left[1-\frac{s-2}{2s}\frac{x}{1-e^{-x}}\right].}
\]
The expression is nonnegative for $0<s\le2$; vacuum exteriors with constant $M_b$ also give positive response density. A uniform core has $s=3$ and changes sign at $x=5.98490123$, or $g_b=35.81904269\,a_0$. For every nonincreasing spherical baryonic density, $s=3\rho_b/\langle\rho_b\rangle\le3$, so
\[
\varrho_{\rm tar}\ge-\zeta_-\rho_b,
\qquad \zeta_-=1.53103554\times10^{-4},
\]
with the limiting minimum at $x=6.98708$. The director correction vanishes in this symmetry. Outward-rising baryonic density can have $s>3$ and is outside this small bound; the formal $s\to\infty$ minimum is $-0.0324475\rho_b$.

Let $N_-(r)=\int_{B_r}[-\varrho_{\rm tar}]_+d^3x$, $M_-=N_-(\mathcal D_P)$, and $\alpha=M_R/(M_R+M_-)$. The projected enclosed stock is $m_{\rm adm}(r)=\alpha[m_{\rm tar}(r)+N_-(r)]$. Its change in total spherical acceleration is
\[
\frac{\delta g}{g_{\rm tar}}
=\frac{\alpha N_-(r)-(1-\alpha)m_{\rm tar}(r)}
{M_b(r)+m_{\rm tar}(r)}.
\]
For a nonincreasing baryonic density, $N_-(r)\le\zeta_-M_b(r)$, and therefore
\[
\left|\frac{\delta g}{g_{\rm tar}}\right|
\le\max\!\left(\zeta_-,\frac{M_-}{M_R+M_-}\right)
\le\zeta_-
\quad\hbox{if }M_R\ge M_b(\mathcal D_P).
\]
Here $g_{\rm tar}$ includes the same outer budget cap. The correction generally persists outside the negative-density shell until the enclosed budget is restored. For a Plummer source with $M_b=10^{10}M_\odot$, $b=0.2\,$kpc and $M_R=14M_b$, the maximum total-force change is $5.54\times10^{-5}$; at $b=0.1\,$kpc it is $1.04\times10^{-4}$. These are static checks of the prescribed projection. They introduce no fitted coefficient and do not establish its microscopic origin.

\paragraph{Kinetic realization of a positive stock.}
A spherical positive finite-mass target has an explicit stationary kinetic realization in its total gravitational potential. Let $M(r)$ include baryons and stock, $\Phi'=GM(r)/r^2\ge0$, and $v_c^2=r\Phi'$. Uniform orbital phases and tangential directions on circular shells give
\[
f_R(r,v_r,v_t)=\frac{\rho_R(r)}{2\pi v_c(r)}
\delta(v_r)\delta(v_t-v_c(r)),
\]
using $d^3v=2\pi v_t\,dv_t\,dv_r$; at $v_c=0$ use the corresponding zero-velocity measure. This nonnegative distribution integrates to $\rho_R$, has $\sigma_r^2=0$ and $\sigma_\theta^2+\sigma_\phi^2=v_c^2$, and satisfies the radial Jeans equation. Individual radial excursions have
\[
\kappa_r^2=\frac{G[M(r)+rM'(r)]}{r^3}\ge0.
\]
This proves existence of a spherical collisionless equilibrium for a nonnegative target. Collective stability, a smooth isotropic distribution, and formation from cosmological data require further analysis. In particular, clipping a negative Eddington inversion does not establish that its density or equilibrium has been preserved.

The kinetic stress supplies the pressure-like correction needed to support the stock. A collisionless realization evolves the positive stock by Vlasov--Poisson dynamics, including its self-gravity. In this regime its gravitational dynamics have the same form as those of collisionless dark matter. The constitutive target specifies the desired equilibrium; it must not be added as a second full halo on top of the kinetic stock. The remaining director polarization is varied separately.

\paragraph{Selection of the stock equilibrium.}
Three constraints fix what the selection can be.
\begin{itemize}
\item \textbf{Stock-only dynamics has no acceleration scale.} Vlasov--Poisson is invariant under $\mathbf x\to\alpha\mathbf x$, $\mathbf v\to\beta\mathbf v$, $M\to\alpha\beta^2M$, under which accelerations scale by $\beta^2/\alpha$. Contacts with a finite work budget add only a velocity scale. The scale $a_0=\epsilon cH$ therefore enters only through the response cell.
\item \textbf{The coupling cannot hold the stock in place.} A penalty energy $\lambda(8\pi G)^{-1}\!\int|\nabla\Phi_R-\mathbf J|^2$ gives, in spherical symmetry, the baryon acceleration $a_b=g_b+g_R+\lambda(g_R-J)J'(g_b)$. A fractional mismatch $d$ in the stock's squared dispersion then gives
\[
\frac{a_b}{g_{\rm tar}}-1=-d\,J'(g_b)\frac{\lambda}{\lambda-1},
\]
with $J'\simeq5$ at $100$\,kpc in a Milky-Way model. The stock therefore supports itself, and the selection acts slowly and carries no stress at its end state.
\item \textbf{The switch is not a function of acceleration} (Section~18).
\end{itemize}

The response cell has gap $\Theta x$, with $\Theta=k_BT_H$ and $x=\sqrt{g_b/a_0}$, and mean occupation $n_*=1/(e^x-1)$. Its only observable effect is the excess field, one baryonic field unit per quantum. With the stock carrying the response, the occupation is the stock field in baryonic units, $n=g_R/g_b$ (along the director in nonspherical systems). The excess free energy of a cell at fixed gap is $\Theta D$, where
\[
D(n\Vert n_*)=n\ln\frac{n}{n_*}-(1+n)\ln\frac{1+n}{1+n_*}\ \ge0 ,
\]
the relative entropy of geometric distributions. It vanishes only at $n_*$ and has curvature $1/[n(1+n)]$. The stock functional is $F=\int\mathcal N\Theta D\,d^3x$. Its Gaussian limit equals the field mismatch $(8\pi G)^{-1}\!\int|\nabla\Phi_R-\mathbf J|^2$ exactly when
\[
\mathcal N\Theta=\frac{g_b^2\,n_*(1+n_*)}{4\pi G},
\]
that is, when the static susceptibility $\chi_0=n_*(1+n_*)/\Theta$ of Section~16 matches the gravitational compliance of the excess field. With $\Theta=\hbar H_0/2\pi$ this gives cells of 13--60\,fm and about $10^{104}$ of them in a galaxy volume.

For phase-mixed stock described by orbit populations $w_i$, $\partial F/\partial w_i=\langle\psi\rangle_i$ with $\psi=\delta F/\delta\rho_R$. Local detailed balance with the horizon bath requires $k_{i\to j}/k_{j\to i}=\exp[-(F_j-F_i)/\Theta]$, and $|\Delta F|/\Theta=|\int\mathcal N\Delta D|\gg1$ for any change of an orbit. Uphill transitions are therefore suppressed, and a contact proceeds when it lowers the orbit-averaged $\psi$ of the pair. The contacts are equal and opposite impulses along the pair separation, which conserve momentum and angular momentum. A finite work register pays the kinetic change. The cells' free-energy decrease is released to the bath as heat. Each cell holds mean energy $\Theta x\,n_*$, so the reservoir density is $\mathcal N\Theta x n_*$; its exact value is $g_ba_0x^3n_*^2(1+n_*)/(4\pi G)$, with Rayleigh--Jeans expansion $g_ba_0[1-x/2+O(x^3)]/(4\pi G)$, and per unit supported stock it tends to $V_f^2$ in the deep regime, 0.85--0.97\,$V_f^2$ to the budget edge of the galaxy models. Identified with the work register, this energy exists only where the cells are.

On a domain of radius $R$ with stock budget $M_R$, the positivity- and budget-constrained minimiser of the spherical mismatch is the weighted isotonic regression of $r^2J/G$ clipped at $M_R$. It equals the exponential law inside the budget edge $R_{\rm cap}$, where $M_\perp(R_{\rm cap})=M_R$, and leaves a Keplerian exterior. It differs from the normalized positive part above only through the negative-response correction bounded by $\zeta_-$.

Spherical three-dimensional runs test this dynamics. They start from abundance-matched NFW stock with an equilibrium Eddington distribution and static spherical-equivalent baryons, and use one to three contact proposals per orbital period. A pair-kinetic rate proportional to the local stock density and pair speed, run in an independent harness, gives the same end state: dwarf, Milky-Way and massive-spiral models lie within $0.04\,$dex of the law after $12\,$Gyr for two amplitudes six times apart. After 8--12\,Gyr, Milky-Way models lie 0.009--0.032\,dex rms from the exponential law over 5--200\,kpc, against 0.10\,dex without contacts. Dwarf models lie 0.02--0.04\,dex from the law, with the mismatch within 5--10\% of its exact constrained minimum, and massive spirals 0.007--0.05\,dex. The full Bose functional and its Gaussian limit give the same results. The stock gains about $0.15V_f^2$ per unit mass from the register. Momentum and angular momentum are conserved to roundoff, and the energy ledger closes to $10^{-3}V_f^2$. The mismatch decreases overall with fluctuations at the particle shot-noise level. The identities used here (relative-entropy positivity, equipartition matching, the zero-temperature limit, a finite H-theorem, pair-impulse conservation and the constrained minimiser) are checked in the Lean module \texttt{Entropic.ContactSelection}.

The selection rests on four premises: the identification $n=g_R/g_b$; equipartition between the cells and the excess field; a pair-kinetic contact rate proportional to the local carried-stock density and pair speed, with one calibrated amplitude; and the identification of a carrier's records with its condensed baryons, of its support with the region those records dominate, of its carried stock with the stock bound to it, and of the work register with the cell energy (Section~18). The record identification is the retained-carrier notion of the release sector (O.12, O.15): the baryons that carry cells are those whose progenitor history remains recoverable, and shock-heated gas mixed in a common well, whose capacity has been released, carries none. One premise therefore serves both sectors; the support edge remains a separate premise. A microscopic exchange operator, live baryons and cosmological assembly remain to be supplied.

\paragraph{Multistream stock and rotating figures.}
A positive phase-space mass measure $f_R(\mathbf x,\mathbf v,t)$ gives an explicit nonspherical stress closure. Put
\[
\rho_R=\int f_R\,d^3v,\qquad
\rho_R u_i=\int v_i f_R\,d^3v,\qquad
\Pi_{ij}=\int(v_i-u_i)(v_j-u_j)f_R\,d^3v.
\]
Collisionless transport in the common gravitational potential is
\[
\partial_t f_R+\mathbf v\cdot\boldsymbol\nabla_xf_R
-\boldsymbol\nabla\Phi\cdot\boldsymbol\nabla_vf_R=0.
\]
Its zeroth moment gives stock continuity, and its first moment gives
\[
\partial_t(\rho_Ru_i)+\partial_j(\rho_Ru_iu_j+\Pi_{ij})
=-\rho_R\partial_i\Phi.
\]
The velocity stress supplies mechanical support while transport preserves mass and positivity. In the live realization of Section~16, $\Phi$ contains baryons and the stock itself, with a rotationally invariant interaction giving opposite disk and stock torques. This is a selected effective closure for the multistream stress. The constrained single-clock action does not uniquely determine it, and it supplies no relaxation law toward the static RAR target; the contact selection above supplies a conditional one. Finite director and polarization forces must enter the common dynamics through their consistent variation.

A rotating candidate is constructed using
\[
H_J=\frac{|\mathbf v|^2}{2}+\Phi-\boldsymbol\Omega\cdot(\mathbf x\times\mathbf v),
\qquad \boldsymbol\tau_A=\boldsymbol\Omega\times\mathbf L_A.
\]
The second relation is the angular-momentum balance of each steadily rotating component. For any test observable $\psi$, stationarity of a phase measure requires $\int f_R\mathcal L_{\rm rot}\psi=0$. On finite orbit segments with nonnegative weights $w_j$, this yields linear constraints
\[
\sum_jw_j\frac{\psi_j(T_j)-\psi_j(0)}{T_j}=0.
\]
The tilted candidate satisfies 60 mass, torque, spin, and weak-stationarity equalities with 1072 positive orbit weights. Its fitted density mismatch is $0.469\%$; a fourth orbit-time window excluded from the fit gives $4.40\%$ and a $7.49\%$ torque shortfall. These finite constraints support construction of a candidate but do not imply stationarity for every phase-space observable. Section~16 tests its subsequent live evolution.

The angular deformation can exert torque with little radial-force change. For a local scale-free profile $\rho_0=v_h^2/(4\pi Gr^2)$, the source $A\rho_0P_2(\hat{\mathbf n}\cdot\hat{\mathbf e})$ has the particular potential
\[
\Phi_2=-\frac{A v_h^2}{6}P_2(\hat{\mathbf n}\cdot\hat{\mathbf e}).
\]
Its radial derivative vanishes while its angular gradient remains. A capped galaxy adds radial structure and boundary terms; the live force checks measure those effects directly.


\paragraph{Budget normalization.}
The ensemble full-release budget gives
\[
M_R^{\max}=\frac{e^{-\eta_*}}{\epsilon}f_bM_h
\simeq e^{-\eta_*}(1-f_b)M_h.
\]
This assigns the homogeneous cohort's surviving allocation to its Lagrangian baryons. A particular carrier also obeys its local exchange and transport history; the homogeneous factor is not a universal local solution in a void or multistream halo. The isolated-lens calculations in Section~16 used the linear-depletion budget $(1-\eta_*)(1-f_b)M_h$. The exponential budget is larger by $0.0455\%$ at full release. Those likelihood values retain their stated benchmark normalization. Both prescriptions assign the budget to Lagrangian baryons. The bound per present baryon, $C_R=(1-\eta_*)/\epsilon=5.16$, is excluded by isolated-lens weak lensing. The carried part of the budget is limited by the support rule of Section~18, and the support-limited stock matches $\Lambda$CDM halos on the isolated-lens profiles (Section~15). For $M_h=10^{12}M_\odot$ and $M_b=6\times10^{10}M_\odot$, the budget is approximately $14M_b$.

The $u\to R$ exchange is nearly local. For $M_{200}=1$--$1.5\times10^{12}M_\odot$, the released NFW-like stock and the RAR response capped at that stock enclose masses that agree to 5--20\% inside $50\,$kpc. The largest mismatch, 19--31\% of the stock, lies near the cap radius at $120$--$175\,$kpc. The auxiliary field energy is of order $\Phi/c^2\sim10^{-6}$ of the phantom mass, so $R$ is the stock's effective stress and not field energy.

To hold an equilibrium against its own gravity the stock needs a pressure-like stress $T_R^{ij}\neq0$; this is the $O(\sigma_Y^2/c^2)$ correction $\Delta T_C$ above. The response stress splits as $T_\perp=T_R+T_{\rm pol}$. The positive stock carries $T_R$, and the polarization stress $T_{\rm pol}$ of the director sector carries the interface and defect structure of Appendix~N.8.

Because the stock moves with its carrier, it supplies the inertia that the Milky Way's reflex toward the LMC requires (Section~15). The LMC's budget, about $10^{11}M_\odot$, is comparable to a $\Lambda$CDM LMC of $1.5\times10^{11}M_\odot$, whose reflex is $23$--$34\ {\rm km\,s^{-1}}$ in the same crude infall estimate. The moving-LMC calculation summarized in \emph{Empirical Results} evolves an inertial stock and gives a reflex scale near $35\,{\rm km\,s^{-1}}$ over the last gigayear. It uses an approximate pair force and is not yet a solution of the complete reciprocal carrier/director system or a matched reflex likelihood. The stock potential is also the stiffness amplitude of the multi-carrier director completion (Appendix~N.8).

\paragraph{Homogeneous normalization and local domain.}
For $\rho_{\Lambda,i}=0$, define the weighted release integral
\[
I_{\eta,\rm rel}(a)=\int_0^a a'^{-3}e^{-\eta_*F_{\rm rel}(a')}\,dF_{\rm rel}(a').
\]
The density and flatness equations are
\[
\begin{aligned}
\bar\rho_C(a)&=B a^{-3}e^{-\eta_*F_{\rm rel}(a)},\\
\rho_\Lambda(a)&=\eta_* B I_{\eta,\rm rel}(a),\\
\omega_{\Lambda,0}&=\eta_*\omega_c^{\rm wb}I_{\eta,\rm rel}(1),\\
h^2&=\omega_r+\omega_b+\omega_c^{\rm wb}e^{-\eta_*F_{\rm rel}(1)}
+\eta_*\omega_c^{\rm wb}I_{\eta,\rm rel}(1).
\end{aligned}
\]
The free-nucleon/helium convention of O.15 sets $\omega_c^{\rm wb}=5.333\omega_b$; the bare convention uses $\omega_b/\epsilon$. The actual present ratio is $5.333e^{-\eta_*F_{\rm rel}(1)}$, approximately $5.2284$ for the reference history. The reciprocal acceleration--abundance identity refers to the bare pre-release ratio.

Release timing and vacuum accumulation weight the same history differently. By $z=2$, the reference history has reached $F(z=2)/F(z=0)=0.5193$. The cumulative comoving release $\int a^3\bar q\,dt$ has reached $52.17\%$ of its present value, while $\rho_\Lambda$ has reached $96.82\%$. The former fraction is $(1-e^{-\eta_*F(z=2)})/(1-e^{-\eta_*F(z=0)})$; the vacuum integral additionally weights release by $a^{-3}$, favouring early epochs. The comoving drain between $z=1$ and today is
\[
\frac{\Delta(a^3\bar C)_{\rm lost,\,z<1}}{B}
=e^{-\eta_*F(z=1)}-e^{-\eta_*F(z=0)}
=0.0049383.
\]
This is $0.4938\%$ of the initial stock, or $0.5012\%$ of the stock remaining at $z=1$. A small late background drain still requires a local budget check in an expanding underdensity.

Positive internal allocation does not remove the total-stock boundary. Along the clock, with proper-volume element $J$ satisfying $\dot J=\theta J$,
\[
JC=J_iC_i-\int_{t_i}^t J(t')Q(t')\,dt'.
\]
This quantity must remain nonnegative. A prescribed positive drain continues past exhaustion unless its physical domain or a boundary completion is imposed. Under common baryon--capacity flow and conserved baryonic mass, put $\Delta_b=\rho_b/\bar\rho_b$. Then
\[
C=\frac{\omega_c^{\rm wb}}{\omega_b}\rho_b(1-\mathcal B),\qquad
\mathcal B=\eta_*\int\frac{e^{-\eta_*F}\,dF}{\Delta_b},\qquad
\Delta_C=\frac{\Delta_b(1-\mathcal B)}{e^{-\eta_*F}}.
\]
A density floor $\Delta_b\ge d$ throughout release gives
$\mathcal B\le(1-e^{-\eta_*F})/d$. For the reference history $F(1)=0.6633922061$, the numerator is $0.01961844$. A floor above $1.961844\%$ of the mean suffices through today under these assumptions; floors of $0.1$ and $0.2$ bound the loss by $19.62\%$ and $9.81\%$. These are bounds on a density history at the chosen coarse scale. A present tracer density alone does not establish them. Pressure, distinct velocities, and shell crossing require the corresponding local budget and momentum equations.

A nonlinear spherical check locates this boundary for specified histories. Evolve a pressureless, single-stream interior with common baryon--capacity flow, homogeneous radiation and vacuum, and the reference exponential-allocation background. At $a_i=0.01$, set equal baryon and capacity contrasts $\delta_i$, radius ratio $y=b/a=1$, and linear growing-mode initial expansion $y'=-f_i\delta_i/3$, where a prime denotes $d/d\ln a$ and $f_i=0.98071$ includes the radiation correction. Histories ending at $\Delta_C=0.2$ and $0.1$ retain positive capacity throughout; they have lost $4.98\%$ and $8.23\%$ of their initial comoving stock, respectively. The history with $\delta_i\simeq-0.35187$ first exhausts its stock today at $\Delta_b\simeq0.006835$. This endpoint belongs to the specified initial-expansion prescription and release history; it is not a universal density threshold. Doubling the background grid and tightening the integration tolerances shifts the endpoint in $\Delta_b$ by $1.1\times10^{-8}$. Direct density evolution independently reproduces the two positive remnants to better than $1.5\times10^{-7}$ in $\Delta_C$. Evolution of the positive-stock branch ends at exhaustion.

The linear-depletion benchmark defines instead $\bar q=B a^{-3}\dot F$, $a^3\bar C=B(1-\eta_*F)$, and $I_{\rm rel}=\int a^{-3}dF$. Its hydrodynamic integral retains that convention. The exponential allocation has its own matched background-and-perturbation fit in O.12, giving $H_0=68.9389\,\mathrm{km\,s^{-1}\,Mpc^{-1}}$ and $f=0.274245$. For that fitted history, $F(1)=0.6635663$, the surviving comoving fraction is $0.9803765$, the vacuum reaches half its present density at $z=7.9278$, and $\rho_\Lambda(2)/\rho_\Lambda(0)=0.9682403$. The comoving drain since $z=1$ is $0.00493467B$. The preceding void trajectories retain their separately specified reference history. In the standard configuration of Section~20 the linear-depletion direct fit has $H_0=68.458\,\mathrm{km\,s^{-1}\,Mpc^{-1}}$, $F(1)=0.66185$ and $I_{\rm rel}=92.33$, $2.0\%$ below the reference-history integral $94.2116$; the integral fixes $\omega_{\Lambda,0}$, so the fitted and reference histories agree at that level. Since every parcel is released at $a\le1$, some release occurs at $a^3\le F(1)/I_{\rm rel}$, before $z\simeq4.2$ for both (\path{Entropic/ReleaseLedger.lean}).

\paragraph{A constitutive completion before exhaustion.}
\label{sec:depletioncompletion}
The positive-stock domain has an explicit local continuation. In units
$c=1$, choose a constant energy floor $K>V(\phi)$ and put
\[
L=3(K-V)>0,\qquad \lambda\ge L/2,\qquad
P=\lambda X-e(\lambda,\phi)-V(\phi),\qquad
e=\frac\lambda2+\frac{2[L-\lambda]_+^3}{3L^2}.
\]
For $\lambda>L$ this action equals the constrained dust action exactly.
The cubic correction has a $C^2$ join there. It is a selected constitutive
law; the join and positivity conditions admit other functions. On the
regular pressure branch, $r=\lambda/L\in(1/2,1)$ gives
\[
X=e_\lambda=\tfrac12-2(1-r)^2,\quad
v^2=2X=(2r-1)(3-2r),\quad j=\lambda v,
\]
\[
\dot j+\theta j=-V_\phi F(r),\quad
F(r)=-1+6r^2-4r^3=A(r)v^2,\quad
A(r)=\frac{1+2r-2r^2}{3-2r}\in[3/4,1].
\]
The moving threshold contributes to $F$. Consequently
\[
\dot j=-\left[\theta+\frac{V_\phi A(r)j}{\lambda^2}\right]j.
\]
If the flow remains regular and $K-V$ stays bounded away from zero, the
fractional rate is regular at $j=0$. A positive future-directed current
cannot cross zero in finite regular evolution. In the homogeneous
depletion limit, the source clock slows relative to proper time, with
\[
j\longrightarrow0,\qquad v\longrightarrow0,\qquad
\rho_{\rm dark}\longrightarrow K,\qquad p_{\rm dark}\longrightarrow-K.
\]
The constant minimum removes the scalar force at this resting state.
This is a statement about the source field; it supplies no universal
slowing law for atomic clocks in voids.

The identity $e/L-1/3=(2-r)(2r-1)^2/6$ gives an observer-independent
lower energy bound $T_{00}\ge K$. The null-energy contraction is
$\lambda(k^\mu\nabla_\mu\phi)^2\ge0$, and the regular material speed is
\[
c_s^2=\frac{4r(1-r)}{4r(1-r)+2X}\in[0,1].
\]
For spacelike gradients the constrained endpoint gives $P=LX/2-K$.
The reduced $P$ and $P_X$ match at $X=0$, while $P_{XX}$ jumps; interfaces
with a nonzero null gradient require further analysis. The local
inequalities do not prove global evolution through shocks or caustics.

For the fitted exponential background,
\[
V_\infty-V(1)\le B[e^{-\eta_*F(1)}-e^{-\eta_*}],\qquad
K=V(1)+B[e^{-\eta_*F(1)}-e^{-\eta_*}]
\]
is an explicit floor choice. In the saved background, $K=0.3369964$ in
units $3H_{100}^2/(8\pi G)$ and $\max_{a\le1}L/\bar C=0.0299937$.
The homogeneous solution and an open neighborhood of its linear
perturbations therefore remain on the unchanged dust branch. The
completion inherits that cold adiabatic linear likelihood. Deep voids
enter a different pressure regime whose abundance, expansion, and
lensing still require nonlinear calculation. The floor changes their
allowed low-density states and is not an invisible boundary convention.

\paragraph{Finite work, recharge, and warm stress.}
A regular signed exchange power can be chosen as $\mathcal P=v^2 B_c$,
where $B_c$ is bounded and may have either sign. The current equation gains
$B_c/\lambda$ in its fractional rate. Depletion switches off further
withdrawal while allowing recharge on the positive branch. A finite
transition $E\to E'$ toward a positive target $h$ can retain an explicit
work and heat ledger,
\[
Q=\Theta\ln(h'/h),\quad W=E'-E+Q,\quad
U'=U-W,\quad E_B'=E_B+Q.
\]
Then $E+U+E_B$ is conserved. A symmetric edge gate $v(U)v(U-W)$,
defined only for endpoints above the floor, preserves detailed balance
and suppresses inaccessible transitions. Central transitions can retain
angular labels. For a fixed spherical potential their stationary family
has the form $f=a(\ell)h(E,\ell)$. The target, transition rate, and
self-consistent evolving potential remain physical inputs. Rescaling a
detailed-balance generator preserves its equilibrium and changes its
relaxation time, so equilibrium alone cannot select that rate.

A fluid extension stores dissipated energy explicitly. If $R_0(n,\phi)$
is the cold Hamiltonian, $\mu_0=\partial_nR_0$ and $p_0=n\mu_0-R_0$, set
\[
\rho=R_0+e_T,\quad e_T=z n^\gamma\ge0,\quad
p=p_0+(\gamma-1)e_T,\quad 1<\gamma\le2.
\]
With $h_0=n\mu_0$, $h_T=\gamma e_T$, its adiabatic speed is
\[
c_s^2=\frac{h_0c_0^2+h_T(\gamma-1)}{h_0+h_T}\in[0,1].
\]
The floor persists. The bounded gate
$g_c=n\mu/(\lambda_0+h_T)\in[0,1]$ and
$\mathcal P=g_cB_c$ retain a regular fractional drain and reduce to
$v_0^2B_c$ when cold. The tested fluid chooses $\gamma=4/3$; the value
is an effective assumption. In planar compression tests from 128 to
2048 cells, mean energy drift stays below $4.9\times10^{-15}$ and the
current remains positive. Mean generated thermal energy decreases from
$0.0641$ to $0.0140$ under refinement, so its physical dissipation has
not converged. The cold calculation has a late energy deficit of
$0.00319$, even though its early smooth solution converges at second
order. The conservative warm formulation closes the energy ledger in
these tests; nonlinear interfaces, entropy production, and a coupled
multistream stress still require validation. The cluster realization in
Section~18 tests reciprocal finite-work particle contact separately.


\paragraph{Sign and effective equation of state.}
Under the one-signed history orientation of Section~22, the homogeneous
marked share accumulates monotonically. An observer fitting $\rho_\Lambda(a)$
as a separately conserved component infers
\[
\boxed{
 w_{\rm eff}(a)=-1-
 \frac{a^{-3}e^{-\eta_*F_{\rm rel}}\,dF_{\rm rel}/d\ln a}
 {3I_{\eta,\rm rel}(a)}.}
\]
Thus active one-signed release gives an effective $w<-1$ and approaches
$-1$ as release completes; a persistent $w_{\rm eff}>-1$ for this same reconstructed vacuum
component would require the bookkeeping flow to reverse. A conserved-CDM
$w_0,w_a$ fit is not that reconstruction.

\paragraph{Clock and metric perturbations.}
Write the scalar metric and clock as
\[
ds^2=-(1+2\Psi)dt^2+a^2(1-2\Phi)d\mathbf x^2,
\qquad
\phi=t+\delta\phi.
\]
The constraint and potential give
\[
\boxed{\delta\dot\phi=\Psi,
\qquad \delta\rho_\Lambda=\delta V_F=\dot V_F\,\delta\phi.}
\]
Thus $\delta\rho_\Lambda=0$ on capacity-clock slices, while another gauge retains the clock perturbation required by covariance. Define
\[
Q:=\dot V_F=\eta_*\bar q,
\qquad
\gamma_C=\frac{Q}{\bar\rho_C},
\qquad
\Theta_C=\frac1a\nabla\!\cdot\mathbf v_C,
\qquad
\mathbf v_C=-\frac1a\nabla\delta\phi.
\]
Linearizing the action gives
\[
\boxed{
\dot\delta_C=-\Theta_C+3\dot\Phi
+\gamma_C\delta_C-\gamma_C\Psi
-\frac{\dot Q}{\bar\rho_C}\,\delta\phi,}
\]
\[
\boxed{
\dot{\mathbf v}_C+H\mathbf v_C
=-\frac1a\nabla\Psi
+O\!\left(\frac{\sigma_Y^2}{c^2},k^2R_Y^2\right).}
\]
The absence of a $-\gamma_C\mathbf v_C$ term follows from the geodesic capacity clock: the exchange current is parallel to $u^\mu$ and transfers energy without transverse momentum. In capacity-comoving synchronous gauge, $\delta\phi=\Psi=0$, and the same system reduces to
\[
\boxed{\dot\delta_C=-\tfrac12\dot h_s+\gamma_C\delta_C,
\qquad \delta\rho_\Lambda=0.}
\]
The action fixes the relativistic scalar perturbations for a supplied $V_F$ on the positive-density domain stated above. It introduces no independent sound speed, anisotropic stress, or $k^2$-enhanced force. A full likelihood still requires the physical release history and the standard radiation, baryon, and metric sectors evolved with these equations.

\paragraph{Conserved-fluid representation of the perturbations.}
Choose a reference dust density $\rho_{\rm ref}=\rho_{{\rm ref},0}a^{-3}$ on the same geodesic flow and define
\[
\rho_D=C-\rho_{\rm ref}+V_F,\qquad p_D=-V_F,
\qquad \rho_D+p_D=C-\rho_{\rm ref}>0.
\]
The combined capacity--vacuum stress is exactly the sum of this conserved fluid and the reference dust. On common-clock slices $\delta p_D=0$, so its rest-frame sound speed is zero. Choosing the reference abundance below the minimum comoving capacity keeps its enthalpy positive. This representation implements the direct scalar perturbations for either supplied background law. Changing the reference split leaves the total stress and its observables invariant when the initial perturbations are matched. The finite numerical checks and matched linear-lensing likelihood are stated in O.12.

\paragraph{Radial adjustment with retained angular information.}
A stock can adjust its radial profile while preserving selected angular data. At finite resolution let $w_{ia}\ge0$ be stock in radial state $i$ and angular sector $a$, with $m_a=\sum_iw_{ia}$. For a normalized positive target $p_i$, the update
\[
w_{ia}^{+}=(1-s)w_{ia}+s p_i m_a,\qquad 0\le s\le1,
\]
preserves every $m_a$, preserves positivity and total mass, and contracts $w_{ia}-p_i m_a$ by $1-s$. A local radial realization is
\[
\dot w_{ia}=\sum_jc_{ij}\left(\frac{w_{ja}}{p_j}-\frac{w_{ia}}{p_i}\right),
\qquad c_{ij}=c_{ji}\ge0.
\]
For a fixed target and no flux at the radial boundaries, put $x_{ia}=w_{ia}/(m_ap_i)$ for $m_a>0$. Its relative entropy obeys
\[
\frac{d}{dt}\sum_{i,a}w_{ia}\ln x_{ia}
=-\frac12\sum_{i,j,a}m_ac_{ij}(x_{ia}-x_{ja})
(\ln x_{ia}-\ln x_{ja})\le0.
\]
On a connected radial graph the stationary family is $w_{ia}=p_im_a$. The radial part of a quadrupole can therefore relax while its angular marginal is retained. A radial flux realization is
\[
\partial_t\rho_R=\frac1{r^2}\partial_r\!\left[
r^2D\rho_*\partial_r(\rho_R/\rho_*)\right],
\]
with no boundary flux on each ray. This is a coarse relaxation model; a finite propagation completion is needed outside its transport regime.

A possible microscopic contact uses one-body central jump operators $A(r,p_r,L^2)$ that commute with each $L_k$. Such a contact preserves angular-momentum observables while allowing energy redistribution. Rotational symmetry or conservation of total angular momentum alone does not supply this stronger selection rule, and the record-projector identities do not establish it for halo orientation. In a spherical orbital description, diffusion in energy at fixed $L$ can be written
\[
\partial_tw(E,L)=\partial_E\!\left[D_Ew_*(E,L)
\partial_E(w/w_*)\right].
\]
The no-flux boundary preserves each $L$-bin mass. The finite-work contact above supplies a consistent mechanical ledger; physical selection of $w_*$ remains open. If a retained angular-sum map is $P_A$ and transition incidence is $B_A$ with $P_AB_A=0$, a noise covariance $N=B_AWB_A^\dagger$, $W\ge0$, obeys $P_ANP_A^\dagger=0$. The scalar quantum thermal factor of N.8 preserves this null space. This proves compatibility of protected observables with noise; it does not fix a physical bath spectrum or energy supply.

Two numerical constructions illustrate the permitted separation. With a fixed spherical baryonic profile, a fixed total stock budget, a chosen outer taper, and an instantaneous RAR radial target evolving with $a_0\propto H$, a chosen $0.5\,$Gyr adjustment time gives a present stock-force discrepancy of at most $2.21\%$ over $5$--$80\,$kpc from formation at $z=2.5$; the largest earlier discrepancy is $12.89\%$. Separately, positive orbit weights at fixed initial angular-momentum-bin masses realize a present spherical endpoint with at most $1.226\%$ force error and $1.573\%$ radial mass $L^1$ error under independent quadrature. The optimizer reached its iteration limit, so this is a feasible witness. Its mechanical energy changes by about $1.94\times10^{-7}$ of the stock rest energy. The density evolution and orbital endpoint are separate constructions. The finite-work contact above supplies an explicit effective transition and energy ledger. Its physical selection, target, and coupled nonspherical equilibrium remain to be established. Gravitational torques can change orientation even when the direct contact torque vanishes.

\subsection*{O.4 What fixes $\eta$}

The Hilbert-space dimensions do not: the isometry fixes information
routing, not energy allocation, so $\eta=9/16$ is not implied. The
clean definition requires the infrared energy operator of the marked
register. Let $U_{\rm rel}$ be the marked source-clock release map and
$\Pi_{\rm mark}$ the marked projector; then
\[
\eta=\frac{\operatorname{Tr}\!\left[U_{\rm rel}\rho_{\rm comm}U_{\rm rel}^\dagger\,
\Pi_{\rm mark}H_{\rm out}\Pi_{\rm mark}\right]}
{\operatorname{Tr}\!\left[\rho_{\rm comm}H_{\rm comm}\right]}.
\]
Write the physical operators as
\[
H_{\rm comm}=\varepsilon_K\widehat K^2,
\qquad
H_{\rm rec}=\frac{\varepsilon_\xi}{2}
\sum_{a=1}^3\widehat\xi_a^2.
\]
Appendix~O.13 then gives the exact conditional relation
\[
\boxed{
\eta=\frac{\varepsilon_\xi}{\varepsilon_K}
\frac{3/2}{\langle K^2\rangle_{\eta_*}}
=\frac{\varepsilon_\xi}{\varepsilon_K}\eta_*.}
\]
The selected common conversion normalization $\varepsilon_\xi=\varepsilon_K$
therefore fixes $\eta=\eta_*$ when the energetic ledger debits the
post-refresh closure state. H.11 makes independence from the renewed $K^2$
register exact under that ordering. The common conversion and ledger ordering
are physical premises; Hilbert-space dimensions alone do not supply them. The
marked fugacity $\zeta_*=0.00512$ belongs to the charged recurrence and is not
a cosmological event probability.

\subsection*{O.5 The initial constant}

The accumulation law determines changes in $\rho_\Lambda$ and leaves an initial integration constant. At fixed primitive event count $N$, a factor $e^{-CN}$ cancels from normalized probabilities. The centering prescription of Section~25 therefore removes sensitivity of those probabilities to an extensive microscopic energy origin.

This cancellation does not determine the gravitational integration constant. At fixed four-volume, actions differing by a volume term have the same normalized fixed-volume weights. Their equivalence on that restricted measure leaves the conjugate constant unmeasured. A common geometry--capacity substrate also permits several local operators in its effective action; identifying their substrate does not fix their coefficients.

The cosmological branch adopts the initial matching condition
\[
\boxed{\rho_{\Lambda,i}=0}
\]
for the empty equilibrium state and then evolves the O.3 release law. A derivation of this value requires a boundary condition or a microscopic prescription that fixes the gravitational volume coefficient. The finite-regulator centering identity remains exact, while the initial zero is a separate physical matching condition.

\subsection*{O.6 Commitment epoch and the recruitment reading}

Two readings of the committed initial conditions in Section~20 are possible. The
recruitment narrative (allocation
$\rho_c=\rho_b/\epsilon$ pointwise at a recombination-era commitment
epoch) gives the committed field the Silk-damped baryon transfer
function. A first-crossing computation on that field returns
$\sigma_0\simeq0.34$ for the unsmoothed committed field today under
standard growth, with first-caustic probability
$P(\lambda_1>1/\sigma_0)<5\times10^{-8}$: no caustics, no dark energy,
and no cosmic web. This registered negative result falsifies the
late, pointwise-tracking reading. The Einstein--Boltzmann check of
Section~20 already used a standard cold component, which is consistent
only with commitment completing before the acoustic evolution of the
fitted modes: primordial adiabaticity, then cold propagation. The CMB
sector and the dark-energy sector therefore both favor early
(pre-acoustic) commitment, and the recombination-era saturation estimate becomes a consistency condition at that earlier epoch. Appendix~O.15 supplies the formation epoch.

\subsection*{O.7 Identity-free source labels cannot record stream assignment}

Per substrate tick \(\tau_*=L_*/c\), every unpinned closure register is
replaced by an independent draw,
\[
p_*(b)=Z_*^{-1}e^{-\eta_*K^2(b)},
\qquad
K_*(b,b')=p_*(b'),
\]
and the old register is exported to history.  A committed allocation
\(c_{A,c}\) carries a Lagrangian source label \(A\) and a continuous mass
weight, but no additional particle or stream identity.  Its local
environment can therefore record occupancy and the presence of the pinned
source, not a permutation assigning identical incoming mass elements to
outgoing streams.

To state this precisely, let \(i\) and \(j\) denote two fine stream
descriptions of one source that differ only by which past occupancy is
declared to continue into which present occupancy.  The exported
environmental record is a function of the occupancy history
\(\{b_c(n)\}\).  Under the identity-free source convention the two
descriptions are the same physical history, so
\[
\boxed{
|e_i\rangle=|e_j\rangle
\quad\Longrightarrow\quad
\gamma_{ij}=\langle e_j|e_i\rangle=1.}
\]
Consequently
\[
\boxed{
C_A^{\rm rec}
=\sum_{ij}p_ip_j|\gamma_{ij}|^2=1,
\qquad
f_A^{\rm rec}=0}
\]
for every fold and at every well depth.  A density record can determine
the net flux, but not a multivalued velocity whose branches differ only by
an exchange of identity-free mass.  Thus the committed capacity sector
cannot write a which-stream record.  This is an exact consequence of the
declared source ontology and the one-register update; it would fail if an
additional persistent identity label were added.

\subsection*{O.8 Cosmology excludes a persistent stream identity on this branch}

The alternative fails cosmologically.  Suppose a committed allocation carried
a persistent identity that neighboring records could resolve.  Two
streams with relative velocity \(\Delta v\) separate by one primitive
length after
\[
t_{\rm sep}=\frac{L_*}{\Delta v}
=1.62\times10^{-39}\ {\rm s}
\left(\frac{10\ {\rm km\,s^{-1}}}{\Delta v}\right).
\]
Let $p$ be the unconditioned neighboring marginal and $q$ the pinned marginal of P.3. Their relative entropy is $D_{\rm KL}(q\Vert p)\simeq0.1499$ nats. For record outputs with amplitudes $\sqrt p$ and $\sqrt q$, the single-tick overlap and its decay exponent are
\[
B=\sum_b\sqrt{p_bq_b}\simeq0.92792003,
\qquad
D_{1/2}=-2\ln B\simeq0.14961945.
\]
Independent outputs then give $|\gamma_n|=B^n=\exp(-nD_{1/2}/2)$. The threshold $|\gamma_n|<e^{-5}$ requires
\[
n>\frac{10}{D_{1/2}}\simeq66.84,
\qquad
t_{\gamma<e^{-5}}\simeq67\tau_*\simeq3.6\times10^{-42}\ {\rm s}.
\]
The overlap satisfies $B\ge\exp[-D_{\rm KL}(q\Vert p)/2]$. At this recording rate, persistent identity makes every microscopic fold effectively fully recorded.

That limit fails the early-release bound.  At \(z=15\), the stated
first-crossing estimate for \(M>10^6M_\odot\) gives
\(\delta F\simeq5.5\times10^{-3}\), whereas requiring the induced
change in the cosmological density to remain below \(0.01\) gives
\[
\delta F<
\frac{0.01}{0.00783(1+z)^3}
=3.1\times10^{-4}
\qquad (z=15).
\]
The identity-carrying branch overshoots by a factor of about \(18\).
This condition excludes the identity-carrying cosmological branch, given the
displayed first-crossing estimate and the absence of a separate compensating
suppression. Within the minimal branch it selects identity-free committed
allocations, after which O.7 applies.

The completed \(48^3\) phase-relaxation harness serves only as a
protocol check.  It sets \(\gamma_{ij}=1\), heals a smooth phase bump by
construction, and types the proposed gate into its marked output.  It
does not test stream recording, the trigger, or the post-refresh ledger.
Those become tests only in an ungated multibranch source simulation.

\subsection*{O.9 The trigger cannot be a local capacity gate}

The displayed capacity action has one closure register per cell.  When
allocations from different sources \(A\ne A'\) occupy the same cell, its
local register can either preserve the source labels, release an
allocation without a potential selector, or merge the coarse labels.
The last operation is Many-Pasts coarsening: it removes source
distinguishability rather than writing a new which-stream bit.  The
potential enters the existing vertex smoothly through
\[
\delta S=-2S_\infty\frac{\Phi}{c^2};
\]
the action contains no local source-resolving operator whose spectrum
produces the \(O(10^{-9})\) photoheating threshold.  Therefore a sharp
retention switch is not generated inside the present one-register
capacity action.  In this minimal completion the supporting condition
must live in the baryonic carrier.

The two roles sometimes assigned to \(a_0\) must remain separate.
The running relation
\[
a_0(z)=\epsilon cH(z)
\]
governs the transverse thermal response and its redshift test; source-clock
release is governed by loss of the baryonic carrier and contains no
\(a_0\).

Three candidate triggers are thereby excluded.

\begin{enumerate}
\item A generic three-stream fold is not recorded by identity-free
capacity.  Fold multipliers, cell-wise ($0.588$) or source-wide ($0.625$),
therefore do not enter the cosmological normalization.
\item A single carrier atom or ionization state cannot be the supporting
condition.  If reionization rerouted the intergalactic allocation, the
formal increments would be
\[
\Delta\Omega_\Lambda\simeq
0.00783(1+z_{\rm re})^3x_{\rm ion}
=2.2,\ 4.4,\ 6.6
\]
for \(z_{\rm re}=6,\,7.7,\,9\), respectively, already excluding that
reading by factors of several.
\item Ordinary entropy production, star formation, metallicity change,
and environmental measurement records do not by themselves invoke the
marked drain.  P.7a restricts it to loss of the particular baryonic
source-clock carrier.
\end{enumerate}

Horizon contact gives \(a_0(z)\propto H(z)\). A trigger tied to this scale
evolves with the Hubble rate.

\subsection*{O.10 Collective trigger and the information observable}

What survives is collective.  The clock \(\theta_A\) belongs to source
\(A\) as a distinct Lagrangian baryonic system.  Its pin is supported
while matter records continue to distinguish that progenitor.  The
supporting condition fails when gas from several sources is irreversibly
mixed into one retained bound object and the baryonic record no longer
resolves progenitor membership.  The committed allocation is still
labelled \(A\), is collisionless, and did not itself mix; it is therefore
orphaned when its clock carrier ceases to exist as a distinguishable
matter subsystem.

This makes the required astrophysical quantity precise.  Label each gas
resolution element before assembly by its progenitor \(A\).  At a later
time divide the bound gas into resolved regions \(Y_t\), and let
\(x_A(y)=P(A\mid Y_t=y)\).  Define
\[
H(A)=-\sum_Ap_A\ln p_A,
\qquad
H(A\mid Y_t)
=-\int dP(y)\sum_Ax_A(y)\ln x_A(y),
\]
and
\[
\boxed{
\epsilon_{\rm mix}(t)
:=\frac{H(A\mid Y_t)}{H(A)}
=1-\frac{I(A;Y_t)}{H(A)}.}
\]
Thus \(\epsilon_{\rm mix}=0\) when the resolved baryonic state still
determines the progenitor, and \(\epsilon_{\rm mix}=1\) when it retains
no recoverable progenitor information.  This is not a metallicity-mixing
fraction or a star-formation efficiency.

Mixing becomes durable only in a common well that retains the participating
gas.  The sharp astrophysical proxy is
\[
V_{\rm thr}(z)\simeq
\begin{cases}
17\ {\rm km\,s^{-1}},&\text{before reionization},\\
30\ {\rm km\,s^{-1}},&\text{after reionization},
\end{cases}
\]
representing internal atomic-cooling heating before reionization and the
photoheated ultraviolet background afterward.  These velocities are
matter-sector inputs, not capacity constants
\citep{Greif2008FirstGalaxies,BarkanaLoeb1999}.

Let \(dF_Z(t_f;V_{\rm thr})\) be the mass-weighted first-assembly measure
for systems satisfying that collective retention condition.  First assembly
starts the common-well mixing clock; it is not itself the carrier-failure
event. Define the cumulative mixing kernel as the fraction of source-record
stock whose progenitor identity is no longer recoverable after delay $\tau$,
\[
\boxed{
K_{\rm mix}(\tau;M,z)
:=1-\frac{I(A;Y_\tau)}{H(A)}.}
\]
The release history is then the delayed convolution
\[
\boxed{
F_{\rm rel}(t)=\int^tK_{\rm mix}(t-t_f;M,z)\,dF_Z(t_f;V_{\rm thr}),
\qquad
\omega_{\Lambda,0}=\eta_*\frac{\omega_b}{\epsilon}I_{\eta,\rm rel}(1),}
\]
with
\[
\boxed{I_{\eta,\rm rel}(1):=\int_0^1a^{-3}e^{-\eta_*F_{\rm rel}(a)}\,dF_{\rm rel}(a).}
\]
The local form $dF_{\rm rel}=\epsilon_{\rm mix}(a)dF_Z$ is the short-delay
limit of this convolution. Finite-time turbulent mixing shifts release to
later, less strongly $a^{-3}$-weighted epochs. The cosmological normalization is
computed from the tracer-derived $K_{\rm mix}$.

\subsection*{O.11 Full rerouting at carrier failure}

The nine-state Gram matrix fixes the marked amplitude \(1/3=1/\sqrt9\);
it does not fix how much of a failed source allocation changes stock.
That coefficient follows instead from renewal.  For the stationary memory
family
\[
K_a(b,b')
=a\,\delta_{bb'}+(1-a)p_*(b'),
\]
faithful full-support resolution gives
\[
H(B'\mid B)
=g_{\rm share,eff}-I(B;B')
<g_{\rm share,eff}
\qquad(a>0),
\]
so H.4 admits only \(a=0\).  A stable mark obeys
\[
\nabla_\mu J^\mu_{c,A}
=\Gamma_{{\rm form},A}-\Gamma_{{\rm rel},A}=0
\]
until its supporting baryonic condition fails.  At such a failure the
only renewal-compatible stock update is a complete fresh draw:
\[
\boxed{
\begin{gathered}
\mathsf E_A
=\mathbf1[\text{carrier of }A\text{ has lost progenitor identity}],\\
f_{\rm reroute}=1\qquad\text{for each failed carrier}.
\end{gathered}}
\]
The coefficient is therefore unity within the selected renewal branch for
every elementary record block whose carrier fails. At coarse resolution the
fraction of such stock is measured by $K_{\rm mix}$; no additional microscopic
routing fraction is available.


\subsection*{O.12 First-entry assembly history and hydrodynamic mixing kernel}



The carrier-loss observable is
\[
\epsilon_{\rm mix}(t)=1-\frac{I(A;Y_t)}{H(A)}.
\]
This quantity is used as a fractional stock law. Let
$\mathcal L=H(A)-I(A;Y_t)$ be the amount of progenitor-record information
absent from the retained matter algebra. Mass--Entropy Equivalence assigns the
same mass-equivalent cost to equal independent record increments. Continuity,
additivity under a partition of the source record, and the normalization that
complete loss releases the complete source-labelled stock then fix
\[
\frac{E_{\rm rel}}{E_{\rm stock}}=\frac{\mathcal L}{H(A)}
=\epsilon_{\rm mix}.
\]
Thus $\epsilon_{\rm mix}$ prices the fraction of record stock rerouted at its common current energy scale. Under the composition-preserving drain, the cumulative fraction of initial energy converted is obtained by integrating the surviving stock as in O.3; it is
not assumed to be the probability of a single all-or-nothing erasure. Faithful
renewal acts with unit coefficient on each released elementary block. The
Poisson estimate
$1-\exp[-N_{\rm eff}M_{\rm mix}/(f_bM_h)]=0.8529$ measures geometrical
coverage only. A parcel covered by one progenitor's mixing region can still
retain essentially complete progenitor information, so coverage is not the
observable used in the release law.

\paragraph{Absorbing first-entry history.}
The assembly history is fixed by an excursion-set first-entry construction
rather than by taking a running maximum of an instantaneous collapsed
fraction. The linear variance $S(M)=\sigma^2(M)$ is computed from the smooth
Eisenstein--Hu baryon-suppressed transfer function \citep{EisensteinHu1998} normalized to
$\sigma_8=0.810$. The retained-well threshold is
\[
V_{\rm crit}(z)=
\begin{cases}
17\ {\rm km\,s^{-1}},&z>z_{\rm re},\\
30\ {\rm km\,s^{-1}},&z<z_{\rm re},
\end{cases}
\qquad z_{\rm re}=7.66,
\]
with the standard virial relation \citep{BryanNorman1998} used to obtain $M_{\rm crit}(z)$; the two thresholds are the atomic-cooling scale before reionization \citep{BarkanaLoeb1999,Greif2008FirstGalaxies} and the photoheated-gas retention scale found in ultraviolet-background simulations afterwards \citep{ThoulWeinberg1996,OkamotoGaoTheuns2008}. Both thresholds measure retention of gas that can cool; the normalization test below takes the pre-reionization gate from feedback-resolved simulations. Define
\[
B(z)=\frac{\delta_c}{D(z)},\qquad
S_{\rm crit}(z)=\sigma^2[M_{\rm crit}(z)],
\]
and for a Markov excursion-set trajectory let
$\mathcal M(S)=\max_{0<s<S}\delta(s)$. The cumulative fraction that has ever
entered an eligible retained well by redshift $z$ is
\[
\boxed{
F_Z(z)=P\!\left[
\max_{z'\ge z}
\left\{\mathcal M[S_{\rm crit}(z')]-B(z')\right\}\ge0
\right].}
\]
The crossing is absorbing: the rise of the retention threshold at
reionization cannot ``unassemble'' a trajectory that entered an eligible
well earlier. The saved $3\times10^5$-trajectory implementation with
Brownian-bridge crossing corrections gives $F_Z=0.0694$ at $z=6$,
$0.2411$ at $z=3$, and $0.6700$ today. This fixed reference assembly
table is the input to the current likelihood; a self-consistent assembly
history across fitted cosmologies remains to be calculated.

\paragraph{Finite mixing delay.}
First entry starts the common-well mixing clock. It does not cause
instantaneous release. The correct release history is the convolution
\[
\boxed{
F_{\rm rel}(t)=\int^t
K_{\rm mix}(t-t_f;M,z)\,dF_Z(t_f).}
\]
The local product $dF_{\rm rel}=\epsilon_{\rm mix}\,dF_Z$ is its short-delay
limit. When the kernel depends on $M$ and entry redshift, the displayed
convolution abbreviates an integral against the corresponding mass-weighted
joint first-entry measure.

\paragraph{Monotonicity for a closed mixing cohort.}
The passive-tracer observable is a cumulative law for a closed-cohort version
of the dynamics used to measure it. Let $\rho$ obey mass continuity and let
$c_A$ be the local mass
fraction carrying progenitor label $A$:
\[
\partial_t\rho+\nabla\!\cdot(\rho\mathbf v)=0,
\qquad
\partial_t(\rho c_A)+\nabla\!\cdot(\rho c_A\mathbf v)
=\nabla\!\cdot(\rho\mathbf D\nabla c_A),
\]
where the same positive-semidefinite diffusion tensor $\mathbf D$ acts on
every label. For a closed Lagrangian cohort with periodic or no-flux boundary,
the label weights
\[
\pi_A=\frac1{M_{\rm gas}}\int\rho c_A\,dV
\]
are constant. The mutual information between label and position is
\[
I(A;X)=\frac1{M_{\rm gas}}
\sum_A\int\rho c_A\ln\frac{c_A}{\pi_A}\,dV.
\]
Mass conservation, integration by parts, and the common diffusion operator give
\[
\boxed{
\frac{dI(A;X)}{dt}
=-\frac1{M_{\rm gas}}\sum_A\int
\frac{\rho}{c_A}(\nabla c_A)^T\mathbf D(\nabla c_A)\,dV\leq0.}
\]
Consequently
\[
\boxed{K_{\rm mix}^{X}=1-\frac{I(A;X)}{H(A)}\in[0,1],
\qquad \dot K_{\rm mix}^{X}\geq0}
\]
whenever the initial labels are fully resolved. The limiting value can remain
below one if the flow does not mix completely. The theorem applies to a fixed
retained observable and a closed material cohort. Eulerian boundary fluxes,
label-dependent diffusion, source creation, selective escape, or a changing
coarse partition add terms whose sign must be measured in the cosmological
zoom. Even a fixed coarse partition need not inherit monotonicity from the
resolved observable: on a four-site periodic cycle with labels
$(0,0,1,1)$ and two fixed adjacent-site bins, successive reversible one-site
translations give normalized information losses $0,1,0$. The physical
$K_{\rm mix}(Y)$ must therefore be measured with the actual retained
observable and checked for the cumulative behavior required by the release law.

\paragraph{Complete-reset audit.}
The information-stock law can be distinguished from the probability of a
literal input-independent reset. For a measured channel $\mathcal K(y|A)$,
define
\[
c(y)=\min_A\mathcal K(y|A),
\qquad
e(\mathcal K)=\sum_yc(y).
\]
For $0<e<1$ the exact decomposition
\[
\mathcal K=e\,r+(1-e)\widetilde{\mathcal K},
\qquad
r(y)=\frac{c(y)}e,
\qquad
\widetilde{\mathcal K}(y|A)
=\frac{\mathcal K(y|A)-c(y)}{1-e}
\]
shows that $e(\mathcal K)$ is the largest input-independent reset component
compatible with the channel. It is monotone under further processing and obeys
\[
e(\mathcal K\mathcal L)\geq e(\mathcal K),
\qquad
e(\mathcal K)\leq\frac{H(A|Y)}{H(A)}=K_{\rm mix}.
\]
The inequality can be strict. A binary symmetric channel with flip probability
$0.1$ has $K_{\rm mix}=0.468996$ and $e=0.2$, and the same channel can be
implemented as a random reversible bit flip with no physical reset. A measured
stochastic channel therefore does not determine its microscopic instrument.
The cosmological branch uses additive information-stock pricing; interpreting
$K_{\rm mix}$ as a complete-reset probability would be a different physical
model.

\paragraph{Fitted mixing kernel.}
The likelihood uses a pure-exponential approximation to unresolved transport,
\[
K_{\rm mix}=1-e^{-\tau/[f\,t_{\rm dyn}(z_f)]},\qquad
t_{\rm dyn}=\frac{R_{\rm vir}}{V_{\rm vir}}
=\frac1{H(z_f)\sqrt{\Delta_c(z_f)/2}},
\]
where $z_f$ is first entry and $\Delta_c$ is the virial overdensity
\citep{BryanNorman1998}. The dimensionless $f$ is fitted; the
microstate count fixes $\eta_*$ but does not fix the mixing time,
retention thresholds, or label resolution. The measured kernel below
is compared through the full convolution.

\paragraph{Background convention for the likelihoods.}
The primary direct-perturbation likelihood below uses the composition-preserving allocation $a^3\bar C=B e^{-\eta_*F}$ and its weighted $I_{\eta,\rm rel}$. The approximate PPF likelihood and hydrodynamic integral comparisons use the linear-depletion benchmark $a^3\bar C=B(1-\eta_*F)$ and $I_{\rm rel}=\int a^{-3}dF$. Their numerical integrals retain that calibration.

\paragraph{The release integral in the PPF benchmark.}
With the abundance tied at the formation value, $h$ follows from $\omega_b$ and $I_{\rm rel}$, so the acoustic scale of the CMB fixes the release integral. With the O.12 history at the reference gates and $f$ free, the Planck 2018 high-$\ell$ TT,TE,EE spectra, DESI DR2 BAO, and DES-Dovekie with a free supernova offset below $z=0.1$ give $\chi^2=2227.33$ with $A_{L,{\rm eff}}=0.979$, against $2226.77$ for $\Lambda$CDM with one more free parameter, at $f=0.28\pm0.03$, $I_{\rm rel}=94.2$, $\Omega_{\Lambda,0}=0.707$, and $h=0.689$. Fixing $f=0.5$ or $0.75$ raises $\chi^2$ by $45$ or $189$. Within this approximate release-model pipeline the data constrain $I_{\rm rel}=94.2\pm0.8$. The gates move the required kernel strongly: $f=0.28$ at the reference gates, $0.63$ with a pre-reionization gate of $15\ {\rm km\,s^{-1}}$, $0.51$ with a post-reionization gate of $25\ {\rm km\,s^{-1}}$, and $0.12$--$0.13$ with a post-reionization gate of $35\ {\rm km\,s^{-1}}$ or $z_{\rm re}=8.4$.

The supernova archive in this pipeline is DES-Dovekie, with the absolute magnitude and a separate offset for the 201 supernovae below $z=0.1$ marginalized in both models. The background
is implemented as constant CDM plus a PPF effective fluid, with fixed
$A_{L,{\rm eff}}=0.979$. The assembly table is held at the reference
cosmology. This fit does not directly evolve the interacting perturbations
or include a joint CMB-lensing reconstruction likelihood.

\paragraph{Direct perturbations and matched linear lensing.}
The conserved-fluid representation in O.3 implements the geodesic capacity clock in a linear Einstein--Boltzmann calculation. The reference dust and residual fluid share their initial velocity and total perturbation. Disabling late radiation truncation is necessary to retain invariance under the arbitrary reference split at high wavenumber. Changing the reference fraction from $0.80$ to $0.95$ changes sampled Weyl power by at most $5.72\times10^{-7}$ and lensing power by $1.16\times10^{-7}$ over $10\le L\le1000$. An independently integrated subhorizon growth system reproduces the Weyl-power ratio within $9.2\times10^{-7}$.

One harness evaluates both models on the same likelihood: Planck 2018 plik\_lite TT,TE,EE with two low-$\ell$ temperature bins and the optical-depth prior $\tau=0.0506\pm0.0086$, DESI DR2 BAO, and DES-Dovekie with its absolute magnitude marginalized. The geodesic release uses the linear-depletion history with $A_L=1$ and lensing from its own perturbations. It reproduces the $\Lambda$CDM value $2226.767$ of the PPF pipeline above and $2228.359$ with nonlinear lensing disabled. Gauss--Newton minima are converged to $0.005$ in $\chi^2$:
\[
\begin{array}{l|rrr}
\text{Supernova likelihood, neutrinos, lensing}&\Lambda\text{CDM}&\text{Geodesic release}&\Delta\chi^2\\\hline
\text{low-}z\text{ offset},\ \textstyle\sum m_\nu=0,\ \text{linear}&2228.359&2227.777&-0.58\\
\text{low-}z\text{ offset},\ \textstyle\sum m_\nu=0,\ \text{nonlinear}&2226.767&2226.380&-0.39\\
\text{Dovekie},\ \textstyle\sum m_\nu=0.06\,\text{eV},\ \text{nonlinear}&2239.031&2239.880&+0.85
\end{array}
\]
The last row also frees the Planck calibration with the prior $1\pm0.0025$. Its difference splits into $-1.54$ from the CMB, $+1.98$ from BAO, $+0.83$ from supernovae and $-0.42$ from the priors. The BAO cost follows from the tied abundance: with $\omega_c=5.333\,\omega_b$ the release model cannot offset the neutrino mass by lowering $\omega_c$. The release fit has $f=0.29$--$0.36$ across the rows. The five fitted parameters are $(\omega_b,f,\ln(10^{10}A_s),n_s,\tau)$, with the formation ratio fixed at $5.333$; the reference has six. Separate archived linear-depletion and exponential fits with massless neutrinos, fixed calibration, linear lensing and a low-redshift offset give $\Delta\chi^2=-0.5186$ and $-0.5152$, respectively. Their relative scores differ by about $0.0034$; the archives use different numerical accuracy settings, so this difference includes numerical effects. Their $\Lambda$CDM reference is $2228.316$, against $2228.359$ for the one-harness table above, whose first row gives $-0.58$ for the corresponding geodesic fit. The differences are statistical ties, and the parameter count gives $\Delta\mathrm{AIC}=-1.15$ on the last row. ACT primary-spectrum and lensing checks are reported separately in Section~27. The assembly history and its astrophysical prior, a validated nonlinear response, and a full joint likelihood remain to be determined. The low-density completion of O.3 preserves the fitted cold linear branch.

\paragraph{A first self-gravitating measurement.}
A cosmological SPH run with Gadget-4 \citep{Springel2021} follows a $0.75\,h^{-1}\,$Mpc box with $2\times96^3$ particles from $z=63$ to $z=6$, with primordial cooling \citep{KatzWeinbergHernquist1996}, the effective-pressure star-formation model of \citet{SpringelHernquist2003}, and no ultraviolet background above $z=6$; each gas particle carries $6.5\times10^3h^{-1}M_\odot$. In each host whose main progenitor first exceeds the gate, the closed baryonic cohort is labelled by its friends-of-friends progenitor one dynamical time earlier, or by the octant of its position about the progenitor if it belonged to none, and $K_{\rm mix}$ is measured on shells and octants and on cubes of $R_{\rm vir}/2$ and $R_{\rm vir}/4$, with a shuffled-label baseline subtracted. The measured kernels share one form,
\[
K_{\rm mix}=1-(1-K_0)\,e^{-\tau/(f_{\rm tail}\,t_{\rm dyn})},\qquad K_0\simeq0.5,\qquad f_{\rm tail}\simeq2\text{--}6 .
\]
Here $f_{\rm tail}$ describes decay after the nonzero entry loss.
It cannot be equated with the pure-exponential $f$ of either fit above;
the full convolved histories are the comparable observables.
Convolved with the first-entry history at the same gate, they give $I_{\rm rel}=76$ for three hosts at $17\ {\rm km\,s^{-1}}$, $93$--$97$ for ten hosts at $14\ {\rm km\,s^{-1}}$, and $108$--$110$ for sixteen hosts at $12\ {\rm km\,s^{-1}}$, against the required $94.2$. Labelling two dynamical times before entry raises the $17\ {\rm km\,s^{-1}}$ value to $82$--$83$; dropping the labels of smoothly accreted gas lowers it to $57$--$65$. The measured kernel therefore reproduces the release integral required by the data when the effective retention gate is $14\ {\rm km\,s^{-1}}$, a virial temperature of about $1.4\times10^4\,$K for neutral primordial gas ($\mu=1.22$), or $7\times10^3\,$K for ionized gas ($\mu=0.59$), and falls short by about a fifth at the reference gate. The fitted release scale becomes a statement about the retention gate and the progenitor labels, both of which the theory must specify. Choosing the gate after comparing its integral with the fitted target does not independently predict the normalization. The run lacks molecular cooling, feedback, subgrid mixing, and resolution beyond about $10^3$ gas particles per host, and its hosts are few; a resolved zoom with these processes, and with the gate measured from gas retention, remains the required test.

\paragraph{The retention gate from the carrier criterion.}
Section~18 identifies a carrier's records with its condensed baryons and assigns none to shock-heated gas whose cooling time exceeds its dynamical time. A newly assembled halo therefore enters the release history when its virialized gas can cool within one dynamical time. Take primordial gas with hydrogen mass fraction $0.76$ in collisional ionization equilibrium, with the atomic hydrogen and helium rates of \citet{KatzWeinbergHernquist1996}, Compton cooling on the microwave background, and neither molecular hydrogen nor an ultraviolet background. The condition $t_{\rm cool}=t_{\rm dyn}=R_{\rm vir}/V_{\rm vir}$ at the mean virial gas density then gives $T_{\rm gate}=(0.99$--$1.07)\times10^4\,$K for $30\ge z\ge8$. Collisional excitation of atomic hydrogen fixes this value: its rate rises by three orders of magnitude between $8\times10^3$ and $1.5\times10^4\,$K, so the threshold moves little with density or redshift. With $T_{\rm vir}=\mu m_pV_{\rm vir}^2/2k_B$ and the nearly neutral gas at the threshold, $\mu=1.22$, the gate is $V_{\rm gate}=11.6$--$12.0\,{\rm km\,s^{-1}}$. The reference gate of $17\,{\rm km\,s^{-1}}$ corresponds to the same temperature for fully ionized gas, $\mu=0.59$.

The Gadget-4 run measures the gate directly. Its friends-of-friends catalogues record the star-forming gas and stellar mass of every group at each of the 48 snapshots. The fraction of groups that hold condensed baryons rises from below $0.1$ to unity across a logistic width of $0.03\,$dex in virial velocity. It reaches one half at $8.6$, $8.9$, $9.3$ and $10.2\,{\rm km\,s^{-1}}$ for $z=11$--$15$, $9$--$11$, $7.5$--$9$ and $6$--$7.5$. Central gas densities exceed the mean virial density, so gas condenses below the mean-density threshold. Friends-of-friends masses slightly exceed virial masses, so these velocities are upper limits.

\paragraph{The vacuum density from the cooling threshold.}
With the abundance tied, the linear-depletion background gives $\omega_{\Lambda,0}=\eta_*\omega_c^{\rm wb}I_{\rm rel}$. Flatness with $\omega_b=0.02238$ then fixes $h$ and $\Omega_{\Lambda,0}$ from $I_{\rm rel}$ alone. The display below recomputes the first-entry history at each pre-reionization gate, with the post-reionization gate at $30\,{\rm km\,s^{-1}}$ and $z_{\rm re}=7.66$. Each history is convolved with the six kernels measured on the ten hosts at $14\,{\rm km\,s^{-1}}$ and the sixteen hosts at $12\,{\rm km\,s^{-1}}$. At fixed gate the kernels change $I_{\rm rel}$ by at most $6\%$, so the gate controls the result. The three hosts at $17\,{\rm km\,s^{-1}}$ have a slower tail and give the lower integral quoted above.
\begingroup
\setlength{\arraycolsep}{4pt}
\[
\begin{array}{l|c|c|c|c}
\text{Origin of the pre-reionization gate}&V_{\rm gate}\ ({\rm km\,s^{-1}})&I_{\rm rel}&h&\Omega_{\Lambda,0}\\\hline
\text{Run, condensed baryons, }z=9\text{--}15&9&128.6\text{--}135.7&0.77\text{--}0.79&0.766\text{--}0.775\\
\text{Run, condensed baryons, }z=6\text{--}9&10&118.9\text{--}125.4&0.75\text{--}0.77&0.752\text{--}0.761\\
\text{Cooling criterion, neutral gas, }\mu=1.22&12&103.7\text{--}109.5&0.71\text{--}0.73&0.725\text{--}0.736\\
\text{Gate selected by the likelihood}&14&92.4\text{--}97.5&0.69\text{--}0.70&0.702\text{--}0.713\\
\text{Cooling criterion, ionized gas, }\mu=0.59&17&80.9\text{--}85.3&0.65\text{--}0.67&0.673\text{--}0.685
\end{array}
\]
\endgroup
The likelihood requires $I_{\rm rel}=94.2\pm0.8$, $h=0.689$ and $\Omega_{\Lambda,0}=0.707$.

Both gates computed without feedback lie below the gate that the likelihood selects. They give $\Omega_{\Lambda,0}=0.73$ from the cooling criterion and $0.75$--$0.77$ from the run. Supernova feedback and photoheating inside ionized regions before the global reionization redshift remove gas from the smallest halos and raise the effective gate \citep{ThoulWeinberg1996,BarkanaLoeb1999,OkamotoGaoTheuns2008}. Neither process is present in the run. Gates from $12$ to $17\,{\rm km\,s^{-1}}$ span the cooling threshold under the neutral and ionized conventions. They give $\Omega_{\Lambda,0}=0.67$--$0.74$ and a vacuum density between $14\%$ below and $16\%$ above the required value.

The cooling-criterion calculation uses $\eta_*$, the formation abundance, excursion-set assembly of the measured fluctuation spectrum, the measured mixing kernel and the atomic-hydrogen cooling threshold. It gives $I_{\rm rel}=103.7$--$109.5$, $10$--$16\%$ above the archived target $94.2$, and $\Omega_{\Lambda,0}=0.725$--$0.736$ against $0.707$. Condensation in the no-feedback run gives $I_{\rm rel}=118.9$--$135.7$. Neither calculation includes feedback. The $14\,{\rm km\,s^{-1}}$ gate, the value selected by the target, corresponds to about $5\times10^7M_\odot$ at $z\simeq10$.

\paragraph{The vacuum density from feedback-resolved retention.}
A retained well is a halo that keeps gas able to cool, in which progenitor labels mix. The post-reionization gate measures this retention for photoheated gas \citep{OkamotoGaoTheuns2008}, and the same criterion holds before reionization. Stars carry records and are collisionless, so the advection--diffusion identity above gives them no label mixing. A halo that forms a few stars and then loses its gas holds records without a mixing medium and lies outside the gate. The fraction of halos holding star-forming gas therefore measures the pre-reionization gate. The Renaissance simulations follow Population~III and metal-enriched star formation with radiative transfer, radiative and supernova feedback, metal cooling and Lyman--Werner backgrounds in rare-peak, normal and void regions to $z=15$, $12.5$ and $8$ \citep{OShea2015Renaissance,Xu2016Renaissance}. About $10\%$ of halos of $10^7M_\odot$ host star formation younger than $20\,$Myr, about half of those of $10^8M_\odot$ do, and all above $10^{8.5}M_\odot$ do. The fraction vanishes near $7\times10^6M_\odot$ and is nearly independent of redshift and region. Let this occupation $f(M)$ be the probability that the gate of a trajectory lies below $M$. The release integral is linear in $F_Z$, so
\[
I_{\rm rel}=\int I_{\rm rel}(M_g)\,df(M_g),
\]
where $I_{\rm rel}(M_g)$ is the integral for a sharp gate at mass $M_g$ with the post-reionization gate and $z_{\rm re}$ unchanged. No input in this chain is adjusted to the vacuum density. The six measured kernels give
\begingroup
\setlength{\arraycolsep}{4pt}
\[
\begin{array}{l|c|c|c}
\text{Pre-reionization gate}&I_{\rm rel}&h&\Omega_{\Lambda,0}\\\hline
\text{Sharp, }10^{7}M_\odot&137\text{--}145&0.81&0.78\\
\text{Sharp, }10^{7.5}M_\odot&104\text{--}109&0.72&0.73\\
\text{Sharp, }10^{8}M_\odot&80\text{--}84&0.66&0.68\\
\text{Renaissance occupation}&92.3\text{--}97.4&0.693&0.709
\end{array}
\]
\endgroup
The kernels move the occupation result by $\pm2.5$ in $I_{\rm rel}$. Varying the zero point, the $10\%$ level, the half-occupation mass and the saturation mass of the occupation curve moves it by $+4/-2.4$. The reionization redshift $z_{\rm re}=6.94$--$8.40$ of the CMB optical depth moves it by $\pm5.5$, and a post-reionization gate of $25$--$35\,{\rm km\,s^{-1}}$ by $\pm4.3$. In quadrature,
\[
\boxed{I_{\rm rel}=95.5\pm8,\qquad\Omega_{\Lambda,0}=0.709^{+0.017}_{-0.019},\qquad h=0.693\pm0.021,}
\]
against $I_{\rm rel}=94.2\pm0.8$, $\Omega_{\Lambda,0}=0.707$ and $h=0.689$ required by the likelihood. Every term of the flatness sum is fixed once $\omega_b$ is measured, so the boxed $h$ is a prediction of $H_0=69.3\pm2.1\,\mathrm{km\,s^{-1}\,Mpc^{-1}}$. Each term scales with $\omega_b$, so $h$ varies nearly as $\omega_b^{1/2}$. The reionization redshift and the post-reionization gate dominate the error. Each trajectory draws one gate; independent draws at each epoch would raise entry. The Milky Way satellite census places the half-occupation peak mass for any star formation near $10^{7.5}M_\odot$ \citep{Nadler2020Satellites}. Those faint galaxies form a few thousand solar masses of stars and lose their gas, so they lie below the retention gate. For comparison, the standard quantum estimate of the vacuum energy exceeds the measured value by about $10^{120}$, and $\Lambda$CDM fits $\Omega_\Lambda$ as a free constant. The cold-gas retention fraction measured at $z\ge8$, a tighter reionization redshift and a feedback-resolved mixing kernel sharpen the test. A converged integral that differs from the target by more than its uncertainty falsifies the normalization under the premises of O.1--O.15.

\paragraph{The frozen-history likelihood.}
The occupation gate, the six measured kernels and the post-reionization gate fix the release history once the reionization redshift is given. The optical depth measures that redshift, so the likelihood takes $z_{\rm re}$ as the midpoint of the reionization history that the fitted $\tau$ implies and interpolates release histories computed at $z_{\rm re}=7.0$--$8.8$. The mixing time $f$ is then absent, and the release model has four cosmological parameters, $(\omega_b,\ln10^{10}A_s,n_s,\tau)$, against six for $\Lambda$CDM. In the one-harness likelihood above, with the linear-depletion background and lensing from its own perturbations,
\[
\begin{array}{l|rrr}
\text{Supernova likelihood, neutrinos, lensing}&\Lambda\text{CDM}&\text{Frozen release}&\Delta\chi^2\\\hline
\text{Dovekie},\ \textstyle\sum m_\nu=0.06\,\text{eV},\ \text{nonlinear}&2239.031&2240.100&+1.07\\
\text{low-}z\text{ offset},\ \textstyle\sum m_\nu=0,\ \text{nonlinear}&2226.767&2226.483&-0.28
\end{array}
\]
The first row also frees the Planck calibration and places $z_{\rm re}=8.12$ at $\tau=0.0587$, the optical depth of the $\Lambda$CDM fit, with $I_{\rm rel}=92.0$ and $H_0=68.42\,\mathrm{km\,s^{-1}\,Mpc^{-1}}$. The second gives $z_{\rm re}=7.89$, $I_{\rm rel}=93.7$ and $H_0=68.89\,\mathrm{km\,s^{-1}\,Mpc^{-1}}$. The parameter count gives $\Delta{\rm AIC}=-2.9$ and $-4.3$. In the first configuration the extreme kernels give $\chi^2=2240.16$ and $2240.32$, and post-reionization gates of $35$ and $25\,{\rm km\,s^{-1}}$ give $2240.19$ and $2241.66$, the second at $\tau=0.067$. Holding $z_{\rm re}$ at $7.66$ independently of the fitted $\tau$ gives $2259.18$ and $2230.74$, because the integral then exceeds the value the data select by about $3.5$. A larger optical depth places reionization earlier, closes small halos to entry sooner and lowers the vacuum density, so the frozen history predicts a correlation between $\tau$ and $\Omega_{\Lambda,0}$ that $\Lambda$CDM lacks. With $a_0=\epsilon cH_0$, the first fit gives $a_0=1.249\times10^{-10}\,\mathrm{m\,s^{-2}}$, and the astrophysical chain without cosmological data gives $a_0=(1.27\pm0.04)\times10^{-10}\,\mathrm{m\,s^{-2}}$, against $(1.20\pm0.02\pm0.24)\times10^{-10}\,\mathrm{m\,s^{-2}}$ from rotation curves \citep{McGaugh2016}. Neither value uses galactic data.

\paragraph{Large-scale release bias.}
The same first-entry construction fixes the response of release locations to
a long-wavelength density mode. Peak-background split shifts the excursion
barrier as $B(z)\to B(z)-\Delta_L$. Propagating the shifted first-entry
history through the same frozen mixing kernel gives
\[
\Gamma=\frac{\dot F_{\rm rel}}{1-F_{\rm rel}},\qquad
\boxed{
{\cal R}_\Gamma(a)=\frac1{D(a)}
\left.\frac{\partial\Gamma(a|\Delta_L)}{\partial\Delta_L}
\right|_{0},\qquad
\delta\Gamma={\cal R}_\Gamma(a)\Delta_m.}
\]
No free assembly-bias coefficient is introduced. With the ensemble potential held fixed as in O.2--O.3, this spatial hazard response
redistributes the internal source-resolved and reallocated capacity ledgers;
it cancels from the leading total metric-visible capacity perturbation for
$k\ne0$.

\paragraph{Cosmological zoom test.}
The remaining astrophysical task is a converged extension of the first
self-gravitating cosmological progenitor-tagged run summarized above. It must label
retained gas before atomic-cooling assembly, include the relevant infall,
cooling, feedback, and reionization physics, state the partition $Y_t$ on
which the carrier is resolved (finer regions retain more progenitor
information, and the theory has not yet fixed that scale), and report
\[
\boxed{
\epsilon_{\rm mix}(t,M,z)=1-\frac{I(A;Y_t)}{H(A)},\qquad
K_{\rm mix}(\tau;M,z).}
\]
It must also report the ratio $f=\tau_{\rm mix}/t_{\rm dyn}$ and the effective
retention gate. The measured kernel is to be frozen before the full CMB, DESI,
Dovekie, lensing, ISW, and growth likelihood is evaluated with direct interacting perturbations. The linear-depletion PPF benchmark constrains the unweighted integral to $I_{\rm rel}=94.2\pm0.8$ under its reference inputs. The matched direct-perturbation likelihood with the composition-preserving allocation selects $f=0.274245$, and a frozen kernel whose convolved history misses the weighted integral $I_{\eta,\rm rel}$ of that fit would falsify the allocation. Substantial pre-gate identity loss, excessive early release, or a frozen physical history inconsistent with that likelihood would falsify this extension without changing the static gravity sector.

\subsection*{O.13 Marked dilation, unit event coefficient, and energy ledger}

Let \(\mathsf E_A\) be the carrier-failure effect in O.11.  Let
\(|\alpha\rangle_M\), \(\alpha=1,\ldots,9\), be the orthonormal marked
states of \(V_1^P\otimes V_1^H\), and let \(W_{\alpha A}\) be mutually
orthogonal history embeddings.  The release isometry is
\[
\boxed{
V_A^{\rm rel}|\psi\rangle
=\sqrt{I-\mathsf E_A}\,|\psi\rangle|0\rangle_M
+\frac13\sum_{\alpha=1}^{9}
|\alpha\rangle_MW_{\alpha A}
\sqrt{\mathsf E_A}\,|\psi\rangle.}
\]
Because
\(W_{\alpha A}^\dagger W_{\beta A}=\delta_{\alpha\beta}I\),
\[
(V_A^{\rm rel})^\dagger V_A^{\rm rel}
=I-\mathsf E_A+\frac19(9\mathsf E_A)=I.
\]
Thus the isometry transfers the failed register into its history outputs with unit conditional probability. Identifying the transferred register with the full allocated mass stock is the energy-ledger prescription; the
\(1/3\) amplitude normalizes the nine equal outputs.  The
charged recurrence fugacity \(\zeta_*\) is not an event probability.

H.11 gives exact reference decoupling after renewal,
\[
(\operatorname{id}_R\otimes\mathcal E_*)(\rho_{R,{\rm ren}})
=\rho_R\otimes\rho_*.
\]
If the energetic ledger reads the closure budget after that refresh,
conditioning on a carrier failure leaves
\[
\rho_{{\rm closure}\mid{\rm rel}}=\rho_*,
\qquad
\langle K^2\rangle_{\rm rel}
=\langle K^2\rangle_{\eta_*},
\]
and therefore
\[
\boxed{
\eta_{\rm drain}
=\frac{\varepsilon_\xi}{\varepsilon_K}
\frac{3/2}{\langle K^2\rangle_{\eta_*}}
=\frac{\varepsilon_\xi}{\varepsilon_K}\eta_*.}
\]
The selected common conversion normalization
$\varepsilon_\xi=\varepsilon_K$ gives
$\eta_{\rm drain}=\eta_*=0.0298668443935$.
The event coefficient \(f_{\rm reroute}=1\) and the energy share
\(\eta_*\) do different jobs. The former follows from faithful register
renewal for each failed record block; the latter uses the common conversion
normalization and the post-refresh closure-energy identification.


\subsection*{O.14 Conditional normalization and status}

In the bare baryonic-mass convention, the normalization is expressed through
a supplied delayed release history. A prediction requires that history to
be fixed independently of the cosmological fit:
\[
\boxed{
\omega_{\Lambda,0}^{\rm pred}
=\eta_*\frac{\omega_b}{\epsilon}I_{\eta,\rm rel}(1),\qquad
I_{\eta,\rm rel}(1)=\int_0^1a^{-3}e^{-\eta_*F_{\rm rel}}\,dF_{\rm rel},\qquad
\Omega_{\Lambda,0}^{\rm pred}=\frac{\omega_{\Lambda,0}^{\rm pred}}{h^2}.}
\]
The value of $h$ is solved simultaneously from the O.3 flatness equation. For the free-nucleon/helium mass convention of the empirical fit, replace $\omega_b/\epsilon$ by $5.333\omega_b$ as specified in O.3.

The logical grade is:
\begin{itemize}
\item Identity-free capacity cannot record stream assignment, and faithful
renewal makes rerouting complete when the baryonic progenitor carrier fails.
\item The constrained capacity-clock action assigns only
$\eta_*\mathbb E_{\rm bg}[\Gamma_{\rm rel}\rho_u]$ to the homogeneous
cosmological term. It reproduces the background drain and fixes the full
linear scalar exchange: $\delta\rho_\Lambda=0$ on capacity-clock slices,
while lapse and clock perturbations appear covariantly in other gauges.
\item Carrier loss creates no new material or scalar component. The
metric-visible remainder is reallocated source-associated capacity of the
retained coarse carrier; the orthogonal marked-history output has no
independent zeroth-derivative metric stress.
\item The total source-associated capacity obeys
$\bar\rho_C=\rho_{c,0}^{\rm wb}a^{-3}e^{-\eta_*F_{\rm rel}}$. Internal release-hazard fluctuations cancel between the source-resolved and reallocated ledgers. The ensemble clock drain obeys the local positivity domain and fitted conditional margin stated in O.3. The constrained action fixes the clock and lapse dependence of $Q_\Lambda^\nu=-\nabla^\nu\rho_\Lambda$ and leaves the capacity flow geodesic.
\item The assembly input $F_Z$ is fixed by the absorbing first-entry
excursion-set algorithm of O.12. The self-gravitating progenitor-tagged
run measures the delayed mutual-information kernel and demonstrates
sensitivity to the retention gate and label definition. Additive
record pricing fixes this normalized information loss as the fractional stock
law. For a closed Lagrangian cohort with common advection--diffusion, the
displayed dissipation identity proves that the fully resolved position loss is
bounded and nondecreasing. The chosen coarse physical kernel needs its own
monotonicity check.
\item At the reference retention gates the linear-depletion PPF benchmark
selects $f=0.282706$ and $I_{\rm rel}=94.2116$. The measured kernel
has nonzero loss at entry and a slower tail. Its complete convolution
reaches the fitted integral near a $14\,{\rm km\,s^{-1}}$ gate and
undershoots it at $17\,{\rm km\,s^{-1}}$. The feedback-resolved occupation
of O.12 fixes the gate independently and reaches the fitted integral within
its uncertainty; the label definition must be fixed independently.

\item A first self-gravitating progenitor-tagged run and an approximate
fit to CMB, DESI, and Dovekie data are available (O.12). The
remaining task is to fix a converged $K_{\rm mix}(\tau;M,z)$ and its gate
independently, then test the resulting fixed history against the direct positive-allocation likelihood and growth, lensing, and ISW observables.
\item The post-refresh energy ledger, common conversion normalization,
identification of the progenitor record with source-clock stock,
fixed-measure continuum premise, and initial condition
$\Lambda_i=0$ are the remaining physical premises of the normalization.
\end{itemize}

\paragraph{Kernel normalization and a catalog retention test.}
For the measured form
$K(\tau)=K_0+(1-K_0)K_{\rm tail}(\tau)$, the linear-depletion integral is
affine in the entry loss:
\[
I_{\rm rel}=I_{\rm tail}+K_0(I_{\rm immediate}-I_{\rm tail}).
\]
Matching a cosmological normalization can therefore select $K_0$ when the
target lies between these endpoints; it does not independently determine
that loss. On the same saved reference histories, replacing $dF$ by
$e^{-\eta_*F}dF$ changes the measured-kernel integrals by only
$0.19$--$0.21\%$. This fixed-background comparison does not transfer the
linear-benchmark target $I_{\rm rel}=94.2$ to the direct positive-allocation likelihood.
Raw coarse information-loss estimates also have downward steps, the largest
about $0.09$. Sampling noise and information return can both contribute;
the monotone exponential fit alone does not establish monotonicity of the
physical coarse observable.

The 48 saved friends-of-friends catalogs permit an operational retention
test independent of the dark-energy target. Define retained fraction as
$(M_{\rm gas}+M_\star)/M_{\rm total}$ divided by the simulated cosmic baryon
fraction, and look for an upward crossing of its median through one half.
Require at least 256 stock particles per group and eight groups per
velocity bin, with no minimum gas-count cut. Among the 31 snapshots above
the reference $z_{\rm re}=7.66$, 17 have qualifying bins. Their medians all
exceed three quarters, so neither the one-quarter, one-half nor three-quarter
criterion yields a measured crossing; the other 14 epochs lack sufficient
bin statistics. At a 1024-particle cut no epoch has the required populated
bins. A separately measured one-half star-forming fraction crosses at
$9.30$--$9.80\,{\rm km\,s^{-1}}$ in eight epochs, but it is a different
observable and is not substituted for retention. Repeated epochs do not
constitute independent halo samples. Group association is not a binding-energy
criterion, and the run's resolution and baryonic-physics limitations remain.
These catalogs consequently do not select the physical release gate; the
source-record definition and a resolved mixing history must fix it before
the cosmological comparison.

Once the physical mixing kernel is frozen, convolving it with the fixed
first-entry history and evolving the homogeneous drain and one-metric
perturbations determines the likelihood without a phenomenological release
amplitude or perturbation-bias parameter. Substantial pre-gate identity
loss, excessive early release, $w_{\rm eff}>-1$ reconstructed for this
same vacuum component with matter transfer accounted for, or a frozen history inconsistent with
cosmological data falsifies this extension without changing the static
gravity sector.

\subsection*{O.15 Formation, the label ledger, and no late recruitment}

\paragraph{Label ledger.}
Conserved baryonic charge carries source labels. A primordial source $A$ is
partitioned into elementary record blocks $\beta$ with ancestry bits
$a_{A\beta}\in\{0,1\}$, and its remaining cap is
\[
M^{\max}_{c,A}=\sum_\beta a_{A\beta}\,\frac{M_{A\beta}}{\epsilon}.
\]
Carrier loss of block $\beta$ sets $a_{A\beta}:1\to0$, reroutes the block's
allocation through the O.13 isometry, and assigns its baryonic carrier to the
coarse label of the retained carrier $Y$, whose cap equals its retained stock
$R_Y$. Release removes cap together with allocation, the surviving blocks of $A$ retain their source identities, and the reallocated stock has no source other than release. The displayed label cap is its formation allocation; the common fractional drain of O.3 scales the surviving source-resolved and reallocated energy budgets without restoring a lost label:
\[
\boxed{
\text{formation feeds }M_c,\qquad
\text{release feeds }R\text{ and }\Lambda,\qquad
R\not\to M_c .}
\]
The local label operations of O.9 preserve, release, or merge labels, and
descendants inherit the label of their parent carrier (Section~15). No
operation removes every label from conserved charge.

\paragraph{Formation map.}
Let $n_\varnothing$ be the density of conserved baryonic charge that carries
no source label. Formation acts only on $n_\varnothing$. For a Lagrangian
element $A$ of that charge it copies the element's identity into the capacity
ledger, assigns the cap $M_A/\epsilon$ with $M_A$ evaluated on the formation
slice $\Sigma_0$, starts the pinned count at $N_A=0$, and draws its energy
from the thermal plasma:
\[
\boxed{
\nabla_\mu T_c^{\mu\nu}=Q_{\rm form}u^\nu,\qquad
\nabla_\mu T_{\rm rad}^{\mu\nu}=-Q_{\rm form}u^\nu,\qquad
\int Q_{\rm form}\,d\tau=\frac{\rho_b}{\epsilon}.}
\]
The map mirrors the release isometry of O.13. Let $P_\varnothing$ project onto
unlabelled conserved charge for which the plasma supplies $\rho_b/\epsilon$, and
let $W^{\rm form}_\alpha$ write the label, the cap, the initial count, and the
plasma debit, with
$(W^{\rm form}_\alpha)^\dagger W^{\rm form}_\beta=\delta_{\alpha\beta}P_\varnothing$.
Then
\[
\boxed{
V^{\rm form}|\psi\rangle
=(I-P_\varnothing)|\psi\rangle|0\rangle_M
+\frac13\sum_{\alpha=1}^{9}|\alpha\rangle_M\,W^{\rm form}_\alpha P_\varnothing|\psi\rangle,
\qquad
(V^{\rm form})^\dagger V^{\rm form}=I .}
\]
A label-stripping map of the same form is equally isometric, so the ledger
rule above carries independent content: O.9 contains no such map.
The content of a Lagrangian element is time independent only for a conserved
current, so the cap requires a conserved source. Thermal electron--positron
pairs carry no net conserved charge and receive no allocation; at
$k_BT=m_ec^2$ they outnumber the net charge by $4\times10^8$ and annihilate
$10^{17}$ times per Hubble time. The homogeneous vacuum cannot supply the
formation energy, because by O.2 the cosmological term changes only through
the ensemble mean of the one-signed marked drain.

Formation acts on the first conserved baryonic charge that admits a label,
which lies between sphaleron freeze-out and the QCD crossover. There the
radiation debit $\rho_b/(\epsilon\rho_{\rm rad})$ is between
$4\times10^{-12}$ and $4\times10^{-9}$, and every mode used in the CMB and
Lyman-$\alpha$ likelihoods is superhorizon, so the formation slice imposes
$\delta_c=\delta_b$. Filling the cap before the $\ell\simeq2500$ modes enter
the horizon requires a per-tick formation weight above
$e^{-16.1g_{\rm share,eff}}$. One marked recurrence, $e^{-7g_{\rm share,eff}}$,
exceeds this by a factor $2.1\times10^{29}$, so the cap fills within a small
fraction of a Hubble time and $\sigma_A=1$ holds exactly. For formation at the
QCD crossover the cap uses the free-nucleon mass, which exceeds the present
mean baryonic mass by $0.24\%$ after helium binding, and the abundance
prediction becomes $\Omega_c/\Omega_b=5.333$.

\paragraph{Clock deficit.}
On each smooth committed stream branch $b$, $X=\tfrac12$ defines the unit
clock $u_\mu=-\nabla_\mu\phi_b$. For an allocation with proper velocity
$v_A^\mu$ and pinned count $N_A$, define
\[
\Delta_A=\frac{\phi_b-\phi_0}{\tau_*}-N_A .
\]
The transfer ceiling gives $dN_A/d\tau_A\le1/\tau_*$, and
$d\phi_b/d\tau_A=-u_b\cdot v_A=\gamma_A\ge1$, so
\[
\boxed{
\frac{d\Delta_A}{d\tau_A}\ge\frac{\gamma_A-1}{\tau_*}\ge0 .}
\]
An allocation first formed at an event $p$ to the future of $\Sigma_0$ has
$N_A(p)=0$ and $\Delta_A(p)=[\phi_b(p)-\phi_0]/\tau_*>0$, and subsequent
evolution keeps $\Delta_A>0$. Relative motion increases the deficit.

\begin{quote}
\textbf{No-late-recruitment theorem.} After formation, $n_\varnothing=0$.
The one-step vertex $\mathcal A_v$ acts on the geometric factor, the
replaceable closure register, and the marked fiber, and conserves $n_d$. The
fresh state $A_*(b)$ carries no clock argument, and history export is forward
only. The release isometry never raises $c_{A,c}$. The label operations of
O.9 contain no operation that creates a label or strips every label from
conserved charge. Net baryon number is conserved below the sphaleron
freeze-out temperature. Hence $n_\varnothing$ remains zero, the formation map
has no support to the future of $\Sigma_0$ at any later temperature, and
committed stock persists or decreases through the one-signed marked drain.
The result is exact for the specified transfer together with the ledger rule
above. Its completion requires a formation vertex supported on
$n_\varnothing$.
\end{quote}

Pair creation, carrier loss, galaxy mergers, supernova cores,
gamma-ray-burst fireballs, and heavy-ion collisions therefore cannot
regenerate primordial stock. Released gas in galaxies is excluded twice: it
carries the coarse label of its retained carrier, and its local thermal and
nuclear energy falls short of the required $m_pc^2/\epsilon=4.99$~GeV per
baryon by more than two orders of magnitude. The released share $\eta_*$
cannot grow back, which protects the late-time deficit of the committed
density and the Lagrangian stock budget $C_R$. Records written in ordinary matter by rare thermal
fluctuations do not involve committed stock; Appendix~P.10 bounds them.

\section*{Appendix P: Record Dynamics, Organization, Computation, and the Growth of Complexity}
Organized complexity accumulates through three distinct processes. Renewal
exports each previous state into history, beginning from a minimal-record
boundary. For the canonical defect, the source cost also measures how quickly
its history becomes distinguishable from vacuum. Gravitational growth builds
correlations among retained records. The primitive orientation-matching
interaction generates arbitrary local unitaries and composes across a
connected network. Finally, collective loss of a baryonic source-clock
carrier gives a one-signed release law for the marked record share of committed capacity. Renewal orients time; the history it writes cannot be recycled and has a minimal-record root; and the released vacuum cannot re-enter the committed phase within the specified dynamics (P.8--P.10).

The grades differ across this chain. Exact recoverability forces fresh logical history capacity (P.9), and its physical source is open. Relaxation of a recorded subsystem requires an irreducible record-writing coupling, established case by case (P.8); purposeful computation further requires a structured gate measure (P.6). Interpreting marked
release as coarse-grained record loss also inherits the source-clock premise,
progenitor-mixing history, post-refresh ledger, and history orientation of
Appendix~O.13 and Section~20.

Here a \emph{record} normally means a state written to the history register
by the export isometry in Appendix~G's Stinespring dilation. P.3 also uses a
weaker definition: any persistent physical distinction that separates a
defect history from vacuum history.
\subsection*{P.1 Grading conventions}
The status terms follow Section~26. A \emph{closed} result is derived
from the stated postulates and ensemble. A \emph{conditional} result
also requires a named physical reading or an unfinished microscopic
completion. Numerical verification is tracked separately. Section~P.11
distinguishes reproduced calculations, single-instance results awaiting a
fresh check, blind calculations against sealed targets, and exploratory
calculations that carry no evidential weight.
These results refine the theory's internal structure; they do not change
any of the probability ledgers defined elsewhere in the paper.
\subsection*{P.2 The reversible-record theorem pair}
\paragraph{One-step preservation.}
The local renewal channel replaces the current cell state by the fixed
state $\rho_*$,
\[
\mathcal E_*(\rho)=\rho_*\operatorname{Tr}\rho.
\]
Viewed only on the renewed cell, this map appears to erase the input.
Its reversible dilation transfers the input to the complementary
history channel $\mathcal C:A\to H$. That channel is an isometric
embedding, $\mathcal C(\rho)=V\rho V^\dagger$, and therefore has an
exact recovery map satisfying
$\mathcal R\circ\mathcal C=\mathrm{id}_A$. Trace distance, fidelity,
relative entropy, and every other isometry-invariant distinguishability
measure are preserved in the history register. One renewal step
transfers the old state together with all of its information.
For coherent spatial bonds, the exported register is first expressed
in the common graph frame of Q.11a. Joint renewal then prepares the
correlated vacuum while preserving the local replacement marginal.
The exact intertwining identity preserves retained graph coherence;
exporting graph-dependent bond states directly would leave a graph
address in history. H.14 separately checks background-wide renewal
on a fixed geometry: it acts identically on retained position and
internal state. Path-conditioned extra resets can export
distinguishable records even with identical fresh marginals.
The graph-frame identity does not by itself equate history records
for geometries with different numbers of physical events.

\paragraph{Reconstruction from descendant registers.}
The present register retains no ancestral information. The faithful-history kernel $K_*(b,b')=p_*(b')$ has identical rows, so on the $1680$-state space
\[
\operatorname{spec}K_*=\{1,0,\ldots,0\},
\qquad
\lambda_2=0 .
\]
For the symmetric three-state and ferromagnetic four-state broadcasting processes on regular and Poisson trees, ancestral information survives to arbitrary depth exactly when $d\lambda^2>1$, with nonreconstruction at equality \citep{OpenAIThreeState2026,OpenAIPottsReconstruction2026}. The renewal kernel has rank one, so each descendant present register is independent of its ancestor and the Kesten--Stigum quantity $d\lambda_2^2$ vanishes for every finite branching $d$, while the history channel $\mathcal C$ preserves every distinguishability measure exactly. The canonical construction therefore has no intermediate reconstruction phase:
\[
\boxed{
\begin{gathered}
\text{present channel: }\lambda_2=0,\ \text{complete renewal},\\
\text{history channel: isometry, complete retention}.
\end{gathered}}
\]
Broadcasting thresholds become relevant once copying is imperfect, exported records are noisy, or history is coarse-grained with loss. The second singular value $\lambda_2$ of that transfer channel then enters a Kesten--Stigum-type criterion $b|\lambda_2|^2\gtrless1$, which the cited theorems make exact in the three- and four-state cases.
\paragraph{Cumulative capacity.}
Suppose update $n$ writes into a fresh independent factor of dimension
$d_n$. After $N$ updates,
\[
\dim\mathcal H_{{\rm hist},N}
=\dim\mathcal H_{{\rm hist},0}\prod_{n=1}^{N}d_n.
\]
The capacity is therefore at least exponential when every $d_n\ge2$.
For the nine-state marked register it is $9^N$ before gauge constraints.
H.11 and P.6 show that, when these are simultaneously retained independent
factors and their projected bridges survive the host constraints, this count
is realized as the quantum composition
$\mathcal H_N=(\mathbb C^9)^{\otimes N}$ rather than merely as a storage
dimension.
The cumulative statement holds whenever every earlier record stays exactly recoverable. P.9 shows that exact recoverability itself forbids destructive recycling, so fresh logical capacity is forced; its physical source remains an open completion.
\subsection*{P.3 The visibility--source identity}
\paragraph{The fixed-type identity.}
The canonical elementary defect excludes one channel label $m$ from one
face $f$. Let $P_*$ be the vacuum distribution on boundary states and
let $A_m=\{b:b_f\neq m\}$ be the allowed subset after the exclusion.
The defect distribution is the conditioned distribution
$P_{\rm def}=P_*(\,\cdot\mid A_m)$. Direct substitution gives
\[
\boxed{\;D(P_{\rm def}\Vert P_*)=-\ln P_*(A_m)=\Delta F_m.\;}
\]
The three expressions describe the same likelihood ratio. Spatially,
$\Delta F_m$ is the defect's free-energy insertion into the vacuum
ensemble and hence the source cost used in Appendix~C.5. Temporally,
the relative entropy is the Stein exponent: for many independent
records, it is the asymptotic error exponent for distinguishing defect
history from vacuum history. This interpretation uses the memoryless
dressing of Postulate~III. On the certified $1680$-state ensemble at
$\eta_*$, the seven channel values lie between $0.1493$ and $0.1582$
nats.
\paragraph{Isotropic benchmark and Jensen structure.}
Write $p_m=P_*(b_f=m)$ for the one-face marginal. Since
$\sum_m p_m=1$, averaging the allowed probabilities before taking the
logarithm gives
$-\ln[\tfrac17\sum_m(1-p_m)]=\ln\tfrac76$, which is the isotropic
source insertion in C.5. Averaging the seven log costs instead gives
\[
\bar D=-\tfrac17\sum_m\ln(1-p_m)\;\ge\;\ln\tfrac76,
\]
by Jensen's inequality, with equality only for a uniform marginal. On
the closed ensemble,
$\bar D-\ln(7/6)=5.7\times10^{-6}$ nats. The closed ensemble has
maximum face-marginal deviation $\max_m|p_m-1/7|=0.00414405$. The annealed/quenched cost difference is $5.7353\times10^{-6}$ nats, so the source cost and the
record-discrimination rate agree to that accuracy. The exact identity
\[
H(P_*)-H(P_*(\,\cdot\mid A))
=D(P_*(\,\cdot\mid A)\Vert P_*)
+\eta_*\!\left[\langle K^2\rangle_{P_*}
-\langle K^2\rangle_{P_*(\,\cdot\mid A)}\right]
\]
shows that the small variation among channels comes from
closure-energy anisotropy rather than numerical error.
\paragraph{Sufficiency.}
The likelihood ratio depends only on the indicator
$X(b)=\mathbf 1_{A_m}(b)$, so the one-bit defect statistic is
sufficient: collapsing the $1680$-state configuration to the scalar
``constraint satisfied or not'' loses no relative information about
this defect. The scalar source projection in C.5 is therefore lossless
for the canonical exclusion defect.
\paragraph{Projective form of the same source cost.}
Embed a probability distribution by its square-root vector,
\[
|\Omega(P)\rangle=\sum_b\sqrt{P(b)}\,|b\rangle.
\]
For the conditioned defect $P_{\rm def}=P_*(\cdot\mid A_m)$,
\[
\langle\Omega_{\rm def}|\Omega_*\rangle=\sqrt{P_*(A_m)}.
\]
The visibility--source identity therefore has the equivalent projective form
\[
\boxed{
\Delta F_m
=-\ln|\langle\Omega_{\rm def}|\Omega_*\rangle|^2
=-2\ln\cos\theta_m,}
\qquad
\theta_m=d_{\rm FS}(\Omega_{\rm def},\Omega_*).
\]
For the isotropic exclusion $P_*(A_m)=6/7$, this gives
$\theta_m=\arccos\sqrt{6/7}=0.38760$ radians.  For $N$ independent records,
the squared overlap is $e^{-N\Delta F_m}$.  The gravitational source cost is
therefore also the exact history-orthogonalization exponent for this canonical
defect family.
\paragraph{Contrapositive and scope.}
Full support implies $P_*(A)<1$ for every genuine exclusion, so the
relative entropy is positive. Conversely, $D=0$ implies
$P_{\rm def}=P_*$ and hence no exclusion defect. Within this defect
class, a record-invisible excitation has no relative defect information,
no scalar source insertion, and no elementary defect mass. A coherent
defect outside the exclusion family could evade the one-bit sufficiency
result. Its exclusion still depends on the existing GFT singlet-gap
gate.

A committed cell that is exempt from refresh leaves its record on the
perimeter. Joint face admissibility changes the neighboring marginals
by $0.0020$ nats per adjacent slot for an exclusion and by $0.1499$ nats
for a pinned face. The same neighbor coupling both writes this perimeter
record and produces the capacity strain that gravitates. A defect that
left no record in any neighbor would therefore have no gravitational
source, contrary to Postulate~II.

\paragraph{Entanglement of the closure vacuum and of its defects.}
The square-root vacuum $|\Omega_*\rangle=\sum_b\sqrt{P_*(b)}\,|b\rangle$ on the $840$ injective face assignments, with $P_*(b)\propto e^{-\eta_*K^2(b)}$, is a state on four seven-dimensional face factors; the independent orientation doubling leaves every face reduced state unchanged. Its face entropies are
\[
S_*(1)=1.452283,
\qquad
S_*(12)=1.888710,
\qquad
I_*(1{:}2)=2S_*(1)-S_*(12)=1.015857\ \text{nats},
\]
so the closure vacuum is strongly entangled across every face cut. Conditioning on the elementary exclusion $m_1\neq m$ and renormalizing gives, in nats,
\[
\begin{array}{l|c|c}
\text{quantity}&m=0&\text{range over the seven exclusions}\\ \hline
\text{source information }\Delta F_m&0.15819&0.14933\text{ to }0.15819\\
\text{marked-face entropy change}&-0.10562&-0.10562\text{ to }-0.09559\\
\text{unmarked-face entropy change}&-0.00685&-0.00685\text{ to }-0.00627\\
\text{pair mutual-information change}&-0.08953&-0.08953\text{ to }-0.08123\\
\text{unmarked-face classical KL}&0.00241&0.00202\text{ to }0.00241\\
\text{unmarked-face quantum relative entropy}&0.00682&0.00629\text{ to }0.00682
\end{array}
\]
The classical KL row is the $0.0020$-nat perimeter signal quoted above, resolved by label. The reduced quantum relative entropy of an unmarked face is about three times its classical KL divergence, so the coherent vacuum carries neighbor information that the diagonal distribution does not register. The source cost $\Delta F_m\simeq0.15$ nats exceeds the marked-face entanglement loss of about $0.10$ nats. The source quantity of C.5 is the vacuum-relative information $\Delta F_m$, and the local entanglement loss, the mutual-information loss, and the small neighboring deficit are its responses inside the microscopic state:
\[
\boxed{
\text{defect}\rightarrow\Delta F_m\rightarrow
\text{marked-face entanglement loss}
+\text{neighboring information deficit}.}
\]
The global pure vacuum and defect states have distinct rank-one supports, so their quantum relative entropy is infinite. The finite source is the dephased relative information, equal to the overlap exponent in the projective form above, and every regional reduced relative entropy is finite. On an extended network the same calculation gives the regional change $\Delta S_A=S_A^{\rm def}-S_A^{\rm vac}$ as a function of distance from the defect; agreement of its linear response with the Green operator of $q$ would identify the regional information deficit with the gravitational response field.

For the canonical defect, the source strength in C.5 and the rate at
which its history separates from vacuum history are the same relative
entropy. The weak-field source calculation and Postulate~III use one
shared constant.
\subsection*{P.4 Boundary, arrow, and well-foundedness}
\paragraph{The minimal-record boundary.}
Many-Pasts contains neither a second time coordinate nor an archive
outside the physical records. If no record survives, the statements
``an earlier era occurred and every trace vanished'' and ``no earlier
era occurred'' describe the same physical state. The initial boundary
is therefore the state with the least record content. A state already
containing many records cannot serve as that boundary because the
records themselves specify a past. Unrecorded prehistory is not a separate variable in this ontology. P.9 derives this boundary as the root of the record order from exact recoverability, memoryless renewal and the conditioning rule of G.6.
\paragraph{The boundary measure in G.7.}
Without a boundary condition, the counting formula in G.7 favors
high-entropy pasts and encounters the usual fluctuation objection. The
minimal-record identification supplies the missing domain: the boundary
measure $\mu_i$ is defined on the first physically distinguishable
record-bearing state. The arrow of time is then a typicality statement
conditioned on that state. For the canonical defect class, the
contrapositive in P.3 also gives a simple matter sector at the boundary.
This supplies G.7 with a named ontological condition for its
low-entropy boundary.
\paragraph{The well-foundedness bound.}
Non-veridical durable records require spontaneous re-entry into the committed phase, which vanishes in the specified dynamics (P.10). For records written in matter, define the content of a specified record by its equilibrium surprisal, $K_R=-\ln\pi_{\rm eq}(R)$; once a recorded subsystem has relaxed to its equilibrium measure (P.8), the per-trial probability of a specified false record is $e^{-K_R}$. A union bound then supplies the estimate below, and P.10 adds the independent thermal bound. Taking the substrate tick
to be the induced Planck time gives a total trial budget
$\ln N_{\rm trials}\simeq5.6\times10^{2}$ nats (convention spread a few
nats), giving
\[
P(\text{false durable record of content }K)\;\lesssim\;
e^{\,\ln N_{\rm trials}-K},
\qquad K_{\rm crit}\sim5.6\times10^{2}\ \text{nats}
\;\approx\;8\times10^{2}\ \text{bits}.
\]
This threshold corresponds to roughly $0.4$~GeV of committed
mass--entropy. Macroscopic records
($K\gtrsim10^{25}$) are veridical at odds indistinguishable from
certainty. The result remains conditional on the pinned/synchrony reading of
Section~20.
\subsection*{P.5 Organization: what actually grows}
\paragraph{Conserved production.}
For fixed total mass, the identity in P.3 fixes total record production
per tick. Complexity growth cannot come from an increasing production
rate in this bookkeeping. It must appear as changing correlations among
the records already being produced. This separates record amount from
record organization, as in the information-growth literature
\citep{Hidalgo2015}.
\paragraph{Native locality.}
An exploratory exact-matching model measures how quickly the native
injectivity channel loses record-stream correlation. Its mutual
information is $1.95$ nats per tick for adjacent cells, $0.155$ with
one intervening cell, and $0.0023$ with two. The corresponding channel
has $\lambda_2=0.1723$ and correlation length
$\xi_{\rm hist}=0.57$ cells. This is a single-instance calculation and
is excluded from evidential claims.

The calculation isolates the question for a completed model. The
capacity/history channel is short-ranged by itself, whereas the gapless strain
mode carries correlations over long distances. Gravity can therefore
correlate records written in separated regions. Ordinary matter interactions
and shared causes can also correlate them; the result is specific to the
history sector's gravitational channel, not exclusive to gravity.
\paragraph{The monotonicity theorem.}
Assume that the one-stream marginals are stationary to leading order. Write
the connected joint distortion of a pair of retained record streams as
\[
p_{ij}(x,y;a)=p_i(x)p_j(y)
\left[1+D_+(a)\,\varepsilon_{ij}(x,y)\right],
\]
with the marginal constraints
$\sum_y p_j(y)\varepsilon_{ij}(x,y)=
\sum_x p_i(x)\varepsilon_{ij}(x,y)=0$ and
$|D_+\varepsilon_{ij}|\ll1$.  This is the precise sense in which the
long-range strain mode changes the cross-correlation while the local record
alphabet and its vacuum marginals remain fixed. Expanding the relative
entropy of the joint state from the product state gives
\[
I_{ij}(a)
=D_{\rm KL}(p_{ij}\Vert p_ip_j)
=\frac{D_+^2(a)}{2}
\sum_{x,y}p_i(x)p_j(y)\varepsilon_{ij}(x,y)^2
+O(\varepsilon^3).
\]
The sign also holds exactly throughout the positive interval of this affine family. Write $p_{xy}(d)=p_xp_y(1+d\varepsilon_{xy})$ with $d\ge0$ and $1+d\varepsilon_{xy}>0$. The fixed-marginal constraints give
\[
\frac{dI}{dd}=\sum_{x,y}p_xp_y\varepsilon_{xy}
\ln(1+d\varepsilon_{xy})\ge0,
\qquad
\frac{d^2I}{dd^2}=\sum_{x,y}
\frac{p_xp_y\varepsilon_{xy}^2}{1+d\varepsilon_{xy}}\ge0.
\]
Each term in the first sum is nonnegative because the logarithm has the sign of $d\varepsilon_{xy}$. Consequently $\mathcal I_H=\sum_{i<j}I_{ij}$ is nondecreasing whenever $D_+'\ge0$ and this family remains positive. Its weak-distortion form is
\[
\mathcal I_H=D_+^2\mathcal I_H^{(2)}+O(\varepsilon^3),
\qquad \mathcal I_H^{(2)}\ge0.
\]
The coefficient $\mathcal I_H^{(2)}$ is a sum of squares and is nonnegative for every spectral shape. In a translationally invariant representation the same fact appears as $C_S^{(2)}=\tilde C_S*\tilde C_S$, a convolution of the nonnegative strain power with itself, so the sign holds for any spectrum. The exact inequality concerns a fixed-marginal affine history family. Applying it to nonlinear physical evolution requires that evolution to preserve the family or satisfy its own information inequality. The rooted-propagation calculation retains the registered consistency targets $G(r)\sim1/r$ and $\mathcal I_{\rm hist}(r)\sim1/r^2$.

Under the square-root embedding of the joint record distribution, the same
expansion gives
\[
d_{\rm FS}^2(p_{ij},p_ip_j)
=\frac{D_+^2}{4}
\sum_{x,y}p_i(x)p_j(y)\varepsilon_{ij}(x,y)^2
+O(\varepsilon^3),
\]
and hence
\[
\boxed{I_{ij}=2d_{\rm FS}^2+O(\varepsilon^3).}
\]
Thus the growing gravitational mode moves the classical record ray away from
the product-state (Segre) submanifold; the increase in mutual information is
twice the squared projective distance at leading order.

Within the stated weak-coupling regime, P.3 fixes the production rate
and the result above fixes the growth of organization. Both statements
describe the gravitational sector in history variables.
\subsection*{P.6 Realized complexity}
Circuit complexity is the number of elementary gates needed to prepare
the realized history state. Sustained growth requires a gate set that is
both controllable and sampled from a nondegenerate measure. The fixed routing
factors of H.9 commute, but they are not the complete one-mark dynamics.  Once
the marked fiber is identified with the retained fusion complement, its frame
generators and inherited Lorentzian mismatch spectrum generate a full local
algebra in each polar fiber; H.11's exact symplectic propagator carries that
algebra unitarily between fibers. Algebraic capacity and physical sampling
remain separate questions.

The primitive geometric interaction passes the algebraic controllability
test exactly. Complexifying the spin-$3/2$ pair generators
$\{iJ_i\otimes J_j\}$ replaces each one-slot Cartesian basis by the
rational $\mathfrak{sl}(2)$ basis $\{E,F,H\}$. Closing the resulting nine
pair generators under commutators gives rank $255$ over both
$\mathbb F_{1009}$ and $\mathbb F_{10007}$. Full rank modulo either prime
certifies rank at least $255$ over $\mathbb Q$; traceless $16\times16$
matrices have dimension at most $255$. The complex closure is therefore
$\mathfrak{sl}(16)$, and the original compact real closure is
$\mathfrak{su}(16)$. The smaller $\mathfrak{so}(16)$ and
$\mathfrak{sp}(16)$ alternatives are excluded by the same dimension
certificate.

\paragraph{Primitive gluing generators.}
For $j_0=3/2$ the finite-width matching operator of Appendix~B may be
written, after transporting the right face into the left relational frame,
\[
Q_f(g)=\frac94 I-\mathbf J_L\!\cdot\!\widetilde{\mathbf J}_R(g).
\]
As the relative holonomy $g$ varies, the transported Cartesian components
span the nine bilinears $J_{L,i}\otimes J_{R,j}$ (the fixed orientation sign
is already contained in $\widetilde{\mathbf J}_R$).  The rank-$255$
certificate above therefore applies to the orientation-matching interaction
itself, without an auxiliary computational Hamiltonian. Full Lie rank is locally
stable: a finite list of nested commutators contains a nonzero $255\times255$
minor, and that minor remains nonzero under sufficiently small perturbations
of the generator coefficients. The universal algebra therefore persists in
a neighborhood of the stated coefficients. Its theory-specific input is the
capacity-selected representation $j_0=3/2$ and the associated $16\to7+9$
fusion structure.

\paragraph{The retained one-mark algebra.}
Let $L:Q\to\mathcal H_{\rm mark}$ be the incidence isometry of H.9.  Moving
the relative face frame compresses the primitive rotations to
\[
A_i=P_Q(I\otimes J_i)P_Q\big|_Q,
\]
whose exact closure is $\mathfrak{so}(9)$.  The primitive mismatch operator
restricts to
\[
Q_f|_Q=6P_0+5P_1+3P_2.
\]
This operator fixes mass-squared differences, not the physical phase
generator itself.  On the stable retained quadratic branch, H.11 constructs
the positive-frequency complex structure from the real Hamiltonian and gives
the co-moving frequency operator
\[
\Omega_{\rm ref}
=\omega_0P_0+\omega_1P_1+\omega_2P_2,
\qquad
\omega_J
=\sqrt{k^2+\frac{M_c^2}{Z}+\frac{\kappa}{Z}q_J},
\quad
(q_0,q_1,q_2)=(6,5,3).
\]
The physical relative-phase generator is therefore
\[
S_\omega
=\Omega_{\rm ref}-\frac{\operatorname{Tr}\Omega_{\rm ref}}9I,
\]
a nonzero real symmetric traceless operator.  Its
$\mathfrak{so}(9)$ orbit spans $\operatorname{Sym}_0(9)$, so in the
positive-frequency coordinates derived in H.11,
\[
\boxed{
\operatorname{Lie}\!\left(\mathfrak{so}(9),iS_\omega\right)
=\mathfrak{su}(9).}
\]
The common frequency restores $\mathfrak u(9)$. Thus the local
controllability statement does not assume a complex marked Hilbert space:
the positive quadratic dynamics first selects the initial polarization, the
full symplectic flow transports it exactly, and the frame and relative-phase
generators close the 80-dimensional unitary algebra in every transported
fiber.

\paragraph{Composition across a network.}
Pair controllability alone does not guarantee network controllability. A
diagonal bridge shows the obstruction. Start with the full
$\mathfrak{sl}(16)$ algebra on slots $(1,2)$ and couple slot $2$ to a third
spin-$3/2$ slot only through a diagonal bridge $H_2\otimes H_3$.  Every
generated operator preserves the four eigenspaces of $H_3$, so the closure
is contained in the traceless block-diagonal algebra
\[
\mathfrak s\!\left(\mathfrak{gl}(16)^{\oplus4}\right),
\qquad
\dim=4(16^2)-1=1023.
\]
The exact modular calculation saturates this bound.  A diagonal bridge can
transport correlation while remaining incapable of rotating the new slot.

The physical network supplies the missing local algebra without a new
interaction: slot $3$ belongs to another face pair, and that pair's certified
$\mathfrak{sl}(16)$ contains $I\otimes\mathfrak{sl}(4)_3$ as a subalgebra.
The following tensor-product lemma then closes the composability question.
Let $A\in\mathfrak{sl}(m)$ and $B\in\mathfrak{sl}(n)$ be nonzero.  Then
\[
\operatorname{Lie}\!\left(
\mathfrak{sl}(m)\otimes I,
I\otimes\mathfrak{sl}(n),
A\otimes B
\right)=\mathfrak{sl}(mn).
\]
Indeed, the adjoint representation of the simple algebra
$\mathfrak{sl}(m)$ spans all of $\mathfrak{sl}(m)$ from any nonzero
$A$, and similarly for $B$.  Nested commutators therefore generate
$\mathfrak{sl}(m)\otimes\mathfrak{sl}(n)$; together with the two local
summands this is the standard vector-space decomposition of
$\mathfrak{sl}(mn)$.  With $m=16$, $n=4$ the three-slot algebra is thus
\[
\boxed{\mathfrak{sl}(64)},\qquad \dim=64^2-1=4095.
\]
A deterministic exact modular breadth-first closure using the physical
set $\{\mathfrak{sl}(16)_{12},\mathfrak{sl}(4)_3,
H_2\otimes H_3\}$ independently reaches rank $4095$ over
$\mathbb F_{1009}$. This calculation checks the analytic tensor-product
lemma. A second-prime implementation would strengthen its computational
provenance.

The same lemma gives the network theorem by induction on a spanning tree.
If a connected cluster of $r$ spin-$3/2$ slots already carries
$\mathfrak{sl}(4^r)$ and an adjacent physical face supplies a universal
pair algebra on one cluster slot and a new slot, the inherited local
$\mathfrak{sl}(4)$ and any nonzero entangling bilinear generate
$\mathfrak{sl}(4^{r+1})$. Hence on any finite connected face network,
provided every physical edge realizes the same two-slot controllability
family,
\[
\boxed{
\mathfrak g_{\rm network}=\mathfrak{su}(4^N)
}
\]
on the compact real form. Every connected slot participates in a certified
pair algebra, so the result propagates over the spanning tree. This is an
algebraic composability theorem; it does not determine which gates the
physical path measure samples.

\paragraph{Composition of retained marks.}
The four-slot construction uses the unreduced retained carrier
\[
\mathcal H_{2,{\rm retained}}=Q_{12}\otimes Q_{34},
\]
with no additional hard-core or gluing reduction. Conditional on this product
surviving the final host projection, the primitive bridge acts on two
neighboring marked complements. For pairs $(1,2)$ and $(3,4)$, define
\[
H=2J_z=\operatorname{diag}(3,1,-1,-3),
\]
and
\[
\begin{aligned}
H_{QQ}
&=(P_Q^{12}\otimes P_Q^{34})(H_2\otimes H_3)
(P_Q^{12}\otimes P_Q^{34})\\
&=A_R\otimes A_L,
\end{aligned}
\]
with
\[
A_R=P_Q(I\otimes H)P_Q\big|_Q,
\qquad
A_L=P_Q(H\otimes I)P_Q\big|_Q.
\]
Exact rational spin-$3/2$ evaluation gives
\[
\operatorname{rank}A_R=\operatorname{rank}A_L=8,
\quad
\operatorname{Tr}A_R=\operatorname{Tr}A_L=0,
\quad
\operatorname{Tr}A_R^2=\operatorname{Tr}A_L^2=38.
\]
The interaction is therefore nonzero and cannot be reduced to a sum of local
terms.  Applying the tensor-product lemma with $m=n=9$ gives
\[
\boxed{
\operatorname{Lie}\!\left(
\mathfrak{su}(9)_1,\mathfrak{su}(9)_2,
iA_R\otimes A_L\right)=\mathfrak{su}(81).}
\]
By induction, a connected marked network has
\[
\boxed{
\mathcal H_N=(\mathbb C^9)^{\otimes N},
\qquad
\mathfrak g_N=\mathfrak{su}(9^N),}
\]
provided adjacent retained marks coexist. The geometric gauge projection acts as the identity on the capacity factor (Q.4, P.8), so the projected bridge survives it. With $\|A_R\otimes A_L\|=41/5$ the bridge is a bounded finite-range interaction, so the network also carries the quasi-local infinite-volume algebra, Lieb--Robinson dynamics, and GNS representations of H.11.

\paragraph{Commitment and the gate measure.}
Algebraic universality establishes computational capacity. A realized
computation also requires a structured, non-equilibrium gate measure. The
framework contains a parameter-free candidate for this bias: commitment
persistently conditions the local capacity draws. Let $M(dg\mid b)$ be the
as-yet-uncomputed
microscopic kernel taking a local capacity state $b$ to the relative
holonomy/gate variable $g$.  The equilibrium and canonical-defect gate
measures are the push-forwards
\[
\mu_*(dg)=\sum_bP_*(b)M(dg\mid b),
\qquad
\mu_{\rm def}(dg)=\sum_bP_{\rm def}(b)M(dg\mid b).
\]
The data-processing inequality and P.3 give the exact bound
\[
\boxed{
0\le D(\mu_{\rm def}\Vert\mu_*)
\le D(P_{\rm def}\Vert P_*)
=\Delta F_m.
}
\]
Thus the same likelihood ratio that prices source strength and record
visibility is the maximum information available to bias the local gate
measure.  The remaining test is sharp: one must derive $M$ and show
$\mu_{\rm def}\ne\mu_*$ (equivalently, nonzero information between the
defect indicator and the gate variable). If that test passes, committed
matter acts both as a record and as a persistent bias on the universal
gluing interaction. If it fails, the controllability theorem remains valid
but commitment does not program the gate distribution. Universality alone
makes no claim of purposeful computation.

The primitive pair interaction can approximate an arbitrary two-slot
unitary when its coefficients vary independently or are sampled in a
sufficiently broad sequence. This supplies the algebraic capacity for
general information processing. A derivation of chemistry would also need
the low-energy spectrum, stable composites, gauge interactions, and physical
control supplied by the continuum matter completion.

The remaining physical questions concern the gate measure and its control. Controllable dynamics requires positive weight on a gate family whose generators close to the controllable algebra. Purposeful computation further requires the commitment kernel above to deform the gate measure, which remains to be shown. Existing circuit-complexity bounds
would then constrain realized growth, but algebraic controllability alone
does not prove exact linear complexity growth or identify a useful
algorithm.

At the native-cell level, the pinned initial configuration has an
$O(N)$ preparation circuit and is therefore low-complexity. The
preparation complexity of the full geometric condensate remains open with
the physical gate-measure problem. P.3 derives record writing by canonical
defects, P.5 derives organization by a specified growing strain mode, and
P.6 derives the available computational algebra. Its physical sampling measure remains open for relaxation and for structured computation.
\subsection*{P.7 The self-limiting loop and collective carrier loss}

P.5 shows how the growing gravitational mode organizes retained records:
to leading nontrivial order their mutual information grows as \(D_+^2\).
The cosmological feedback uses a different operation: irreversible loss
of the particular baryonic record that supports a Lagrangian source
clock.

Identity-free committed capacity cannot record which stream it occupies.
Two descriptions that differ only by a permutation of identical
stream assignments export the same occupancy record,
\[
|e_i\rangle=|e_j\rangle
\quad\Longrightarrow\quad
\gamma_{ij}=1,\quad C_A^{\rm rec}=1.
\]
Shell crossing and the fold-record conventions therefore supply no
marked event.  The release effect is instead
\[
\boxed{
\mathsf E_A
=\mathbf1[\text{carrier of }A\text{ has lost progenitor identity}].}
\]
When gas from several progenitors becomes irreversibly mixed in one
retained bound object, matter no longer distinguishes the source label
while the collisionless committed allocation still carries it.  The
orphaned allocation is rerouted through the O.13 isometry.

Faithful full-support renewal admits no partial-memory stock update at
that event, so \(f_{\rm reroute}=1\). At coarse resolution the source record
is partitioned into elementary blocks, and the same additive information
pricing used by Mass--Entropy Equivalence assigns released stock in proportion
to the fraction of progenitor information absent from matter. The history
dependence is therefore entirely in
\[
\epsilon_{\rm mix}(t)
=1-\frac{I(A;Y_t)}{H(A)},
\]
the fraction of progenitor information absent from the resolved baryonic
state. Appendix~O.12 proves that its fully resolved position counterpart is
bounded and monotone for a closed cohort under common advection--diffusion,
shows why a coarse observable needs an independent check, and distinguishes
information loss from the maximal complete-reset component of a general stochastic channel. The marked
accumulation law is delayed rather than instantaneous,
\[
\boxed{
\begin{aligned}
F_{\rm rel}(t)&=\int^tK_{\rm mix}(t-t_f)\,dF_Z(t_f),\\
\omega_{\Lambda,0}&=\eta_*\frac{\omega_b}{\epsilon}
\int_0^1a^{-3}e^{-\eta_*F_{\rm rel}(a)}\,dF_{\rm rel}(a),\\
\Omega_{\Lambda,0}&=\frac{\omega_{\Lambda,0}}{h^2}.
\end{aligned}}
\]
The linear-depletion PPF benchmark selects $f=0.282706$ and $I_{\rm rel}=94.2116$ for the reference first-entry history. The
self-gravitating progenitor-tagged run supplies a kernel whose integral
depends on the retention gate and progenitor labels (O.12). With the gate
from the feedback-resolved occupation it gives $I_{\rm rel}=95.5\pm8$; a
converged label and kernel measurement sharpens the test.

This also fixes the self-limiting direction without confusing the two
\(a_0\) roles.  The relation \(a_0(z)=\epsilon cH(z)\) controls the
transverse thermal response.  Release contains no \(a_0\): structure
formation builds retained, turbulently mixed common wells; accumulated
\(\Lambda\) then suppresses later growth and hence further carrier loss.
\subsection*{P.7a Three laws of the retained history register and the Landauer reading}
The following ``laws'' are bookkeeping statements internal to the history
sector and stand alongside the thermodynamic first, second, and
third laws.  In them, production, organization, and irreversible
coarse-grained loss, often grouped under the one word ``information,'' are
represented by three different quantities.

\paragraph{Law I: record production is priced by the same quantity that sources mass.}
For the canonical exclusion defect, P.3 proves exactly
\[
D(P_{\rm def}\Vert P_*)=\Delta F_m.
\]
The left side is the asymptotic per-record distinguishability rate and the
right side is the source insertion.  Through Mass--Entropy Equivalence the
same committed entropy is read inertially as mass.  Thus, at fixed committed
defect inventory and refresh cadence, the record-production budget is fixed
by the mass/source budget; complexity growth cannot be obtained merely by
postulating an ever-increasing rate of record creation.  The statement is
exact for the canonical exclusion class and does not assert a universal rate
for every possible composite matter state.

\paragraph{Law II: gravitational growth monotonically organizes retained records.}
Under the fixed-marginal weak-correlation hypothesis of P.5, the connected
record distortion is linear in the growing strain mode whereas mutual
information begins quadratically.  Hence
\[
\boxed{
\mathcal I_H(a)=D_+^2(a)\mathcal I_H^{(2)}+O(\varepsilon^3),
\qquad
D_+'(a)\ge0\Longrightarrow
\frac{d\mathcal I_H}{da}\ge0
}
\]
to leading nontrivial order.  Matter correlations have several sources;
the theorem identifies the history sector's own
long-range strain channel and the sign of the organization it builds.

\paragraph{Law III: loss of a source-clock carrier has a one-signed energetic residue.}
Appendix~O.13 supplies the operator identity behind the third statement.
Let
\[
\widehat H_{\rm rel}=\widehat K^2,
\qquad
\langle\widehat H_{\rm rel}\rangle_*
=\langle K^2\rangle_{\eta_*}=50.2229154254,
\]
and let the normalized marked record strand be generated by
\[
\widehat H_{\rm rec}
=\frac12\sum_{a=1}^3\widehat\xi_a^2,
\qquad
\langle\widehat H_{\rm rec}\rangle=\frac32.
\]
The dark-energy branch uses three stated conditions: (i) common conversion
normalization \(\varepsilon_K=\varepsilon_\xi\); (ii) a post-refresh energy
ledger, which makes the closure state \(\rho_*\) independent of the slow
carrier-failure effect; and (iii) one-signed routing of the marked share
out of the retained coarse stress budget.  For an event selected by
\(\mathsf E_A\), these give
\[
\boxed{
\eta_{\rm drain}
=\frac{\varepsilon_\xi}{\varepsilon_K}
\frac{\langle\widehat H_{\rm rec}\rangle}
{\langle\widehat H_{\rm rel}\rangle_*}
=\frac{\varepsilon_\xi}{\varepsilon_K}
\frac{3/2}{\langle K^2\rangle_{\eta_*}}
=\eta_*=0.0298668443935,}
\]
where the last equality uses the selected common conversion normalization
$\varepsilon_\xi=\varepsilon_K$.
Faithful renewal separately fixes the stock coefficient to
\(f_{\rm reroute}=1\).  Therefore, for physical released energy
\(E_{\rm rel}\ge0\),
\[
\boxed{E_{\rm drain}=\eta_*E_{\rm rel}\ge0,}
\qquad
\boxed{
E_\Lambda(N)
=E_\Lambda(0)+\sum_{j=1}^N\eta_*E_{{\rm rel},j}}
\]
is monotone.  Microscopically the \(16\to7+9\) update remains isometric.
\lq\lq Erasure'' means that progenitor identity is no longer recoverable
from the retained baryonic algebra and the marked complement is absent
from the local stress bookkeeping; it is not fundamental nonunitarity.

\paragraph{Landauer interpretation, with the scale dependence explicit.}
The energetic statement admits a Landauer-like interpretation but does
not by itself prove a thermal Landauer process.  Mass--Entropy Equivalence
already supplies
\[
\boxed{\epsilon_{\rm rec}(\ell)=\kappa_m(\ell)c^2}
\]
at the scale \(\ell\) of the committed record.  At the electron anchor,
\[
\epsilon_{\rm rec}(\lambda_e)
=\frac{m_ec^2}{\ln2}
\simeq0.737\ {\rm MeV\,nat^{-1}}.
\]
A literal thermodynamic identification \(Q=k_BT\Delta S\)
\citep{Landauer1961} would additionally require a reservoir and a
fluctuation--dissipation relation; neither is used in the cosmological
normalization.

\paragraph{Three ledgers and the scope falsifier.}
Ordinary thermal reset, quantum measurement, and particle reactions do
not source the cosmological drain merely because records or entropy
change.  Only the marked source-clock channel selected by
\[
\mathsf E_A
=\mathbf1[\text{the baryonic carrier has lost progenitor identity}]
\]
enters O.3.  This prevents the false inference that computation,
annihilation, ionization, or generic environmental decoherence sources
dark energy.

\paragraph{Observational test.}
O.7 proves that the identity-free capacity register cannot carry primitive
which-stream information. The self-gravitating progenitor-tagged run in
O.12 measures the mutual-information kernel on several coarse partitions.
The feedback-resolved occupation of O.12 fixes the retention gate. The
remaining test is to fix the label rules through converged astrophysical
calculations, freeze the resulting kernel, and
evaluate its $a^{-3}$-weighted release history against the cosmological
likelihood. The coarse observable requires its own monotonicity check.

The three laws govern different quantities: record production,
gravitational organization, and coarse-grained loss. The first two are exact
in their stated defect and weak-correlation domains. For the third, additive
record pricing and the closed-cohort mixing theorem fix the fractional stock
law, while renewal and the marked isometry fix the fate of each released
block. The result remains conditional on identifying the baryonic progenitor
record with the source-clock stock, on the physical mixing history, on the
common energy conversion and post-refresh ledger, and on the one-signed
history orientation.
\subsection*{P.8 The reduced channel and the thermodynamic arrow}

\paragraph{Relative-entropy contraction.}
Let $\mathcal T$ be the one-tick channel on a finite connected retained patch,
obtained from the one-step vertex $\mathcal A_v$ of Appendix~Q.9 by preparing
the fresh registers in $\rho_*$ and tracing out the exported history factors.
If $\rho_{\rm eq}$ is a faithful state with $\mathcal T(\rho_{\rm eq})=\rho_{\rm eq}$,
monotonicity of quantum relative entropy gives
\[
\boxed{
D\big(\mathcal T(\rho)\,\big\Vert\,\rho_{\rm eq}\big)\le
D(\rho\Vert\rho_{\rm eq}).}
\]
On the renewed register itself $\mathcal T=\mathcal E_*$, and
\[
D\big(\mathcal E_*(\rho)\,\big\Vert\,\rho_*\big)=0\le D(\rho\Vert\rho_*)
\]
after one tick. Local renewal is therefore an exact relative-entropy
H-theorem, while the export isometry of P.2 keeps the global evolution
reversible. On a microcanonical shell, $\rho_{\rm eq}=\Pi/d$ and
$D(\rho\Vert\rho_{\rm eq})=\ln d-S(\rho)$, so $S(\mathcal T\rho)\ge S(\rho)$.

\paragraph{The source of the orientation.}
Relative-entropy contraction holds for every channel with a fixed point; the
orientation lies in the conditions that make $\mathcal T$ a channel. Forward
in time each tick consumes a fresh register prepared in $\rho_*$ and
uncorrelated with the present, and exports the previous register into
history. The reduced map on the present is then completely positive and trace
preserving. The time-reversed description draws its inputs from history
factors that are correlated with the present, and its reduced map is not a
channel on the present alone. The stationary kernel is detailed-balanced
(G.7); the arrow comes from the export structure together with the
minimal-record root of P.9.

\paragraph{Gauge projection.}
Projecting the four spin-$3/2$ face slots of a cell onto their $SU(2)$
singlets leaves the intertwiner space $\operatorname{Inv}(V_{3/2}^{\otimes4})$
of dimension four. Between such states every transported face bilinear
$J_{A,i}\,R(g)_{ij}\,J_{B,j}$ vanishes for every holonomy (largest matrix
element $3\times10^{-31}$), so a bridge acting on gauge-projected slots would
vanish identically. Appendix~Q excludes that reading. Q.3 shows that the
capacity alphabet cannot be the geometric Gauss data, since the Gauss Casimir
vanishes on gauge-invariant states while $K^2\ge122/3$. Q.5 and Q.6 exclude a
capacity factor without frames, since the lock of Postulate~I and the matching
operator of Appendix~B compare neighboring frames through the transported
holonomy. The remaining reading is the factorization of Q.4,
\[
U_{\rm gauge}=U_{\rm geom}\otimes I_{\rm cap}\otimes I_{\rm orient}.
\]
The geometric gauge projection therefore acts as the identity on the capacity
factor, and the bridge depends on geometry only through the gauge-invariant
relative orientation of the locked frames. The network results of P.6 hold
without a separate gauge-projection condition on the bridge.

\paragraph{Which factors relax.}
The one-tick channel acts differently on the factors of the retained patch. On the replaceable closure register it is $\mathcal E_*$, which reaches $\rho_*$ in one tick from every input. On the persistent geometry--defect system the renewal identity of H.11 gives
\[
(U_S\otimes I_R)(\operatorname{id}_S\otimes\mathcal E_*)(\rho_{SR})(U_S^\dagger\otimes I_R)=U_S\operatorname{Tr}_R(\rho_{SR})U_S^\dagger\otimes\rho_*,
\]
and the no-which-path condition keeps the exported output independent of the defect's path. The reduced channel on that system is the unitary $U_S$. It preserves relative entropy and every superposition, as the unitary transport of H.11 and the interference account of G.10 require. A subsystem relaxes when its interactions write records into exported factors. Its reduced channel then contracts relative entropy, and it reaches a unique stationary state when its Kraus operators generate the full matrix algebra of its sector and one of them is positive definite. The canonical dilation of H.9 supplies a positive-definite no-event branch with eigenvalues $1$ and $\alpha=\sqrt{1-u}$, $u=8\eta_*/49$. Irreducibility depends on the particular record-writing coupling and is established subsystem by subsystem.

\paragraph{Code-preserving gates on the closure code.}
Gates acting on the capacity labels must preserve injectivity. A single-port rotation does not: $J_+$ on the first port maps $|0,1,2,3\rangle$ to a multiple of $|1,1,2,3\rangle$. Projected generators $P_{\rm cap}HP_{\rm cap}$ preserve the code and generate unitary gates. Each moves one label by one unit, so the four labels of a cell behave as hard-core particles on seven sites and keep their order. On two cells glued along all four faces, face bilinears divide the $840^2$ code states into $576=24^2$ sectors, and intra-cell exchange joins them into one. Renewal redraws the complete register from $p_{\eta_*}$ each tick, and this supplies the mixing of the closure code. The commutant calculations for spin-$3/2$ pairs and gluing algebras concern gate families on unrestricted tensor products.

\paragraph{The thermodynamic arrow.}
The arrow has three parts with different grades. Within one tick, renewal is an exact H-theorem on the closure register, and the export structure fixes the orientation. Every reduced channel with a faithful fixed point contracts relative entropy to it. The macroscopic second law follows from conditional typicality (G.7): records accumulate in the exported history, the minimal-record root of P.9 fixes the boundary, and coarse-grained entropy grows. Coherent sectors evolve unitarily, and subsystems continually recorded by their surroundings relax. A unique stationary state for a given recorded subsystem requires its own irreducibility calculation. Conserved quantities such as energy and charge restrict any relaxation to their sectors.

\subsection*{P.9 History capacity, the record root, and the de~Sitter budget}

\paragraph{Non-reuse.}
Composing the export isometries of P.2 over $N$ updates keeps every earlier
record exactly recoverable. If update $n$ exports $d_n>1$ distinguishable
alternatives, the product alternatives $|a_1\rangle\otimes\cdots\otimes|a_N\rangle$
remain orthogonal under the composed isometry, so
\[
\boxed{
\dim\mathcal H_{\rm hist}(N)\ge\prod_{n=1}^Nd_n,
\qquad
\ln\dim\mathcal H_{\rm hist}(N)\ge\sum_{n=1}^N\ln d_n.}
\]
Overwriting an earlier record in place while keeping it recoverable violates
injectivity. Hardware can be reused only after its content has been moved
elsewhere, which consumes the same logical capacity. Exact renewal at every
tick therefore requires unbounded fresh logical history capacity, and
destructive recycling contradicts the dilation theorem. The physical source of
that capacity is open.

\paragraph{Finite-energy reset and exported history.}
\phantomsection\label{sec:historyenergy}
An unbounded record supply need not require an ever-growing energy reserve
in the present register. An explicit dilation separates these requirements.
Let $|r\rangle$ have native probabilities $p_r$ and register energies $E_r$.
Choose $E_{\max}=\max_r E_r$ and an auxiliary register with states
$|B_r\rangle$ of nonnegative energy $E_{\max}-E_r$.
Every code state $|r,B_r\rangle$ has the same total energy. Define
\[
W|r,B_r\rangle
=\sum_s\sqrt{p_s}\,|s,B_s\rangle\otimes|r,s\rangle_{\rm hist},
\qquad H_{\rm hist}=0.
\]
Orthogonal record indices and $\sum_s p_s=1$ give $W^\dagger W=I$.
For every density operator on the code,
\[
\operatorname{Tr}_{\rm hist}(W\rho W^\dagger)
=\sum_s p_s|s,B_s\rangle\langle s,B_s|\,\operatorname{Tr}\rho.
\]
Thus the native register has exactly its Gibbs marginal after the reset,
while the isometry retains the input information in the exported factor.
Each term has energy $E_{\max}$; the auxiliary register exchanges
$E_{\rm old}-E_{\rm new}$ and remains in a fixed finite band.
In the dimensionless convention $E_r=K_r^2$, that band is $[0,16]$ because
$122/3\le K_r^2\le170/3$. The conversion to physical energy remains the
common energy-conversion premise of O.3.

The map can be tensored with the persistent geometry--defect evolution of
H.10. Its coefficients and exported address must be independent of that
system's path to satisfy the no-which-path condition. This construction
assumes a fresh $d^2$-dimensional history factor at each step, where $d=1680$;
successive record pairs share their intermediate label, so a fixed initial
input has $d^N$ populated alternatives after $N$ resets. A finite-step
isometry admits a unitary extension with a blank ancilla in the same energy
sector. Neither that extension nor the finite auxiliary energy band supplies
the unlimited logical storage. The zero-leading-stress interpretation also
requires a metric- and source-independent history coupling; a zero
Hamiltonian in one frame is insufficient. The physical origin of this non-geometric resource
remains distinct from the reset's energy bookkeeping.

\paragraph{The record root.}
The history registers belong to the present state, and non-reuse makes
their dimension grow with every recoverable layer. A bound on present
capacity that includes them is equivalent to a finite past and cannot
establish one. The probability rule of Many-Pasts bounds the past
directly. It conditions on the realized present record $R$ through
$p(h\mid R)=p(h,R)/p(R)$ (G.6), so $p(R)>0$. Exact recoverability keeps
every past renewal outcome as an orthogonal record in the present, and
memoryless renewal draws each outcome independently from $p_{\eta_*}$.
On the $1680$ states, $122/3\le K^2\le170/3$ gives
\[
p_{\max}=\max_b p_{\eta_*}(b)\le\frac{e^{16\eta_*}}{1680},
\qquad
\ln\frac1{p_{\max}}\ge\ln1680-16\eta_*=6.95 .
\]
The realized record refines the event that the $N_{\rm past}$ recovered
outcomes took their recorded values, so
\[
0<p(R)\le p_{\max}^{N_{\rm past}},
\qquad
N_{\rm past}\le\frac{K_R}{\ln(1/p_{\max})},
\qquad
K_R=-\ln p(R).
\]
The surprisal of the present record bounds the number of recoverable past
layers. An unbounded past of recoverable renewals would give every
candidate present zero weight, which the conditioning rule excludes. A
kernel with memory gives the same bound with $p_{\max}$ replaced by its
largest conditional probability, which full support keeps below one. The
record order is therefore well founded. Its minimal elements carry no
record of a distinguishable predecessor, and Many-Pasts identifies an
unrecorded predecessor with none (G.6). Exact recoverability, memoryless
renewal and the conditioning rule thus imply a minimal-record root: the
set of minimal elements of the record order, which is the boundary of P.4
and G.7. The bound requires no capacity premise and does not single out
one initial event. The root has the fewest records; its gravitational
entropy is not derived here.

\paragraph{The de~Sitter budget.}
The trial budget of P.4, $\ln N_{\rm trials}\simeq5.6\times10^2$, gives
$N_{\rm trials}\approx10^{243}$. Memoryless renewal makes every trial an
independent draw whose exported state stays recoverable, so the history
already written carries at least of order $10^{243}$ nats. A de~Sitter patch
at the observed cosmological constant has
\[
S_{\rm dS}=\frac{\pi c^5}{G\hbar H_\Lambda^2}\approx3.3\times10^{122}.
\]
These order estimates put the written history about $3\times10^{120}$ times above that capacity. Two
consequences follow. The history output cannot be counted inside the
area-law capacity of the patch: it belongs to the non-geometric history factor, which carries no independent leading metric stress (O.3). A completion in which record time ends once the patch
capacity is exhausted is excluded by the theory's own count. The far future is therefore decided dynamically, by whether the released vacuum can re-enter the committed phase (P.10).

\paragraph{What the de~Sitter bounds constrain.}
The N-bound limits the entropy that an observer in a spacetime with a positive cosmological constant can measure to $S_{\rm dS}$ \citep{Bousso2000}, and the observable algebra of a static patch with an observer is of type II$_1$, with empty de~Sitter space as its maximum-entropy state \citep{ChandrasekaranEtAl2023}. Both statements concern information that can be prepared, stored, and retrieved inside the patch. Three entropies of the construction enter the comparison. The horizon cut carries $\ln2$ per channel face of area $a_{\rm ch}=4\ln2\,L_*^2$, which equals $S_{\rm dS}=3.31\times10^{122}$ by the channel-area normalization of Appendix~F.5. The renewed closure registers of the $2.6\times10^{183}$ cells in a Hubble volume carry $g_{\rm share,eff}$ each, $1.9\times10^{184}$ nats in all; renewal resets every register to $\rho_*$ at every tick independently of the state, so this entropy is a state-independent constant, stores no retrievable alternative, and drops out of type II$_1$ entropy differences. The exported history carries $g_{\rm share,eff}N_{\rm trials}$ nats, of order $10^{244}$ for the trial budget above. The one-step vertex acts on the geometric factor, the replaceable closure register, and the marked fiber (Q.9), and history export is forward only, so after export every operator of the completed local action commutes with the history factor. History therefore enters no local observable, no stress tensor, and no light sheet. The N-bound, the covariant entropy bound, and the type II$_1$ maximum constrain the retrievable present factor, which obeys them.

\paragraph{The denied conjecture and its consequences.}
The construction is incompatible with the proposal that the full Hilbert space of a de~Sitter patch has dimension $e^{S_{\rm dS}}$ \citep{Banks2000}. Three consequences follow. The global state never recurs, since the history factor grows at every tick. Recorded subsystems of the present factor relax (P.8) and return to low-entropy configurations only through the fluctuation bounds of P.10, so arguments that rely on Poincar\'e recurrence in a finite de~Sitter Hilbert space, including recurrence-driven dominance of fluctuated observers \citep{DysonKlebanSusskind2002}, do not apply. Horizon thermodynamics concerns the present factor and is unchanged. The finite-energy reset supplies a compatible energy ledger, while the physical source of fresh history capacity remains an independent completion.

\subsection*{P.10 Late-time re-entry and false records}

\paragraph{Re-entry in the specified dynamics.}
The dressing Hamiltonian of H.2 contains the channel occupations only through
$n_m=c_m^\dagger c_m$, with every other factor acting on the cloud sector. On
the seven-channel fermionic Fock space,
\[
\boxed{[H,n_m]=0\qquad(m=-3,\ldots,3)}
\]
exactly. Renewal acts on the cloud and closure registers and leaves the
occupations unchanged, and release is one-signed (O.13, P.7a). The zero-defect
released sector is therefore absorbing: within the specified dynamics it
cannot re-enter the committed phase or form false durable records.

\paragraph{Committed stock cannot be regenerated.}
A completed matter sector may admit number-changing processes. Matter created
in such processes carries no net conserved charge and receives no allocation,
and released matter keeps the coarse label of its retained carrier (O.15). The
cosmological committed stock therefore cannot be regenerated at any later
temperature.

\paragraph{Thermal bound for records written in matter.}
Define the content of a specified record by its equilibrium surprisal,
$K_R=-\ln\pi_{\rm eq}(R)$. For a recorded subsystem that has relaxed to its equilibrium measure (P.8), a
specified false record forms in one independent mixing window with
probability $e^{-K_R}$; the physical content of this statement is that the
equilibrium measure is uniform over the record's microscopic realizations. A
durable record of content $K$ written in matter also requires matter defects
of energy at least $KE_{\rm nat}$, with $E_{\rm nat}\simeq0.74$~MeV per nat
from Mass--Entropy Equivalence and the conversion of P.4. Detailed balance
with a bath at temperature $T$ bounds that formation by
$e^{-KE_{\rm nat}/k_BT}$. The event must satisfy both bounds, so
\[
\boxed{
P_{\rm false}(K)\le\min\!\left(e^{-K},\,e^{-KE_{\rm nat}/k_BT}\right).}
\]
Multiplying the two factors would count the same multiplicity twice (G.7). The
thermal factor dominates below $T\simeq8.6\times10^9$~K. At $2.7$~K,
$E_{\rm nat}/k_BT\simeq3\times10^9$; at the asymptotic de~Sitter temperature
$T_{\rm dS}=2.2\times10^{-30}$~K, $E_{\rm nat}/k_BT_{\rm dS}\simeq3.9\times10^{39}$.

\paragraph{Status of the far future.}
In the specified dynamics late-time re-entry vanishes exactly. If the
completed matter sector admits number-changing processes, the lightest
allowed number-changing commitment event sets the rate of matter re-entry,
which is horizon-thermal with an exponent of at least $3.9\times10^{39}$ per
nat; for electron--positron creation the event weight is
$e^{-5.4\times10^{39}}$. This suppression is finite, so exact exclusion over
an unbounded future requires either the number conservation of the specified
dynamics or a record measure restricted to histories with durable records.
Regeneration of the committed stock is excluded in both cases.

\subsection*{P.11 Verification register}
\begin{itemize}
\item \emph{Reproduced twice by independent code, with all printed
digits matched:} the $1680$-state spectrum and multiplicities;
$\eta_*$, $g_{\rm share,eff}$, $\langle K^2\rangle$, and
${\rm Var}(K^2)$; and the rooted $140/69$ reduction.
\item \emph{Checked by an independent fresh instance:} the P.3
relative-entropy identity, Jensen excess, sufficiency result, and exact
decomposition.
\item \emph{Exact pair controllability certificate:} the P.6
spin-$3/2$ pair-generator closure has rank $255$ over both
$\mathbb F_{1009}$ and $\mathbb F_{10007}$, certifying
$\mathfrak{su}(16)$ controllability. The calculation uses integer
$\mathfrak{sl}(2)$ generators and exact finite-field elimination; an
independent implementation would provide a further check.
\item \emph{Network composability:} the tensor-product Lie-algebra lemma
proves the three-slot and connected-network result analytically once the
edge $\mathfrak{su}(16)$ certificate is granted.  A deterministic modular
BFS on the first nontrivial triple reaches rank $4095$ over
$\mathbb F_{1009}$. This numerical check is not a premise of the analytic
theorem. The diagonal-bridge control saturates the exact block-diagonal bound
$1023$.
\item \emph{Exact affine-information theorem:} P.5 proves monotonicity throughout the positive fixed-marginal affine family. The finite logarithmic sign inequality and the analytic KL derivative are kernel checked in the finite-statistics and information modules. The P.4
trial budget remains a single-instance result requiring a fresh check before
ledger entry.
\item \emph{Analytic mixing and reset audit:} for a closed cohort under a
common positive-semidefinite advection--diffusion operator, direct
differentiation gives $\dot I(A;X)\le0$ and hence
$\dot K_{\rm mix}^{X}\ge0$. The common-overlap decomposition independently fixes
the maximal input-independent reset component and proves
$e(\mathcal K)\le K_{\rm mix}$; the two quantities need not coincide.
\item \emph{Finite shell recurrence and overlap:} stationary return-time
averaging gives the exact arithmetic factor $720^N$ on the full-support shell
register, and the characteristic function of one centered Gaussian gives the
positive coherent overlap $(2/7)^{N^2}$.
\item \emph{Excluded triggers:} a threshold fixed at the present $a_0$, a
single carrier atom or ionization state, and the $0.588/0.625$ fold
multipliers (O.9).  None implements collective loss of a baryonic progenitor
record.
\item \emph{Self-gravitating progenitor-tagged measurement (O.12):}
the Gadget-4 run measures nonzero information loss at entry and a slower
mixing tail. Convolved normalization depends on the retention gate and
label definition. The gate comes from the feedback-resolved occupation of
O.12; resolution and missing gas physics in the kernel remain open.

\item \emph{Unrestricted-tensor-product gate calculations (P.8):}
on the spin-$3/2$ pair with all nine $J_i\otimes J_j$ generators at fixed
angles, the commutant is one-dimensional, Kraus products span $M_{16}$ after
three ticks, and the gap is $1.5\times10^{-3}$ per tick; the isotropic and
$J_z\otimes J_z$ families leave $8$ and $48$ conserved directions. On
$\operatorname{Inv}(V_{3/2}^{\otimes4})$, of dimension four, every transported
face bilinear vanishes to $3\times10^{-31}$ and the node angle operators
generate $\mathfrak{su}(4)$. One face pair with the four native gluings
conserves only the slot exchange, for spin $3/2$ and for $j_{\rm eff}=3$; with
the $K^2$ cross term between two face pairs the commutant is trivial
(normalized commutator singular values $1$, $0.91$, $0.81$).
\item \emph{Exact occupation conservation (P.10):} $\max_m\|[H,n_m]\|=0$ on
the seven-channel fermionic Fock space of H.2.
\item \emph{Label-ledger enumeration (O.15):} breadth-first search over every
sequence of formation, release, mixing, merging, heating, and cooling on three
carriers (two conserved charges and one neutral pair) reaches $96$
post-formation states, none with unlabelled conserved charge and none with
allocation on the neutral carrier. An added label-stripping isometry makes
unlabelled charge reachable and produces $12$ allocation-raising transitions
when the plasma supplies the debit; with the plasma held cold it produces none.
\item \emph{Tied-abundance likelihood (Section~20 and O.12):}
the current saved fit uses Planck 2018 plik\_lite TT,TE,EE, two low-$\ell$
temperature bins and a Gaussian optical-depth prior, DESI DR2 BAO,
and DES-Dovekie with a free offset below $z=0.1$. It gives $\chi^2=2227.329$ against $2226.77$ with
five versus six cosmological parameters. The background-matched PPF
fluid and fixed $A_{L,{\rm eff}}=0.979$ are approximations; the fit
does not include a joint ACT-lensing reconstruction likelihood.

\item \emph{Direct linear cosmology (O.3 and O.12):} reference-split invariance and an independent growth calculation verify the geodesic-clock implementation. With matched lensing the direct fit gives $\Delta\chi^2=+0.85$ for DES-Dovekie and $\sum m_\nu=0.06\,$eV, and $-0.39$ with a free low-redshift supernova offset; the depletion completion leaves its cold linear branch unchanged.
\item \emph{Common formal proofs (H.10, H.15, N.8, O.3, and R.1):} isometric reduction, source-adjoint reconstruction, positive allocation, finite Gaussian integration and source elimination, and the retained angular-marginal update are checked in one Lean project. Its analytic and numerical certificates retain the stated physical and continuum hypotheses.

\item \emph{Registered and sealed:} the conjecture relating
lepton-ladder records. For generation $N$, it predicts mutual information
between channels equal to the ladder's quadratic term,
$N^2\ln(7/2)$ per pass. The muon supplies a zero-freedom pass/fail test.
\item \emph{Acceleration-normalization fork accounting (Appendix~L):} the eleven standard compact volumes, the four abundance maps, the conditional hit probability $0.05$, and the scenario Bayes factors, computed in one script.
\item \emph{Reduced vertical model (Section~16):} on NGC~6946 it reproduces the full director and QUMOND solutions to $0.02$ in $\Upsilon$ before its application to NGC~628.
\item \emph{Halo-shape combination and debris tilt (Section~16):} a random-effects combination with leave-one-out checks, and test-particle integrations in six potentials with two orbital anisotropies, including tilted-halo and disk-free controls. Three fixed-count stock phase realizations test eight-Gyr live retention with reciprocal disk torque. Separate diagnostics use matched stellar moments, local streaming subtraction, higher multipoles, density reconstruction, and radial and vertical force comparisons. The phase seeds and retention criterion were fixed before their runs; the matched-shape diagnostics were supplementary. Full force convergence and a survey likelihood remain open.
\item \emph{Dirac frame flips (Appendix~H.9a):} the parity, time-flip, and inversion bilinears and the incompatibility of a compensating physical spinor map have kernel-checked finite-matrix proofs. Random-frame evaluations provide separate numerical controls.
\item \emph{O.12 release history and kernel (Appendix~O.12):} an independent implementation of the first-entry history reproduces the tabulated $F_Z$ to $1.5\%$ and $I_{\rm rel}$ to $1.3\%$; the full-likelihood profile in $f$ uses the Section~20 pipeline; the SPH kernel analysis subtracts a shuffled-label baseline from the plug-in mutual information and reports three partitions, three gates, and two label definitions.
\item \emph{Native lattice fermion and orientation (Appendix~H.9a):} zero-mode scans for spin $0$, $\tfrac12$, and $\tfrac32$; tracing of the three nodal loops and their windings; Weyl charges after a Zeeman splitting, with slice Chern numbers; the second-neighbor and handedness-weighted terms; the lift-class representative on random frames; and the spectrum of a $7^3$-cell network with coherent and refreshed orientations; the first-order splitting on the diamond graph; the balanced body-centred disphenoid complex and barycentric subdivisions of periodic Poisson--Delaunay triangulations, with manifold, coloring, and bipartiteness checks; the color-drift identity; zero-mode counts, speeds, and splittings; and dispersion under twisted boundary conditions; the Dirac--K\"ahler operator on periodic complexes, with Betti-number zero-mode counts, sector speeds under Voronoi--Delaunay and regular-cell weights, odd-in-$k$ fits, self-averaging over samples, and constant per-degree rescalings.
\item \emph{de~Sitter counts (P.9):} $S_{\rm dS}$, the channel-area cut entropy, the closure-register and history totals.
\item \emph{Vacuum entanglement and defect response (P.3):} the face entropies and mutual information of the square-root closure vacuum, the seven exclusion responses, and the unmarked-face KL and reduced relative entropies, reproduced by two independent implementations to all printed digits.
\item \emph{Bridge spectrum and Gauss projection (H.10, H.11):} exact symbolic evaluation of $\operatorname{spec}A_R$ reproduces $\operatorname{rank}A_R=8$ and $\operatorname{Tr}A_R^2=38$; on $\operatorname{Inv}(V_3^{\otimes4})$, $P_cJ_{cf}P_c$ and $P_cJ_{cf}^aJ_{cf}^bP_c-4\delta^{ab}P_c$ vanish to $10^{-12}$, and the shape spectrum is $\{-12,-11,-9,-6,-2,3,9\}$.
\item \emph{Carrier Green function (Section~24.7):} the $128^3$ periodic-Laplacian amplitude and the smeared-Coulomb shell fit, reproduced by two independent implementations.
\end{itemize}
The remaining complexity claims reduce to two open tasks: the structured GFT gate measure and the register-support scale. None of these
results changes the inside, reviewer, or outside probability books.


\section*{Appendix Q: Gauge Invariance and the Geometry--Capacity Completion}

The ordinary gravitational branch uses a microscopic geometry--capacity
completion with two representation factors. The finite capacity alphabet
cannot serve as a geometric Gauss index, so geometry and capacity occupy
separate factors of the same substrate. An equivariant lock relates the
selected $V_3$ capacity channel to geometric spin and propagates under
maximal blocking.

The construction also supplies the finite matching and marked-transfer
operators that set the physical scale. Its coefficient-free dilation handles
collective loss of a baryonic source-clock carrier. Identity-free capacity
cannot write a which-stream record, and faithful renewal fixes unit rerouting
after carrier failure. The progenitor-mixing history remains an astrophysical
input to be measured. In the equilibrium vacuum, the capacity factor adds no
graviton polarization. On a boundary cut, it fixes the conversion between one
capacity channel and geometric area.

The chain is
\[
\boxed{
\begin{gathered}
V_{3/2}^{L}\otimes V_{3/2}^{R}
\xrightarrow{\;Q_f\;}
V_3^{\rm cap}\oplus(V_0\oplus V_1\oplus V_2)
\\[-1mm]
\text{capacity data}
\quad\not\equiv\quad
\text{geometric Gauss data}
\\[-1mm]
\mathcal H_{\rm geom}\otimes\mathcal H_{\rm cap}\otimes\mathbb C^2_{\rm orient}
\xrightarrow{\;L_3\;}
\text{one relationally locked geometry--capacity description}
\\[-1mm]
V_3^{\otimes N}\xrightarrow{\;F_N\;}V_{3N},
\qquad
F_N^{\rm geom}L_3^{\otimes N}
=e^{i\phi_N}L_{3N}F_N^{\rm cap}
\\[-1mm]
\text{large-spin geometric witness}
\longrightarrow
\text{Einstein/TT semiclassical branch},
\end{gathered}}
\]
On the selected metric-Regge branch, this chain has a closed massless-TT infrared theorem and an exact perfect-action primitive realization. The connected normalized patches of Q.11b give a concrete Euclidean CDT gluing tensor, and H.15 supplies a finite joint geometry--orbital operator. The host-underdetermination theorem shows that the capacity premises do not uniquely fix these geometric tensors or their measure. Factorization and the one-cell lock therefore define a selected action-level realization of Postulate~I. The underdetermination result, Schur uniqueness, maximal-channel propagation, equilibrium rigging-map lift, and the conditional area normalization are established below or by the cited exact lemmas.

\subsection*{Q.1 Two different meanings of closure}

Two structures called ``closure'' appear in the microscopic theory and must be kept separate.

\emph{Gauss closure} is the exact constraint associated with the local gauge redundancy of the geometric factor.  Its generator annihilates gauge-invariant geometric states.  It removes redundant descriptions and carries the standard spin-network/spin-foam representation data.

\emph{Capacity closure}, or admissibility closure, is the physical statistical weighting of the finite routed capacity alphabet defined in Section~6.  Its invariant $K^2$ measures the mismatch of a routed channel configuration from the regular tetrahedral closure pattern.  It is not a constraint: every state in the nonempty ensemble has strictly positive $K^2$, and different allowed states receive different Gibbs weights $e^{-\eta K^2}$.

Their roles are
\[
\boxed{
\text{Gauss closure: redundancy/constraint},
\qquad
\text{capacity closure: physical weighting}.}
\]
Identifying them would erase the nonuniform capacity ensemble and the
stationarity equation that fixes $\eta_*$. Q.4 gives the selected minimal
separation.

\subsection*{Q.2 Primitive face matching and the selected seven-state channel}

The capacity factor entering the completion is the transmitted multiplet already selected by the primitive oriented face-matching problem of Appendix~B.  Two descriptions of one shared face each carry a primitive spin
\[
j_0=\frac32,
\qquad
\mathcal H_f
=V_{3/2}\otimes V_{3/2}
=V_0\oplus V_1\oplus V_2\oplus V_3.
\]
After transporting the two local frames into one relative orientation, the positive mismatch operator is
\[
Q_f
=\frac12\left[
(\mathbf J_L-\widetilde{\mathbf J}_R)^2
-2j_0I
\right].
\]
On the coupled irrep $V_J$,
\[
Q_f\big|_{V_J}
=\frac12\left[12-J(J+1)\right]I_{V_J},
\]
so for $j_0=3/2$,
\[
\boxed{
Q_f=6P_0+5P_1+3P_2,
\qquad
\ker Q_f=V_3.}
\]
Equivalently the finite-width matching transfer
\[
T_s=e^{-sQ_f}
=e^{-6s}P_0+e^{-5s}P_1+e^{-3s}P_2+P_3
\]
converges to $P_3$ as $s\to\infty$.  Thus the transmitted capacity channel is
\[
\boxed{V_3^{\rm cap},\qquad \dim V_3=7,}
\]
while the complete nonmaximal information is
\[
\boxed{
Q_{\rm cap}=V_0\oplus V_1\oplus V_2,
\qquad \dim Q_{\rm cap}=9.}
\]
The independent marked present/history construction of Appendix~H.9 produces exactly the same multiplicity-free representation
\[
\mathcal H_{\rm mark}=V_1^P\otimes V_1^H=V_0\oplus V_1\oplus V_2.
\]
This representation identity accounts for the recurring seven-state
transmitted channel and nine-state marked complement. It does not identify
either object with geometric magnetic data.

The value $j_0=3/2$ has the grade stated in Appendix~B: it is uniquely forced
within the selected sharp reversible carrier incidence; for a different
microscopic vertex it remains the unique surviving point of the stated
discrete audit. No smallest-spin rule is used here.

\subsection*{Q.3 Same-factor impossibility theorem}

Suppose one tries to reuse the geometric Gauss representation factor itself as the capacity alphabet by identifying the capacity operator with the Casimir of the Gauss generator.  The Gauss generator vanishes on gauge-invariant physical states, so its Casimir must vanish there.  The certified capacity spectrum instead satisfies
\[
K^2\ge\frac{122}{3}>0
\]
on every state of the nonempty ensemble.  Hence
\[
\boxed{K^2_{\rm admissibility}\neq C^2_{\rm Gauss}.}
\]
Applying the geometric gauge projector only after inserting the Gibbs weight does not evade the contradiction: on physical states a Gauss-Casimir factor would reduce to unity, erasing the nonuniform distribution and removing the stationarity equation that defines $\eta_*$.  The numerical spine therefore requires capacity data that are physical but gauge inert with respect to the local geometric redundancy.

This theorem rules out the literal reading in which the capacity labels were identified with the open Peter--Weyl magnetic index of the geometric vertex.  It does \emph{not} by itself dictate a unique alternative.  The factorized construction below is the selected minimal realization because it leaves standard geometric gluing intact while preserving every finite capacity calculation.

\subsection*{Q.4 Selected factorization: one substrate, two bookkeeping factors}

As the selected action-level realization of Postulate~I, the one-cell bookkeeping space is
\[
\boxed{
\mathcal H_{\rm cell}
=\mathcal H_{\rm geom}
\otimes\mathcal H_{\rm cap}
\otimes\mathbb C^2_{\rm orient}.}
\]
Here $\mathcal H_{\rm geom}$ carries the ordinary gauge-reduced simplicial GFT/spin-foam data: geometric spins, magnetic indices, intertwiners, and the variables required by the geometric vertex.  The factor $\mathcal H_{\rm cap}$ carries the finite routed capacity alphabet, its admissibility operator, marks, and renewal variables.  The final doublet carries the two microscopic orientations.

The tensor product separates operator roles within one substrate. Postulate~I
is implemented by the equivariant lock on the two relational frames in Q.5.
The coupled theory must also pass the no-extra-mode and Ward-identity audits
of Appendix~H.10. The capacity factor is not a second freely propagating
geometry.

The four capacity ports draw without replacement from the seven transmitted weights,
\[
b_i\in\{-3,-2,-1,0,1,2,3\},
\qquad
b_i\neq b_j\quad(i\neq j),
\]
so
\[
\dim\mathcal B=P(7,4)=840,
\qquad
\boxed{\Omega=2P(7,4)=1680}
\]
after orientation doubling.  No geometric intertwiner multiplicity enters this count.  Consequently the EPRL/GFT factor is free to retain the full intertwiner structure required by its dynamics; the finite $1680$ count belongs only to the capacity/orientation code.

The local geometric gauge action is
\[
U_{\rm gauge}(h)
=U_{\rm geom}(h)\otimes I_{\rm cap}\otimes I_{\rm orient}.
\]
Every capacity observable therefore has the form
\[
O_{\rm cap}=I_{\rm geom}\otimes O\otimes I_{\rm orient},
\]
and, by the tensor-product form of the factorization, obeys the strong commutation relation
\[
\boxed{[O_{\rm cap},U_{\rm gauge}(h)]=0.}
\]
Thus the injectivity predicate, $K^2$, $Z$, $\langle K^2\rangle$,
$\mathrm{Var}(K^2)$, the Gibbs weights, and $\eta_*$ survive geometric gauge
projection unchanged. The factorization therefore leaves the finite
numerical spine unchanged.

\subsection*{Q.5 The one-cell equivariant lock and what it means physically}

The two bookkeeping factors must describe the same relational orientation on the physical branch.  Positive primitive fusion supplies $V_3^{\rm cap}$; the selected geometric branch carries the corresponding $j=3$ rotational representation $V_3^{\rm geom}$.  The action-level lock is a nonzero $SU(2)$-equivariant isometry
\[
L_3:V_3^{\rm cap}\longrightarrow V_3^{\rm geom}
\]
satisfying
\[
\boxed{
L_3U_{\rm cap}(g)=U_{\rm geom}(g)L_3
\qquad(g\in SU(2)).}
\]
The \emph{existence} of this identification is the selected realization choice.  Its form is then rigid.  Since both sides are irreducible copies of $V_3$,
\[
\operatorname{Hom}_{SU(2)}(V_3,V_3)\cong\mathbb C,
\]
and Schur's lemma gives, after choosing a common relational basis,
\[
\boxed{L_3=e^{i\phi}I.}
\]
The common phase is conventional unless another microscopic operator compares the two factors coherently.  Therefore the lock introduces no continuous physical coupling.

The lock acts on the relational representation data entering the coupled
vertex. It does not project the whole tensor product onto a diagonal span
$|m\rangle_{\rm geom}|m\rangle_{\rm cap}$. The dynamical test is whether the
locked interaction preserves the metric constraints without creating an
independent light capacity pole. Appendix~H.10 states those audits.

The marked complement has the same multiplicity-free structure,
\[
Q_{\rm cap}=V_0\oplus V_1\oplus V_2.
\]
Each irrep occurs once, so any equivariant lock to the corresponding spatial-rotation content of the geometric response is unique blockwise up to phases.  The matching of representations is exact, and the construction uses only this matching.

\subsection*{Q.6 Maximal blocking propagates the lock uniquely}

Appendix~B supplies exact canonical maximal-channel maps that preserve the
one-cell lock under coarse graining:
\[
F_N^{\rm cap}:(V_3^{\rm cap})^{\otimes N}\to V_{3N}^{\rm cap},
\qquad
F_N^{\rm geom}:(V_3^{\rm geom})^{\otimes N}\to V_{3N}^{\rm geom}.
\]
The highest-spin representation $V_{3N}$ occurs with multiplicity one in $V_3^{\otimes N}$.  Therefore the one-cell lock induces a unique block lock.

\paragraph{Lemma (equivariant lock under maximal blocking).}
There exists a normalized equivariant map
\[
L_{3N}:V_{3N}^{\rm cap}\to V_{3N}^{\rm geom}
\]
and a phase $e^{i\phi_N}$ such that
\[
\boxed{
F_N^{\rm geom}L_3^{\otimes N}
=e^{i\phi_N}L_{3N}F_N^{\rm cap}.}
\]

\emph{Proof.} Both sides are normalized $SU(2)$-equivariant maps from $(V_3^{\rm cap})^{\otimes N}$ into the unique highest-spin copy $V_{3N}^{\rm geom}$.  Multiplicity one implies that their restrictions to the maximal irrep are proportional.  Normalization makes the proportionality constant a phase.  $\square$

Thus
\[
V_3^{\rm cap}\stackrel{L_3}{\longleftrightarrow}V_3^{\rm geom}
\]
propagates automatically to
\[
\boxed{
V_{3N}^{\rm cap}\stackrel{L_{3N}}{\longleftrightarrow}V_{3N}^{\rm geom}}
\]
on the selected maximal-channel ray.  For spin coherent states,
\[
F_N|3,\mathbf n\rangle^{\otimes N}=|3N,\mathbf n\rangle,
\]
so the same relational normal is preserved by coherent blocking.

This representation-theory result does not show that every interacting RG
trajectory projects dynamically onto the maximal channel. The primitive
positive matching and the diagnostic gaps of Appendix~B select that ray in
the stated homogeneous/coherent model. It is therefore an exact canonical
maximal-channel blocking ray, not a unique full RG flow.

\subsection*{Q.7 Relational-frame covariance and the 840/1680 bookkeeping}

A fixed channel basis is a relational-frame convention, distinct from the geometric Gauss redundancy.  Let $R\in SU(2)$ denote the relational orientation of a cell and define
\[
P_{\rm cap}(R)
=U_{\rm cap}(R)P_{\rm cap}U_{\rm cap}(R)^\dagger,
\qquad
K^2(R)=U_{\rm cap}(R)K^2U_{\rm cap}(R)^\dagger.
\]
Under a common global rotation $G$,
\[
R\mapsto GR,
\qquad
K^2(GR)=U_{\rm cap}(G)K^2(R)U_{\rm cap}(G)^\dagger.
\]
Rank, spectrum, traces, and all finite state-counting invariants are therefore frame independent.  On the capacity factor alone,
\[
\operatorname{rank}P_{\rm cap}(R)=840,
\qquad
\operatorname{spec}K^2(R)=\operatorname{spec}K^2.
\]
Including the independent orientation doublet,
\[
P_{\rm code}(R)=P_{\rm cap}(R)\otimes I_{\rm orient},
\qquad
\boxed{\operatorname{rank}P_{\rm code}(R)=1680.}
\]
The weighted trace satisfies
\[
\operatorname{Tr}\!\left[P_{\rm code}(R)e^{-\eta K^2(R)}\right]
=\operatorname{Tr}\!\left[P_{\rm code}e^{-\eta K^2}\right],
\]
where $K^2$ acts trivially on orientation.  Choosing $R=I$ in the finite enumeration therefore fixes a relational frame without breaking physical rotational covariance.  Neighboring cells compare those frames through the transported face holonomy $g_{LR}$ already used in Appendix~B.

The integer tetrahedral witness makes the same point without representation language.  For
\[
v_1=(1,1,1),\quad v_2=(1,-1,-1),\quad
v_3=(-1,1,-1),\quad v_4=(-1,-1,1),
\]
one has
\[
\left|\sum_i b_iv_i\right|^2
=4\sum_i b_i^2-\left(\sum_i b_i\right)^2,
\]
which is $3|\mathbf c(b)|^2$, three times the mean nonclosure of Appendix~B. $K^2$ adds the label variance $\sum_i[j(j+1)-m_i^2]$, so the witness reproduces the classical part of the closure polynomial.  It is a relational capacity invariant, not a geometric Gauss generator.

\subsection*{Q.8 Division of labor between geometry and capacity}

The factorization is easiest to read as a division of operator roles.  The two columns are two sets of variables used by the selected microscopic realization of one substrate.

\begin{center}
\begin{tabular}{p{0.42\textwidth}p{0.48\textwidth}}
\hline
\textbf{Geometric factor} & \textbf{Capacity/record factor}\\
\hline
local Gauss redundancy and gauge projection & physical admissibility weights $e^{-\eta K^2}$\\
spin/intertwiner data and simplicial gluing & seven transmitted channel labels and four-port injectivity\\
Regge phase, metric constraints, and TT propagation & defects, exclusion, renewal, marks, and record bookkeeping\\
geometric area/flux observable & cut entropy and finite capacity weights\\
all-state geometric rigging-map problem & exact stationary renewal projector on the empty register\\
\hline
\end{tabular}
\end{center}

The equilibrium capacity sector has an exact consequence for the physical vacuum support.  Its renewed quantum state is the faithful mixed state
\[
\rho_*=\sum_b p_{\eta_*}(b)|b\rangle\langle b|,
\]
not a pure rank-one state.  On the $\rho_*$-GNS/operator Hilbert space define
\[
\omega_*(A)=\operatorname{Tr}(\rho_*A),
\qquad
\Pi_*(A)=\omega_*(A)I.
\]
Then
\[
\boxed{
\Pi_*^2=\Pi_*,
\qquad
\Pi_*^\dagger=\Pi_*,
\qquad
\operatorname{Ran}\Pi_*=\mathbb CI.}
\]
Thus the replaceable equilibrium closure register contains exactly one stationary transfer direction.  This does not erase the full-support mixed state $\rho_*$ and does not exclude separately retained or conservation-protected collective capacity fields.

If
\[
\eta_g:D_g\to D_g^*
\]
is any positive geometric rigging map, the empty capacity decoration lifts it exactly by
\[
\boxed{
\eta_{\rm full}(\phi\otimes A)[\psi\otimes B]
=\eta_g(\phi)[\psi]\,
\langle A,\Pi_*B\rangle_{\omega_*}.}
\]
After quotienting null states and completing,
\[
\boxed{
\mathcal H_{\rm phys}^{\rm full,vac}
\cong\mathcal H_{\rm phys}^{g}.}
\]
So the factorized capacity decoration does not create an additional vacuum constraint solution or graviton polarization once a geometric rigging map exists.  Persistent defect/marked sectors are outside this vacuum theorem and remain subject to the coupled interaction audits.

The same distinction applies to vacuum amplitudes.  After homogeneous centering, the normalized reference amplitude obeys
\[
\mathcal A_{\rm full}^{\rm ref}=\mathcal A_{\rm geom}^{\rm ref},
\]
but for a generic empty geometry $T$,
\[
\mathcal A_{\rm full}[T]
=\mathcal A_{\rm geom}[T]\,\widehat Z_{\rm cap}[T].
\]
Centering removes the source-independent homogeneous extensive term; it does not force every geometry-dependent finite capacity contribution to unity.


\subsection*{Q.9 Action-level realization: from matching operator to physical scale}

The capacity/record factor has an explicit action-level realization. The
finite ensemble fixes the admissible states, the matching operator fixes
their relative cost, and a local field action selects the homogeneous phase.
The marked vertex then specifies the microscopic update. Its positive
transfer spectrum supplies an energy gap, and the electron anchor converts
the dimensionless dynamics into $L_*$ and $G_*$. In schematic form,
\[
\boxed{
\begin{gathered}
\text{finite ensemble}\to\text{operator}\to\text{action}\to\text{selected phase}\\
\to\text{transfer spectrum}\to\text{physical scale}
\end{gathered}}
\]
Appendices~B and H contain the derivations and countervertex audits. The
equations needed for the action-level chain are collected here.

\paragraph{Phase selection from the primitive matching operator.}
For $j_0=3/2$, the positive face mismatch derived in Q.2 is
\[
Q_f=6P_0+5P_1+3P_2,
\qquad
\ker Q_f=V_3.
\]
Promote the complete primitive fusion space to a field $\Phi$ and consider the minimal homogeneous action
\[
\boxed{
\Gamma_f
=
\int d^dx\left[
Z|\partial\Phi|^2
+r\,\Phi^\dagger\Phi
+\kappa\,\Phi^\dagger Q_f\Phi
+\frac{u}{2}(\Phi^\dagger\Phi)^2
\right],
\qquad Z,\kappa,u>0.}
\]
Writing $\rho_J$ for the occupation of the coupled irrep $V_J$ and $\rho=\sum_J\rho_J$, its homogeneous potential is
\[
V
=r\rho+\frac{u}{2}\rho^2
+6\kappa\rho_0+5\kappa\rho_1+3\kappa\rho_2.
\]
For $r<0$ the unique support of the homogeneous minimum is therefore
\[
\boxed{
\rho_3=-\frac{r}{u},
\qquad
\rho_0=\rho_1=\rho_2=0,}
\]
with positive relative masses
\[
\boxed{
m_{0|3}^2=6\kappa,
\qquad
m_{1|3}^2=5\kappa,
\qquad
m_{2|3}^2=3\kappa.}
\]
Within this minimal action, the same positive mismatch operator that defines
coherent face matching selects the seven-state $V_3$ sector as the unique
homogeneous phase. The displayed action controls phase selection in the
representation-space/GFT sector. It does not supply the complete interacting
Lorentzian geometric GFT parent.

\paragraph{Fresh-state preparation and the marked record fiber.}
The stationary capacity measure fixes the fresh-state amplitude
\[
A_*(b)=Z^{-1/2}\exp\!\left[-\frac{\eta_*}{2}K^2(b)\right],
\qquad
|A_*(b)|^2=p_*(b).
\]
Its Gaussian linearization supplies the normalized three-component closure strand used by the marked present/history register.  The nine marked states are represented by the bilinear incidence
\[
\boxed{
V_Q
=\sum_{a,b=1}^{3}
 d_{ab}^\dagger\,\xi_a^P\xi_b^H+\mathrm{h.c.}}
\]
with $V_1^P\otimes V_1^H=V_0\oplus V_1\oplus V_2$.  The reduced Gram matrix is $I_9$, so equal marked weights follow from the two-strand contraction rather than from the dimension count alone.  This is the explicit dynamical realization of the same nine-state complement that appeared representation-theoretically in Q.2.

\paragraph{Inherited one-mark quadratic action.}
Take the invariant real form
$Q_{\mathbb R}=V_{0,\mathbb R}\oplus V_{1,\mathbb R}\oplus
V_{2,\mathbb R}$ and let $L:Q_{\mathbb R}\to
\mathcal H_{{\rm mark},\mathbb R}$ be the exact carrier-incidence isometry.
Set $d=L(P_Q\Phi_{\mathbb R})$. Because the mark is the retained complement
itself, the quadratic marked Hessian is the real push-forward of the
$Q$-restricted Hessian,
\[
\boxed{
\mathcal H_D=L\mathcal H_QL^T,
\qquad
\mathcal H_Q=\mathcal K_{\rm kin}+M_c^2I_Q
+\kappa(6P_0+5P_1+3P_2).}
\]
This follows by changing variables in $\delta^2\Gamma_f$; an independent
defect-only quadratic term would be a different, nonminimal action.  For a
local relative frame with $P_Q(x)=L(x)^TL(x)$ and $P_3(x)=I-P_Q(x)$, the same substitution gives
\[
\boxed{
A_\mu=L\partial_\mu L^T\in\mathfrak{so}(9),
\qquad
B_\mu=P_3\partial_\mu L^T,}
\]
and the exact orthogonal decomposition is
\[
\boxed{
Z\,\partial(P_Q\Phi_{\mathbb R})^T\partial(P_Q\Phi_{\mathbb R})
=Z(D_\mu d)^TD^\mu d+Z\,d^TB_\mu^TB^\mu d.}
\]
The second-order Born--Huang matrix vanishes for a fixed $P_Q$ and is nonzero along frame directions
with $Q\leftrightarrow V_3$ leakage. H.11 includes it in the full real
symplectic generator rather than dropping it in a tangent-only projection.
The $6{:}5{:}3$ mismatch eigenvalues determine mass-squared differences.
Appendix~H.11 takes the positive square root of the real Hamiltonian, obtains
the initial polar complex structure $J_0=-K_0(-K_0^2)^{-1/2}$, transports it
exactly as $J(t)=S(t,t_0)J_0S(t,t_0)^{-1}$ under the full symplectic flow, and
identifies the physical relative-phase generator as
$S_\omega=\Omega_{\rm ref}-(\operatorname{Tr}\Omega_{\rm ref}/9)I$.
Then
$\operatorname{Lie}(\mathfrak{so}(9),iS_\omega)=\mathfrak{su}(9)$,
giving the one-mark $\mathbb C^9/\mathbb{CP}^8$ state geometry without
assuming complex quantum kinematics at the start.

For each occupied unordered channel pair $e=(m,m')$, the two directed scalar returns have contraction norm
\[
u=2\eta_*\left(\frac{2}{7}\right)^2
=\frac{8\eta_*}{49},
\qquad
\alpha=\sqrt{1-u}.
\]
The square root is fixed at amplitude level by the canonical reversible dilation; using $1-u$ in the amplitude would double count the no-event probability.  On the 21 unordered pairs of the seven-channel branch this produces the exact internal factor $\alpha^{21}$.

\paragraph{Finite marked-transfer action and one-step vertex.}
Using hard-core marked occupation $n_d$ and auxiliary present/history Grassmann strands $\Gamma_{mm'}$, Appendix~H.9 gives the finite marked action
\[
\boxed{
\begin{aligned}
S_{\rm mark}
={}&g_{\rm share,eff}n_d+\Lambda_dn_d(n_d-1)\\
&+\frac12\sum_{m<m'}
\Gamma_{mm'}^T
\left[1+n_dn_mn_{m'}(\alpha-1)\right]J
\Gamma_{mm'},
\end{aligned}}
\qquad \Lambda_d\to+\infty .
\]
The auxiliary Majorana variables, equivalently represented by Grassmann integration, encode the antisymmetric present/history return; they are not additional propagating particles.  Let $W_*$ prepare the fresh closure state and reversibly export the previous replaceable register.  Let $\Pi_Q$ project onto the unmarked plus nine marked response states, and let $\widetilde U_e$ be the controlled unitary dilation on pair $e$.  The complete one-step finite transfer is
\[
\boxed{
\mathcal A_v
=\mathcal G_vW_*\Pi_Q
\exp\!\left[-g_{\rm share,eff}n_d-\Lambda_dn_d(n_d-1)\right]
\prod_{m<m'}\widetilde U_{mm'}.}
\]
Here $\mathcal G_v$ denotes the geometric gluing tensor supplied by the selected host.  The factors following $\mathcal G_v$ are the explicitly specified finite capacity/record decoration.  Their internal trace gives
\[
\boxed{
\zeta_*
=9e^{-g_{\rm share,eff}}
\left(1-\frac{8\eta_*}{49}\right)^{21/2}
=0.005123584484947\ldots .}
\]
The quantity $\zeta_*$ is the finite marked weight of this reversible vertex,
not a phenomenological correction added after the calculation. The complete
edge-Hessian audit in Appendix~H.9 gives a conditioned kernel proportional to
$I_{21}$ on the internal incidence labels. Alternative contractions split
the spectrum and define distinct countertheories.

\paragraph{Positive transfer spectrum and the physical scale.}
Faithful full-support renewal gives the state-weighted seven-channel determinant
\[
r=e^{-7g_{\rm share,eff}},
\]
and the positive one-bit survival transfer
\[
T_{\rm surv}=1-r.
\]
Its Euclidean transfer Hamiltonian therefore has the exact raw positive gap
\[
E_{\rm raw}
=-\frac{\hbar}{\tau_*}\ln(1-r).
\]
The marked routing graph dresses the charged mode by
\[
\boxed{
Z_e=(1+\zeta_*)(1+7\zeta_*^2).}
\]
The additive-generator prescription of H.9 defines the effective charged transfer $T_e=\exp[(3/2)Z_e\ln(1-r)n_e]$. Its charged energy is
\[
E_e=\frac32Z_eE_{\rm raw}.
\]
Equating this lowest charged transfer gap with the electron rest energy $m_ec^2$ fixes the cadence,
\[
\boxed{
\tau_*
=-\frac32Z_e\frac{\hbar}{m_ec^2}
\ln\!\left(1-e^{-7g_{\rm share,eff}}\right),}
\]
and hence
\[
\boxed{
L_*=c\tau_*
=-\frac32Z_e\lambda_e
\ln\!\left(1-e^{-7g_{\rm share,eff}}\right).}
\]
The ordinary gravitational normalization then follows from the same length,
\[
\boxed{
G_*\equiv\frac{c^3L_*^2}{\hbar}.}
\]
The logical chain is therefore explicit:
\[
\boxed{
Q_f
\Rightarrow
\Gamma_f
\Rightarrow
V_3\text{ phase}
\oplus Q\text{ retained dynamics}
\Rightarrow
\mathcal A_v
\Rightarrow
\zeta_*,Z_e
\Rightarrow
T_{\rm surv},E_e
\Rightarrow
\tau_*,L_*
\Rightarrow
G_* .}
\]
No additional continuous coefficient enters between the finite marked vertex and the electron-anchored scale.

\paragraph{Scope and remaining condition.}
The action-level chain specifies the finite capacity/record microdynamics
used for scale setting: stationary renewal, marked response, reversible
routing, the positive charged transfer gap, and the conversion to $L_*$.
The remaining problem lies in the coupled host. Its geometric tensor
$\mathcal G_v$ must preserve the metric constraints, Ward identity, physical
inner product, and causal structure without adding a source-coupled light
pole. Appendix~H.10 states these audits. The status is
\[
\boxed{
\begin{gathered}
\text{finite capacity/record action and transfer: \textbf{specified and audited}},\\
\text{$V_3$ homogeneous phase of the minimal action: \textbf{proved}},\\
\text{$Q$ one-mark }\mathbb C^9/\mathbb{CP}^8\text{: \textbf{proved in the stable quadratic sector}},\\
\text{exact moving-frame symplectic transport: \textbf{proved}},\\
\text{finite Born/instrument/history reconstruction: \textbf{proved given complete physical records}},\\
\text{full coupled Lorentzian geometric parent: \textbf{conditional/open}}.
\end{gathered}}
\]

\subsection*{Q.9a Marked GFT dilation and collective carrier loss}

The GFT action already contains the finite nine-state marked complement and
the exact memoryless renewal map.  What it cannot contain on the
identity-free capacity factor is a which-stream trigger.  A permutation of
identical stream assignments leaves the occupancy history unchanged, so
the capacity environment exports the same state and
\[
|e_i\rangle=|e_j\rangle,\qquad\gamma_{ij}=1.
\]
Thus neither the prepared-ray complement nor a fold-purity factor is a
cosmological control.

The slow source factor must instead include the retained baryonic state
that carries the Lagrangian label \(A\).  On the enlarged
geometry--capacity--matter space define
\[
\boxed{
\mathsf E_A
=\mathbf1[\text{the retained baryonic state no longer resolves
progenitor }A].}
\]
The effect is collective and present-tense.  Atomic-cooling and
photoheated-gas retention enter the matter dynamics because irreversible
mixing requires a common retained well; they are not local step functions
of the primitive capacity register.  Ionization alone, shell crossing,
and ordinary environmental record production leave \(\mathsf E_A=0\).

Faithful full-support renewal fixes the stock change once
\(\mathsf E_A=1\).  The memory family
\[
K_a(b,b')=a\delta_{bb'}+(1-a)p_*(b')
\]
has stationary averaged conditional entropy $\sum_b p_*(b)H(K_a(b,\cdot))<g_{\rm share,eff}$ for $0<a\le1$;
hence the only admissible failed-carrier update is \(a=0\), a complete
fresh draw.  The corresponding event coefficient is
\[
\boxed{f_{\rm reroute}=1.}
\]

Let \(|\alpha\rangle_M\), \(\alpha=1,\ldots,9\), denote the orthonormal
marked basis and let \(W_{\alpha A}\) be orthogonal history embeddings.
The normalized GFT/source dilation is
\[
\boxed{
V_A^{\rm rel}|\psi\rangle
=\sqrt{I-\mathsf E_A}\,|\psi\rangle|0\rangle_M
+\frac13\sum_{\alpha=1}^{9}
|\alpha\rangle_MW_{\alpha A}\sqrt{\mathsf E_A}\,|\psi\rangle.}
\]
It obeys
\[
(V_A^{\rm rel})^\dagger V_A^{\rm rel}
=I-\mathsf E_A+\frac19\sum_{\alpha=1}^{9}\mathsf E_A=I.
\]
The factor \(1/3\) normalizes nine marked outputs; it is not the released
stock fraction.

Under refresh-first ordering, H.11 gives
\[
(\operatorname{id}_R\otimes\mathcal E_*)(\rho_{R,{\rm ren}})
=\rho_R\otimes\rho_*,
\]
so conditioning on carrier failure leaves
\[
\langle K^2\rangle_{\rm rel}
=\langle K^2\rangle_{\eta_*}
\]
and the energy share is
\[
\boxed{
\eta_{\rm drain}
=\frac{\varepsilon_\xi}{\varepsilon_K}
\frac{3/2}{\langle K^2\rangle_{\eta_*}}
=\eta_*}
\]
under the selected common conversion $\varepsilon_\xi=\varepsilon_K$.
The common conversion and post-refresh ledger are premises. They are
logically separate from the exact unit rerouting coefficient.

The action-level sequence is therefore
\[
\boxed{
\begin{array}{c}
\text{an elementary baryonic progenitor record is lost}\\
\Downarrow\\
\mathsf E_A=1\\
\Downarrow\\
\text{complete renewed stock update}\\
\Downarrow\\
\text{normalized }16\to7+9\text{ marked dilation}\\
\Downarrow\\
\text{coarse drain }E_{\rm drain}=\eta_*E_{\rm rel}.
\end{array}}
\]
The GFT action fixes the last three arrows. Additive record pricing assembles
the elementary events into a coarse released-stock fraction measured by
\[
\epsilon_{\rm mix}=1-\frac{I(A;Y)}{H(A)}.
\]
Appendix~O.12 supplies an absorbing first-entry assembly history, a measured
passive-tracer information observable, and the closed-cohort resolved-position
monotonicity theorem together with the coarse-observable caveat. O.2--O.3 specify the capacity-clock projection and its constrained
covariant action. The finite release isometry fixes conditional register
probabilities; the source-record identification, common conversion,
full-allocation energy ledger, and gravitational stress contact remain
physical matching prescriptions. A native calculation must reproduce these
assignments and the galaxy-scale mixing kernel.

\subsection*{Q.10 From the locked block to the Einstein witness}

The representation-theory results above explain how the fixed capacity representation reaches the large-spin geometric regime without running the microscopic spin.  On the maximal ray,
\[
J_{\rm block}=3N\longrightarrow\infty,
\]
while the primitive face remains $j=3$.  The block-lock lemma ensures that the capacity and geometric coarse representations continue to carry the same relational orientation.

Three external results support different parts of the geometric branch.
Proper-EPRL constructions remove the unwanted asymptotic sectors at the
vertex level and retain the Regge exponential in the controlled nondegenerate
sector \citep{EngleProper2013,EngleVilenskyZipfel2016}. Han-type spin-foam
constructions provide a coupled large-spin/refinement witness approaching
Einstein--Regge dynamics. In the controlled linearized construction, the
low-energy excitations exhaust the smooth solutions of linearized Einstein
gravity and give the two helicity-$2$ graviton modes
\citep{Han2017,HanHuangZipfel2019}. Area-Regge continuum analysis gives the
same leading graviton dynamics, with the first length-metric correction at
order $a^2C^2$ \citep{DittrichKogios2023}. These are separate compatibility
results, not one theorem about a decorated vertex.

The selected metric-Regge branch therefore has controlled semiclassical
Einstein witnesses at block scale. H.9a gives the full finite-cell-to-Einstein
derivation. H.10 proves the strong-resolvent TT theorem \((H_{\rm TT})\) and
its exact gapped fixed-\(j=3\) perfect-action pullback. It also proves the
complementary underdetermination result: because the selected host supplies \(\mathcal G_v\)
independently, the finite capacity premises cannot determine its geometric
measure or nonperturbative phase. Q.8 shows that the equilibrium capacity
factor creates no separate vacuum obstruction.

The retained side of the same $7+9$ split is equally explicit.  The
maximal $V_3=\ker Q_f$ channel carries the Einstein TT witness, whereas the
positive-mismatch complement $Q=V_0\oplus V_1\oplus V_2$ carries the marked
record whose inherited dynamics gives $\mathbb C^9/\mathbb{CP}^8$ in H.11.
The two results are complementary consequences of one fusion event and
describe the two sides of its $7+9$ split.

Schematically, on the controlled linearized witness branch,
\[
\mathbb K_{\rm phys}(k)
\sim Z_Gk^2P_{\rm TT}
+O(a^2k^4)
+O(\Lambda_{\rm sf}^{-1}k^2),
\]
This expression summarizes the common low-energy content of the separate continuum results. It is not a theorem for one decorated proper-EPRL model.

\subsection*{Q.11 Coherent capacity bonds and channel area}
\label{aqg:bonds}

The local capacity ensemble admits a coherent spatial realization within
its existing state space. Separate the unordered set of four labels, its
routing to the four ports, and the orientation doublet:
\[
\mathcal H_{\rm cap,o}\cong
\mathbb C^{35}\otimes\mathbb C[S_4]\otimes\mathbb C^2.
\]
The closure value for an unordered set \(s\) is
\[
K^2(s)=48-\frac{(\sum_{m\in s}m)^2-\sum_{m\in s}m^2}{3}.
\]
It is constant on the 48 routing and orientation states of that set.
An equivariant encoding identifies
\[
\mathbb C[S_4]\otimes\mathbb C^2
\cong\mathcal F(\mathbb C^4)\otimes V_{\rm std},
\qquad
\mathcal H_{\rm cap,o}\cong\mathcal F_4\otimes\mathbb C^{105}.
\]
Here \(V_{\rm std}\) is the three-dimensional standard representation of
\(S_4\). The Fock occupations are virtual port observables encoded in the
routing register. The original injective seven-label alphabet is retained.

For the conjugacy classes \(1,(12),(12)(34),(123),(1234)\), the character
of the exterior-algebra representation is
\((16,0,0,4,0)\), and that of \(V_{\rm std}\) is
\((3,1,-1,0,-1)\). Their product, \((48,0,0,0,0)\), is the character
of two regular representations. This proves the encoding. Fermionic
permutation signs enter the character identity. A group average followed by
polar normalization gives an explicit unitary intertwiner, with covariance
residual \(1.8\times10^{-14}\) in the finite calculation.

Write
\[
P_s=\frac{e^{-\eta_*K^2(s)}}{\sum_{s'}e^{-\eta_*K^2(s')}},
\qquad
\sigma=\bigoplus_s\frac{P_s}{3}I_3.
\]
The original local density matrix becomes
\[
\rho_{\rm cell}=\frac{I_{16}}{16}\otimes\sigma.
\]
Every original routing microstate still has probability \(P_s/48\).
Thus the closure spectrum, \(\eta_*=0.0298668443935\), the sharing entropy,
and the local charged-gap traces are unchanged.

On a four-port graph \(\Gamma\), a selected coherent capacity vacuum is
\[
|\Omega_\Gamma\rangle
=\prod_{e=(L,R)}
\frac{c^\dagger_{L,e}+c^\dagger_{R,e}}{\sqrt2}|0\rangle,
\qquad
\rho_{\Gamma,*}
=|\Omega_\Gamma\rangle\langle\Omega_\Gamma|
\otimes\bigotimes_v\sigma_v.
\]
Each virtual mode belongs to one link. A consistent fermionic ordering
fixes the signs. Every cell has the local state displayed above, while a
cut through \(N_\partial\) bonds has conditional mode entropy
\[
S_{\rm cut}=N_\partial\ln2.
\]
The local entropy separates as
\[
g_{\rm share,eff}=4\ln2+S(\sigma)
=2.77258872224+4.64721130133.
\]
The residual registers are independent between cells in this state.
Their entropy is a volume contribution. Original label observables and
virtual cut observables act on the same encoded register, so their
entropies cannot be added as independent boundary species.

The coherent bond is the unique ground state of
\[
H_e=(n_L+n_R-1)^2+n_L+n_R-2n_Ln_R
-(c_L^\dagger c_R+c_R^\dagger c_L),
\]
whose spectrum is \(0,1,1,2\) in the specified bond units. This selects
a spatial state beyond the local Gibbs marginal. A separable mixture of
the two one-copy occupations has the same local marginals and a different
entanglement structure; single-copy occupation alone does not select the
coherent vacuum.

\paragraph{Renewal of the correlated state.}
Independent resets of the two ends erase the bond. Joint renewal instead
uses
\[
\mathcal R_\Gamma(\rho)=\rho_{\Gamma,*}\operatorname{Tr}\rho.
\]
Its isometry prepares \(|\Omega_\Gamma\rangle\), prepares local
purifications \(\sum_r\sqrt{\sigma_r}|r\rangle|r\rangle_{\rm fresh}\),
and exports the complete old register into history. The marginal map at
each cell remains the required memoryless replacement kernel. The retained
orbital branch is acted on identically by this preparation. Local cadence
and rest-gap traces therefore survive. H.14 checks a background-wide
schedule with commuting disjoint resets, and H.15 supplies a joint
orbital transfer. Renewal across different geometric event records
and a common continuum matter cone remain dynamical requirements.

\paragraph{Absolute angular transfer.}
The half-bond state is \(\rho_h=I_2/2\). Specify a positive continuous
angular semigroup whose unnormalized full-turn transfer equals this state:
\[
K_h(2\pi)=\rho_h,\qquad
K_h(\theta)=e^{-b\theta}I_2,\qquad
b=\frac{\ln2}{2\pi}.
\]
The full-turn equality fixes the absolute amplitude. Normalization of a
density matrix after tracing would leave that amplitude undetermined.
At a metric hinge with total wedge angle \(2\pi-\delta_h\),
\[
\operatorname{Tr}\prod_{w\supset h}K_h(\theta_{hw})
=e^{b\delta_h},\qquad
W_h=e^{bN_h\delta_h}.
\]
At replica angle \(2\pi n\), the same tensor gives
\(\operatorname{Tr}K_h(2\pi n)=2^{1-n}=\operatorname{Tr}\rho_h^n\).
For a cell the residual factor is
\(\operatorname{Tr}\rho_{\rm cell}^n
=2^{4(1-n)}\operatorname{Tr}\sigma^n\).
At fixed cell count and residual state its explicit volume factor can be
removed when reading the conditional cut entropy.
Gravitational corner constructions provide related continuum
interpretations \citep{AQGTakayanagi2020,AQGBotta2020}; the full-turn
normalization here is a specified microscopic rule.

Let \(A_h=A_{\rm ch}N_h\) on the selected metric continuation. The Regge
identity \(\int R\,dV=2\sum_h A_h\delta_h\) then gives
\[
\frac{I_E}{\hbar}=-\frac{b}{2A_{\rm ch}}\int R\,dV,\qquad
\boxed{\frac{\hbar G_0}{c^3}
=\frac{A_{\rm ch}}{8\pi b}
=\frac{A_{\rm ch}}{4\ln2}.}
\]
The curvature and replica coefficients are evaluations of one tensor.
The bare value \(G_0\) consequently controls both at this level. Establishing
the same relation for \(G_{\rm R}\) requires a common renormalized state and
effective action. The angular rule gives no independent freedom to absorb
a bulk/cut discrepancy in another measure.

\paragraph{Additive area and the matched branch.}
The standard kinematic puncture eigenvalue is
\(8\pi\gamma_{\rm BI}L_*^2\sqrt{J(J+1)}\).
For the coherent maximal-fusion ray \(J_N=3N\), use the additive flux area
\[
A_N=24\pi\gamma_{\rm BI}NL_*^2,\qquad
Z_A(N)=\sqrt{\frac{3N}{3N+1}}.
\]
Here \(A_N=Z_A(N)A_{\rm LQG}(3N)\); the primitive spin is fixed.
Combining this area with the angular rule yields
\[
\frac{G_0}{G_*}=\frac{6\pi\gamma_{\rm BI}}{\ln2}.
\]
The matched choice \(G_0=G_*\) gives
\(\gamma_{\rm BI}=\ln2/(6\pi)\) and
\(A_{\rm ch}=4\ln2\,L_*^2\). The event-volume calibration of C.8
instead supplies \(A_{\rm ch}\) before imposing that equality. Its
provisional result is \(G_0/G_*=1.03002\). These are distinct calibration
conditions for the same relation. The response number \(n_{\rm hor}\)
in F.5 remains a scalar normalization quantity.

\subsection*{Q.11a Graph reconnections and collective area coordinates}
\label{aqg:reconnections}

A binary switch reconnects two bonds among four active ports on existing
cells. For \(M\) switches let \(n\) count the switched patterns. Retain every
microscopic pattern and use the normalized symmetric count states
\[
|n\rangle=\binom Mn^{-1/2}\sum_{x:n(x)=n}|x\rangle.
\]
The collective raising operator satisfies
\[
S^+|n\rangle=\sqrt{(n+1)(M-n)}|n+1\rangle,\qquad
T^+=S^+\big[(n+1)(M-n)\big]^{-1/2}.
\]
The denominator acts first, and the terminal value is zero. Thus \(T^+\)
is a unit count shift on the retained sector. Products of two such
operators implement paired reconnections.

Independent microscopic flips \(t\sum_i(1-X_i)\) select the binomial
count wavefunction. A flat count wavefunction requires the additional
diagonal term
\[
t\left[\sqrt{n(M-n+1)}+\sqrt{(n+1)(M-n)}-M\right].
\]
Its microscopic amplitudes are proportional to
\(\binom Mn^{-1/2}\). Positive edge factorizations verify this statement
on the complete switch space. The first count gap obeys
\(M\Delta_M/t=4.3666,4.4710,4.4973,4.5039\) for
\(M=32,128,512,2048\). This is motion in count configuration space;
spatial wave dispersion requires the inter-cell action.

For a TT count coordinate with
\(\Delta q=4\sqrt2/N\), the quadratic metric action supplies oscillator
mass \(V_B/(32\pi G_0)\). At \(\omega\sqrt{NA_{\rm ch}}=0.2\), the
finite-count gap divided by \(\omega\) is
\(0.999998653,0.999999663,0.999999916\) for
\(N=16384,65536,262144\). This is a quantization of the selected quadratic
action. Independent switch-pattern correlations can carry extensive
entropy, even when the entropy of the aggregate count grows logarithmically.

\paragraph{Coherent transport between graph frames.}
Let \(U_\Gamma|{\rm ref}\rangle=|\Omega_\Gamma\rangle\), and define
\[
\mathcal D=\sum_\Gamma|\Gamma\rangle\langle\Gamma|\otimes U_\Gamma.
\]
Directly exporting a graph-dependent bond state multiplies graph
coherences by
\(\langle\Omega_{\Gamma'}|\Omega_\Gamma\rangle\).
For the two tested matchings this overlap vanishes; the purity of a
prepared graph superposition falls to \(0.5338\).
The frame-covariant renewal first applies \(\mathcal D^\dagger\), exports
and renews the common reference capacity state, evolves the retained graph,
and reapplies \(\mathcal D\). If
\(\mathcal J|\Gamma\rangle=|\Gamma\rangle|\Omega_\Gamma\rangle\),
\[
U_{\rm phys}=\mathcal D(U_{\rm geom}\otimes I)\mathcal D^\dagger,
\qquad U_{\rm phys}\mathcal J=\mathcal J U_{\rm geom}.
\]
The finite intertwining residual is \(7.7\times10^{-16}\). Persistent
geometric orientation and refreshed capacity orientation have separate
operator roles. The coherent connection \(U_\Gamma\), its locality, and its
extension to marked paths are specified geometric data.

\paragraph{Nonlinear count amplitudes.}
The six-area reconstruction of N.10 supplies \(K^{AB}(q)\). Decompose it
into directions \(e_A\) and \(e_A+e_B\):
\[
c_{AB}=K^{AB},\qquad
c_A=K^{AA}-\sum_{B\ne A}K^{AB}.
\]
The midpoint conductance on direction \(v\) is
\(c_K\mu(q_{\rm mid})c_v(q_{\rm mid})/(\Delta q)^2\), where
\(c_K=16\pi G_0/V_{\rm ref}\) in tick units and \(\mu\) is the measure
of N.11. With the symmetric divergence matrix \(L\) and
\(W=\operatorname{diag}\mu\), the selected flat-count operator is
\[
H_{\rm count}=W^{-1/2}LW^{-1/2}
-c_K\operatorname{diag}\frac{21}{16\sqrt{\det g}}+V_R.
\]
Replace each count shift by \(T_A^+\) or \(T_A^+T_B^+\) and every diagonal
coefficient by its count function. The resulting microscopic operator
intertwines exactly with \(H_{\rm count}\). Two switches for each of six
areas give 729 count states in all 4096 microscopic patterns, with six
single and fifteen paired move directions. The relative residual is
\(6.4\times10^{-17}\); the finite operator is Hermitian. Negative
conductances are allowed because the unreduced DeWitt kinetic form is
indefinite. The graph amplitudes are constructed from the specified metric
action, area map, and ordering.


\subsection*{Q.11b Connected primitive patches and the geometric history measure}
\label{aqg:connected}

The angular rule has a connected realization in networks of the primitive
spin-three geometric cells. In the Euclidean CDT regulator, the spatial
and mixed triangle areas are
\[
A_{\rm SL}=\frac{\sqrt3}{4}a^2,\qquad
A_{\rm mix}=\frac{\sqrt{4\alpha-1}}4a^2.
\]
If these areas contain \(m\) and \(n\) binary channels,
\[
\alpha=\frac{1+3(n/m)^2}{4}.
\]
For a four-port patch with \(V\) cells, \(E\) internal edges, and \(B\)
external ports, \(B=4V-2E\) is even. A mixed tetrahedron with face counts
\((1,n,n,n)\) consequently requires odd \(n\). After the isotropic choice,
the first allowed value in this patch construction is \(n=3\), giving
\[
A_{\rm SL}=A_{\rm ch},\qquad A_{\rm mix}=3A_{\rm ch},\qquad \alpha=7.
\]
This selects a Euclidean regulator branch. It does not identify mixed
Euclidean facet counts with a Lorentzian timelike area spectrum.
The calibration of C.8 remains
\(a=2.56812077553L_*\), \(A_{\rm ch}=2.85582456165L_*^2\),
and \(G_0/G_*=1.03002098319\) at the illustrative \(m_4=2.35\).

The required tetrahedron spaces are
\[
\begin{array}{c|c|c}
\text{slice split}&\text{face spins}&\text{intertwiner dimension}\\ \hline
4\!-\!0&(3,3,3,3)&7\\
3\!-\!1&(3,9,9,9)&7\\
2\!-\!2&(9,9,9,9)&19
\end{array}
\]
Spin nine is the maximal fusion of three primitive facets.
Let \(B_k\), \(k=0,\ldots,6\), be an orthonormal basis of
\(\mathcal I_3=\operatorname{Inv}(V_3^{\otimes4})\).
Contract internal edges with
\(E_{mm'}=(-1)^{3-m}\delta_{m,-m'}\).
The normalized three-facet fusion has magnetic-basis entries
\[
(F_3)_{R,r_1r_2r_3}
=\delta_{R,r_1+r_2+r_3}
\sqrt{\frac{\binom6{r_1}\binom6{r_2}\binom6{r_3}}{\binom{18}R}},
\qquad r_i=0,\ldots,6.
\]

\paragraph{Connected interfaces.}
Four cells arranged as
\((a,b,c,x),(x,d,e,y),(y,f,g,z),(z,h,i,j)\), with external bundles
\(a;(b,d,f);(c,e,h);(g,i,j)\), define a map to
\(\operatorname{Inv}(V_3\otimes V_9^{\otimes3})\).
Its rank remains seven after both consecutive pairs are restricted to
the 13-dimensional sectors of H.12; the input dimension is 169.
For the second interface, six cells form a cycle with external face
pairs \((0,1),(0,2),(1,2),(1,3),(2,3),(0,3)\).
Each coarse face receives three primitive facets.
Restricting three consecutive pairs gives
\[
R_{22}:(\mathbb C^{13})^{\otimes3}
\longrightarrow\operatorname{Inv}(V_9^{\otimes4}),
\qquad \operatorname{rank}R_{22}=19.
\]
Eighty coherent-state evaluations determine this map and sixteen
independent evaluations check it, with maximum relative residual
\(3.29\times10^{-15}\). These are numerical full-rank and contraction
checks. The interpretation of coherent intertwiners as finite-spin
polyhedral states retains their shape uncertainty
\citep{AQGBianchi2011}.

These interfaces are connected spin-network patches.
A finer tetrahedral ball with four connected boundary face disks is a
stronger geometric requirement. The tested six- and eight-tetrahedron
balls have rank thirteen, and the full-rank cycle has not been
identified with such a ball. The construction uses the more general
relational graph permitted by the microscopic ontology.

\paragraph{Normalized gluing.}
For either surjective patch map \(R_e\), define
\[
F_e=(R_eR_e^\dagger)^{-1/2}R_e,\qquad
I_e=F_e^\dagger,\qquad I_e^\dagger I_e=I.
\]
The Gram condition numbers are approximately \(9.68\) and \(25.81\);
the corresponding isometry errors are \(6.82\times10^{-15}\) and
\(2.02\times10^{-14}\).
The normalization fixes the complete matrix, including its shape
dependence. Inside the paired ground space, add
\[
H_{\rm block}=\Delta(P_{\rm pair}-P_{\rm code}),\qquad
P_{\rm code}=I_eI_e^\dagger\le P_{\rm pair},\qquad \Delta>0.
\]
This positive finite-patch interaction selects the desired code.
Indiscriminate overlapping copies would define a different interaction.
Geometric observables on the code are transported coarse operators;
equality to sums of unmodified fine volume operators is unproved.

For a normalized coherent shape packet \(u_e\), set
\(\widetilde u_e=I_eu_e\). Choose a four-simplex tensor consisting of its
boundary packets times
\[
w_\sigma=\exp\left[-b\sum_{h\subset\sigma}N_h\theta_{h\sigma}
-\lambda V_\sigma\right],\qquad b=\frac{\ln2}{2\pi}.
\]
Shared tetrahedra use adjoint packets with the same ordering, so each
interior packet contraction equals one.
An open hinge carries the half-turn normalization
\[
B_h(\theta)=\rho_h^{-1/2}K_h(\theta)
=e^{b(\pi-\theta)}I_2,\qquad
\operatorname{Tr}[\rho_h B_h(\theta_1)B_h(\theta_2)]
=e^{b(2\pi-\theta_1-\theta_2)}.
\]
Interior full-turn and boundary half-turn factors therefore yield
\[
\boxed{W(\mathcal T)=\exp\left[
b\sum_{h\in{\rm int}}N_h\delta_h
+b\sum_{h\in\partial}N_h\psi_h
-\lambda\sum_\sigma V_\sigma\right].}
\]
Here \(\psi_h\) is the boundary exterior angle.
The identity holds for every allowed coarse history by factorization,
with curvature coefficient \(1/(8\pi G_0)\).
On 1,920 causal slabs its maximum log-weight discrepancy is
\(1.46\times10^{-13}\). External scalar-source and boundary-sewing
checks agree at comparable precision.
Raw patch contractions would insert shape-dependent norms at internal
tetrahedra and change this measure.
The chosen packets and block normalization give an explicit Euclidean
existence construction; they are distinct from the standard EPRL
vertex and do not establish its Lorentzian representation lift.

\paragraph{Causal history sum.}
Specify the geometric counting measure by summing compatible causal
triangulations with their relabeling quotient and no additional
geometric weight:
\[
T(g_1,g_0;z_{41},z_{32})
=\sum_{[\mathcal T]\in\mathcal C(g_1,g_0)}
\frac{z_{41}^{N_{41}}z_{32}^{N_{32}}}{C(\mathcal T)}
W(\mathcal T).
\]
The domain, volume treatment, and symmetry factors are part of the
selected host. Uniform event placement within a given geometry does
not select this measure over geometries.
Boundary sewing gives
\[
T^{(2)}(g_2,g_0)=\sum_{g_1}C(g_1)T(g_2,g_1)T(g_1,g_0),\qquad
\widetilde T(g',g)=\sqrt{C(g')C(g)}\,T(g',g).
\]
At a finite symmetric cutoff, \(\widetilde T\) is real symmetric and
\(\widetilde T^\dagger\widetilde T\) is positive semidefinite.
One-step positivity and a continuum Lorentzian phase require separate
arguments. The microscopic isometry gives
\(T_{\rm enc}=J\widetilde TJ^\dagger\) and
\(T_{\rm enc}^n=J\widetilde T^nJ^\dagger\) for \(n\ge1\).
The finite capacity state is attached once; Monte Carlo move sequences
are sampling algorithms and supply no additional history multiplicity.

The canonical count kernel of N.11--N.14 and this covariant sum remain
distinct selected realizations. The signed-loop test in N.14 excludes
their simplest entrywise-positive identification in one basis.
Both can be tested without first proving that they are equivalent.

\subsection*{Q.12 Microscopic acceptance conditions}

A complete decorated geometric vertex must satisfy all of the following conditions.
\begin{enumerate}
\item \emph{Geometric branch.} The geometric factor retains its Gauss/intertwiner structure and a controlled metric limit. H.10 gives an EPRL/Regge witness; Q.11b gives a separately specified Euclidean CDT tensor realization. Their representation maps and Lorentzian limits must not be identified without a derivation.
\item \emph{Capacity branch.}  The finite factor carries the selected $V_3$ seven-state multiplet, four single-copy relational ports, the same $K^2$, and the same stationary measure used by the numerical spine.
\item \emph{Equivariant lock.}  Geometry--capacity interactions enforce the selected relational lock and its block propagation; they do not identify $K^2$ with a Gauss Casimir.
\item \emph{Decoration locality.}  Exclusion, renewal, marks, and record return act on the capacity/record factor, while geometric gluing retains its standard magnetic/intertwiner data.
\item \emph{Hard-core occupation.}  The four ports draw from one single-copy capacity multiplet, preserving the injective $P(7,4)$ count.
\item \emph{Orientation/parity.}  The two orientation sectors are exchanged by an exact microscopic $\mathbb Z_2$ symmetry.
\item \emph{Measure normalization.} The local capacity marginal retains its normalized closure weights and the homogeneous centering of Section~25. A coherent spatial state, its kinematic metric measure, and its physical rigging form are distinct objects. The graph connection and renewal must preserve retained coherence.
\item \emph{Source control.}  The slow baryonic source factor must carry a distinguishable Lagrangian progenitor label.  Irreversible loss of that label in the retained baryonic state defines $\mathsf E_A$ and enters the normalized nine-state dilation.  Prepared-ray phase failure, stream permutation, ionization alone, and generic environmental records are not sufficient controls.
\item \emph{Coupled spectrum and Ward structure.} The full decorated action introduces no extra source-coupled zero pole and has a controlled metric constraint limit. Its bulk response, cut coefficient, and matter coupling must yield one renormalized Newton coefficient. Matching that coefficient to $G_*$ is a physical test.
\item \emph{Charged orbital propagation.}  The retained defect has a local, gauge-covariant spatial transfer whose small-momentum pole has rest gap $\epsilon_*$ and invariant speed $c$. Its discarded renewal output carries no path address.
\item \emph{Physical-inner-product/causal structure.}  A Euclidean transfer realization must satisfy the relevant complete-transfer positivity/reflection condition; a directly Lorentzian realization must instead possess a positive rigging form, consistent causal gluing, and healthy retarded poles.
\end{enumerate}
The coherent angular construction supplies the bond state, pair
interaction, connected normalized packets, and count lift.
H.14--H.15 supply a finite joint orbital and clock realization.
The Einstein/ADM continuation, covariant host measure, and inverse
cochain metric are specified completion data. Their common
interacting physical state, causal continuum, and renormalized
normalization must satisfy the final three conditions.

\subsection*{Q.13 Status ledger}

The geometry--capacity chain has the following logical grades:
\begin{description}
\item[Primitive maximal-channel selection.] Proved for the stated $Q_f$; the $j_0=3/2$ identification carries the Appendix~B grade.
\item[Same-factor Gauss identification.] \textbf{Excluded.}
\item[Tensor-factor realization.] \textbf{Selected microscopic completion.}
\item[Existence of the one-cell relational lock.] \textbf{Selected realization condition.}
\item[Form and uniqueness of a nonzero lock.] \textbf{Proved by Schur's lemma.}
\item[Maximal block-lock propagation.] \textbf{Proved.}
\item[Gauge-invariant capacity count.] \textbf{Proved: $840$ capacity states and $1680$ after orientation doubling.}
\item[Empty capacity rigging-map lift.] \textbf{Proved exactly.}
\item[Selected-branch Einstein TT theorem.] \textbf{Closed}; the metric-Regge symbol converges to the massless two-helicity TT operator and the perfect-action pullback gives an exact gapped fixed-\(j=3\) realization.
\item[Selected-branch matter--gravity compatibility.] \textbf{Closed as a leading-order EFT compatibility statement}; one Einstein--Dirac action carries the shell mass operator, relativistic propagation, metric sourcing, and the $G_*$ normalization, while the joint perfect-action and renewal identities preserve the coupled persistent sector at finite regulator.
\item[Channel--area normalization.] Exact bare relation $\hbar G_0/c^3=A_{\rm ch}/(4\ln2)$ within the absolute angular rule. Physical area calibration and renormalized bulk/cut matching remain conditional.
\item[Finite capacity/record marked-transfer vertex.] \textbf{Specified and finitely audited.}
\item[Charged spatial transfer.] \textbf{Specified finite realization}; the joint geometry--cochain model preserves the native gap, positive clock, and retained coherence on the stated renewal background. The general spinor, gauge, species, and continuum-cone requirements remain open. On the native complex a common speed requires the Hodge weights of a flat realization (Appendix~H.9a, paragraph~2g).
\item[Source-controlled marked dilation.] \textbf{Specified and isometric}; additive record pricing fixes normalized progenitor-information loss as fractional released stock, faithful renewal fixes unit rerouting for every failed elementary block, and common closed-cohort advection--diffusion makes the fully resolved position loss monotone. The coarse physical kernel and its monotonicity, source-record identification, common energy conversion, and post-refresh ledger remain inputs.
\item[Factorized vertex.] \textbf{Capacity-side spectral and causal audits closed}; alternative nonfactorized realizations must preserve them.
\item[Unique geometric host from capacity premises.] \textbf{Excluded}; \(\mathcal G_v\) and its measure are independent host data.
\item[Coherent capacity vacuum.] Exact equivariant encoding and local-marginal preservation; the spatial bond state is selected.
\item[Primitive pair interaction.] Exact positive quartic parent and fusion kernel. The connected all-edge extension has a soft-mode upper bound and is excluded as the sole relativistic propagation mechanism.
\item[Metric and graph realization.] Six areas reconstruct one metric; normalized single and paired reconnections give an exact finite operator lift. The Einstein/ADM continuation supplies the amplitudes.
\item[Connected primitive geometry.] Full-rank chain and cycle maps, Gram-normalized boundary packets, and exact sewing reproduce the selected \(\alpha=7\) CDT weight. A literal fine tetrahedral ball and a Lorentzian representation lift are additional requirements.
\item[Joint orbital operator.] A positive enumerated geometry transfer and retained cochain operator preserve the rest gap and mass-derivative clock. The common isometric theorem preserves the retained algebra, trace, positivity, and polynomial dynamics; its source-family derivative fixes the additional moving-code terms. Exact code refinement is established for the chosen inverse metrics.
\item[Positive source and thermal map.] Exact mass, adjoint-coupling, covariance, and common-temperature quantum thermal identities for the stated reconstruction. Autonomous dynamics require generator closure.
\item[Positive capacity allocation.] Exact finite positivity and conservation, uniquely selected by composition preservation on the common flow. O.3 supplies the homogeneous exponential solution, local total-stock domain, and a selected positive multistream exchange with an expanding beam realization. Physical record pricing, the ensemble history, and the joint bound-galaxy dynamics remain inputs or further requirements.

\item[Quantum ordering.] Exact DeWitt measure and transported metric Weyl correction within the chosen lapse convention. Continuous-configuration commutators exactly quantize the regulated classical brackets.
\item[Constraint limits.] Numerical convergence on smooth nonlinear probes; exact finite spatial closure and uniform physical-state convergence are open.
\item[History energy and capacity.] An explicit energy-conserving reset uses a finite auxiliary energy band and zero-energy exported records (P.9). Unlimited fresh logical capacity and its metric-independent physical realization remain assumptions.
\item[Physical state and host phase.] A positive homogeneous volume--shape rigging form is constructed. The inhomogeneous physical constraints, metricity boundary, common gravitational refinement, interacting matter phase, and renormalized normalization remain open.
\end{description}

The Lean results for the same-factor Gauss contradiction of Q.3 and the finite $1680/K^2$ arithmetic concern the unchanged capacity factor, and the coherent completion preserves those hypotheses. The formal verification in R.1 covers the finite capacity arithmetic and the common reduction, reconstruction, allocation, and statistical identities. It preserves the distinction between mathematical consequences of a specified map and physical selection of that map. The coherent completion also has conditional continuum reductions and numerical constraint tests. A common interacting physical Lorentzian continuum remains to be established.

\section*{Appendix R: Reproduction Code}

This appendix contains the executable reproduction block for the manuscript's numerical spine, using the definitions stated in the paper. Appendix~S gives the electromagnetic extension and identifies its accompanying formal and numerical checks.

\subsection*{R.1 Microscopic geometry and constraint calculations}

The companion directory
\path{Matter sector/chain_clock_2026-10-13/unified_micro_model/}
contains the finite geometric construction. Running
\texttt{python run.py} executes nineteen checks and records their
outputs and a manifest. They use Python, NumPy, SciPy, and SymPy.
The groups correspond to the derivations as follows.

\begin{center}
\small
\begin{tabular}{L{0.24\textwidth}L{0.47\textwidth}L{0.19\textwidth}}
\hline
Calculation & Modules (with \texttt{.py} suffix) & Manuscript\\
\hline
Capacity state and normalization &
\path{capacity_bonds}, \path{renewal_event}, \path{corner_model},
\path{volume_state} & C.8, Q.11\\
Pair and metric response &
\path{collective_parent}, \path{regge_response},
\path{metric_fourier} & H.12--H.13\\
Graph and count dynamics &
\path{graph_frame}, \path{count_transfer},
\path{reconnection_count}, \path{reconnection_lift} & Q.11a\\
Local constraint complex &
\path{local_constraints}, \path{finite_count_ward},
\path{canonical_count} & H.13\\
Nonlinear geometry and quantization &
\path{area_metric_dynamics}, \path{metric_measure},
\path{nonlinear_spatial_constraints}, \path{metric_ordering},
\path{quantum_constraint_current} & N.10--N.12\\
\hline
\end{tabular}
\end{center}

The checks include controls that fail the desired condition: the
all-edge parent has soft modes, direct graph-dependent renewal loses
coherence, naive finite-count mobility violates linear constraints,
and the small Weyl-ordered reflecting box has no zero constraint state.
Successful execution reproduces these limits as well as the stated
identities. It is not a physical-state or phase proof.
The fluid/Dirac calculation and native profile/cut comparison are
separate companion calculations with the normalizations stated in
N.10 and J.11.

Further companion calculations extend these checks without changing
the nineteen-module reproduction set. Within the same directory,
\path{volume_shape_rigging.py} gives the homogeneous spectral modes
and original-coordinate residuals of N.13.
\path{reconnection_histories.py},
\path{conformal_loop_audit.py},
\path{conformal_source_response.py},
\path{conformal_integral.py}, and
\path{microscopic_history_measure.py} implement the history and
contour checks of N.14.
\path{clocked_reconnection_step.py} and
\path{local_clock_consistency.py} implement the auxiliary clock
and local-deformation checks of N.15.
\path{clock_volume_matching.py},
\path{covariant_causal_host.py}, and
\path{slab_boundary_sewing.py} supply the host dictionaries and
finite sewing diagnostics.

Three adjacent companion directories contain the connected and
joint constructions:
\begingroup\raggedright
\begin{itemize}
\item \path{causal_primitive_realization_2026-10-16/}:\newline
\path{connected_patch.py}, \path{cycle_patch.py},
\path{slab_realization.py}, and \path{ball_patch.py} implement the
patch maps, full Gram normalization, slab contractions, and
tetrahedral-ball rank controls of Q.11b.
\item \path{clock_matter_bridge_2026-10-16/}:\newline
\path{marked_clock.py}, \path{renewal_coherence.py},
\path{positive_clock_transfer.py}, and
\path{boundary_clock_diagnostic.py} implement the native source,
reset counterexample, background renewal, and positive-clock
checks of H.14.
\item \path{orbital_transfer_2026-10-16/}:\newline
\path{causal_boundary_transfer.py},
\path{harmonic_dressing.py}, \path{joint_orbital_clock.py},
\path{cochain_block_completion.py},
\path{joint_microscopic_lift.py}, and \path{whitney_orbital.py}
implement the enumerated geometric transfer, protected rest gap,
joint clock, exact cochain lifts, and dispersion controls of H.15.
\end{itemize}
\endgroup
Their saved results distinguish exact operator identities, numerical
rank and residual checks, and physical continuum conditions.
The held-source diagnostic of J.13 has separate archived inputs
and coverage requirements. None of these finite checks supplies
an equilibrated infrared force measurement or transfers a result
between different geometric ensembles.

\paragraph{Formal proofs and extension checks.}
The companion project, \url{https://github.com/jakechinitz/Entropic-EFT-Lean}, uses Lean 4.35.0 with a pinned mathlib revision. Its audit permits only \texttt{propext}, \texttt{Classical.choice}, and \texttt{Quot.sound}, with no admitted proofs, custom axioms, or \texttt{native\_decide}. Versioned build records give the library inventory; the manuscript-level scope is recorded separately in the claim ledger.

The upstream ledger contains 1{,}437 statements: 882 marked proved (six given a cited classical result), 90 proved in part, 59 numerical entries, 64 data comparisons, 152 physical premises, 136 not formalized, 19 simulation outputs, and 35 other entries. The citation audit resolves 1{,}028 referenced Lean names. Six extension modules add 119 explicit theorem declarations covering empirical bounds, the depletion constitutive identities, finite-work contact, warm sound-speed bounds, merger conservation, and conditional cluster inventories. The integrated 4.35 build, source scan, axiom audit, and ledger citation check pass; the six extensions also pass standalone Lean/mathlib 4.19 checks. Neither proof set certifies an observational likelihood or the physical selection of an effective contact.

The checked mathematics includes the following connections.
\begin{itemize}
\item \emph{Closure and scale:} finite counting, Gibbs differentiation, the unique global evidence maximum, and rational exponential/logarithm bounds give $0.0298668443934<\eta_*<0.0298668443936$ and $7.4198000220<g_{\rm share,eff}<7.4198000250$. The supplied marked and shell formulas give certified mass-ratio intervals and the exhaustive reduced-alphabet cutoff.
\item \emph{Marked transfer and records:} explicit spin-$3/2$ matrices realize the $7+9$ projector split, positive matching semigroup, and mismatch spectrum. Gaussian Fourier and Hermite-Gram identities realize the three-component response. Finite replacement, recovery, retained coherence, and history-capacity identities are proved under their displayed hypotheses.
\item \emph{Gaussian response and geometry:} the positive finite-dimensional Gaussian integral is derived from Lebesgue measure. Determinant normalization, the quadratic-source Legendre relation, source elimination, isometric pullbacks, and finite staggered constraint identities are checked. The physical Hessian and common continuum support remain model data.
\item \emph{Galactic and thermal response:} the spherical RAR is monotone and obeys the bounds in Section~15; its uncapped finite-mass limit gives the baryonic Tully--Fisher normalization. The unit-director density is nonnegative, radial cancellation selects $\kappa=1$, and the causal two-pole response has the stated static sum rule, including coincident rates. Carrier selection and thermal/metric contact remain physical premises.
\item \emph{Capacity, clocks, and release:} multiplicative lapse classification, reduced spherical identities, and horizon thermodynamic algebra are checked after their metric and area inputs. Positive internal allocation, integrating-factor balance, conserved-fluid algebra, finite information loss, and angular-marginal-preserving relaxation are proved. The extension checks establish the stated algebraic floor, regular-current, and conservative-contact identities. Constitutive selection, nonlinear spatial evolution, kinetic formation, and the compatible flux domain of F.7 remain physical requirements.
\end{itemize}
Compatibility supplements construct explicit nonempty effective models sharing the native closure parameter, positive stock budgets and expanding release fields. An infinite product measure realizes the native Gibbs histories; the reversible 720-state return witness is a separate process. Further certificates cover a positive inhomogeneous expanding clock/stock family, the multistream source and finite budget identities of O.3, and the finite-energy reset identities of P.9. The supplements have standalone Lean/mathlib 4.19 checks with the same axiom restrictions; the integrated build records specify their 4.35 coverage. Metric tensor reductions and finite-interval beam integrations are checked separately. These witnesses establish coexistence of their stated assumptions; they do not derive the metric from individual histories or construct a bound galaxy, cosmological release and native interacting gravity in one physical model. The detailed proof inventories remain in the reproduction materials.

The typed certificates of the v0.8 library (\texttt{Entropic.EFT}) for closure analysis, numerical bounds, and the orthogonal split have explicit instances. Structures containing only proposition labels carry no independent proof credit. Pseudo-Riemannian variational calculus, general unbounded-operator limits, unrestricted $SU(2)$ representation theory, and several Lie-closure and continuum constructions remain outside the formalization. The finite results do not establish the interacting Einstein phase, a universal matter cone, halo assembly, or the independent cosmological release history. Saved numerical studies and likelihoods retain their separate evidential status.

\subsection*{R.2 Local capacity and particle coefficients}

The script below recomputes the numerical consequences of the selected definitions and matching prescriptions: the $1680$-state spectrum, closure solution, sharing entropy, marked vertex, dressed scale, Joyce constant, edge-kernel chain, dressed charged-lepton ratios, slot mutual information, closed routing covariance, primitive route spectrum, adjoint channel, causal color-blocking factor, and alphabet-independent generation cutoff. It also verifies the algebraic normalization identity $n_{\rm hor}S_\infty^{\rm cell}=1/4$. CODATA constants are comparison inputs; no observational target is fitted inside the script. It runs under Python~3 with \texttt{numpy} and \texttt{scipy}. The complete Hessian and graph enumeration are supplied in the accompanying audit script.

{\footnotesize\begin{verbatim}
# ============================================================
# Reproduction script for the Entropic Scalar EFT manuscript.
# Recomputes the full numerical spine from stated definitions:
# the 1680-state K^2 spectrum, eta*, g_share,eff, marked
# vertex, dressed L* and G*, the
# Joyce constant, the 1/4 horizon identity, the edge-kernel
# chain, epsilon, sigma*, the lepton ladder, slot MI, and
# the conditional color-route algebra.
# Dependencies: Python 3, numpy, scipy.
# Printed values follow the stated definitions; assertions check identities.
# ============================================================
import itertools, math
from math import log, exp, pi, gamma, sqrt
from collections import Counter
import numpy as np
from scipy.optimize import brentq

# 1. Enumerate 1680 states: injective 4-of-7 labels m in -3..3, x2 parity
labels = range(-3,4)
states = list(itertools.permutations(labels,4))
print("P(7,4) =", len(states), " total with parity:", 2*len(states))

def K2(ms):
    S = sum(ms); Sig2 = sum(m*m for m in ms)
    return 48 - (S*S - Sig2)/3

spec = Counter()
for st in states:
    spec[round(K2(st)*3)] += 2   # parity doubling; key = 3*K2 to keep exact
print("Distinct K2 values and multiplicities (K2 as fraction /3):")
for k in sorted(spec): print(f"  K2={k}/3 = {k/3:.4f}  mult={spec[k]}")
print("Total multiplicity:", sum(spec.values()))

# 2. Solve closure condition <K2> = 3/(2 eta)
Ks = np.array([k/3 for k in sorted(spec)])
ns = np.array([spec[k] for k in sorted(spec)], dtype=float)
def avgK2(eta):
    w = ns*np.exp(-eta*Ks)
    return (w*Ks).sum()/w.sum()
f = lambda eta: avgK2(eta) - 3/(2*eta)
eta_star = brentq(f, 1e-4, 1.0, xtol=1e-16)
print("\neta* =", eta_star, " (paper: 0.02986684439352237)")

w = ns*np.exp(-eta_star*Ks); Z = w.sum(); p_class = w/Z
p_micro = np.exp(-eta_star*Ks)/Z
H = -(ns*p_micro*np.log(p_micro)).sum()
print("g_share,eff =", H, " (paper: 7.419800023570903)")
avg = avgK2(eta_star)
var = (p_class*(Ks-avg)**2).sum()
print("<K2> =", avg, " (paper 50.2229154254)   3/(2eta*) =", 3/(2*eta_star))
print("Var(K2) =", var, " (paper 15.6889750078)   a_UV = 1/Var =", 1/var)
print("C_cl = eta* <K2> =", eta_star*avg)

# 3. Decorated marked vertex, L*, G*
hbar=6.62607015e-34/(2*pi)  # exact SI h; hbar evaluated numerically
c=299792458.0; me=9.1093837139e-31  # CODATA 2022 electron anchor
lam_e = hbar/(me*c)
print("\nlambda_e =", lam_e, " (paper 3.86159267435238e-13)")
r = exp(-7*H)
L0 = -1.5*lam_e*math.log1p(-r)
q = 2/7
zeta = 9*exp(-H)*(1-eta_star*8/49)**(21/2)
Zmu = 1+zeta
Ze = (1+zeta)*(1+7*zeta*zeta)
Ztau2 = 1+q*zeta
Lstar = Ze*L0
print("fewer-channel energy bound =", exp(H)/Ze)
print("return-block log width =", math.log1p(7*zeta*zeta))
# Leading relative difference; the omitted term is O(r**2).
print("finite-event relative energy shift =", r*(Ze-1)/2)
print("e^{7g} =", exp(7*H), " (paper 3.602860521062804e22)")
print("(2/3)e^{7g} =", (2/3)*exp(7*H), " (paper 2.401907014041869e22)")
print("zeta* =", zeta, " (paper 0.005123584484947)")
print("Z_mu, Z_e, Z_tau2 =", Zmu, Ze, Ztau2)
print("L0 =", L0, " (paper 1.60771947114394e-35)")
print("L* =", Lstar, " (paper 1.61625370238286e-35)")
Gstar = c**3*Lstar**2/hbar
print("G* =", Gstar, " (paper 6.67428908131019e-11)")
G_codata = 6.67430e-11
LP = sqrt(hbar*G_codata/c**3)
print("L* vs Planck length:", (LP-Lstar)/LP*100, "%")
print("G* vs CODATA:", (G_codata-Gstar)/G_codata*100, "%")

# 4. Joyce constant
Gtet = 3*gamma(1/3)**6/(2**(14/3)*pi**4)
print("\nG_tet(0) =", Gtet, " (paper 0.4482203943883814)")
print("ln(7/6)/(pi*Gtet) =", log(7/6)/(pi*Gtet), " (paper 0.109472228)")
print("S_inf_cell = 3ln2/(32 pi Gtet) =", 3*log(2)/(32*pi*Gtet), " (paper 0.0461482516)")
print("n_hor*S_inf_cell =", (8*pi*Gtet/(3*log(2))) * (3*log(2)/(32*pi*Gtet)))
# Factorized j=3 horizon embedding: coherent/Regge channel-area normalization
gamma_BI = log(2)/(6*pi)
a_ch_over_L2 = 24*pi*gamma_BI  # equivalently 4 ln 2
ZA1 = sqrt(3/4)
print("gamma_BI = ln2/(6pi) =", gamma_BI, " (paper 0.0367726000254)")
print("a_ch/L*^2 =", a_ch_over_L2, " (paper 2.77258872224)")
print("Z_A(N=1) = sqrt(3/4) =", ZA1)
for N in (1, 2, 10, 1000):
    ZA = sqrt((3*N)/(3*N+1))
    Aphys_over_L2 = ZA * 8*pi*gamma_BI*sqrt((3*N)*(3*N+1))
    assert abs(Aphys_over_L2 - 4*N*log(2)) < 1e-10
print("area-blocking identity checked for N=1,2,10,1000")

# 5. Edge kernel chain
Jbare = 2/3*eta_star
Jtree = Jbare/3
Sig = 7+2/9
cloop = 1/(1+Jtree*Sig)
Jren = Jtree*cloop
print("\nJ_bare =", Jbare, " (paper 0.0199112296)")
print("J_tree =", Jtree, " (paper 0.0066370765)")
print("c_loop =", cloop, " (paper 0.95426)")
print("J_ren =", Jren, " (paper 0.00633348)")
print("gamma coeff 4 J_ren/(3 pi^2) =", 4*Jren/(3*pi**2), " (paper 8.556e-4)")

# 6. epsilon and cluster numbers
epsv = H/(4*pi**2)
print("\nepsilon = g/(4pi^2) =", epsv, " 1-eps =", 1-epsv)
print("committed-to-baryonic cap =", 1/epsv)
print("sigma* = pi/g =", pi/H, " (paper 0.42340665)")

# 7. Dressed charged-lepton ladder
mmu0 = 720*q; mtau0 = 720**2*q**4
mmu = mmu0*Zmu; mtau = mtau0*Zmu*Ztau2
print("\nbaseline m_mu/m_e =", mmu0, " baseline m_tau/m_e =", mtau0)
print("m_mu/m_e =", mmu, " (paper 206.768280237)")
print("m_tau/m_e =", mtau, " (paper 3477.343310)")
me_MeV=me*c*c/1.602176634e-13
print("PDG tau residual sigma =", (mtau*me_MeV-1776.93)/0.09)
print("historical baseline alpha_mu =", log(206.7683)/log(mmu0),
      " historical baseline alpha_tau =", log(1776.93/me_MeV)/log(mtau0))

# 8. mixing-time ratio
tau_star = Lstar/c; tau_e = hbar/(me*c**2)
print("\ntau* =", tau_star, " tau_e =", tau_e, " ratio =", tau_e/tau_star)

# 9. slot mutual information (marginal of slot under admissibility)
joint = Counter()
for st in states:
    wgt = exp(-eta_star*K2(st))
    joint[(st[0],st[1])] += 2*wgt
tot = sum(joint.values())
pj = {k:v/tot for k,v in joint.items()}
p0 = Counter(); p1 = Counter()
for (a,b),v in pj.items(): p0[a]+=v; p1[b]+=v
I = sum(v*log(v/(p0[a]*p1[b])) for (a,b),v in pj.items())
H0 = -sum(v*log(v) for v in p0.values())
print("\nslot MI =", I, " (paper 0.1545)   H(slot) =", H0, " I/H =", I/H0, " (paper 0.079)")

# 10. Conditional color-route checks
route_rows = []
route_weights = []
for st in states:
    m1,m2,m3,m4 = st
    route_rows.append([
        m1*m2 + m3*m4,
        m1*m3 + m2*m4,
        m1*m4 + m2*m3,
    ])
    route_weights.append(2*exp(-eta_star*K2(st)))  # parity doubling
route_rows = np.asarray(route_rows, dtype=float)
route_weights = np.asarray(route_weights, dtype=float)
route_weights /= route_weights.sum()
route_mean = (route_weights[:,None]*route_rows).sum(axis=0)
centered = route_rows-route_mean
Croute = (route_weights[:,None,None]
          * centered[:,:,None]*centered[:,None,:]).sum(axis=0)
P1 = np.ones((3,3))/3
P2 = np.eye(3)-P1
lam1 = np.trace(P1@Croute)
lam2 = np.trace(P2@Croute)/2
print("\nroute covariance =\n", Croute)
print("route singlet/doublet eigenvalues =", lam1, lam2)
print("(3/4) Var(K2) =", 0.75*var)
assert np.allclose(Croute, lam1*P1+lam2*P2, atol=1e-10)
assert abs(lam1-0.75*var) < 1e-10

Groute = np.eye(3)- (np.ones((3,3))-np.eye(3))/20
Graw = np.linalg.eigvalsh(Groute)
Kroute = (20/21)*Groute
Lroute = sqrt(13)/7*P1
print("primitive route Gram eigenvalues =", Graw)
print("lambda2/lambda1 =", Graw[-1]/Graw[0], " (paper 7/6)")
print("dilation residual =",
      np.linalg.norm(Kroute.T@Kroute + Lroute.T@Lroute-np.eye(3)))

Utest = np.diag(np.exp(1j*np.array([0.2, -0.5, 0.3])))
trU = np.trace(Utest)
return_projective = 1.5*(1-abs(trU)**2/9)
return_adjoint = (4/3)*(1-(abs(trU)**2-1)/8)
phase = np.exp(0.37j)
return_rephased = 1.5*(1-abs(np.trace(phase*Utest))**2/9)
print("projective/adjoint return costs =",
      return_projective, return_adjoint, return_rephased)
assert abs(return_projective-return_adjoint) < 1e-14
assert abs(return_projective-return_rephased) < 1e-14

t8 = 13/14
g_step = (2/3)*log(14/13)
block = 2*sqrt(6)
g_graph = block*g_step
print("t8 =", t8, " -ln(t8) =", -log(t8))
print("g_HK,step^2 =", g_step)
print("diamond blocking =", block, " g_graph,HK^2 =", g_graph)

# 11. Alphabet-independent generation cutoff
for excluded_count in (1, 2, 3):
    exclusion_variances = []
    for excluded in itertools.combinations(labels, excluded_count):
        reduced = [m for m in labels if m not in excluded]
        reduced_K2 = [K2(st) for st in itertools.permutations(reduced, 4)]
        exclusion_variances.append(float(np.var(reduced_K2)))
    nonzero = sum(value > 1e-12 for value in exclusion_variances)
    zero = sum(abs(value) <= 1e-12 for value in exclusion_variances)
    print(
        f"excluded labels = {excluded_count}: cases = {len(exclusion_variances)}, "
        f"nonzero = {nonzero}, zero = {zero}, "
        f"variance range = [{min(exclusion_variances)}, {max(exclusion_variances)}]"
    )
    if excluded_count < 3:
        assert nonzero == len(exclusion_variances)
    else:
        assert zero == len(exclusion_variances)
\end{verbatim}}

\clearpage
\section*{Appendix S: Electromagnetic Extension}
\label{app:electromagnetic}

The retained coherent matter sector of H.11--H.15 admits an explicit
electromagnetic extension with conserved charge, compact link variables, an
exact Gauss constraint, and the two transverse modes of the quadratic Maxwell
theory. The extension is specified as additional operators. The $1680$-state
closure ensemble and its gravitational coefficient chain stay fixed; the extra
sector's backreaction and microscopic capacity cost belong to a common
completion. Electric charge and the electromagnetic coupling are inputs of the
extension, and Sections S.4--S.10 identify what a native selection of them
would require.

The route transport of I.2 is the projective channel $X\mapsto UXU^\dagger$,
invariant under $U\mapsto e^{i\theta}U$. Its common phase becomes
electromagnetic data once a physical charged amplitude compares different
paths, and that amplitude, together with its field energy, is the dynamical
input the extension supplies.

\subsection*{S.1 Conserved matter and a physical phase connection}

On a fixed finite relational network let $\psi_x$ denote a fermionic
annihilation operator, with internal retained labels understood. A
number-preserving second quantization of a supplied one-particle operator
$h$ has
\[
H_{\rm matter}=\sum_{xy}\psi_x^\dagger h_{xy}\psi_y,
\qquad N=\sum_x n_x,\qquad n_x=\psi_x^\dagger\psi_x,
\qquad [N,H_{\rm matter}]=0.
\]
Indeed $[N,\psi_x^\dagger]=\psi_x^\dagger$ and
$[N,\psi_y]=-\psi_y$ make each bilinear neutral. Identifying this generator
with electric charge is a physical assignment, and every added interaction
must preserve it. The free H.15 completion conserves number without fixing
that assignment.

Choose an orientation for each spatial link $\ell=(x,y)$ and a unitary
compact link $U_{xy}$, with $U_{yx}=U_{xy}^\dagger$. The minimal hopping
interaction is
\[
H_{\rm hop}=-t\sum_{\langle xy\rangle}
\left(\psi_x^\dagger U_{xy}\psi_y+
\psi_y^\dagger U_{xy}^\dagger\psi_x\right),\qquad t\in\mathbb R.
\]
The transformations
\[
\psi_x\mapsto e^{i\alpha_x}\psi_x,
\qquad U_{xy}\mapsto e^{i\alpha_x}U_{xy}e^{-i\alpha_y}
\]
leave each hopping term unchanged. Opposite charges use the conjugate link
representation. An electron--positron reading further requires the
particle/antiparticle and species identifications that H.15 leaves open.

Nearest-neighbor hopping is the simplest gauge construction. H.15's positive
square root is generally nonlocal; to retain it, supply a gauge-covariant
self-adjoint cochain operator $D_U$ and set $\Omega_U=\sqrt{D_U^2+m^2I}$.
Unitary covariance of $D_U$ passes to $\Omega_U$ by functional calculus, so
the positive orbital branch survives, conditional on the charged cochain
placement, the link assignment, and the source-dependent geometric coupling.
With nonzero curvature the charged coboundary need not square to zero.

\subsection*{S.2 Electric flux, Gauss's law, and finite capacity}

For the exact compact rotor take one link space $\ell^2(\mathbb Z)$,
\[
E_\ell|n\rangle=n|n\rangle,\qquad
U_\ell|n\rangle=|n+1\rangle,\qquad
[E_\ell,U_{\ell'}]=\delta_{\ell\ell'}U_\ell.
\]
In units $\hbar=c=1$, a regular spatial regulator of spacing $a>0$ has the
compact lattice Hamiltonian \citep{KogutSusskind1975EM}
\[
\boxed{H_{\rm EM}=\frac{g^2}{2a}\sum_\ell E_\ell^2
 +\frac{1}{g^2a}\sum_p\left[1-\frac{W_p+W_p^\dagger}{2}\right]},
\qquad W_p=\prod_{\ell\in\partial p}U_\ell^{\sigma_{p\ell}},
\]
where $g>0$, $\sigma_{p\ell}=\pm1$ records orientation, and
$W_p=e^{i\Phi_p}$ in the phase representation. An irregular native complex
replaces the uniform coefficients by geometric weights. On a finite network
with finitely many fermionic modes the hopping and plaquette terms are
bounded and the electric term is nonnegative, so the extension is a
lower-bounded Hamiltonian on each finite regulator.

Let $b_{x\ell}$ be $+1$ on outgoing and $-1$ on incoming links. With signed
matter charge $q_x$ and an optional fixed background charge
$\rho_x^{\rm bg}$, define
\[
\boxed{G_x=\sum_\ell b_{x\ell}E_\ell-q_x+\rho_x^{\rm bg}},
\qquad G_x|\mathrm{phys}\rangle=0.
\]
For a unit-charge species $q_x=n_x$. The generator signs agree with the
displayed transformations: $[G_x,\psi_y]=\delta_{xy}\psi_y$ and
$[G_x,U_{yz}]=(\delta_{xy}-\delta_{xz})U_{yz}$. The three weights of
$\psi_y^\dagger U_{yz}\psi_z$ cancel. Along a closed plaquette every vertex
has one entry and one exit, so $[G_x,W_p]=0$. Electric energies and
charge-preserving on-site terms also commute. Consequently
\[
[G_x,H]=0
\]
for the combined Hamiltonian, and its evolution preserves the Gauss sector.
Any further interaction must be checked against $G_x$ separately.

On a closed periodic network $\sum_x b_{x\ell}=0$, hence
\[
\sum_xG_x=-\sum_xq_x+\sum_x\rho_x^{\rm bg}.
\]
The physical sector has zero total charge relative to the background. A
finite example with net positive charge needs compensating background
charge, opposite charges, or outward boundary flux.

The exact rotor has an infinite-dimensional link space and therefore lies
outside any finite matrix algebra of the closure register. For invertible
$U$ in finite dimension $d>0$, $[E,U]=U$ would imply
\[
U^{-1}EU-E=I,\qquad
0=\operatorname{Tr}(U^{-1}EU-E)=\operatorname{Tr}I=d,
\]
a contradiction. A finite quantum-link realization keeps $[E,U]=U$ and exact
Gauss symmetry with a nonunitary raising operator
\citep{ChandrasekharanWiese1997EM}. An electric cutoff is another finite
control; a cyclic cutoff changes the commutator at its boundary. Either route
recovers rotor physics and a Coulomb phase only through a separate limit or
phase argument. The exact compact extension keeps the closure-state count as
a tensor factor, and a strictly finite capacity per cell or link requires one
of these finite realizations.

\subsection*{S.3 Two transverse modes of the quadratic theory}

In the small-flux sector, $1-\cos\Phi_p=\Phi_p^2/2+O(\Phi_p^4)$.
For a three-dimensional periodic control lattice write
\[
s_i=e^{ik_i a}-1,\qquad \lambda(\mathbf k)=\sum_i|s_i|^2,
\qquad M_{ij}(\mathbf k)=\lambda\delta_{ij}-s_i s_j^*.
\]
The magnetic energy is the squared lattice curl. Direct multiplication gives
\[
M\mathbf s=0,\qquad M\mathbf v=\lambda\mathbf v
\quad\hbox{if }\mathbf s^\dagger\mathbf v=0,
\qquad M^2=\lambda M.
\]
For $\mathbf s\ne0$ the orthogonal complement has complex dimension two, so
the eigenvalues are
\[
\boxed{0,\quad 4\sum_i\sin^2(k_i a/2),\quad
4\sum_i\sin^2(k_i a/2)}.
\]
In the neutral free sector Gauss's law removes the longitudinal electric mode
and gauge equivalence removes its coordinate. With the electric and magnetic
coefficients of S.2, Hamilton's equations give
\[
\boxed{\omega(\mathbf k)=\frac{\sqrt{\lambda(\mathbf k)}}a
=\frac2a\sqrt{\sum_i\sin^2(k_i a/2)}
\longrightarrow |\mathbf k|\quad (|\mathbf k|a\to0)}.
\]
At zero lattice momentum the quadratic magnetic symbol vanishes and the
torus carries global holonomy and electric-flux sectors; the two-polarization
count applies at nonzero momentum. Reality pairs the $\mathbf k$ and
$-\mathbf k$ amplitudes.

For the control $\mathbf k a=(0.4,0.8,1.6)$ the computed eigenvalues are
\[
(-4.44\times10^{-16},\quad 2.82286363790248,\quad 2.82286363790248),
\]
with $\|M\mathbf s\|_\infty<5.0\times10^{-16}$; the first entry is numerical
zero. These checks fix the quadratic symbol on the chosen regulator. A photon
of the interacting compact model further requires its deconfined Coulomb
phase, and on the native fluctuating geometry an isotropic coarse symbol, a
speed common to matter and gravity, and control of Lorentz-violating terms.

\subsection*{S.4 A conditional route from coherent alignment to flux energy}

The retained nine-state fiber in H.11 has normalized representatives
$z_x\in\mathbb C^9$ with projective state space $\mathbb{CP}^8$.
Suppose a specified coherent gluing energy favors fidelity,
$-\kappa|z_x^\dagger z_y|^2$, with $\kappa>0$ in the dimensionless Euclidean
weight. For $A=z_x^\dagger z_y$ the exact complex Gaussian identity is
\[
e^{\kappa|A|^2}=\frac1{\pi\kappa}\int_{\mathbb C}d^2u\,
 \exp\!\left[-\frac{|u|^2}{\kappa}+uA+u^*A^*\right],
\]
which follows by completing the square
$-|u-\kappa A^*|^2/\kappa+\kappa|A|^2$.
Independent changes $z_x\mapsto e^{i\alpha_x}z_x$ require
$u_{xy}\mapsto e^{i(\alpha_x-\alpha_y)}u_{xy}$, the transformation of a
charged link. The change of variables is exact and adds no propagating
degree of freedom; treating $u$ as a compact phase requires control of its
amplitude, zeros, and collective dynamics. The construction uses the
theory's own retained projective fiber and needs no separate qubit--qutrit
assignment.

\paragraph{The native retained bridge.}
The bridge of H.11 and P.6 gives a conditional source for the fidelity
interaction. Use its common intertwiner to identify the two retained fibers
and write $H_{xy}=\eta A_R\otimes A_L$, with
$A_R=P_Q(I\otimes2J_z)P_Q|_Q$ and
$A_L=P_Q(2J_z\otimes I)P_Q|_Q$. Besides
$\operatorname{Tr}A_R^2=\operatorname{Tr}A_L^2=38$, the primitive
spin-$3/2$ compression gives
\[
\operatorname{Tr}A_R=\operatorname{Tr}A_L=0,
\qquad \operatorname{Tr}(A_RA_L)=-14.
\]
Let $F$ exchange the two nine-state factors. A common Haar average over
$U\otimes U$, $U\in SU(9)$, has the form $aI+bF$, and preserving its trace
and its trace against $F$ gives $81a+9b=0$ and $9a+81b=-14$. Consequently
\[
\boxed{\mathcal T_{9}(A_R\otimes A_L)
=\frac7{360}I-\frac7{40}F,}
\qquad
\langle z_x,z_y|\mathcal T_{9}(H_{xy})|z_x,z_y\rangle
=\frac{7\eta}{360}-\frac{7\eta}{40}|z_x^\dagger z_y|^2.
\]
A positive bridge scale therefore gives the favorable fidelity sign with
coefficient $7\eta/40$ whenever this Hamiltonian average is physically
realized. The Euclidean alignment parameter is
$\kappa=7\eta\Delta\tau/(40\hbar)$ for a step $\Delta\tau$; it is an
alignment coefficient, distinct from the electromagnetic charge.

The averaging assumption carries weight. The $\mathfrak{su}(9)$ control
algebra of H.11 specifies no Haar path measure, the native $J=0,1,2$ gaps are
unequal, and the common average presupposes a relative amplitude
identification between neighbors that the independent block phases of H.9
leave free. Independent endpoint averages erase the traceless bridge.
Averaging only simultaneous spatial rotations gives
$\mathcal T_{\rm rot}(H_{xy})=(\eta/3)\sum_a A_R^a\otimes A_L^a$, where
$A_{R,L}^a$ compress $2J_a$, and for the dimensionless bridge
\[
\left\|\mathcal T_{\rm rot}(A_R\otimes A_L)
-\mathcal T_9(A_R\otimes A_L)\right\|_{\rm HS}^2
=\frac{28733}{60}>0.
\]
Averaging a finite transfer and exponentiating the averaged Hamiltonian also
differ, by the further terms of the logarithm.

The isolated favorable average selects an aligned ferromagnet. Writing
$J_{\rm al}=7\eta/40>0$, its exchange Hamiltonian is
$-J_{\rm al}\sum_{\langle xy\rangle}F_{xy}$ up to a constant. Every globally
aligned product saturates the bond lower bounds, since $I-F_{xy}\ge0$
annihilates it. Above such a ground state, one change to any of the other
eight internal colors has the exact Hamiltonian $J_{\rm al}$ times the graph
Laplacian; on a cubic control,
\[
\hbar\omega_\alpha(\mathbf k)
=4J_{\rm al}\sum_i\sin^2(k_i a/2),\qquad \alpha=1,\ldots,8.
\]
These are eight quadratic exchange branches, the Goldstone modes of the
ordered phase of the native gluing. S.9 examines its phase evidence; a native
photon has not been demonstrated.

\paragraph{Induced curvature from a gapped boson.}
A specified gapped complex-boson control gives curvature energy directly.
Let its positive quadratic kernel be $K_U=mI-tA_U$, where $A_U$ is the
Hermitian adjacency matrix with unitary link weights, $m>0$, and
$\|rA_U\|<1$ for $r=t/m>0$. Gaussian integration gives
$S_{\rm ind}=\log\det K_U$ up to a field-independent constant, so
\[
S_{\rm ind}=-\sum_{n\ge1}\frac{r^n}{n}\operatorname{Tr}A_U^n
 +\text{constant}.
\]
Each trace sums closed walks. For a simple elementary loop of length $L$,
the $L$ starting vertices in each orientation give $2L\cos\Phi_L$, and
relative to zero flux the first winding contribution is
\[
\boxed{\Delta S_L=2r^L(1-\cos\Phi_L)+\cdots},
\]
positive and of Maxwell curvature form. Multiplicities, backtracking
decorations, other loops, and geometric weights enter the full network
coefficient. On a fixed spatial network this gives magnetic stiffness;
electric dynamics further require temporal transport and a positive transfer
construction.

The isolated-cycle controls use $r=0.14$ and $L=4,6$. Their zero-flux
curvatures $S_L''(0)$ follow from
\[
S_L(\Phi)=\sum_{j=0}^{L-1}
\log\!\left[1-2r\cos\frac{2\pi j+\Phi}{L}\right]:
\]
\begin{center}
\begin{tabular}{lrr}
\hline
Loop & Exact $S_L''(0)$ & Leading $2r^L$\\
\hline
Four edges & $0.000833680556$ & $0.000768320000$\\
Six edges & $0.0000170000085$ & $0.0000150590720$\\
\hline
\end{tabular}
\end{center}
Their exact ratio is $49.0400$, with the control ratio $r$ specified. At flux
$\Phi=0.7$, independent determinant differences agree with the closed-walk
series through order 32 to better than $2\times10^{-16}$.

The four-link mixed loops and six-link spatial loops match the minimal
product diamond regulator of I.2. Their bare ratio enters the continuum
electric-to-magnetic ratio together with areas, spacings, multiplicities,
and the temporal transfer, and the common light cone is matched after these
factors are included (S.7). The determinant is bosonic. The theory's intended
charged elementary defects are fermionic, with their own Grassmann
determinant, spin transport, occupied-state prescription, and regulator;
S.7 computes one filled fermionic control. This unconstrained Gaussian field is also distinct from the constrained
$\mathbb{CP}^8$ model above.

\paragraph{The empty fermion vacuum.}
The normal-ordered, number-conserving fermionic completion of S.1 has an
inert empty vacuum. For positive $h(U)$ and
$H(U)=\sum_{xy}\psi_x^\dagger h_{xy}(U)\psi_y$,
\[
H(U)|0\rangle=0,
\qquad Z_\beta(U)=\det(1+e^{-\beta h(U)})\longrightarrow1
\]
on a fixed finite regulator with a positive gap, so its vacuum energy and
flux stiffness vanish. A filled sector, a pair-producing relativistic
completion, or a specified vacuum-energy contribution responds to flux once
its state and regulator are derived; S.7 computes one filled case.

\subsection*{S.5 What the native projective data leave open}

The Berry connection of the retained projective state is composite. In the
local section $z=(1,w)/\sqrt{1+w^\dagger w}$,
\[
\mathcal A_\mu=-iz^\dagger\partial_\mu z
=\frac{\operatorname{Im}(w^\dagger\partial_\mu w)}{1+w^\dagger w}.
\]
Scaling $w\mapsto\varepsilon w$ makes the numerator quadratic in
$\varepsilon$, and its first variation at $w=0$ vanishes. A smooth nearly
constant projective state therefore carries no independent linear photon in
this connection; a collective deconfined gauge mode needs the additional
dynamics of S.4.

For any phases $a_{xy}$,
\[
U'_{xy}=e^{ia_{xy}}U_{xy}
\quad\Longrightarrow\quad
U'_{xy}XU_{xy}'{}^\dagger=U_{xy}XU_{xy}^\dagger,
\]
so every projective route channel is unchanged, while an amplitude lift
around a closed path obeys
\[
W'_C=e^{i\sum_C a_{xy}}W_C.
\]
Vertex-gradient phases cancel and general link phases change the loop
holonomy; already $+1$ and $-1$ give the same conjugation channel and
different amplitude phases. A phase-sensitive loop readout therefore needs a
physical lift to charged amplitudes, or an equivalent coherent charged link
operator, beyond the channels.

With that lift, every positive coefficient $\beta$ in
$\beta\sum_p(1-\cos\Phi_p)$ respects the same gauge symmetry, and the closure
count and projective channel leave $\beta$ free. The physical electric charge
and fine-structure constant further require normalization of the charge
generator, scale-dependent matching, and renormalization.

\subsection*{S.6 Charge representations and independent central transport}

The phase freedom survives fixing the color transfer coefficient. Consider a
positive product heat evolution on $PU(3)\times U(1)$, with representation
multipliers
\[
t_{R,q}(s,u)=\exp[-sC_R-uq^2],\qquad s,u>0.
\]
Every neutral projective observable has $q=0$, so its transfer is independent
of $u$. Fixing $t_{\mathbf8,0}=13/14$ fixes $3s=\ln(14/13)$ and leaves every
positive central heat time available; charged multipliers with $q\ne0$
distinguish them. The product is an explicit extension family. The native
global quotient and allowed charge representations are further physical
data. At the Lie-algebra level the same freedom is the pair of independent
invariant quadratic forms on $\mathfrak{su}(3)$ and its central $\mathfrak
u(1)$ extension, so positivity, a transfer semigroup, and the full neutral
color spectrum leave the central normalization free.

The protected two-dimensional subspace of the route qutrit cannot carry an
independent commuting weak doublet. If $Q$ commutes with all off-diagonal
matrix units $E_{ij}$ of the full color algebra, equating entries of
$E_{ij}Q=QE_{ij}$ makes the off-diagonal entries vanish and the diagonal
entries agree, so $Q=cI_3$. A nontrivial commuting $\mathfrak{su}(2)$ thus
needs additional multiplicity spaces or a larger unified representation. The
commutant of the full retained $M_9(\mathbb C)$ algebra is likewise scalar;
this restricts charges that commute with every retained control, whether or
not every control appears in the physical Hamiltonian.

The native spinor transport of H.9a has links
$U_j=-i\hat n_j\cdot\boldsymbol\sigma$ of determinant one. They supply
spin-frame transport, and a spin lift adds a discrete sign. A charged
background couples by
\[
U_{xy}\longmapsto e^{iq\int_x^y A\cdot dl}U_{xy},
\]
a well-defined probe because the quadratic spinor theory conserves number.
Promoting the probe to a dynamical field, identifying its generator with
electric charge in the interacting theory, and fixing its microscopic measure
are separate steps.

\subsection*{S.7 Native fermion response and continuum normalization}

A phase-sensitive response is calculable once the charged lift and state are
specified. With the background held fixed, define
\[
Z[A]=\int d\mu(\xi)\,e^{-S(\xi,A)},\qquad
\Gamma[A]=-\ln Z[A],
\]
where $\xi$ collects matter and any designated fluctuating geometry or fast
modes, the background is held fixed, and any background dependence of the
regulator or measure belongs in $S$. For a differentiable finite
positive-weight model, write $J=\partial_\Phi S$ and $K=\partial_\Phi^2S$.
Direct differentiation gives
\[
\boxed{\Gamma''(0)=\langle K\rangle
 -\bigl(\langle J^2\rangle-\langle J\rangle^2\bigr).}
\]
Both the contact term and the connected current response enter. For a
quantum thermal state the second term is the imaginary-time ordered current
correlator. A phase-blind action with phase-independent measure has constant
$Z$ and zero central response; this holds for the projective pure-gauge
construction, and an independently specified charged matter lift responds.

\paragraph{A finite response of the native diamond spinor.}
The H.9a Hamiltonian supplies a fermionic control. In cubic units $a=t=1$,
use its four bond transports and handed Wilson term, with a positive mass
regulator:
\[
h(\mathbf k)=
\begin{pmatrix}M(\mathbf k)I_2&Q(\mathbf k)\\
Q(\mathbf k)^\dagger&-M(\mathbf k)I_2\end{pmatrix},\qquad
M(\mathbf k)=m+r_W\sum_{j<l}
 [1-\cos\mathbf k\cdot(\mathbf d_j-\mathbf d_l)].
\]
Fill the negative-energy bands and apply the unit-charge background
$\mathbf A(\mathbf x)=\Phi\mathbf e\cos(\mathbf q\cdot\mathbf x)$,
where $\mathbf e\perp\mathbf q$ and $\mathbf q$ is one smallest reciprocal
vector of an $N^3$ periodic primitive-cell regulator. Each nearest-neighbor
hop receives the exact bond integral and each Wilson hop the integral along
its two-bond path, which preserves local gauge covariance. For occupied
states $a$, empty states $b$, and a gap at zero energy, finite perturbation
theory gives
\[
\boxed{E_0''(0)=\sum_{a\,{\rm occ}}\langle a|h''|a\rangle
 +2\sum_{\substack{a\,{\rm occ}\\b\,{\rm empty}}}
 \frac{|\langle a|h'|b\rangle|^2}{\epsilon_a-\epsilon_b}.}
\]
The dimensionless background-energy curvatures are:
\begin{center}
\begin{tabular}{rrrrrr}
\hline
$N$ & $m/t$ & $r_W/t$ & Contact & Current response & $E_0''(0)$\\
\hline
3 & 0.3 & 0.2 & 2.319569 & $-2.058025$ & 0.261544\\
4 & 0.3 & 0.2 & 5.622040 & $-5.268218$ & 0.353822\\
4 & 0.6 & 0.2 & 5.212003 & $-4.937010$ & 0.274993\\
4 & 0.3 & 0.1 & 6.237258 & $-5.976798$ & 0.260460\\
\hline
\end{tabular}
\end{center}
The real-space spectra agree with the Bloch operator to $2.2\times10^{-14}$;
finite differences agree with the curvatures to $4\times10^{-7}$; pure-gauge
curvatures cancel to $1.5\times10^{-14}$; local gauge covariance holds to
$5.2\times10^{-16}$; and doubling the probe charge multiplies the curvature
by four. The filled native-spinor control thus has a positive magnetic
response. Its filling, mass, Wilson coefficient, charged lift, and fixed
geometry are inputs, and the physical vacuum and electron mass ratio remain
to be selected.

A time-dependent uniform background has a positive quadratic frequency
response. With $\Delta_{ba}=\epsilon_b-\epsilon_a>0$, its coefficient after
subtracting the static response is
\[
K(\omega)-K(0)
=\omega^2\sum_{a,b}
 \frac{2|\langle a|h'|b\rangle|^2}
 {\Delta_{ba}(\Delta_{ba}^2+\omega^2)},\qquad
Z_E^{(N)}=\frac1V\sum_{a,b}
 \frac{2|\langle a|h'|b\rangle|^2}{\Delta_{ba}^3},
\]
in continuous Hamiltonian time. At $m/t=0.3$, $r_W/t=0.2$, $Z_E^{(N)}$ is
$0.94478$ for $N=3$ and $0.42975$ for $N=4$, while the magnetic
finite-momentum estimator $2E_0''/(V|\mathbf q|^2)$ is $0.005889$ and
$0.005975$. The strong finite-volume dependence of the electric response
leaves the common cone and the continuum coupling to an extrapolation on
larger regulators once the physical charged theory and state are selected.

\paragraph{Geometric matching.}
For local electric and magnetic plaquette coefficients $\kappa_E$ and
$\kappa_M$, the I.2 area tensors and volume per site give, with
$r_t=a_t/d$,
\[
S^{(2)}=\frac12\int d^4x\,(Z_E\mathbf E^2+Z_B\mathbf B^2),\qquad
Z_E=\frac{\sqrt3}{4}\kappa_E r_t,\qquad
Z_B=\frac{4\sqrt3}{3r_t}\kappa_M.
\]
After continuation to real time the transverse speed is
$c_\gamma^2=Z_B/Z_E$. In the time coordinate of I.2's equal-transfer
completion, $r_t^2=32/9$, equality with its invariant speed requires
\[
\boxed{\kappa_E/\kappa_M=3/2.}
\]
The isolated-cycle boson coefficients cannot meet this condition through
their common hopping ratio. Setting $x=(t/m)^2\in(0,1/4)$, their exact
determinants give
\[
\kappa_4=\frac{2x^2}{1-4x},\qquad
\kappa_6=\frac{2x^3}{(1-x)^2(1-4x)},\qquad
\frac{\kappa_4}{\kappa_6}=\frac{(1-x)^2}{x}>\frac94,
\]
from $\det(I-rA_\Phi)=2r^L[\cosh(L\gamma)-\cos\Phi]$ with
$\cosh\gamma=1/(2r)$ for an isolated cycle; the strict bound is
$4(1-x)^2-9x=(4-x)(1-4x)>0$. At $t/m=0.14$, assigning these cycle
coefficients directly to the native plaquettes gives
$c_\gamma^2=0.0305873$, which excludes that assignment with the fixed
clock. A full graph determinant, different matter, temporal dynamics, or
derived geometric weights change the coefficients and remain open.

A global flat holonomy can have zero curvature and a nonzero finite-volume
matter response, so the continuum target is the low-momentum transverse
response to a field with curvature, with contact terms, geometry, and
temporal normalization included. In a Lorentz-invariant limit, for a
connection in which the chosen elementary defect has unit
covariant-derivative charge,
\[
\Gamma[A]\supset\frac{Z_F(\mu)}4\int d^4x\,F_{\mu\nu}F^{\mu\nu},
\qquad \boxed{\alpha(\mu)=\frac1{4\pi Z_F(\mu)}}.
\]
A bare local $F^2$ term adds to the induced response, and gauge symmetry and
the neutral channel leave it free; its value is a microscopic boundary
condition to be derived. A physical prediction needs the charge
normalization, charged dynamics and state, field measure or collective
phase, and scale matching within one construction.

\subsection*{S.8 Coupling trials and their scale test}

The color heat-kernel coordinate of I.2 is
\[
g_{\rm graph,HK}^2=\frac{4\sqrt6}{3}\ln\frac{14}{13}
=0.242035623532,
\qquad \alpha_{3,\rm graph}^{-1}=\frac{4\pi}{g_{\rm graph,HK}^2}
=51.9195085.
\]
An independent unified five-dimensional representation has
\[
Y=\operatorname{diag}(-1/3,-1/3,-1/3,1/2,1/2),
\qquad \operatorname{Tr}Y^2=5/6,
\]
with $Q=T_3+Y$ (conventions using $Y/2$ rename this generator).
Normalizing its trace to $1/2$ gives $T_Y=\sqrt{3/5}\,Y$. If both gauge
kinetic terms descend from a common coupling $g_U$,
\[
g_2^2=g_U^2,\quad g_Y^2=\frac35g_U^2,\quad
\frac1{e^2}=\frac1{g_2^2}+\frac1{g_Y^2},\quad e^2=\frac38g_U^2,
\]
the conventional unified normalization \citep{PDGGUT2025EM}. Identifying
$g_U$ with the graph coordinate gives
\[
\alpha_{\rm trial}=\frac{\sqrt6}{8\pi}\ln\frac{14}{13},\qquad
\alpha_{\rm trial}^{-1}=138.4520227,
\]
$1.0333\%$ above the CODATA 2022 Thomson-limit inverse coupling
$137.035999177(21)$ \citep{CODATA2022Astra}. The five-dimensional
representation, the common kinetic normalization, and the graph-to-continuum
identification are inputs beyond the qutrit and its protected subspace.

The trial fails a scale test. As an explicit control, set
$\mu_*=1.22\times10^{19}\,\mathrm{GeV}$, $M_Z=91.1876\,\mathrm{GeV}$, and
impose equal continuum couplings $\alpha_i^{-1}(\mu_*)=u=51.9195085$; the
scale and identity matching are choices of the control. With Standard Model
content, one-loop coefficients $b_i=(41/10,-19/6,-7)$ in the convention
$d\alpha_i^{-1}/d\ln\mu=-b_i/(2\pi)$ give \citep{PDGGUT2025EM}
\[
\alpha_i^{-1}(M_Z)=u+\frac{b_i}{2\pi}\ln\frac{\mu_*}{M_Z}
=(77.65227,\ 32.04461,\ 7.98553),
\]
and
\[
\boxed{\alpha_{\rm EM}^{-1}(M_Z)
=\frac53\alpha_1^{-1}(M_Z)+\alpha_2^{-1}(M_Z)
=\frac83u+\frac{11}{6\pi}\ln\frac{\mu_*}{M_Z}
=161.46506,}
\]
with $\sin^2\theta_W(M_Z)=0.19846$. The empirical electroweak inverse
coupling in the $\overline{\rm MS}$ convention is about $128$
\citep{PDGEW2024EM}. For every $\mu_*\ge M_Z$ the one-loop correction is
nonnegative, so no choice of unification scale lowers $138.452$ to the
electroweak value; additional thresholds or a different charged theory would
be new inputs. The proximity to $137$ therefore carries no predictive weight.

The color comparison is closer to the calculated sector and also depends on
running. From the PDG 2025 reference $\alpha_s(M_Z)=0.1180$
\citep{PDGQCD2025EM}, the controls at the same $\mu_*$ are
\begin{center}
\begin{tabular}{lr}
\hline
Prescription & $\alpha_3^{-1}(\mu_*)$\\
\hline
Graph heat-kernel coordinate & 51.91951\\
One loop, six flavors throughout & 52.40856\\
One loop, five to six flavors at $173\,\mathrm{GeV}$ & 52.47650\\
Two-loop QCD, same continuous threshold & 53.01898\\
\hline
\end{tabular}
\end{center}
The last row integrates
$d\alpha_s/d\ln\mu=-\beta_0\alpha_s^2/(2\pi)
-\beta_1\alpha_s^3/(8\pi^2)$, with
$\beta_0=11-2n_f/3$ and $\beta_1=102-38n_f/3$
\citep{PDGQCD2025EM}, as a sensitivity control without finite decoupling,
electroweak terms, further thresholds, or graph-to-continuum matching. The
fractional discrepancy in inverse coupling grows from $0.93\%$ in the
simplest six-flavor comparison to $2.07\%$ in the last row. These
comparisons motivate a scheme-matched color calculation and leave the
electromagnetic coupling undetermined.

\subsection*{S.9 Phases of the native gluing and its charged spectrum}

By the identity of S.4, the native fidelity gluing is a lattice $\mathbb{CP}^8$
model with quartic coupling and no bare plaquette term,
\[
Z=\int\prod_x dz_x\,\exp\Bigl(\kappa\sum_{\langle xy\rangle}|z_x^\dagger z_y|^2\Bigr),
\qquad z_x\in\mathbb{CP}^8,
\]
whose local redundancy $z_x\mapsto e^{i\alpha_x}z_x$ is carried by the
auxiliary link $u_{xy}$. Its phases decide whether that link propagates.

In the ordered branch the composite $P_x=z_xz_x^\dagger$ orders, the link
is Higgsed, and the $2(N-1)=16$ real Goldstone fields are the eight quadratic
exchange branches of S.4. Uniform averages on $\mathbb{CP}^{N-1}$ give
$\delta\langle P\rangle=H/[N(N+1)]+O(H^2)$ for a small traceless Hermitian
field $H$. With each undirected bond counted once, the mean field is
$H=\kappa z_c Q$, where $Q=\langle P\rangle-I/N$. The disordered mean-field
spinodal is therefore $\kappa_{\rm sp}=N(N+1)/z_c=45/4=11.25$ for $N=9$ on a
four-dimensional Euclidean cubic lattice ($z_c=8$). The order parameter
admits the cubic invariant $\operatorname{Tr}Q^3$ for $N\ge3$, permitting a
first-order mean-field transition. Its coexistence coupling requires a
nonlinear free-energy comparison. Three-dimensional lattice studies of
$\mathbb{CP}^{N-1}$ models without a gauge kinetic term find first-order
transitions for every $N>2$ up to $N=100$, with gauge correlations of one
lattice spacing on the disordered side
\citep{PelissettoVicariCPN2019,PelissettoVicariLargeN2020}.

A Metropolis scan of the four-dimensional model with $N=9$ on a periodic
$6^4$ lattice, with 300 thermalization and 300 measurement sweeps at each
coupling in steps of $0.5$, gives a hysteresis loop in the order parameter
$\langle\operatorname{Tr}\bar Q^2\rangle N/(N-1)$, where
$\bar Q=V^{-1}\sum_x(P_x-I/N)$. From a disordered start it stays
below $0.005$ up to $\kappa=5.5$ and rises to $0.21$ and $0.53$ at
$\kappa=6.0$ and $6.5$, while the link energy
$\langle|z_x^\dagger z_y|^2\rangle$ jumps from $0.21$ to $0.44$. From an ordered
start it remains $0.34$ at $\kappa=5.0$ and falls to $0.11$ and $0.001$ at
$\kappa=4.5$ and $4.0$. The loop between $\kappa\simeq4.5$ and $6$ is
consistent with a first-order transition. Establishing its order and
coexistence coupling requires finite-size and equilibration tests in the
four-dimensional model.

A one-loop estimate for gapped charged fiber quanta gives a Maxwell term.
Each charged complex scalar contributes $\Delta(1/e^2)=\ln(\Lambda/m)/(24\pi^2)$
at one loop, so the nine fiber components give
\[
\frac1{e^2_{\rm ind}}=\frac{9}{24\pi^2}\ln\frac{\Lambda}{m}
=0.0380\,\ln\frac{\Lambda}{m}.
\]
Four-dimensional compact $U(1)$ with the Wilson action has its Coulomb phase
for $1/e^2>\beta_c=1.0111331(21)$ \citep{Panero2005U1}. The displayed one-loop
coefficient reaches this reference value at $\ln(\Lambda/m)\simeq26.6$, a
fiber mass below $3\times10^{-12}$ of the cutoff. The phase boundary of the native induced action also depends on its
additional loop terms and interactions.

The native charged states are the three charged leptons, with the cell scale
$\Lambda=m_P=1.22089\times10^{19}\,\mathrm{GeV}$ as cutoff. With a fully induced
coupling, $\alpha^{-1}=(2/3\pi)\sum_fN_cq_f^2\ln(\Lambda/m_f)$ at one loop. The
leptons give
\[
\sum_\ell\ln\frac{m_P}{m_\ell}=51.53+46.20+43.37=141.10,\qquad
\frac1{e^2}=2.38,\qquad \alpha^{-1}=29.9 .
\]
The induced stiffness exceeds the pure-Wilson reference threshold, and the coupling is
four and a half times the observed one. The observed $137.036$ needs a bare
stiffness $1/e^2=8.52$ at the cell scale in addition. As a control, the full
Standard Model charged fermions give $\sum_fN_cq_f^2\ln(\Lambda/m_f)=364$ and
$\alpha^{-1}=77.2$ ($75.4$ with $0.33$ and $0.5\,\mathrm{GeV}$ light-quark
thresholds; charged gauge bosons and hadronic running are omitted), which
still needs a bare $1/e^2=4.76$.

\subsection*{S.10 Electroweak stiffness: location, form and targets}

The photon is the collective mode $A=\sin\theta_W\,W^3+\cos\theta_W\,B$ of the
electroweak links, with $1/e^2=1/g_2^2+1/g_Y^2$, so both stiffnesses must be
present at the cell scale. Loops of the known charged states supply neither.
With three generations and one Higgs doublet the $SU(2)_L$ coefficient is
$b_2=-22/3+4n_g/3+n_H/6=-19/6$ in the convention of S.8, so $\alpha_2^{-1}$
grows toward the ultraviolet and stays positive at every scale above $M_Z$;
an induced coupling requires $\alpha_2^{-1}=0$ at the cutoff and is excluded.
Two-loop gauge running with the top Yukawa coupling, from the
$\overline{\rm MS}$ inputs $\alpha^{-1}(M_Z)=127.95$,
$\sin^2\theta_W=0.23122$ and $\alpha_s=0.1180$
\citep{PDGEW2024EM,PDGQCD2025EM} without threshold corrections, gives at
$m_P$
\[
\alpha_Y^{-1}=55.08,\qquad \alpha_2^{-1}=49.15,\qquad \alpha_3^{-1}=52.90,
\]
with $\alpha_Y=g_Y^2/4\pi$ in the normalization $Q=T_3+Y$. Hypercharge is
therefore far from a Landau pole at the cell scale. An induced $U(1)_Y$ would
need $55$ units of $(2/3\pi)\sum N_cY^2\ln(\Lambda/m)$ from hypercharged
states beyond the Standard Model, which the native charged spectrum of S.9
does not contain. The photon is a composite of the substrate links, and its
cell-scale stiffness must therefore be native to the substrate, as the color
stiffness of I.2 is; the electroweak multipliers below are required values
and are not yet derived.

The theory's one native mechanism that fixes a gauge stiffness is the
diamond heat-kernel transfer of I.2, with $t_R=e^{-sC_R}$ per primitive step
and $g^2=4\sqrt6\,s$ on the diamond regulator. A reversible transfer on a
$d$-dimensional record with direct amplitude $k$,
$K_P=I-(1-k)P$ and $L_P=\sqrt{1-k^2}\,P$ for a rank-one projector $P$, averaged
over $P$ with $\int P=I/d$ and $\int PXP=(X+\operatorname{Tr}X\,I)/[d(d+1)]$,
multiplies every traceless operator by
\[
t_{\rm adj}=1-\frac{2(1-k)}d+\frac{(1-k)^2+(1-k^2)}{d(d+1)}
=1-\frac{2(1-k)}{d+1}.
\]
For the qutrit route record, $d=3$ and $k=6/7$ give $t_8=13/14$, the color
multiplier of I.2. The electroweak factors take the same form on records that
carry weak isospin ($d=2$) and hypercharge, and a $U(1)$ transfer multiplies a
record of charge $q$ by $e^{-sq^2}$. Projective channels cannot fix a central
heat time (S.5, S.6), so the hypercharge transfer acts on the amplitude of a
persistent charged record.

The running couplings fix what these transfers must give:
\begin{center}
\begin{tabular}{lcccc}
\hline
Factor & $\alpha^{-1}(m_P)$ & $s=g^2/(4\sqrt6)$ & required multiplier & native\\
\hline
$SU(3)$ & 52.90 & 0.02425 & $t_8=e^{-3s}=0.92985$ & $13/14=0.92857$\\
$SU(2)_L$ & 49.15 & 0.02610 & $t_3=e^{-2s}=0.94915$ & not derived\\
$U(1)_Y$ & 55.08 & 0.02328 & $t_{Y=1}=e^{-s}=0.97698$ & not derived\\
\hline
\end{tabular}
\end{center}
The native color value $\alpha_3^{-1}=51.92$ lies $1.9\%$ below the
requirement, within the spread of the controls in S.8. A weak-isospin
doublet record would need $k=0.9237$. One common heat time
$s=\ln(14/13)/3$ for all three factors, run down at two loops, gives
$\alpha^{-1}(M_Z)=127.58$, $\sin^2\theta_W=0.2537$ and $\alpha_s=0.1340$,
deviations of $-0.3\%$, $+9.7\%$ and $+13.6\%$, and fails a pre-registered
$3\%$ criterion. The near agreement of $\alpha^{-1}$ comes from errors of
$-3.2\%$ in $\alpha_Y^{-1}$ and $+9.4\%$ in $\alpha_2^{-1}$ that offset in the
sum and carries no weight.

\paragraph{A candidate record.} The theory contains a native qutrit, the route record of I.2, and a native qubit, the spinor lift of H.9a. A maximally entangled link joining them has identification algebra $\mathfrak{su}(3)\oplus\mathfrak{su}(2)\oplus\mathfrak u(1)$, the block-diagonal traceless generators on $\mathbb C^2\oplus\mathbb C^3$. Its centre $Q=i\,\mathrm{diag}(3,3,-2,-2,-2)$ acts on the link as the scalar $i$, so $Y=Q/6i$ assigns hypercharge $1/6$, and the group is $S(U(3)\times U(2))\cong(SU(3)\times SU(2)\times U(1))/\mathbb Z_6$ \citep{BaezHuerta2010}. The central phase acts on the pair amplitude itself, so the link is the persistent charged record that S.5 and S.6 require. A cube root of unity $\omega 1_3$ has determinant one, acts on the link as $\omega$ and leaves every conjugation channel unchanged, so the link also carries the triality that the projective channel of I.2 cannot see. Every positive map between the ends of a link is decomposable exactly when $d_1d_2\le6$ \citep{Stormer1963,Woronowicz1976,Choi1975}, and $(2,3)$ is the only such type with unequal ends; \citet{WissnerGross2026} selects the link by this property. The link fixes the location and algebra of the record. Its stiffness depends on the frame alignment between the link and the route record. Transported isotropically through the color channel, the qubit record has $t_3=13/14$ and $\alpha_2^{-1}(m_P)=34.6$. Pinned to the unit-amplitude doublet of the route Gram matrix, it has $t_3=1$ and no coupling. The required $t_3=0.94915$ lies strictly between these limits, and the alignment that would fix it remains to be derived. The link carries the quark-doublet hypercharge; leptons and right-handed fields need records of their own.

Electromagnetism then follows from native weak-isospin and hypercharge
records, computed as color's transfer is: each record specified in full, its
$k$ computed, then compared with the table. An equivalent premise is a
substrate Gauss constraint tying a link flux to the charged label $n_e$, with
ring-exchange stiffness, as in the emergent photon of quantum spin ice
\citep{HermeleFisherBalents2004}. The capacity labels carry only the vector
closure that Appendix~Q.3 excludes as an electromagnetic reading and leave the
scalar sum $S=\sum_im_i$ free; under the native weights $e^{-\eta_*K^2}$ the
probability of $S=0$ is $0.13$, so the present capacity data supply no
scalar Gauss law. The obstacles that look hardest have known constructions.
The Nielsen--Ninomiya theorem assumes an exactly chirally symmetric local
bilinear action \citep{NielsenNinomiya1981}. Overlap fermions satisfy the
Ginsparg--Wilson relation, keep an exact lattice chiral symmetry and are
exponentially local \citep{Neuberger1998}; anomaly-free chiral gauge
theories have lattice constructions with exact gauge invariance for abelian
groups and to all orders of perturbation theory for nonabelian groups
\citep{Luscher1999,Luscher2000}. The overlap operator is built from a
Wilson-type kernel, which the handedness-weighted term of H.9a, paragraph~2d,
supplies on the diamond graph. Given one generation of chiral content,
cancellation of the gauge and mixed gravitational anomalies fixes the
hypercharges up to normalization. The remaining items are new premises: the
weak-isospin and hypercharge records or the scalar Gauss constraint, and
electroweak symmetry breaking, for which the theory has no native scalar
doublet or condensate. Triality-carrying quark defects remain open with no
known obstruction (I.2). Deriving these premises is the electroweak
matter-sector problem.

\subsection*{S.11 Verification scope and physical selection}

The \texttt{Entropic.Electromagnetism} modules prove the finite gauge and
charge algebra, closed-graph neutrality, the rotor shift identity on
sequences, the transverse symbol identities and the two-dimensional
transverse subspace, projective phase freedom, the finite-matrix rotor
obstruction, Gaussian square completion, the vanishing linear term of the
composite connection, and the leading positive loop coefficient under the
displayed hypotheses. The native-selection module checks the exact
bridge-coordinate compressions and traces, the coefficients given the two
twirl trace constraints, aligned-product swap identities, the exchange
dispersion's zero slope, and annihilation of the empty vacuum by normally
ordered bilinears. The coupling module checks the unified trace
normalization, the one-loop algebra and its scale bound, the matrix
commutants, central-transfer freedom within the stated multiplier family,
finite phase-blind partition identities, the response algebra, and the
diamond matching condition with the isolated-cycle ratio bound. The
stiffness module proves the adjoint transfer eigenvalue
$1-2(1-k)/(d+1)$ from the two Haar averages and its color instance $13/14$,
the cubic invariant, the spinodal substitution $9\cdot10/8=45/4$, the bound
$9\ell/(24\pi^2)<\beta_c$ for $\ell\le26$, the lepton bound
$\alpha^{-1}<33$, the coefficient $b_2=-19/6$, and the positivity of an
asymptotically free inverse coupling above its reference scale. The readable-link module proves that the centre of a $(d_1,d_2)$ link acts as $d_2-d_1$, the link hypercharge $1/6$, the triality action and its invisibility to conjugation channels, the compressed isotropic multiplier, and $13/14<0.94915<1$, with the read cap as a named hypothesis. The modules
compile in the shared Lean~4.35.0 library, whose axiom audit admits only the
standard axioms. The group average and its physical realization, unbounded
rotor dynamics, Gaussian functional integration, determinant-series
convergence, the lattice phase diagram, and an interacting continuum phase
lie outside these statements.

The numerical script checks random local gauges, projective phase
cancellation, a finite two-site matter--flux Hamiltonian, a closed plaquette,
the photon symbol, and the two bosonic loop determinants. Its finite flux
link is nonunitary, preserving the gauge commutator and testing Gauss
conservation within the finite-rotor obstruction. The coupling and response
scripts reproduce the scale tests, the exact isolated-cycle curvatures, and
the filled native-spinor controls of S.7, with Bloch-spectrum, local gauge
covariance, pure-gauge cancellation, charge scaling, and finite-difference
checks. The phase scan of S.9 and the two-loop running of S.10 are archived
with their inputs. Input values, residuals, and source hashes are archived
with the reproduction materials.

The remaining physical target is
\[
\boxed{\text{native weak-isospin and hypercharge records}
\ \longrightarrow\ \text{physical charged }U(1)\text{ links}}.
\]
Closing it requires a conserved native charge and amplitude lift, a
finite-capacity realization or a justified emergent rotor limit, native
records whose transfers give the stiffnesses of S.10, or an equivalent
substrate Gauss constraint, the common causal cone, and physical coupling
matching. The compact gauge action gives the Gauss constraint and two
transverse modes in its quadratic symbol. The one-loop lepton contribution gives
$\alpha^{-1}\simeq30$; running fixes the electroweak stiffnesses required
of the native records.
A native derivation of electromagnetism from the microscopic vertex remains
open.

\begin{thebibliography}{99}

\bibitem[Mohr et al.(2025)]{CODATA2022Astra}
P.~J. Mohr, D.~B. Newell, B.~N. Taylor, and E.~Tiesinga,
``CODATA recommended values of the fundamental physical constants: 2022,''
\emph{Reviews of Modern Physics} \textbf{97}, 025002 (2025).
\url{https://physics.nist.gov/cuu/pdf/wall_2022.pdf}.

\bibitem[Takahashi et al.(2026)]{PDG2026Astra}
F.~Takahashi et al. (Particle Data Group),
\emph{International Journal of Modern Physics A} \textbf{41}, 2630011 (2026).
Tau mass listing, accessed 4 September 2026:
\url{https://pdglive.lbl.gov/Particle.action?node=S035}.


\bibitem[McGaugh et al.(2016)]{McGaugh2016}
S.~S. McGaugh, F. Lelli, and J.~M. Schombert,
``Radial acceleration relation in rotationally supported galaxies,''
\emph{Physical Review Letters} \textbf{117}, 201101 (2016).

\bibitem[Clowe et al.(2006)]{Clowe2006}
D. Clowe et al.,
``A direct empirical proof of the existence of dark matter,''
\emph{The Astrophysical Journal Letters} \textbf{648}, L109--L113 (2006).

\bibitem[Bertotti et al.(2003)]{Bertotti2003}
B. Bertotti, L. Iess, and P. Tortora,
``A test of general relativity using radio links with the Cassini spacecraft,''
\emph{Nature} \textbf{425}, 374--376 (2003).

\bibitem[Tian et al.(2020)]{Tian2020}
Y. Tian et al.,
``The radial acceleration relation in CLASH galaxy clusters,''
arXiv:2001.08340 (2020).

\bibitem[Guttmann(2010)]{Guttmann2010}
A.~J. Guttmann,
``Lattice Green functions in all dimensions,''
\emph{Journal of Physics A: Mathematical and Theoretical} \textbf{43}, 305205 (2010).

\bibitem[Tsukerman and Veomett(2016)]{TsukermanVeomett2016}
E. Tsukerman and E. Veomett,
``A simple proof of Cauchy's surface area formula,''
arXiv:1604.05815 [math.MG] (2016).

\bibitem[Wilson(1974)]{Wilson1974}
K.~G. Wilson,
``Confinement of quarks,''
\emph{Physical Review D} \textbf{10}, 2445--2459 (1974).

\bibitem[Joyce(1994)]{Joyce1994}
G.~S. Joyce,
``On the cubic lattice Green functions,''
\emph{Proceedings of the Royal Society of London A} \textbf{445}, 463--477 (1994).

\bibitem[Arnowitt et al.(1962)]{ADM1962}
R. Arnowitt, S. Deser, and C.~W. Misner,
``The dynamics of general relativity,''
in \emph{Gravitation: An Introduction to Current Research},
ed. L. Witten (Wiley, New York, 1962).

\bibitem[Bekenstein(1973)]{Bekenstein1973}
J.~D. Bekenstein,
``Black holes and entropy,''
\emph{Physical Review D} \textbf{7}, 2333--2346 (1973).

\bibitem[Hawking(1975)]{Hawking1975}
S.~W. Hawking,
``Particle creation by black holes,''
\emph{Communications in Mathematical Physics} \textbf{43}, 199--220 (1975).

\bibitem[Gibbons and Hawking(1977)]{GibbonsHawking1977}
G.~W. Gibbons and S.~W. Hawking,
``Action integrals and partition functions in quantum gravity,''
\emph{Physical Review D} \textbf{15}, 2752--2756 (1977).

\bibitem[Regge and Wheeler(1957)]{ReggeWheeler1957}
T. Regge and J.~A. Wheeler,
``Stability of a Schwarzschild singularity,''
\emph{Physical Review} \textbf{108}, 1063--1069 (1957).

\bibitem[Zerilli(1970)]{Zerilli1970}
F.~J. Zerilli,
``Effective potential for even-parity Regge--Wheeler gravitational perturbation equations,''
\emph{Physical Review Letters} \textbf{24}, 737--738 (1970).

\bibitem[Sanderson et al.(2003)]{Sanderson2003}
A.~J.~R. Sanderson, T.~J. Ponman, A. Finoguenov, E.~J. Lloyd-Davies, and M. Markevitch,
``The Birmingham--CfA cluster scaling project --- I. Gas fraction and the $M$--$T_X$ relation,''
\emph{Monthly Notices of the Royal Astronomical Society} \textbf{340}, 989--1010 (2003).

\bibitem[Vikhlinin et al.(2006)]{Vikhlinin2006}
A. Vikhlinin, A. Kravtsov, W. Forman, C. Jones, M. Markevitch, S.~S. Murray, and L. Van Speybroeck,
``Chandra sample of nearby relaxed galaxy clusters: mass, gas fraction, and mass--temperature relation,''
\emph{The Astrophysical Journal} \textbf{640}, 691--709 (2006).

\bibitem[Eckert et al.(2022)]{Eckert2022}
D. Eckert, S. Ettori, E. Pointecouteau, R.~F.~J. van der Burg, and S.~I. Loubser,
``The gravitational field of X-COP galaxy clusters,''
\emph{Astronomy \& Astrophysics} \textbf{662}, A123 (2022).

\bibitem[Li et al.(2023)]{Li2023}
P. Li, Y. Tian, M.~P. J\'ulio, M.~S. Pawlowski, F. Lelli, S.~S. McGaugh, J.~M. Schombert, J.~I. Read, P.-C. Yu, and C.-M. Ko,
``Measuring galaxy cluster mass profiles into the low-acceleration regime with galaxy kinematics,''
\emph{Astronomy \& Astrophysics} \textbf{677}, A24 (2023).

\bibitem[Milgrom(1983)]{Milgrom1983}
M. Milgrom,
``A modification of the Newtonian dynamics as a possible alternative to the hidden mass hypothesis,''
\emph{Astrophysical Journal} \textbf{270}, 365 (1983).

\bibitem[Milgrom(1999)]{Milgrom1999}
M. Milgrom,
``The modified dynamics as a vacuum effect,''
\emph{Physics Letters A} \textbf{253}, 273 (1999),
doi:10.1016/S0375-9601(99)00077-8.

\bibitem[Milgrom(2010)]{Milgrom2010}
M. Milgrom,
``Quasi-linear formulation of MOND,''
\emph{Monthly Notices of the Royal Astronomical Society} \textbf{403}, 886 (2010),
doi:10.1111/j.1365-2966.2009.16184.x.

\bibitem[Bekenstein(2004)]{Bekenstein2004}
J.~D. Bekenstein,
``Relativistic gravitation theory for the modified Newtonian dynamics paradigm,''
\emph{Physical Review D} \textbf{70}, 083509 (2004).

\bibitem[Famaey and McGaugh(2012)]{FamaeyMcGaugh2012}
B. Famaey and S.~S. McGaugh,
``Modified Newtonian dynamics (MOND): observational phenomenology and relativistic extensions,''
\emph{Living Reviews in Relativity} \textbf{15}, 10 (2012).

\bibitem[Jacobson(1995)]{Jacobson1995}
T. Jacobson,
``Thermodynamics of spacetime: the Einstein equation of state,''
\emph{Physical Review Letters} \textbf{75}, 1260 (1995).

\bibitem[Dorau and Much(2026)]{DorauMuch2026}
P. Dorau and A. Much,
``From quantum relative entropy to the semiclassical Einstein equations,''
\emph{Physical Review Letters} \textbf{136}, 091602 (2026),
doi:10.1103/lmq8-nsty.

\bibitem[Chandrasekaran et al.(2023)]{ChandrasekaranEtAl2023}
V. Chandrasekaran, R. Longo, G. Penington, and E. Witten,
``An algebra of observables for de Sitter space,''
\emph{Journal of High Energy Physics} \textbf{2023}(2), 082 (2023),
doi:10.1007/JHEP02(2023)082.

\bibitem[Bianconi(2025)]{Bianconi2025}
G. Bianconi,
``Gravity from entropy,''
\emph{Physical Review D} \textbf{111}, 066001 (2025),
doi:10.1103/PhysRevD.111.066001.

\bibitem[Bianconi(2026)]{Bianconi2026}
G. Bianconi,
``Thermodynamics of the gravity from entropy theory,''
\emph{Physical Review D} \textbf{114}, 024042 (2026),
doi:10.1103/26kn-thgp.

\bibitem[Caticha(2021)]{CatichaHK2021}
A. Caticha,
``Quantum mechanics as Hamilton--Killing flows on a statistical manifold,''
\emph{Physical Sciences Forum} \textbf{3}, 12 (2021),
doi:10.3390/psf2021003012.

\bibitem[Caticha(2025)]{CatichaQED2025}
A. Caticha,
``Entropic dynamics approach to quantum electrodynamics,''
\emph{Entropy} \textbf{27}, 1247 (2025),
doi:10.3390/e27121247.

\bibitem[Tonomura et al.(1989)]{Tonomura1989}
A. Tonomura, J. Endo, T. Matsuda, T. Kawasaki, and H. Ezawa,
``Demonstration of single-electron buildup of an interference pattern,''
\emph{American Journal of Physics} \textbf{57}, 117--120 (1989),
doi:10.1119/1.16104.

\bibitem[Verlinde(2011)]{Verlinde2011}
E. Verlinde,
``On the origin of gravity and the laws of Newton,''
\emph{Journal of High Energy Physics} \textbf{2011}(4), 29 (2011).

\bibitem[Verlinde(2017)]{Verlinde2017}
E. Verlinde,
``Emergent gravity and the dark universe,''
\emph{SciPost Physics} \textbf{2}, 016 (2017).

\bibitem[Skordis and Z\l{}o\'snik(2021)]{Skordis2021}
C. Skordis and T. Z\l{}o\'snik,
``New relativistic theory for modified Newtonian dynamics,''
\emph{Physical Review Letters} \textbf{127}, 161302 (2021).

\bibitem[Collins et al.(2004)]{Collins2004}
J. Collins, A. Perez, D. Sudarsky, L. Urrutia, and H. Vucetich,
``Lorentz invariance and quantum gravity: an additional fine-tuning problem?,''
\emph{Physical Review Letters} \textbf{93}, 191301 (2004).

\bibitem[Bombelli et al.(1987)]{Bombelli1987}
L. Bombelli, J. Lee, D. Meyer, and R.~D. Sorkin,
``Space-time as a causal set,''
\emph{Physical Review Letters} \textbf{59}, 521 (1987).

\bibitem[Dowker et al.(2004)]{Dowker2004}
F. Dowker, J. Henson, and R.~D. Sorkin,
``Quantum gravity phenomenology, Lorentz invariance and discreteness,''
\emph{Modern Physics Letters A} \textbf{19}, 1829 (2004).

\bibitem[Pittordis et al.(2025)]{Pittordis2025}
C. Pittordis, W. Sutherland, and P. Shepherd,
``Wide binaries from Gaia DR3: testing GR versus MOND with realistic triple modelling,''
\emph{The Open Journal of Astrophysics} \textbf{8} (2025),
doi:10.33232/001c.142887.

\bibitem[Chae(2025)]{Chae2025}
K.-H. Chae,
``Low-acceleration gravitational anomaly from Bayesian 3D modeling of wide binary orbits: methodology and results with Gaia Data Release 3,''
\emph{Astrophysical Journal} \textbf{985}, 210 (2025),
doi:10.3847/1538-4357/adce09.

\bibitem[Chae and Yoon(2026)]{Chae2026}
K.-H. Chae and Y. Yoon,
``Revisiting data quality control and multiple-star modeling in wide binary gravity tests: confirmation of MOND-type gravitational anomaly at low acceleration,''
arXiv:2607.14450 (2026).

\bibitem[Wang et al.(2026)]{Wang2026MW}
Z.-L. Wang, Y.-L.~S. Tsai, L. Zhang, Y. Wu, H. Li, X.-X. Xue, H. Zhao, and Y.-Z. Fan,
``Milky Way Dynamics Favor Dark Matter over Modified Gravity Models,''
arXiv:2605.10857 (2026).

\bibitem[McKee et al.(2015)]{McKee2015}
C.~F. McKee, A. Parravano, and D.~J. Hollenbach,
``Stars, Gas, and Dark Matter in the Solar Neighborhood,''
\emph{The Astrophysical Journal} \textbf{814}, 13 (2015).

\bibitem[Palau and Miralda-Escud\'e(2023)]{Palau2023}
C.~G. Palau and J. Miralda-Escud\'e,
``The oblateness of the Milky Way dark matter halo from the stellar streams of NGC 3201, M68, and Palomar 5,''
\emph{Monthly Notices of the Royal Astronomical Society} \textbf{524}, 2124--2147 (2023).

\bibitem[Chemaly et al.(2026)]{Chemaly2026}
D. Chemaly, E. Sola, S. Koposov, H. Zhang, V. Belokurov, D. Erkal, et al.,
``Constraints on the population level distribution of nearby dark matter halo shapes with extragalactic streams,''
arXiv:2607.05510 (2026).

\bibitem[Wegg et al.(2019)]{Wegg2019}
C. Wegg, O. Gerhard, and M. Bieth,
``The gravitational force field of the Galaxy measured from the kinematics of RR Lyrae in Gaia,''
\emph{Monthly Notices of the Royal Astronomical Society} \textbf{485}, 3296--3316 (2019).

\bibitem[Prada et al.(2019)]{Prada2019}
J. Prada, J.~E. Forero-Romero, R.~J.~J. Grand, R. Pakmor, and V. Springel,
``Dark matter halo shapes in the Auriga simulations,''
\emph{Monthly Notices of the Royal Astronomical Society} \textbf{490}, 4877 (2019).

\bibitem[Chua et al.(2019)]{Chua2019}
K.~T.~E. Chua, A. Pillepich, M. Vogelsberger, and L. Hernquist,
``Shape of dark matter haloes in the Illustris simulation: effects of baryons,''
\emph{Monthly Notices of the Royal Astronomical Society} \textbf{484}, 476 (2019).

\bibitem[Martinsson et al.(2013a)]{Martinsson2013a}
T.~P.~K. Martinsson, M.~A.~W. Verheijen, K.~B. Westfall, M.~A. Bershady, D.~R. Andersen, and R.~A. Swaters,
``The DiskMass Survey. VI. Gas and stellar kinematics in spiral galaxies from PPak integral-field spectroscopy,''
\emph{Astronomy \& Astrophysics} \textbf{557}, A130 (2013).

\bibitem[Martinsson et al.(2013b)]{Martinsson2013b}
T.~P.~K. Martinsson, M.~A.~W. Verheijen, K.~B. Westfall, M.~A. Bershady, D.~R. Andersen, and R.~A. Swaters,
``The DiskMass Survey. VII. The distribution of luminous and dark matter in spiral galaxies,''
\emph{Astronomy \& Astrophysics} \textbf{557}, A131 (2013).

\bibitem[Angus et al.(2015)]{Angus2015}
G.~W. Angus, G. Gentile, R. Swaters, B. Famaey, A. Diaferio, S.~S. McGaugh, and K.~J. van der Heyden,
``Mass models of disc galaxies from the DiskMass Survey in modified Newtonian dynamics,''
\emph{Monthly Notices of the Royal Astronomical Society} \textbf{451}, 3551 (2015).

\bibitem[Aniyan et al.(2016)]{Aniyan2016}
S. Aniyan, K.~C. Freeman, O.~E. Gerhard, M. Arnaboldi, and C. Flynn,
``The influence of a kinematically cold young component on disc--halo decompositions in spiral galaxies: insights from solar neighbourhood K-giants,''
\emph{Monthly Notices of the Royal Astronomical Society} \textbf{456}, 1484 (2016).

\bibitem[Aniyan et al.(2021)]{Aniyan2021}
S. Aniyan, A.~A. Ponomareva, K.~C. Freeman, M. Arnaboldi, O.~E. Gerhard, et al.,
``Resolving the disc--halo degeneracy -- II: NGC 6946,''
\emph{Monthly Notices of the Royal Astronomical Society} \textbf{500}, 3579 (2021).

\bibitem[Park et al.(2026)]{Park2026}
R.~S. Park, A. Hees, B. Famaey, H. Desmond, and A. Durakovic,
``Improved constraints on modified Newtonian gravity from Cassini radio tracking data,''
\emph{Physical Review D} \textbf{114}, 024066 (2026),
doi:10.1103/r7n8-kw38.

\bibitem[J\'ulio et al.(2025)]{Julio2025}
M.~P. J\'ulio, J.~I. Read, M.~S. Pawlowski, P. Li, D. Vaz, J. Brinchmann, M.~P. Rey, O. Agertz, and T. Holmes,
``The radial acceleration relation at the EDGE of galaxy formation: testing its universality in low-mass dwarf galaxies,''
\emph{Astronomy \& Astrophysics} \textbf{704}, A330 (2025),
doi:10.1051/0004-6361/202557106.

\bibitem[Genzel et al.(2017)]{Genzel2017}
R. Genzel et al.,
``Strongly baryon-dominated disk galaxies at the peak of galaxy formation ten billion years ago,''
\emph{Nature} \textbf{543}, 397 (2017).

\bibitem[Prince and Dunkley(2019)]{PrinceDunkley2019}
H. Prince and J. Dunkley,
``Data compression in cosmology: A compressed likelihood for Planck data,''
\emph{Physical Review D} \textbf{100}, 083502 (2019).

\bibitem[Ciocan et al.(2026)]{Ciocan2026}
B.~I. Ciocan, N.~F. Bouch\'e, J. Fensch, D. Krajnovi\'c, J. Freundlich, H. Desmond, B. Famaey, and R. Techi,
``MUSE-DARK III: The evolution of the radial acceleration relation at intermediate redshifts,''
\emph{Astronomy \& Astrophysics} \textbf{709}, L16 (2026).

\bibitem[Varasteanu et al.(2025)]{Varasteanu2025}
A.~A. V\u{a}r\u{a}\c{s}teanu, M.~J. Jarvis, A.~A. Ponomareva, H. Desmond, I. Heywood, T. Yasin, N. Maddox, M. Glowacki, M. Maksymowicz-Maciata, P.~E.~M. Mancera Pi\~na, and H. Pan,
``MIGHTEE-HI: The radial acceleration relation with resolved stellar mass measurements,''
\emph{Monthly Notices of the Royal Astronomical Society} \textbf{541}, 2366 (2025).

\bibitem[Lang et al.(2017)]{Lang2017}
P. Lang et al.,
``Falling rotation curves of star-forming galaxies at $0.6<z<2.6$,''
\emph{Astrophysical Journal} \textbf{840}, 92 (2017).

\bibitem[Lelli et al.(2016)]{Lelli2016}
F. Lelli, S.~S. McGaugh, and J.~M. Schombert,
``SPARC: mass models for 175 disk galaxies with Spitzer photometry and accurate rotation curves,''
\emph{Astronomical Journal} \textbf{152}, 157 (2016).

\bibitem[Moster et al.(2013)]{Moster2013}
B.~P. Moster, T. Naab, and S.~D.~M. White,
``Galactic star formation and accretion histories from matching galaxies to dark matter haloes,''
\emph{Monthly Notices of the Royal Astronomical Society} \textbf{428}, 3121 (2013).

\bibitem[Petersen and Pe\~narrubia(2021)]{PetersenPenarrubia2021}
M.~S. Petersen and J. Pe\~narrubia,
``Detection of the Milky Way reflex motion due to the Large Magellanic Cloud infall,''
\emph{Nature Astronomy} \textbf{5}, 251--255 (2021).

\bibitem[Chandra et al.(2024)]{Chandra2024}
V. Chandra et al.,
``All-sky kinematics of the distant halo: the reflex response to the LMC,''
arXiv:2406.01676 (2024).

\bibitem[McConnachie(2012)]{McConnachie2012}
A.~W. McConnachie,
``The observed properties of dwarf galaxies in and around the Local Group,''
\emph{The Astronomical Journal} \textbf{144}, 4 (2012).

\bibitem[Caldwell et al.(2017)]{Caldwell2017}
N. Caldwell et al.,
``Crater II: an extremely cold dark matter halo,''
\emph{The Astrophysical Journal} \textbf{839}, 20 (2017).

\bibitem[Lelli et al.(2015)]{Lelli2015}
F. Lelli et al.,
``Gas dynamics in tidal dwarf galaxies: disc formation at $z=0$,''
\emph{Astronomy \& Astrophysics} \textbf{584}, A113 (2015).

\bibitem[Jordi et al.(2009)]{Jordi2009}
K. Jordi et al.,
``Testing fundamental physics with distant star clusters: analysis of observational data on Palomar 14,''
\emph{The Astronomical Journal} \textbf{137}, 4586 (2009).

\bibitem[Brouwer et al.(2021)]{Brouwer2021}
M.~M. Brouwer et al.,
``The weak lensing radial acceleration relation: Constraining modified gravity and cold dark matter theories with KiDS-1000,''
\emph{Astronomy \& Astrophysics} \textbf{650}, A113 (2021).

\bibitem[Fakhouri et al.(2010)]{Fakhouri2010}
O. Fakhouri, C.-P. Ma, and M. Boylan-Kolchin,
``The merger rates and mass assembly histories of dark matter haloes in the two Millennium simulations,''
\emph{Monthly Notices of the Royal Astronomical Society} \textbf{406}, 2267 (2010).

\bibitem[Read et al.(2019)]{Read2019}
J.~I. Read, M.~G. Walker, and P. Steger,
``Dark matter heats up in dwarf galaxies,''
\emph{Monthly Notices of the Royal Astronomical Society} \textbf{484}, 1401 (2019).

\bibitem[Mistele et al.(2024)]{Mistele2024}
T. Mistele, S.~S. McGaugh, F. Lelli, J.~M. Schombert, and P. Li,
``Radial acceleration relation of galaxies with joint kinematic and weak-lensing data,''
\emph{Journal of Cosmology and Astroparticle Physics} \textbf{2024}(04), 020 (2024),
doi:10.1088/1475-7516/2024/04/020.

\bibitem[Unruh(1976)]{Unruh1976}
W.~G. Unruh,
``Notes on black-hole evaporation,''
\emph{Physical Review D} \textbf{14}, 870 (1976).

\bibitem[Tully and Fisher(1977)]{TullyFisher1977}
R.~B. Tully and J.~R. Fisher,
``A new method of determining distances to galaxies,''
\emph{Astronomy \& Astrophysics} \textbf{54}, 661 (1977).

\bibitem[McGaugh et al.(2000)]{McGaugh2000}
S.~S. McGaugh, J.~M. Schombert, G.~D. Bothun, and W.~J.~G. de Blok,
``The baryonic Tully--Fisher relation,''
\emph{Astrophysical Journal Letters} \textbf{533}, L99 (2000).

\bibitem[Planck Collaboration(2020)]{Planck2020}
Planck Collaboration,
``Planck 2018 results. VI. Cosmological parameters,''
\emph{Astronomy \& Astrophysics} \textbf{641}, A6 (2020).

\bibitem[Riess et al.(2022)]{Riess2022}
A.~G. Riess et al.,
``A comprehensive measurement of the local value of the Hubble constant with 1 km s$^{-1}$ Mpc$^{-1}$ uncertainty from the Hubble Space Telescope and the SH0ES team,''
\emph{Astrophysical Journal Letters} \textbf{934}, L7 (2022).

\bibitem[Freedman et al.(2025)]{Freedman2025}
W.~L. Freedman, B.~F. Madore, I.~S. Jang, et al.,
``Status report on the Chicago--Carnegie Hubble Program (CCHP): measurement of the Hubble constant using the Hubble and James Webb Space Telescopes,''
\emph{The Astrophysical Journal} (2025), doi:10.3847/1538-4357/adce78.

\bibitem[Di Valentino et al.(2021)]{DiValentino2021}
E. Di Valentino et al.,
``In the realm of the Hubble tension --- a review of solutions,''
\emph{Classical and Quantum Gravity} \textbf{38}, 153001 (2021).

\bibitem[Will(2014)]{Will2014}
C.~M. Will,
``The confrontation between general relativity and experiment,''
\emph{Living Reviews in Relativity} \textbf{17}, 4 (2014).

\bibitem[Brans and Dicke(1961)]{BransDicke1961}
C. Brans and R.~H. Dicke,
``Mach's principle and a relativistic theory of gravitation,''
\emph{Physical Review} \textbf{124}, 925 (1961).

\bibitem[Adelberger et al.(2003)]{Adelberger2003}
E.~G. Adelberger, B.~R. Heckel, and A.~E. Nelson,
``Tests of the gravitational inverse-square law,''
\emph{Annual Review of Nuclear and Particle Science} \textbf{53}, 77 (2003).

\bibitem[Padmanabhan(2010)]{Padmanabhan2010}
T. Padmanabhan,
``Thermodynamical aspects of gravity: new insights,''
\emph{Reports on Progress in Physics} \textbf{73}, 046901 (2010).

\bibitem[Sanders and McGaugh(2002)]{SandersMcGaugh2002}
R.~H. Sanders and S.~S. McGaugh,
``Modified Newtonian dynamics as an alternative to dark matter,''
\emph{Annual Review of Astronomy and Astrophysics} \textbf{40}, 263 (2002).

\bibitem[Oriti(2009)]{Oriti2009}
D. Oriti,
``The group field theory approach to quantum gravity,''
in \emph{Approaches to Quantum Gravity}, edited by D. Oriti (Cambridge University Press, 2009),
arXiv:gr-qc/0607032.

\bibitem[Oriti et al.(2016)]{OritiSindoniWilsonEwing2016}
D. Oriti, L. Sindoni, and E. Wilson-Ewing,
``Emergent Friedmann dynamics with a quantum bounce from quantum gravity condensates,''
\emph{Classical and Quantum Gravity} \textbf{33}, 224001 (2016),
doi:10.1088/0264-9381/33/22/224001.

\bibitem[Oriti(2016)]{Oriti2016}
D. Oriti,
``Group field theory as the second quantization of loop quantum gravity,''
\emph{Classical and Quantum Gravity} \textbf{33}, 085005 (2016).

\bibitem[Gielen et al.(2013)]{Gielen2013}
S. Gielen, D. Oriti, and L. Sindoni,
``Cosmology from group field theory formalism for quantum gravity,''
\emph{Physical Review Letters} \textbf{111}, 031301 (2013).

\bibitem[Barbieri(1998)]{Barbieri1998}
A. Barbieri,
``Quantum tetrahedra and simplicial spin networks,''
\emph{Nuclear Physics B} \textbf{518}, 714 (1998).

\bibitem[Baez and Barrett(1999)]{BaezBarrett1999}
J.~C. Baez and J.~W. Barrett,
``The quantum tetrahedron in 3 and 4 dimensions,''
\emph{Advances in Theoretical and Mathematical Physics} \textbf{3}, 815 (1999).

\bibitem[Rovelli and Smolin(1995)]{RovelliSmolin1995}
C. Rovelli and L. Smolin,
``Discreteness of area and volume in quantum gravity,''
\emph{Nuclear Physics B} \textbf{442}, 593 (1995).

\bibitem[Gleason(1957)]{Gleason1957}
A.~M. Gleason,
``Measures on the closed subspaces of a Hilbert space,''
\emph{Journal of Mathematics and Mechanics} \textbf{6}, 885--893 (1957).

\bibitem[Hartle(2016)]{Hartle2016}
J.~B. Hartle,
``Decoherent histories quantum mechanics starting with records of what happens,''
arXiv:1608.04145 (2016).

\bibitem[Jaynes(1980)]{Jaynes1980}
E.~T. Jaynes,
``The minimum entropy production principle,''
\emph{Annual Review of Physical Chemistry} \textbf{31}, 579--601 (1980),
doi:10.1146/annurev.pc.31.100180.003051.

\bibitem[Stinespring(1955)]{Stinespring1955}
W.~F. Stinespring,
``Positive functions on $C^*$-algebras,''
\emph{Proceedings of the American Mathematical Society} \textbf{6}, 211--216 (1955),
doi:10.1090/S0002-9939-1955-0069403-4.

\bibitem[Wyner and Ziv(1989)]{WynerZiv1989}
A.~D. Wyner and J. Ziv,
``Some asymptotic properties of the entropy of a stationary ergodic data source with applications to data compression,''
\emph{IEEE Transactions on Information Theory} \textbf{35}, 1250--1258 (1989).

\bibitem[Ornstein and Weiss(1993)]{OrnsteinWeiss1993}
D.~S. Ornstein and B. Weiss,
``Entropy and data compression schemes,''
\emph{IEEE Transactions on Information Theory} \textbf{39}, 78--83 (1993).

\bibitem[Halliwell(1994)]{Halliwell1994}
J.~J. Halliwell,
``A review of the decoherent histories approach to quantum mechanics,''
arXiv:gr-qc/9407040 (1994).

\bibitem[Eckert et al.(2019)]{Eckert2019}
D. Eckert et al.,
``Non-thermal pressure support in X-COP galaxy clusters,''
\emph{Astronomy \& Astrophysics} \textbf{621}, A40 (2019).

\bibitem[Dupourqu\'e et al.(2023)]{Dupourque2023}
S. Dupourqu\'e, N. Clerc, E. Pointecouteau, D. Eckert, S. Ettori, and F. Vazza,
``Investigating the turbulent hot gas in X-COP galaxy clusters,''
\emph{Astronomy \& Astrophysics} \textbf{673}, A91 (2023).

\bibitem[XRISM Collaboration(2025)]{XRISM2025}
XRISM Collaboration,
``XRISM reveals low nonthermal pressure in the core of the hot, relaxed galaxy cluster Abell 2029,''
\emph{Astrophysical Journal Letters} \textbf{982}, L5 (2025).

\bibitem[Chamseddine \& Mukhanov(2013)]{ChamseddineMukhanov2013}
A.~H. Chamseddine and V. Mukhanov,
``Mimetic dark matter,''
\emph{Journal of High Energy Physics} \textbf{1311}, 135 (2013).

\bibitem[Chamseddine, Mukhanov, and Vikman(2014)]{ChamseddineMukhanovVikman2014}
A.~H. Chamseddine, V. Mukhanov, and A. Vikman,
``Cosmology with mimetic matter,''
\emph{Journal of Cosmology and Astroparticle Physics} \textbf{06}, 017 (2014).

\bibitem[Kac(1947)]{Kac1947}
M. Kac,
``On the notion of recurrence in discrete stochastic processes,''
\emph{Bulletin of the American Mathematical Society} \textbf{53},
1002--1010 (1947).

\bibitem[Lewis et al.(2000)]{Lewis2000}
A. Lewis, A. Challinor, and A. Lasenby,
``Efficient computation of cosmic microwave background anisotropies in closed Friedmann--Robertson--Walker models,''
\emph{Astrophysical Journal} \textbf{538}, 473 (2000).

\bibitem[More et al.(2015)]{More2015}
S. More, B. Diemer, and A. Kravtsov,
``The splashback radius as a physical halo boundary and the growth of halo mass,''
\emph{Astrophysical Journal} \textbf{810}, 36 (2015).

\bibitem[Volovik(2005)]{VolovikVacuum2005}
G.~E. Volovik,
``Cosmological constant and vacuum energy,''
\emph{Annalen der Physik} \textbf{14}, 165--176 (2005).

\bibitem[Regge(1961)]{Regge1961}
T. Regge,
``General relativity without coordinates,''
\emph{Il Nuovo Cimento} \textbf{19}, 558--571 (1961).

\bibitem[Ambj{\o}rn et al.(2004)]{AmbjornJurkiewiczLoll2004}
J. Ambj{\o}rn, J. Jurkiewicz, and R. Loll,
``Emergence of a 4D world from causal quantum gravity,''
\emph{Physical Review Letters} \textbf{93}, 131301 (2004).

\bibitem[Ambj{\o}rn et al.(2012)]{AmbjornGJL2012}
J. Ambj{\o}rn, A. G{\"o}rlich, J. Jurkiewicz, and R. Loll,
``Nonperturbative quantum gravity,''
\emph{Physics Reports} \textbf{519}, 127--210 (2012).

\bibitem[Ambj{\o}rn et al.(2001)]{AmbjornJurkiewiczLoll2001}
J. Ambj{\o}rn, J. Jurkiewicz, and R. Loll,
``Dynamically triangulating Lorentzian quantum gravity,''
\emph{Nuclear Physics B} \textbf{610}, 347--382 (2001).

\bibitem[Ambj{\o}rn et al.(2017)]{AmbjornNewPhase2017}
J. Ambj{\o}rn, J. Gizbert-Studnicki, A. G{\"o}rlich, J. Jurkiewicz, N. Klitgaard, and R. Loll,
``Characteristics of the new phase in CDT,''
\emph{European Physical Journal C} \textbf{77}, 152 (2017).

\bibitem[Ambj{\o}rn et al.(2008)]{AmbjornGJL2008dS}
J. Ambj{\o}rn, A. G{\"o}rlich, J. Jurkiewicz, and R. Loll,
``The nonperturbative quantum de Sitter universe,''
\emph{Physical Review D} \textbf{78}, 063544 (2008).

\bibitem[Thoul and Weinberg(1996)]{ThoulWeinberg1996}
A.~A. Thoul and D.~H. Weinberg,
``Hydrodynamic simulations of galaxy formation. II. Photoionization and the formation of low-mass galaxies,''
\emph{Astrophysical Journal} \textbf{465}, 608 (1996).

\bibitem[Okamoto et al.(2008)]{OkamotoGaoTheuns2008}
T. Okamoto, L. Gao, and T. Theuns,
``Mass loss of galaxies due to an ultraviolet background,''
\emph{Monthly Notices of the Royal Astronomical Society} \textbf{390}, 920--928 (2008).

\bibitem[O'Shea et al.(2015)]{OShea2015Renaissance}
B.~W. O'Shea, J.~H. Wise, H. Xu, and M.~L. Norman,
``Probing the ultraviolet luminosity function of the earliest galaxies with the Renaissance Simulations,''
\emph{Astrophysical Journal Letters} \textbf{807}, L12 (2015), arXiv:1503.01110.

\bibitem[Xu et al.(2016)]{Xu2016Renaissance}
H. Xu, J.~H. Wise, M.~L. Norman, K. Ahn, and B.~W. O'Shea,
``Galaxy properties and UV escape fractions during the epoch of reionization: results from the Renaissance Simulations,''
\emph{Astrophysical Journal} \textbf{833}, 84 (2016), arXiv:1604.07842.

\bibitem[Nadler et al.(2020)]{Nadler2020Satellites}
E.~O. Nadler et al. (DES Collaboration),
``Milky Way satellite census. II. Galaxy--halo connection constraints including the impact of the Large Magellanic Cloud,''
\emph{Astrophysical Journal} \textbf{893}, 48 (2020), arXiv:1912.03303.

\bibitem[Ambj{\o}rn et al.(2011)]{AmbjornSecondOrder2011}
J. Ambj{\o}rn, S. Jordan, J. Jurkiewicz, and R. Loll,
``A second-order phase transition in CDT,''
\emph{Physical Review Letters} \textbf{107}, 211303 (2011).

\bibitem[Ambj{\o}rn et al.(2019)]{AmbjornCritical2019}
J. Ambj{\o}rn, D. Coumbe, J. Gizbert-Studnicki, A. G{\"o}rlich, and J. Jurkiewicz,
``Critical phenomena in causal dynamical triangulations,''
\emph{Classical and Quantum Gravity} \textbf{36}, 224001 (2019).

\bibitem[Perez(2013)]{PerezSpinFoam2013}
A. Perez,
``The spin-foam approach to quantum gravity,''
\emph{Living Reviews in Relativity} \textbf{16}, 3 (2013).

\bibitem[Han(2019)]{HanSmallSpin2019}
M. Han,
``Semiclassical behavior of spinfoam amplitude with small spins and entanglement entropy,''
\emph{Physical Review D} \textbf{100}, 084049 (2019).

\bibitem[Han(2026)]{HanStack2026}
M. Han,
``Summation and triangulation independence of Lorentzian spinfoam amplitudes for all LQG,''
\emph{Physical Review D} \textbf{113}, 084034 (2026),
doi:10.1103/n76f-31gf.

\bibitem[Thiemann and Zipfel(2014)]{ThiemannZipfel2014}
T. Thiemann and A. Zipfel,
``Linking covariant and canonical LQG II: spin foam projector,''
\emph{Classical and Quantum Gravity} \textbf{31}, 125008 (2014).

\bibitem[Calcinari and Gielen(2025)]{CalcinariGielen2025}
A. Calcinari and S. Gielen,
``Relational dynamics and Page--Wootters formalism in group field theory,''
\emph{Quantum} \textbf{9}, 1610 (2025).

\bibitem[Han and Ma(2006)]{HanMa2006}
M. Han and Y. Ma,
``Master constraint operators in loop quantum gravity,''
\emph{Physics Letters B} \textbf{635}, 225--231 (2006).

\bibitem[Han and Thiemann(2010)]{HanThiemann2010}
M. Han and T. Thiemann,
``On the relation between rigging inner product and master constraint direct integral decomposition,''
arXiv:0911.3431 (2010).

\bibitem[Fairbairn and Meusburger(2012)]{FairbairnMeusburger2012}
W.~J. Fairbairn and C. Meusburger,
``Quantum deformation of two four-dimensional spin foam models,''
\emph{Journal of Mathematical Physics} \textbf{53}, 022501 (2012).

\bibitem[Han(2017)]{Han2017}
M. Han,
``Einstein equation from covariant loop quantum gravity in semiclassical continuum limit,''
\emph{Physical Review D} \textbf{96}, 024047 (2017).

\bibitem[Han, Huang, and Zipfel(2019)]{HanHuangZipfel2019}
M. Han, Z. Huang, and A. Zipfel,
``Emergent four-dimensional linearized gravity from a spin foam model,''
\emph{Physical Review D} \textbf{100}, 024060 (2019).

\bibitem[Engle(2013)]{EngleProper2013}
J. Engle,
``Proposed proper Engle--Pereira--Rovelli--Livine vertex amplitude,''
\emph{Physical Review D} \textbf{87}, 084048 (2013).

\bibitem[Engle, Vilensky, and Zipfel(2016)]{EngleVilenskyZipfel2016}
J. Engle, I. Vilensky, and A. Zipfel,
``The Lorentzian proper vertex amplitude: Asymptotics,''
\emph{Physical Review D} \textbf{94}, 064025 (2016).

\bibitem[Dittrich and Kogios(2023)]{DittrichKogios2023}
B. Dittrich and A. Kogios,
``From spin foams to area metric dynamics to gravitons,''
\emph{Classical and Quantum Gravity} \textbf{40}, 095011 (2023).

\bibitem[H\"ohn(2015)]{Hoehn2015}
P.~A. H\"ohn,
Canonical linearized Regge calculus: Counting lattice gravitons with
Pachner moves,
\emph{Physical Review D} \textbf{91}, 124034 (2015).

\bibitem[Ro\v{c}ek and Williams(1981)]{RocekWilliams1981}
M. Ro\v{c}ek and R.~M. Williams,
Quantum Regge calculus,
\emph{Physics Letters B} \textbf{104}, 31--37 (1981).

\bibitem[Dittrich, Kami\'nski, and Steinhaus(2014)]{DittrichKaminskiSteinhaus2014}
B. Dittrich, W. Kami\'nski, and S. Steinhaus,
``Discretization independence implies non-locality in 4D discrete quantum gravity,''
\emph{Classical and Quantum Gravity} \textbf{31}, 245009 (2014).

\bibitem[Lovelock(1971)]{Lovelock1971}
D. Lovelock,
``The Einstein tensor and its generalizations,''
\emph{Journal of Mathematical Physics} \textbf{12}, 498--501 (1971).

\bibitem[Deser(1970)]{Deser1970}
S. Deser,
``Self-interaction and gauge invariance,''
\emph{General Relativity and Gravitation} \textbf{1}, 9--18 (1970).

\bibitem[Borissova et al.(2025)]{BorissovaEtAl2025}
J. Borissova, B. Dittrich, A. Eichhorn, and M. Schiffer,
``Renormalization group flows in area-metric gravity,''
arXiv:2507.02034 [gr-qc] (2025).

\bibitem[Bruno et al.(2026)]{BrunoEtAl2026}
M. Bruno, E. Colafranceschi, F. M. Mele, and C. Rovelli,
``The structure of the continuum limit of spin foams,''
arXiv:2603.16999 [gr-qc] (2026).

\bibitem[Engle et al.(2008)]{EPRL2008}
J. Engle, E. Livine, R. Pereira, and C. Rovelli,
``LQG vertex with finite Immirzi parameter,''
\emph{Nuclear Physics B} \textbf{799}, 136--149 (2008).

\bibitem[Gozzini(2021)]{Gozzini2021}
F. Gozzini,
``A high-performance code for EPRL spin foam amplitudes,''
\emph{Classical and Quantum Gravity} \textbf{38}, 225010 (2021).


\bibitem[Unruh(1989)]{Unruh1989}
W.~G. Unruh,
``Unimodular theory of canonical quantum gravity,''
\emph{Physical Review D} \textbf{40}, 1048 (1989).

\bibitem[Henneaux and Teitelboim(1989)]{HenneauxTeitelboim1989}
M.~Henneaux and C.~Teitelboim,
``The cosmological constant and general covariance,''
\emph{Physics Letters B} \textbf{222}, 195 (1989).

\bibitem[Weinberg(1989)]{Weinberg1989}
S.~Weinberg,
``The cosmological constant problem,''
\emph{Reviews of Modern Physics} \textbf{61}, 1 (1989).

\bibitem[Josset et al.(2017)]{JossetPerezSudarsky2017}
T.~Josset, A.~Perez, and D.~Sudarsky,
``Dark energy from violation of energy conservation,''
\emph{Physical Review Letters} \textbf{118}, 021102 (2017).

\bibitem[Perez and Sudarsky(2019)]{PerezSudarsky2019}
A.~Perez and D.~Sudarsky,
``Dark energy from quantum gravity discreteness,''
\emph{Physical Review Letters} \textbf{122}, 221302 (2019).

\bibitem[Ahmed et al.(2004)]{Sorkin2004}
M.~Ahmed, S.~Dodelson, P.~B. Greene, and R.~D. Sorkin,
``Everpresent $\Lambda$,''
\emph{Physical Review D} \textbf{69}, 103523 (2004).

\bibitem[Das et al.(2023)]{DasNasiriYazdi2023}
S.~Das, A.~Nasiri, and Y.~K. Yazdi,
``Everpresent $\Lambda$. Part II. Structural stability,''
\emph{Journal of Cosmology and Astroparticle Physics} \textbf{10}, 047 (2023).

\bibitem[Barrow(2007)]{Barrow2007}
J.~D. Barrow,
``A strong constraint on ever-present lambda,''
\emph{Physical Review D} \textbf{75}, 067301 (2007).

\bibitem[Hsu(2004)]{Hsu2004}
S.~D.~H. Hsu,
``Entropy bounds and dark energy,''
\emph{Physics Letters B} \textbf{594}, 13 (2004).

\bibitem[Li(2004)]{Li2004}
M.~Li,
``A model of holographic dark energy,''
\emph{Physics Letters B} \textbf{603}, 1 (2004).

\bibitem[DESI Collaboration(2025)]{DESI2025}
DESI Collaboration,
``DESI DR2 results II: Measurements of baryon acoustic oscillations and cosmological constraints,''
arXiv:2503.14738 (2025).

\bibitem[Brout et al.(2022)]{Brout2022}
D.~Brout et al.,
``The Pantheon+ analysis: Cosmological constraints,''
\emph{The Astrophysical Journal} \textbf{938}, 110 (2022).

\bibitem[Press and Schechter(1974)]{PressSchechter1974}
W.~H. Press and P.~Schechter,
``Formation of galaxies and clusters of galaxies by self-similar gravitational condensation,''
\emph{The Astrophysical Journal} \textbf{187}, 425 (1974).

\bibitem[Eisenstein and Hu(1998)]{EisensteinHu1998}
D.~J. Eisenstein and W.~Hu,
``Baryonic features in the matter transfer function,''
\emph{The Astrophysical Journal} \textbf{496}, 605 (1998).

\bibitem[Ormondroyd et al.(2026)]{Ormondroyd2026}
A.~N. Ormondroyd, W.~J. Handley, M.~P. Hobson, A.~N. Lasenby, and D.~Yallup,
``Dynamic or systematic? Bayesian model selection between dark energy and supernova biases,''
\emph{Monthly Notices of the Royal Astronomical Society} \textbf{548}, stag615 (2026).

\bibitem[DES Collaboration(2026)]{DESReanalysis2026}
DES Collaboration,
``Dark Energy Survey supernova program: A reanalysis of cosmology results and evidence for evolving dark energy with an updated Type Ia supernova calibration,''
\emph{Monthly Notices of the Royal Astronomical Society} \textbf{548}, stag632 (2026).

\bibitem[Lesci et al.(2026)]{Lesci2026}
G.~F. Lesci, C.~Giocoli, F.~Marulli, M.~Romanello, L.~Moscardini, et al.,
``AMICO galaxy clusters in KiDS-1000: Splashback radius from weak lensing and cluster--galaxy correlation function,''
\emph{Astronomy \& Astrophysics} \textbf{710}, A269 (2026).

\bibitem[Hidalgo(2015)]{Hidalgo2015}
C.~Hidalgo,
\emph{Why Information Grows: The Evolution of Order, from Atoms to Economies}
(Basic Books, New York, 2015).

\bibitem[Landauer(1961)]{Landauer1961}
R.~Landauer,
``Irreversibility and heat generation in the computing process,''
\emph{IBM Journal of Research and Development} \textbf{5}, 183--191 (1961).

\bibitem[Bryan and Norman(1998)]{BryanNorman1998}
G.~L. Bryan and M.~L. Norman,
``Statistical properties of X-ray clusters: Analytic and numerical comparisons,''
\emph{The Astrophysical Journal} \textbf{495}, 80 (1998).

\bibitem[Shin et al.(2019)]{Shin2019}
T.~Shin et al.,
``Measurement of the splashback feature around SZ-selected galaxy clusters with DES, SPT and ACT,''
\emph{Monthly Notices of the Royal Astronomical Society} \textbf{487}, 2900 (2019).

\bibitem[Rana et al.(2023)]{Rana2023}
D.~Rana et al.,
``The eROSITA Final Equatorial-Depth Survey (eFEDS): Splashback radius of X-ray galaxy clusters using galaxies from HSC survey,''
\emph{Monthly Notices of the Royal Astronomical Society} \textbf{522}, 4181 (2023).

\bibitem[Rubin et al.(2023)]{Rubin2023}
D.~Rubin et al.,
``Union through UNITY: Cosmology with 2,000 SNe using a unified Bayesian framework,''
arXiv:2311.12098 (2023).

\bibitem[Greif et al.(2008)]{Greif2008FirstGalaxies}
T.~H. Greif, J.~L. Johnson, R.~S. Klessen, and V. Bromm,
``The first galaxies: assembly, cooling and the onset of turbulence,''
\emph{Monthly Notices of the Royal Astronomical Society} \textbf{387},
1021--1036 (2008).

\bibitem[Greif et al.(2010)]{Greif2010Mixing}
T.~H. Greif, S.~C.~O. Glover, V. Bromm, and R.~S. Klessen,
``The first galaxies: chemical enrichment, mixing, and star formation,''
\emph{The Astrophysical Journal} \textbf{716}, 510--520 (2010).

\bibitem[Pan and Scannapieco(2010)]{PanScannapieco2010}
L. Pan and E. Scannapieco,
``Mixing in supersonic turbulence,''
\emph{The Astrophysical Journal} \textbf{721}, 1765--1782 (2010).

\bibitem[Barkana and Loeb(1999)]{BarkanaLoeb1999}
R. Barkana and A. Loeb,
``The photoevaporation of dwarf galaxies during reionization,''
\emph{The Astrophysical Journal} \textbf{523}, 54--65 (1999).

\bibitem[Ji et al.(2015)]{JiFrebelBromm2015}
A.~P. Ji, A. Frebel, and V. Bromm,
``Preserving chemical signatures of primordial star formation in the first low-mass stars,''
\emph{Monthly Notices of the Royal Astronomical Society} \textbf{454},
659--674 (2015).

\bibitem[Griffen et al.(2018)]{GriffenEtAl2018}
B.~F. Griffen, G.~A. Dooley, A.~P. Ji, B.~W. O'Shea, F.~A. Gomez, and A. Frebel,
``Tracing the first stars and galaxies of the Milky Way,''
\emph{Monthly Notices of the Royal Astronomical Society} \textbf{474},
443--459 (2018).

\bibitem[Bombelli et al.(2009)]{AQGBombelli2009}
L.~Bombelli, J.~Henson, and R.~D.~Sorkin,
``Discreteness without symmetry breaking: a theorem,''
\emph{Mod. Phys. Lett. A} \textbf{24}, 2579--2587 (2009),
\href{https://arxiv.org/abs/gr-qc/0605006}{arXiv:gr-qc/0605006}.

\bibitem[Takayanagi and Tamaoka(2020)]{AQGTakayanagi2020}
T.~Takayanagi and K.~Tamaoka,
``Gravity Edges Modes and Hayward Term,''
\emph{JHEP} \textbf{02}, 167 (2020),
\href{https://arxiv.org/abs/1912.01636}{arXiv:1912.01636}.

\bibitem[Botta-Cantcheff et al.(2020)]{AQGBotta2020}
M.~Botta-Cantcheff, P.~J.~Martinez, and J.~F.~Zarate,
``R\'enyi entropies and area operator from gravity with Hayward term,''
\emph{JHEP} \textbf{07}, 227 (2020),
\href{https://arxiv.org/abs/2005.11338}{arXiv:2005.11338}.

\bibitem[Dittrich and H\"ohn(2010)]{AQGDittrich2010}
B.~Dittrich and P.~A.~H\"ohn,
``From covariant to canonical formulations of discrete gravity,''
\emph{Class. Quantum Grav.} \textbf{27}, 155001 (2010),
\href{https://arxiv.org/abs/0912.1817}{arXiv:0912.1817}.

\bibitem[Obukhov et al.(2011)]{AQGObukhov2011}
Y.~N.~Obukhov, A.~J.~Silenko, and O.~V.~Teryaev,
``Dirac fermions in strong gravitational fields,''
\emph{Phys. Rev. D} \textbf{84}, 024025 (2011),
\href{https://arxiv.org/abs/1106.0173}{arXiv:1106.0173}.

\bibitem[Bianchi et al.(2011)]{AQGBianchi2011}
E.~Bianchi, P.~Don\`a, and S.~Speziale,
\lq\lq Polyhedra in loop quantum gravity,\rq\rq\
\emph{Phys. Rev. D} \textbf{83}, 044035 (2011),
\href{https://arxiv.org/abs/1009.3402}{arXiv:1009.3402}.

\bibitem[Louko and Martinez-Pascual(2011)]{AQGLouko2011}
J.~Louko and E.~Martinez-Pascual,
\lq\lq Constraint rescaling in refined algebraic quantisation:
momentum constraint,\rq\rq\
\emph{J. Math. Phys.} \textbf{52}, 123504 (2011),
\href{https://arxiv.org/abs/1107.1092}{arXiv:1107.1092}.


\bibitem[Christ, Friedberg, and Lee(1982)]{ChristFriedbergLee1982}
N.~H. Christ, R. Friedberg, and T.~D. Lee,
``Random lattice field theory: general formulation,''
\emph{Nuclear Physics B} \textbf{202}, 89--125 (1982).

\bibitem[Kosteleck\'y and Russell(2011)]{KosteleckyRussell2011}
V.~A. Kosteleck\'y and N. Russell,
``Data tables for Lorentz and CPT violation,''
\emph{Reviews of Modern Physics} \textbf{83}, 11--31 (2011).

\bibitem[Maccione et al.(2007)]{Maccione2007}
L. Maccione, S. Liberati, A. Celotti, and J.~G. Kirk,
``New constraints on Planck-scale Lorentz violation in QED from the Crab Nebula,''
\emph{Journal of Cosmology and Astroparticle Physics} \textbf{2007}(10), 013 (2007).

\bibitem[Becher and Joos(1982)]{BecherJoos1982}
P. Becher and H. Joos,
``The Dirac--K\"ahler equation and fermions on the lattice,''
\emph{Zeitschrift f\"ur Physik C} \textbf{15}, 343--365 (1982).

\bibitem[Rabin(1982)]{Rabin1982}
J.~M. Rabin,
``Homology theory of lattice fermion doubling,''
\emph{Nuclear Physics B} \textbf{201}, 315--332 (1982).

\bibitem[Nielsen and Ninomiya(1981)]{NielsenNinomiya1981}
H.~B. Nielsen and M. Ninomiya,
``Absence of neutrinos on a lattice (I). Proof by homotopy theory,''
\emph{Nuclear Physics B} \textbf{185}, 20--40 (1981).

\bibitem[Fu, Kane, and Mele(2007)]{FuKaneMele2007}
L. Fu, C.~L. Kane, and E.~J. Mele,
``Topological insulators in three dimensions,''
\emph{Physical Review Letters} \textbf{98}, 106803 (2007).

\bibitem[L\"uscher(1999)]{Luscher1999}
M. L\"uscher,
``Abelian chiral gauge theories on the lattice with exact gauge invariance,''
\emph{Nuclear Physics B} \textbf{549}, 295--334 (1999).

\bibitem[Wang and You(2022)]{WangYou2022}
J. Wang and Y.-Z. You,
``Symmetric mass generation,''
\emph{Symmetry} \textbf{14}, 1475 (2022).

\bibitem[Springel et al.(2021)]{Springel2021}
V. Springel, R. Pakmor, O. Zier, and M. Reinecke,
``Simulating cosmic structure formation with the GADGET-4 code,''
\emph{Monthly Notices of the Royal Astronomical Society} \textbf{506}, 2871--2949 (2021).

\bibitem[Springel and Hernquist(2003)]{SpringelHernquist2003}
V. Springel and L. Hernquist,
``Cosmological smoothed particle hydrodynamics simulations: a hybrid multiphase model for star formation,''
\emph{Monthly Notices of the Royal Astronomical Society} \textbf{339}, 289--311 (2003).

\bibitem[Katz, Weinberg, and Hernquist(1996)]{KatzWeinbergHernquist1996}
N. Katz, D.~H. Weinberg, and L. Hernquist,
``Cosmological simulations with TreeSPH,''
\emph{The Astrophysical Journal Supplement Series} \textbf{105}, 19--35 (1996).

\bibitem[Aniyan et al.(2018)]{Aniyan2018}
S. Aniyan, K.~C. Freeman, M. Arnaboldi, O.~E. Gerhard, L. Coccato, et al.,
``Resolving the disc--halo degeneracy -- I: a look at NGC 628,''
\emph{Monthly Notices of the Royal Astronomical Society} \textbf{476}, 1909--1930 (2018).

\bibitem[Bovy et al.(2016)]{Bovy2016}
J. Bovy, A. Bahmanyar, T.~K. Fritz, and N. Kallivayalil,
``The shape of the inner Milky Way halo from observations of the Pal 5 and GD-1 stellar streams,''
\emph{The Astrophysical Journal} \textbf{833}, 31 (2016).

\bibitem[Malhan and Ibata(2019)]{MalhanIbata2019}
K. Malhan and R.~A. Ibata,
``Constraining the Milky Way halo potential with the GD-1 stellar stream,''
\emph{Monthly Notices of the Royal Astronomical Society} \textbf{486}, 2995--3005 (2019).

\bibitem[Ibata et al.(2024)]{Ibata2024}
R. Ibata, K. Malhan, W. Tenachi, et al.,
``Charting the Galactic acceleration field. II. A global mass model of the Milky Way from the STREAMFINDER atlas of stellar streams detected in Gaia DR3,''
\emph{The Astrophysical Journal} \textbf{967}, 89 (2024).

\bibitem[Woudenberg and Helmi(2024)]{Woudenberg2024}
H.~C. Woudenberg and A. Helmi,
``First measurement of the triaxiality of the inner dark matter halo of the Milky Way,''
\emph{Astronomy \& Astrophysics} \textbf{691}, A277 (2024).

\bibitem[Nibauer and Bonaca(2025)]{NibauerBonaca2025}
J. Nibauer and A. Bonaca,
``Galactic accelerations from the GD-1 stream suggest a tilted dark matter halo,''
\emph{The Astrophysical Journal Letters} (2025), doi:10.3847/2041-8213/add0a9; arXiv:2504.07187.

\bibitem[Dillamore and Sanders(2026)]{DillamoreSanders2025}
A.~M. Dillamore and J.~L. Sanders,
``Geometry of the Milky Way's dark matter from dynamical models of the tilted stellar halo,''
\emph{Monthly Notices of the Royal Astronomical Society} \textbf{546}, stag226 (2026); arXiv:2510.00095.

\bibitem[Zhu et al.(2025)]{Zhu2025}
L. Zhu, R. Cai, X. Kang, X.-X. Xue, C. Yang, L. Zhang, S. Mao, and C. Liu,
``A vertically orientated dark matter halo marks a flip of the Galactic disc,''
arXiv:2510.08684 (2025).

\bibitem[Bousso(2000)]{Bousso2000}
R. Bousso,
``Positive vacuum energy and the N-bound,''
\emph{Journal of High Energy Physics} \textbf{0011}, 038 (2000).

\bibitem[Banks(2000)]{Banks2000}
T. Banks,
``Cosmological breaking of supersymmetry or little Lambda goes back to the future II,''
\emph{International Journal of Modern Physics A} \textbf{16}, 910 (2001); arXiv:hep-th/0007146.

\bibitem[Dyson, Kleban, and Susskind(2002)]{DysonKlebanSusskind2002}
L. Dyson, M. Kleban, and L. Susskind,
``Disturbing implications of a cosmological constant,''
\emph{Journal of High Energy Physics} \textbf{0210}, 011 (2002).

\bibitem[Barrett et al.(2009)]{BarrettEtAl2009}
J.~W. Barrett, R.~J. Dowdall, W.~J. Fairbairn, H. Gomes, and F. Hellmann,
``Asymptotic analysis of the EPRL four-simplex amplitude,''
\emph{Journal of Mathematical Physics} \textbf{50}, 112504 (2009).

\bibitem[Kovachy et al.(2015)]{Kovachy2015}
T. Kovachy, P. Asenbaum, C. Overstreet, C.~A. Donnelly, S.~M. Dickerson, A. Sugarbaker, J.~M. Hogan, and M.~A. Kasevich,
``Quantum superposition at the half-metre scale,''
\emph{Nature} \textbf{528}, 530--533 (2015).

\bibitem[Fein et al.(2019)]{Fein2019}
Y.~Y. Fein, P. Geyer, P. Zwick, F. Kia{\l}ka, S. Pedalino, M. Mayor, S. Gerlich, and M. Arndt,
``Quantum superposition of molecules beyond 25 kDa,''
\emph{Nature Physics} \textbf{15}, 1242--1245 (2019).

\bibitem[OpenAI(2026a)]{OpenAIHeisenberg2026}
OpenAI,
``Spontaneous magnetization in the quantum Heisenberg ferromagnet,''
preprint (September 24, 2026), OpenAI mathematics manuscript collection, result~271,
\url{https://github.com/openai/math}; main theorem formalized in Lean.

\bibitem[OpenAI(2026b)]{OpenAIBloch2026}
OpenAI,
``Bloch's law for finite-range Heisenberg ferromagnets in three dimensions,''
preprint (October 5, 2026), OpenAI mathematics manuscript collection, result~271,
\url{https://github.com/openai/math}.

\bibitem[OpenAI(2026c)]{OpenAIAreaLaw2026}
OpenAI,
``A two-dimensional area law from a global spectral gap,''
preprint (September 24, 2026), OpenAI mathematics manuscript collection, result~265,
\url{https://github.com/openai/math}.

\bibitem[OpenAI(2026d)]{OpenAIPEPS2026}
OpenAI,
``Polynomial PEPS approximation of gapped square-grid ground states,''
preprint (September 24, 2026), OpenAI mathematics manuscript collection, result~265,
\url{https://github.com/openai/math}.

\bibitem[OpenAI(2026e)]{OpenAIFKSpectral2026}
OpenAI,
``Spectral convergence for critical FK--Ising planar maps,''
preprint (October 5, 2026), OpenAI mathematics manuscript collection, result~211,
\url{https://github.com/openai/math}.

\bibitem[OpenAI(2026f)]{OpenAIO3Continuum2026}
OpenAI,
``The canonical massive continuum limit of the two-dimensional $O(3)$ model,''
preprint (October 4, 2026), OpenAI mathematics manuscript collection, result~215,
\url{https://github.com/openai/math}.

\bibitem[OpenAI(2026g)]{OpenAIKadison2026}
OpenAI,
``Kadison's similarity theorem through uniform derivation estimates,''
preprint (September 23, 2026), OpenAI mathematics manuscript collection, result~288,
\url{https://github.com/openai/math}; main theorem and hyperreflexivity constant formalized in Lean.

\bibitem[OpenAI(2026h)]{OpenAIThreeState2026}
OpenAI,
``The exact reconstruction threshold for the three-state symmetric channel,''
preprint (September 25, 2026), OpenAI mathematics manuscript collection, result~229,
\url{https://github.com/openai/math}.

\bibitem[OpenAI(2026i)]{OpenAIPottsReconstruction2026}
OpenAI,
``The reconstruction threshold for the ferromagnetic four-state Potts model,''
preprint (October 5, 2026), OpenAI mathematics manuscript collection, result~229,
\url{https://github.com/openai/math}.

\bibitem[OpenAI(2026j)]{OpenAIScaleConformal2026}
OpenAI,
``Scale and conformal symmetry in four-dimensional operational quantum field theory,''
preprint (September 26, 2026), OpenAI mathematics manuscript collection, result~282,
\url{https://github.com/openai/math}.

\bibitem{Alpha7Checkpoint2026}
J. Chinitz, connected $\alpha=7$ CDT registration and numerical evidence,
17--20 October 2026. Original registration and pre-data amendment,
stage-2 verdict with A5, separate registered trace verdict, and development
readouts; preserved evidence packet in the companion CDT project,
\nolinkurl{evidence/20261020_alpha7_checkpoint/INPUT_MANIFEST.json}.


\bibitem[Clerk et al.(2010)]{Clerk2010}
A.~A. Clerk, M.~H. Devoret, S.~M. Girvin, F. Marquardt, and R.~J. Schoelkopf,
``Introduction to quantum noise, measurement, and amplification,''
\emph{Reviews of Modern Physics} \textbf{82}, 1155--1208 (2010),
doi:10.1103/RevModPhys.82.1155.

\bibitem[Martinelli et al.(2019)]{Martinelli2019}
M. Martinelli, N.~B. Hogg, S. Peirone, M. Bruni, and D. Wands,
``Constraints on the interacting vacuum--geodesic CDM scenario,''
\emph{Monthly Notices of the Royal Astronomical Society} \textbf{488},
3423--3438 (2019), doi:10.1093/mnras/stz1915.


\bibitem{Ettori2019Cluster} S. Ettori et al., ``Hydrostatic mass profiles in X-COP galaxy clusters,'' A\&A 621, A39 (2019), \href{https://arxiv.org/abs/1805.00035}{arXiv:1805.00035}.
\bibitem{Herbonnet2020Cluster} R. Herbonnet et al., ``CCCP and MENeaCS: (updated) weak-lensing masses for 100 galaxy clusters,'' MNRAS 497, 4684 (2020), \href{https://arxiv.org/abs/1912.04414}{arXiv:1912.04414}.
\bibitem{Rihtarsic2026Cluster} G. Rihtar\v{s}i\v{c} et al., ``Mapping dark matter in the Bullet Cluster using JWST imaging and spectroscopy'' (2026), \href{https://arxiv.org/html/2601.22245v2}{arXiv:2601.22245v2}; \href{https://niriss.github.io/lensing.html}{CANUCS Bullet v1 maps}.
\bibitem{Sarkar2026Cluster} A. Sarkar et al., ``XRISM Observations of Abell 1795: Evidence for Low Turbulence and Resonant Scattering'' (2026), \href{https://arxiv.org/abs/2606.08097}{arXiv:2606.08097}.
\bibitem{XRISM2025Cluster} XRISM Collaboration, ``Constraining gas motion and non-thermal pressure beyond the core of the Abell 2029 galaxy cluster with XRISM'' (2025), \href{https://arxiv.org/abs/2505.06533}{arXiv:2505.06533}.

\bibitem[Kogut and Susskind(1975)]{KogutSusskind1975EM}
J.~Kogut and L.~Susskind,
``Hamiltonian formulation of Wilson's lattice gauge theories,''
\emph{Physical Review D} \textbf{11}, 395--408 (1975),
\href{https://doi.org/10.1103/PhysRevD.11.395}{doi:10.1103/PhysRevD.11.395}.

\bibitem[Chandrasekharan and Wiese(1997)]{ChandrasekharanWiese1997EM}
S.~Chandrasekharan and U.-J.~Wiese,
``Quantum link models: A discrete approach to gauge theories,''
\emph{Nuclear Physics B} \textbf{492}, 455--474 (1997),
\href{https://arxiv.org/abs/hep-lat/9609042}{arXiv:hep-lat/9609042}.


\bibitem[Nagata et al.(2025)]{PDGGUT2025EM}
N.~Nagata, A.~Hebecker, and J.~Hisano,
``Grand Unified Theories,'' in S.~Navas et al. (Particle Data Group),
\emph{Physical Review D} \textbf{110}, 030001 (2024), 2025 update,
Sec.~92.3 and Eq.~(92.4).
\url{https://pdg.lbl.gov/2025/reviews/rpp2025-rev-guts.pdf}.

\bibitem[Huston et al.(2025)]{PDGQCD2025EM}
J.~Huston, K.~Rabbertz, and G.~Zanderighi,
``Quantum Chromodynamics,'' in the Particle Data Group 2025 update,
Eqs.~(9.3)--(9.4), (9.25).
\url{https://pdg.lbl.gov/2025/reviews/rpp2025-rev-qcd.pdf}.

\bibitem[Erler and Freitas(2024)]{PDGEW2024EM}
J.~Erler and A.~Freitas,
``Electroweak Model and Constraints on New Physics,'' in the Particle
Data Group 2024 review, Sec.~10.2.2.
\url{https://pdg.lbl.gov/2024/reviews/rpp2024-rev-standard-model.pdf}.

\bibitem[Mayer et al.(2022)]{Mayer2022}
A.~C. Mayer, A.~F. Teklu, K. Dolag and R.-S. Remus, ``$\Lambda$CDM with baryons vs.\ MOND: the time evolution of the universal acceleration scale in the Magneticum simulations,'' arXiv:2206.04333 (2022).

\bibitem[Lelli et al.(2017)]{Lelli2017RAR}
F. Lelli, S.~S. McGaugh, J.~M. Schombert and M.~S. Pawlowski, ``One law to rule them all: the radial acceleration relation of galaxies,'' Astrophys.\ J.\ 836, 152 (2017), arXiv:1610.08981.

\bibitem[Humphrey et al.(2011)]{Humphrey2011N720}
P.~J. Humphrey, D.~A. Buote, C.~R. Canizares, A.~C. Fabian and J.~M. Miller, ``A census of baryons and dark matter in an isolated, Milky Way sized elliptical galaxy,'' Astrophys.\ J.\ 729, 53 (2011), arXiv:1010.6078.

\bibitem[Ludlow et al.(2017)]{Ludlow2017}
A.~D. Ludlow et al., ``Mass-discrepancy acceleration relation: a natural outcome of galaxy formation in cold dark matter halos,'' Phys.\ Rev.\ Lett.\ 118, 161103 (2017), arXiv:1610.07663.

\bibitem[Pelissetto and Vicari(2019)]{PelissettoVicariCPN2019}
A. Pelissetto and E. Vicari,
``Three-dimensional ferromagnetic CP(N-1) models,''
arXiv:1905.03307 (2019).

\bibitem[Pelissetto and Vicari(2020)]{PelissettoVicariLargeN2020}
A. Pelissetto and E. Vicari,
``Large-N behavior of three-dimensional lattice CP(N-1) models,''
arXiv:1912.04597 (2020).

\bibitem[Panero(2005)]{Panero2005U1}
M. Panero,
``A numerical study of confinement in compact QED,''
\emph{Journal of High Energy Physics} \textbf{05}, 066 (2005), arXiv:hep-lat/0503024.

\bibitem[Neuberger(1998)]{Neuberger1998}
H. Neuberger,
``Exactly massless quarks on the lattice,''
\emph{Physics Letters B} \textbf{417}, 141--144 (1998), arXiv:hep-lat/9707022.

\bibitem[L\"uscher(2000)]{Luscher2000}
M. L\"uscher,
``Lattice regularization of chiral gauge theories to all orders of perturbation theory,''
\emph{Journal of High Energy Physics} \textbf{06}, 028 (2000), arXiv:hep-lat/0006014.

\bibitem[Hermele et al.(2004)]{HermeleFisherBalents2004}
M. Hermele, M.~P.~A. Fisher and L. Balents,
``Pyrochlore photons: the U(1) spin liquid in a S=1/2 three-dimensional frustrated magnet,''
\emph{Physical Review B} \textbf{69}, 064404 (2004), arXiv:cond-mat/0305401.

\bibitem[Gastaldello et al.(2007)]{Gastaldello2007}
F. Gastaldello, D.~A. Buote, P.~J. Humphrey, L. Zappacosta, J.~S. Bullock, F. Brighenti and W.~G. Mathews,
``Probing the dark matter and gas fraction in relaxed galaxy groups with X-ray observations from Chandra and XMM-Newton,''
\emph{Astrophysical Journal} \textbf{669}, 158 (2007), arXiv:astro-ph/0610134.

\bibitem[Mistele et al.(2024b)]{MisteleFlat2024}
T. Mistele, S.~S. McGaugh, F. Lelli, J.~M. Schombert and P. Li,
``Indefinitely flat circular velocities and the baryonic Tully--Fisher relation from weak lensing,''
arXiv:2406.09685 (2024).

\bibitem[Baez and Huerta(2010)]{BaezHuerta2010}
J. Baez and J. Huerta,
``The algebra of grand unified theories,''
\emph{Bulletin of the American Mathematical Society} \textbf{47}, 483 (2010).

\bibitem[St{\o}rmer(1963)]{Stormer1963}
E. St{\o}rmer,
``Positive linear maps of operator algebras,''
\emph{Acta Mathematica} \textbf{110}, 233 (1963).

\bibitem[Woronowicz(1976)]{Woronowicz1976}
S.~L. Woronowicz,
``Positive maps of low dimensional matrix algebras,''
\emph{Reports on Mathematical Physics} \textbf{10}, 165 (1976).

\bibitem[Choi(1975)]{Choi1975}
M.-D. Choi,
``Positive semidefinite biquadratic forms,''
\emph{Linear Algebra and its Applications} \textbf{12}, 95 (1975).

\bibitem[Wissner-Gross(2026)]{WissnerGross2026}
A.~D. Wissner-Gross,
``A parameter-free candidate for the fine-structure constant from entangled qudits,''
preprint (23 October 2026), \url{https://alexwg.org/publications/ChainTheoryAlpha.pdf}.
\bibitem[Murphy et al.(2022)]{Murphy2022ESPRESSO}
M.~T. Murphy, P. Molaro, A.~C.~O. Leite, G. Cupani, S. Cristiani et al.,
``Fundamental physics with ESPRESSO: precise limit on variations in the fine-structure constant towards the bright quasar HE 0515$-$4414,''
\emph{Astronomy \& Astrophysics} \textbf{658}, A123 (2022).
\end{thebibliography}

\end{document}
